\input{preamble}

% OK, start here.
%
\begin{document}

\title{Cohomologie \'etale}


\maketitle

\phantomsection
\label{section-phantom}

\tableofcontents


\section{Introduction}
\label{section-introduction}

\noindent
Ce chapitre est le premier d'une s\'erie de chapitres sur la cohomologie
\'etale des sch\'emas. Nous y pr\'esentons les rudiments de la topologie
\'etale et de la cohomologie des faisceaux ab\'eliens pour cette topologie. Une
grande partie des sujets abord\'es peut sans inconv\'enient \^etre omise lors
d'une premi\`ere lecture; on consultera la section suivante pour d\'ecider des
passages \`a laisser de c\^ot\'e.

\medskip\noindent
La version initiale de ce chapitre a \`a l'origine \'et\'e constitu\'ee par les notes
de la premi\`ere partie d'un cours de cohomologie \'etale donn\'e par Johan de Jong
\`a l'Universit\'e Columbia \`a l'automne 2009. Les premiers auteurs des notes
\'etaient Thibaut Pugin, Zachary Maddock et Min Lee. La seconde partie du
cours se trouve dans le chapitre sur la formule des traces; voir
La formule des traces, section \ref{trace-section-introduction}.




\section{Quelles sections omettre lors d'une premi\`ere lecture?}
\label{section-skip}

\noindent
Nous voulons utiliser le contenu de ce chapitre pour d\'evelopper la th\'eorie
des espaces alg\'ebriques, des champs de Deligne--Mumford, des champs
alg\'ebriques, etc. Nous avons donc ajout\'e \`a l'expos\'e initial de la
cohomologie \'etale des sch\'emas certains d\'eveloppements assez techniques.
Le lecteur les reconna\^itra \`a la fr\'equence du mot ``topos'', aux discussions
relevant de la th\'eorie des ensembles ou encore aux d\'emonstrations portant
sur des propri\'et\'es tr\`es g\'en\'erales des morphismes de sch\'emas. Certaines de
ces discussions peuvent \^etre omises lors d'une premi\`ere lecture.

\medskip\noindent
En particulier, nous sugg\'erons au lecteur d'omettre les sections suivantes:
\begin{enumerate}
\item Comparaison des grands et petits topos,
section \ref{section-compare}.
\item Reconstruction des morphismes,
section \ref{section-morphisms}.
\item Image directe et image inverse,
section \ref{section-monomorphisms}.
\item Propri\'et\'e (A),
section \ref{section-A}.
\item Propri\'et\'e (B),
section \ref{section-B}.
\item Propri\'et\'e (C),
section \ref{section-C}.
\item Invariance topologique du petit site \'etale,
section \ref{section-topological-invariance}.
\item Morphismes entiers universellement injectifs,
section \ref{section-integral-universally-injective}.
\item Grands sites et image directe,
section \ref{section-big}.
\item Exactitude du grand foncteur image directe \`a support propre,
section \ref{section-exactness-lower-shriek}.
\end{enumerate}
Outre ces sections, quelques r\'esultats isol\'es peuvent \^etre omis; le lecteur
les reconna\^itra aux mots-cl\'es indiqu\'es ci-dessus.



%9.08.09
\section{Prologue}
\label{section-prologue}

\noindent
Ces le\c{c}ons portent sur une autre th\'eorie cohomologique. Il faut tout
d'abord observer que la topologie de Zariski n'est pas enti\`erement
satisfaisante. L'une des principales raisons pour lesquelles elle ne donne pas
les r\'esultats souhait\'es est que, si $X$ est une vari\'et\'e complexe et
$\mathcal{F}$ un faisceau constant, alors
$$
H^i(X, \mathcal{F}) = 0, \quad \text{ pour tout } i > 0.
$$
En voici la raison. Dans un sch\'ema irr\'eductible (en particulier dans une
vari\'et\'e), deux ouverts non vides quelconques se rencontrent, de sorte que les
applications de restriction d'un faisceau constant sont surjectives. On dit
que le faisceau est {\it flasque}. Dans ce cas, tous les groupes de cohomologie
de {\v C}ech sup\'erieurs sont nuls, de m\^eme que tous les groupes de
cohomologie de Zariski sup\'erieurs. Autrement dit, la topologie de Zariski ne
poss\`ede ``pas assez'' d'ouverts pour d\'etecter cette cohomologie sup\'erieure.

\medskip\noindent
En revanche, si $X$ est une vari\'et\'e complexe projective lisse, alors
$$
H_{Betti}^{2 \dim X}(X (\mathbf{C}), \Lambda) = \Lambda \quad \text{ pour }
\Lambda = \mathbf{Z}, \ \mathbf{Z}/n\mathbf{Z},
$$
o\`u $X(\mathbf{C})$ d\'esigne l'ensemble des points complexes de $X$. Il serait
souhaitable de retrouver cette propri\'et\'e en g\'eom\'etrie alg\'ebrique, tout
particuli\`erement en caract\'eristique positive.




\section{La topologie \'etale}
\label{section-etale-topology}

\noindent
Il est tr\`es difficile de simplement ``ajouter'' des ouverts suppl\'ementaires
pour raffiner la topologie de Zariski. Une mani\`ere efficace de d\'efinir une
topologie consiste \`a consid\'erer non seulement des ouverts, mais aussi certains
sch\'emas au-dessus de ceux-ci. Pour d\'efinir la topologie \'etale, on consid\`ere
tous les morphismes \'etales $\varphi : U \to X$. Si $X$ est une vari\'et\'e
projective lisse sur $\mathbf{C}$, cela signifie que
\begin{enumerate}
\item $U$ est une r\'eunion disjointe de vari\'et\'es lisses, et
\item $\varphi$ est localement un isomorphisme (analytiquement).
\end{enumerate}
Le mot ``analytiquement'' renvoie \`a la topologie usuelle (transcendante) sur
$\mathbf{C}$. La seconde condition signifie donc que la diff\'erentielle de
$\varphi$ est partout de rang maximal (et, en particulier, que toutes les
composantes de $U$ ont la m\^eme dimension que $X$).

\medskip\noindent
Un rev\^etement double --- d\'efini de mani\`ere informelle comme une application
finie de degr\'e $2$ entre vari\'et\'es --- par exemple
$$
\Spec(\mathbf{C}[t])
\longrightarrow
\Spec(\mathbf{C}[t]),
\quad t \longmapsto t^2
$$
n'est pas un morphisme \'etale s'il poss\`ede une fibre r\'eduite \`a un seul point.
Dans cet exemple, cela se produit pour $t = 0$. Pour qu'une application finie
entre vari\'et\'es sur $\mathbf{C}$ soit \'etale, toutes ses fibres doivent avoir le
m\^eme nombre de points. En retirant le point $t = 0$ de la source de
l'application dans l'exemple, on rend le morphisme \'etale. Mais on peut aussi
retirer d'autres points de la source et le morphisme demeure \'etale. Pour
consid\'erer la topologie \'etale, il faut examiner tous ces morphismes. \`A la
diff\'erence de la topologie de Zariski, ils ne sont pas n\'ecessairement de
simples ouverts de $X$, bien que leurs images le soient toujours.

\begin{definition}
\label{definition-etale-covering-initial}
Une famille de morphismes $\{ \varphi_i : U_i \to X\}_{i \in I}$ est appel\'ee
un {\it recouvrement \'etale} si chaque $\varphi_i$ est un morphisme \'etale et
si leurs images recouvrent $X$, c'est-\`a-dire si
$X = \bigcup_{i \in I} \varphi_i(U_i)$.
\end{definition}

\noindent
Ceci ``d\'efinit'' la topologie \'etale. Autrement dit, nous pouvons maintenant
pr\'eciser ce que sont les faisceaux. Un {\it faisceau \'etale} $\mathcal{F}$
d'ensembles (resp.\ de groupes ab\'eliens, d'espaces vectoriels, etc.) sur $X$
est la donn\'ee de:
\begin{enumerate}
\item pour chaque morphisme \'etale $\varphi : U \to X$, un ensemble
(resp.\ un groupe ab\'elien, un espace vectoriel, etc.) $\mathcal{F}(U)$;
\item pour chaque paire $U, \ U'$ de sch\'emas \'etales sur $X$, et chaque
morphisme $U \to U'$ au-dessus de $X$ (qui est automatiquement \'etale), une
application de restriction
$\rho^{U'}_U : \mathcal{F}(U') \to \mathcal{F}(U)$
\end{enumerate}
Ces donn\'ees doivent satisfaire les conditions $\rho^U_U = \text{id}$ pour le
morphisme identique $U \to U$ et
$\rho^{U'}_U \circ \rho^{U''}_{U'} = \rho^{U''}_U$ pour des morphismes
$U \to U' \to U''$ de sch\'emas \'etales sur $X$, ainsi que l'{\it axiome de
faisceau} suivant:
\begin{itemize}
\item[(*)] pour tout recouvrement \'etale $\{ \varphi_i : U_i \to U\}_{i \in
I}$, le diagramme
$$
\xymatrix{
\mathcal{F} (U) \ar[r] &
\Pi_{i \in I} \mathcal{F} (U_i) \ar@<1ex>[r] \ar@<-1ex>[r] &
\Pi_{i, j \in I} \mathcal{F} (U_i \times_U U_j)
}
$$
est un diagramme d'\'egalisateur dans la cat\'egorie des ensembles
(resp.\ des groupes ab\'eliens, des espaces vectoriels, etc.).
\end{itemize}

\begin{remark}
\label{remark-i-is-j}
Dans le dernier \'enonc\'e, il est essentiel de ne pas oublier le cas $i = j$,
qui constitue en g\'en\'eral une condition loin d'\^etre triviale (contrairement
\`a ce qui se passe pour la topologie de Zariski). De fait, des recouvrements
importants ne comportent souvent qu'un seul \`el\'ement.
\end{remark}

\noindent
Puisque l'identit\'e est un morphisme \'etale, nous pouvons calculer les sections
globales d'un faisceau \'etale, et la cohomologie sera simplement donn\'ee par les
foncteurs d\'eriv\'es \`a droite correspondants. Autrement dit, une fois la th\'eorie
davantage d\'evelopp\'ee et les \'enonc\'es rendus pr\'ecis, rien ne s'opposera \`a la
d\'efinition de la cohomologie.




\section{Prouesses de la topologie \'etale}
\label{section-feats}

\noindent
Pour tout entier naturel $n \in \mathbf{N} = \{1, 2, 3, 4, \ldots\}$, on a
$$
H_\etale^2 (\mathbf{P}^1_\mathbf{C}, \mathbf{Z}/n\mathbf{Z}) =
\mathbf{Z}/n\mathbf{Z}.
$$
Plus g\'en\'eralement, si $X$ est une vari\'et\'e complexe, ses nombres de Betti
\'etales \`a coefficients dans un corps fini co\"incident avec les nombres de
Betti usuels de $X(\mathbf{C})$, c'est-\`a-dire que
$$
\dim_{\mathbf{F}_q} H_\etale^{2i} (X, \mathbf{F}_q) =
\dim_{\mathbf{F}_q} H_{Betti}^{2i} (X(\mathbf{C}), \mathbf{F}_q).
$$
C'est tout \`a fait satisfaisant. Toutefois, ces \'egalit\'es ne valent que pour des
coefficients de torsion, et non en g\'en\'eral. \`A coefficients entiers, on a
$$
H_\etale^2 (\mathbf{P}^1_\mathbf{C}, \mathbf{Z}) = 0.
$$
En revanche,
$H_{Betti}^2(\mathbf{P}^1(\mathbf{C}), \mathbf{Z}) = \mathbf{Z}$, car l'espace
topologique $\mathbf{P}^1(\mathbf{C})$ est hom\'eomorphe \`a une sph\`ere de
dimension $2$. Il existe des moyens de retrouver des coefficients non de
torsion \`a partir de coefficients de torsion par un proc\'ed\'e de passage \`a la
limite, que nous aborderons bient\^ot.




\section{Un calcul}
\label{section-computation}

\noindent
Comment calculer la cohomologie de $\mathbf{P}^1_\mathbf{C}$ \`a coefficients
dans $\Lambda = \mathbf{Z}/n\mathbf{Z}$?
Nous utilisons la cohomologie de {\v C}ech. Un recouvrement de
$\mathbf{P}^1_\mathbf{C}$ est donn\'e par les deux ouverts standards $U_0, U_1$,
tous deux isomorphes \`a $\mathbf{A}^1_\mathbf{C}$, et dont l'intersection est
isomorphe \`a
$\mathbf{A}^1_\mathbf{C} \setminus \{0\} = \mathbf{G}_{m, \mathbf{C}}$.
Il se trouve que la suite de Mayer--Vietoris vaut en cohomologie \'etale. On en
d\'eduit une suite exacte
$$
H_\etale^{i-1}(U_0\cap U_1, \Lambda) \to
H_\etale^i(\mathbf{P}^1_\mathbf{C}, \Lambda) \to
H_\etale^i(U_0, \Lambda) \oplus
H_\etale^i(U_1, \Lambda) \to H_\etale^i(U_0\cap U_1,
\Lambda).
$$
Pour obtenir la r\'eponse attendue, il faudrait montrer que la somme directe
figurant dans le troisi\`eme terme est nulle. En fait, comme pour la topologie
usuelle, on a bien
$$
H_\etale^q (\mathbf{A}^1_\mathbf{C}, \Lambda) = 0
\quad \text{ pour } q \geq 1,
$$
et
$$
H_\etale^q (\mathbf{A}^1_\mathbf{C} \setminus \{0\}, \Lambda) = \left\{
\begin{matrix}
\Lambda & \text{ si }q = 1\text{, et} \\
0 & \text{ pour }q \geq 2.
\end{matrix}
\right.
$$
Ces r\'esultats sont d\'ej\`a assez difficiles (quelle en serait une preuve
\'el\'ementaire?). Expliquons comment effectuer ce calcul une fois que nous
disposerons du formalisme de la cohomologie \'etale.

\medskip\noindent
{\bf Cohomologie sup\'erieure.} Elle est trait\'ee par le fait g\'en\'eral suivant:
si $X$ est une courbe affine sur $\mathbf{C}$, alors
$$
H_\etale^q (X, \mathbf{Z}/n\mathbf{Z}) = 0 \quad \text{ pour } q \geq 2.
$$
On le d\'emontre en consid\'erant le point g\'en\'erique de la courbe et en faisant
un peu de cohomologie galoisienne. Il ne reste donc \`a examiner que la
cohomologie en degr\'e 1.

\medskip\noindent
{\bf Cohomologie en degr\'e 1.} Nous utilisons les identifications suivantes:
\begin{eqnarray*}
H_\etale^1 (X, \mathbf{Z}/n\mathbf{Z}) = \left\{
\begin{matrix}
\text{faisceaux d'ensembles }\mathcal{F}\text{ sur le site \'etale }X_\etale
\text{ munis d'une} \\
\text{action }\mathbf{Z}/n\mathbf{Z} \times \mathcal{F} \to \mathcal{F}
\text{ telle que }\mathcal{F}\text{ soit un }\mathbf{Z}/n\mathbf{Z}\text{-torseur.}
\end{matrix}
\right\}
\Big/ \cong
\\
 = \left\{
\begin{matrix}
\text{morphismes }Y \to X\text{ finis \'etales, munis} \\
\text{ d'une action libre de }\mathbf{Z}/n\mathbf{Z}\text{ tels que }
X = Y/(\mathbf{Z}/n\mathbf{Z}).
\end{matrix}
\right\}
\Big/ \cong.
\end{eqnarray*}
La premi\`ere identification est tr\`es g\'en\'erale (elle vaut pour toute th\'eorie
cohomologique sur un site) et n'a rien \`a voir avec la topologie \'etale. La
seconde identification est une cons\'equence de la th\'eorie de la descente. Le
dernier ensemble d\'ecrit une collection d'objets g\'eom\'etriques que nous pouvons
manipuler concr\`etement.

\medskip\noindent
La courbe $\mathbf{A}^1_\mathbf{C}$ n'a aucun rev\^etement \'etale fini non
trivial; par cons\'equent,
$H_\etale^1 (\mathbf{A}^1_\mathbf{C}, \mathbf{Z}/n\mathbf{Z}) = 0$.
On peut le voir soit topologiquement, soit au moyen de l'argument du paragraphe
suivant.

\medskip\noindent
D\'ecrivons les rev\^etements \'etales finis
$\varphi : Y \to \mathbf{A}^1_\mathbf{C} \setminus \{0\}$.
Il suffit de consid\'erer le cas o\`u $Y$ est connexe, ce que nous supposons. Nous
allons d\'eterminer les possibilit\'es pour $Y$ en appliquant la formule de
Riemann--Hurwitz (c'est bien s\^ur un peu superflu, et le lecteur peut passer la
section suivante s'il le souhaite). Supposons que ce morphisme soit de degr\'e
$n$ et consid\'erons une compactification projective
$$
\xymatrix{
{Y\ } \ar@{^{(}->}[r] \ar[d]^\varphi &
{\bar Y} \ar[d]^{\bar\varphi} \\
{\mathbf{A}^1_\mathbf{C} \setminus \{0\}} \ar@{^{(}->}[r] &
{\mathbf{P}^1_\mathbf{C}}
}
$$
Bien que $\varphi$ soit \'etale et non ramifi\'e, $\bar{\varphi}$ peut \^etre
ramifi\'e en 0 et en $\infty$. Supposons que les ant\'ec\'edents de 0 soient les
points $y_1, \ldots, y_r$, d'indices de ramification $e_1, \ldots e_r$, et que
les ant\'ec\'edents de $\infty$ soient les points $y_1', \ldots, y_s'$, d'indices
de ramification $d_1, \ldots d_s$. En particulier,
$\sum e_i = n = \sum d_j$. La formule de Riemann--Hurwitz donne
$$
2 g_Y - 2 = -2n + \sum (e_i - 1) + \sum (d_j - 1)
$$
et donc $g_Y = 0$, $r = s = 1$ et $e_1 = d_1 = n$.
Ainsi $Y \cong {\mathbf{A}^1_\mathbf{C} \setminus \{0\}}$, et l'on voit
facilement que $\varphi(z) = \lambda z^n$ pour un certain
$\lambda \in \mathbf{C}^*$. Apr\`es avoir reparam\'etr\'e $Y$, nous pouvons
supposer $\lambda = 1$. Notre rev\^etement s'obtient donc en extrayant une
racine $n$-i\`eme de la coordonn\'ee sur
$\mathbf{A}^1_{\mathbf{C}} \setminus \{0\}$.

\medskip\noindent
Rappelons qu'il nous faut classifier les rev\^etements de
${\mathbf{A}^1_\mathbf{C} \setminus \{0\}}$ munis d'une action libre de
$\mathbf{Z}/n\mathbf{Z}$. Dans notre cas, toute action de ce type correspond
\`a un automorphisme de $Y$ envoyant $z$ sur $\zeta_n z$, o\`u $\zeta_n$ est une
racine primitive $n$-i\`eme de l'unit\'e. Il existe $\phi(n)$ telles actions (ici,
$\phi(n)$ d\'esigne l'indicatrice d'Euler). Il existe donc exactement $\phi(n)$
rev\^etements \'etales finis connexes munis d'une action libre donn\'ee de
$\mathbf{Z}/n\mathbf{Z}$, chacun correspondant \`a une racine primitive
$n$-i\`eme de l'unit\'e. Nous laissons au lecteur le soin de v\'erifier que les
rev\^etements \'etales finis non connexes de degr\'e $n$ de
$\mathbf{A}^1_{\mathbf{C}} \setminus \{0\}$, munis d'une action libre donn\'ee
de $\mathbf{Z}/n\mathbf{Z}$, correspondent bijectivement aux racines
$n$-i\`emes non primitives de $1$. Autrement dit, ce calcul montre que
$$
H_\etale^1 (\mathbf{A}^1_\mathbf{C} \setminus \{0\},
\mathbf{Z}/n\mathbf{Z}) =
\Hom(\mu_n(\mathbf{C}), \mathbf{Z}/n\mathbf{Z}) \cong \mathbf{Z}/n\mathbf{Z}.
$$
La premi\`ere identification est canonique, la seconde ne l'est pas; voir la
remarque \ref{remark-normalize-H1-Gm}.
Puisque la preuve de la formule de Riemann--Hurwitz ne fait pas appel au calcul
de la cohomologie, ce qui pr\'ec\`ede constitue effectivement une preuve (\`a
condition de compl\'eter les d\'etails concernant l'annulation, etc.).




\section{Coefficients non de torsion}
\label{section-nontorsion}

\noindent
Pour \'etudier les coefficients non de torsion, on pose la d\'efinition suivante:
$$
H_\etale^i (X, \mathbf{Q}_\ell) :=
\left( \lim_n H_\etale^i(X, \mathbf{Z}/\ell^n\mathbf{Z}) \right)
\otimes_{\mathbf{Z}_\ell} \mathbf{Q}_\ell.
$$
Le symbole $\lim_n$ d\'esigne la {\it limite} du syst\`eme de groupes de
cohomologie $H_\etale^i(X, \mathbf{Z}/\ell^n\mathbf{Z})$ index\'e par $n$; voir
Cat\'egories, section \ref{categories-section-posets-limits}.
Nous devrons donc \'etudier des syst\`emes de faisceaux satisfaisant certaines
conditions de compatibilit\'e.




\section{Th\'eorie des faisceaux}
\label{section-sheaf-theory}
%9.10.09

\noindent
Nous abordons maintenant v\'eritablement les sites et les faisceaux.
Une quantit\'e impressionnante de r\'esultats abstraits utiles pourrait trouver place
dans les quelques sections qui suivent. Une partie de ces r\'esultats est d\'evelopp\'ee dans des
chapitres ant\'erieurs, tels que ceux sur les sites, les modules sur les sites et la cohomologie
sur les sites. Nous nous effor\c{c}ons de ne pas ajouter trop de mati\`ere ici, mais seulement
ce qu'il faut pour que la suite de ce chapitre soit compr\'ehensible.




\section{Pr\'efaisceaux}
\label{section-presheaves}

\noindent
Une r\'ef\'erence pour cette section est
Sites, section \ref{sites-section-presheaves}.

\begin{definition}
\label{definition-presheaf}
Soit $\mathcal{C}$ une cat\'egorie. Un {\it pr\'efaisceau d'ensembles} (respectivement, un
{\it pr\'efaisceau ab\'elien}) sur $\mathcal{C}$ est un foncteur $\mathcal{C}^{opp} \to
\textit{Ensembles}$ (resp.\ $\textit{Ab}$).
\end{definition}

\noindent
{\bf Terminologie.} Si $U \in \Ob(\mathcal{C})$, les \'el\'ements de
$\mathcal{F}(U)$ sont appel\'es les {\it sections} de $\mathcal{F}$ au-dessus de
$U$. Pour $\varphi : V \to U$ dans $\mathcal{C}$, l'application
$\mathcal{F}(\varphi) : \mathcal{F}(U) \to \mathcal{F}(V)$
est appel\'ee l'{\it application de restriction} et est souvent not\'ee $s \mapsto s|_V$
ou parfois $s \mapsto \varphi^*s$. La notation $s|_V$ est ambigu\"e
car l'application de restriction d\'epend de $\varphi$, mais il s'agit d'un abus
de notation usuel. Nous employons aussi la notation
$\Gamma(U, \mathcal{F}) = \mathcal{F}(U)$.

\medskip\noindent
Dire que $\mathcal{F}$ est un foncteur signifie que, si
$W \to V \to U$ sont des morphismes de $\mathcal{C}$ et si
$s \in \Gamma(U, \mathcal{F})$, alors
$(s|_V)|_W = s |_W$, avec l'abus de
notation que nous venons de voir. De plus, les applications de restriction correspondant aux
morphismes identit\'e $\text{id}_U : U \to U$ sont l'identit\'e.

\medskip\noindent
La cat\'egorie des pr\'efaisceaux d'ensembles (respectivement, des pr\'efaisceaux ab\'eliens) sur
$\mathcal{C}$ est not\'ee $\textit{PSh} (\mathcal{C})$ (resp. $\textit{PAb}
(\mathcal{C})$). C'est la cat\'egorie des foncteurs de $\mathcal{C}^{opp}$ vers
$\textit{Ensembles}$ (resp. $\textit{Ab}$), c'est-\`a-dire que les morphismes de
pr\'efaisceaux sont les transformations naturelles de foncteurs. Nous ne consid\'erons les
cat\'egories $\textit{PSh}(\mathcal{C})$ et $\textit{PAb}(\mathcal{C})$
que lorsque la cat\'egorie $\mathcal{C}$ est petite. (Notre convention est qu'une cat\'egorie
est petite sauf mention contraire et, si elle ne l'est pas, elle doit \^etre
r\'epertori\'ee dans Cat\'egories, remarque \ref{categories-remark-big-categories}.)

\begin{example}
\label{example-representable-presheaf}
\'Etant donn\'e un objet $X \in \Ob(\mathcal{C})$, consid\'erons la r\`egle
$$
\begin{matrix}
U & \longmapsto & h_X(U) = \Mor_\mathcal{C}(U, X) \\
V \xrightarrow{\varphi} U & \longmapsto &
(\psi \mapsto \psi \circ \varphi) : h_X(U) \to h_X(V).
\end{matrix}
$$
Elle d\'efinit un foncteur $h_X : \mathcal{C}^{opp} \to \textit{Ensembles}$
et donc un pr\'efaisceau. On l'appelle le
{\it pr\'efaisceau repr\'esentable associ\'e \`a $X$.}
Il n'est pas vrai que les pr\'efaisceaux repr\'esentables soient des faisceaux pour toute topologie sur
tout site.
\end{example}

\begin{lemma}[Yoneda]
\label{lemma-yoneda}
\begin{slogan}
Les morphismes entre objets sont en bijection avec les transformations naturelles
entre les foncteurs qu'ils repr\'esentent.
\end{slogan}
Soit $\mathcal{C}$ une cat\'egorie, et soient $X, Y \in
\Ob(\mathcal{C})$. Il existe une bijection naturelle
$$
\begin{matrix}
\Mor_\mathcal{C}(X, Y) &
\longrightarrow &
\Mor_{\textit{PSh}(\mathcal{C})} (h_X, h_Y) \\
\psi &
\longmapsto &
h_\psi = \psi \circ - : h_X \to h_Y.
\end{matrix}
$$
\end{lemma}

\begin{proof}
Voir
Cat\'egories, lemme \ref{categories-lemma-yoneda}.
\end{proof}




\section{Sites}
\label{section-sites}


\begin{definition}
\label{definition-family-morphisms-fixed-target}
Soit $\mathcal{C}$ une cat\'egorie. Une {\it famille de morphismes \`a cible fix\'ee}
$\mathcal{U} = \{\varphi_i : U_i \to U\}_{i\in I}$ est la donn\'ee de
\begin{enumerate}
\item un objet $U \in \mathcal{C}$,
\item un ensemble $I$ (\'eventuellement vide), et
\item pour tout $i\in I$, un morphisme $\varphi_i : U_i \to U$ de $\mathcal{C}$
de cible $U$.
\end{enumerate}
\end{definition}

\noindent
Il existe une notion de {\it morphisme de familles de morphismes \`a cible
fix\'ee}. Un cas particulier est la notion de {\it raffinement}.
On pourra consulter \`a ce sujet
Sites, section \ref{sites-section-refinements}.

\begin{definition}
\label{definition-site}
Un {\it site}\footnote{Ce que nous appelons un site est appel\'e une cat\'egorie munie d'une
pr\'etopologie dans \cite[Expos\'e II, D\'efinition 1.3]{SGA4}.
Dans \cite{ArtinTopologies}, on l'appelle une cat\'egorie munie d'une topologie de
Grothendieck.} est la donn\'ee d'une cat\'egorie $\mathcal{C}$ et d'un ensemble
$\text{Cov}(\mathcal{C})$ form\'e de familles de morphismes \`a cible fix\'ee
appel\'ees {\it recouvrements}, qui satisfont aux axiomes suivants :
\begin{enumerate}
\item (isomorphisme) si $\varphi : V \to U$ est un isomorphisme dans $\mathcal{C}$,
alors $\{\varphi : V \to U\}$ est un recouvrement,
\item (localit\'e) si $\{\varphi_i : U_i \to U\}_{i\in I}$ est un recouvrement et si
pour tout $i \in I$ on se donne un recouvrement
$\{\psi_{ij} : U_{ij} \to U_i \}_{j\in I_i}$, alors
$$
\{
\varphi_i \circ \psi_{ij} : U_{ij} \to U
\}_{(i, j)\in \prod_{i\in I} \{i\} \times I_i}
$$
est aussi un recouvrement, et
\item (changement de base) si $\{U_i \to U\}_{i\in I}$
est un recouvrement et si $V \to U$ est un morphisme dans $\mathcal{C}$, alors
\begin{enumerate}
\item pour tout $i \in I$, le produit fibr\'e
$U_i \times_U V$ existe dans $\mathcal{C}$, et
\item $\{U_i \times_U V \to V\}_{i\in I}$ est un recouvrement.
\end{enumerate}
\end{enumerate}
\end{definition}

\noindent
Pour nous, la cat\'egorie sous-jacente \`a un site est toujours ``petite'', c.-\`a-d. que sa
collection d'objets forme un ensemble et que la collection des recouvrements d'un
site est aussi un ensemble (comme dans la d\'efinition ci-dessus). Dans ce chapitre,
nous omettrons le plus souvent les arguments qui ram\`enent \`a un ensemble la collection
des objets et des recouvrements. Pour une discussion plus approfondie, voir
Sites, remarque \ref{sites-remark-no-big-sites}.

\begin{example}
\label{example-site-topological-space}
Si $X$ est un espace topologique, il poss\`ede un site associ\'e $X_{Zar}$
d\'efini comme suit : les objets de $X_{Zar}$ sont les ouverts de $X$,
les morphismes entre eux sont les applications d'inclusion, et les recouvrements sont
les recouvrements topologiques usuels (surjectifs). Observons que si
$U, V \subset W \subset X$ sont des ouverts, alors $U \times_W V = U \cap V$
existe : cette cat\'egorie admet les produits fibr\'es. Toutes les v\'erifications sont imm\'ediates et
tout fonctionne comme pr\'evu.
\end{example}




\section{Faisceaux}
\label{section-sheaves}

\begin{definition}
\label{definition-sheaf}
Un pr\'efaisceau $\mathcal{F}$ d'ensembles (resp. un pr\'efaisceau ab\'elien) sur un site
$\mathcal{C}$ est dit {\it s\'epar\'e} si, pour tout recouvrement
$\{\varphi_i : U_i \to U\}_{i\in I} \in \text{Cov} (\mathcal{C})$,
l'application
$$
\mathcal{F}(U) \longrightarrow \prod\nolimits_{i\in I} \mathcal{F}(U_i)
$$
est injective. Ici, l'application est $s \mapsto (s|_{U_i})_{i\in I}$.
Le pr\'efaisceau $\mathcal{F}$ est un {\it faisceau} si, pour tout recouvrement
$\{\varphi_i : U_i \to U\}_{i\in I} \in \text{Cov} (\mathcal{C})$, le
diagramme
\begin{equation}
\label{equation-sheaf-axiom}
\xymatrix{
\mathcal{F}(U) \ar[r] &
\prod_{i\in I} \mathcal{F}(U_i) \ar@<1ex>[r] \ar@<-1ex>[r] &
\prod_{i, j \in I} \mathcal{F}(U_i \times_U U_j),
}
\end{equation}
o\`u la premi\`ere application est $s \mapsto (s|_{U_i})_{i\in I}$ et les deux
applications de droite sont
$(s_i)_{i\in I} \mapsto (s_i |_{U_i \times_U U_j})$ et
$(s_i)_{i\in I} \mapsto (s_j |_{U_i \times_U U_j})$,
est un diagramme \'egalisateur dans la cat\'egorie des ensembles (resp.\ des groupes ab\'eliens).
\end{definition}

\begin{remark}
\label{remark-empty-covering}
Pour le recouvrement vide (o\`u $I = \emptyset$), ceci implique que
$\mathcal{F}(\emptyset)$ est un produit vide, qui est un objet final de la
cat\'egorie correspondante (un singleton, tant pour $\textit{Ensembles}$ que pour
$\textit{Ab}$).
\end{remark}

\begin{example}
\label{example-sheaf-site-space}
En explicitant cette d\'efinition pour le site $X_{Zar}$ associ\'e \`a un espace
topologique, voir l'exemple \ref{example-site-topological-space}, on retrouve la notion usuelle
de faisceau.
\end{example}

\begin{definition}
\label{definition-category-sheaves}
Nous notons $\Sh(\mathcal{C})$ (resp.\ $\textit{Ab}(\mathcal{C})$)
la sous-cat\'egorie pleine de $\textit{PSh}(\mathcal{C})$
(resp.\ $\textit{PAb}(\mathcal{C})$) dont les objets sont les faisceaux. C'est la
{\it cat\'egorie des faisceaux d'ensembles} (resp.\ celle des {\it faisceaux ab\'eliens}) sur
$\mathcal{C}$.
\end{definition}




\section{L'exemple des G-ensembles}
\label{section-G-sets}

\noindent
Soit $G$ un groupe et d\'efinissons un site $\mathcal{T}_G$ comme suit : la cat\'egorie
sous-jacente est la cat\'egorie des $G$-ensembles, c'est-\`a-dire que ses objets sont des ensembles munis
d'une action \`a gauche de $G$ et que les morphismes sont les applications \'equivariantes ; les
recouvrements de $\mathcal{T}_G$ sont les familles
$\{\varphi_i : U_i \to U\}_{i\in I}$ satisfaisant
$U = \bigcup_{i\in I} \varphi_i(U_i)$.

\medskip\noindent
Il existe un objet particulier du site $\mathcal{T}_G$, \`a savoir le $G$-ensemble $G$
muni de son action naturelle par translations \`a gauche. Nous le notons ${}_G G$.
Observons qu'il existe un isomorphisme naturel de groupes
$$
\begin{matrix}
\rho : & G^{opp} & \longrightarrow & \text{Aut}_{G\textit{-Ensembles}}({}_G G) \\
& g & \longmapsto & (h \mapsto hg).
\end{matrix}
$$
En particulier, pour tout pr\'efaisceau $\mathcal{F}$, l'ensemble $\mathcal{F}({}_G G)$
h\'erite d'une action de $G$ au moyen de $\rho$. (Notons que, par contravariance de
$\mathcal{F}$, l'ensemble $\mathcal{F}({}_G G)$ est encore un $G$-ensemble \`a gauche.) En fait,
le foncteur
$$
\begin{matrix}
\Sh(\mathcal{T}_G) & \longrightarrow & G\textit{-Ensembles} \\
\mathcal{F} & \longmapsto & \mathcal{F}({}_G G)
\end{matrix}
$$
est une \'equivalence de cat\'egories. Son quasi-inverse est le foncteur $X \mapsto
h_X$. Sans donner la d\'emonstration compl\`ete (que l'on trouvera dans
Sites, section \ref{sites-section-example-sheaf-G-sets}),
essayons d'expliquer pourquoi il en est ainsi.
\begin{enumerate}
\item
Si $S$ est un $G$-ensemble, nous pouvons le d\'ecomposer en orbites $S = \coprod_{i\in I}
O_i$. L'axiome des faisceaux appliqu\'e au recouvrement $\{O_i \to S\}_{i\in I}$ affirme que
$$
\xymatrix{
\mathcal{F}(S) \ar[r] &
\prod_{i\in I} \mathcal{F}(O_i) \ar@<1ex>[r] \ar@<-1ex>[r] &
\prod_{i, j \in I} \mathcal{F}(O_i \times_S O_j)
}
$$
est un diagramme \'egalisateur. Comme les produits fibr\'es dans $G\textit{-Ensembles}$
sont induits par les produits fibr\'es dans $\textit{Ensembles}$, et comme
$\mathcal{F}(\emptyset)$ est un $G$-ensemble r\'eduit \`a un point, nous obtenons
$$
\prod_{i, j \in I} \mathcal{F}(O_i \times_S O_j) = \prod_{i \in I}
\mathcal{F}(O_i)
$$
et les deux applications ci-dessus sont en fait identiques. Ainsi, l'axiome des faisceaux affirme simplement
que $\mathcal{F}(S) = \prod_{i\in I} \mathcal{F}(O_i)$.
\item
Si $S$ est le $G$-ensemble $S= G/H$ et si $\mathcal{F}$ est un faisceau sur $\mathcal{T}_G$,
nous affirmons que
$$
\mathcal{F}(G/H) = \mathcal{F}({}_G G)^H
$$
et en particulier que $\mathcal{F}(\{*\}) = \mathcal{F}({}_G G)^G$. Pour le voir,
appliquons l'axiome des faisceaux au recouvrement $\{ {}_G G \to G/H \}$ de $S$. Nous
avons
\begin{eqnarray*}
{}_G G \times_{G/H} {}_G G & \cong & G \times H \\
(g_1, g_2) & \longmapsto & (g_1, g_1^{-1} g_2)
\end{eqnarray*}
Le membre de gauche est une r\'eunion disjointe de copies de ${}_G G$ (comme $G$-ensemble). L'axiome des faisceaux
s'\'ecrit donc
$$
\xymatrix{
\mathcal{F} (G/H) \ar[r] &
\mathcal{F}({}_G G) \ar@<1ex>[r] \ar@<-1ex>[r] &
\prod_{h\in H} \mathcal{F}({}_G G)
}
$$
o\`u les deux applications de droite sont $s \mapsto (s)_{h \in H}$ et $s \mapsto
(hs)_{h \in H}$. Par cons\'equent, $\mathcal{F}(G/H) = \mathcal{F}({}_G G)^H$, comme
annonc\'e.
\end{enumerate}
Ceci ne d\'emontre pas tout \`a fait l'\'equivalence de cat\'egories annonc\'ee, mais montre au
moins qu'un faisceau $\mathcal{F}$ est enti\`erement d\'etermin\'e par ses sections au-dessus de
${}_G G$. On trouvera les d\'etails (ainsi que des remarques ensemblistes) dans
Sites, section \ref{sites-section-example-sheaf-G-sets}.




\section{Passage au faisceau associ\'e}
\label{section-sheafification}

\begin{definition}
\label{definition-0-cech}
Soit $\mathcal{F}$ un pr\'efaisceau sur le site $\mathcal{C}$ et soit
$\mathcal{U} = \{U_i \to U\} \in \text{Cov} (\mathcal{C})$.
Nous d\'efinissons le {\it groupe de cohomologie de {\v C}ech de degr\'e z\'ero} de
$\mathcal{F}$ relativement \`a $\mathcal{U}$ par
$$
\check H^0 (\mathcal{U}, \mathcal{F}) =
\left\{
(s_i)_{i\in I} \in \prod\nolimits_{i\in I }\mathcal{F}(U_i)
\text{ tels que }
s_i|_{U_i \times_U U_j} = s_j |_{U_i \times_U U_j}
\right\}.
$$
\end{definition}

\noindent
Il existe une application canonique
$\mathcal{F}(U) \to \check H^0 (\mathcal{U}, \mathcal{F})$,
$s \mapsto (s |_{U_i})_{i\in I}$.
Nous appelons {\it morphisme de recouvrements} d'un recouvrement
$\mathcal{V} = \{V_j \to V\}_{j \in J}$ vers $\mathcal{U}$ un triplet
$(\chi, \alpha, \chi_j)$, o\`u
$\chi : V \to U$ est un morphisme,
$\alpha : J \to I$ est une application d'ensembles, et, pour tout
$j \in J$, le morphisme $\chi_j$ s'ins\`ere dans un diagramme commutatif
$$
\xymatrix{
V_j \ar[rr]_{\chi_j} \ar[d] & & U_{\alpha(j)} \ar[d] \\
V \ar[rr]^\chi & & U.
}
$$
\`A partir des donn\'ees $\chi, \alpha, \{\chi_j\}_{j \in J}$, nous d\'efinissons
\begin{eqnarray*}
\check H^0(\mathcal{U}, \mathcal{F}) & \longrightarrow &
\check H^0(\mathcal{V}, \mathcal{F}) \\
(s_i)_{i\in I} & \longmapsto &
\left(\chi_j^*\left(s_{\alpha(j)}\right)\right)_{j\in J}.
\end{eqnarray*}
Nous affirmons alors que
\begin{enumerate}
\item l'application est bien d\'efinie, et
\item elle ne d\'epend que de $\chi$ et est ind\'ependante du choix de
$\alpha, \{\chi_j\}_{j \in J}$.
\end{enumerate}
Nous omettons la d\'emonstration du premier point.
Pour voir (2), consid\'erons un autre triplet $(\psi, \beta, \psi_j)$ tel que
$\chi = \psi$. On a alors le diagramme commutatif
$$
\xymatrix{
V_j \ar[rrr]_{(\chi_j, \psi_j)} \ar[dd] & & &
U_{\alpha(j)} \times_U U_{\beta(j)} \ar[dl] \ar[dr] \\
& & U_{\alpha(j)} \ar[dr] & &
U_{\beta(j)} \ar[dl] \\
V \ar[rrr]^{\chi = \psi} & & & U.
}
$$
\'Etant donn\'ee une section $s \in \mathcal{F}(\mathcal{U})$, son image dans
$\mathcal{F}(V_j)$ par l'application d\'efinie par
$(\chi, \alpha, \{\chi_j\}_{j \in J})$
est $\chi_j^*s_{\alpha(j)}$, et
son image par l'application d\'efinie par $(\psi, \beta, \{\psi_j\}_{j \in J})$
est $\psi_j^*s_{\beta(j)}$. Ces deux \'el\'ements
sont \'egaux car, par hypoth\`ese, $s \in \check H^0(\mathcal{U}, \mathcal{F})$ ;
ils sont donc tous deux \'egaux au tir\'e en arri\`ere de la valeur commune
$$
s_{\alpha(j)}|_{U_{\alpha(j)} \times_U U_{\beta(j)}} =
s_{\beta(j)}|_{U_{\alpha(j)} \times_U U_{\beta(j)}}
$$
par l'application $(\chi_j, \psi_j)$ du diagramme.

\begin{theorem}
\label{theorem-sheafification}
Soient $\mathcal{C}$ un site et $\mathcal{F}$ un pr\'efaisceau sur $\mathcal{C}$.
\begin{enumerate}
\item La r\`egle
$$
U \mapsto \mathcal{F}^+(U) :=
\colim_{\mathcal{U} \text{ recouvrement de }U}
\check H^0(\mathcal{U}, \mathcal{F})
$$
est un pr\'efaisceau. De plus, la limite inductive est filtrante.
\item Il existe un morphisme canonique de pr\'efaisceaux $\mathcal{F} \to \mathcal{F}^+$.
\item Si $\mathcal{F}$ est un pr\'efaisceau s\'epar\'e, alors $\mathcal{F}^+$ est un faisceau
et le morphisme de (2) est injectif.
\item $\mathcal{F}^+$ est un pr\'efaisceau s\'epar\'e.
\item $\mathcal{F}^\# = (\mathcal{F}^+)^+$ est un faisceau, et le morphisme canonique
induit un isomorphisme fonctoriel
$$
\Hom_{\textit{PSh}(\mathcal{C})}(\mathcal{F}, \mathcal{G}) =
\Hom_{\Sh(\mathcal{C})}(\mathcal{F}^\#, \mathcal{G})
$$
pour tout $\mathcal{G} \in \Sh(\mathcal{C})$.
\end{enumerate}
\end{theorem}

\begin{proof}
Voir Sites, th\'eor\`eme \ref{sites-theorem-plus}.
\end{proof}

\noindent
Autrement dit, le passage au faisceau associ\'e, donn\'e par le morphisme naturel
$\mathcal{F} \to \mathcal{F}^\#$, est adjoint \`a gauche au foncteur d'oubli
$\Sh(\mathcal{C}) \to \textit{PSh}(\mathcal{C})$.




\section{Cohomologie}
\label{section-cohomology}

\noindent
Le r\'esultat fondamental suivant permet de d\'efinir la cohomologie
des faisceaux ab\'eliens sur les sites.

\begin{theorem}
\label{theorem-enough-injectives}
La cat\'egorie des faisceaux ab\'eliens sur un site est une cat\'egorie ab\'elienne
qui poss\`ede assez d'injectifs.
\end{theorem}

\begin{proof}
Voir
Modules sur les sites, lemme \ref{sites-modules-lemma-abelian-abelian} et
Injectifs, th\'eor\`eme \ref{injectives-theorem-sheaves-injectives}.
\end{proof}

\noindent
On peut donc d\'efinir la cohomologie au moyen des foncteurs d\'eriv\'es \`a droite du foncteur
des sections : si $U \in \Ob(\mathcal{C})$ et
$\mathcal{F} \in \textit{Ab}(\mathcal{C})$,
$$
H^p(U, \mathcal{F}) :=
R^p\Gamma(U, \mathcal{F}) =
H^p(\Gamma(U, \mathcal{I}^\bullet))
$$
o\`u $\mathcal{F} \to \mathcal{I}^\bullet$ est une r\'esolution injective. Pour ce faire,
il faut v\'erifier que le foncteur $\Gamma(U, -)$ est exact \`a gauche. C'est
le cas, et c'est en partie pourquoi la cat\'egorie $\textit{Ab}(\mathcal{C})$ est ab\'elienne ;
voir
Modules sur les sites, lemme \ref{sites-modules-lemma-abelian-abelian}.
Pour une \'etude plus g\'en\'erale de la cohomologie sur les sites (y compris le
foncteur des sections globales et ses foncteurs d\'eriv\'es \`a droite), voir
Cohomologie sur les sites, section \ref{sites-cohomology-section-cohomology-sheaves}.



\section{La topologie fpqc}
\label{section-fpqc}
%9.15.09

\noindent
Avant d'aborder la cohomologie \'etale, nous \'etudions quelque peu la topologie fpqc, car
elle se comporte bien avec les faisceaux quasi-coh\'erents.

\begin{definition}
\label{definition-fpqc-covering}
Soit $T$ un sch\'ema. Un {\it recouvrement fpqc} de $T$ est une famille
$\{ \varphi_i : T_i \to T\}_{i \in I}$ telle que
\begin{enumerate}
\item chaque $\varphi_i$ est un morphisme plat et
$\bigcup_{i\in I} \varphi_i(T_i) = T$, et
\item pour tout ouvert affine $U \subset T$, il existe un ensemble fini
$K$, une application $\mathbf{i} : K \to I$ et des ouverts affines
$U_{\mathbf{i}(k)} \subset T_{\mathbf{i}(k)}$ tels que
$U = \bigcup_{k \in K} \varphi_{\mathbf{i}(k)}(U_{\mathbf{i}(k)})$.
\end{enumerate}
\end{definition}

\begin{remark}
\label{remark-fpqc}
La premi\`ere condition correspond \`a fp, sigle fran\c{c}ais de
{\it fid\`element plat}, qui signifie ``faithfully flat'' en anglais, et
la seconde \`a qc, {\it quasi-compact}.
La seconde partie de la premi\`ere condition est superflue lorsque la seconde est satisfaite.
\end{remark}

\begin{example}
\label{example-fpqc-coverings}
Voici des exemples de recouvrements fpqc.
\begin{enumerate}
\item Tout recouvrement ouvert de Zariski de $T$ est un recouvrement fpqc.
\item La famille $\{\Spec(B) \to \Spec(A)\}$ est un recouvrement fpqc
si et seulement si $A \to B$ est un morphisme d'anneaux fid\`element plat.
\item Si $f: X \to Y$ est plat, surjectif et quasi-compact, alors $\{ f: X\to
Y\}$ est un recouvrement fpqc.
\item Le morphisme
$\varphi :
\coprod_{x \in \mathbf{A}^1_k} \Spec(\mathcal{O}_{\mathbf{A}^1_k, x})
\to \mathbf{A}^1_k$,
o\`u $k$ est un corps, est plat et surjectif. Il n'est pas quasi-compact, et
en fait la famille $\{\varphi\}$ n'est pas un recouvrement fpqc.
\item \'Ecrivons
$\mathbf{A}^2_k = \Spec(k[x, y])$. Notons $i_x : D(x) \to \mathbf{A}^2_k$
et $i_y : D(y) \to \mathbf{A}^2_k$ les immersions ouvertes standard.
Alors les familles
$\{i_x, i_y, \Spec(k[[x, y]]) \to \mathbf{A}^2_k\}$
et
$\{i_x, i_y, \Spec(\mathcal{O}_{\mathbf{A}^2_k, 0}) \to \mathbf{A}^2_k\}$
sont des recouvrements fpqc.
\end{enumerate}
\end{example}

\begin{lemma}
\label{lemma-site-fpqc}
La collection des recouvrements fpqc de la cat\'egorie des sch\'emas
satisfait les axiomes d'un site.
\end{lemma}

\begin{proof}
Voir Topologies, lemme \ref{topologies-lemma-fpqc}.
\end{proof}

\noindent
Il semble que ce lemme permette de d\'efinir le site fpqc de la cat\'egorie
des sch\'emas. Cependant, un probl\`eme ensembliste surgit lorsqu'on
consid\`ere la topologie fpqc ; voir
Topologies, section \ref{topologies-section-fpqc}.
Il provient de ce que nous exigeons que les sites soient ``petits'', alors qu'aucune
petite cat\'egorie de sch\'emas ne peut contenir un syst\`eme cofinal de recouvrements fpqc d'un
sch\'ema non vide donn\'e. Bien que cela ne nous emp\^eche pas, \`a proprement parler,
de d\'efinir des sites fpqc ``partiels'',
il ne semble pas prudent de le faire. La solution consiste \`a autoriser
la notion de faisceau pour la topologie fpqc (voir ci-dessous), mais \`a interdire de
consid\'erer la cat\'egorie de tous les faisceaux fpqc.

\begin{definition}
\label{definition-sheaf-property-fpqc}
Soit $S$ un sch\'ema. La cat\'egorie des sch\'emas sur $S$ est not\'ee
$\Sch/S$. Consid\'erons un foncteur
$\mathcal{F} : (\Sch/S)^{opp} \to \textit{Ensembles}$, autrement dit
un pr\'efaisceau d'ensembles. Nous dirons que $\mathcal{F}$
{\it v\'erifie la propri\'et\'e de faisceau pour la topologie fpqc}
si, pour tout recouvrement fpqc $\{U_i \to U\}_{i \in I}$ de sch\'emas sur $S$,
le diagramme (\ref{equation-sheaf-axiom}) est un diagramme \'egalisateur.
\end{definition}

\noindent
Nous dirons de m\^eme que $\mathcal{F}$
{\it v\'erifie la propri\'et\'e de faisceau pour la topologie de Zariski} si, pour
tout recouvrement ouvert $U = \bigcup_{i \in I} U_i$, le diagramme
(\ref{equation-sheaf-axiom}) est un diagramme \'egalisateur. Voir
Sch\'emas, d\'efinition \ref{schemes-definition-representable-by-open-immersions}.
Il est clair que cela revient \`a dire que, pour tout sch\'ema $T$ sur $S$, la
restriction de $\mathcal{F}$ aux ouverts de $T$ est un faisceau (au sens usuel).

\begin{lemma}
\label{lemma-fpqc-sheaves}
Soit $\mathcal{F}$ un pr\'efaisceau sur $\Sch/S$. Alors
$\mathcal{F}$ v\'erifie la propri\'et\'e de faisceau pour la topologie fpqc
si et seulement si
\begin{enumerate}
\item $\mathcal{F}$ v\'erifie la propri\'et\'e de faisceau relativement \`a la
topologie de Zariski, et
\item pour tout morphisme fid\`element plat $\Spec(B) \to \Spec(A)$
de sch\'emas affines sur $S$, l'axiome de faisceau est v\'erifi\'e pour le recouvrement
$\{\Spec(B) \to \Spec(A)\}$. Plus pr\'ecis\'ement, cela signifie que
$$
\xymatrix{
\mathcal{F}(\Spec(A)) \ar[r] &
\mathcal{F}(\Spec(B)) \ar@<1ex>[r] \ar@<-1ex>[r] &
\mathcal{F}(\Spec(B \otimes_A B))
}
$$
est un diagramme \'egalisateur.
\end{enumerate}
\end{lemma}

\begin{proof}
Voir Topologies, lemme \ref{topologies-lemma-sheaf-property-fpqc}.
\end{proof}

\noindent
Une autre mani\`ere de consid\'erer un pr\'efaisceau $\mathcal{F}$ sur
$\Sch/S$ satisfaisant la condition de faisceau pour la
topologie fpqc consiste \`a se donner les donn\'ees suivantes :
\begin{enumerate}
\item pour tout $T/S$, un faisceau usuel (c'est-\`a-dire pour la topologie de Zariski) $\mathcal{F}_T$ sur
$T_{Zar}$,
\item pour tout morphisme $f : T' \to T$ sur $S$, un morphisme de restriction
$f^{-1}\mathcal{F}_T \to \mathcal{F}_{T'}$
\end{enumerate}
satisfaisant
\begin{enumerate}
\item[(a)] les morphismes de restriction sont fonctoriels,
\item[(b)] si $f : T' \to T$ est une immersion ouverte, alors le morphisme de restriction
$f^{-1}\mathcal{F}_T \to \mathcal{F}_{T'}$ est un isomorphisme, et
\item[(c)] pour tout morphisme fid\`element plat
$\Spec(B) \to \Spec(A)$ sur $S$, le diagramme
$$
\xymatrix{
\mathcal{F}_{\Spec(A)}(\Spec(A)) \ar[r] &
\mathcal{F}_{\Spec(B)}(\Spec(B)) \ar@<1ex>[r] \ar@<-1ex>[r] &
\mathcal{F}_{\Spec(B \otimes_A B)}(\Spec(B \otimes_A B))
}
$$
est un diagramme \'egalisateur.
\end{enumerate}
Les donn\'ees (1) et (2), avec les conditions (a), (b), d\'efinissent un pr\'efaisceau
sur $\Sch/S$ satisfaisant la condition de faisceau pour la topologie de Zariski.
D'apr\`es le lemme \ref{lemma-fpqc-sheaves}, la condition (c) suffit alors pour obtenir la
condition de faisceau pour la topologie fpqc.

\begin{example}
\label{example-quasi-coherent}
Consid\'erons le pr\'efaisceau
$$
\begin{matrix}
\mathcal{F} : & (\Sch/S)^{opp} & \longrightarrow & \textit{Ab} \\
& T/S & \longmapsto & \Gamma(T, \Omega_{T/S}).
\end{matrix}
$$
La compatibilit\'e des diff\'erentielles avec la localisation entra\^ine que
$\mathcal{F}$ est un faisceau sur le site de Zariski.
Il ne v\'erifie toutefois pas la condition de faisceau pour la topologie fpqc.
Consid\'erons en effet le cas
$S = \Spec(\mathbf{F}_p)$ et le morphisme
$$
\varphi :
V = \Spec(\mathbf{F}_p[v])
\to
U = \Spec(\mathbf{F}_p[u])
$$
d\'efini par l'envoi de $u$ sur $v^p$. La famille $\{\varphi\}$ est un recouvrement fpqc,
mais le morphisme de restriction
$\mathcal{F}(U) \to \mathcal{F}(V)$
envoie le g\'en\'erateur $\text{d}u$ sur $\text{d}(v^p) = 0$, donc
c'est le morphisme nul, et le diagramme
$$
\xymatrix{
\mathcal{F}(U) \ar[r]^{0} &
\mathcal{F}(V) \ar@<1ex>[r] \ar@<-1ex>[r] &
\mathcal{F}(V \times_U V)
}
$$
n'est pas un diagramme \'egalisateur. Nous verrons plus loin que $\mathcal{F}$ donne en fait
lieu \`a un faisceau sur les sites \'etale et lisse.
\end{example}

\begin{lemma}
\label{lemma-representable-sheaf-fpqc}
Tout pr\'efaisceau repr\'esentable sur $\Sch/S$ satisfait \`a la
condition de faisceau pour la topologie fpqc.
\end{lemma}

\begin{proof}
Voir
Descente, lemme \ref{descent-lemma-fpqc-universal-effective-epimorphisms}.
\end{proof}

\noindent
Nous y reviendrons plus loin, car la d\'emonstration de ce fait utilise
la descente des faisceaux quasi-coh\'erents, que nous \'etudierons dans la
section suivante. Une mani\`ere \'el\'egante d'exprimer le lemme consiste \`a dire que
{\it la topologie fpqc est moins fine que la topologie canonique}, ou
que la topologie fpqc est {\it sous-canonique}. Dans le cadre des sites,
ceci est \'etudi\'e dans
Sites, section \ref{sites-section-representable-sheaves}.

\begin{remark}
\label{remark-fpqc-finest}
La topologie fpqc est plus fine que les topologies de Zariski, \'etale,
lisse, syntomique et fppf. Par cons\'equent, tout pr\'efaisceau
satisfaisant \`a la condition de faisceau pour la topologie fpqc est un
faisceau sur les sites de Zariski, \'etale, lisse, syntomique et fppf.
En particulier, les pr\'efaisceaux repr\'esentables sont des faisceaux sur
le site \'etale d'un sch\'ema,
par exemple.
\end{remark}

\begin{example}
\label{example-additive-group-sheaf}
Soit $S$ un sch\'ema.
Consid\'erons le sch\'ema en groupes additif $\mathbf{G}_{a, S} = \mathbf{A}^1_S$
sur $S$, voir
Groupo\"ides, exemple \ref{groupoids-example-additive-group}.
Le pr\'efaisceau repr\'esentable associ\'e est donn\'e par
$$
h_{\mathbf{G}_{a, S}}(T) =
\Mor_S(T, \mathbf{G}_{a, S}) =
\Gamma(T, \mathcal{O}_T).
$$
D'apr\`es ce qui pr\'ec\`ede, nous savons qu'il s'agit d'un pr\'efaisceau
d'ensembles satisfaisant \`a la condition de faisceau pour la topologie fpqc. Il s'agit manifestement
aussi d'un pr\'efaisceau d'anneaux. Nous pouvons donc le consid\'erer comme le foncteur
$$
\begin{matrix}
\mathcal{O} : &
(\Sch/S)^{opp} &
\longrightarrow &
\textit{Rings} \\
&
T/S &
\longmapsto &
\Gamma(T, \mathcal{O}_T)
\end{matrix}
$$
qui satisfait \`a la condition de faisceau pour la topologie fpqc.
On dispose de m\^eme d'une notion de $\mathcal{O}$-module, et ainsi de
suite.
\end{example}




\section{Descente fid\`element plate}
\label{section-fpqc-descent}

\noindent
Dans cette section, nous \'etudions la descente fid\`element plate des modules
quasi-coh\'erents. Plus pr\'ecis\'ement, nous d\'emontrerons que les modules quasi-coh\'erents
satisfont \`a la descente effective relativement aux recouvrements fpqc.

\begin{definition}
\label{definition-descent-datum}
Soit $\mathcal{U} = \{ t_i : T_i \to T\}_{i \in I}$ une famille de
morphismes de sch\'emas \`a cible fix\'ee. Une {\it donn\'ee de descente} pour les
faisceaux quasi-coh\'erents relativement \`a $\mathcal{U}$ est une famille
$((\mathcal{F}_i)_{i \in I}, (\varphi_{ij})_{i, j \in I})$ telle que
\begin{enumerate}
\item $\mathcal{F}_i$ est un faisceau quasi-coh\'erent sur $T_i$, et
\item $\varphi_{ij} : \text{pr}_0^* \mathcal{F}_i \to
\text{pr}_1^* \mathcal{F}_j$ est un isomorphisme de modules
sur $T_i \times_T T_j$,
\end{enumerate}
et que la {\it condition de cocycle} soit satisfaite : les diagrammes
$$
\xymatrix{
\text{pr}_0^*\mathcal{F}_i \ar[dr]_{\text{pr}_{02}^*\varphi_{ik}}
\ar[rr]^{\text{pr}_{01}^*\varphi_{ij}} & &
\text{pr}_1^*\mathcal{F}_j \ar[dl]^{\text{pr}_{12}^*\varphi_{jk}} \\
& \text{pr}_2^*\mathcal{F}_k
}
$$
commutent sur $T_i \times_T T_j \times_T T_k$.
Cette donn\'ee de descente est dite {\it effective} s'il existe un faisceau
quasi-coh\'erent $\mathcal{F}$ sur $T$ et des isomorphismes de $\mathcal{O}_{T_i}$-modules
$\varphi_i : t_i^* \mathcal{F} \cong \mathcal{F}_i$ compatibles avec
les applications $\varphi_{ij}$, \`a savoir
$$
\varphi_{ij} = \text{pr}_1^* (\varphi_j) \circ \text{pr}_0^* (\varphi_i)^{-1}.
$$
\end{definition}

\noindent
Dans cette section et la suivante, nous \'etudions quelques ingr\'edients de la
d\'emonstration du th\'eor\`eme ci-dessous, ainsi que des r\'esultats connexes.

\begin{theorem}
\label{theorem-descent-quasi-coherent}
Si $\mathcal{V} = \{T_i \to T\}_{i\in I}$ est un recouvrement fpqc, alors toutes
les donn\'ees de descente pour les faisceaux quasi-coh\'erents relativement \`a $\mathcal{V}$
sont effectives.
\end{theorem}

\begin{proof}
Voir
Descente, proposition \ref{descent-proposition-fpqc-descent-quasi-coherent}.
\end{proof}

\noindent
Autrement dit, la cat\'egorie fibr\'ee des faisceaux quasi-coh\'erents est un champ
sur le site fpqc.
La d\'emonstration du th\'eor\`eme comporte deux \'etapes. La premi\`ere consiste \`a observer
que, pour les recouvrements de Zariski, le r\'esultat est facile (ou bien connu), par
le recollement usuel des faisceaux (voir
Faisceaux, section \ref{sheaves-section-glueing-sheaves})
et le caract\`ere local de la quasi-coh\'erence. La seconde \'etape traite le cas d'un
recouvrement fpqc de la forme $\{\Spec(B) \to \Spec(A)\}$,
o\`u $A \to B$ est un homomorphisme d'anneaux fid\`element plat.
Il s'agit d'un lemme d'alg\`ebre, que nous \'enon\c{c}ons maintenant.

\medskip\noindent
{\bf Descente des modules.}
Si $A \to B$ est un homomorphisme d'anneaux, consid\'erons le complexe
$$
(B/A)_\bullet : B \to B \otimes_A B \to B \otimes_A B \otimes_A B \to \ldots
$$
o\`u $B$ est en degr\'e 0, $B \otimes_A B$ en degr\'e 1, etc., les applications \'etant
donn\'ees par
\begin{eqnarray*}
b & \mapsto & 1 \otimes b - b \otimes 1, \\
b_0 \otimes b_1 & \mapsto & 1 \otimes b_0 \otimes b_1 - b_0 \otimes 1 \otimes
b_1 + b_0 \otimes b_1 \otimes 1, \\
& \text{etc.}
\end{eqnarray*}

\begin{lemma}
\label{lemma-algebra-descent}
Si $A \to B$ est fid\`element plat, alors le complexe $(B/A)_\bullet$ est exact en
degr\'es strictement positifs, et $H^0((B/A)_\bullet) = A$.
\end{lemma}

\begin{proof}
Voir Descente, lemme \ref{descent-lemma-ff-exact}.
\end{proof}

\noindent
Grothendieck d\'emontre ceci en trois \'etapes. Premi\`erement, il suppose que l'application $A
\to B$ admet une section, et construit une homotopie explicite avec le complexe dont
$A$ est l'unique terme non nul, en degr\'e 0. Deuxi\`emement, il observe que, pour d\'emontrer
le r\'esultat, il suffit de le faire apr\`es un changement de base fid\`element plat $A \to
A'$, en rempla\c{c}ant $B$ par $B' = B \otimes_A A'$. Troisi\`emement, il effectue le
changement de base fid\`element plat $A \to A' = B$ et remarque que l'application
$A' = B \to B' = B \otimes_A B$ poss\`ede une section naturelle.

\medskip\noindent
La m\^eme strat\'egie d\'emontre le lemme suivant.

\begin{lemma}
\label{lemma-descent-modules}
Si $A \to B$ est fid\`element plat et si $M$ est un $A$-module, alors le
complexe $(B/A)_\bullet \otimes_A M$ est exact en degr\'es strictement positifs, et
$H^0((B/A)_\bullet \otimes_A M) = M$.
\end{lemma}

\begin{proof}
Voir Descente, lemme \ref{descent-lemma-ff-exact}.
\end{proof}

\begin{definition}
\label{definition-descent-datum-modules}
Soient $A \to B$ un homomorphisme d'anneaux et $N$ un $B$-module. Une {\it donn\'ee de descente} pour
$N$ relativement \`a $A \to B$ est un isomorphisme
$\varphi : N \otimes_A B \cong B \otimes_A N$ de $B \otimes_A B$-modules tel
que le diagramme de $B \otimes_A B \otimes_A B$-modules
$$
\xymatrix{
{N \otimes_A B \otimes_A B} \ar[dr]_{\varphi_{02}} \ar[rr]^{\varphi_{01}} & &
{B \otimes_A N \otimes_A B} \ar[dl]^{\varphi_{12}} \\
& {B \otimes_A B \otimes_A N}
}
$$
commute, o\`u $\varphi_{01} = \varphi \otimes \text{id}_B$, et de m\^eme
pour $\varphi_{12}$ et $\varphi_{02}$.
\end{definition}

\noindent
Si $N' = B \otimes_A M$ pour un certain $A$-module M, alors il est muni d'une donn\'ee de descente
canonique donn\'ee par l'application
$$
\begin{matrix}
\varphi_\text{can}: & N' \otimes_A B & \to & B \otimes_A N' \\
& b_0 \otimes m \otimes b_1 & \mapsto & b_0 \otimes b_1 \otimes m.
\end{matrix}
$$

\begin{definition}
\label{definition-effective-modules}
Une donn\'ee de descente $(N, \varphi)$ est dite {\it effective} s'il existe un
$A$-module $M$ tel que $(N, \varphi) \cong (B \otimes_A M,
\varphi_\text{can})$, avec la notion \'evidente d'isomorphisme de donn\'ees de descente.
\end{definition}

\noindent
Le th\'eor\`eme \ref{theorem-descent-quasi-coherent} est une cons\'equence du
r\'esultat suivant.

\begin{theorem}
\label{theorem-descent-modules}
Si $A \to B$ est fid\`element plat, alors les donn\'ees de descente relativement \`a $A\to B$
sont effectives.
\end{theorem}

\begin{proof}
Voir
Descente, proposition \ref{descent-proposition-descent-module}.
Voir aussi
Descente, remarque \ref{descent-remark-homotopy-equivalent-cosimplicial-algebras},
pour un autre point de vue sur la d\'emonstration.
\end{proof}

\begin{remarks}
\label{remarks-theorem-modules-exactness}
Les r\'esultats sur la descente des modules ont plusieurs applications :
\begin{enumerate}
\item L'exactitude du complexe de {\v C}ech en degr\'es strictement positifs pour
le recouvrement $\{\Spec(B) \to \Spec(A)\}$, o\`u $A \to B$ est
fid\`element plat. Ceci donnera des r\'esultats d'annulation en cohomologie.
\item Si $(N, \varphi)$ est une donn\'ee de descente relativement \`a un homomorphisme
fid\`element plat $A \to B$, alors l'$A$-module correspondant est donn\'e par
$$
M = \Ker \left(
\begin{matrix}
N & \longrightarrow & B \otimes_A N \\
n & \longmapsto & 1 \otimes n - \varphi(n \otimes 1)
\end{matrix}
\right).
$$
Voir
Descente, proposition \ref{descent-proposition-descent-module}.
\end{enumerate}
\end{remarks}




%9.17.09
\section{Faisceaux quasi-coh\'erents}
\label{section-quasi-coherent}

\noindent
Nous pouvons appliquer la descente des modules \`a l'\'etude des faisceaux quasi-coh\'erents.

\begin{proposition}
\label{proposition-quasi-coherent-sheaf-fpqc}
Pour tout faisceau quasi-coh\'erent $\mathcal{F}$ sur $S$, le pr\'efaisceau
$$
\begin{matrix}
\mathcal{F}^a : & \Sch/S & \to & \textit{Ab}\\
& (f: T \to S) & \mapsto & \Gamma(T, f^*\mathcal{F})
\end{matrix}
$$
est un $\mathcal{O}$-module qui satisfait \`a la condition de faisceau pour la
topologie fpqc.
\end{proposition}

\begin{proof}
Ceci est d\'emontr\'e dans
Descente, lemme \ref{descent-lemma-sheaf-condition-holds}.
Indiquons ici la d\'emonstration. Comme \'etabli dans le
lemme \ref{lemma-fpqc-sheaves},
il suffit de v\'erifier la propri\'et\'e de faisceau
sur les recouvrements de Zariski et les morphismes fid\`element plats de sch\'emas affines. La
propri\'et\'e de faisceau pour les recouvrements de Zariski rel\`eve de la th\'eorie usuelle des sch\'emas, puisque
$\Gamma(U, i^\ast \mathcal{F}) = \mathcal{F}(U)$ lorsque
$i : U \hookrightarrow S$ est une immersion ouverte.

\medskip\noindent
Pour $\left\{\Spec(B)\to \Spec(A)\right\}$, o\`u $A\to B$ est fid\`element
plat et o\`u
$\mathcal{F}|_{\Spec(A)} = \widetilde{M}$,
cela correspond au fait que
$M = H^0\left((B/A)_\bullet \otimes_A M \right)$, c'est-\`a-dire que
\begin{align*}
0 \to M \to B \otimes_A M \to B \otimes_A B \otimes_A M
\end{align*}
est exacte d'apr\`es le
lemme \ref{lemma-descent-modules}.
\end{proof}

\noindent
Il existe une notion abstraite de faisceau quasi-coh\'erent sur un site annel\'e.
Nous l'introduisons bri\`evement ici. Pour plus d'informations, consulter
Modules sur les sites, section \ref{sites-modules-section-local}.
Soient $\mathcal{C}$ une cat\'egorie et $U$ un objet de $\mathcal{C}$.
Alors $\mathcal{C}/U$ d\'esigne la cat\'egorie des objets au-dessus de $U$, voir
Cat\'egories, exemple \ref{categories-example-category-over-X}.
Si $\mathcal{C}$ est un site, alors $\mathcal{C}/U$ est \'egalement un site :
les recouvrements de $V/U$ sont les familles $\{V_i/U \to V/U\}$ de morphismes
de $\mathcal{C}/U$ \`a cible fix\'ee telles que
$\{V_i \to V\}$ soit un recouvrement de $\mathcal{C}$. De plus, pour tout
faisceau $\mathcal{F}$ sur $\mathcal{C}$, la {\it restriction}
$\mathcal{F}|_{\mathcal{C}/U}$ (d\'efinie de mani\`ere \'evidente)
est encore un faisceau. Voir
Sites, section \ref{sites-section-localize},
pour plus de d\'etails.

\begin{definition}
\label{definition-ringed-site}
Soit $\mathcal{C}$ un {\it site annel\'e}, c'est-\`a-dire un site muni d'un
faisceau d'anneaux $\mathcal{O}$. Un faisceau de $\mathcal{O}$-modules $\mathcal{F}$ sur
$\mathcal{C}$ est dit {\it quasi-coh\'erent} si, pour tout
$U \in \Ob(\mathcal{C})$, il existe un recouvrement
$\{U_i \to U\}_{i\in I}$ de $\mathcal{C}$ tel que la restriction
$\mathcal{F}|_{\mathcal{C}/U_i}$ soit isomorphe au conoyau d'une
application $\mathcal{O}$-lin\'eaire entre $\mathcal{O}$-modules libres
$$
\bigoplus\nolimits_{k \in K} \mathcal{O}|_{\mathcal{C}/U_i}
\longrightarrow
\bigoplus\nolimits_{l \in L} \mathcal{O}|_{\mathcal{C}/U_i}.
$$
La somme directe index\'ee par $K$ est le faisceau associ\'e au pr\'efaisceau
$V \mapsto \bigoplus_{k \in K} \mathcal{O}(V)$, et de m\^eme pour l'autre.
\end{definition}

\noindent
Bien qu'il soit utile de pouvoir donner une d\'efinition g\'en\'erale comme ci-dessus,
cette notion ne se comporte pas bien en g\'en\'eral.

\begin{remark}
\label{remark-final-object}
Dans le cas o\`u $\mathcal{C}$ poss\`ede un objet final, par exemple $S$, il
suffit de v\'erifier la condition de la d\'efinition pour
$U = S$ dans l'\'enonc\'e ci-dessus. Voir
Modules sur les sites, lemme \ref{sites-modules-lemma-local-final-object}.
\end{remark}

\begin{theorem}[M\'eta-th\'eor\`eme sur les faisceaux quasi-coh\'erents]
\label{theorem-quasi-coherent}
Soit $S$ un sch\'ema.
Soit $\mathcal{C}$ un site. Supposons que
\begin{enumerate}
\item la cat\'egorie sous-jacente $\mathcal{C}$ soit une
sous-cat\'egorie pleine de $\Sch/S$,
\item tout recouvrement de Zariski de $T \in \Ob(\mathcal{C})$
puisse \^etre raffin\'e par un recouvrement de $\mathcal{C}$,
\item $S/S$ soit un objet de $\mathcal{C}$,
\item tout recouvrement de $\mathcal{C}$ soit un recouvrement fpqc de sch\'emas.
\end{enumerate}
Alors le pr\'efaisceau $\mathcal{O}$ est un faisceau sur $\mathcal{C}$ et
tout $\mathcal{O}$-module quasi-coh\'erent sur $(\mathcal{C}, \mathcal{O})$
est de la forme $\mathcal{F}^a$ pour un certain faisceau quasi-coh\'erent
$\mathcal{F}$ sur $S$.
\end{theorem}

\begin{proof}
Apr\`es quelques arguments formels, c'est exactement le th\'eor\`eme
\ref{theorem-descent-quasi-coherent}. Nous omettons les d\'etails. Dans
Descente, Proposition \ref{descent-proposition-equivalence-quasi-coherent},
nous d\'emontrons une version plus pr\'ecise du th\'eor\`eme pour les
grands sites de Zariski, fppf, \'etale, lisse et syntomique de $S$,
ainsi que pour les petits sites de Zariski et \'etale de $S$.
\end{proof}

\noindent
Autrement dit, il n'y a aucune diff\'erence entre les modules quasi-coh\'erents
sur le sch\'ema $S$ et les $\mathcal{O}$-modules quasi-coh\'erents
sur les sites $\mathcal{C}$ comme dans le th\'eor\`eme. On trouvera des \'enonc\'es plus pr\'ecis
pour les grands et petits sites $(\Sch/S)_{fppf}$, $S_\etale$, etc.,
dans Descente, sections
\ref{descent-section-quasi-coherent-sheaves},
\ref{descent-section-quasi-coherent-cohomology}, et
\ref{descent-section-quasi-coherent-sheaves-bis}.
Dans ce chapitre, nous parlerons parfois d'un
``site comme dans le th\'eor\`eme \ref{theorem-quasi-coherent}''
afin d'\'enoncer commod\'ement des r\'esultats valables dans chacune de ces
situations.






\section{Cohomologie de {\v C}ech}
\label{section-cech-cohomology}

\noindent
Notre prochain objectif est d'utiliser la th\'eorie de la descente pour montrer que
$H^i(\mathcal{C}, \mathcal{F}^a) = H_{Zar}^i(S, \mathcal{F})$
pour tout faisceau quasi-coh\'erent $\mathcal{F}$ sur $S$, et
tout site $\mathcal{C}$ comme dans le th\'eor\`eme \ref{theorem-quasi-coherent}.
\`A cette fin, nous introduisons la cohomologie de {\v C}ech sur les sites.
Voir \cite{ArtinTopologies} et
Cohomologie sur les sites, sections \ref{sites-cohomology-section-cech},
\ref{sites-cohomology-section-cech-functor}
et \ref{sites-cohomology-section-cech-cohomology-cohomology}
pour plus de d\'etails.

\begin{definition}
\label{definition-cech-complex}
Soit $\mathcal{C}$ une cat\'egorie. Soit
$\mathcal{U} = \{U_i \to U\}_{i \in I}$ une famille de morphismes de
$\mathcal{C}$ \`a cible fix\'ee. Supposons que tous les produits fibr\'es
$U_{i_0} \times_U \ldots \times_U U_{i_p}$ existent dans $\mathcal{C}$.
Soit $\mathcal{F} \in \textit{PAb}(\mathcal{C})$ un
pr\'efaisceau ab\'elien. Nous d\'efinissons le {\it complexe de {\v C}ech}
$\check{\mathcal{C}}^\bullet(\mathcal{U}, \mathcal{F})$ par
$$
\prod_{i_0 \in I} \mathcal{F}(U_{i_0}) \to
\prod_{i_0, i_1 \in I} \mathcal{F}(U_{i_0} \times_U U_{i_1}) \to
\prod_{i_0, i_1, i_2 \in I}
\mathcal{F}(U_{i_0} \times_U U_{i_1} \times_U U_{i_2}) \to \ldots
$$
o\`u le premier terme est en degr\'e 0 et les applications sont usuelles.
Les {\it groupes de cohomologie de {\v C}ech} sont d\'efinis par
$$
\check{H}^p(\mathcal{U}, \mathcal{F}) =
H^p(\check{\mathcal{C}}^\bullet(\mathcal{U}, \mathcal{F})).
$$
\end{definition}

\noindent
Dans la d\'efinition ci-dessus, il est essentiel d'autoriser les indices
$(i_0, \ldots, i_p)$ \`a se r\'ep\'eter.

\begin{lemma}
\label{lemma-cech-presheaves}
Notations et hypoth\`eses comme dans la d\'efinition \ref{definition-cech-complex}.
Le foncteur $\check{\mathcal{C}}^\bullet(\mathcal{U}, -)$
est exact sur la cat\'egorie $\textit{PAb}(\mathcal{C})$.
\end{lemma}

\noindent
Autrement dit, si $0\to \mathcal{F}_1\to \mathcal{F}_2\to \mathcal{F}_3\to 0$
est une suite exacte courte de pr\'efaisceaux de groupes ab\'eliens, alors
$$
0 \to \check{\mathcal{C}}^\bullet\left(\mathcal{U}, \mathcal{F}_1\right)
\to\check{\mathcal{C}}^\bullet(\mathcal{U}, \mathcal{F}_2) \to
\check{\mathcal{C}}^\bullet(\mathcal{U}, \mathcal{F}_3)\to 0
$$
est une suite exacte courte de complexes.

\begin{proof}
Cela r\'esulte imm\'ediatement de la d\'efinition d'une suite exacte courte de
pr\'efaisceaux. En effet, puisque la cat\'egorie des pr\'efaisceaux ab\'eliens est la cat\'egorie des
foncteurs d'une certaine cat\'egorie \`a valeurs dans $\textit{Ab}$, c'est automatiquement une
cat\'egorie ab\'elienne : une suite $\mathcal{F}_1\to \mathcal{F}_2\to \mathcal{F}_3$
est exacte dans $\textit{PAb}$ si et seulement si, pour tout
$U \in \Ob(\mathcal{C})$, la suite
$\mathcal{F}_1(U) \to \mathcal{F}_2(U) \to \mathcal{F}_3(U)$ est exacte dans
$\textit{Ab}$. Ainsi, le complexe ci-dessus est simplement un produit de suites exactes courtes
en chaque degr\'e. Voir aussi
Cohomologie sur les sites, lemme \ref{sites-cohomology-lemma-cech-exact-presheaves}.
\end{proof}

\noindent
Cela montre que $\check{H}^\bullet(\mathcal{U}, -)$ est un $\delta$-foncteur.
Nous allons maintenant montrer qu'il s'agit d'un $\delta$-foncteur universel. Il nous faut donc
montrer que c'est un foncteur {\it effa\c{c}able}. Commen\c{c}ons par rappeler le
lemme de Yoneda.

\begin{lemma}[Lemme de Yoneda]
\label{lemma-yoneda-presheaf}
Pour tout pr\'efaisceau $\mathcal{F}$ sur une cat\'egorie $\mathcal{C}$
et tout $U \in \Ob(\mathcal{C})$, il existe un isomorphisme fonctoriel
$$
\Hom_{\textit{PSh}(\mathcal{C})}(h_U, \mathcal{F}) =
\mathcal{F}(U).
$$
\end{lemma}

\begin{proof}
Voir Catégories, lemme \ref{categories-lemma-yoneda}.
\end{proof}

\noindent
Pour un ensemble $E$, nous notons (dans cette section)
$\mathbf{Z}[E]$ le groupe ab\'elien libre engendr\'e par $E$. Sous forme de formule,
$\mathbf{Z}[E] = \bigoplus_{e \in E} \mathbf{Z}$, c'est-\`a-dire que $\mathbf{Z}[E]$ est
un $\mathbf{Z}$-module libre dont une base est form\'ee des \'el\'ements de $E$.
Avec cette notation, nous introduisons le pr\'efaisceau ab\'elien libre engendr\'e par un
pr\'efaisceau d'ensembles.

\begin{definition}
\label{definition-free-abelian-presheaf}
Soit $\mathcal{C}$ une cat\'egorie.
\'Etant donn\'e un pr\'efaisceau d'ensembles $\mathcal{G}$, nous d\'efinissons le
{\it pr\'efaisceau ab\'elien libre engendr\'e par $\mathcal{G}$},
not\'e $\mathbf{Z}_\mathcal{G}$, par la r\`egle
$$
\mathbf{Z}_\mathcal{G}(U)
=
\mathbf{Z}[\mathcal{G}(U)]
$$
pour $U \in \Ob(\mathcal{C})$,
les applications de restriction \'etant induites par celles de $\mathcal{G}$.
Dans le cas particulier o\`u $\mathcal{G} = h_U$, nous \'ecrivons simplement
$\mathbf{Z}_U = \mathbf{Z}_{h_U}$.
\end{definition}

\noindent
Le foncteur $\mathcal{G} \mapsto \mathbf{Z}_\mathcal{G}$ est adjoint \`a gauche du
foncteur d'oubli $\textit{PAb}(\mathcal{C}) \to \textit{PSh}(\mathcal{C})$.
Ainsi, pour tout pr\'efaisceau $\mathcal{F}$, il existe un isomorphisme canonique
$$
\Hom_{\textit{PAb}(\mathcal{C})}(\mathbf{Z}_U, \mathcal{F})
=
\Hom_{\textit{PSh}(\mathcal{C})}(h_U, \mathcal{F})
=
\mathcal{F}(U)
$$
o\`u la derni\`ere \'egalit\'e r\'esulte du lemme de Yoneda. En particulier, nous obtenons le
r\'esultat suivant.

\begin{lemma}
\label{lemma-cech-complex-describe}
Notations et hypoth\`eses comme dans la d\'efinition \ref{definition-cech-complex}.
Le complexe de {\v C}ech $\check{\mathcal{C}}^\bullet(\mathcal{U}, \mathcal{F})$
peut se d\'ecrire explicitement comme suit :
\begin{eqnarray*}
\check{\mathcal{C}}^\bullet(\mathcal{U}, \mathcal{F})
& = &
\left(
\prod_{i_0 \in I}
\Hom_{\textit{PAb}(\mathcal{C})}(\mathbf{Z}_{U_{i_0}}, \mathcal{F}) \to
\prod_{i_0, i_1 \in I}
\Hom_{\textit{PAb}(\mathcal{C})}(
\mathbf{Z}_{U_{i_0} \times_U U_{i_1}}, \mathcal{F}) \to \ldots
\right) \\
& = &
\Hom_{\textit{PAb}(\mathcal{C})}\left(
\left(
\bigoplus_{i_0 \in I} \mathbf{Z}_{U_{i_0}} \leftarrow
\bigoplus_{i_0, i_1 \in I} \mathbf{Z}_{U_{i_0} \times_U U_{i_1}} \leftarrow
\ldots
\right), \mathcal{F}\right)
\end{eqnarray*}
\end{lemma}

\begin{proof}
Cela r\'esulte de la formule ci-dessus. Voir
Cohomologie sur les sites, lemme \ref{sites-cohomology-lemma-cech-map-into}.
\end{proof}

\noindent
Nous sommes ainsi ramen\'es \`a l'\'etude du seul complexe situ\'e dans le premier argument du
dernier $\Hom$.

\begin{lemma}
\label{lemma-exact}
Notations et hypoth\`eses comme dans la d\'efinition \ref{definition-cech-complex}.
Le complexe de pr\'efaisceaux ab\'eliens
\begin{align*}
\mathbf{Z}_\mathcal{U}^\bullet \quad : \quad
\bigoplus_{i_0 \in I} \mathbf{Z}_{U_{i_0}} \leftarrow
\bigoplus_{i_0, i_1 \in I} \mathbf{Z}_{U_{i_0} \times_U U_{i_1}} \leftarrow
\bigoplus_{i_0, i_1, i_2 \in I}
\mathbf{Z}_{U_{i_0} \times_U U_{i_1} \times_U U_{i_2}} \leftarrow
\ldots
\end{align*}
est exact en tout degr\'e sauf en degr\'e $0$ dans $\textit{PAb}(\mathcal{C})$.
\end{lemma}

\begin{proof}
Pour tout $V\in \Ob(\mathcal{C})$, le complexe de groupes ab\'eliens
$\mathbf{Z}_\mathcal{U}^\bullet(V)$ est
$$
\begin{matrix}
\mathbf{Z}\left[
\coprod_{i_0\in I} \Mor_\mathcal{C}(V, U_{i_0})\right]
\leftarrow
\mathbf{Z}\left[
\coprod_{i_0, i_1 \in I}
\Mor_\mathcal{C}(V, U_{i_0} \times_U U_{i_1})\right]
\leftarrow \ldots = \\
\bigoplus_{\varphi : V \to U}
\left(
\mathbf{Z}\left[\coprod_{i_0 \in I} \Mor_\varphi(V, U_{i_0})\right]
\leftarrow
\mathbf{Z}\left[\coprod_{i_0, i_1\in I} \Mor_\varphi(V, U_{i_0}) \times
\Mor_\varphi(V, U_{i_1})\right]
\leftarrow
\ldots
\right)
\end{matrix}
$$
o\`u
$$
\Mor_{\varphi}(V, U_i)
=
\{ V \to U_i \text{ tels que } V \to U_i \to U \text{ est \'egal \`a } \varphi \}.
$$
Posons $S_\varphi = \coprod_{i\in I} \Mor_\varphi(V, U_i)$, de sorte que
$$
\mathbf{Z}_\mathcal{U}^\bullet(V)
=
\bigoplus_{\varphi : V \to U}
\left(
\mathbf{Z}[S_\varphi] \leftarrow
\mathbf{Z}[S_\varphi \times S_\varphi] \leftarrow
\mathbf{Z}[S_\varphi \times S_\varphi \times S_\varphi] \leftarrow
\ldots \right).
$$
Il suffit donc de montrer que, pour chaque $S = S_\varphi$, le complexe
\begin{align*}
\mathbf{Z}[S] \leftarrow
\mathbf{Z}[S \times S] \leftarrow
\mathbf{Z}[S \times S \times S] \leftarrow \ldots
\end{align*}
est exact en degr\'es n\'egatifs. Pour le voir, nous pouvons donner une homotopie explicite.
Fixons $s\in S$ et d\'efinissons $K: n_{(s_0, \ldots, s_p)} \mapsto n_{(s, s_0,
\ldots, s_p)}.$ On v\'erifie ais\'ement que $K$ est une homotopie \`a z\'ero pour l'op\'erateur
$$
\delta :
\eta_{(s_0, \ldots, s_p)}
\mapsto
\sum\nolimits_{i = 0}^p (-1)^p \eta_{(s_0, \ldots, \hat s_i, \ldots, s_p)}.
$$
Voir
Cohomologie sur les sites, lemme \ref{sites-cohomology-lemma-homology-complex},
pour plus de d\'etails.
\end{proof}

\begin{lemma}
\label{lemma-hom-injective}
Notations et hypoth\`eses comme dans la d\'efinition \ref{definition-cech-complex}.
Si $\mathcal{I}$ est un objet injectif de $\textit{PAb}(\mathcal{C})$,
alors $\check H^p(\mathcal{U}, \mathcal{I}) = 0$ pour tout $p > 0$.
\end{lemma}

\begin{proof}
Le complexe de {\v C}ech s'obtient en appliquant le foncteur
$\Hom_{\textit{PAb}(\mathcal{C})}(-, \mathcal{I}) $ au complexe $
\mathbf{Z}^\bullet_\mathcal{U} $, c'est-\`a-dire
$$
\check H^p(\mathcal{U}, \mathcal{I}) = H^p
(\Hom_{\textit{PAb}(\mathcal{C})} (\mathbf{Z}^\bullet_\mathcal{U},
\mathcal{I})).
$$
Mais nous venons de voir que $\mathbf{Z}^\bullet_\mathcal{U}$ est exact en
degr\'es n\'egatifs, et le foncteur $\Hom_{\textit{PAb}(\mathcal{C})}(-,
\mathcal{I})$ est exact ; par cons\'equent, $\Hom_{\textit{PAb}(\mathcal{C})}
(\mathbf{Z}^\bullet_\mathcal{U}, \mathcal{I})$ est exact en degr\'es positifs.
\end{proof}

\begin{theorem}
\label{theorem-cech-derived}
Notations et hypoth\`eses comme dans la d\'efinition \ref{definition-cech-complex}.
Sur $\textit{PAb}(\mathcal{C})$, les foncteurs $\check{H}^p(\mathcal{U}, -)$ sont
les foncteurs d\'eriv\'es \`a droite de $\check{H}^0(\mathcal{U}, -)$.
\end{theorem}

\begin{proof}
D'apr\`es le lemme \ref{lemma-hom-injective}, les foncteurs
$\check H^p(\mathcal{U}, -)$ sont des
$\delta$-foncteurs universels, puisqu'ils sont effa\c{c}ables.
Il en va de m\^eme des foncteurs d\'eriv\'es \`a droite de $\check H^0(\mathcal{U}, -)$. Puisqu'ils
co\"incident en degr\'e $0$, ils co\"incident par la propri\'et\'e universelle des
$\delta$-foncteurs universels. Pour plus de d\'etails, voir
Cohomologie sur les sites,
lemme \ref{sites-cohomology-lemma-cech-cohomology-derived-presheaves}.
\end{proof}

\begin{remark}
\label{remark-presheaves-no-topology}
Notons que tous les \'enonc\'es pr\'ec\'edents portent sur les pr\'efaisceaux ; nous n'avons donc
pas encore utilis\'e la topologie.
\end{remark}




\section{La suite spectrale de {\v C}ech \`a la cohomologie}
\label{section-cech-ss}

\noindent
Cette suite spectrale est essentielle pour \`etablir les r\'esultats fondamentaux sur la
cohomologie des faisceaux.

\begin{lemma}
\label{lemma-forget-injectives}
Le foncteur d'oubli $\textit{Ab}(\mathcal{C})\to \textit{PAb}(\mathcal{C})$
transforme les objets injectifs en objets injectifs.
\end{lemma}

\begin{proof}
C'est formel, puisque le foncteur d'oubli admet un adjoint \`a gauche,
\`a savoir le passage au faisceau associ\'e, qui est un foncteur exact. Pour plus de d\'etails, voir
Cohomologie sur les sites,
lemme \ref{sites-cohomology-lemma-injective-abelian-sheaf-injective-presheaf}.
\end{proof}

\begin{theorem}
\label{theorem-cech-ss}
Soit $\mathcal{C}$ un site. Pour tout recouvrement
$\mathcal{U} = \{U_i \to U\}_{i \in I}$ de $U \in \Ob(\mathcal{C})$
et tout faisceau ab\'elien $\mathcal{F}$ sur $\mathcal{C}$,
il existe une suite spectrale
$$
E_2^{p, q}
=
\check H^p(\mathcal{U}, \underline{H}^q(\mathcal{F}))
\Rightarrow
H^{p+q}(U, \mathcal{F}),
$$
o\`u $\underline{H}^q(\mathcal{F})$ est le pr\'efaisceau ab\'elien
$V \mapsto H^q(V, \mathcal{F})$.
\end{theorem}

\begin{proof}
Choisissons une r\'esolution injective $\mathcal{F}\to \mathcal{I}^\bullet$ dans
$\textit{Ab}(\mathcal{C})$ et consid\'erons le bicomplexe
$\check{\mathcal{C}}^\bullet(\mathcal{U}, \mathcal{I}^\bullet)$
et les applications
$$
\xymatrix{
\Gamma(U, I^\bullet) \ar[r] &
\check{\mathcal{C}}^\bullet(\mathcal{U}, \mathcal{I}^\bullet) \\
& \check{\mathcal{C}}^\bullet(\mathcal{U}, \mathcal{F}) \ar[u]
}
$$
Ici, l'application horizontale est l'application naturelle
$\Gamma(U, I^\bullet) \to
\check{\mathcal{C}}^0(\mathcal{U}, \mathcal{I}^\bullet)$
vers la colonne de gauche, et l'application verticale est induite par
$\mathcal{F}\to \mathcal{I}^0$ et aboutit \`a la ligne inf\'erieure.
Par hypoth\`ese, $\mathcal{I}^\bullet$ est un complexe d'injectifs dans
$\textit{Ab}(\mathcal{C})$ ; par le
lemme \ref{lemma-forget-injectives}, c'est donc un complexe d'injectifs dans
$\textit{PAb}(\mathcal{C})$. Ainsi, les lignes du bicomplexe sont
exactes en degr\'es positifs (lemme \ref{lemma-hom-injective}), et
le noyau de $\check{\mathcal{C}}^0(\mathcal{U}, \mathcal{I}^\bullet)
\to \check{\mathcal{C}}^1(\mathcal{U}, \mathcal{I}^\bullet)$
est \'egal \`a
$\Gamma(U, \mathcal{I}^\bullet)$, puisque $\mathcal{I}^\bullet$
est un complexe de faisceaux. En particulier, la cohomologie du complexe total
est la cohomologie usuelle
du foncteur des sections globales $H^0(U, \mathcal{F})$.

\medskip\noindent
Dans la direction verticale, le $q$-i\`eme groupe de cohomologie de la $p$-i\`eme colonne est
$$
\prod_{i_0, \ldots, i_p}
H^q(U_{i_0} \times_U \ldots \times_U U_{i_p}, \mathcal{F})
=
\prod_{i_0, \ldots, i_p}
\underline{H}^q(\mathcal{F})(U_{i_0} \times_U \ldots \times_U U_{i_p})
$$
dans le terme $E_1^{p, q}$. Il s'agit donc d'une suite spectrale standard de bicomplexe,
et sa page $E_2$ est celle indiqu\'ee. Pour plus de d\'etails, voir
Cohomologie sur les sites,
lemme \ref{sites-cohomology-lemma-cech-spectral-sequence}.
\end{proof}

\begin{remark}
\label{remark-grothendieck-ss}
C'est une suite spectrale de Grothendieck pour la compos\'ee des foncteurs
$$
\textit{Ab}(\mathcal{C}) \longrightarrow
\textit{PAb}(\mathcal{C}) \xrightarrow{\check H^0} \textit{Ab}.
$$
\end{remark}








\section{Grands et petits sites des sch\'emas}
\label{section-big-small}

\noindent
Soit $S$ un sch\'ema.
Soit $\tau$ l'une des topologies dont nous allons parler.
Ainsi, $\tau \in \{fppf, syntomic, smooth, \etale, Zariski\}$.
Bien entendu, si seule la topologie \'etale vous int\'eresse, vous pouvez
simplement supposer $\tau = \etale$ partout. En outre, nous traiterons
plus en d\'etail des morphismes \'etales, des recouvrements \'etales et des sites \'etales
\`a partir de la section \ref{section-etale-site}.
Pour poursuivre l'\'etude de la cohomologie des
faisceaux quasi-coh\'erents, il est commode d'introduire le
grand $\tau$-site et, lorsque $\tau \in \{\etale, Zariski\}$, le
petit $\tau$-site de $S$. Pour ce faire, nous introduisons d'abord
la notion de $\tau$-recouvrement.

\begin{definition}
\label{definition-tau-covering}
(Voir
Topologies, définitions
\ref{topologies-definition-fppf-covering},
\ref{topologies-definition-syntomic-covering},
\ref{topologies-definition-smooth-covering},
\ref{topologies-definition-etale-covering}, et
\ref{topologies-definition-zariski-covering}.)
Soit $\tau \in \{fppf, syntomic, smooth, \etale, Zariski\}$.
Une famille de morphismes de sch\'emas $\{f_i : T_i \to T\}_{i \in I}$ \`a cible
fix\'ee est appel\'ee un {\it $\tau$-recouvrement} si et seulement si
chaque $f_i$ est plat et de pr\'esentation finie, syntomique, lisse, \'etale,
respectivement une immersion ouverte, et si $\bigcup f_i(T_i) = T$.
\end{definition}

\noindent
La classe de tous les $\tau$-recouvrements satisfait aux axiomes
(1), (2) et (3) de la
d\'efinition \ref{definition-site} (notre d\'efinition d'un site), voir
Topologies, lemmes
\ref{topologies-lemma-fppf},
\ref{topologies-lemma-syntomic},
\ref{topologies-lemma-smooth},
\ref{topologies-lemma-etale}, et
\ref{topologies-lemma-zariski}.

\medskip\noindent
Introduisons les sites avec lesquels nous travaillerons.
Contrairement \`a ce qui se passe dans
\cite{SGA4}, nous ne voulons pas choisir un univers. Nous choisissons plut\^ot un ``univers
partiel'' (qui est un ensemble suffisamment grand comme dans
Ensembles, section \ref{sets-section-sets-hierarchy}) et consid\'erons tous les sch\'emas
contenus dans cet ensemble. Bien entendu, nous veillons \`a ce que notre sch\'ema de base privil\'egi\'e
$S$ appartienne \`a l'univers partiel. Apr\`es avoir choisi la cat\'egorie sous-jacente,
nous choisissons un ensemble suffisamment grand de $\tau$-recouvrements qui en fait un site.
Les d\'etails se trouvent dans le chapitre sur les topologies des sch\'emas ; les choix
laissent une grande libert\'e, mais, en fin de compte, les choix effectivement faits n'auront
aucune incidence sur la cohomologie \'etale (ou autre) de $S$ (tout comme, dans \cite{SGA4},
le choix effectif de l'univers n'a finalement aucune importance). En outre, la mani\`ere dont le
texte est r\'edig\'e permet au lecteur qui accepte les cardinaux fortement inaccessibles
(c'est-\`a-dire les univers) de les utiliser en remplacement.

\begin{definition}
\label{definition-tau-site}
Soit $S$ un sch\'ema.
Soit $\tau \in \{fppf, syntomic, smooth, \etale, \linebreak[0] Zariski\}$.
\begin{enumerate}
\item Un {\it grand $\tau$-site de $S$} est l'un quelconque des sites
$(\Sch/S)_\tau$ construits comme expliqu\'e ci-dessus et, plus en d\'etail, dans
Topologies, définitions
\ref{topologies-definition-big-small-fppf},
\ref{topologies-definition-big-small-syntomic},
\ref{topologies-definition-big-small-smooth},
\ref{topologies-definition-big-small-etale}, et
\ref{topologies-definition-big-small-Zariski}.
\item Si $\tau \in \{\etale, Zariski\}$, alors le
{\it petit $\tau$-site de $S$}
est la sous-cat\'egorie pleine $S_\tau$ de $(\Sch/S)_\tau$ dont les objets
sont les sch\'emas $T$ sur $S$ dont le morphisme structural $T \to S$ est \'etale,
respectivement une immersion ouverte. Un recouvrement dans $S_\tau$ est un recouvrement
$\{U_i \to U\}$ dans $(\Sch/S)_\tau$
tel que $U$ soit un objet de $S_\tau$.
\end{enumerate}
\end{definition}

\noindent
La cat\'egorie sous-jacente au site $(\Sch/S)_\tau$ poss\`ede des propri\'et\'es
raisonnables de ``stabilit\'e'' : \'etant donn\'e un sch\'ema $T$ qui lui appartient, tout sous-sch\'ema
localement ferm\'e de $T$ est isomorphe \`a un objet de $(\Sch/S)_\tau$.
Parmi les autres propri\'et\'es de ce type : stabilit\'e par produits fibr\'es de sch\'emas,
par sommes disjointes d\'enombrables,
par passage aux sch\'emas de type fini sur un sch\'ema donn\'e ; \'etant donn\'e un sch\'ema affine
$\Spec(R)$, on peut compl\'eter ou localiser $R$, ou en prendre le quotient
par un id\'eal, tout en restant dans la cat\'egorie, etc.
En revanche, par exemple, des sommes disjointes arbitraires
de sch\'emas de $(\Sch/S)_\tau$ feront sortir de cette cat\'egorie.
Notons aussi que, pour un objet $T$ de $(\Sch/S)_\tau$, il existe des
$\tau$-recouvrements $\{T_i \to T\}_{i \in I}$ (comme dans la
d\'efinition \ref{definition-tau-covering})
qui ne sont pas des recouvrements dans $(\Sch/S)_\tau$, par exemple
parce que les sch\'emas $T_i$ ne sont pas des objets de la cat\'egorie
$(\Sch/S)_\tau$. Mais notre choix des sites $(\Sch/S)_\tau$
est tel qu'il existe toujours
un recouvrement $\{U_j \to T\}_{j \in J}$ de $(\Sch/S)_\tau$ qui raffine
le recouvrement $\{T_i \to T\}_{i \in I}$, voir
Topologies, lemmes
\ref{topologies-lemma-fppf-induced},
\ref{topologies-lemma-syntomic-induced},
\ref{topologies-lemma-smooth-induced},
\ref{topologies-lemma-etale-induced}, et
\ref{topologies-lemma-zariski-induced}.
Dans ce chapitre, nous ignorerons pour l'essentiel ces questions.

\medskip\noindent
Si $\mathcal{F}$ est un faisceau sur $(\Sch/S)_\tau$ ou $S_\tau$, alors
nous notons
$$
H^p_\tau(U, \mathcal{F}), \text{ en particulier }
H^p_\tau(S, \mathcal{F})
$$
les groupes de cohomologie de $\mathcal{F}$ sur l'objet $U$ du site, voir
la section \ref{section-cohomology}. Nous avons ainsi
$H^p_{fppf}(S, \mathcal{F})$,
$H^p_{syntomic}(S, \mathcal{F})$,
$H^p_{smooth}(S, \mathcal{F})$,
$H^p_\etale(S, \mathcal{F})$, et
$H^p_{Zar}(S, \mathcal{F})$. Les deux derniers sont potentiellement ambigus, puisqu'ils
peuvent renvoyer soit au grand, soit au petit site \'etale ou de Zariski. Cependant,
cette ambigu\"it\'e est sans cons\'equence gr\^ace au lemme suivant.

\begin{lemma}
\label{lemma-compare-cohomology-big-small}
Soit $\tau \in \{\etale, Zariski\}$.
Si $\mathcal{F}$ est un faisceau ab\'elien d\'efini sur
$(\Sch/S)_\tau$, alors
les groupes de cohomologie de $\mathcal{F}$ sur $S$ co\"incident avec les groupes de cohomologie
de $\mathcal{F}|_{S_\tau}$ sur $S$.
\end{lemma}

\begin{proof}
D'apr\`es
Topologies, lemmes \ref{topologies-lemma-at-the-bottom} et
\ref{topologies-lemma-at-the-bottom-etale},
les foncteurs $S_\tau \to (\Sch/S)_\tau$
satisfont aux hypoth\`eses de
Sites, lemme \ref{sites-lemma-bigger-site}.
Notre lemme r\'esulte donc de
Cohomologie sur les sites, lemme \ref{sites-cohomology-lemma-cohomology-bigger-site}.
\end{proof}

\noindent
La cat\'egorie des faisceaux sur le grand ou le petit site \'etale de $S$ ne d\'epend que
de la sous-cat\'egorie pleine de $(\Sch/S)_\etale$ ou de $S_\etale$ form\'ee des
sch\'emas affines, et il suffit de consid\'erer entre eux les recouvrements \'etales standards
(d\'efinis ci-dessous). Cela donne les sites
$(\textit{Aff}/S)_\etale$ et $S_{affine, \etale}$, voir
Topologies, définition \ref{topologies-definition-big-small-etale}.
Les r\'esultats de comparaison sont d\'emontr\'es dans Topologies,
lemmes \ref{topologies-lemma-affine-big-site-etale} et
\ref{topologies-lemma-alternative}. Voici notre d\'efinition des
recouvrements standards pour certaines des topologies consid\'er\'ees dans ce chapitre.

\begin{definition}
\label{definition-standard-tau}
(Voir
Topologies, définitions
\ref{topologies-definition-standard-fppf},
\ref{topologies-definition-standard-syntomic},
\ref{topologies-definition-standard-smooth},
\ref{topologies-definition-standard-etale}, et
\ref{topologies-definition-standard-Zariski}.)
Soit $\tau \in \{fppf, syntomic, smooth, \etale, Zariski\}$.
Soit $T$ un sch\'ema affine.
Un {\it $\tau$-recouvrement standard} de $T$ est une famille
$\{f_j : U_j \to T\}_{j = 1, \ldots, m}$ o\`u chaque $U_j$ est affine,
et chaque $f_j$ est plat et de pr\'esentation finie,
syntomique standard, lisse standard, \'etale, respectivement l'immersion d'un
ouvert principal standard dans $T$, et $T = \bigcup f_j(U_j)$.
\end{definition}

\begin{lemma}
\label{lemma-tau-affine}
Soit $\tau \in \{fppf, syntomic, smooth, \etale, Zariski\}$.
Tout $\tau$-recouvrement d'un sch\'ema affine peut \^etre raffin\'e par un
$\tau$-recouvrement standard.
\end{lemma}

\begin{proof}
Voir
Topologies, lemmes
\ref{topologies-lemma-fppf-affine},
\ref{topologies-lemma-syntomic-affine},
\ref{topologies-lemma-smooth-affine},
\ref{topologies-lemma-etale-affine}, et
\ref{topologies-lemma-zariski-affine}.
\end{proof}

\noindent
Par souci d'exhaustivit\'e, nous \'enon\c{c}ons et d\'emontrons l'invariance, par rapport au choix de l'univers
partiel, des groupes de cohomologie que nous consid\'erons. Nous d\'emontrerons l'invariance
du petit topos \'etale dans le
lemme \ref{lemma-etale-topos-independent-partial-universe} ci-dessous.
Pour les notations et la terminologie employ\'ees dans ce lemme, nous renvoyons \`a
Topologies, section \ref{topologies-section-change-alpha}.

\begin{lemma}
\label{lemma-cohomology-enlarge-partial-universe}
Soit $\tau \in \{fppf, syntomic, smooth, \etale, Zariski\}$.
Soit $S$ un sch\'ema.
Soient $(\Sch/S)_\tau$ et $(\Sch'/S)_\tau$ deux
grands $\tau$-sites de $S$, et supposons que le premier soit contenu dans le second.
Dans ce cas,
\begin{enumerate}
\item pour tout faisceau ab\'elien $\mathcal{F}'$ d\'efini sur $(\Sch'/S)_\tau$ et
tout objet $U$ de $(\Sch/S)_\tau$, nous avons
$$
H^p_\tau(U, \mathcal{F}'|_{(\Sch/S)_\tau}) =
H^p_\tau(U, \mathcal{F}')
$$
Autrement dit : la cohomologie de $\mathcal{F}'$ sur $U$, calcul\'ee dans le site le plus grand,
co\"incide avec celle de la restriction de $\mathcal{F}'$ au site le plus petit,
sur $U$.
\item pour tout faisceau ab\'elien $\mathcal{F}$ sur $(\Sch/S)_\tau$, il existe un
faisceau ab\'elien $\mathcal{F}'$ sur $(\Sch/S)_\tau'$ dont la restriction \`a
$(\Sch/S)_\tau$ est isomorphe \`a $\mathcal{F}$.
\end{enumerate}
\end{lemma}

\begin{proof}
D'apr\`es Topologies, lemme \ref{topologies-lemma-change-alpha}, le foncteur d'inclusion
$(\Sch/S)_\tau \to (\Sch'/S)_\tau$ satisfait aux hypoth\`eses de
Sites, lemme \ref{sites-lemma-bigger-site}. Cela implique (2), et (1)
r\'esulte de
Cohomologie sur les sites, lemme \ref{sites-cohomology-lemma-cohomology-bigger-site}.
\end{proof}




\section{Le topos \'etale}
\label{section-etale-topos}

\noindent
Un {\it topos} est la cat\'egorie des faisceaux d'ensembles sur un site, voir
Sites, définition \ref{sites-definition-topos}. Il est donc d'usage
d'employer l'expression ``topos \'etale d'un sch\'ema'' pour d\'esigner
la cat\'egorie des faisceaux sur le petit site \'etale d'un sch\'ema.
En voici la d\'efinition formelle.

\begin{definition}
\label{definition-etale-topos}
Soit $S$ un sch\'ema.
\begin{enumerate}
\item Le {\it topos \'etale}, ou le {\it petit topos \'etale}
de $S$, est la cat\'egorie $\Sh(S_\etale)$ des faisceaux d'ensembles sur
le petit site \'etale de $S$.
\item Le {\it topos de Zariski}, ou le {\it petit topos de Zariski}
de $S$, est la cat\'egorie $\Sh(S_{Zar})$ des faisceaux d'ensembles sur le
petit site de Zariski de $S$.
\item Pour $\tau \in \{fppf, syntomic, smooth, \etale, Zariski\}$, un
{\it grand $\tau$-topos} est la cat\'egorie des faisceaux d'ensembles sur un
grand $\tau$-topos de $S$.
\end{enumerate}
\end{definition}

\noindent
Notons que le petit topos de Zariski de $S$ est simplement la cat\'egorie des faisceaux
d'ensembles sur l'espace topologique sous-jacent \`a $S$, voir
Topologies, lemme \ref{topologies-lemma-Zariski-usual}.
Alors que le petit topos \'etale ne d\'epend pas des choix faits lors de la
construction du petit site \'etale, les grands topos d\'ependent en g\'en\'eral
de ces choix.

\medskip\noindent
Il se trouve que le grand ou le petit topos \'etale ne d\'epend que de la
sous-cat\'egorie pleine de $(\Sch/S)_\etale$ ou de $S_\etale$ form\'ee des
sch\'emas affines, voir
Topologies, lemmes \ref{topologies-lemma-affine-big-site-etale} et
\ref{topologies-lemma-alternative}.
Nous l'utiliserons par exemple dans la d\'emonstration du lemme suivant.

\begin{lemma}
\label{lemma-etale-topos-independent-partial-universe}
Soit $S$ un sch\'ema. Le topos \'etale de $S$ est ind\'ependant
(\`a \'equivalence canonique pr\`es) de la construction du petit
site \'etale de la d\'efinition \ref{definition-tau-site}.
\end{lemma}

\begin{proof}
Nous devons montrer qu'\'etant donn\'es deux grands sites \'etales
$\Sch_\etale$ et $\Sch_\etale'$ contenant
$S$, on a $\Sh(S_\etale) \cong \Sh(S_\etale')$
avec les notations \'evidentes. D'apr\`es Topologies, lemme \ref{topologies-lemma-contained-in},
nous pouvons supposer $\Sch_\etale \subset \Sch_\etale'$.
D'apr\`es Ensembles, lemme \ref{sets-lemma-what-is-in-it},
tout sch\'ema affine \'etale sur $S$ est isomorphe \`a un objet
\`a la fois de $\Sch_\etale$ et de $\Sch_\etale'$.
Ainsi, le foncteur induit
$S_{affine, \etale} \to S_{affine, \etale}'$
est une \'equivalence. En outre, il est clair que ce foncteur
et un quasi-inverse transforment les recouvrements \'etales standards en
recouvrements \'etales standards.
Le r\'esultat d\'ecoule donc de
Topologies, lemme \ref{topologies-lemma-alternative}.
\end{proof}





\section{Cohomologie des faisceaux quasi-coh\'erents}
\label{section-cohomology-quasi-coherent}
%9.22.09

\noindent
Commen\c{c}ons par un lemme simple (qui reste vrai dans une g\'en\'eralit\'e plus grande que
celle de l'\'enonc\'e). Il affirme que le complexe de {\v C}ech d'un recouvrement standard est
\'egal au complexe de {\v C}ech d'un recouvrement fpqc de la forme
$\{\Spec(B) \to \Spec(A)\}$ o\`u $A \to B$ est fid\`element plat.

\begin{lemma}
\label{lemma-cech-complex}
Soit $\tau \in \{fppf, syntomic, smooth, \etale, Zariski\}$.
Soit $S$ un sch\'ema.
Soit $\mathcal{F}$ un faisceau ab\'elien sur $(\Sch/S)_\tau$, ou sur
$S_\tau$ dans le cas o\`u $\tau = \etale$, et soit
$\mathcal{U} = \{U_i \to U\}_{i \in I}$
un $\tau$-recouvrement standard de ce site.
Soit $V = \coprod_{i \in I} U_i$. Alors
\begin{enumerate}
\item $V$ est un sch\'ema affine,
\item $\mathcal{V} = \{V \to U\}$ est un recouvrement fpqc
et aussi un $\tau$-recouvrement, sauf si $\tau = Zariski$,
\item les complexes de {\v C}ech
$\check{\mathcal{C}}^\bullet (\mathcal{U}, \mathcal{F})$ et
$\check{\mathcal{C}}^\bullet (\mathcal{V}, \mathcal{F})$ co\"incident.
\end{enumerate}
\end{lemma}

\begin{proof}
La d\'efinition d'un $\tau$-recouvrement standard est donn\'ee dans
Topologies, définition
\ref{topologies-definition-standard-Zariski},
\ref{topologies-definition-standard-etale},
\ref{topologies-definition-standard-smooth},
\ref{topologies-definition-standard-syntomic}, et
\ref{topologies-definition-standard-fppf}.
Par d\'efinition, chacun des sch\'emas
$U_i$ est affine et $I$ est un ensemble fini. Ainsi, $V$ est un sch\'ema affine.
Il est clair que $V \to U$ est plat et surjectif ; par cons\'equent,
$\mathcal{V}$ est un recouvrement fpqc, voir
l'exemple \ref{example-fpqc-coverings}.
Sauf dans le cas de Zariski, le recouvrement $\mathcal{V}$
est aussi un $\tau$-recouvrement, voir
Topologies, définition
\ref{topologies-definition-etale-covering},
\ref{topologies-definition-smooth-covering},
\ref{topologies-definition-syntomic-covering}, et
\ref{topologies-definition-fppf-covering}.

\medskip\noindent
Notons que $\mathcal{U}$ est un raffinement de $\mathcal{V}$ ;
il existe donc un morphisme de complexes de {\v C}ech
$\check{\mathcal{C}}^\bullet (\mathcal{V}, \mathcal{F}) \to
\check{\mathcal{C}}^\bullet (\mathcal{U}, \mathcal{F})$, voir
Cohomologie sur les sites,
Equation (\ref{sites-cohomology-equation-map-cech-complexes}).
Ensuite, observons que si $T = \coprod_{j \in J} T_j$ est une
somme disjointe de sch\'emas du site sur lequel $\mathcal{F}$ est d\'efini,
alors la famille de morphismes \`a cible fix\'ee
$\{T_j \to T\}_{j \in J}$ est un recouvrement de Zariski, et donc
\begin{equation}
\label{equation-sheaf-coprod}
\mathcal{F}(T) =
\mathcal{F}(\coprod\nolimits_{j \in J} T_j) =
\prod\nolimits_{j \in J} \mathcal{F}(T_j)
\end{equation}
d'apr\`es la condition de faisceau pour $\mathcal{F}$.
Cela implique que le morphisme de complexes de {\v C}ech ci-dessus est un isomorphisme
en chaque degr\'e, car
$$
V \times_U \ldots \times_U V
=
\coprod\nolimits_{i_0, \ldots i_p} U_{i_0} \times_U \ldots \times_U U_{i_p}
$$
en tant que sch\'emas.
\end{proof}

\noindent
Notons que l'\'egalit\'e (\ref{equation-sheaf-coprod})
est fausse pour un pr\'efaisceau g\'en\'eral. M\^eme pour les faisceaux, elle n'est pas vraie sur tout
site, car les coproduits peuvent ne pas donner des recouvrements et ne pas \^etre disjoints.
Mais elle est vraie pour tous les sites usuels (du moins pour tous ceux que nous \'etudierons).

\begin{remark}
\label{remark-refinement}
Dans l'\'enonc\'e du lemme \ref{lemma-cech-complex}, le recouvrement $\mathcal{U}$
est un raffinement de $\mathcal{V}$, mais non l'inverse. Les recouvrements
de la forme $\{V \to U\}$ ne forment pas une sous-cat\'egorie initiale de la
cat\'egorie de tous les recouvrements de $U$. Il reste pourtant vrai que
nous pouvons calculer la cohomologie de {\v C}ech $\check H^n(U, \mathcal{F})$ (qui
est d\'efinie comme la limite inductive index\'ee par la cat\'egorie oppos\'ee \`a celle des
recouvrements $\mathcal{U}$ de $U$ des groupes de cohomologie de {\v C}ech de
$\mathcal{F}$ relativement \`a $\mathcal{U}$) au moyen des recouvrements
$\{V \to U\}$. Nous formulerons un lemme pr\'ecis (il ne vaut que pour les faisceaux)
et l'ajouterons ici si nous en avons un jour besoin.
\end{remark}

\begin{lemma}[Localit\'e de la cohomologie]
\label{lemma-locality-cohomology}
Soient $\mathcal{C}$ un site, $\mathcal{F}$ un faisceau ab\'elien sur $\mathcal{C}$,
$U$ un objet de $\mathcal{C}$, $p > 0$ un entier et $\xi \in
H^p(U, \mathcal{F})$. Alors il existe un recouvrement
$\mathcal{U} = \{U_i \to U\}_{i \in I}$ de $U$ dans $\mathcal{C}$
tel que $\xi |_{U_i} = 0$ pour tout $i \in I$.
\end{lemma}

\begin{proof}
Choisissons une r\'esolution injective $\mathcal{F} \to \mathcal{I}^\bullet$. Alors
$\xi$ est repr\'esent\'ee par un cocycle $\tilde{\xi} \in \mathcal{I}^p(U)$
tel que $d^p(\tilde{\xi}) = 0$. Par hypoth\`ese, la suite
$\mathcal{I}^{p - 1} \to \mathcal{I}^p \to \mathcal{I}^{p + 1}$ est exacte dans
$\textit{Ab}(\mathcal{C})$, ce qui signifie qu'il existe un recouvrement
$\mathcal{U} = \{U_i \to U\}_{i \in I}$ tel que
$\tilde{\xi}|_{U_i} = d^{p - 1}(\xi_i)$ pour un certain
$\xi_i \in \mathcal{I}^{p-1}(U_i)$. Puisque
la classe de cohomologie $\xi|_{U_i}$ est repr\'esent\'ee par le cocycle
$\tilde{\xi}|_{U_i}$, qui est un cobord, elle s'annule.
Pour plus de d\'etails, voir
Cohomologie sur les sites,
lemme \ref{sites-cohomology-lemma-kill-cohomology-class-on-covering}.
\end{proof}

\begin{theorem}
\label{theorem-zariski-fpqc-quasi-coherent}
Soit $S$ un sch\'ema et soit $\mathcal{F}$ un $\mathcal{O}_S$-module quasi-coh\'erent.
Soit $\mathcal{C}$ l'un des sites $(\Sch/S)_\tau$ pour
$\tau \in \{fppf, syntomic, smooth, \etale, Zariski\}$, ou
$S_\etale$. Alors
$$
H^p(S, \mathcal{F}) = H^p_\tau(S, \mathcal{F}^a)
$$
pour tout $p \geq 0$, o\`u
\begin{enumerate}
\item le membre de gauche d\'esigne la cohomologie usuelle du faisceau
$\mathcal{F}$ sur l'espace topologique sous-jacent au sch\'ema $S$, et
\item le membre de droite d\'esigne la cohomologie
du faisceau ab\'elien $\mathcal{F}^a$ (voir
la proposition \ref{proposition-quasi-coherent-sheaf-fpqc})
sur le site $\mathcal{C}$.
\end{enumerate}
\end{theorem}

\begin{proof}
Nous allons montrer que
$H^p(U, f^*\mathcal{F}) = H^p_\tau(U, \mathcal{F}^a)$
pour tout objet $f : U \to S$ du site $\mathcal{C}$.
Le r\'esultat est vrai pour $p = 0$ par la propri\'et\'e de faisceau.

\medskip\noindent
Supposons que $U$ soit affine. Nous voulons alors prouver que
$H^p_\tau(U, \mathcal{F}^a) = 0$ pour tout $p > 0$. Nous raisonnons par r\'ecurrence sur $p$.
\begin{enumerate}
\item[p = 1]
Prenons $\xi \in H^1_\tau(U, \mathcal{F}^a)$.
D'apr\`es le lemme \ref{lemma-locality-cohomology},
il existe un recouvrement fpqc $\mathcal{U} = \{U_i \to U\}_{i \in I}$
tel que $\xi|_{U_i} = 0$ pour tout $i \in I$. Quitte \`a raffiner
$\mathcal{U}$, nous pouvons supposer que $\mathcal{U}$ est un
$\tau$-recouvrement standard. En appliquant la suite spectrale du
th\'eor\`eme \ref{theorem-cech-ss},
nous voyons que $\xi$ provient d'une classe de cohomologie
$\check \xi \in \check H^1(\mathcal{U}, \mathcal{F}^a)$.
Consid\'erons le recouvrement $\mathcal{V} = \{\coprod_{i\in I} U_i \to U\}$. D'apr\`es le
lemme \ref{lemma-cech-complex},
$\check H^\bullet(\mathcal{U}, \mathcal{F}^a) =
\check H^\bullet(\mathcal{V}, \mathcal{F}^a)$.
D'autre part, puisque $\mathcal{V}$ est un recouvrement de la forme
$\{\Spec(B) \to \Spec(A)\}$ et $f^*\mathcal{F} = \widetilde{M}$
pour un certain $A$-module $M$, nous voyons que le complexe de {\v C}ech
$\check{\mathcal{C}}^\bullet(\mathcal{V}, \mathcal{F})$
n'est autre que le complexe $(B/A)_\bullet \otimes_A M$.
D'apr\`es le lemme \ref{lemma-descent-modules},
$H^p((B/A)_\bullet \otimes_A M) = 0$ pour $p > 0$ ; ainsi, $\check \xi = 0$
et donc $\xi = 0$.
\item[p > 1]
Prenons $\xi \in H^p_\tau(U, \mathcal{F}^a)$. D'apr\`es le
lemme \ref{lemma-locality-cohomology},
il existe un recouvrement fpqc $\mathcal{U} = \{U_i \to U\}_{i \in I}$
tel que $\xi|_{U_i} = 0$ pour tout $i \in I$. Quitte \`a raffiner
$\mathcal{U}$, nous pouvons supposer que $\mathcal{U}$ est un
$\tau$-recouvrement standard. Nous appliquons la suite spectrale du
th\'eor\`eme \ref{theorem-cech-ss}.
Observons que les intersections $U_{i_0} \times_U \ldots \times_U U_{i_p}$
sont affines ; par l'hypoth\`ese de r\'ecurrence, les groupes de cohomologie
$$
E_2^{p, q} = \check H^p(\mathcal{U}, \underline{H}^q(\mathcal{F}^a))
$$
s'annulent donc pour tout $0 < q < p$. Nous voyons que $\xi$ doit provenir d'un
$\check \xi \in \check H^p(\mathcal{U}, \mathcal{F}^a)$. En rempla\c{c}ant
$\mathcal{U}$ par le recouvrement $\mathcal{V}$ qui ne contient qu'un seul morphisme
et en utilisant de nouveau le lemme \ref{lemma-descent-modules},
nous voyons que la classe de cohomologie de {\v C}ech $\check \xi$ doit \^etre nulle,
donc $\xi = 0$.
\end{enumerate}
Supposons ensuite que $U$ soit s\'epar\'e. Choisissons des ouverts affines
$U = \bigcup_{i \in I} U_i$ formant un recouvrement de $U$. La famille
$\mathcal{U} = \{U_i \to U\}_{i \in I}$ est alors un recouvrement fpqc,
et toutes les intersections
$U_{i_0} \times_U \ldots \times_U U_{i_p}$ sont affines,
puisque $U$ est s\'epar\'e. Toutes les lignes de la suite spectrale du
th\'eor\`eme \ref{theorem-cech-ss}
sont donc nulles, sauf la ligne de degr\'e z\'ero. Par cons\'equent,
$$
H^p_\tau(U, \mathcal{F}^a) =
\check H^p(\mathcal{U}, \mathcal{F}^a) =
\check H^p(\mathcal{U}, \mathcal{F}) = H^p(U, \mathcal{F})
$$
o\`u la derni\`ere \'egalit\'e r\'esulte de la th\'eorie usuelle des sch\'emas, voir
Cohomologie des schémas, lemme
\ref{coherent-lemma-cech-cohomology-quasi-coherent}.

\medskip\noindent
Le cas g\'en\'eral est technique et (pour prolonger la d\'emonstration donn\'ee ici)
n\'ecessite une discussion sur les morphismes de suites spectrales ; nous ne le traiterons donc pas.
Il r\'esulte de
Descente, Proposition \ref{descent-proposition-same-cohomology-quasi-coherent},
(dont la d\'emonstration adopte une approche l\'eg\`erement diff\'erente), combin\'ee avec
Cohomologie sur les sites, lemme \ref{sites-cohomology-lemma-cohomology-of-open}.
\end{proof}

\begin{remark}
\label{remark-right-derived-global-sections}
Commentaire sur le th\'eor\`eme \ref{theorem-zariski-fpqc-quasi-coherent}.
Puisque $S$ est un objet final de la cat\'egorie $\mathcal{C}$, les groupes de cohomologie
du membre de droite sont simplement les foncteurs d\'eriv\'es \`a droite du
foncteur des sections globales. En fait, la d\'emonstration montre que
$H^p(U, f^*\mathcal{F}) = H^p_\tau(U, \mathcal{F}^a)$
pour tout objet $f : U \to S$ du site $\mathcal{C}$.
\end{remark}





\section{Exemples de faisceaux}
\label{section-examples-sheaves}

\noindent
Soient $S$ et $\tau$ comme dans la section \ref{section-big-small}.
Nous avons d\'ej\`a vu que tout pr\'efaisceau repr\'esentable est un faisceau sur
$(\Sch/S)_\tau$ ou $S_\tau$, voir
le lemme \ref{lemma-representable-sheaf-fpqc}
et la
remarque \ref{remark-fpqc-finest}.
Voici quelques cas particuliers.

\begin{definition}
\label{definition-additive-sheaf}
Sur l'un quelconque des sites $(\Sch/S)_\tau$ ou $S_\tau$ de la
section \ref{section-big-small} :
\begin{enumerate}
\item Le faisceau $T \mapsto \Gamma(T, \mathcal{O}_T)$ est not\'e
$\mathcal{O}_S$, ou $\mathbf{G}_a$, ou $\mathbf{G}_{a, S}$ si nous
voulons indiquer le sch\'ema de base.
\item De m\^eme, le faisceau
$T \mapsto \Gamma(T, \mathcal{O}^*_T)$ est not\'e $\mathcal{O}_S^*$, ou
$\mathbf{G}_m$, ou $\mathbf{G}_{m, S}$ si nous voulons
indiquer le sch\'ema de base.
\item Le {\it faisceau constant} $\underline{\mathbf{Z}/n\mathbf{Z}}$ sur un
site quelconque est le faisceau associ\'e au pr\'efaisceau constant
$U \mapsto \mathbf{Z}/n\mathbf{Z}$.
\end{enumerate}
\end{definition}

\noindent
Le premier est un faisceau, par exemple d'apr\`es le
th\'eor\`eme \ref{theorem-quasi-coherent}.
Le deuxi\`eme est un sous-pr\'efaisceau du premier, dont on v\'erifie facilement
qu'il est lui-m\^eme un faisceau. Le troisi\`eme est un faisceau par d\'efinition.
Notons que chacun de ces faisceaux est repr\'esentable.
Le premier et le deuxi\`eme le sont par les sch\'emas $\mathbf{G}_{a, S}$ et
$\mathbf{G}_{m, S}$, voir
Groupoïdes, section \ref{groupoids-section-group-schemes}.
Le troisi\`eme l'est par le sch\'ema en groupes fini \'etale $\mathbf{Z}/n\mathbf{Z}_S$,
parfois not\'e $(\mathbf{Z}/n\mathbf{Z})_S$,
qui consiste simplement en $n$ copies de $S$ munies
de l'\'evidente structure de sch\'ema en groupes sur $S$, voir
Groupoïdes, exemple \ref{groupoids-example-constant-group}
et la remarque suivante.

\begin{remark}
\label{remark-constant-locally-constant-maps}
Soit $G$ un groupe abstrait.
Sur l'un quelconque des sites $(\Sch/S)_\tau$ ou $S_\tau$ de la
section \ref{section-big-small},
le faisceau associ\'e $\underline{G}$
au pr\'efaisceau constant associ\'e \`a $G$ pour la
{\it topologie de Zariski} du site donne d\'ej\`a
$$
\Gamma(U, \underline{G}) =
\{\text{applications localement constantes pour la topologie de Zariski }U \to G\}
$$
Ce faisceau de Zariski est repr\'esentable par le sch\'ema en groupes $G_S$ d'apr\`es
Groupoïdes, exemple \ref{groupoids-example-constant-group}.
D'apr\`es le
lemme \ref{lemma-representable-sheaf-fpqc},
tout pr\'efaisceau repr\'esentable satisfait aussi la condition de faisceau pour la
topologie $\tau$ ; nous concluons donc que le faisceau associ\'e de Zariski
$\underline{G}$ ci-dessus est aussi le faisceau associ\'e pour la topologie $\tau$.
\end{remark}

\begin{definition}
\label{definition-structure-sheaf}
Soit $S$ un sch\'ema. Le {\it faisceau structural} de $S$ est le faisceau d'anneaux
$\mathcal{O}_S$
sur l'un quelconque des sites $S_{Zar}$, $S_\etale$ ou $(\Sch/S)_\tau$
consid\'er\'es ci-dessus.
\end{definition}

\noindent
S'il peut y avoir une ambigu\"it\'e quant au site sur lequel nous travaillons,
nous l'indiquerons au moyen d'indices. Par exemple, nous pouvons utiliser
$\mathcal{O}_{S_\etale}$ pour souligner que nous travaillons sur le
petit site \'etale de $S$.

\begin{remark}
\label{remark-special-case-fpqc-cohomology-quasi-coherent}
Dans la terminologie introduite ci-dessus, un cas particulier du
th\'eor\`eme \ref{theorem-zariski-fpqc-quasi-coherent}
est
$$
H_{fppf}^p(X, \mathbf{G}_a) =
H_\etale^p(X, \mathbf{G}_a) =
H_{Zar}^p(X, \mathbf{G}_a) =
H^p(X, \mathcal{O}_X)
$$
pour tout $p \geq 0$. De plus, nous pourrions utiliser la notation
$H^p_{fppf}(X, \mathcal{O}_X)$ pour d\'esigner la cohomologie du
faisceau structural sur le grand site fppf de $X$.
\end{remark}




\section{Groupes de Picard}
\label{section-picard-groups}

\noindent
Le th\'eor\`eme suivant est parfois appel\'e ``th\'eor\`eme 90 de Hilbert''.

\begin{theorem}
\label{theorem-picard-group}
Pour tout sch\'ema $X$, nous avons des identifications canoniques
\begin{align*}
H_{fppf}^1(X, \mathbf{G}_m) & = H^1_{syntomic}(X, \mathbf{G}_m) \\
& = H^1_{smooth}(X, \mathbf{G}_m) \\
& = H_\etale^1(X, \mathbf{G}_m) \\
& = H^1_{Zar}(X, \mathbf{G}_m) \\
& = \Pic(X) \\
& = H^1(X, \mathcal{O}_X^*)
\end{align*}
\end{theorem}

\begin{proof}
Soit $\tau$ l'une des topologies consid\'er\'ees dans la
section \ref{section-big-small}.
D'apr\`es
Cohomologie sur les sites, lemme
\ref{sites-cohomology-lemma-h1-invertible},
nous voyons que
$H^1_\tau(X, \mathbf{G}_m) =
H^1_\tau(X, \mathcal{O}_\tau^*) =
\Pic(\mathcal{O}_\tau)$,
o\`u $\mathcal{O}_\tau$ est le faisceau structural du site
$(\Sch/X)_\tau$. Or un $\mathcal{O}_\tau$-module inversible
est un $\mathcal{O}_\tau$-module quasi-coh\'erent.
D'apr\`es le th\'eor\`eme \ref{theorem-quasi-coherent}, ou plus pr\'ecis\'ement
Descente, Proposition \ref{descent-proposition-equivalence-quasi-coherent},
nous voyons que $\Pic(\mathcal{O}_\tau) = \Pic(X)$.
La derni\`ere \'egalit\'e se d\'emontre de la m\^eme mani\`ere.
\end{proof}






\section{Le site \'etale}
\label{section-etale-site}

\noindent
Nous commen\c{c}ons maintenant \`a examiner plus en d\'etail le site \'etale d'un
sch\'ema. Dans un premier temps, nous examinons bri\`evement la notion de
morphisme \'etale.





\section{Morphismes \'etales}
\label{section-etale-morphism}

\noindent
Pour plus de d\'etails, voir
Morphismes, section \ref{morphisms-section-etale}
pour la d\'efinition formelle, et
Morphismes étales, sections
\ref{etale-section-etale-morphisms},
\ref{etale-section-structure-etale-map},
\ref{etale-section-etale-smooth},
\ref{etale-section-topological-etale},
\ref{etale-section-functorial-etale}, et
\ref{etale-section-properties-permanence}
pour un aper\c{c}u de propri\'et\'es int\'eressantes des morphismes \'etales.

\medskip\noindent
Rappelons qu'une alg\`ebre $A$ sur un corps alg\'ebriquement clos $k$ est
{\it lisse} si elle est de type fini et si le module des diff\'erentielles
$\Omega_{A/k}$ est localement libre de type fini, de rang \`egal \`a la dimension.
Un sch\'ema $X$ sur $k$ est {\it lisse} sur $k$ s'il est localement de type
fini et si chaque ouvert affine est le spectre d'une $k$-alg\`ebre lisse.
Si $k$ n'est pas alg\'ebriquement clos, alors une $k$-alg\`ebre $A$ est
une $k$-alg\`ebre lisse si $A \otimes_k \overline{k}$ est une
$\overline{k}$-alg\`ebre lisse. Un morphisme d'anneaux $A \to B$ est lisse s'il est
plat, de pr\'esentation finie et si, pour tout id\'eal premier $\mathfrak p \subset A$,
l'anneau de la fibre $\kappa(\mathfrak p) \otimes_A B$ est lisse sur le corps
r\'esiduel $\kappa(\mathfrak p)$. Plus g\'en\'eralement, un morphisme de sch\'emas est
{\it lisse} s'il est plat, localement de pr\'esentation finie et si ses
fibres g\'eom\'etriques sont lisses.

\medskip\noindent
Pour ces faits, voir
Morphismes, section \ref{morphisms-section-smooth}.
Cela nous permet de d\'efinir un morphisme \'etale comme suit.

\begin{definition}
\label{definition-etale-morphism}
Un morphisme de sch\'emas est {\it \'etale} s'il est lisse de dimension relative 0.
\end{definition}

\noindent
En particulier, un morphisme de sch\'emas $X \to S$ est \'etale s'il est lisse
et si $\Omega_{X/S} = 0$.

\begin{proposition}
\label{proposition-etale-morphisms}
Propri\'et\'es des morphismes \'etales.
\begin{enumerate}
\item Soit $k$ un corps. Un morphisme de sch\'emas $U \to \Spec(k)$ est
\'etale si et seulement si $U \cong \coprod_{i \in I} \Spec(k_i)$,
o\`u, pour chaque $i \in I$,
le corps $k_i$ est une extension finie s\'eparable de $k$.
\item Soit $\varphi : U \to S$ un morphisme de sch\'emas. Les
conditions suivantes sont \`equivalentes :
\begin{enumerate}
\item $\varphi$ est \'etale,
\item $\varphi$ est localement de pr\'esentation finie, plat, et toutes ses fibres sont
\'etales,
\item $\varphi$ est plat, non ramifi\'e et localement de pr\'esentation finie.
\end{enumerate}
\item Un morphisme d'anneaux $A \to B$ est \'etale si et seulement si
$B \cong A[x_1, \ldots, x_n]/(f_1, \ldots, f_n)$
tel que $\Delta = \det \left( \frac{\partial f_i}{\partial x_j} \right)$
soit inversible dans $B$.
\item Le changement de base d'un morphisme \'etale est \'etale.
\item Les compos\'ees de morphismes \'etales sont \'etales.
\item Les produits fibr\'es et les produits de morphismes \'etales sont \'etales.
\item Un morphisme \'etale a pour dimension relative 0.
\item Soit $Y \to X$ un morphisme \'etale.
Si $X$ est r\'eduit (respectivement r\'egulier), alors $Y$ l'est aussi.
\item Les morphismes \'etales sont ouverts.
\item Si $X \to S$ et $Y \to S$ sont \'etales, alors tout
$S$-morphisme $X \to Y$ est lui aussi \'etale.
\end{enumerate}
\end{proposition}

\begin{proof}
Nous avons d\'emontr\'e ces faits (et davantage) dans les chapitres pr\'ec\'edents.
Voici une liste de r\'ef\'erences :
(1) Morphismes, lemme \ref{morphisms-lemma-etale-over-field}.
(2) Morphismes, lemmes \ref{morphisms-lemma-etale-flat-etale-fibres}
et \ref{morphisms-lemma-flat-unramified-etale}.
(3) Algèbre, lemme \ref{algebra-lemma-etale-standard-smooth}.
(4) Morphismes, lemme \ref{morphisms-lemma-base-change-etale}.
(5) Morphismes, lemme \ref{morphisms-lemma-composition-etale}.
(6) D\'ecoule formellement de (4) et (5).
(7) Morphismes, lemmes \ref{morphisms-lemma-etale-locally-quasi-finite}
et \ref{morphisms-lemma-locally-quasi-finite-rel-dimension-0}.
(8) Voir Algèbre, lemmes \ref{algebra-lemma-reduced-goes-up} et
\ref{algebra-lemma-Rk-goes-up} ; voir aussi d'autres r\'esultats de ce type
dans Morphismes étales, section \ref{etale-section-properties-permanence}.
(9) Voir Morphismes, lemme \ref{morphisms-lemma-fppf-open} et
\ref{morphisms-lemma-etale-flat}.
(10) Voir Morphismes, lemme \ref{morphisms-lemma-etale-permanence}.
\end{proof}

\begin{definition}
\label{definition-standard-etale}
Un morphisme d'anneaux $A \to B$ est dit {\it \'etale standard} si
$B \cong \left(A[t]/(f)\right)_g$ avec $f, g \in A[t]$, o\`u $f$ est unitaire,
et si $\text{d}f/\text{d}t$ est inversible dans $B$.
\end{definition}

\noindent
Il est vrai qu'un morphisme d'anneaux \'etale standard est \'etale. En effet, supposons
que $B = \left(A[t]/(f)\right)_g$ avec $f, g \in A[t]$, o\`u $f$ est unitaire,
et que $\text{d}f/\text{d}t$ soit inversible dans $B$. Alors $A[t]/(f)$ est un
$A$-module libre de rang \'egal au degr\'e du polyn\^ome unitaire $f$.
Par cons\'equent, $B$, en tant que localisation de cette alg\`ebre qui est libre comme module, est de pr\'esentation finie
et plat sur $A$. Pour achever la preuve que $B$ est \'etale, il suffit
de montrer que les anneaux des fibres
$$
\kappa(\mathfrak p) \otimes_A B
\cong
\kappa(\mathfrak p) \otimes_A (A[t]/(f))_g
\cong
\kappa(\mathfrak p)[t, 1/\overline{g}]/(\overline{f})
$$
sont des produits finis d'extensions de corps finies s\'eparables.
Ici, $\overline{f}, \overline{g} \in \kappa(\mathfrak p)[t]$ sont
les images de $f$ et $g$. Soit
$$
\overline{f} = \overline{f}_1 \ldots \overline{f}_a
\overline{f}_{a + 1}^{e_1} \ldots \overline{f}_{a + b}^{e_b}
$$
la factorisation de $\overline{f}$ en puissances de facteurs unitaires irr\'eductibles
deux \`a deux distincts $\overline{f}_i$, avec $e_1, \ldots, e_b > 1$.
Par hypoth\`ese, $\text{d}\overline{f}/\text{d}t$ est inversible dans
$\kappa(\mathfrak p)[t, 1/\overline{g}]$. Nous voyons donc que
au moins les $\overline{f}_i$, $i > a$, sont inversibles. Nous concluons
que
$$
\kappa(\mathfrak p)[t, 1/\overline{g}]/(\overline{f})
\cong
\prod\nolimits_{i \in I} \kappa(\mathfrak p)[t]/(\overline{f}_i)
$$
o\`u $I \subset \{1, \ldots, a\}$ est le sous-ensemble des indices $i$ tels que
$\overline{f}_i$ ne divise pas $\overline{g}$. De plus, l'image de
$\text{d}\overline{f}/\text{d}t$ dans le facteur
$\kappa(\mathfrak p)[t]/(\overline{f}_i)$ est clairement \'egale au produit d'une
unit\'e par $\text{d}\overline{f}_i/\text{d}t$. Nous en concluons que
$\kappa_i = \kappa(\mathfrak p)[t]/(\overline{f}_i)$ est une extension de corps
finie de $\kappa(\mathfrak p)$, engendr\'ee par un \'el\'ement dont le
polyn\^ome minimal est s\'eparable, c'est-\`a-dire que l'extension de corps
$\kappa_i/\kappa(\mathfrak p)$ est finie s\'eparable, comme souhait\'e.

\medskip\noindent
Il se trouve que tout morphisme d'anneaux \'etale est localement \'etale standard.
Pour formuler cela, introduisons la notation suivante.
Un morphisme d'anneaux $A \to B$ est {\it \'etale en un id\'eal premier $\mathfrak q$} de $B$ s'il
existe $h \in B$, $h \not \in \mathfrak q$, tel que $A \to B_h$ soit \'etale.
Voici le r\'esultat.

\begin{theorem}
\label{theorem-standard-etale}
Un morphisme d'anneaux $A \to B$ est \'etale en un id\'eal premier $\mathfrak q$ si et seulement s'il
existe $g \in B$, $g \not \in \mathfrak q$, tel que $B_g$ soit \'etale
standard sur $A$.
\end{theorem}

\begin{proof}
Voir
Algèbre, Proposition \ref{algebra-proposition-etale-locally-standard}.
\end{proof}





\section{Recouvrements \'etales}
\label{section-etale-covering}

\noindent
Rappelons la d\'efinition.

\begin{definition}
\label{definition-etale-covering}
Un {\it recouvrement \'etale} d'un sch\'ema $U$ est une famille de morphismes
de sch\'emas
$\{\varphi_i : U_i \to U\}_{i \in I}$ telle que
\begin{enumerate}
\item chaque $\varphi_i$ est un morphisme \'etale,
\item les $U_i$ recouvrent $U$, c'est-\`a-dire $U = \bigcup_{i\in I}\varphi_i(U_i)$.
\end{enumerate}
\end{definition}

\begin{lemma}
\label{lemma-etale-fpqc}
Tout recouvrement \'etale est un recouvrement fpqc.
\end{lemma}

\begin{proof}
(Voir aussi
Topologies,
lemme \ref{topologies-lemma-zariski-etale-smooth-syntomic-fppf-fpqc}.)
Soit $\{\varphi_i : U_i \to U\}_{i \in I}$ un recouvrement \'etale.
Puisqu'un morphisme \'etale est plat et que les \'el\'ements du recouvrement doivent
recouvrir sa cible, la propri\'et\'e fp (fid\`element plat) est satisfaite.
Pour v\'erifier la propri\'et\'e qc (quasi-compacit\'e), soit $V \subset U$ un ouvert
affine, et \'ecrivons $\varphi_i^{-1}(V) = \bigcup_{j \in J_i} V_{ij}$
pour certains ouverts affines $V_{ij} \subset U_i$. Puisque $\varphi_i$ est ouvert
(car les morphismes \'etales sont ouverts), nous voyons que
$V = \bigcup_{i\in I} \bigcup_{j \in J_i} \varphi_i(V_{ij})$
est un recouvrement ouvert de $V$.
De plus, puisque $V$ est quasi-compact, ce recouvrement admet un
raffinement fini.
\end{proof}

\noindent
Ainsi, tout \'enonc\'e vrai pour les recouvrements fpqc
reste vrai {\it a fortiori} pour les recouvrements \'etales. Par
exemple, le site \'etale est sous-canonique.

\begin{definition}
\label{definition-big-etale-site}
(Pour plus de d\'etails, voir la section \ref{section-big-small}, ou
Topologies, section \ref{topologies-section-etale}.)
Soit $S$ un sch\'ema.
Le {\it grand site \'etale de $S$} est le site
$(\Sch/S)_\etale$, voir
la d\'efinition \ref{definition-tau-site}.
Le {\it petit site \'etale de $S$} est le site $S_\etale$, voir
la d\'efinition \ref{definition-tau-site}.
Nous d\'efinissons de m\^eme le {\it grand} et le {\it petit site de Zariski} de $S$,
not\'es $(\Sch/S)_{Zar}$ et $S_{Zar}$.
\end{definition}

\noindent
Grosso modo, le grand site \'etale de $S$ est constitu\'e des sch\'emas sur $S$
et ses recouvrements sont les recouvrements \'etales. Le petit site \'etale de $S$ est constitu\'e
des sch\'emas \'etales sur $S$, avec pour recouvrements les recouvrements \'etales.
En fait, tout morphisme entre objets de $S_\etale$ est \'etale, en
vertu de la
proposition \ref{proposition-etale-morphisms} ;
par cons\'equent, pour v\'erifier que $\{U_i \to U\}_{i \in I}$ dans $S_\etale$
est un recouvrement, il suffit de v\'erifier que $\coprod U_i \to U$ est surjectif.

\medskip\noindent
Le petit site \'etale a moins d'objets que le grand site \'etale ; il
ne contient que les ``ouverts'' de la topologie \'etale sur $S$. C'est une sous-cat\'egorie
pleine du grand site \'etale, et sa topologie est induite par la
topologie du grand site. Ainsi, le foncteur de restriction
du grand site \'etale au petit est exact et transforme les objets injectifs en
objets injectifs. Cela a la cons\'equence suivante.

\begin{proposition}
\label{proposition-cohomology-restrict-small-site}
Soit $S$ un sch\'ema et soit $\mathcal{F}$ un faisceau ab\'elien sur
$(\Sch/S)_\etale$.
Alors $\mathcal{F}|_{S_\etale}$ est un faisceau sur $S_\etale$ et
$$
H^p_\etale(S, \mathcal{F}|_{S_\etale}) =
H^p_\etale(S, \mathcal{F})
$$
pour tout $p \geq 0$.
\end{proposition}

\begin{proof}
C'est un cas particulier du lemme \ref{lemma-compare-cohomology-big-small}.
\end{proof}

\noindent
Conform\'ement \`a la notation g\'en\'erale introduite dans la
section \ref{section-big-small},
nous notons $H_\etale^p(S, \mathcal{F})$ le groupe de cohomologie ci-dessus.





%9.24.09
\section{Th\'eorie de Kummer}
\label{section-kummer}

\noindent
Soit $n \in \mathbf{N}$ et consid\'erons le foncteur $\mu_n$ d\'efini par
$$
\begin{matrix}
\Sch^{opp} & \longrightarrow & \textit{Ab} \\
S & \longmapsto &
\mu_n(S)
=
\{t \in \Gamma(S, \mathcal{O}_S^*) \mid t^n = 1 \}.
\end{matrix}
$$
D'apr\`es
Groupoïdes, exemple \ref{groupoids-example-roots-of-unity},
ce foncteur est repr\'esentable, et le sch\'ema qui le repr\'esente
est lui aussi not\'e $\mu_n$. D'apr\`es le
lemme \ref{lemma-representable-sheaf-fpqc},
ce foncteur satisfait \`a la condition de faisceau pour la topologie fpqc
(en particulier, il satisfait aussi \`a la condition de faisceau pour les
topologies \'etale, de Zariski, etc.).

\begin{lemma}
\label{lemma-kummer-sequence}
Si $n\in \mathcal{O}_S^*$, alors
$$
0 \to
\mu_{n, S} \to
\mathbf{G}_{m, S} \xrightarrow{(\cdot)^n}
\mathbf{G}_{m, S} \to 0
$$
est une suite exacte courte de faisceaux aussi bien sur le petit que sur le
grand site \'etale de $S$.
\end{lemma}

\begin{proof}
Par d\'efinition, le faisceau $\mu_{n, S}$ est le noyau de l'application
$(\cdot)^n$. Il suffit donc de montrer que la derni\`ere application est surjective.
Soit $U$ un sch\'ema sur $S$. Soit
$f \in \mathbf{G}_m(U) = \Gamma(U, \mathcal{O}_U^*)$.
Nous devons montrer que nous pouvons trouver un recouvrement \'etale de
$U$ tel que la restriction de $f$ \`a chacun de ses membres soit une puissance $n$-i\`eme.
Posons
$$
U' =
\underline{\Spec}_U(\mathcal{O}_U[T]/(T^n-f))
\xrightarrow{\pi}
U.
$$
(Voir
Constructions, section \ref{constructions-section-spec-via-glueing} ou
\ref{constructions-section-spec}
pour une pr\'esentation du spectre relatif.)
Soit $\Spec(A) \subset U$ un ouvert affine, et supposons que $f|_{\Spec(A)}$ corresponde
\`a l'unit\'e $a \in A^*$. Alors $\pi^{-1}(\Spec(A)) = \Spec(B)$ avec
$B = A[T]/(T^n - a)$. Le morphisme d'anneaux $A \to B$ est fini libre de rang $n$,
donc fid\`element plat, et nous en concluons que
$\Spec(B) \to \Spec(A)$ est surjectif. Comme cela vaut pour tout
ouvert affine de $U$, nous concluons que $\pi$ est surjectif.
De plus, $n$ et $T^{n - 1}$ sont inversibles dans $B$, donc
$nT^{n-1} \in B^*$ et le morphisme d'anneaux $A \to B$ est \'etale standard,
en particulier \'etale. Comme cela vaut pour tout ouvert affine de $U$,
nous concluons que $\pi$ est \'etale. Par cons\'equent,
$\mathcal{U} = \{\pi : U' \to U\}$ est un recouvrement \'etale.
De plus, $f|_{U'} = (f')^n$, o\`u $f'$ est la classe de $T$
dans $\Gamma(U', \mathcal{O}_{U'}^*)$ ; ainsi, $\mathcal{U}$ poss\`ede la propri\'et\'e voulue.
\end{proof}

\begin{remark}
\label{remark-no-kummer-sequence-zariski}
Le lemme \ref{lemma-kummer-sequence} est faux lorsque ``\'etale'' est remplac\'e
par ``Zariski''.
Comme la topologie \'etale est moins fine que la topologie lisse, voir
Topologies, lemme \ref{topologies-lemma-zariski-etale-smooth},
il s'ensuit que la suite est aussi exacte pour la topologie lisse.
\end{remark}

\noindent
D'apr\`es le
th\'eor\`eme \ref{theorem-picard-group},
le
lemme \ref{lemma-kummer-sequence}
et les propri\'et\'es g\'en\'erales de la cohomologie, nous obtenons
la suite exacte longue de cohomologie
$$
\xymatrix{
0 \ar[r] &
H_\etale^0(S, \mu_{n, S}) \ar[r] &
\Gamma(S, \mathcal{O}_S^*) \ar[r]^{(\cdot)^n} &
\Gamma(S, \mathcal{O}_S^*) \ar@(rd, ul)[rdllllr]
\\
& H_\etale^1(S, \mu_{n, S}) \ar[r] &
\Pic(S) \ar[r]^{(\cdot)^n} &
\Pic(S) \ar@(rd, ul)[rdllllr] \\
& H_\etale^2(S, \mu_{n, S}) \ar[r] &
\ldots
}
$$
du moins si $n$ est inversible sur $S$. Lorsque $n$ n'est pas inversible sur $S$,
nous pouvons appliquer le lemme suivant.

\begin{lemma}
\label{lemma-kummer-sequence-syntomic}
Pour tout $n \in \mathbf{N}$, la suite
$$
0 \to
\mu_{n, S} \to
\mathbf{G}_{m, S} \xrightarrow{(\cdot)^n}
\mathbf{G}_{m, S} \to 0
$$
est une suite exacte courte de faisceaux sur les sites
$(\Sch/S)_{fppf}$ et $(\Sch/S)_{syntomic}$.
\end{lemma}

\begin{proof}
Par d\'efinition, le faisceau $\mu_{n, S}$ est le noyau de l'application
$(\cdot)^n$. Il suffit donc de montrer que la derni\`ere application est surjective.
Comme la topologie syntomique est moins fine que la topologie fppf, voir
Topologies, lemme \ref{topologies-lemma-zariski-etale-smooth-syntomic-fppf},
il suffit de le d\'emontrer pour la topologie syntomique.
Soit $U$ un sch\'ema sur $S$. Soit
$f \in \mathbf{G}_m(U) = \Gamma(U, \mathcal{O}_U^*)$.
Nous devons montrer que nous pouvons trouver un recouvrement syntomique de
$U$ tel que la restriction de $f$ \`a chacun de ses membres soit une puissance $n$-i\`eme.
Posons
$$
U' =
\underline{\Spec}_U(\mathcal{O}_U[T]/(T^n-f))
\xrightarrow{\pi}
U.
$$
(Voir
Constructions, section \ref{constructions-section-spec-via-glueing} ou
\ref{constructions-section-spec}
pour une pr\'esentation du spectre relatif.)
Soit $\Spec(A) \subset U$ un ouvert affine, et supposons que $f|_{\Spec(A)}$ corresponde
\`a l'unit\'e $a \in A^*$. Alors $\pi^{-1}(\Spec(A)) = \Spec(B)$ avec
$B = A[T]/(T^n - a)$. Le morphisme d'anneaux $A \to B$ est fini libre de rang $n$,
donc fid\`element plat, et nous en concluons que
$\Spec(B) \to \Spec(A)$ est surjectif. Comme cela vaut pour tout
ouvert affine de $U$, nous concluons que $\pi$ est surjectif.
De plus, $B$ est une intersection compl\`ete globale relative sur $A$, donc
le morphisme d'anneaux $A \to B$ est syntomique standard,
en particulier syntomique. Comme cela vaut pour tout ouvert affine de $U$,
nous concluons que $\pi$ est syntomique. Par cons\'equent,
$\mathcal{U} = \{\pi : U' \to U\}$ est un recouvrement syntomique.
De plus, $f|_{U'} = (f')^n$, o\`u $f'$ est la classe de $T$
dans $\Gamma(U', \mathcal{O}_{U'}^*)$ ; ainsi, $\mathcal{U}$ poss\`ede la propri\'et\'e voulue.
\end{proof}

\begin{remark}
\label{remark-no-kummer-sequence-smooth-etale-zariski}
Le lemme \ref{lemma-kummer-sequence-syntomic}
est faux pour les topologies lisse, \'etale ou de Zariski.
\end{remark}

\noindent
D'apr\`es le
th\'eor\`eme \ref{theorem-picard-group},
le
lemme \ref{lemma-kummer-sequence-syntomic}
et les propri\'et\'es g\'en\'erales de la cohomologie, nous obtenons
la suite exacte longue de cohomologie
$$
\xymatrix{
0 \ar[r] &
H_{fppf}^0(S, \mu_{n, S}) \ar[r] &
\Gamma(S, \mathcal{O}_S^*) \ar[r]^{(\cdot)^n} &
\Gamma(S, \mathcal{O}_S^*) \ar@(rd, ul)[rdllllr]
\\
& H_{fppf}^1(S, \mu_{n, S}) \ar[r] &
\Pic(S) \ar[r]^{(\cdot)^n} &
\Pic(S) \ar@(rd, ul)[rdllllr] \\
& H_{fppf}^2(S, \mu_{n, S}) \ar[r] &
\ldots
}
$$
pour tout sch\'ema $S$ et tout entier $n$. Bien entendu, il existe une suite analogue
en cohomologie syntomique.

\medskip\noindent
Soit $n \in \mathbf{N}$ et soit $S$ un sch\'ema quelconque.
Il existe une autre fa\c{c}on, plus directe, de d\'ecrire le premier groupe de cohomologie \`a
valeurs dans $\mu_n$. Consid\'erons les paires
$(\mathcal{L}, \alpha)$ o\`u $\mathcal{L}$ est un faisceau inversible sur $S$
et $\alpha : \mathcal{L}^{\otimes n} \to \mathcal{O}_S$ est une trivialisation
de la $n$-i\`eme puissance tensorielle de $\mathcal{L}$.
Soit $(\mathcal{L}', \alpha')$ une seconde paire de ce type.
Un isomorphisme $\varphi : (\mathcal{L}, \alpha) \to (\mathcal{L}', \alpha')$
est un isomorphisme $\varphi : \mathcal{L} \to \mathcal{L}'$ de faisceaux
inversibles tel que le diagramme
$$
\xymatrix{
\mathcal{L}^{\otimes n} \ar[d]_{\varphi^{\otimes n}} \ar[r]_\alpha &
\mathcal{O}_S \ar[d]^1 \\
(\mathcal{L}')^{\otimes n} \ar[r]^{\alpha'} &
\mathcal{O}_S \\
}
$$
soit commutatif. Nous avons donc
\begin{equation}
\label{equation-isomorphisms-pairs}
\mathit{Isom}_S((\mathcal{L}, \alpha), (\mathcal{L}', \alpha'))
=
\left\{
\begin{matrix}
\emptyset & \text{si} & \text{elles ne sont pas isomorphes} \\
H^0(S, \mu_{n, S})\cdot \varphi & \text{si} &
\varphi \text{ est un isomorphisme de paires}
\end{matrix}
\right.
\end{equation}
De plus, \'etant donn\'ees deux paires $(\mathcal{L}, \alpha)$, $(\mathcal{L}', \alpha')$,
le produit tensoriel
$$
(\mathcal{L}, \alpha) \otimes (\mathcal{L}', \alpha')
=
(\mathcal{L} \otimes \mathcal{L}', \alpha \otimes \alpha')
$$
est encore une paire. La paire $(\mathcal{O}_S, 1)$ est un \'el\'ement neutre pour cette
op\'eration de produit tensoriel, et un inverse est donn\'e par
$$
(\mathcal{L}, \alpha)^{-1} = (\mathcal{L}^{\otimes -1}, \alpha^{\otimes -1}).
$$
Ainsi, l'ensemble des classes d'isomorphisme de paires forme un groupe ab\'elien.
Notons que
$$
(\mathcal{L}, \alpha)^{\otimes n}
=
(\mathcal{L}^{\otimes n}, \alpha^{\otimes n})
\xrightarrow{\alpha}
(\mathcal{O}_S, 1)
$$
est un isomorphisme ;
par cons\'equent, tout \'el\'ement de ce groupe est d'ordre divisant $n$. Nous avertissons le lecteur
que ce groupe n'est en g\'en\'eral {\bf pas} la $n$-torsion de $\Pic(S)$.

\begin{lemma}
\label{lemma-describe-h1-mun}
Soit $S$ un sch\'ema. Il existe une identification canonique
$$
H_\etale^1(S, \mu_n) =
\text{groupe des paires }(\mathcal{L}, \alpha)\text{ \`a isomorphisme pr\`es comme ci-dessus}
$$
si $n$ est inversible sur $S$. En g\'en\'eral, nous avons
$$
H_{fppf}^1(S, \mu_n) =
\text{groupe des paires }(\mathcal{L}, \alpha)\text{ \`a isomorphisme pr\`es comme ci-dessus}.
$$
Le m\^eme r\'esultat vaut en rempla\c{c}ant fppf par syntomique.
\end{lemma}

\begin{proof}
Nous d\'emontrons d'abord le second isomorphisme.
Soit $(\mathcal{L}, \alpha)$ une paire comme ci-dessus.
Choisissons un recouvrement de $S = \bigcup U_i$ par des ouverts affines tel que
$\mathcal{L}|_{U_i} \cong \mathcal{O}_{U_i}$. Choisissons un g\'en\'erateur $s_i \in \mathcal{L}(U_i)$.
Alors $\alpha(s_i^{\otimes n}) = f_i \in \mathcal{O}_S^*(U_i)$.
En \'ecrivant $U_i = \Spec(A_i)$, nous voyons qu'il existe une
intersection compl\`ete globale relative $A_i \to B_i = A_i[T]/(T^n - f_i)$
telle que $f_i$ devienne une puissance $n$-i\`eme dans $B_i$. Autrement dit, en posant
$V_i = \Spec(B_i)$, nous obtenons un recouvrement syntomique
$\mathcal{V} = \{V_i \to S\}_{i \in I}$ et des trivialisations
$\varphi_i : (\mathcal{L}, \alpha)|_{V_i} \to (\mathcal{O}_{V_i}, 1)$.

\medskip\noindent
Nous utiliserons ce r\'esultat (l'existence du recouvrement $\mathcal{V}$)
pour associer \`a cette paire une classe de cohomologie dans
$H^1_{syntomic}(S, \mu_{n, S})$. Nous donnons deux constructions (\'equivalentes).

\medskip\noindent
Premi\`ere construction : au moyen de la cohomologie de {\v C}ech.
Sur les intersections doubles $V_i \times_S V_j$, nous avons l'isomorphisme
$$
(\mathcal{O}_{V_i \times_S V_j}, 1)
\xrightarrow{\text{pr}_0^*\varphi_i^{-1}}
(\mathcal{L}|_{V_i \times_S V_j}, \alpha|_{V_i \times_S V_j})
\xrightarrow{\text{pr}_1^*\varphi_j}
(\mathcal{O}_{V_i \times_S V_j}, 1)
$$
de paires. D'apr\`es (\ref{equation-isomorphisms-pairs}), il est donn\'e par un
\'el\'ement $\zeta_{ij} \in \mu_n(V_i \times_S V_j)$. Nous omettons la v\'erification
que ces $\zeta_{ij}$ forment un $1$-cocycle, c'est-\`a-dire qu'ils donnent
un \'el\'ement $(\zeta_{i_0i_1}) \in \check C(\mathcal{V}, \mu_n)$
tel que $d(\zeta_{i_0i_1}) = 0$. Sa classe est donc un \'el\'ement de
$\check H^1(\mathcal{V}, \mu_n)$ et, d'apr\`es le
th\'eor\`eme \ref{theorem-cech-ss},
elle a pour image une classe de cohomologie dans $H^1_{syntomic}(S, \mu_{n, S})$.

\medskip\noindent
Seconde construction : au moyen des torseurs. Consid\'erons le pr\'efaisceau
$$
\mu_n(\mathcal{L}, \alpha) :
U
\longmapsto
\mathit{Isom}_U((\mathcal{O}_U, 1), (\mathcal{L}, \alpha)|_U)
$$
sur $(\Sch/S)_{syntomic}$.
Nous pouvons le voir comme un sous-pr\'efaisceau de
$\SheafHom_\mathcal{O}(\mathcal{O}, \mathcal{L})$ (faisceau Hom interne,
voir
Modules sur les sites, section \ref{sites-modules-section-internal-hom}).
Comme les conditions qui d\'efinissent ce sous-pr\'efaisceau sont locales, nous voyons que c'est
un faisceau.
D'apr\`es (\ref{equation-isomorphisms-pairs}), ce faisceau est muni d'une action libre du
faisceau $\mu_{n, S}$. Il reste donc seulement \`a v\'erifier qu'il
admet localement des sections. C'est vrai gr\^ace \`a l'existence du
recouvrement trivialisant $\mathcal{V}$. Ainsi, $\mu_n(\mathcal{L}, \alpha)$
est un $\mu_{n, S}$-torseur et, d'apr\`es
Cohomologie sur les sites, lemme \ref{sites-cohomology-lemma-torsors-h1},
nous obtenons un \'el\'ement correspondant de $H^1_{syntomic}(S, \mu_{n, S})$.

\medskip\noindent
Il nous reste maintenant \`a montrer les points suivants :
\begin{enumerate}
\item Les deux constructions donnent la m\^eme classe de cohomologie.
\item Des paires isomorphes donnent la m\^eme classe de cohomologie.
\item La classe de cohomologie de
$(\mathcal{L}, \alpha) \otimes (\mathcal{L}', \alpha')$
est la somme des classes de cohomologie de
$(\mathcal{L}, \alpha)$ et $(\mathcal{L}', \alpha')$.
\item Si la classe de cohomologie est triviale, alors la paire est triviale.
\item Tout \'el\'ement de $H^1_{syntomic}(S, \mu_{n, S})$ est la
classe de cohomologie d'une paire.
\end{enumerate}
Nous omettons la preuve de (1). Le point (2) ressort clairement de la seconde construction,
puisque des torseurs isomorphes donnent la m\^eme classe de cohomologie.
Le point (3) ressort clairement de la premi\`ere construction, puisque les classes de
{\v C}ech obtenues s'additionnent. Le point (4) ressort clairement de la seconde construction,
puisqu'un torseur est trivial si et seulement s'il poss\`ede une section globale, voir
Cohomologie sur les sites, lemme \ref{sites-cohomology-lemma-trivial-torsor}.

\medskip\noindent
Le point (5) se voit comme suit (bien qu'une preuve directe soit
pr\'ef\'erable). Supposons que $\xi \in H^1_{syntomic}(S, \mu_{n, S})$.
Alors $\xi$ s'envoie sur un \'el\'ement
$\overline{\xi} \in H^1_{syntomic}(S, \mathbf{G}_{m, S})$
tel que $n \overline{\xi} = 0$. D'apr\`es le
th\'eor\`eme \ref{theorem-picard-group},
nous voyons que $\overline{\xi}$ correspond \`a un faisceau inversible $\mathcal{L}$
dont la $n$-i\`eme puissance tensorielle est isomorphe \`a $\mathcal{O}_S$.
Il existe donc une paire $(\mathcal{L}, \alpha')$ dont la classe de cohomologie
$\xi'$ a la m\^eme image $\overline{\xi'}$ dans
$H^1_{syntomic}(S, \mathbf{G}_{m, S})$. Il suffit donc de montrer
que $\xi - \xi'$ est la classe d'une paire. Par construction et d'apr\`es la
suite exacte longue de cohomologie ci-dessus, nous voyons que
$\xi - \xi' = \partial(f)$ pour un certain $f \in H^0(S, \mathcal{O}_S^*)$.
Consid\'erons la paire $(\mathcal{O}_S, f)$. Nous omettons la v\'erification
que la classe de cohomologie de cette paire est $\partial(f)$, ce qui
ach\`eve la preuve de la premi\`ere identification (avec fppf remplac\'ee
par syntomique).

\medskip\noindent
Pour voir la premi\`ere identification, notons que si $n$ est inversible sur $S$, alors le
recouvrement $\mathcal{V}$ construit dans la premi\`ere partie de la preuve
est en fait un recouvrement \'etale (comparer avec la preuve du
lemme \ref{lemma-kummer-sequence}). Le reste de la preuve est ind\'ependant
de la topologie, \`a l'exception du tout dernier argument, qui utilise l'exactitude de
la suite de Kummer, c'est-\`a-dire le lemme \ref{lemma-kummer-sequence}.
\end{proof}






\section{Voisinages, fibres et points}
\label{section-stalks}

\noindent
Nous pouvons associer \`a tout point g\'eom\'etrique de $S$ un foncteur fibre qui est
exact. Un morphisme de faisceaux sur $S_\etale$ est un isomorphisme si et seulement
s'il
induit un isomorphisme sur toutes ces fibres. Un complexe de faisceaux ab\'eliens est
exact si et seulement si le complexe des fibres est exact en tout point g\'eom\'etrique.
Ces faits, pris ensemble, signifient que le petit site \'etale d'un sch\'ema $S$
poss\`ede assez de points. Il se trouve aussi que tout point du petit topos \'etale
de $S$ (notion abstraite) est donn\'e par un point g\'eom\'etrique.
Ainsi, en un certain sens, le petit topos \'etale de $S$ peut se comprendre en
termes de points g\'eom\'etriques et de voisinages.

\begin{definition}
\label{definition-geometric-point}
Soit $S$ un sch\'ema.
\begin{enumerate}
\item Un {\it point g\'eom\'etrique} de $S$ est un morphisme
$\Spec(k) \to S$ o\`u $k$ est alg\'ebriquement clos.
Un tel point se note habituellement $\overline{s}$, autrement dit au moyen d'une
lettre minuscule surmont\'ee d'une barre. Nous employons souvent $\overline{s}$ pour d\'esigner aussi bien le sch\'ema
$\Spec(k)$ que le morphisme, et nous employons $\kappa(\overline{s})$
pour d\'esigner $k$.
\item Nous disons que $\overline{s}$ {\it est au-dessus de} $s$
pour indiquer que $s \in S$ est l'image de $\overline{s}$.
\item Un {\it voisinage \'etale} d'un point g\'eom\'etrique $\overline{s}$
de $S$ est un diagramme commutatif
$$
\xymatrix{
& U \ar[d]^\varphi \\
{\overline{s}} \ar[r]^{\overline{s}} \ar[ur]^{\bar u} & S
}
$$
o\`u $\varphi$ est un morphisme \'etale de sch\'emas.
Nous \'ecrivons $(U, \overline{u}) \to (S, \overline{s})$.
\item Un {\it morphisme de voisinages \'etales}
$(U, \overline{u}) \to (U', \overline{u}')$
est un $S$-morphisme $h: U \to U'$
tel que $\overline{u}' = h \circ \overline{u}$.
\end{enumerate}
\end{definition}

\begin{remark}
\label{remark-etale-between-etale}
Puisque $U$ et $U'$ sont \'etales sur $S$, tout $S$-morphisme
entre eux est lui aussi \'etale, voir la
proposition \ref{proposition-etale-morphisms}.
En particulier, tous les morphismes de voisinages \'etales sont \'etales.
\end{remark}

\begin{remark}
\label{remark-etale-neighbourhoods}
Soient $S$ un sch\'ema et $s \in S$ un point. Dans
Compléments sur les morphismes,
définition \ref{more-morphisms-definition-etale-neighbourhood},
nous avons d\'efini la notion de voisinage \'etale $(U, u) \to (S, s)$
de $(S, s)$. Si $\overline{s}$ est un point g\'eom\'etrique de $S$ au-dessus de
$s$, alors tout voisinage \'etale $(U, \overline{u}) \to (S, \overline{s})$
donne lieu \`a un voisinage \'etale $(U, u)$ de $(S, s)$ en prenant pour
$u \in U$ l'unique point de $U$ tel que $\overline{u}$
soit au-dessus de $u$. R\'eciproquement, \'etant donn\'e un voisinage \'etale $(U, u)$
de $(S, s)$, l'extension de corps r\'esiduels $\kappa(u)/\kappa(s)$
est finie s\'eparable (voir la
proposition \ref{proposition-etale-morphisms}) ;
nous pouvons donc trouver un plongement $\kappa(u) \subset \kappa(\overline{s})$
qui induit l'identit\'e sur $\kappa(s)$. Autrement dit, nous pouvons trouver un point g\'eom\'etrique
$\overline{u}$ de $U$ au-dessus de $u$ tel que $(U, \overline{u})$
soit un voisinage \'etale de $(S, \overline{s})$.
Nous utiliserons ces observations pour passer d'un type de
voisinage \'etale \`a l'autre.
\end{remark}

\begin{lemma}
\label{lemma-cofinal-etale}
Soit $S$ un sch\'ema, et soit $\overline{s}$ un point g\'eom\'etrique de $S$.
La cat\'egorie des voisinages \'etales est cofiltrante. Plus pr\'ecis\'ement :
\begin{enumerate}
\item Soient $(U_i, \overline{u}_i)_{i = 1, 2}$ deux voisinages \'etales de
$\overline{s}$ dans $S$. Alors il existe un troisi\`eme voisinage \'etale
$(U, \overline{u})$ et des morphismes
$(U, \overline{u}) \to (U_i, \overline{u}_i)$, $i = 1, 2$.
\item Soient $h_1, h_2: (U, \overline{u}) \to (U', \overline{u}')$ deux
morphismes entre voisinages \'etales de $\overline{s}$. Alors il existe un
voisinage \'etale $(U'', \overline{u}'')$ et un morphisme
$h : (U'', \overline{u}'') \to (U, \overline{u})$
qui \'egalise $h_1$ et $h_2$, c'est-\`a-dire tel que
$h_1 \circ h = h_2 \circ h$.
\end{enumerate}
\end{lemma}

\begin{proof}
Pour le point (1), consid\'erons le produit fibr\'e $U = U_1 \times_S U_2$.
Il est \'etale sur $U_1$ et sur $U_2$, puisque les morphismes \'etales sont
pr\'eserv\'es par changement de base, voir la
proposition \ref{proposition-etale-morphisms}.
Le morphisme $\overline{s} \to U$ d\'efini par $(\overline{u}_1, \overline{u}_2)$
lui conf\`ere la structure d'un voisinage \'etale muni de morphismes vers
$U_1$ et $U_2$. Pour le point (2), d\'efinissons $U''$ comme le produit fibr\'e
$$
\xymatrix{
U'' \ar[r] \ar[d] & U \ar[d]^{(h_1, h_2)} \\
U' \ar[r]^-\Delta & U' \times_S U'.
}
$$
Puisque les compos\'ees de $\overline{u}$ et $\overline{u}'$ vers $S$ sont toutes deux \'egales \`a $\overline{s}$,
nous voyons que $\overline{u}'' = (\overline{u}, \overline{u}')$ est un point g\'eom\'etrique
de $U''$. En particulier, $U'' \not = \emptyset$.
De plus, puisque $U'$ est \'etale sur $S$, le produit fibr\'e
$U'\times_S U'$ l'est aussi (voir la
proposition \ref{proposition-etale-morphisms}).
Ainsi, la fl\`eche verticale $(h_1, h_2)$ est \'etale d'apr\`es la
remarque \ref{remark-etale-between-etale} ci-dessus.
Par cons\'equent, $U''$ est \'etale sur $U'$ par changement de base, et donc aussi
\'etale sur $S$ (car les compos\'ees de morphismes \'etales sont \'etales).
Ainsi, $(U'', \overline{u}'')$ est une solution au probl\`eme.
\end{proof}

\begin{lemma}
\label{lemma-geometric-lift-to-cover}
Soit $S$ un sch\'ema.
Soit $\overline{s}$ un point g\'eom\'etrique de $S$.
Soit $(U, \overline{u})$ un voisinage \'etale de $\overline{s}$.
Soit $\mathcal{U} = \{\varphi_i : U_i \to U \}_{i\in I}$ un recouvrement \'etale.
Alors il existe $i \in I$ et $\overline{u}_i : \overline{s} \to U_i$
tels que $\varphi_i : (U_i, \overline{u}_i) \to (U, \overline{u})$
soit un morphisme de voisinages \'etales.
\end{lemma}

\begin{proof}
Comme $U = \bigcup_{i\in I} \varphi_i(U_i)$, le produit fibr\'e
$\overline{s} \times_{\overline{u}, U, \varphi_i} U_i$ n'est pas vide
pour un certain $i$. Consid\'erons alors le diagramme cart\'esien
$$
\xymatrix{
\overline{s} \times_{\overline{u}, U, \varphi_i} U_i
\ar[d]^{\text{pr}_1} \ar[r]_-{\text{pr}_2} & U_i
\ar[d]^{\varphi_i} \\
\Spec(k) = \overline{s} \ar@/^1pc/[u]^\sigma
\ar[r]^-{\overline{u}} & U
}
$$
La projection $\text{pr}_1$ est le changement de base d'un morphisme \'etale ; elle
est donc \'etale, voir la
proposition \ref{proposition-etale-morphisms}.
Par cons\'equent, $\overline{s} \times_{\overline{u}, U, \varphi_i} U_i$
est une somme disjointe d'extensions finies s\'eparables de $k$, d'apr\`es la
proposition \ref{proposition-etale-morphisms}. Ici,
$\overline{s} = \Spec(k)$. Mais $k$ est alg\'ebriquement clos, donc toutes
ces extensions sont triviales, et il existe une section $\sigma$ de
$\text{pr}_1$. La compos\'ee
$\text{pr}_2 \circ \sigma$ donne un morphisme compatible avec $\overline{u}$.
\end{proof}

\begin{definition}
\label{definition-stalk}
Soit $S$ un sch\'ema.
Soit $\mathcal{F}$ un pr\'efaisceau sur $S_\etale$.
Soit $\overline{s}$ un point g\'eom\'etrique de $S$.
La {\it fibre} de $\mathcal{F}$ en $\overline{s}$ est
$$
\mathcal{F}_{\overline{s}}
=
\colim_{(U, \overline{u})} \mathcal{F}(U)
$$
o\`u $(U, \overline{u})$ parcourt tous les voisinages
\'etales de $\overline{s}$ dans $S$.
\end{definition}

\noindent
D'apr\`es le lemme \ref{lemma-cofinal-etale}, cette limite inductive est index\'ee par une cat\'egorie
filtrante, \`a savoir l'oppos\'ee de la cat\'egorie des voisinages \'etales.
Autrement dit, un \'el\'ement de $\mathcal{F}_{\overline{s}}$ peut \'etre
consid\'er\'e comme un triplet $(U, \overline{u}, \sigma)$ o\`u
$\sigma \in \mathcal{F}(U)$. Deux triplets
$(U, \overline{u}, \sigma)$, $(U', \overline{u}', \sigma')$
d\'efinissent le m\^eme \'el\'ement de la fibre s'il existe un troisi\`eme
voisinage \'etale $(U'', \overline{u}'')$ et des morphismes de voisinages
\'etales $h : (U'', \overline{u}'') \to (U, \overline{u})$,
$h' : (U'', \overline{u}'') \to (U', \overline{u}')$ tels que
$h^*\sigma = (h')^*\sigma'$ dans $\mathcal{F}(U'')$. Voir
Catégories, section \ref{categories-section-directed-colimits}.

\begin{lemma}
\label{lemma-stalk-gives-point}
Soient $S$ un sch\'ema et $\overline{s}$ un point g\'eom\'etrique de $S$.
Consid\'erons le foncteur
\begin{align*}
u : S_\etale & \longrightarrow \textit{Ensembles}, \\
U & \longmapsto
|U_{\overline{s}}|
=
\{\overline{u} \text{ tel que }(U, \overline{u})
\text{ soit un voisinage \'etale de }\overline{s}\}.
\end{align*}
Ici, $|U_{\overline{s}}|$ d\'esigne l'ensemble sous-jacent \`a la fibre g\'eom\'etrique.
Alors $u$ d\'efinit un point $p$ du site $S_\etale$
(Sites, définition \ref{sites-definition-point})
et le foncteur fibre associ\'e $\mathcal{F} \mapsto \mathcal{F}_p$
(Sites, Equation \ref{sites-equation-stalk})
est le foncteur $\mathcal{F} \mapsto \mathcal{F}_{\overline{s}}$
d\'efini ci-dessus.
\end{lemma}

\begin{proof}
Dans la preuve du
lemme \ref{lemma-geometric-lift-to-cover},
nous avons vu que le sch\'ema $U_{\overline{s}}$ est une somme disjointe de
sch\'emas isomorphes \`a $\overline{s}$. Nous pouvons donc aussi consid\'erer
$|U_{\overline{s}}|$ comme l'ensemble des points g\'eom\'etriques de $U$ situ\'es au-dessus de
$\overline{s}$, c'est-\`a-dire comme l'ensemble des morphismes
$\overline{u} : \overline{s} \to U$ qui s'inscrivent dans le diagramme de la
d\'efinition \ref{definition-geometric-point}.
Il s'ensuit que $u(S)$ est un singleton, et que
$u(U \times_V W) = u(U) \times_{u(V)} u(W)$
chaque fois que $U \to V$ et $W \to V$ sont des morphismes dans $S_\etale$.
Et, \'etant donn\'e un recouvrement $\{U_i \to U\}_{i \in I}$ dans $S_\etale$,
nous voyons que $\coprod u(U_i) \to u(U)$ est surjectif par le
lemme \ref{lemma-geometric-lift-to-cover}.
Par cons\'equent,
Sites, Proposition \ref{sites-proposition-point-limits}
s'applique, de sorte que $p$ est un point du site $S_\etale$.
Enfin, notre foncteur $\mathcal{F} \mapsto \mathcal{F}_{\overline{s}}$
est donn\'e par exactement la m\^eme limite inductive que le foncteur
$\mathcal{F} \mapsto \mathcal{F}_p$ associ\'e \`a $p$ dans
Sites, Equation \ref{sites-equation-stalk},
ce qui d\'emontre la derni\`ere assertion.
\end{proof}

\begin{remark}
\label{remark-map-stalks}
Soit $S$ un sch\'ema et soient $\overline{s} : \Spec(k) \to S$
et $\overline{s}' : \Spec(k') \to S$ deux points g\'eom\'etriques de
$S$. Un {\it morphisme $a : \overline{s} \to \overline{s}'$ de points g\'eom\'etriques}
est simplement un morphisme $a : \Spec(k) \to \Spec(k')$ tel que
$\overline{s}' \circ a = \overline{s}$. \'Etant donn\'e un tel morphisme, nous obtenons
un foncteur de la cat\'egorie des voisinages \'etales de $\overline{s}'$
vers la cat\'egorie des voisinages \'etales de $\overline{s}$ par la r\`egle
$(U, \overline{u}') \mapsto (U, \overline{u}' \circ a)$. Nous obtenons donc
une application canonique
$$
\mathcal{F}_{\overline{s}'}
=
\colim_{(U, \overline{u}')} \mathcal{F}(U)
\longrightarrow
\colim_{(U, \overline{u})} \mathcal{F}(U)
=
\mathcal{F}_{\overline{s}}
$$
donn\'ee par Catégories, lemme \ref{categories-lemma-functorial-colimit}. En utilisant la
description des \'el\'ements des fibres par des triplets, elle envoie l'\'el\'ement de
$\mathcal{F}_{\overline{s}'}$ repr\'esent\'e par le triplet
$(U, \overline{u}', \sigma)$ sur l'\'el\'ement de $\mathcal{F}_{\overline{s}}$
repr\'esent\'e par le triplet $(U, \overline{u}' \circ a, \sigma)$.
Comme le foncteur ci-dessus est manifestement une \'equivalence, nous concluons que cette
application canonique est un isomorphisme de foncteurs fibres.

\medskip\noindent
Assurons-nous que l'application de fibres correspondant \`a $a$ est orient\'ee
dans le bon sens. Notons que ce qui pr\'ec\`ede signifie, d'apr\`es
Sites, définition \ref{sites-definition-morphism-points},
que $a$ d\'efinit un morphisme $a : p \to p'$ entre les points $p, p'$ du
site $S_\etale$ associ\'es \`a $\overline{s}, \overline{s}'$ par le
lemme \ref{lemma-stalk-gives-point}. Il existe des morphismes plus g\'en\'eraux de
points (correspondant \`a des sp\'ecialisations de points de $S$), que nous
d\'ecrirons plus loin et qui ne seront pas des isomorphismes, voir la section
\ref{section-specialization}.
\end{remark}

\begin{lemma}
\label{lemma-stalk-exact}
Soit $S$ un sch\'ema. Soit $\overline{s}$ un point g\'eom\'etrique de $S$.
\begin{enumerate}
\item Le foncteur fibre
$\textit{PAb}(S_\etale) \to \textit{Ab}$,
$\mathcal{F}  \mapsto  \mathcal{F}_{\overline{s}}$
est exact.
\item On a $(\mathcal{F}^\#)_{\overline{s}} = \mathcal{F}_{\overline{s}}$
pour tout pr\'efaisceau d'ensembles $\mathcal{F}$ sur $S_\etale$.
\item Le foncteur
$\textit{Ab}(S_\etale) \to \textit{Ab}$,
$\mathcal{F} \mapsto \mathcal{F}_{\overline{s}}$ est exact.
\item De m\^eme, les foncteurs
$\textit{PSh}(S_\etale) \to \textit{Ensembles}$ et
$\Sh(S_\etale) \to \textit{Ensembles}$ donn\'es par le foncteur fibre
$\mathcal{F} \mapsto \mathcal{F}_{\overline{s}}$ sont exacts (voir
Catégories, définition \ref{categories-definition-exact})
et commutent aux limites inductives arbitraires.
\end{enumerate}
\end{lemma}

\begin{proof}
Avant d'indiquer comment le d\'emontrer par des arguments directs,
notons que le r\'esultat d\'ecoule de la th\'eorie g\'en\'erale de
Modules sur les sites, section \ref{sites-modules-section-stalks}.
Cela est vrai parce que $\mathcal{F} \mapsto \mathcal{F}_{\overline{s}}$
provient d'un point du petit site \'etale de $S$, voir le
lemme \ref{lemma-stalk-gives-point}.
Nous ne donnerons une preuve directe que de (1), (2) et (3), et omettrons
une preuve directe de (4).

\medskip\noindent
L'exactitude comme foncteur sur $\textit{PAb}(S_\etale)$ d\'ecoule formellement du
fait que les limites inductives filtrantes commutent \`a toutes les limites inductives et aux
limites finies. L'identification des fibres dans (2) se fait au moyen de l'application
$$
\kappa :
\mathcal{F}_{\overline{s}}
\longrightarrow
(\mathcal{F}^\#)_{\overline{s}}
$$
induite par le morphisme naturel $\mathcal{F}\to \mathcal{F}^\#$, voir le
th\'eor\`eme \ref{theorem-sheafification}.
Nous affirmons que cette application est un isomorphisme de groupes ab\'eliens. Nous montrerons
l'injectivit\'e et omettrons la preuve de la surjectivit\'e.

\medskip\noindent
Soit $\sigma\in \mathcal{F}_{\overline{s}}$.
Il existe un voisinage \'etale
$(U, \overline{u})\to (S, \overline{s})$ tel que $\sigma$ soit l'image d'une
section $s \in \mathcal{F}(U)$. Si $\kappa(\sigma) = 0$ dans
$(\mathcal{F}^\#)_{\overline{s}}$, alors il existe un morphisme de voisinages \'etales
$(U', \overline{u}')\to (U, \overline{u})$ tel que
$s|_{U'}$ soit nul dans $\mathcal{F}^\#(U')$. Il s'ensuit qu'il
existe un recouvrement \'etale
$\{U_i'\to U'\}_{i\in I}$ tel que $s|_{U_i'}=0$ dans
$\mathcal{F}(U_i')$ pour tout $i$. Par le lemme \ref{lemma-geometric-lift-to-cover},
il existe $i \in I$ et un morphisme
$\overline{u}_i': \overline{s} \to U_i'$ tels que
$(U_i', \overline{u}_i') \to (U', \overline{u}')\to (U, \overline{u})$
soient des morphismes de voisinages \'etales. Ainsi $\sigma = 0$
puisque $(U_i', \overline{u}_i') \to (U, \overline{u})$
est un morphisme de voisinages \'etales tel que
l'on ait $s|_{U'_i}=0$. Cela d\'emontre que $\kappa$ est injectif.

\medskip\noindent
Pour montrer que le foncteur $\textit{Ab}(S_\etale) \to \textit{Ab}$ est
exact, consid\'erons une suite exacte courte quelconque dans $\textit{Ab}(S_\etale)$ :
$
0\to \mathcal{F}\to \mathcal{G}\to \mathcal H \to 0.
$
Elle nous donne la suite exacte de pr\'efaisceaux
$$
0 \to \mathcal{F} \to \mathcal{G} \to \mathcal H \to
\mathcal H/^p\mathcal{G} \to 0,
$$
o\`u $/^p$ d\'esigne le quotient dans $\textit{PAb}(S_\etale)$.
En prenant les fibres en
$\overline{s}$, nous voyons que $(\mathcal H /^p\mathcal{G})_{\bar{s}} =
(\mathcal H /\mathcal{G})_{\bar{s}} = 0$, puisque le faisceau associ\'e \`a
$\mathcal H/^p\mathcal{G}$ est $0$.
Par cons\'equent,
$$
0\to \mathcal{F}_{\overline{s}} \to \mathcal{G}_{\overline{s}} \to
\mathcal{H}_{\overline{s}} \to 0 = (\mathcal H/^p\mathcal{G})_{\overline{s}}
$$
est exacte, puisque prendre les fibres est exact comme foncteur sur les pr\'efaisceaux.
\end{proof}

\begin{theorem}
\label{theorem-exactness-stalks}
Soit $S$ un sch\'ema.
Un morphisme $a : \mathcal{F} \to \mathcal{G}$ de faisceaux d'ensembles est injectif
(resp.\ surjectif) si et seulement si le morphisme sur les fibres
$a_{\overline{s}} : \mathcal{F}_{\overline{s}} \to \mathcal{G}_{\overline{s}}$
est injectif (resp.\ surjectif) pour tout point g\'eom\'etrique de $S$.
Une suite de faisceaux ab\'eliens sur $S_\etale$ est exacte
si et seulement si elle est exacte sur toutes les fibres aux points g\'eom\'etriques de $S$.
\end{theorem}

\begin{proof}
La n\'ecessit\'e de l'exactitude sur les fibres d\'ecoule du
lemme \ref{lemma-stalk-exact}.
R\'eciproquement, il suffit de montrer qu'un morphisme de faisceaux est surjectif
(respectivement injectif) si et seulement s'il est surjectif (respectivement
injectif) sur toutes les fibres. Nous le d\'emontrons dans le cas de la surjectivit\'e et omettons
la preuve dans le cas de l'injectivit\'e.

\medskip\noindent
Soit $\alpha : \mathcal{F} \to \mathcal{G}$ un morphisme de faisceaux tel
que $\mathcal{F}_{\overline{s}} \to \mathcal{G}_{\overline{s}}$
soit surjectif pour tout point g\'eom\'etrique. Fixons
$U \in \Ob(S_\etale)$
et $s \in \mathcal{G}(U)$. Pour tout $u \in U$, choisissons un
$\overline{u} \to U$ situ\'e au-dessus de $u$ et un voisinage \'etale
$(V_u , \overline{v}_u) \to (U, \overline{u})$ tel que
$s|_{V_u} = \alpha(s_{V_u})$ pour un certain
$s_{V_u} \in \mathcal{F}(V_u)$.
C'est possible puisque $\alpha$ est surjectif sur les
fibres. Alors $\{V_u \to U\}_{u \in U}$
est un recouvrement \'etale sur lequel les restrictions de $s$
appartiennent \`a l'image du morphisme $\alpha$.
Ainsi, $\alpha$ est surjectif, voir
Sites, section \ref{sites-section-sheaves-injective}.
\end{proof}

\begin{remarks}
\label{remarks-enough-points}
Quelques remarques sur les points des sites g\'eom\'etriques.
\begin{enumerate}
\item Le th\'eor\`eme \ref{theorem-exactness-stalks} dit que la famille de points
de $S_\etale$ donn\'ee par les points g\'eom\'etriques de $S$
(lemme \ref{lemma-stalk-gives-point}) est conservative, voir
Sites, définition \ref{sites-definition-enough-points}.
En particulier, $S_\etale$ a assez de points.
\item Supposons que $\mathcal{F}$ soit un faisceau sur le grand site \'etale
\label{item-stalks-big}
de $S$. Soit $T \to S$ un objet du grand site \'etale de $S$,
et soit $\overline{t}$ un point g\'eom\'etrique de $T$. Nous d\'efinissons alors
$\mathcal{F}_{\overline{t}}$ comme la fibre
de la restriction $\mathcal{F}|_{T_\etale}$ de $\mathcal{F}$
au petit site \'etale de $T$. Autrement dit, nous pouvons d\'efinir
la fibre de $\mathcal{F}$ en tout point g\'eom\'etrique de tout
sch\'ema $T/S \in \Ob((\Sch/S)_\etale)$.
\item Le grand site \'etale de $S$ a lui aussi assez de points, en
consid\'erant tous les points g\'eom\'etriques de tous les objets de ce site, voir
(\ref{item-stalks-big}).
\end{enumerate}
\end{remarks}

\noindent
Le lemme suivant devrait \^etre omis en premi\`ere lecture.

\begin{lemma}
\label{lemma-points-small-etale-site}
Soit $S$ un sch\'ema.
\begin{enumerate}
\item Soit $p$ un point du petit site \'etale
$S_\etale$ de $S$ donn\'e par un foncteur
$u : S_\etale \to \textit{Ensembles}$.
Alors il existe un point g\'eom\'etrique $\overline{s}$ de $S$ tel que
$p$ soit isomorphe au point de $S_\etale$ associ\'e \`a
$\overline{s}$ dans le
lemme \ref{lemma-stalk-gives-point}.
\item Soit $p : \Sh(pt) \to \Sh(S_\etale)$ un point
du petit topos \'etale de $S$. Alors $p$ provient d'un point g\'eom\'etrique
de $S$, c'est-\`a-dire que le foncteur fibre $\mathcal{F} \mapsto \mathcal{F}_p$
est isomorphe \`a un foncteur fibre tel que d\'efini dans la
d\'efinition \ref{definition-stalk}.
\end{enumerate}
\end{lemma}

\begin{proof}
D'apr\`es
Sites, lemme \ref{sites-lemma-point-site-topos},
il existe une correspondance bijective entre les points du site et les points
du topos associ\'e ; il suffit donc de d\'emontrer (1).
D'apr\`es
Sites, Proposition \ref{sites-proposition-point-limits},
le foncteur $u$ a les propri\'et\'es suivantes :
(a) $u(S) = \{*\}$, (b) $u(U \times_V W) = u(U) \times_{u(V)} u(W)$, et
(c) si $\{U_i \to U\}$ est un recouvrement \'etale, alors
$\coprod u(U_i) \to u(U)$ est surjectif.
En particulier, si $U' \subset U$ est un sous-sch\'ema ouvert,
alors $u(U') \subset u(U)$. De plus, d'apr\`es
Sites, lemme \ref{sites-lemma-point-site-topos},
nous pouvons \'ecrire $u(U) = p^{-1}(h_U^\#)$ ; autrement dit, $u(U)$ est la
fibre du faisceau repr\'esentable $h_U$. Si
$U = V \amalg W$, alors nous voyons que $h_U = (h_V \amalg h_W)^\#$ et obtenons
$u(U) = u(V) \amalg u(W)$ puisque $p^{-1}$ est exact.

\medskip\noindent
Consid\'erons la restriction de $u$ \`a $S_{Zar}$. D'apr\`es
Sites, Exemples \ref{sites-example-point-topological} et
\ref{sites-example-point-topology},
il existe un unique point $s \in S$ tel que, pour $S' \subset S$ ouvert, nous
ayons $u(S') = \{*\}$ si $s \in S'$ et $u(S') = \emptyset$ si $s \not \in S'$.
Notons que si $\varphi : U \to S$ est un objet de $S_\etale$, alors
$\varphi(U) \subset S$ est ouvert (voir la
proposition \ref{proposition-etale-morphisms})
et $\{U \to \varphi(U)\}$ est un recouvrement \'etale. Nous concluons donc que
$u(U) = \emptyset \Leftrightarrow s \not \in \varphi(U)$.

\medskip\noindent
Choisissons un point g\'eom\'etrique $\overline{s} : \overline{s} \to S$ au-dessus de $s$, voir la
d\'efinition \ref{definition-geometric-point}
pour l'abus de notation usuel. Supposons que $\varphi : U \to S$ soit un objet
de $S_\etale$ avec $U$ affine. Notons que $\varphi$ est s\'epar\'e et
que la fibre $U_s$ de $\varphi$ au-dessus de $s$ est un sch\'ema affine sur
$\Spec(\kappa(s))$ et qu'elle est le spectre d'un produit fini
d'extensions finies s\'eparables $k_i$ de $\kappa(s)$. Nous pouvons donc appliquer
Morphismes étales, lemme \ref{etale-lemma-etale-etale-local-technical}
pour obtenir un voisinage \'etale $(V, \overline{v})$ de $(S, \overline{s})$
tel que
$$
U \times_S V = U_1 \amalg \ldots \amalg U_n \amalg W
$$
o\`u $U_i \to V$ est un isomorphisme et $W$ n'a aucun point au-dessus de
$\overline{v}$. Nous concluons donc que
$$
u(U) \times u(V) =
u(U \times_S V) =
u(U_1) \amalg \ldots \amalg u(U_n) \amalg u(W)
$$
et, bien s\^ur, aussi $u(U_i) = u(V)$. Apr\`es avoir un peu r\'etr\'eci $V$, nous pouvons
supposer que $V$ a exactement un point au-dessus de $s$ et que, par cons\'equent, $W$ n'a aucun
point au-dessus de $s$. D'apr\`es ce qui pr\'ec\`ede, cela donne alors $u(W) = \emptyset$.
Nous obtenons donc
$$
u(U) \times u(V) =
u(U_1) \amalg \ldots \amalg u(U_n) =
\coprod\nolimits_{i = 1, \ldots, n} u(V)
$$
o\`u les signes d'\'egalit\'e sont compatibles avec les applications vers l'ensemble $u(V)$.
Notons que $u(V) \not = \emptyset$ puisque $s$ est dans l'image de $V \to S$.
En particulier, nous voyons que, dans cette situation, $u(U)$ est un
ensemble fini \`a $n$ \'el\'ements : en effet, le membre de gauche a pour fibre $u(U)$ au-dessus de
tout \'el\'ement de $u(V)$ et le membre de droite a une fibre de cardinal $n$.

\medskip\noindent
Consid\'erons la limite
$$
\lim_{(V, \overline{v})} u(V)
$$
sur la cat\'egorie des voisinages \'etales $(V, \overline{v})$ de
$\overline{s}$. Il est clair que nous obtenons la m\^eme valeur en prenant
la limite sur la sous-cat\'egorie des $(V, \overline{v})$ tels que $V$ soit affine.
D'apr\`es le paragraphe pr\'ec\'edent (appliqu\'e en \'echangeant les r\^oles de $V$ et $U$),
nous voyons que, dans ce cas, $u(V)$ est toujours un ensemble fini non vide.
De plus, la cat\'egorie d'indexation est cofiltrante, voir le
lemme \ref{lemma-cofinal-etale}.
Par cons\'equent, d'apr\`es
Catégories, section \ref{categories-section-codirected-limits},
la limite est non vide. Choisissons un \'el\'ement $x$ de cette limite.
Cela signifie que nous obtenons un $x_{V, \overline{v}} \in u(V)$ pour
chaque voisinage \'etale $(V, \overline{v})$ de $(S, \overline{s})$
tel que, pour tout morphisme de voisinages \'etales
$\varphi : (V', \overline{v}') \to (V, \overline{v})$, on ait
$u(\varphi)(x_{V', \overline{v}'}) = x_{V, \overline{v}}$.

\medskip\noindent
Nous utiliserons le choix de $x$ pour construire une application bijective fonctorielle
$$
c : |U_{\overline{s}}| \longrightarrow u(U)
$$
pour $U \in \Ob(S_\etale)$, ce qui conclura la preuve. Voir le
lemme \ref{lemma-stalk-gives-point}
et sa preuve pour une description de $|U_{\overline{s}}|$.
Nous affirmons d'abord qu'il suffit de construire l'application pour $U$ affine.
Nous omettons la preuve de cette affirmation.
Supposons que $U \to S$ soit un objet de $S_\etale$, avec $U$ affine, et soit
$\overline{u} : \overline{s} \to U$ un \'el\'ement de $|U_{\overline{s}}|$.
Choisissons $(V, \overline{v})$ tel que $U \times_S V$ se d\'ecompose
comme dans le troisi\`eme paragraphe de la preuve.
Alors le couple $(\overline{u}, \overline{v})$ donne un point g\'eom\'etrique de
$U \times_S V$ au-dessus de $\overline{v}$ et d\'etermine l'une des
composantes $U_i$ de $U \times_S V$. Plus pr\'ecis\'ement, il existe
une section $\sigma : V \to U \times_S V$ de la projection $\text{pr}_U$
telle que $(\overline{u}, \overline{v}) = \sigma \circ \overline{v}$. Posons
$c(\overline{u}) = u(\text{pr}_U)(u(\sigma)(x_{V, \overline{v}})) \in u(U)$.
Nous devons v\'erifier que cela est ind\'ependant du choix de $(V, \overline{v})$. D'apr\`es le
lemme \ref{lemma-cofinal-etale},
la cat\'egorie des voisinages \'etales est cofiltrante.
Il suffit donc de montrer
que, \'etant donn\'es un morphisme de voisinages \'etales
$\varphi : (V', \overline{v}') \to (V, \overline{v})$ et le choix d'une
section $\sigma' : V' \to U \times_S V'$ de la projection telle que
$(\overline{u}, \overline{v'}) = \sigma' \circ \overline{v}'$,
on a $u(\sigma')(x_{V', \overline{v}'}) = u(\sigma)(x_{V, \overline{v}})$.
Consid\'erons le diagramme
$$
\xymatrix{
V' \ar[d]^{\sigma'} \ar[r]_\varphi & V \ar[d]^\sigma \\
U \times_S V' \ar[r]^{1 \times \varphi} &
U \times_S V
}
$$
Il se peut que ce diagramme ne soit pas commutatif. La raison en est
que les sch\'emas $V'$ et $V$ ne sont peut-\^etre pas connexes, et donc que
les d\'ecompositions utilis\'ees ci-dessus pour construire $\sigma'$ et $\sigma$ ne sont
peut-\^etre pas uniques. Mais nous savons que
$\sigma \circ \varphi \circ \overline{v}' =
(1 \times \varphi) \circ \sigma' \circ \overline{v}'$
par construction. Par cons\'equent, puisque $U \times_S V$ est \'etale sur $S$,
il existe un voisinage ouvert
$V'' \subset V'$ de $\overline{v'}$ tel que le diagramme
soit commutatif apr\`es restriction \`a $V''$, voir
Morphismes, lemme \ref{morphisms-lemma-value-at-one-point}.
Cela signifie que nous pouvons compl\'eter le diagramme ci-dessus en
$$
\xymatrix{
V'' \ar[r] \ar[d]^{\sigma'|_{V''}} &
V' \ar[d]^{\sigma'} \ar[r]_\varphi &
V \ar[d]^\sigma \\
U \times_S V'' \ar[r] &
U \times_S V' \ar[r]^{1 \times \varphi} &
U \times_S V
}
$$
de sorte que le carr\'e de gauche et le rectangle ext\'erieur soient commutatifs.
Comme $u$ est un foncteur, cela implique que
$x_{V'', \overline{v}'}$ s'envoie sur le m\^eme \'el\'ement de
$u(U \times_S V)$ quel que soit le chemin suivi dans le
diagramme. D'autre part, il s'envoie sur les \'el\'ements
$x_{V', \overline{v}'}$ et $x_{V, \overline{v}}$ dans
$u(V')$ et $u(V)$. Cela implique l'\'egalit\'e voulue
$u(\sigma')(x_{V', \overline{v}'}) = u(\sigma)(x_{V, \overline{v}})$.

\medskip\noindent
De mani\`ere analogue, on d\'emontre que la construction
$c : |U_{\overline{s}}| \to u(U)$ est fonctorielle en $U$ ;
nous omettons les d\'etails. Enfin, d'apr\`es les r\'esultats du
troisi\`eme paragraphe, il est clair que l'application $c$ est bijective,
ce qui termine la preuve du lemme.
\end{proof}







\section{Points dans d'autres topologies}
\label{section-points-topologies}

\noindent
Dans cette section, nous discutons bri\`evement de l'existence de points pour certains
sites autres que le site \'etale d'un sch\'ema. Nous renvoyons \`a
Sites, section \ref{sites-section-sites-enough-points}
et \`a
Topologies, section \ref{topologies-section-procedure} et suivantes
pour la terminologie employ\'ee dans cette section.
Tous les sites g\'eom\'etriques ont assez de points.

\begin{lemma}
\label{lemma-points-fppf}
Soit $S$ un sch\'ema. Tous les sites suivants ont assez de points :
$S_{affine, Zar}$, $S_{Zar}$, $S_{affine, \etale}$, $S_\etale$,
$(\Sch/S)_{Zar}$, $(\textit{Aff}/S)_{Zar}$,
$(\Sch/S)_\etale$, $(\textit{Aff}/S)_\etale$,
$(\Sch/S)_{smooth}$, $(\textit{Aff}/S)_{smooth}$,
$(\Sch/S)_{syntomic}$, $(\textit{Aff}/S)_{syntomic}$,
$(\Sch/S)_{fppf}$ et $(\textit{Aff}/S)_{fppf}$.
\end{lemma}

\begin{proof}
Pour chacun des grands sites, le topos associ\'e est \'equivalent au
topos d\'efini par le site $(\textit{Aff}/S)_\tau$, voir
Topologies, lemmes \ref{topologies-lemma-affine-big-site-Zariski},
\ref{topologies-lemma-affine-big-site-etale},
\ref{topologies-lemma-affine-big-site-smooth},
\ref{topologies-lemma-affine-big-site-syntomic} et
\ref{topologies-lemma-affine-big-site-fppf}.
Le r\'esultat pour les sites $(\textit{Aff}/S)_\tau$ d\'ecoule imm\'ediatement
du r\'esultat de Deligne,
Sites, lemme \ref{sites-lemma-criterion-points}.

\medskip\noindent
Le r\'esultat pour $S_{Zar}$ est clair. Celui pour $S_{affine, Zar}$
d\'ecoule du r\'esultat de Deligne. Le r\'esultat pour $S_\etale$
d\'ecoule soit (de la preuve) du
th\'eor\`eme \ref{theorem-exactness-stalks},
soit de Topologies, lemme \ref{topologies-lemma-alternative}
et du r\'esultat de Deligne appliqu\'e \`a $S_{affine, \etale}$.
\end{proof}

\noindent
Le lemme ci-dessus garantit l'existence de points, mais ne nous
dit pas \`a quoi ils ressemblent. Nous pouvons construire explicitement
{\it certains} points comme suit.
Supposons que $\overline{s} : \Spec(k) \to S$ soit un point
g\'eom\'etrique, avec $k$ alg\'ebriquement clos. Consid\'erons le foncteur
$$
u : (\Sch/S)_{fppf} \longrightarrow \textit{Ensembles},
\quad
u(U) = U(k) = \Mor_S(\Spec(k), U).
$$
Notons que $U \mapsto U(k)$ commute aux limites finies puisque
$S(k) = \{\overline{s}\}$ et
$(U_1 \times_U U_2)(k) = U_1(k) \times_{U(k)} U_2(k)$.
De plus, si $\{U_i \to U\}$ est un recouvrement fppf, alors
$\coprod U_i(k) \to U(k)$ est surjectif.
D'apr\`es
Sites, Proposition \ref{sites-proposition-point-limits},
nous voyons que $u$ d\'efinit un point $p$ de $(\Sch/S)_{fppf}$ dont les
fibres sont
$$
\mathcal{F}_p = \colim_{(U, x)} \mathcal{F}(U)
$$
o\`u la limite inductive porte, comme d'habitude, sur les couples $U \to S$, $x \in U(k)$.
Mais... cette cat\'egorie a un objet initial, \`a savoir
$(\Spec(k), \text{id})$ ; nous voyons donc que
$$
\mathcal{F}_p = \mathcal{F}(\Spec(k))
$$
ce qui n'est gu\`ere int\'eressant ! En fait, en g\'en\'eral, ces points ne
forment pas une famille conservative de points. Un type de point plus int\'eressant
est d\'ecrit dans la remarque suivante.

\begin{remark}
\label{remark-points-fppf-site}
\begin{reference}
Ceci est discut\'e dans \cite{Schroeer}.
\end{reference}
Soit $S = \Spec(A)$ un sch\'ema affine. Soit $(p, u)$ un point du
site $(\textit{Aff}/S)_{fppf}$, voir
Sites, sections \ref{sites-section-points} et
\ref{sites-section-construct-points}. Soit $B = \mathcal{O}_p$ la fibre
du faisceau structural au point $p$. Rappelons que
$$
B = \colim_{(U, x)} \mathcal{O}(U) =
\colim_{(\Spec(C), x_C)} C
$$
o\`u $x_C \in u(\Spec(C))$. Il peut arriver que
$\Spec(B)$ soit un objet de $(\textit{Aff}/S)_{fppf}$
et qu'il existe un \'el\'ement $x_B \in u(\Spec(B))$ qui s'envoie sur
le syst\`eme compatible $x_C$. Dans ce cas, le syst\`eme de voisinages
a un objet initial et il s'ensuit que
$\mathcal{F}_p = \mathcal{F}(\Spec(B))$ pour tout faisceau $\mathcal{F}$
sur $(\textit{Aff}/S)_{fppf}$. Il est imm\'ediat
de voir que, si $\mathcal{F} \mapsto \mathcal{F}(\Spec(B))$ d\'efinit un point
de $\Sh((\textit{Aff}/S)_{fppf})$, alors
$B$ doit \^etre une $A$-alg\`ebre locale telle que, pour tout morphisme d'anneaux
fid\`element plat et de pr\'esentation finie $B \to B'$, il existe une section $B' \to B$.
R\'eciproquement, pour toute $A$-alg\`ebre $B$ de ce type, le foncteur
$\mathcal{F} \mapsto \mathcal{F}(\Spec(B))$ est le foncteur fibre
d'un point. Nous omettons les d\'etails. On ne sait pas clairement \`a quoi ressemble un point quelconque du
site $(\textit{Aff}/S)_{fppf}$.
\end{remark}











\section{Supports des faisceaux ab\'eliens}
\label{section-support}

\noindent
Nous commen\c{c}ons par les supports des sections locales.

\begin{lemma}
\label{lemma-support-subsheaf-final}
Soit $S$ un sch\'ema. Soit $\mathcal{F}$ un sous-faisceau de l'objet
final du topos \'etale de $S$ (voir
Sites, exemple \ref{sites-example-singleton-sheaf}).
Alors il existe un unique ouvert
$W \subset S$ tel que $\mathcal{F} = h_W$.
\end{lemma}

\begin{proof}
La condition signifie que $\mathcal{F}(U)$ est un singleton ou est
vide pour tout $\varphi : U \to S$ dans $\Ob(S_\etale)$.
En particulier, les sections locales se recollent toujours. Si
$\mathcal{F}(U) \not = \emptyset$, alors
$\mathcal{F}(\varphi(U)) \not = \emptyset$ parce que
$\{\varphi : U \to \varphi(U)\}$ est un recouvrement.
Nous pouvons donc prendre
$W = \bigcup_{\varphi : U \to S, \mathcal{F}(U) \not = \emptyset} \varphi(U)$.
\end{proof}

\begin{lemma}
\label{lemma-zero-over-image}
Soit $S$ un sch\'ema.
Soit $\mathcal{F}$ un faisceau ab\'elien sur $S_\etale$.
Soit $\sigma \in \mathcal{F}(U)$ une section locale.
Il existe un ouvert $W \subset U$ tel que
\begin{enumerate}
\item $W \subset U$ soit le plus grand ouvert de Zariski de $U$ tel
que $\sigma|_W = 0$,
\item pour tout $\varphi : V \to U$ dans $S_\etale$, on ait
$$
\sigma|_V = 0 \Leftrightarrow \varphi(V) \subset W,
$$
\item pour tout point g\'eom\'etrique $\overline{u}$ de $U$, on ait
$$
(U, \overline{u}, \sigma) = 0\text{ dans }\mathcal{F}_{\overline{s}}
\Leftrightarrow
\overline{u} \in W
$$
o\`u $\overline{s} = (U \to S) \circ \overline{u}$.
\end{enumerate}
\end{lemma}

\begin{proof}
Comme $\mathcal{F}$ est un faisceau pour la topologie \'etale, la restriction de
$\mathcal{F}$ \`a $U_{Zar}$ est un faisceau sur $U$ pour la topologie de Zariski.
Il existe donc un ouvert de Zariski $W$ ayant la propri\'et\'e (1), voir
Modules, lemme \ref{modules-lemma-support-section-closed}. Soit
$\varphi : V \to U$ un morphisme de $S_\etale$. Notons que
$\varphi(V) \subset U$ est un ouvert et que
$\{V \to \varphi(V)\}$ est un recouvrement \'etale. Ainsi, si
$\sigma|_V = 0$, alors la condition de faisceau pour $\mathcal{F}$ montre
que $\sigma|_{\varphi(V)} = 0$. Cela d\'emontre (2).
Pour d\'emontrer (3), nous devons montrer que si $(U, \overline{u}, \sigma)$
d\'efinit l'\'el\'ement nul de $\mathcal{F}_{\overline{s}}$, alors
$\overline{u} \in W$. Cela est vrai parce que l'hypoth\`ese signifie
qu'il existe un morphisme de voisinages \'etales
$(V, \overline{v}) \to (U, \overline{u})$ tel que
$\sigma|_V = 0$. D'apr\`es (2), nous voyons donc que $V \to U$ est \`a valeurs dans $W$, et
par cons\'equent que $\overline{u} \in W$.
\end{proof}

\noindent
Soit $S$ un sch\'ema. Soit $s \in S$.
Soit $\mathcal{F}$ un faisceau sur $S_\etale$. D'apr\`es la
remarque \ref{remark-map-stalks},
la classe d'isomorphisme de la fibre du faisceau $\mathcal{F}$
en un point g\'eom\'etrique au-dessus de $s$ est bien d\'efinie.

\begin{definition}
\label{definition-support}
Soit $S$ un sch\'ema.
Soit $\mathcal{F}$ un faisceau ab\'elien sur $S_\etale$.
\begin{enumerate}
\item Le {\it support de $\mathcal{F}$} est l'ensemble des
points $s \in S$ tels que $\mathcal{F}_{\overline{s}} \not = 0$
pour tout (ou un) point g\'eom\'etrique $\overline{s}$ au-dessus de $s$.
\item Soit $\sigma \in \mathcal{F}(U)$ une section.
Le {\it support de $\sigma$} est le ferm\'e $U \setminus W$, o\`u
$W \subset U$ est le plus grand ouvert de $U$ sur lequel $\sigma$
est nulle (voir le
lemme \ref{lemma-zero-over-image}).
\end{enumerate}
\end{definition}

\noindent
En g\'en\'eral, le support d'un faisceau ab\'elien n'est pas ferm\'e.
Par exemple, supposons que $S = \mathbf{A}^1_{\mathbf{C}}$.
Soit $i_t : \Spec(\mathbf{C}) \to S$ l'inclusion du
point $t \in \mathbf{C}$.
Nous verrons plus loin que $\mathcal{F}_t = i_{t, *}(\mathbf{Z}/2\mathbf{Z})$
est un faisceau ab\'elien dont le support est exactement $\{t\}$, voir la
section \ref{section-closed-immersions}.
Alors
$$
\bigoplus\nolimits_{n \in \mathbf{N}} \mathcal{F}_n
$$
est un faisceau ab\'elien de support $\{1, 2, 3, \ldots\} \subset S$.
Cela est vrai parce que prendre les fibres commute aux limites inductives, voir le
lemme \ref{lemma-stalk-exact}.
Nous avons ainsi un exemple de faisceau ab\'elien dont le support n'est pas ferm\'e.
Voici quelques faits fondamentaux sur les supports des faisceaux et des sections.

\begin{lemma}
\label{lemma-support-section-closed}
Soit $S$ un sch\'ema.
Soit $\mathcal{F}$ un faisceau ab\'elien sur $S_\etale$.
Soient $U \in \Ob(S_\etale)$ et $\sigma \in \mathcal{F}(U)$.
\begin{enumerate}
\item Le support de $\sigma$ est ferm\'e dans $U$.
\item Le support de $\sigma + \sigma'$ est contenu dans la r\'eunion des
supports de $\sigma, \sigma' \in \mathcal{F}(U)$.
\item Si $\varphi : \mathcal{F} \to \mathcal{G}$ est un morphisme de
faisceaux ab\'eliens sur $S_\etale$, alors le support de $\varphi(\sigma)$
est contenu dans le support de $\sigma \in \mathcal{F}(U)$.
\item Le support de $\mathcal{F}$ est la r\'eunion des images des
supports de toutes les sections locales de $\mathcal{F}$.
\item Si $\mathcal{F} \to \mathcal{G}$ est surjectif, alors le support
de $\mathcal{G}$ est un sous-ensemble du support de $\mathcal{F}$.
\item Si $\mathcal{F} \to \mathcal{G}$ est injectif, alors le support
de $\mathcal{F}$ est un sous-ensemble du support de $\mathcal{G}$.
\end{enumerate}
\end{lemma}

\begin{proof}
L'assertion (1) est vraie par d\'efinition.
Les assertions (2) et (3) sont vraies parce qu'elles le sont pour les restrictions de
$\mathcal{F}$ et $\mathcal{G}$ \`a $U_{Zar}$, voir
Modules, lemme \ref{modules-lemma-support-section-closed}.
L'assertion (4) est une cons\'equence directe de la partie (3) du
lemme \ref{lemma-zero-over-image}.
Les assertions (5) et (6) d\'ecoulent des autres.
\end{proof}

\begin{lemma}
\label{lemma-support-sheaf-rings-closed}
Le support d'un faisceau d'anneaux sur $S_\etale$ est ferm\'e.
\end{lemma}

\begin{proof}
Cela est vrai parce que (selon nos conventions)
un anneau est $0$ si et seulement si
$1 = 0$ ; par cons\'equent, le support d'un faisceau d'anneaux
est le support de la section unit\'e.
\end{proof}




\section{Anneaux hens\'eliens}
\label{section-henselian-ring}

\noindent
Nous commen\c{c}ons par \'enoncer un th\'eor\`eme qui a d\'ej\`a \'et\'e utilis\'e de nombreuses fois
dans le projet Champs. Il existe de nombreuses versions de ce r\'esultat ; ici, nous
n'\'enon\c{c}ons que la version alg\'ebrique.

\begin{theorem}
\label{theorem-quasi-finite-etale-locally}
Soient $A\to B$ un morphisme d'anneaux de type fini et $\mathfrak p \subset A$ un
id\'eal premier. Alors il existe un morphisme d'anneaux \'etale $A \to A'$ et un id\'eal premier
$\mathfrak p' \subset A'$ au-dessus de $\mathfrak p$ tels que
\begin{enumerate}
\item
$\kappa(\mathfrak p) = \kappa(\mathfrak p')$,
\item
$ B \otimes_A A' = B_1\times \ldots \times B_r \times C$,
\item
$ A'\to B_i$ est fini et il existe un unique id\'eal premier $q_i\subset B_i$ au-dessus
de $\mathfrak p'$, et
\item toutes les composantes irr\'eductibles de la fibre
$\Spec(C \otimes_{A'} \kappa(\mathfrak p'))$ de $C$ au-dessus de $\mathfrak p'$
sont de dimension au moins 1.
\end{enumerate}
\end{theorem}

\begin{proof}
Voir Algèbre, lemme \ref{algebra-lemma-etale-makes-quasi-finite-finite}, ou
voir \cite[Th\'eor\`eme 18.12.1]{EGA4}. Pour de nombreuses versions en termes de
morphismes de sch\'emas, voir
Compléments sur les morphismes, section \ref{more-morphisms-section-etale-localization}.
\end{proof}

\noindent
Rappelons le lemme de Hensel.
Il existe de nombreuses versions de ce lemme. En voici deux :
\begin{enumerate}
\item[(f)] si $f\in \mathbf{Z}_p[T]$ est unitaire et
$f \bmod p = g_0 h_0$ avec $gcd(g_0, h_0) = 1$, alors $f$ se factorise
sous la forme $f = gh$ avec $\bar g = g_0$ et $\bar h = h_0$,
\item[(r)] si $f \in \mathbf{Z}_p[T]$ est unitaire, $a_0 \in \mathbf{F}_p$,
$\bar f(a_0) =0$ mais $\bar f'(a_0) \neq 0$,
alors il existe $a \in \mathbf{Z}_p$ tel que
$f(a) = 0$ et $\bar a = a_0$.
\end{enumerate}
Les deux versions sont vraies (nous le verrons plus loin). La premi\`ere version
demande des rel\`evements de factorisations en facteurs premiers entre eux,
et la seconde demande des rel\`evements de racines simples
modulo l'id\'eal maximal. Il se trouve que, pour un anneau local g\'en\'eral,
ces deux conditions sont \'equivalentes, et qu'elles sont
\'equivalentes \`a beaucoup d'autres conditions. Nous utilisons la propri\'et\'e de rel\`evement
des racines comme d\'efinition d'un anneau local hens\'elien, car c'est
souvent la plus facile \`a v\'erifier.

%10.01.09
\begin{definition}
\label{definition-henselian}
(Voir Algèbre, définition \ref{algebra-definition-henselian}.)
Un anneau local $(R, \mathfrak m, \kappa)$ est dit
{\it hens\'elien} si, pour tout
$f \in R[T]$ unitaire et tout $a_0 \in \kappa$ tel que
$\bar f(a_0) = 0$ et $\bar f'(a_0) \neq 0$, il existe
un $a \in R$ tel que $f(a) = 0$ et $a \bmod \mathfrak m = a_0$.
\end{definition}

\noindent
Les anneaux locaux complets constituent un bon exemple
d'anneaux locaux hens\'eliens \`a garder \`a l'esprit.
Rappelons
(Algèbre, définition \ref{algebra-definition-complete-local-ring})
qu'un anneau local complet est un anneau local $(R, \mathfrak m)$ tel que
$R \cong \lim_n R/\mathfrak m^n$, c'est-\`a-dire complet et s\'epar\'e
pour la topologie $\mathfrak m$-adique.

\begin{theorem}
\label{theorem-hensel}
Les anneaux locaux complets sont hens\'eliens.
\end{theorem}

\begin{proof}
M\'ethode de Newton. Voir
Algèbre, lemme \ref{algebra-lemma-complete-henselian}.
\end{proof}

\begin{theorem}
\label{theorem-henselian}
Soit $(R, \mathfrak m, \kappa)$ un anneau local. Les propri\'et\'es suivantes sont \'equivalentes :
\begin{enumerate}
\item $R$ est hens\'elien,
\item pour tout $f\in R[T]$ et toute factorisation $\bar f = g_0 h_0$ dans
$\kappa[T]$ avec $\gcd(g_0, h_0)=1$, il existe une factorisation $f = gh$ dans
$R[T]$ avec $\bar g = g_0$ et $\bar h = h_0$,
\item toute $R$-alg\`ebre finie $S$ est isomorphe \`a un produit fini
d'anneaux locaux finis sur $R$,
\item toute $R$-alg\`ebre de type fini $A$ est isomorphe \`a un produit
$A \cong A' \times C$ o\`u $A' \cong A_1 \times \ldots \times A_r$
est un produit de $R$-alg\`ebres locales finies et toutes les composantes
irr\'eductibles de $C \otimes_R \kappa$ sont de dimension au moins 1,
\item si $A$ est une $R$-alg\`ebre \'etale et $\mathfrak n$ un id\'eal maximal de
$A$ au-dessus de $\mathfrak m$ tel que $\kappa \cong A/\mathfrak n$, alors il
existe un isomorphisme $\varphi : A \cong R \times A'$ tel que
$\varphi(\mathfrak n) = \mathfrak m \times A' \subset R \times A'$.
\end{enumerate}
\end{theorem}

\begin{proof}
Il s'agit simplement d'une partie des r\'esultats de
Algèbre, lemme \ref{algebra-lemma-characterize-henselian}.
Notons que la partie (5) ci-dessus correspond \`a la partie (8) de
Algèbre, lemme \ref{algebra-lemma-characterize-henselian},
mais qu'elle est formul\'ee l\'eg\`erement diff\'eremment.
\end{proof}

\begin{lemma}
\label{lemma-finite-over-henselian}
Si $R$ est hens\'elien et $A$ est une $R$-alg\`ebre finie, alors $A$ est un produit
fini d'anneaux locaux hens\'eliens.
\end{lemma}

\begin{proof}
Voir
Algèbre, lemme \ref{algebra-lemma-finite-over-henselian}.
\end{proof}

\begin{definition}
\label{definition-strictly-henselian}
Un anneau local $R$ est dit {\it strictement hens\'elien} s'il est hens\'elien et si son
corps r\'esiduel est s\'eparablement clos.
\end{definition}

\begin{example}
\label{example-powerseries}
Dans le cas $R = \mathbf{C}[[t]]$, les $R$-alg\`ebres \'etales sont des produits finis
de l'extension triviale $R \to R$ et des extensions
$R \to R[X, X^{-1}]/(X^n-t)$.
Ces derni\`eres se factorisent par l'ouvert $D(t) \subset \Spec(R)$, de sorte que tout
recouvrement \'etale peut \^etre raffin\'e par le recouvrement
$\{\text{id} : \Spec(R) \to \Spec(R)\}$. Nous verrons ci-dessous que
c'est un fait assez g\'en\'eral sur les recouvrements \'etales des spectres d'anneaux
hens\'eliens. Cela montrera que la cohomologie \'etale sup\'erieure du spectre d'un anneau
strictement hens\'elien est nulle.
\end{example}

\begin{theorem}
\label{theorem-henselization}
Soient $(R, \mathfrak m, \kappa)$ un anneau local et
$\kappa\subset\kappa^{sep}$ une cl\^oture alg\'ebrique s\'eparable.
Il existe des morphismes locaux plats canoniques d'anneaux $R \to R^h \to R^{sh}$ tels que
\begin{enumerate}
\item $R^h$, $R^{sh}$ soient des limites inductives filtrantes de $R$-alg\`ebres \'etales,
\item $R^h$ soit hens\'elien, $R^{sh}$ soit strictement hens\'elien,
\item $\mathfrak m R^h$ (resp.\ $\mathfrak m R^{sh}$) soit
l'id\'eal maximal de $R^h$ (resp.\ de $R^{sh}$), et
\item $\kappa = R^h/\mathfrak m R^h$, et
$\kappa^{sep} = R^{sh}/\mathfrak m R^{sh}$ comme extensions de $\kappa$.
\end{enumerate}
\end{theorem}

\begin{proof}
La structure de $R^h$ et de $R^{sh}$ est d\'ecrite dans
Algèbre, lemmes \ref{algebra-lemma-henselization} et
\ref{algebra-lemma-strict-henselization}.
\end{proof}

\noindent
Les anneaux construits dans le th\'eor\`eme \ref{theorem-henselization}
sont appel\'es respectivement {\it l'hens\'elisation} et
{\it l'hens\'elisation stricte} de l'anneau local $R$, voir
Algèbre, définition \ref{algebra-definition-henselization}.
De nombreuses propri\'et\'es de $R$ se refl\`etent dans son hens\'elisation (stricte),
voir Compléments d'algèbre,
section \ref{more-algebra-section-permanence-henselization}.




\section{Fibres du faisceau structural}
\label{section-stalks-structure-sheaf}

\noindent
Dans cette section, nous identifions la fibre du faisceau structural en un point
g\'eom\'etrique \`a l'hens\'elisation stricte de l'anneau local au point
``usuel'' correspondant.

\begin{lemma}
\label{lemma-describe-etale-local-ring}
\begin{slogan}
La fibre du faisceau structural d'un sch\'ema
pour la topologie \'etale est l'hens\'elisation stricte.
\end{slogan}
Soit $S$ un sch\'ema.
Soit $\overline{s}$ un point g\'eom\'etrique de $S$ au-dessus de $s \in S$.
Posons $\kappa = \kappa(s)$ et notons
$\kappa \subset \kappa^{sep} \subset \kappa(\overline{s})$
la cl\^oture alg\'ebrique s\'eparable de $\kappa$ dans $\kappa(\overline{s})$.
Alors il existe une identification canonique
$$
(\mathcal{O}_{S, s})^{sh}
\cong
(\mathcal{O}_S)_{\overline{s}}
$$
o\`u le membre de gauche est l'hens\'elisation stricte de l'anneau local
$\mathcal{O}_{S, s}$ telle qu'elle est d\'ecrite dans le
th\'eor\`eme \ref{theorem-henselization},
et le membre de droite est la fibre du faisceau structural
$\mathcal{O}_S$ sur $S_\etale$ au
point g\'eom\'etrique $\overline{s}$.
\end{lemma}

\begin{proof}
Soit $\Spec(A) \subset S$ un voisinage affine de $s$.
Soit $\mathfrak p \subset A$ l'id\'eal premier correspondant \`a $s$.
Avec ces choix, nous avons des isomorphismes canoniques
$\mathcal{O}_{S, s} = A_{\mathfrak p}$ et $\kappa(s) = \kappa(\mathfrak p)$.
Nous avons donc
$\kappa(\mathfrak p) \subset \kappa^{sep} \subset \kappa(\overline{s})$.
Rappelons que
$$
(\mathcal{O}_S)_{\overline{s}} =
\colim_{(U, \overline{u})} \mathcal{O}(U)
$$
o\`u la limite inductive porte sur les voisinages \'etales de $(S, \overline{s})$.
Un syst\`eme cofinal est donn\'e par les voisinages \'etales $(U, \overline{u})$
tels que $U$ soit affine et que $U \to S$ se factorise par $\Spec(A)$.
Autrement dit, nous voyons que
$$
(\mathcal{O}_S)_{\overline{s}} = \colim_{(B, \mathfrak q, \phi)} B
$$
o\`u la limite inductive porte sur les $A$-alg\`ebres \'etales $B$ munies d'un id\'eal premier
$\mathfrak q$ au-dessus de $\mathfrak p$ et d'un
morphisme de $\kappa(\mathfrak p)$-alg\`ebres
$\phi : \kappa(\mathfrak q) \to \kappa(\overline{s})$.
Notons que, puisque $\kappa(\mathfrak q)$ est une extension finie s\'eparable de
$\kappa(\mathfrak p)$, l'image de $\phi$ est contenue dans $\kappa^{sep}$.
Moyennant ces identifications, le r\'esultat du lemme est \'equivalent
au r\'esultat de
Algèbre, lemme \ref{algebra-lemma-strict-henselization-different}.
\end{proof}

\begin{definition}
\label{definition-etale-local-rings}
Soit $S$ un sch\'ema. Soit $\overline{s}$ un point g\'eom\'etrique de $S$
au-dessus du point $s \in S$.
\begin{enumerate}
\item L'{\it anneau local \'etale de $S$ en $\overline{s}$}
est la fibre du faisceau structural $\mathcal{O}_S$ sur $S_\etale$
en $\overline{s}$. Nous appelons parfois cet anneau
{\it l'hens\'elisation stricte de $\mathcal{O}_{S, s}$} relative
au point g\'eom\'etrique $\overline{s}$.
Notation utilis\'ee : $\mathcal{O}_{S, \overline{s}}^{sh}$.
\item L'{\it hens\'elisation de $\mathcal{O}_{S, s}$} d\'esigne
l'hens\'elisation de l'anneau local de $S$ en $s$. Voir
Algèbre, définition \ref{algebra-definition-henselization},
et le
th\'eor\`eme \ref{theorem-henselization}.
Notation : $\mathcal{O}_{S, s}^h$.
\item L'{\it hens\'elisation stricte de $S$ en $\overline{s}$}
est le sch\'ema $\Spec(\mathcal{O}_{S, \overline{s}}^{sh})$.
\item L'{\it hens\'elisation de $S$ en $s$} est le sch\'ema
$\Spec(\mathcal{O}_{S, s}^h)$.
\end{enumerate}
\end{definition}

\noindent
Soit $f : T \to S$ un morphisme de sch\'emas. Soit $\overline{t}$
un point g\'eom\'etrique de $T$ d'image $\overline{s}$ dans $S$.
Soient $t \in T$ et $s \in S$ leurs images.
Nous obtenons alors un diagramme commutatif canonique
$$
\xymatrix{
\Spec(\mathcal{O}^{sh}_{T, \overline{t}}) \ar[r] \ar[d] &
\Spec(\mathcal{O}^h_{T, t}) \ar[r] \ar[d] &
T \ar[d]^f \\
\Spec(\mathcal{O}^{sh}_{S, \overline{s}}) \ar[r] &
\Spec(\mathcal{O}^h_{S, s}) \ar[r] &
S
}
$$
des hens\'elisations et hens\'elisations strictes de $T$ et de $S$. On peut
le d\'emontrer en choisissant des voisinages affines de $t$ et de $s$ et en utilisant
la fonctorialit\'e des hens\'elisations (strictes) donn\'ee par
Algèbre, lemmes \ref{algebra-lemma-henselian-functorial-improve} et
\ref{algebra-lemma-strictly-henselian-functorial-improve}.

\begin{lemma}
\label{lemma-describe-henselization}
Soit $S$ un sch\'ema. Soit $s \in S$. Alors
$$
\mathcal{O}_{S, s}^h =
\colim_{(U, u)} \mathcal{O}(U)
$$
o\`u la limite inductive porte sur la cat\'egorie filtrante des
voisinages \'etales $(U, u)$ de $(S, s)$ tels que
$\kappa(s) = \kappa(u)$.
\end{lemma}

\begin{proof}
Ce lemme est une copie de
Compléments sur les morphismes, lemme \ref{more-morphisms-lemma-describe-henselization}.
\end{proof}

\begin{remark}
\label{remark-henselization-Noetherian}
Soit $S$ un sch\'ema. Soit $s \in S$.
Si $S$ est localement noeth\'erien, alors $\mathcal{O}_{S, s}^h$
est lui aussi noeth\'erien et il a le m\^eme compl\'et\'e :
$$
\widehat{\mathcal{O}_{S, s}} \cong \widehat{\mathcal{O}_{S, s}^h}.
$$
En particulier,
$\mathcal{O}_{S, s} \subset
\mathcal{O}_{S, s}^h \subset
\widehat{\mathcal{O}_{S, s}}$.
L'hens\'elisation de $\mathcal{O}_{S, s}$ est en g\'en\'eral beaucoup
plus petite que son compl\'et\'e et h\'erite de nombre de ses propri\'et\'es.
Par exemple, si $\mathcal{O}_{S, s}$ est r\'eduit, alors
$\mathcal{O}_{S, s}^h$ l'est aussi, mais ce n'est en g\'en\'eral pas vrai du compl\'et\'e.
Ins\'erer ici des r\'ef\'erences futures.
\end{remark}

\begin{lemma}
\label{lemma-etale-site-locally-ringed}
Soit $S$ un sch\'ema. Le petit site \'etale $S_\etale$, muni de
son faisceau structural $\mathcal{O}_S$, est un site localement annel\'e, voir
Modules sur les sites, définition \ref{sites-modules-definition-locally-ringed}.
\end{lemma}

\begin{proof}
Cela r\'esulte du fait que les fibres
$(\mathcal{O}_S)_{\overline{s}} = \mathcal{O}^{sh}_{S, \overline{s}}$ sont
des anneaux locaux, et du fait que $S_\etale$ a assez de points, voir
le lemme \ref{lemma-describe-etale-local-ring},
le th\'eor\`eme \ref{theorem-exactness-stalks},
et les
remarques \ref{remarks-enough-points}.
Voir
Modules sur les sites, lemmes \ref{sites-modules-lemma-locally-ringed-stalk} et
\ref{sites-modules-lemma-ringed-stalk-not-zero}
pour le fait que cela implique que le petit site \'etale est localement annel\'e.
\end{proof}






\section{Fonctorialit\'e du petit topos \'etale}
\label{section-functoriality}

\noindent
Jusqu'ici, nous n'avons pas abord\'e la fonctorialit\'e du site \'etale,
autrement dit ce qui se passe lorsque l'on se donne un morphisme de sch\'emas. Une
discussion formelle pr\'ecise se trouve dans
Topologies, section \ref{topologies-section-etale}.
Dans cette section et les suivantes, nous examinons bri\`evement ces questions,
plus pr\'ecis\'ement dans le cadre des petits sites \'etales.

\medskip\noindent
Soit $f : X \to Y$ un morphisme de sch\'emas. Nous obtenons un foncteur
\begin{equation}
\label{equation-functorial}
u : Y_\etale \longrightarrow X_\etale, \quad
V/Y \longmapsto X \times_Y V/X.
\end{equation}
Ce foncteur poss\`ede les propri\'et\'es importantes suivantes :
\begin{enumerate}
\item $u(\text{objet final}) = \text{objet final}$,
\item $u$ pr\'eserve les produits fibr\'es,
\item si $\{V_j \to V\}$ est un recouvrement dans $Y_\etale$, alors
$\{u(V_j) \to u(V)\}$ est un recouvrement dans $X_\etale$.
\end{enumerate}
Chacune d'elles se v\'erifie facilement (nous omettons la v\'erification). Par cons\'equent, nous obtenons ce
que l'on appelle un {\it morphisme de sites}
$$
f_{small} : X_\etale \longrightarrow Y_\etale,
$$
voir
Sites, définition \ref{sites-definition-morphism-sites}
et
Sites, Proposition \ref{sites-proposition-get-morphism}.
Il n'est pas n\'ecessaire de conna\^itre en d\'etail la notion abstraite
pour travailler avec les faisceaux \'etales et la cohomologie \'etale.
Il suffit habituellement de savoir qu'il existe des foncteurs
$f_{small, *}$ (image directe) et $f_{small}^{-1}$ (image inverse)
sur les faisceaux \'etales, et de conna\^itre quelques-unes de leurs propri\'et\'es simples.
Nous examinerons ces propri\'et\'es dans les sections suivantes, mais nous nous
r\'ef\'ererons parfois au cadre plus abstrait pour les d\'emonstrations, car
c'est souvent le cadre naturel pour les \`etablir.


\section{Images directes}
\label{section-direct-image}

\noindent
D\'efinissons l'image directe d'un pr\'efaisceau.

\begin{definition}
\label{definition-direct-image-presheaf}
Soit $f: X\to Y$ un morphisme de sch\'emas.
Soit $\mathcal{F} $ un pr\'efaisceau d'ensembles sur $X_\etale$.
L'{\it image directe}, ou {\it pouss\'ee en avant}, de $\mathcal{F}$
(par $f$) est
$$
f_*\mathcal{F} : Y_\etale^{opp} \longrightarrow \textit{Ensembles}, \quad
(V/Y) \longmapsto \mathcal{F}(X \times_Y V/X).
$$
Nous \`ecrivons parfois $f_* = f_{small, *}$ pour distinguer ce foncteur des autres
foncteurs d'image directe (comme l'image directe usuelle pour la topologie de Zariski ou $f_{big, *}$).
\end{definition}

\noindent
C'est bien un pr\'efaisceau \'etale, puisque le changement de base d'un morphisme
\'etale est encore \'etale. Une fa\c{c}on plus cat\'egorique de l'exprimer consiste \`a dire que
$f_*\mathcal{F}$ est le foncteur compos\'e $\mathcal{F} \circ u$
o\`u $u$ est d\'efini par l'\'equation (\ref{equation-functorial}). Cela montre
clairement que la construction d\'epend fonctoriellement du pr\'efaisceau
$\mathcal{F}$ et nous obtenons donc un foncteur
$$
f_* = f_{small, *} :
\textit{PSh}(X_\etale)
\longrightarrow
\textit{PSh}(Y_\etale)
$$
Notons que si $\mathcal{F}$ est un pr\'efaisceau de groupes ab\'eliens, alors
$f_*\mathcal{F}$ est aussi un pr\'efaisceau de groupes ab\'eliens et nous obtenons
$$
f_* = f_{small, *} :
\textit{PAb}(X_\etale)
\longrightarrow
\textit{PAb}(Y_\etale)
$$
comme pr\'ec\'edemment (c'est-\`a-dire d\'efini par exactement la m\^eme r\`egle).

\begin{remark}
\label{remark-direct-image-sheaf}
Nous affirmons que l'image directe d'un faisceau est un faisceau.
En effet, si $\{V_j \to V\}$ est un recouvrement \'etale dans $Y_\etale$,
alors $\{X \times_Y V_j \to X \times_Y V\}$ est un recouvrement \'etale dans
$X_\etale$. Par cons\'equent, la condition de faisceau pour $\mathcal{F}$ relativement
\`a $\{X \times_Y V_i \to X \times_Y V\}$
est \`equivalente \`a la condition de faisceau pour $f_*\mathcal{F}$ relativement \`a
$\{V_i \to V\}$. Ainsi, si $\mathcal{F}$ est un faisceau,
$f_*\mathcal{F}$ l'est aussi.
\end{remark}

\begin{definition}
\label{definition-direct-image-sheaf}
Soit $f: X\to Y$ un morphisme de sch\'emas.
Soit $\mathcal{F} $ un faisceau d'ensembles sur $X_\etale$.
L'{\it image directe}, ou {\it pouss\'ee en avant}, de $\mathcal{F}$
(par $f$) est
$$
f_*\mathcal{F} : Y_\etale^{opp} \longrightarrow \textit{Ensembles}, \quad
(V/Y) \longmapsto \mathcal{F}(X \times_Y V/X)
$$
qui est un faisceau d'apr\`es la
remarque \ref{remark-direct-image-sheaf}.
Nous \`ecrivons parfois $f_* = f_{small, *}$ pour distinguer ce foncteur des autres
foncteurs d'image directe (comme l'image directe usuelle pour la topologie de Zariski ou $f_{big, *}$).
\end{definition}

\noindent
La discussion pr\'ec\'edente s'applique mot pour mot et nous obtenons des foncteurs
$$
f_* = f_{small, *} :
\Sh(X_\etale)
\longrightarrow
\Sh(Y_\etale)
$$
et
$$
f_* = f_{small, *} :
\textit{Ab}(X_\etale)
\longrightarrow
\textit{Ab}(Y_\etale)
$$
\`a nouveau appel\'es {\it foncteurs d'image directe}.

\medskip\noindent
Le foncteur $f_*$ sur les faisceaux ab\'eliens est exact \`a gauche. (Voir
Homology, section \ref{homology-section-functors}
pour savoir ce que signifie qu'un foncteur entre cat\'egories ab\'eliennes est exact \`a gauche.)
En effet, si
$0 \to \mathcal{F}_1 \to \mathcal{F}_2 \to \mathcal{F}_3$
est exacte sur $X_\etale$, alors, pour tout
$U/X \in \Ob(X_\etale)$,
la suite de groupes ab\'eliens
$0 \to \mathcal{F}_1(U) \to \mathcal{F}_2(U) \to \mathcal{F}_3(U)$
est exacte. Par cons\'equent, pour tout $V/Y \in \Ob(Y_\etale)$,
la suite de groupes ab\'eliens
$0 \to f_*\mathcal{F}_1(V) \to f_*\mathcal{F}_2(V) \to f_*\mathcal{F}_3(V)$
est exacte, car c'est la suite pr\'ec\'edente avec $U = X \times_Y V$.

\begin{definition}
\label{definition-higher-direct-images}
Soit $f: X \to Y$ un morphisme de sch\'emas.
Les foncteurs d\'eriv\'es droits $\{R^pf_*\}_{p \geq 1}$ de
$f_* : \textit{Ab}(X_\etale) \to \textit{Ab}(Y_\etale)$
sont appel\'es {\it images directes sup\'erieures}.
\end{definition}

\noindent
Les images directes sup\'erieures et leurs variantes dans la cat\'egorie d\'eriv\'ee
sont examin\'ees plus en d\'etail dans (ins\'erer ici une r\'ef\'erence future).



\section{Image inverse}
\label{section-inverse-image}

\noindent
Dans cette section, nous examinons bri\`evement l'image inverse des faisceaux sur les petits
sites \'etales. Sa construction pr\'ecise se trouve dans
Topologies, section \ref{topologies-section-etale}.

\begin{definition}
\label{definition-inverse-image}
Soit $f: X\to Y$ un morphisme de sch\'emas. Les foncteurs d'{\it image inverse}, ou
de {\it tir\'e en arri\`ere}\footnote{Nous utilisons la notation $f^{-1}$ pour les images inverses de
faisceaux d'ensembles ou de groupes ab\'eliens, et nous r\'eservons $f^*$ aux
images inverses de faisceaux de modules par un morphisme de sites ou de topos annel\'es.}
sont les foncteurs
$$
f^{-1} = f_{small}^{-1} :
\Sh(Y_\etale)
\longrightarrow
\Sh(X_\etale)
$$
et
$$
f^{-1} = f_{small}^{-1} :
\textit{Ab}(Y_\etale)
\longrightarrow
\textit{Ab}(X_\etale)
$$
qui sont adjoints \`a gauche de $f_* = f_{small, *}$. Ainsi,
$f^{-1}$ est caract\'eris\'e par le fait que
$$
\Hom_{{\Sh(X_\etale)}} (f^{-1}\mathcal{G}, \mathcal{F})
=
\Hom_{\Sh(Y_\etale)} (\mathcal{G}, f_*\mathcal{F})
$$
fonctoriellement, pour tout $\mathcal{F} \in \Sh(X_\etale)$ et tout
$\mathcal{G} \in \Sh(Y_\etale)$. De m\^eme, nous avons
$$
\Hom_{{\textit{Ab}(X_\etale)}} (f^{-1}\mathcal{G}, \mathcal{F})
=
\Hom_{\textit{Ab}(Y_\etale)} (\mathcal{G}, f_*\mathcal{F})
$$
pour $\mathcal{F} \in \textit{Ab}(X_\etale)$ et
$\mathcal{G} \in \textit{Ab}(Y_\etale)$.
\end{definition}

\noindent
L'existence d'un tel adjoint n'est pas triviale.
En revanche, il existe dans un cadre assez g\'en\'eral, voir la
remarque \ref{remark-functoriality-general}
ci-dessous. Le formalisme g\'en\'eral montre que $f^{-1}\mathcal{G}$
est le faisceau associ\'e au pr\'efaisceau
\begin{equation}
\label{equation-pullback}
U/X
\longmapsto
\colim_{U \to X \times_Y V} \mathcal{G}(V/Y)
\end{equation}
o\`u la limite inductive est prise sur la cat\'egorie des couples
$(V/Y, \varphi : U/X \to X \times_Y V/X)$.
Pour le voir, appliquons
Sites, Proposition \ref{sites-proposition-get-morphism}
au foncteur $u$ de l'\'equation (\ref{equation-functorial})
et utilisons la description de $u_s = (u_p\ )^\#$ dans
Sites, sections \ref{sites-section-continuous-functors} et
\ref{sites-section-functoriality-PSh}.
Nous utiliserons occasionnellement cette formule pour l'image inverse
afin de d\'emontrer certaines de ses propri\'et\'es fondamentales.

\begin{lemma}
\label{lemma-stalk-pullback}
Soit $f : X \to Y$ un morphisme de sch\'emas.
\begin{enumerate}
\item Le foncteur
$f^{-1} : \textit{Ab}(Y_\etale) \to \textit{Ab}(X_\etale)$
est exact.
\item Le foncteur
$f^{-1} : \Sh(Y_\etale) \to \Sh(X_\etale)$
est exact, c'est-\`a-dire qu'il commute aux limites et colimites finies, voir
Catégories, définition \ref{categories-definition-exact}.
\item Soit $\overline{x} \to X$ un point g\'eom\'etrique.
Soit $\mathcal{G}$ un faisceau sur $Y_\etale$.
Il existe alors une identification canonique
$$
(f^{-1}\mathcal{G})_{\overline{x}} = \mathcal{G}_{\overline{y}}.
$$
o\`u $\overline{y} = f \circ \overline{x}$.
\item Pour tout $V \to Y$ \'etale, nous avons $f^{-1}h_V = h_{X \times_Y V}$.
\end{enumerate}
\end{lemma}

\begin{proof}
L'exactitude de $f^{-1}$ sur les faisceaux d'ensembles r\'esulte de
Sites, Proposition \ref{sites-proposition-get-morphism}
appliqu\'ee \`a notre foncteur $u$ de l'\'equation (\ref{equation-functorial}).
En fait, l'exactitude de l'image inverse fait partie de la d\'efinition d'un
morphisme de topos (ou de sites, si l'on pr\'ef\`ere). Ainsi, (2) est v\'erifi\'e.
Cela implique (1), car, pour tout faisceau ab\'elien $\mathcal{G}$ sur
$Y_\etale$,
le faisceau d'ensembles sous-jacent \`a $f^{-1}\mathcal{F}$ est le m\^eme
que l'image par $f^{-1}$ du faisceau d'ensembles sous-jacent \`a $\mathcal{F}$, voir
Sites, section \ref{sites-section-sheaves-algebraic-structures}.
Voir aussi
Modules sur les sites, lemme \ref{sites-modules-lemma-flat-pullback-exact}.
Dans la litt\'erature, (1) et (2) sont parfois d\'eduits de (3) au moyen du
th\'eor\`eme \ref{theorem-exactness-stalks}.

\medskip\noindent
L'assertion (3) est un fait g\'en\'eral sur les fibres des images inverses, voir
Sites, lemme \ref{sites-lemma-point-morphism-sites}.
Nous allons aussi d\'emontrer directement (3) comme suit. Notons, d'apr\`es le
lemme \ref{lemma-stalk-exact},
que la prise des fibres commute \`a la faisceautisation.
Rappelons maintenant que $f^{-1}\mathcal{G}$ est le faisceau
associ\'e au pr\'efaisceau
$$
U \longrightarrow \colim_{U \to X \times_Y V} \mathcal{G}(V),
$$
voir l'\'equation (\ref{equation-pullback}).
Ainsi, nous avons
\begin{align*}
(f^{-1}\mathcal{G})_{\overline{x}}
& = \colim_{(U, \overline{u})} f^{-1}\mathcal{G}(U) \\
& = \colim_{(U, \overline{u})}
\colim_{a : U \to X \times_Y V} \mathcal{G}(V) \\
& = \colim_{(V, \overline{v})} \mathcal{G}(V) \\
& = \mathcal{G}_{\overline{y}}
\end{align*}
Dans la troisi\`eme \`egalit\'e, le couple $(U, \overline{u})$ et le morphisme
$a : U \to X \times_Y V$ correspondent au couple $(V, a \circ \overline{u})$.

\medskip\noindent
L'assertion (4) se d\'emontre de mani\`ere analogue en identifiant les limites inductives
qui d\'efinissent $f^{-1}h_V$. On peut aussi utiliser
le lemme de Yoneda (Catégories, lemme \ref{categories-lemma-yoneda})
et les \`egalit\'es fonctorielles
$$
\Mor_{\Sh(X_\etale)}(f^{-1}h_V, \mathcal{F}) =
\Mor_{\Sh(Y_\etale)}(h_V, f_*\mathcal{F}) =
f_*\mathcal{F}(V) = \mathcal{F}(X \times_Y V)
$$
ainsi que le fait que les pr\'efaisceaux repr\'esentables sont des faisceaux. Voir aussi
Sites, lemme \ref{sites-lemma-pullback-representable-sheaf}
pour un r\'esultat enti\`erement g\'en\'eral.
\end{proof}

\noindent
Le couple de foncteurs $(f_*, f^{-1})$ d\'efinit un morphisme de petits topos \'etales
$$
f_{small} :
\Sh(X_\etale)
\longrightarrow
\Sh(Y_\etale)
$$
De nombreux r\'esultats g\'en\'eraux sur la cohomologie des faisceaux valent pour les topos et les
morphismes de topos. Nous nous efforcerons de signaler ceux qui sont
g\'en\'eraux et ceux qui sont propres au topos \'etale.

\begin{remark}
\label{remark-functoriality-general}
Plus g\'en\'eralement, soient $\mathcal{C}_1, \mathcal{C}_2$ des sites, et
supposons qu'ils poss\`edent des objets finals et des produits fibr\'es. Soit
$u: \mathcal{C}_2 \to \mathcal{C}_1$ un foncteur satisfaisant aux conditions suivantes :
\begin{enumerate}
\item si $\{V_i \to V\}$ est un recouvrement dans $\mathcal{C}_2$, alors
$\{u(V_i) \to u(V)\}$ est un recouvrement dans $\mathcal{C}_1$ (nous
disons que $u$ est {\it continu}), et
\item $u$ commute aux limites finies (c'est-\`a-dire que $u$ est exact \`a gauche, ou encore que
$u$ pr\'eserve les produits fibr\'es et les objets finals).
\end{enumerate}
On peut alors d\'efinir
$f_*: \Sh(\mathcal{C}_1) \to \Sh(\mathcal{C}_2)$
par $ f_* \mathcal{F}(V) = \mathcal{F}(u(V))$.
De plus, il existe un foncteur exact $f^{-1}$ qui
est adjoint \`a gauche de $f_*$, voir
Sites, définition \ref{sites-definition-morphism-sites} et
Proposition \ref{sites-proposition-get-morphism}.
Attention : demander simplement que $u$ soit continu
et commute aux produits fibr\'es ne suffit pas pour obtenir un morphisme de topos.
\end{remark}




\section{Fonctorialit\'e des grands topos}
\label{section-functoriality-big-topoi}

\noindent
Pour tout morphisme de sch\'emas $f : X \to Y$, on dispose de toute une famille de
morphismes de topos associ\'es \`a $f$, voir
Topologies, section \ref{topologies-section-change-topologies}
pour une liste. Les plus utilis\'es sont peut-\^etre les morphismes de topos
$$
f_{big} = f_{big, \tau} :
\Sh((\Sch/X)_\tau)
\longrightarrow
\Sh((\Sch/Y)_\tau)
$$
o\`u $\tau \in \{Zariski, \etale, smooth, syntomic, fppf\}$.
Chacun d'eux correspond \`a un foncteur continu
$$
(\Sch/Y)_\tau \longrightarrow (\Sch/X)_\tau, \quad
V/Y \longmapsto X \times_Y V/X
$$
qui pr\'eserve les objets finals, les produits fibr\'es et les recouvrements, et d\'efinit donc
un morphisme de sites
$$
f_{big} : (\Sch/X)_\tau \longrightarrow (\Sch/Y)_\tau.
$$
Voir
Topologies, sections \ref{topologies-section-zariski},
\ref{topologies-section-etale},
\ref{topologies-section-smooth},
\ref{topologies-section-syntomic}, et
\ref{topologies-section-fppf}.
En particulier, l'image directe par $f_{big}$ est donn\'ee par la r\`egle
$$
(f_{big, *}\mathcal{F})(V/Y) = \mathcal{F}(X \times_Y V/X)
$$
Il se trouve que ces morphismes de topos poss\`edent un foncteur image
inverse $f_{big}^{-1}$ tr\`es facile \`a d\'ecrire.
Plus pr\'ecis\'ement, nous avons
$$
(f_{big}^{-1}\mathcal{G})(U/X) = \mathcal{G}(U/Y)
$$
o\`u le morphisme structural de $U/Y$ est la compos\'ee du
morphisme structural $U \to X$ et de $f$, voir
Topologies, lemmes \ref{topologies-lemma-morphism-big},
\ref{topologies-lemma-morphism-big-etale},
\ref{topologies-lemma-morphism-big-smooth},
\ref{topologies-lemma-morphism-big-syntomic}, et
\ref{topologies-lemma-morphism-big-fppf}.







\section{Fonctorialit\'e et faisceaux de modules}
\label{section-morphisms-modules}

\noindent
Dans cette section, nous allons reformuler une partie des r\'esultats expliqu\'es dans
Descente, sections \ref{descent-section-quasi-coherent-sheaves},
\ref{descent-section-quasi-coherent-cohomology}, et
\ref{descent-section-quasi-coherent-sheaves-bis}
dans le cadre des topologies \'etales. Soit $f : X \to Y$ un morphisme de
sch\'emas. Nous avons vu plus haut, voir les
sections \ref{section-functoriality}, \ref{section-direct-image} et
\ref{section-inverse-image},
qu'il induit un morphisme $f_{small}$ de petits sites \'etales. Dans
Descente, remarque \ref{descent-remark-change-topologies-ringed},
nous avons vu que $f$ induit aussi un morphisme naturel
$$
f_{small}^\sharp :
\mathcal{O}_{Y_\etale}
\longrightarrow
f_{small, *}\mathcal{O}_{X_\etale}
$$
de faisceaux d'anneaux sur
$Y_\etale$ tel que $(f_{small}, f_{small}^\sharp)$
soit un morphisme de sites annel\'es. Voir
Modules sur les sites, définition \ref{sites-modules-definition-ringed-site}
pour la d\'efinition d'un morphisme de sites annel\'es.
Rappelons simplement ici que $f_{small}^\sharp$ est d\'efini par le
syst\`eme compatible de morphismes
$$
\text{pr}_V^\sharp : \mathcal{O}(V) \longrightarrow \mathcal{O}(X \times_Y V)
$$
lorsque $V$ parcourt les objets de $Y_\etale$.

\medskip\noindent
Il est clair que cette construction est compatible avec les compositions de
morphismes de sch\'emas. Plus pr\'ecis\'ement, si $f : X \to Y$ et $g : Y \to Z$
sont des morphismes de sch\'emas, alors nous avons
$$
(g_{small}, g_{small}^\sharp)
\circ
(f_{small}, f_{small}^\sharp)
=
((g \circ f)_{small}, (g \circ f)_{small}^\sharp)
$$
comme morphismes de topos annel\'es. De plus, d'apr\`es
Modules sur les sites, définition \ref{sites-modules-definition-pushforward},
nous voyons que tout morphisme de sch\'emas $f : X \to Y$
donne des foncteurs bien d\'efinis d'image inverse et d'image directe
\begin{align*}
f_{small}^* :
\textit{Mod}(\mathcal{O}_{Y_\etale})
\longrightarrow
\textit{Mod}(\mathcal{O}_{X_\etale}), \\
f_{small, *} :
\textit{Mod}(\mathcal{O}_{X_\etale})
\longrightarrow
\textit{Mod}(\mathcal{O}_{Y_\etale})
\end{align*}
qui sont adjoints de la mani\`ere habituelle. Si $g : Y \to Z$ est un autre morphisme
de sch\'emas, alors nous avons
$(g \circ f)_{small}^* = f_{small}^* \circ g_{small}^*$
et $(g \circ f)_{small, *} = g_{small, *} \circ f_{small, *}$
en vertu de ce que nous avons dit sur les compositions.

\medskip\noindent
Il y a une diff\'erence assez importante entre la cat\'egorie
de tous les $\mathcal{O}_X$-modules sur $X$ et la cat\'egorie de tous les
$\mathcal{O}_{X_\etale}$-modules sur $X_\etale$. Mais les
r\'esultats de
Descente, sections \ref{descent-section-quasi-coherent-sheaves},
\ref{descent-section-quasi-coherent-cohomology}, et
\ref{descent-section-quasi-coherent-sheaves-bis}
nous disent qu'il y a peu de diff\'erence entre consid\'erer les modules quasi-coh\'erents
sur $S$ et les modules quasi-coh\'erents sur $S_\etale$.
(Nous l'avons d\'ej\`a vu dans le
th\'eor\`eme \ref{theorem-quasi-coherent},
par exemple.)
En particulier, si $f : X \to Y$ est un morphisme quelconque de sch\'emas, alors
les foncteurs d'image inverse $f_{small}^*$ et $f^*$ s'identifient pour les
faisceaux quasi-coh\'erents, voir
Descente,
Proposition \ref{descent-proposition-equivalence-quasi-coherent-functorial}.
De plus, il en va de m\^eme pour l'image directe pourvu que $f$ soit
quasi-compact et quasi-s\'epar\'e, voir
Descente, lemme \ref{descent-lemma-higher-direct-images-small-etale}.

\medskip\noindent
Disons quelques mots de la fonctorialit\'e du faisceau structural sur les grands sites.
Soit $f : X \to Y$ un morphisme de sch\'emas. Choisissons l'une des
topologies $\tau \in \{Zariski,\linebreak[0]
\etale,\linebreak[0] smooth,\linebreak[0] syntomic,\linebreak[0]
fppf\}$. Alors le morphisme
$f_{big} : (\Sch/X)_\tau \to (\Sch/Y)_\tau$
devient un morphisme de sites annel\'es au moyen d'un morphisme
$$
f_{big}^\sharp : \mathcal{O}_Y \longrightarrow f_{big, *}\mathcal{O}_X
$$
voir Descente, remarque \ref{descent-remark-change-topologies-ringed}.
En fait, il est donn\'e par la m\^eme construction que celle expos\'ee pour les petits
sites plus haut.












\section{Comparaison de topologies}
\label{section-compare-topologies}

\noindent
Dans cette section, nous commen\c{c}ons \`a \`etudier ce qui se produit lorsque l'on compare
les faisceaux relativement \`a diff\'erentes topologies.

\begin{lemma}
\label{lemma-where-sections-are-equal}
Soit $S$ un sch\'ema. Soit $\mathcal{F}$ un faisceau d'ensembles sur $S_\etale$.
Soient $s, t \in \mathcal{F}(S)$. Il existe alors un ouvert $W \subset S$
caract\'eris\'e par la propri\'et\'e suivante : un morphisme $f : T \to S$
se factorise par $W$ si et seulement si $s|_T = t|_T$ (la restriction est
l'image inverse par $f_{small}$).
\end{lemma}

\begin{proof}
Consid\'erons le pr\'efaisceau qui associe \`a $U \in \Ob(S_\etale)$ l'ensemble vide
si $s|_U \not = t|_U$, et un singleton sinon. Il est clair qu'il s'agit
d'un sous-faisceau de l'objet final de $\Sh(S_\etale)$. D'apr\`es le
lemme \ref{lemma-support-subsheaf-final},
nous trouvons un ouvert $W \subset S$ qui repr\'esente ce pr\'efaisceau.
Pour un point g\'eom\'etrique $\overline{x}$ de $S$, nous voyons que $\overline{x} \in W$
si et seulement si les germes de $s$ et de $t$ en $\overline{x}$ sont identiques.
D'apr\`es la description des fibres des images inverses donn\'ee dans le
lemme \ref{lemma-stalk-pullback},
nous voyons que $W$ poss\`ede la propri\'et\'e voulue.
\end{proof}

\begin{lemma}
\label{lemma-describe-pullback}
Soit $S$ un sch\'ema. Soit $\tau \in \{Zariski, \etale\}$. Consid\'erons le morphisme
$$
\pi_S : (\Sch/S)_\tau \longrightarrow S_\tau
$$
de Topologies, lemme \ref{topologies-lemma-at-the-bottom} ou
\ref{topologies-lemma-at-the-bottom-etale}. Soit $\mathcal{F}$ un faisceau sur
$S_\tau$. Alors $\pi_S^{-1}\mathcal{F}$ est donn\'e par la r\`egle
$$
(\pi_S^{-1}\mathcal{F})(T) = \Gamma(T_\tau, f_{small}^{-1}\mathcal{F})
$$
o\`u $f : T \to S$. De plus, $\pi_S^{-1}\mathcal{F}$ satisfait \`a la
condition de faisceau relativement aux recouvrements fpqc.
\end{lemma}

\begin{proof}
Observons que nous avons un morphisme $i_f : \Sh(T_\tau) \to \Sh(\Sch/S)_\tau)$
tel que $\pi_S \circ i_f = f_{small}$ comme morphismes
$T_\tau \to S_\tau$, voir
Topologies, lemmes \ref{topologies-lemma-put-in-T},
\ref{topologies-lemma-morphism-big-small},
\ref{topologies-lemma-put-in-T-etale}, et
\ref{topologies-lemma-morphism-big-small-etale}.
Comme l'image inverse est transitive, nous voyons que
$i_f^{-1} \pi_S^{-1}\mathcal{F} = f_{small}^{-1}\mathcal{F}$ comme voulu.

\medskip\noindent
Soit $\{g_i : T_i \to T\}_{i \in I}$ un recouvrement fpqc.
La derni\`ere assertion signifie ce qui suit : \`etant donn\'e un faisceau $\mathcal{G}$
sur $T_\tau$ et des sections
$s_i \in \Gamma(T_i, g_{i, small}^{-1}\mathcal{G})$ dont les images inverses
sur $T_i \times_T T_j$ concordent, il existe une unique section $s$ de $\mathcal{G}$
sur $T$ dont l'image inverse sur $T_i$ concorde avec $s_i$.

\medskip\noindent
Soit $V \to T$ un objet de $T_\tau$ et soit $t \in \mathcal{G}(V)$.
Pour tout $i$, il existe un plus grand ouvert $W_i \subset T_i \times_T V$
tel que les images inverses de $s_i$ et de $t$ concordent comme sections de l'image
inverse de $\mathcal{G}$ sur $W_i \subset T_i \times_T V$, voir le
lemme \ref{lemma-where-sections-are-equal}.
Comme $s_i$ et $s_j$ concordent sur $T_i \times_T T_j$, nous constatons
que les images inverses de $W_i$ et de $W_j$ sont le m\^eme ouvert au-dessus de
$T_i \times_T T_j \times_T V$. D'apr\`es
Descente, lemme \ref{descent-lemma-open-fpqc-covering},
nous trouvons un ouvert $W \subset V$ dont l'image inverse sur $T_i \times_T V$
redonne $W_i$.

\medskip\noindent
Par construction de $g_{i, small}^{-1}\mathcal{G}$, il existe
un recouvrement pour $\tau$ $\{T_{ij} \to T_i\}_{j \in J_i}$, pour chaque $j$ une
immersion ouverte ou un morphisme \'etale $V_{ij} \to T$, une section
$t_{ij} \in \mathcal{G}(V_{ij})$, et des diagrammes commutatifs
$$
\xymatrix{
T_{ij} \ar[r] \ar[d] & V_{ij} \ar[d] \\
T_i \ar[r] & T
}
$$
tels que $s_i|_{T_{ij}}$ soit l'image inverse de $t_{ij}$. Autrement dit,
apr\`es avoir remplac\'e le recouvrement $\{T_i \to T\}$ par $\{T_{ij} \to T\}$,
nous pouvons supposer qu'il existe des factorisations $T_i \to V_i \to T$ avec
$V_i \in \Ob(T_\tau)$ et des sections $t_i \in \mathcal{G}(V_i)$
dont les images inverses sont $s_i$ sur $T_i$.
D'apr\`es le r\'esultat du paragraphe pr\'ec\'edent, nous trouvons des ouverts $W_i \subset V_i$
tels que $t_i|_{W_i}$ ``concorde avec'' chaque $s_j$ sur $T_j \times_T W_i$.
Notons que $T_i \to V_i$ se factorise par $W_i$.
Ainsi, $\{W_i \to T\}$ est un recouvrement pour $\tau$ et le lemme est d\'emontr\'e.
\end{proof}

\begin{lemma}
\label{lemma-sections-upstairs}
Soit $S$ un sch\'ema. Soit $f : T \to S$ un morphisme tel que
\begin{enumerate}
\item $f$ est plat et quasi-compact, et
\item les fibres g\'eom\'etriques de $f$ sont connexes.
\end{enumerate}
Soit $\mathcal{F}$ un faisceau sur $S_\etale$.
Alors $\Gamma(S, \mathcal{F}) = \Gamma(T, f^{-1}_{small}\mathcal{F})$.
\end{lemma}

\begin{proof}
Il existe un morphisme canonique
$\Gamma(S, \mathcal{F}) \to \Gamma(T, f_{small}^{-1}\mathcal{F})$.
Comme $f$ est surjectif (puisque ses fibres sont connexes), nous voyons que
ce morphisme est injectif.

\medskip\noindent
Pour montrer que le morphisme est surjectif, soit
$\alpha \in \Gamma(T, f_{small}^{-1}\mathcal{F})$.
Comme $\{T \to S\}$ est un recouvrement fpqc, nous pouvons utiliser le
lemme \ref{lemma-describe-pullback} pour voir qu'il suffit de montrer que
les deux images inverses de $\alpha$ sur $T \times_S T$ sont la m\^eme section
par les deux projections. Soit $\overline{s} \to S$ un point g\'eom\'etrique.
Il suffit de montrer l'\'egalit\'e sur $(T \times_S T)_{\overline{s}}$,
car tout point g\'eom\'etrique de $T \times_S T$ est contenu dans l'une de
ces fibres g\'eom\'etriques. Autrement dit, nous voulons montrer que
$\alpha|_{T_{\overline{s}}}$ donne la m\^eme section par les deux projections sur
$$
(T \times_S T)_{\overline{s}} =
T_{\overline{s}} \times_{\overline{s}} T_{\overline{s}}
$$
vers $T_{\overline{s}}$.
Cependant, comme $\mathcal{F}|_{T_{\overline{s}}}$ est l'image
inverse de $\mathcal{F}|_{\overline{s}}$, c'est un faisceau constant
de valeur $\mathcal{F}_{\overline{s}}$. Comme $T_{\overline{s}}$
est connexe par hypoth\`ese, toute section d'un faisceau constant est constante.
Ainsi, $\alpha|_{T_{\overline{s}}}$ correspond \`a un \`el\'ement
de $\mathcal{F}_{\overline{s}}$. Les deux images inverses sur
$(T \times_S T)_{\overline{s}}$ correspondent donc toutes deux
\`a ce m\^eme \`el\'ement, d'o\`u la conclusion.
\end{proof}

\noindent
Voici une variante du lemme \ref{lemma-sections-upstairs}
o\`u nous ne supposons pas que le morphisme est plat.

\begin{lemma}
\label{lemma-sections-upstairs-submersive}
Soit $S$ un sch\'ema. Soit $f : X \to S$ un morphisme tel que
\begin{enumerate}
\item $f$ est submersif, et
\item les fibres g\'eom\'etriques de $f$ sont connexes.
\end{enumerate}
Soit $\mathcal{F}$ un faisceau sur $S_\etale$.
Alors $\Gamma(S, \mathcal{F}) = \Gamma(X, f^{-1}_{small}\mathcal{F})$.
\end{lemma}

\begin{proof}
Il existe un morphisme canonique
$\Gamma(S, \mathcal{F}) \to \Gamma(X, f_{small}^{-1}\mathcal{F})$.
Comme $f$ est surjectif (puisque ses fibres sont connexes), nous voyons que
ce morphisme est injectif.

\medskip\noindent
Pour montrer que le morphisme est surjectif, soit
$\tau \in \Gamma(X, f_{small}^{-1}\mathcal{F})$.
Il suffit de trouver
un recouvrement \'etale $\{U_i \to S\}$ et des
sections $\sigma_i \in \mathcal{F}(U_i)$
telles que l'image inverse de $\sigma_i$ soit $\tau|_{X \times_S U_i}$.
En effet, l'injectivit\'e montr\'ee ci-dessus garantit
que $\sigma_i$ et $\sigma_j$ induisent la m\^eme
section de $\mathcal{F}$ sur $U_i \times_S U_j$.
Nous obtenons donc une unique section $\sigma \in \mathcal{F}(S)$
qui se restreint \`a $\sigma_i$ sur $U_i$.
L'image inverse de $\sigma$ sur $X$ est alors $\tau$,
car c'est vrai localement.

\medskip\noindent
Soit $\overline{x}$ un point g\'eom\'etrique de $X$ d'image $\overline{s}$
dans $S$. Consid\'erons l'image de $\tau$ dans la fibre
$$
(f_{small}^{-1}\mathcal{F})_{\overline{x}} = \mathcal{F}_{\overline{s}}
$$
Voir le lemme \ref{lemma-stalk-pullback}.
Nous pouvons trouver un voisinage \'etale $U \to S$ de $\overline{s}$
et une section $\sigma \in \mathcal{F}(U)$ qui s'envoie sur cette image
dans la fibre. Ainsi, apr\`es avoir remplac\'e $S$ par $U$ et $X$ par $X \times_S U$,
nous pouvons supposer qu'il existe une section $\sigma$ de $\mathcal{F}$ sur $S$
dont l'image dans $(f_{small}^{-1}\mathcal{F})_{\overline{x}}$ est la
m\^eme que celle de $\tau$.

\medskip\noindent
D'apr\`es le lemme \ref{lemma-where-sections-are-equal},
il existe un ouvert maximal $W \subset X$ sur lequel
$f_{small}^{-1}\sigma$ et $\tau$ concordent sur $W$,
et la formation de $W$ commute aux changements de base ult\'erieurs.
Observons que l'image inverse de $\mathcal{F}$ sur la
fibre g\'eom\'etrique $X_{\overline{s}}$ est l'image inverse
de $\mathcal{F}_{\overline{s}}$, consid\'er\'e comme faisceau sur
$\overline{s}$, par $X_{\overline{s}} \to \overline{s}$.
Nous voyons donc que $\tau$ et $\sigma$ donnent des sections
du faisceau constant de valeur $\mathcal{F}_{\overline{s}}$
sur $X_{\overline{s}}$ qui concordent en un point. Comme
$X_{\overline{s}}$ est connexe par hypoth\`ese, nous concluons
que $W$ contient $X_s$. Le m\^eme argument appliqu\'e aux autres
fibres g\'eom\'etriques montre que $W$ contient toute fibre qu'il rencontre.
Comme $f$ est submersif, nous concluons que $W$ est l'image
inverse d'un voisinage ouvert de $s$ dans $S$.
Cela ach\`eve la d\'emonstration.
\end{proof}

\begin{lemma}
\label{lemma-sections-base-field-extension}
Soit $K/k$ une extension de corps o\`u $k$ est
s\'eparablement clos. Soit $S$ un sch\'ema sur $k$. Notons
$p : S_K = S \times_{\Spec(k)} \Spec(K) \to S$ la projection.
Soit $\mathcal{F}$ un faisceau sur $S_\etale$.
Alors $\Gamma(S, \mathcal{F}) = \Gamma(S_K, p^{-1}_{small}\mathcal{F})$.
\end{lemma}

\begin{proof}
Cela r\'esulte du lemme \ref{lemma-sections-upstairs}. En effet, il est clair
que $p$ est plat et quasi-compact comme changement de base de
$\Spec(K) \to \Spec(k)$. D'autre part, si $\overline{s} : \Spec(L) \to S$
est un point g\'eom\'etrique, alors la fibre de $p$ en $\overline{s}$
est le spectre de $K \otimes_k L$, qui est irr\'eductible, donc connexe, d'apr\`es
Algèbre, lemme \ref{algebra-lemma-separably-closed-irreducible}.
\end{proof}









\section{Reconstruction des morphismes}
\label{section-morphisms}

\noindent
Dans cette section, nous d\'emontrons que la r\`egle qui associe \`a un sch\'ema
son petit topos \'etale localement annel\'e est pleinement fid\`ele en un sens
convenable, voir le
th\'eor\`eme \ref{theorem-fully-faithful}.

\begin{lemma}
\label{lemma-morphism-locally-ringed}
Soit $f : X \to Y$ un morphisme de sch\'emas.
Le morphisme de sites annel\'es $(f_{small}, f_{small}^\sharp)$
associ\'e \`a $f$ est un morphisme de sites localement annel\'es, voir
Modules sur les sites,
définition \ref{sites-modules-definition-morphism-locally-ringed-topoi}.
\end{lemma}

\begin{proof}
Notons que l'assertion a un sens, puisque nous avons vu que
$(X_\etale, \mathcal{O}_{X_\etale})$ et
$(Y_\etale, \mathcal{O}_{Y_\etale})$
sont des sites localement annel\'es, voir le
lemme \ref{lemma-etale-site-locally-ringed}.
De plus, nous savons que $X_\etale$ a assez de points, voir le
th\'eor\`eme \ref{theorem-exactness-stalks}
et les
remarques \ref{remarks-enough-points}.
Il suffit donc de d\'emontrer que $(f_{small}, f_{small}^\sharp)$
satisfait \`a la condition (3) de
Modules sur les sites,
lemme \ref{sites-modules-lemma-locally-ringed-morphism}.
Pour le voir, prenons un point $p$ de $X_\etale$. D'apr\`es le
lemme \ref{lemma-points-small-etale-site},
$p$ correspond \`a un point g\'eom\'etrique $\overline{x}$ de $X$.
D'apr\`es le
lemme \ref{lemma-stalk-pullback},
le point $q = f_{small} \circ p$ correspond au
point g\'eom\'etrique $\overline{y} = f \circ \overline{x}$ de $Y$.
L'assertion \`a d\'emontrer est donc que le morphisme induit
sur les fibres
$$
(\mathcal{O}_Y)_{\overline{y}} \longrightarrow (\mathcal{O}_X)_{\overline{x}}
$$
est un morphisme local d'anneaux. Supposons que $a \in (\mathcal{O}_Y)_{\overline{y}}$
soit un \`el\'ement du membre de gauche qui s'envoie sur un \`el\'ement de l'id\'eal
maximal du membre de droite. Supposons que $a$ soit la classe d'\'equivalence
d'un triplet $(V, \overline{v}, a)$ o\`u $V \to Y$ est \'etale,
$\overline{v} : \overline{x} \to V$ au-dessus de $Y$, et $a \in \mathcal{O}(V)$.
Il s'envoie sur la classe d'\'equivalence de
$(X \times_Y V, \overline{x} \times \overline{v}, \text{pr}_V^\sharp(a))$
dans l'anneau local $(\mathcal{O}_X)_{\overline{x}}$. Mais il est clair que
le fait d'appartenir \`a l'id\'eal maximal signifie que l'image de $\text{pr}_V^\sharp(a)$
dans $\kappa(\overline{x})$ est nulle. Par cons\'equent, l'image de
$a$ dans $\kappa(\overline{x})$ est elle aussi nulle. Cela signifie que $a$ appartient \`a
l'id\'eal maximal de $(\mathcal{O}_Y)_{\overline{y}}$.
\end{proof}

\begin{lemma}
\label{lemma-2-morphism}
Soient $X$, $Y$ des sch\'emas. Soit $f : X \to Y$ un morphisme de sch\'emas.
Soit $t$ un $2$-morphisme de $(f_{small}, f_{small}^\sharp)$ vers lui-m\^eme, voir
Modules sur les sites,
définition \ref{sites-modules-definition-2-morphism-ringed-topoi}.
Alors $t = \text{id}$.
\end{lemma}

\begin{proof}
Cela signifie que $t : f^{-1}_{small} \to f^{-1}_{small}$
est une transformation naturelle telle que le diagramme
$$
\xymatrix{
f_{small}^{-1}\mathcal{O}_Y
\ar[rd]_{f_{small}^\sharp}  & &
f_{small}^{-1}\mathcal{O}_Y \ar[ll]^t \ar[ld]^{f_{small}^\sharp} \\
& \mathcal{O}_X
}
$$
soit commutatif. Supposons que $V \to Y$ soit \'etale avec $V$ affine. D'apr\`es
Morphismes, lemme \ref{morphisms-lemma-quasi-affine-finite-type-over-S},
nous pouvons choisir une immersion $i : V \to \mathbf{A}^n_Y$ au-dessus de $Y$.
En termes de faisceaux, cela signifie que $i$ induit une injection
$h_i : h_V \to \prod_{j = 1, \ldots, n} \mathcal{O}_Y$ de faisceaux.
Le changement de base $i'$ de $i$ \`a $X$ est une immersion
(Schémas, lemme \ref{schemes-lemma-base-change-immersion}).
Ainsi, $i' : X \times_Y V \to \mathbf{A}^n_X$ est une immersion, ce qui
signifie \`a son tour que
$h_{i'} : h_{X \times_Y V} \to \prod_{j = 1, \ldots, n} \mathcal{O}_X$
est une injection de faisceaux.
Via l'identification $f_{small}^{-1}h_V = h_{X \times_Y V}$ du
lemme \ref{lemma-stalk-pullback},
le morphisme $h_{i'}$ est \`egal \`a
$$
\xymatrix{
f_{small}^{-1}h_V \ar[r]^-{f^{-1}h_i} &
\prod_{j = 1, \ldots, n} f_{small}^{-1}\mathcal{O}_Y
\ar[r]^{\prod f^\sharp} &
\prod_{j = 1, \ldots, n} \mathcal{O}_X
}
$$
(nous omettons la v\'erification). Cela signifie que le morphisme
$t : f_{small}^{-1}h_V \to f_{small}^{-1}h_V$
s'ins\`ere dans le diagramme commutatif
$$
\xymatrix{
f_{small}^{-1}h_V \ar[r]^-{f^{-1}h_i} \ar[d]^t &
\prod_{j = 1, \ldots, n} f_{small}^{-1}\mathcal{O}_Y
\ar[r]^-{\prod f^\sharp} \ar[d]^{\prod t} &
\prod_{j = 1, \ldots, n} \mathcal{O}_X \ar[d]^{\text{id}}\\
f_{small}^{-1}h_V \ar[r]^-{f^{-1}h_i} &
\prod_{j = 1, \ldots, n} f_{small}^{-1}\mathcal{O}_Y
\ar[r]^-{\prod f^\sharp} &
\prod_{j = 1, \ldots, n} \mathcal{O}_X
}
$$
Le carr\'e de droite est commutatif par notre hypoth\`ese sur $t$
expliqu\'ee ci-dessus.
Comme la compos\'ee des fl\`eches horizontales est injective
d'apr\`es la discussion pr\'ec\'edente, nous concluons que la fl\`eche verticale de gauche
est elle aussi l'identit\'e. Tout faisceau d'ensembles sur
$Y_\etale$ est quotient d'un coproduit (\'enorme) de faisceaux
de la forme $h_V$ avec $V$ affine (combiner
Topologies, lemme \ref{topologies-lemma-alternative}
avec
Sites, lemme \ref{sites-lemma-sheaf-coequalizer-representable}).
Nous concluons donc que $t : f_{small}^{-1} \to f_{small}^{-1}$
est la transformation identit\'e, comme voulu.
\end{proof}

\begin{lemma}
\label{lemma-faithful}
Soient $X$, $Y$ des sch\'emas.
Deux morphismes de sch\'emas $a, b : X \to Y$
pour lesquels il existe un $2$-isomorphisme
$(a_{small}, a_{small}^\sharp) \cong (b_{small}, b_{small}^\sharp)$
dans la $2$-cat\'egorie des topos annel\'es sont \`egaux.
\end{lemma}

\begin{proof}
Donnons l'argument avec soin, car il est quelque peu d\'eroutant.
Soit $t : a_{small}^{-1} \to b_{small}^{-1}$ le $2$-isomorphisme.
Consid\'erons un ouvert quelconque $V \subset Y$. Notons que $h_V$ est un sous-faisceau du
faisceau final $*$. Ainsi, $a_{small}^{-1}h_V = h_{a^{-1}(V)}$
et $b_{small}^{-1}h_V = h_{b^{-1}(V)}$ sont tous deux des sous-faisceaux du faisceau final.
L'isomorphisme
$$
t : a_{small}^{-1}h_V = h_{a^{-1}(V)} \to b_{small}^{-1}h_V = h_{b^{-1}(V)}
$$
est donc n\'ecessairement l'identit\'e, et $a^{-1}(V) = b^{-1}(V)$.
Il s'ensuit que $a$ et $b$ sont \`egaux sur les espaces topologiques sous-jacents.
Prenons ensuite une section $f \in \mathcal{O}_Y(V)$. Elle d\'etermine, et est
d\'etermin\'ee par, un morphisme de faisceaux d'ensembles
$f : h_V \to \mathcal{O}_Y$.
Prenons son image inverse et appliquons $t$ pour obtenir un diagramme commutatif
$$
\xymatrix{
h_{b^{-1}(V)} \ar@{=}[r] &
b_{small}^{-1}h_V \ar[d]^{b_{small}^{-1}(f)} & &
a_{small}^{-1}h_V \ar[d]^{a_{small}^{-1}(f)} \ar[ll]^t &
h_{a^{-1}(V)} \ar@{=}[l]
\\
& b_{small}^{-1}\mathcal{O}_Y
\ar[rd]_{b^\sharp}  & &
a_{small}^{-1}\mathcal{O}_Y \ar[ll]^t \ar[ld]^{a^\sharp} \\
& & \mathcal{O}_X
}
$$
o\`u le triangle est commutatif par d\'efinition d'un $2$-isomorphisme dans
Modules sur les sites, section \ref{sites-modules-section-2-category}.
Nous avons vu plus haut que la compos\'ee des fl\`eches horizontales
sup\'erieures provient de l'identit\'e $a^{-1}(V) = b^{-1}(V)$.
La commutativit\'e du diagramme nous dit donc que
$a_{small}^\sharp(f) = b_{small}^\sharp(f)$ dans
$\mathcal{O}_X(a^{-1}(V)) = \mathcal{O}_X(b^{-1}(V))$.
Comme cela vaut pour tout ouvert $V$ et tout $f \in \mathcal{O}_Y(V)$,
nous concluons que $a = b$ comme morphismes de sch\'emas.
\end{proof}

\begin{lemma}
\label{lemma-morphism-ringed-etale-topoi-affines}
Soient $X$, $Y$ des sch\'emas affines.
Soit
$$
(g, g^\#) :
(\Sh(X_\etale), \mathcal{O}_X)
\longrightarrow
(\Sh(Y_\etale), \mathcal{O}_Y)
$$
un morphisme de topos localement annel\'es. Il existe alors un
unique morphisme de sch\'emas $f : X \to Y$ tel que
$(g, g^\#)$ soit $2$-isomorphe \`a $(f_{small}, f_{small}^\sharp)$,
voir
Modules sur les sites,
définition \ref{sites-modules-definition-2-morphism-ringed-topoi}.
\end{lemma}

\begin{proof}
Dans cette d\'emonstration, nous notons $\mathcal{O}_X$ le faisceau structural
du petit site \'etale $X_\etale$, et de m\^eme pour
$\mathcal{O}_Y$. Posons $Y = \Spec(B)$ et $X = \Spec(A)$. Comme
$B = \Gamma(Y_\etale, \mathcal{O}_Y)$,
$A = \Gamma(X_\etale, \mathcal{O}_X)$,
nous voyons que $g^\sharp$ induit un morphisme d'anneaux $\varphi : B \to A$.
Soit $f = \Spec(\varphi) : X \to Y$ le morphisme correspondant
de sch\'emas affines. Nous allons montrer que ce $f$ convient.

\medskip\noindent
Soit $V \to Y$ un sch\'ema affine \'etale sur $Y$. Nous pouvons donc \`ecrire
$V = \Spec(C)$, o\`u $C$ est une $B$-alg\`ebre \'etale. Nous pouvons \`ecrire
$$
C = B[x_1, \ldots, x_n]/(P_1, \ldots, P_n)
$$
avec des polyn\^omes $P_i$ tels que $\Delta = \det(\partial P_i/ \partial x_j)$
soit inversible dans $C$, voir par exemple
Algèbre, lemme \ref{algebra-lemma-etale-standard-smooth}.
Si $T$ est un sch\'ema sur $Y$, alors un $T$-point de $V$ est donn\'e par
$n$ sections de $\Gamma(T, \mathcal{O}_T)$ qui satisfont aux \`equations polynomiales
$P_1 = 0, \ldots, P_n = 0$. Autrement dit, le faisceau $h_V$
sur $Y_\etale$ est l'\'egalisateur des deux morphismes
$$
\xymatrix{
\prod\nolimits_{i = 1, \ldots, n} \mathcal{O}_Y
\ar@<1ex>[r]^a \ar@<-1ex>[r]_b &
\prod\nolimits_{j = 1, \ldots, n} \mathcal{O}_Y
}
$$
o\`u $b(h_1, \ldots, h_n) = 0$ et
$a(h_1, \ldots, h_n) =
(P_1(h_1, \ldots, h_n), \ldots, P_n(h_1, \ldots, h_n))$.
Comme $g^{-1}$ est exact, nous concluons que la ligne sup\'erieure de la
partie du diagramme commutatif suivant form\'ee par les fl\`eches pleines est elle aussi un diagramme d'\'egalisateur :
$$
\xymatrix{
g^{-1}h_V \ar[r] \ar@{..>}[d] &
\prod\nolimits_{i = 1, \ldots, n} g^{-1}\mathcal{O}_Y
\ar@<1ex>[r]^{g^{-1}a} \ar@<-1ex>[r]_{g^{-1}b} \ar[d]^{\prod g^\sharp} &
\prod\nolimits_{j = 1, \ldots, n} g^{-1}\mathcal{O}_Y \ar[d]^{\prod g^\sharp}\\
h_{X \times_Y V} \ar[r] &
\prod\nolimits_{i = 1, \ldots, n} \mathcal{O}_X
\ar@<1ex>[r]^{a'} \ar@<-1ex>[r]_{b'} &
\prod\nolimits_{j = 1, \ldots, n} \mathcal{O}_X  \\
}
$$
Ici, $b'$ est l'application nulle et $a'$ est l'application d\'efinie par les
images $P'_i = \varphi(P_i) \in A[x_1, \ldots, x_n]$ suivant la m\^eme
r\`egle
$a'(h_1, \ldots, h_n) =
(P'_1(h_1, \ldots, h_n), \ldots, P'_n(h_1, \ldots, h_n))$.
que celle qui d\'efinit $a$. La commutativit\'e du diagramme r\'esulte de
l'\'egalit\'e $\varphi = g^\sharp$ sur les sections globales. La ligne
inf\'erieure est \'egalement un diagramme \'egalisateur, exactement par les m\^emes arguments que
pr\'ec\'edemment, puisque $X \times_Y V$ est le sch\'ema affine
$\Spec(A \otimes_B C)$ et
$A \otimes_B C = A[x_1, \ldots, x_n]/(P'_1, \ldots, P'_n)$.
Nous obtenons donc une unique fl\`eche pointill\'ee
$g^{-1}h_V \to h_{X \times_Y V}$ qui s'ins\`ere dans le diagramme.

\medskip\noindent
Nous affirmons que le morphisme de faisceaux $g^{-1}h_V \to h_{X \times_Y V}$
est un isomorphisme. Comme le petit site \'etale de $X$ poss\`ede suffisamment de points
(Th\'eor\`eme \ref{theorem-exactness-stalks}),
il suffit de le d\'emontrer sur les germes. Soit donc $\overline{x}$ un
point g\'eom\'etrique de $X$, et notons $p$ le point associ\'e du
petit topos \'etale de $X$. Posons $q = g \circ p$. C'est un point du
petit topos \'etale de $Y$. D'apr\`es le
Lemme \ref{lemma-points-small-etale-site},
nous voyons que $q$ correspond \`a un point g\'eom\'etrique $\overline{y}$ de
$Y$. Consid\'erons le morphisme de germes
$$
(g^\sharp)_p :
(\mathcal{O}_Y)_{\overline{y}} =
\mathcal{O}_{Y, q} =
(g^{-1}\mathcal{O}_Y)_p
\longrightarrow
\mathcal{O}_{X, p} =
(\mathcal{O}_X)_{\overline{x}}
$$
Comme $(g, g^\sharp)$ est un morphisme de topos {\it localement} annel\'es,
$(g^\sharp)_p$ est un homomorphisme local entre anneaux locaux strictement
hens\'eliens. En localisant le grand diagramme commutatif ci-dessus et en appliquant
Alg\`ebre, Lemme \ref{algebra-lemma-strictly-henselian-solutions},
nous concluons que $(g^{-1}h_V)_p \to (h_{X \times_Y V})_p$ est un isomorphisme,
comme voulu.

\medskip\noindent
Nous affirmons que les isomorphismes $g^{-1}h_V \to h_{X \times_Y V}$ sont
fonctoriels. Plus pr\'ecis\'ement, supposons que $V_1 \to V_2$ soit un morphisme de sch\'emas affines
\'etales sur $Y$. \'Ecrivons
$V_i = \Spec(C_i)$ avec
$$
C_i = B[x_{i, 1}, \ldots, x_{i, n_i}]/(P_{i, 1}, \ldots, P_{i, n_i})
$$
Le morphisme $V_1 \to V_2$ est donn\'e par un homomorphisme de $B$-alg\`ebres $C_2 \to C_1$
qui, \`a son tour, est donn\'e par des polyn\^omes
$Q_j \in B[x_{1, 1}, \ldots, x_{1, n_1}]$, pour $j = 1, \ldots, n_2$.
Il est alors ais\'e de montrer que le diagramme de faisceaux
$$
\xymatrix{
h_{V_1} \ar[d] \ar[r] & \prod_{i = 1, \ldots, n_1} \mathcal{O}_Y
\ar[d]^{Q_1, \ldots, Q_{n_2}}\\
h_{V_2} \ar[r] & \prod_{i = 1, \ldots, n_2} \mathcal{O}_Y
}
$$
est commutatif et que, par image inverse sur $X_\etale$, nous obtenons le
diagramme commutatif \`a fl\`eches pleines
$$
\xymatrix{
g^{-1}h_{V_1} \ar@{..>}[dd] \ar[rrd] \ar[r] &
\prod_{i = 1, \ldots, n_1} g^{-1}\mathcal{O}_Y
\ar[dd]^{g^\sharp}
\ar[rrd]^{Q_1, \ldots, Q_{n_2}} \\
& & g^{-1}h_{V_2} \ar@{..>}[dd] \ar[r] &
\prod_{i = 1, \ldots, n_2} g^{-1}\mathcal{O}_Y
\ar[dd]^{g^\sharp} \\
h_{X \times_Y V_1} \ar[r] \ar[rrd] &
\prod\nolimits_{i = 1, \ldots, n_1} \mathcal{O}_X
\ar[rrd]^{Q'_1, \ldots, Q'_{n_2}} \\
& & h_{X \times_Y V_2} \ar[r] &
\prod\nolimits_{i = 1, \ldots, n_2} \mathcal{O}_X
}
$$
o\`u $Q'_j \in A[x_{1, 1}, \ldots, x_{1, n_1}]$ est l'image de
$Q_j$ par $\varphi$. Puisque les fl\`eches pointill\'ees existent, rendent les
deux carr\'es commutatifs, et que les fl\`eches horizontales sont injectives,
nous voyons que le diagramme entier est commutatif. Cela prouve la fonctorialit\'e
(ainsi que l'ind\'ependance de la construction de $g^{-1}h_V \to h_{X \times_Y V}$
\`a l'\'egard du choix de la pr\'esentation, bien qu'\`a
strictement parler nous n'ayons pas besoin de le montrer).

\medskip\noindent
\`A ce stade, nous pouvons montrer que $f_{small, *} \cong g_*$.
En effet, soit $\mathcal{F}$ un faisceau sur $X_\etale$. Pour tout objet affine
$V \in \Ob(X_\etale)$, nous avons
\begin{align*}
(g_*\mathcal{F})(V)
& =
\Mor_{\Sh(Y_\etale)}(h_V, g_*\mathcal{F}) \\
& =
\Mor_{\Sh(X_\etale)}(g^{-1}h_V, \mathcal{F}) \\
& =
\Mor_{\Sh(X_\etale)}(h_{X \times_Y V}, \mathcal{F}) \\
& =
\mathcal{F}(X \times_Y V) \\
& =
f_{small, *}\mathcal{F}(V)
\end{align*}
o\`u, dans la troisi\`eme \'egalit\'e, nous utilisons l'isomorphisme
$g^{-1}h_V \cong h_{X \times_Y V}$ construit ci-dessus. Ces isomorphismes
sont manifestement fonctoriels en $\mathcal{F}$ et en $V$,
car les isomorphismes $g^{-1}h_V \cong h_{X \times_Y V}$ sont fonctoriels.
Or tout faisceau sur $Y_\etale$ est d\'etermin\'e par sa restriction
\`a la sous-cat\'egorie des sch\'emas affines
(Topologies, Lemme \ref{topologies-lemma-alternative}) ;
nous obtenons donc un isomorphisme de foncteurs $f_{small, *} \cong g_*$,
comme voulu.

\medskip\noindent
Enfin, nous devons v\'erifier que, par l'isomorphisme
$f_{small, *} \cong g_*$ ci-dessus, les morphismes $f_{small}^\sharp$ et
$g^\sharp$ co\"incident. Par construction, c'est d\'ej\`a le cas sur les
sections globales de $\mathcal{O}_Y$, c'est-\`a-dire sur les \'el\'ements de $B$.
Il suffit de v\'erifier le r\'esultat sur les
sections au-dessus d'un ouvert affine $V$ \'etale sur $Y$ (encore d'apr\`es
Topologies, Lemme \ref{topologies-lemma-alternative}).
En \'ecrivant
$V = \Spec(C)$, $C = B[x_i]/(P_j)$ comme pr\'ec\'edemment, il suffit
de v\'erifier que les fonctions coordonn\'ees $x_i$ sont envoy\'ees sur
les m\^emes sections de $\mathcal{O}_X$ au-dessus de $X \times_Y V$.
Et c'est exactement ce que signifie la commutativit\'e du diagramme
$$
\xymatrix{
g^{-1}h_V \ar[r] \ar@{..>}[d] &
\prod\nolimits_{i = 1, \ldots, n} g^{-1}\mathcal{O}_Y
\ar[d]^{\prod g^\sharp} \\
h_{X \times_Y V} \ar[r] &
\prod\nolimits_{i = 1, \ldots, n} \mathcal{O}_X
}
$$
Le lemme est donc d\'emontr\'e.
\end{proof}

\noindent
Voici une version pour des sch\'emas quelconques.

\begin{theorem}
\label{theorem-fully-faithful}
Soient $X$, $Y$ des sch\'emas. Soit
$$
(g, g^\#) :
(\Sh(X_\etale), \mathcal{O}_X)
\longrightarrow
(\Sh(Y_\etale), \mathcal{O}_Y)
$$
un morphisme de topos localement annel\'es. Alors il existe un unique
morphisme de sch\'emas $f : X \to Y$ tel que
$(g, g^\#)$ soit isomorphe \`a $(f_{small}, f_{small}^\sharp)$.
Autrement dit, la construction
$$
\Sch \longrightarrow \textit{Topos localement annel\'es},
\quad
X \longrightarrow (X_\etale, \mathcal{O}_X)
$$
est pleinement fid\`ele (les morphismes du membre de droite sont consid\'er\'es \`a $2$-isomorphisme pr\`es).
\end{theorem}

\begin{proof}
On peut d\'emontrer ce th\'eor\`eme en adaptant soigneusement les arguments de
la d\'emonstration du
Lemme \ref{lemma-morphism-ringed-etale-topoi-affines}
au cadre global. Nous voulons toutefois indiquer comment
recoller le r\'esultat de ce lemme pour obtenir un morphisme global,
gr\^ace \`a la rigidit\'e fournie par le r\'esultat du
Lemme \ref{lemma-2-morphism}.
Malheureusement, les d\'etails sont quelque peu laborieux.

\medskip\noindent
Montrons l'existence lorsque $Y$ est affine. Choisissons alors un
recouvrement ouvert affine $X = \bigcup U_i$. Pour chaque $i$, le morphisme
d'inclusion $j_i : U_i \to X$ induit un morphisme de topos localement annel\'es
$(j_{i, small}, j_{i, small}^\sharp) :
(\Sh(U_{i, \etale}), \mathcal{O}_{U_i})
\to
(\Sh(X_\etale), \mathcal{O}_X)$
d'apr\`es le
Lemme \ref{lemma-morphism-locally-ringed}.
Nous pouvons le composer avec $(g, g^\sharp)$ pour obtenir un morphisme
de topos localement annel\'es
$$
(g, g^\sharp) \circ (j_{i, small}, j_{i, small}^\sharp) :
(\Sh(U_{i, \etale}), \mathcal{O}_{U_i})
\to
(\Sh(Y_\etale), \mathcal{O}_Y)
$$
voir
Modules sur les sites,
Lemme \ref{sites-modules-lemma-composition-morphisms-locally-ringed-topoi}.
D'apr\`es le
Lemme \ref{lemma-morphism-ringed-etale-topoi-affines},
il existe un unique morphisme de sch\'emas $f_i : U_i \to Y$
et un $2$-isomorphisme
$$
t_i :
(f_{i, small}, f_{i, small}^\sharp)
\longrightarrow
(g, g^\sharp) \circ (j_{i, small}, j_{i, small}^\sharp).
$$
Posons $U_{i, i'} = U_i \cap U_{i'}$, et notons $j_{i, i'} : U_{i, i'} \to U_i$
le morphisme d'inclusion. Puisque
$j_i \circ j_{i, i'} = j_{i'} \circ j_{i', i}$
nous avons
\begin{align*}
(g, g^\sharp) \circ
(j_{i, small}, j_{i, small}^\sharp) \circ
(j_{i, i', small}, j_{i, i', small}^\sharp)
= \\
(g, g^\sharp) \circ
(j_{i', small}, j_{i', small}^\sharp) \circ
(j_{i', i, small}, j_{i', i, small}^\sharp)
\end{align*}
Par unicit\'e (voir le
Lemme \ref{lemma-faithful}),
nous en concluons que
$f_i \circ j_{i, i'} = f_{i'} \circ j_{i', i}$ ; autrement dit, les
morphismes de sch\'emas $f_i = f \circ j_i$ sont les restrictions d'un
morphisme de sch\'emas global $f : X \to Y$. Consid\'erons le diagramme
de $2$-isomorphismes (dans lequel nous omettons les composantes ${}^\sharp$ pour all\'eger la
notation)
$$
\xymatrix{
g \circ j_{i, small} \circ j_{i, i', small}
\ar[rr]^{t_i \star \text{id}_{j_{i, i', small}}}
\ar@{=}[d] & &
f_{small} \circ j_{i, small} \circ j_{i, i', small} \ar@{=}[d] \\
g \circ j_{i', small} \circ j_{i', i, small}
\ar[rr]^{t_{i'} \star \text{id}_{j_{i', i, small}}} & &
f_{small} \circ j_{i', small} \circ j_{i', i, small}
}
$$
La notation $\star$ d\'esigne la composition horizontale ; voir
Catégories, D\'efinition \ref{categories-definition-2-category}
dans le cas g\'en\'eral et
Sites, section \ref{sites-section-2-category}
dans notre cas particulier. D'apr\`es le r\'esultat du
Lemme \ref{lemma-2-morphism},
ce diagramme est commutatif. Ainsi, pour tout faisceau $\mathcal{G}$
sur $Y_\etale$, les isomorphismes
$t_i : f_{small}^{-1}\mathcal{G}|_{U_i} \to g^{-1}\mathcal{G}|_{U_i}$
co\"incident au-dessus de $U_{i, i'}$ et nous obtenons un isomorphisme global
$t : f_{small}^{-1}\mathcal{G} \to g^{-1}\mathcal{G}$.
Il est clair que cet isomorphisme est fonctoriel en $\mathcal{G}$
et compatible avec les morphismes $f_{small}^\sharp$ et $g^\sharp$
(puisqu'il est localement compatible avec ces morphismes).
Cela d\'emontre le th\'eor\`eme lorsque $Y$ est affine.

\medskip\noindent
Dans le cas g\'en\'eral, soit $V \subset Y$ un ouvert affine.
Alors $h_V$ est un sous-faisceau du faisceau final $*$ sur $Y_\etale$.
Comme $g$ est exact, nous voyons que $g^{-1}h_V$ est un sous-faisceau du faisceau
final sur $X_\etale$. Ainsi, d'apr\`es le
Lemme \ref{lemma-support-subsheaf-final},
il existe un sous-sch\'ema ouvert $W \subset X$ tel que $g^{-1}h_V = h_W$. D'apr\`es
Modules sur les sites,
Lemme \ref{sites-modules-lemma-localize-morphism-locally-ringed-topoi},
il existe un diagramme commutatif de morphismes de topos localement
annel\'es
$$
\xymatrix{
(\Sh(W_\etale), \mathcal{O}_W) \ar[r] \ar[d]_{g'} &
(\Sh(X_\etale), \mathcal{O}_X) \ar[d]^g \\
(\Sh(V_\etale), \mathcal{O}_V) \ar[r] &
(\Sh(Y_\etale), \mathcal{O}_Y)
}
$$
o\`u les fl\`eches horizontales sont les morphismes de localisation
(induits par les morphismes d'inclusion $V \to Y$ et $W \to X$)
et o\`u $g'$ est induit par $g$. Le paragraphe pr\'ec\'edent
nous fournit un morphisme de sch\'emas $f' : W \to V$ et
un $2$-isomorphisme
$t : (f'_{small}, (f'_{small})^\sharp) \to (g', (g')^\sharp)$.
Exactement comme pr\'ec\'edemment, ces morphismes $f'$ (pour les divers ouverts affines $V \subset Y$)
co\"incident sur les intersections par unicit\'e, et donnent donc un morphisme $f : X \to Y$.
De plus, les $2$-isomorphismes $t$ sont compatibles sur les intersections, encore d'apr\`es le
Lemme \ref{lemma-2-morphism},
et nous obtenons un $2$-isomorphisme global
$(f_{small}, (f_{small})^\sharp) \to (g, (g)^\sharp)$.
comme voulu. Nous omettons certains d\'etails.
\end{proof}










\section{Image directe et image inverse}
\label{section-monomorphisms}

\noindent
Soit $f : X \to Y$ un morphisme de sch\'emas.
Voici une liste de conditions que nous consid\'ererons dans la suite :
\begin{enumerate}
\item[(A)] Pour tout morphisme \'etale $U \to X$ et tout $u \in U$, il existe
un morphisme \'etale $V \to Y$, une d\'ecomposition en union disjointe
$X \times_Y V = W \amalg W'$ et un morphisme $h : W \to U$ au-dessus de $X$
tel que $u$ appartienne \`a l'image de $h$.
\item[(B)] Pour tout morphisme \'etale $V \to Y$ et tout recouvrement \'etale
$\{U_i \to X \times_Y V\}$, il existe un recouvrement \'etale
$\{V_j \to V\}$ tel que, pour chaque $j$, on ait
$X \times_Y V_j = \coprod W_{ij}$, o\`u $W_{ij} \to X \times_Y V$
se factorise par $U_i \to X \times_Y V$ pour un certain $i$.
\item[(C)] Pour tout morphisme \'etale $U \to X$, il existe un morphisme \'etale $V \to Y$
et un morphisme surjectif $X \times_Y V \to U$ au-dessus de $X$.
\end{enumerate}
Il s'av\`ere que chacune de ces propri\'et\'es s'interpr\`ete en termes du
comportement du foncteur $f_{small, *}$. Nous expliciterons cela
dans les quelques sections suivantes.



\section{Propri\'et\'e (A)}
\label{section-A}

\noindent
Voir la section \ref{section-monomorphisms} pour la d\'efinition de la propri\'et\'e
(A).

\begin{lemma}
\label{lemma-property-A-implies}
Soit $f : X \to Y$ un morphisme de sch\'emas.
Supposons (A).
\begin{enumerate}
\item
$f_{small, *} :
\textit{Ab}(X_\etale)
\to
\textit{Ab}(Y_\etale)$
refl\`ete les injections et les surjections,
\item $f_{small}^{-1}f_{small, *}\mathcal{F} \to \mathcal{F}$
est surjectif pour tout faisceau ab\'elien $\mathcal{F}$ sur $X_\etale$,
\item
$f_{small, *} :
\textit{Ab}(X_\etale)
\to
\textit{Ab}(Y_\etale)$
est fid\`ele.
\end{enumerate}
\end{lemma}

\begin{proof}
Soit $\mathcal{F}$ un faisceau ab\'elien sur $X_\etale$.
Soit $U$ un objet de $X_\etale$. Par hypoth\`ese, nous pouvons trouver un
recouvrement $\{W_i \to U\}$ dans $X_\etale$ tel que chaque $W_i$ soit
un sous-sch\'ema ouvert et ferm\'e de $X \times_Y V_i$ pour un certain objet
$V_i$ de $Y_\etale$. La condition de faisceau montre que
$$
\mathcal{F}(U) \subset \prod \mathcal{F}(W_i)
$$
et que $\mathcal{F}(W_i)$ est un facteur direct de
$\mathcal{F}(X \times_Y V_i) = f_{small, *}\mathcal{F}(V_i)$.
Il est donc clair que $f_{small, *}$ refl\`ete les injections.

\medskip\noindent
Supposons ensuite que $a : \mathcal{G} \to \mathcal{F}$ soit un morphisme de
faisceaux ab\'eliens tel que $f_{small, *}a$ soit surjectif. Soit
$s \in \mathcal{F}(U)$, avec $U$ comme ci-dessus. Avec les $W_i$, $V_i$ ci-dessus,
nous voyons qu'il suffit de montrer que $s|_{W_i}$ est, localement pour la topologie \'etale,
l'image d'une section de $\mathcal{G}$ par $a$. Comme $\mathcal{F}(W_i)$
est un facteur direct de $\mathcal{F}(X \times_Y V_i)$,
il suffit de montrer que, pour tout $V \in \Ob(Y_\etale)$,
tout \'el\'ement $s \in \mathcal{F}(X \times_Y V)$
est, localement pour la topologie \'etale sur $X \times_Y V$, l'image d'une section de
$\mathcal{G}$ par $a$. Puisque
$\mathcal{F}(X \times_Y V) = f_{small, *}\mathcal{F}(V)$
nous voyons par hypoth\`ese qu'il existe un recouvrement $\{V_j \to V\}$ tel que
$s$ soit l'image de
$s_j \in f_{small, *}\mathcal{G}(V_j) = \mathcal{G}(X \times_Y V_j)$.
Cela prouve que $f_{small, *}$ refl\`ete les surjections.

\medskip\noindent
Les assertions (2) et (3) d\'ecoulent formellement de l'assertion (1) ; voir
Modules sur les sites, Lemme \ref{sites-modules-lemma-reflect-surjections}.
\end{proof}

\begin{lemma}
\label{lemma-locally-quasi-finite-A}
Soit $f : X \to Y$ un morphisme de sch\'emas s\'epar\'e et localement quasi-fini.
Alors la propri\'et\'e (A) ci-dessus est satisfaite.
\end{lemma}

\begin{proof}
Soient $U \to X$ un morphisme \'etale et $u \in U$.
L'\'enonc\'e g\'eom\'etrique (A) se r\'eduit directement au cas o\`u $U$ et $Y$
sont des sch\'emas affines. Notons $x \in X$ et $y \in Y$ les
images de $u$. Comme $X \to Y$ est localement quasi-fini et que $U \to X$ est
localement quasi-fini (voir
Morphismes, Lemme \ref{morphisms-lemma-etale-locally-quasi-finite}),
nous voyons que $U \to Y$ est localement quasi-fini (voir
Morphismes, Lemme \ref{morphisms-lemma-composition-quasi-finite}).
De plus, $X \to Y$ et $U \to Y$ sont tous deux s\'epar\'es. Ainsi, le
Compléments sur les morphismes, Lemme
\ref{more-morphisms-lemma-etale-splits-off-quasi-finite-part-technical-variant}
s'applique aux deux morphismes. Nous pouvons donc choisir un voisinage \'etale
$(V, v) \to (Y, y)$ tel que
$$
X \times_Y V = W \amalg R, \quad
U \times_Y V = W' \amalg R'
$$
et des points $w \in W$, $w' \in W'$ tels que
\begin{enumerate}
\item $W$, $R$ sont ouverts et ferm\'es dans $X \times_Y V$,
\item $W'$, $R'$ sont ouverts et ferm\'es dans $U \times_Y V$,
\item $W \to V$ et $W' \to V$ sont finis,
\item $w$, $w'$ s'envoient sur $v$,
\item $\kappa(v) \subset \kappa(w)$ et $\kappa(v) \subset \kappa(w')$
sont purement ins\'eparables, et
\item aucun autre point de $W$ ou de $W'$ ne s'envoie sur $v$.
\end{enumerate}
Voici un diagramme commutatif
$$
\xymatrix{
U \ar[d] & U \times_Y V \ar[l] \ar[d] & W' \amalg R' \ar[d] \ar[l] \\
X \ar[d] & X \times_Y V \ar[l] \ar[d] & W \amalg R \ar[l] \\
Y & V \ar[l]
}
$$
Apr\`es avoir r\'etr\'eci $V$, nous pouvons supposer que $W'$ s'envoie dans $W$ :
il suffit d'enlever l'image de l'image inverse de $R$ dans $W'$ ; c'est
un ferm\'e (puisque $W' \to V$ est fini) qui ne contient pas $v$.
Alors $W' \to W$ est fini, puisque $W \to V$ et $W' \to V$ sont finis.
Ainsi $W' \to W$ est fini \'etale, et la
fibre au-dessus de $w$ contient exactement un point, avec $\kappa(w) = \kappa(w')$. Par cons\'equent, $W' \to W$ est un
isomorphisme sur un voisinage ouvert $W^\circ$ de $w$ ; voir
Morphismes étales, Lemme \ref{etale-lemma-finite-etale-one-point}.
Comme $W \to V$ est fini, l'image de $W \setminus W^\circ$ est un sous-ensemble
ferm\'e $T$ de $V$ qui ne contient pas $v$. Apr\`es avoir remplac\'e $V$ par
$V \setminus T$, nous pouvons donc supposer que $W' \to W$ est un isomorphisme.
La d\'ecomposition $X \times_Y V = W \amalg R$ et le morphisme
$W \to U$ sont alors ceux recherch\'es, ce qui ach\`eve la d\'emonstration.
\end{proof}

\begin{lemma}
\label{lemma-integral-A}
Soit $f : X \to Y$ un morphisme entier de sch\'emas.
Alors la propri\'et\'e (A) est satisfaite.
\end{lemma}

\begin{proof}
Soient $U \to X$ \'etale et $u \in U$ un point.
Nous devons trouver un morphisme \'etale $V \to Y$, une d\'ecomposition en union disjointe
$X \times_Y V = W \amalg W'$ et un $X$-morphisme $W \to U$
tel que $u$ appartienne \`a son image. Nous pouvons r\'etr\'ecir $U$ et $Y$ et supposer
que $U$ et $Y$ sont affines. Alors $X$ est \'egalement affine, puisqu'un
morphisme entier est affine par d\'efinition. \'Ecrivons $Y = \Spec(A)$,
$X = \Spec(B)$ et $U = \Spec(C)$. Alors $A \to B$ est un
homomorphisme d'anneaux entier, et $B \to C$ un homomorphisme d'anneaux \'etale. D'apr\`es
Algèbre, Lemme \ref{algebra-lemma-etale},
nous pouvons trouver une sous-$A$-alg\`ebre finie $B' \subset B$ et un homomorphisme d'anneaux \'etale
$B' \to C'$ tels que $C = B \otimes_{B'} C'$. La question
se ram\`ene donc au morphisme \'etale
$U' = \Spec(C') \to X' = \Spec(B')$
au-dessus du morphisme fini $X' \to Y$. Dans ce cas, le r\'esultat d\'ecoule du
Lemme \ref{lemma-locally-quasi-finite-A}.
\end{proof}

\begin{lemma}
\label{lemma-when-push-pull-surjective}
Soit $f : X \to Y$ un morphisme de sch\'emas. Notons
$f_{small} :
\Sh(X_\etale)
\to
\Sh(Y_\etale)$
le morphisme associ\'e de petits topos \'etales. Supposons qu'au moins l'une
des hypoth\`eses suivantes soit satisfaite :
\begin{enumerate}
\item $f$ est entier, ou
\item $f$ est s\'epar\'e et localement quasi-fini.
\end{enumerate}
Alors le foncteur
$f_{small, *} :
\textit{Ab}(X_\etale)
\to
\textit{Ab}(Y_\etale)$
poss\`ede les propri\'et\'es suivantes :
\begin{enumerate}
\item le morphisme
$f_{small}^{-1}f_{small, *}\mathcal{F} \to \mathcal{F}$
est toujours surjectif,
\item $f_{small, *}$ est fid\`ele, et
\item $f_{small, *}$ refl\`ete les injections et les surjections.
\end{enumerate}
\end{lemma}

\begin{proof}
Il suffit de combiner les
Lemmes \ref{lemma-locally-quasi-finite-A},
\ref{lemma-integral-A}, et
\ref{lemma-property-A-implies}.
\end{proof}



\section{Propri\'et\'e (B)}
\label{section-B}

\noindent
Voir la section \ref{section-monomorphisms} pour la d\'efinition de la propri\'et\'e
(B).

\begin{lemma}
\label{lemma-property-B-implies}
Soit $f : X \to Y$ un morphisme de sch\'emas. Supposons la propri\'et\'e (B) satisfaite.
Alors le foncteur
$f_{small, *} :
\Sh(X_\etale)
\to
\Sh(Y_\etale)$
envoie les morphismes surjectifs sur des morphismes surjectifs.
\end{lemma}

\begin{proof}
Cela d\'ecoule de
Sites, Lemme \ref{sites-lemma-weaker}.
\end{proof}

\begin{lemma}
\label{lemma-simplify-B}
Soit $f : X \to Y$ un morphisme de sch\'emas. Supposons que
\begin{enumerate}
\item $V \to Y$ soit un morphisme \'etale de sch\'emas,
\item $\{U_i \to X \times_Y V\}$ soit un recouvrement \'etale, et
\item $v \in V$ soit un point.
\end{enumerate}
Supposons que, pour toute donn\'ee de ce type, il existe un voisinage \'etale
$(V', v') \to (V, v)$, une d\'ecomposition en union disjointe
$X \times_Y V' = \coprod W'_i$ et des morphismes $W'_i \to U_i$
au-dessus de $X \times_Y V$. Alors la propri\'et\'e (B) est satisfaite.
\end{lemma}

\begin{proof}
D\'emonstration omise.
\end{proof}

\begin{lemma}
\label{lemma-finite-B}
Soit $f : X \to Y$ un morphisme fini de sch\'emas.
Alors la propri\'et\'e (B) est satisfaite.
\end{lemma}

\begin{proof}
Consid\'erons un morphisme \'etale $V \to Y$, un recouvrement \'etale $\{U_i \to X \times_Y V\}$ et
un point $v \in V$. Nous devons trouver $V' \to V$, une d\'ecomposition et des morphismes comme dans le
Lemme \ref{lemma-simplify-B}.
Nous pouvons r\'etr\'ecir $V$ et $Y$, donc supposer que $V$ et $Y$ sont affines.
Comme $X$ est fini sur $Y$, il en r\'esulte aussi que $X$ est affine.
Au cours de la d\'emonstration, nous pouvons remplacer (un nombre fini de fois) $(V, v)$ par un
voisinage \'etale $(V', v')$ et, de mani\`ere correspondante, le recouvrement
$\{U_i \to X \times_Y V\}$ par $\{V' \times_V U_i \to X \times_Y V'\}$.

\medskip\noindent
Comme $X \times_Y V \to V$ est fini, il existe un nombre fini de
points deux \`a deux distincts $x_1, \ldots, x_n \in X \times_Y V$
qui s'envoient sur $v$. Nous pouvons appliquer
Compléments sur les morphismes, Lemme
\ref{more-morphisms-lemma-etale-splits-off-quasi-finite-part-technical-variant}
\`a $X \times_Y V \to V$ et aux points $x_1, \ldots, x_n$ situ\'es au-dessus de
$v$, et trouver un voisinage \'etale $(V', v') \to (V, v)$
tel que
$$
X \times_Y V' = R \amalg \coprod T_a
$$
o\`u $T_a \to V'$ est fini, avec exactement un point $p_a$ au-dessus de $v'$,
o\`u de plus $\kappa(v') \subset \kappa(p_a)$ est purement ins\'eparable, et
tel que la fibre de $R \to V'$ au-dessus de $v'$ soit vide.
Puisque $X \to Y$ est fini, $R \to V'$ l'est aussi. Apr\`es avoir
r\'etr\'eci $V'$, nous pouvons donc supposer que $R = \emptyset$. Ainsi, nous pouvons
supposer que $X \times_Y V = X_1 \amalg \ldots \amalg X_n$, avec
exactement un point $x_l \in X_l$ au-dessus de $v$, et tel que, de plus,
$\kappa(v) \subset \kappa(x_l)$ soit purement ins\'eparable. Remarquons que cette
propri\'et\'e est pr\'eserv\'ee par tout raffinement du voisinage \'etale
$(V, v)$.

\medskip\noindent
Pour chaque $l$, choisissons un $i_l$ et un point $u_l \in U_{i_l}$ qui s'envoie sur $x_l$.
Appliquons maintenant la propri\'et\'e (A) au morphisme fini
$X \times_Y V \to V$, aux morphismes \'etales
$U_{i_l} \to X \times_Y V$ et aux points $u_l$.
Cela est permis par le
Lemme \ref{lemma-integral-A}.
Nous obtenons ainsi un voisinage \'etale $(V', v') \to (V, v)$
et des d\'ecompositions
$$
X \times_Y V' = W_l \amalg R_l
$$
et des $X$-morphismes $a_l : W_l \to U_{i_l}$ dont l'image contient $u_{i_l}$.
Voici une figure :
$$
\xymatrix{
& & & U_{i_l} \ar[d] & \\
W_l \ar[rrru] \ar[r] & W_l \amalg R_l \ar@{=}[r] &
X \times_Y V' \ar[r] \ar[d] &
X \times_Y V \ar[r] \ar[d] & X \ar[d] \\
& & V' \ar[r] & V \ar[r] & Y
}
$$
Apr\`es avoir remplac\'e $(V, v)$ par $(V', v')$, nous concluons que chaque
$x_l$ appartient \`a un voisinage ouvert et ferm\'e $W_l$ tel que
le morphisme d'inclusion $W_l \to X \times_Y V$ se factorise par
$U_i \to X \times_Y V$ pour un certain $i$. En rempla\c{c}ant $W_l$ par $W_l \cap X_l$,
nous voyons que ces ouverts et ferm\'es sont disjoints et, de plus,
que $\{x_1, \ldots, x_n\} \subset W_1 \cup \ldots \cup W_n$.
Comme $X \times_Y V \to V$ est fini, nous pouvons r\'etr\'ecir $V$ et supposer que
$X \times_Y V = W_1 \amalg \ldots \amalg W_n$, comme voulu.
\end{proof}

\begin{lemma}
\label{lemma-integral-B}
Soit $f : X \to Y$ un morphisme entier de sch\'emas.
Alors la propri\'et\'e (B) est satisfaite.
\end{lemma}

\begin{proof}
Consid\'erons un morphisme \'etale $V \to Y$, un recouvrement \'etale $\{U_i \to X \times_Y V\}$ et
un point $v \in V$. Nous devons trouver $V' \to V$, une d\'ecomposition et des morphismes comme dans le
Lemme \ref{lemma-simplify-B}.
Nous pouvons r\'etr\'ecir $V$ et $Y$, donc supposer que $V$ et $Y$ sont affines.
Comme $X$ est entier sur $Y$, il en r\'esulte aussi que $X$ et
$X \times_Y V$ sont affines. Nous pouvons raffiner le recouvrement
$\{U_i \to X \times_Y V\}$ et donc supposer que
$\{U_i \to X \times_Y V\}_{i = 1, \ldots, n}$ est un recouvrement \'etale standard.
\'Ecrivons $Y = \Spec(A)$, $X = \Spec(B)$,
$V = \Spec(C)$ et $U_i = \Spec(B_i)$.
Alors $A \to B$ est un homomorphisme d'anneaux entier, et les $B \otimes_A C \to B_i$ sont des
homomorphismes d'anneaux \'etales. D'apr\`es
Algèbre, Lemme \ref{algebra-lemma-etale},
nous pouvons trouver une sous-$A$-alg\`ebre finie $B' \subset B$ et un homomorphisme d'anneaux \'etale
$B' \otimes_A C \to B'_i$, pour $i = 1, \ldots, n$,
tels que $B_i = B \otimes_{B'} B'_i$. La question
se ram\`ene donc au recouvrement \'etale
$\{\Spec(B'_i) \to X' \times_Y V\}_{i = 1, \ldots, n}$
o\`u $X' = \Spec(B')$ est fini sur $Y$.
Dans ce cas, le r\'esultat d\'ecoule du
Lemme \ref{lemma-finite-B}.
\end{proof}

\begin{lemma}
\label{lemma-what-integral}
Soit $f : X \to Y$ un morphisme de sch\'emas.
Supposons $f$ entier (par exemple fini).
Alors
\begin{enumerate}
\item $f_{small, *}$ envoie les morphismes surjectifs sur des morphismes surjectifs (pour les faisceaux
d'ensembles et pour les faisceaux ab\'eliens),
\item $f_{small}^{-1}f_{small, *}\mathcal{F} \to \mathcal{F}$
est surjectif pour tout faisceau ab\'elien $\mathcal{F}$ sur $X_\etale$,
\item
$f_{small, *} :
\textit{Ab}(X_\etale)
\to
\textit{Ab}(Y_\etale)$
est fid\`ele et refl\`ete les injections et les surjections, et
\item
$f_{small, *} :
\textit{Ab}(X_\etale)
\to
\textit{Ab}(Y_\etale)$
est exact.
\end{enumerate}
\end{lemma}

\begin{proof}
Nous avons vu les assertions (2) et (3) dans le
Lemme \ref{lemma-when-push-pull-surjective}.
L'assertion (1) d\'ecoule des
Lemmes \ref{lemma-integral-B} et \ref{lemma-property-B-implies}.
L'assertion (4) est une cons\'equence de l'assertion (1) ; voir
Modules sur les sites, Lemme \ref{sites-modules-lemma-exactness}.
\end{proof}








\section{Propri\'et\'e (C)}
\label{section-C}

\noindent
Voir la section \ref{section-monomorphisms} pour la d\'efinition de la propri\'et\'e
(C).

\begin{lemma}
\label{lemma-property-C-implies}
Soit $f : X \to Y$ un morphisme de sch\'emas. Supposons (C) satisfaite. Alors le foncteur
$f_{small, *} :
\Sh(X_\etale)
\to
\Sh(Y_\etale)$
refl\`ete les injections et les surjections.
\end{lemma}

\begin{proof}
Cela d\'ecoule de
Sites, Lemme \ref{sites-lemma-cover-from-below}.
Nous omettons la v\'erification du fait que la propri\'et\'e (C) implique que le foncteur
$Y_\etale \to X_\etale$, $V \mapsto X \times_Y V$
satisfait l'hypoth\`ese de
Sites, Lemme \ref{sites-lemma-cover-from-below}.
\end{proof}

\begin{remark}
\label{remark-property-C-strong}
La propri\'et\'e (C) est satisfaite si $f : X \to Y$ est une immersion ouverte. En effet, si
$U \in \Ob(X_\etale)$, nous pouvons aussi consid\'erer $U$ comme un objet
de $Y_\etale$ et $U \times_Y X = U$. Ainsi, la propri\'et\'e (C)
n'implique pas que $f_{small, *}$ soit exact, car ce n'est pas
le cas pour les immersions ouvertes (en g\'en\'eral).
\end{remark}

\begin{lemma}
\label{lemma-property-C-closed-implies}
Soit $f : X \to Y$ un morphisme de sch\'emas. Supposons que,
pour tout morphisme \'etale $V \to Y$, on ait
\begin{enumerate}
\item $X \times_Y V \to V$ poss\`ede la propri\'et\'e (C), et
\item $X \times_Y V \to V$ est ferm\'e.
\end{enumerate}
Alors le foncteur
$Y_\etale \to X_\etale$, $V \mapsto X \times_Y V$
est presque cocontinu ; voir
Sites, D\'efinition \ref{sites-definition-almost-cocontinuous}.
\end{lemma}

\begin{proof}
Soit $V \to Y$ un objet de $Y_\etale$ et soit
$\{U_i \to X \times_Y V\}_{i \in I}$ un recouvrement de $X_\etale$.
Par l'hypoth\`ese (1), pour chaque $i$, nous pouvons trouver un morphisme \'etale
$h_i : V_i \to V$ et un morphisme surjectif $X \times_Y V_i \to U_i$
au-dessus de $X \times_Y V$. Remarquons que $\bigcup h_i(V_i) \subset V$ est un
ouvert qui contient le ferm\'e $Z = \Im(X \times_Y V \to V)$.
Soit $h_0 : V_0 = V \setminus Z \to V$ l'immersion ouverte.
Il est clair que $\{V_i \to V\}_{i \in I \cup \{0\}}$ est un
recouvrement \'etale tel que, pour tout $i \in I \cup \{0\}$, on ait
soit $V_i \times_Y X = \emptyset$ (si $i = 0$), soit
que $V_i \times_Y X \to V \times_Y X$ se factorise par $U_i \to X \times_Y V$
(si $i \not = 0$). Le foncteur $Y_\etale \to X_\etale$ est donc
presque cocontinu.
\end{proof}

\begin{lemma}
\label{lemma-integral-homeo-onto-image-C}
Soit $f : X \to Y$ un morphisme entier de sch\'emas qui d\'efinit
un hom\'eomorphisme de $X$ sur un sous-ensemble ferm\'e de $Y$.
Alors la propri\'et\'e (C) est satisfaite.
\end{lemma}

\begin{proof}
Soit $g : U \to X$ un morphisme \'etale. Nous devons trouver un objet
$V \to Y$ de $Y_\etale$ et un morphisme surjectif $X \times_Y V \to U$
au-dessus de $X$. Supposons que, pour tout $u \in U$, nous puissions trouver un objet
$V_u \to Y$ de $Y_\etale$ et un morphisme $h_u : X \times_Y V_u \to U$
au-dessus de $X$ avec $u \in \Im(h_u)$. Nous pouvons alors prendre $V = \coprod V_u$
et $h = \coprod h_u$, ce qui donne le r\'esultat. Ainsi, \'etant donn\'e un point
$u \in U$, construisons une paire $(V_u, h_u)$ comme ci-dessus. Pour ce faire, nous pouvons
r\'etr\'ecir $U$ et supposer que $U$ est affine. Dans ce cas,
$g : U \to X$ est localement quasi-fini. Posons
$g^{-1}(g(\{u\})) = \{u, u_2, \ldots, u_n\}$. Puisqu'il n'existe aucune
sp\'ecialisations $u_i \leadsto u$, nous pouvons remplacer $U$ par un voisinage affine
de sorte que $g^{-1}(g(\{u\})) = \{u\}$.

\medskip\noindent
L'image $g(U) \subset X$ est ouverte ;
par cons\'equent, $f(g(U))$ est localement ferm\'e dans $Y$. Choisissons un ouvert $V \subset Y$ tel
que $f(g(U)) = f(X) \cap V$. Il en r\'esulte que $g$ se factorise par
$X \times_Y V$ et que le morphisme $\{U \to X \times_Y V\}$ ainsi obtenu est un recouvrement
\'etale. Comme $f$ poss\`ede la propri\'et\'e (B), voir le
Lemme \ref{lemma-integral-B},
nous voyons qu'il existe un recouvrement \'etale $\{V_j \to V\}$ tel que
$X \times_Y V_j \to X \times_Y V$ se factorise par $U$.
Il en r\'esulte que $V' = \coprod V_j$ est \'etale sur $Y$ et qu'il existe un
morphisme $h : X \times_Y V' \to U$ dont l'image
se projette surjectivement sur $g(U)$. Comme $u$ est le seul point de sa fibre, il doit
appartenir \`a l'image de $h$, ce qui donne le r\'esultat.
\end{proof}

\noindent
Nous invitons le lecteur \`a consid\'erer le lemme suivant comme une
\'etape\footnote{Une \'etape est un lieu o\`u l'on s'arr\^ete pour manger
et se reposer au cours d'un long voyage.} sur le chemin qui m\`ene au
r\'esultat ultime concernant $f_{small, *}$ pour les morphismes entiers et universellement
injectifs.

\begin{lemma}
\label{lemma-integral-universally-injective}
Soit $f : X \to Y$ un morphisme de sch\'emas. Supposons que $f$ soit
universellement injectif et entier (par exemple une immersion ferm\'ee).
Alors
\begin{enumerate}
\item
$f_{small, *} :
\Sh(X_\etale)
\to
\Sh(Y_\etale)$
refl\`ete les injections et les surjections,
\item
$f_{small, *} :
\Sh(X_\etale)
\to
\Sh(Y_\etale)$
commute aux sommes amalgam\'ees et aux co\'egalisateurs (et, plus g\'en\'eralement,
aux colimites connexes finies),
\item $f_{small, *}$ envoie les morphismes surjectifs sur des morphismes surjectifs (pour les faisceaux
d'ensembles et pour les faisceaux ab\'eliens),
\item le morphisme
$f_{small}^{-1}f_{small, *}\mathcal{F} \to \mathcal{F}$
est surjectif pour tout faisceau (d'ensembles ou de groupes ab\'eliens)
$\mathcal{F}$ sur $X_\etale$,
\item le foncteur $f_{small, *}$ est fid\`ele (sur les faisceaux d'ensembles et
sur les faisceaux ab\'eliens),
\item
$f_{small, *} :
\textit{Ab}(X_\etale)
\to
\textit{Ab}(Y_\etale)$
est exact, et
\item le foncteur
$Y_\etale \to X_\etale$, $V \mapsto X \times_Y V$ est
presque cocontinu.
\end{enumerate}
\end{lemma}

\begin{proof}
D'apr\`es les
Lemmes \ref{lemma-integral-A},
\ref{lemma-integral-B} et
\ref{lemma-integral-homeo-onto-image-C},
nous savons que le morphisme $f$ poss\`ede les propri\'et\'es (A), (B) et (C).
De plus, d'apr\`es le
Lemme \ref{lemma-property-C-closed-implies},
nous savons que le foncteur $Y_\etale \to X_\etale$ est
presque cocontinu. Nous avons maintenant :
\begin{enumerate}
\item la propri\'et\'e (C) implique (1) d'apr\`es le
Lemme \ref{lemma-property-C-implies},
\item la presque cocontinuit\'e implique (2) d'apr\`es
Sites, Lemme \ref{sites-lemma-morphism-of-sites-almost-cocontinuous},
\item la propri\'et\'e (B) implique (3) d'apr\`es le
Lemme \ref{lemma-property-B-implies}.
\end{enumerate}
Les propri\'et\'es (4), (5) et (6) d\'ecoulent formellement des trois premi\`eres ; voir
Sites, Lemme \ref{sites-lemma-exactness-properties}
et
Modules sur les sites, Lemme \ref{sites-modules-lemma-exactness}.
Nous avons vu la propri\'et\'e (7) ci-dessus.
\end{proof}




\section{Invariance topologique du petit site \'etale}
\label{section-topological-invariance}

\noindent
Dans le th\'eor\`eme suivant, nous montrons que le petit site \'etale est un invariant
topologique au sens suivant : si $f : X \to Y$ est un morphisme de sch\'emas
qui est un hom\'eomorphisme universel, alors $X_\etale \cong Y_\etale$
en tant que sites. Cela am\'eliore le r\'esultat de
Morphismes étales, Th\'eor\`eme \ref{etale-theorem-remarkable-equivalence}.
Nous d\'emontrons d'abord le r\'esultat pour les morphismes, puis nous \'enon\c{c}ons celui
pour les cat\'egories.

\begin{theorem}
\label{theorem-etale-topological}
Soient $X$ et $Y$ deux sch\'emas sur un sch\'ema de base $S$. Soit
$S' \to S$ un hom\'eomorphisme universel.
Notons $X'$ (resp.\ $Y'$) le changement de base \`a $S'$.
Si $X$ est \'etale sur $S$, alors l'application
$$
\Mor_S(Y, X) \longrightarrow \Mor_{S'}(Y', X')
$$
est bijective.
\end{theorem}

\begin{proof}
Apr\`es changement de base par $Y \to S$, nous pouvons supposer que $Y = S$.
Ainsi, nous pouvons et allons supposer que $X$ et $Y$ sont tous deux \'etales sur $S$.
Autrement dit, le th\'eor\`eme affirme que le foncteur de changement de base
est un foncteur pleinement fid\`ele de la cat\'egorie des sch\'emas \'etales
sur $S$ vers la cat\'egorie des sch\'emas \'etales sur $S'$.

\medskip\noindent
Consid\'erons le foncteur d'oubli
\begin{equation}
\label{equation-descent-etale-forget}
\begin{matrix}
\text{donn\'ees de descente }(X', \varphi')\text{ relativement \`a }S'/S \\
\text{ avec }X'\text{ \'etale sur }S'
\end{matrix}
\longrightarrow
\text{sch\'emas }X'\text{ \'etales sur }S'
\end{equation}
Nous affirmons que ce foncteur est une \'equivalence. D'autre part, le
foncteur
\begin{equation}
\label{equation-descent-etale}
\text{sch\'emas }X\text{ \'etales sur }S \longrightarrow
\begin{matrix}
\text{donn\'ees de descente }(X', \varphi')\text{ relativement \`a }S'/S \\
\text{ avec }X'\text{ \'etale sur }S'
\end{matrix}
\end{equation}
est pleinement fid\`ele d'apr\`es Morphismes étales, Lemme
\ref{etale-lemma-fully-faithful-cases}.
L'affirmation entra\^ine donc le th\'eor\`eme.

\medskip\noindent
D\'emonstration de l'affirmation.
Rappelons qu'un hom\'eomorphisme universel est la m\^eme chose qu'un morphisme
entier, universellement injectif et surjectif ; voir
Morphismes, Lemme \ref{morphisms-lemma-universal-homeomorphism}.
En particulier, la diagonale $\Delta : S' \to S' \times_S S'$ est un \'epaississement
d'apr\`es Morphismes, Lemme \ref{morphisms-lemma-universally-injective}.
Ainsi, d'apr\`es Morphismes étales, Th\'eor\`eme
\ref{etale-theorem-etale-topological},
nous voyons que, pour tout $X' \to S'$ \'etale, il existe un unique isomorphisme
$$
\varphi' : X' \times_S S' \to S' \times_S X'
$$
de sch\'emas \'etales sur $S' \times_S S'$ dont l'image inverse par
$\Delta$ est $\text{id} : X' \to X'$ sur $S'$.
Comme $S' \to S' \times_S S' \times_S S'$
est aussi un \'epaississement (c'est une immersion ferm\'ee bijective),
nous concluons que $(X', \varphi')$ est une donn\'ee de descente relativement \`a $S'/S$.
Le caract\`ere canonique de la construction de $\varphi'$ montre
qu'elle est compatible aux morphismes entre sch\'emas \'etales sur $S'$.
Autrement dit, nous obtenons un quasi-inverse
$X' \mapsto (X', \varphi')$ du foncteur
(\ref{equation-descent-etale-forget}). Cela d\'emontre l'affirmation et
ach\`eve la d\'emonstration du th\'eor\`eme.
\end{proof}

\begin{theorem}
\label{theorem-topological-invariance}
\begin{reference}
\cite[IV Theorem 18.1.2]{EGA}
\end{reference}
Soit $f : X \to Y$ un morphisme de sch\'emas.
Supposons $f$ entier, universellement injectif et surjectif
(c'est-\`a-dire que $f$ est un hom\'eomorphisme universel ; voir
Morphismes, Lemme \ref{morphisms-lemma-universal-homeomorphism}).
Le foncteur
$$
V \longmapsto V_X = X \times_Y V
$$
d\'efinit une \'equivalence de cat\'egories
$$
\{
\text{sch\'emas }V\text{ \'etales sur }Y
\}
\leftrightarrow
\{
\text{sch\'emas }U\text{ \'etales sur }X
\}
$$
\end{theorem}

\noindent
Nous donnons deux d\'emonstrations. La premi\`ere utilise l'effectivit\'e de la descente
des morphismes quasi-compacts, s\'epar\'es et \'etales relativement
aux morphismes entiers surjectifs. La seconde utilise les r\'esultats
sur les propri\'et\'es (A), (B) et (C) discut\'es plus haut dans le chapitre.

\begin{proof}[Premi\`ere d\'emonstration]
D'apr\`es le Th\'eor\`eme \ref{theorem-etale-topological},
nous voyons que le foncteur est pleinement fid\`ele.
Il reste \`a montrer que le foncteur est essentiellement surjectif.
Soit $U \to X$ un morphisme \'etale de sch\'emas.

\medskip\noindent
Supposons le r\'esultat vrai lorsque $U$ et $Y$ sont affines.
Dans ce cas, choisissons un recouvrement ouvert affine
$U = \bigcup U_i$ tel que chaque $U_i$ s'envoie
dans un ouvert affine de $Y$. Par hypoth\`ese (cas affine), nous pouvons
trouver des morphismes \'etales $V_i \to Y$ tels que $X \times_Y V_i \cong U_i$
comme sch\'emas sur $X$. Soit $V_{i, i'} \subset V_i$
le sous-sch\'ema ouvert dont l'espace topologique sous-jacent
correspond \`a $U_i \cap U_{i'}$. Puisque nous avons des isomorphismes
$$
X \times_Y V_{i, i'} \cong U_i \cap U_{i'} \cong X \times_Y V_{i', i}
$$
comme sch\'emas sur $X$, la pleine fid\'elit\'e montre que nous
obtenons des isomorphismes
$\theta_{i, i'} : V_{i, i'} \to V_{i', i}$ de sch\'emas sur $Y$.
Nous omettons de v\'erifier que ces isomorphismes satisfont \`a la
condition de cocycle de Schémas, section \ref{schemes-section-glueing-schemes}.
En appliquant Schémas, Lemme \ref{schemes-lemma-glue-schemes}
nous obtenons un sch\'ema $V \to Y$ par
recollement des sch\'emas $V_i$ suivant les identifications $\theta_{i, i'}$.
Il est clair que $V \to Y$ est \'etale et que $X \times_Y V \cong U$
par construction.

\medskip\noindent
Ainsi, il suffit de d\'emontrer le lemme dans le cas o\`u $U$ et $Y$ sont affines.
Rappelons que, dans la d\'emonstration du Th\'eor\`eme \ref{theorem-etale-topological},
nous avons montr\'e que $U$ est muni d'une unique donn\'ee de descente
$(U, \varphi)$ relativement \`a $X/Y$. D'apr\`es
Morphismes étales, Proposition \ref{etale-proposition-effective}
(qui s'applique car $U \to X$ est quasi-compact et s\'epar\'e
ainsi qu'\'etale, d'apr\`es notre r\'eduction au cas affine),
il existe un morphisme \'etale $V \to Y$ tel que
$X \times_Y V \cong U$, ce qui ach\`eve la d\'emonstration.
\end{proof}

\begin{proof}[Seconde d\'emonstration]
D'apr\`es le Th\'eor\`eme \ref{theorem-etale-topological},
nous voyons que le foncteur est pleinement fid\`ele.
Il reste \`a montrer que le foncteur est essentiellement surjectif.
Soit $U \to X$ un morphisme \'etale de sch\'emas.

\medskip\noindent
Supposons le r\'esultat vrai lorsque $U$ et $Y$ sont affines.
Dans ce cas, choisissons un recouvrement ouvert affine
$U = \bigcup U_i$ tel que chaque $U_i$ s'envoie
dans un ouvert affine de $Y$. Par hypoth\`ese (cas affine), nous pouvons
trouver des morphismes \'etales $V_i \to Y$ tels que $X \times_Y V_i \cong U_i$
comme sch\'emas sur $X$. Soit $V_{i, i'} \subset V_i$
le sous-sch\'ema ouvert dont l'espace topologique sous-jacent
correspond \`a $U_i \cap U_{i'}$. Puisque nous avons des isomorphismes
$$
X \times_Y V_{i, i'} \cong U_i \cap U_{i'} \cong X \times_Y V_{i', i}
$$
comme sch\'emas sur $X$, la pleine fid\'elit\'e montre que nous
obtenons des isomorphismes
$\theta_{i, i'} : V_{i, i'} \to V_{i', i}$ de sch\'emas sur $Y$.
Nous omettons de v\'erifier que ces isomorphismes satisfont \`a la
condition de cocycle de Schémas, section \ref{schemes-section-glueing-schemes}.
En appliquant Schémas, Lemme \ref{schemes-lemma-glue-schemes}
nous obtenons un sch\'ema $V \to Y$ par
recollement des sch\'emas $V_i$ suivant les identifications $\theta_{i, i'}$.
Il est clair que $V \to Y$ est \'etale et que $X \times_Y V \cong U$
par construction.

\medskip\noindent
Ainsi, il suffit de d\'emontrer que le foncteur
\begin{equation}
\label{equation-affine-etale}
\{
\text{sch\'emas affines }V\text{ \'etales sur }Y
\}
\leftrightarrow
\{
\text{sch\'emas affines }U\text{ \'etales sur }X
\}
\end{equation}
est essentiellement surjectif lorsque $X$ et $Y$ sont affines.

\medskip\noindent
Soit $U \to X$ un sch\'ema affine \'etale sur $X$.
Nous devons trouver $V \to Y$ \'etale (et affine) tel que $X \times_Y V$
soit isomorphe \`a $U$ sur $X$. Remarquons qu'un morphisme \'etale de sch\'emas affines
a des fibres universellement born\'ees, voir
Morphismes,
Lemmes \ref{morphisms-lemma-etale-locally-quasi-finite} et
\ref{morphisms-lemma-locally-quasi-finite-qc-source-universally-bounded}.
Nous pouvons donc raisonner par r\'ecurrence sur l'entier $n$ qui borne le degr\'e des fibres
de $U \to X$. Voir
Morphismes, Lemme \ref{morphisms-lemma-etale-universally-bounded}
pour une description de cet entier dans le cas d'un morphisme \'etale.
Si $n = 1$, alors $U \to X$ est une immersion ouverte (voir
Morphismes étales, Th\'eor\`eme \ref{etale-theorem-etale-radicial-open}),
et le r\'esultat est clair. Supposons $n > 1$.

\medskip\noindent
D'apr\`es
le Lemme \ref{lemma-integral-homeo-onto-image-C},
il existe un morphisme \'etale de sch\'emas $W \to Y$ et un
morphisme surjectif $W_X \to U$ sur $X$.
Comme $U$ est quasi-compact, nous pouvons remplacer $W$ par une somme disjointe d'un
nombre fini d'ouverts affines de $W$ ; nous pouvons donc supposer que $W$
est lui aussi affine. Voici un diagramme
$$
\xymatrix{
U \ar[d] & U \times_Y W \ar[l] \ar[d] & W_X \amalg R \ar@{=}[l]\\
X \ar[d] & W_X \ar[l] \ar[d] \\
Y & W \ar[l]
}
$$
La d\'ecomposition en somme disjointe provient du fait que, par construction, le
morphisme \'etale de sch\'emas affines $U \times_Y W \to W_X$ admet une section.
Nous voyons alors que le morphisme $R \to X \times_Y W$ est un morphisme \'etale
de sch\'emas affines dont les fibres ont un degr\'e universellement born\'e
par $n - 1$. Par l'hypoth\`ese de r\'ecurrence, il existe donc un sch\'ema
$V' \to W$ \'etale tel que $R \cong W_X \times_W V'$.
En posant $V'' = W \amalg V'$, nous obtenons un sch\'ema $V''$ \'etale sur $W$ dont le
changement de base \`a $W_X$ est isomorphe \`a $U \times_Y W$
sur $X \times_Y W$.

\medskip\noindent
\`A ce stade, nous pouvons utiliser la descente pour trouver $V$ sur $Y$ dont le
changement de base \`a $X$ est isomorphe \`a $U$ sur $X$. En effet, par la pleine
fid\'elit\'e du foncteur (\ref{equation-affine-etale})
correspondant \`a l'hom\'eomorphisme universel
$X \times_Y (W \times_Y W) \to (W \times_Y W)$
il existe un unique isomorphisme $\varphi : V'' \times_Y W \to W \times_Y V''$
dont le changement de base \`a $X \times_Y (W \times_Y W)$ est la donn\'ee de
descente canonique de $U \times_Y W$ sur $X \times_Y W$. En particulier,
$\varphi$ satisfait \`a la condition de cocycle. Ainsi, d'apr\`es
Descente, Lemme \ref{descent-lemma-affine},
nous voyons que $\varphi$ est effective (rappelons que tous les sch\'emas ci-dessus sont affines).
Nous obtenons donc $V \to Y$ et un isomorphisme $V'' \cong W \times_Y V$
tel que la donn\'ee de descente canonique sur $W \times_Y V/W/Y$ co\"incide
avec $\varphi$. Remarquons que $V \to Y$ est \'etale, d'apr\`es
Descente, Lemme \ref{descent-lemma-descending-property-etale}.
De plus, il existe un isomorphisme $V_X \cong U$ obtenu par
descente de l'isomorphisme
$$
V_X \times_X W_X =
X \times_Y V \times_Y W =
(X \times_Y W) \times_W (W \times_Y V) \cong
W_X  \times_W V'' \cong U \times_Y W
$$
dont nous disposons par construction. Nous omettons certains d\'etails.
\end{proof}

\begin{remark}
\label{remark-affine-inside-equivalence}
Dans la situation du
Th\'eor\`eme \ref{theorem-topological-invariance},
il est \`egalement vrai que $V \mapsto V_X$ induit une \'equivalence
entre les morphismes \'etales $V \to Y$ pour lesquels $V$ est affine et
les morphismes \'etales $U \to X$ pour lesquels $U$ est affine.
Cela r\'esulte par exemple de
Limits, Proposition \ref{limits-proposition-affine}.
\end{remark}

\begin{proposition}[Invariance topologique de la cohomologie \'etale]
\label{proposition-topological-invariance}
Soit $X_0 \to X$ un hom\'eomorphisme universel de sch\'emas
(par exemple l'immersion ferm\'ee d\'efinie par un faisceau d'id\'eaux nilpotent).
Alors
\begin{enumerate}
\item les sites \'etales $X_\etale$ et $(X_0)_\etale$ sont isomorphes,
\item les topos \'etales $\Sh(X_\etale)$ et $\Sh((X_0)_\etale)$
sont \'equivalents, et
\item $H^q_\etale(X, \mathcal{F}) = H^q_\etale(X_0, \mathcal{F}|_{X_0})$
pour tout $q$ et
pour tout faisceau ab\'elien $\mathcal{F}$ sur $X_\etale$.
\end{enumerate}
\end{proposition}

\begin{proof}
L'\'equivalence de cat\'egories $X_\etale \to (X_0)_\etale$ est
donn\'ee par le Th\'eor\`eme \ref{theorem-topological-invariance}. Nous omettons
de d\'emontrer que, sous cette \'equivalence, les recouvrements \'etales se correspondent.
Ainsi, (1) est d\'emontr\'e. Les assertions (2) et (3) d\'ecoulent formellement de (1).
\end{proof}






\section{Immersions ferm\'ees et image directe}
\label{section-closed-immersions}

\noindent
Avant d'\'enoncer et de d\'emontrer
Proposition \ref{proposition-closed-immersion-pushforward}
dans toute sa g\'en\'eralit\'e, nous l'\'enon\c{c}ons et la d\'emontrons bri\`evement pour les
immersions ferm\'ees. En effet, certains des arguments pr\'ec\'edents
sont beaucoup plus faciles \`a suivre dans le cas d'une immersion ferm\'ee et
nous les r\'ep\'etons donc ici sous leur forme simplifi\'ee.

\medskip\noindent
Dans la suite de cette section, $i : Z \to X$ est une immersion ferm\'ee.
Le foncteur
$$
\Sch/X \longrightarrow \Sch/Z, \quad
U \longmapsto U_Z = Z \times_X U
$$
sera not\'e $U \mapsto U_Z$, comme indiqu\'e. Comme la propri\'et\'e d'\^etre une immersion ferm\'ee
se conserve par tout changement de base, le sch\'ema $U_Z$ est un sous-sch\'ema ferm\'e
de $U$.

\begin{lemma}
\label{lemma-closed-immersion-almost-full}
Soit $i : Z \to X$ une immersion ferm\'ee de sch\'emas.
Soient $U, U'$ des sch\'emas \'etales sur $X$. Soit $h : U_Z \to U'_Z$
un morphisme sur $Z$. Alors il existe un diagramme
$$
\xymatrix{
U & W \ar[l]_a \ar[r]^b & U'
}
$$
dans $X_\etale$ tel que $a_Z : W_Z \to U_Z$
soit un isomorphisme et $h = b_Z \circ (a_Z)^{-1}$.
\end{lemma}

\begin{proof}
Consid\'erons le sch\'ema $M = U \times_X U'$. Le graphe $\Gamma_h \subset M_Z$
de $h$ est ouvert. Cela r\'esulte, par exemple, de ce que $\Gamma_h$ est l'image d'une
section du morphisme \'etale $\text{pr}_{1, Z} : M_Z \to U_Z$, voir
Morphismes étales, Proposition \ref{etale-proposition-properties-sections}.
Il existe donc un sous-sch\'ema ouvert $W \subset M$ dont l'intersection avec
la partie ferm\'ee $M_Z$ est $\Gamma_h$. Posons $a = \text{pr}_1|_W$
et $b = \text{pr}_2|_W$.
\end{proof}

\begin{lemma}
\label{lemma-closed-immersion-almost-essentially-surjective}
Soit $i : Z \to X$ une immersion ferm\'ee de sch\'emas.
Soit $V \to Z$ un morphisme \'etale de sch\'emas.
Il existe des morphismes \'etales $U_i \to X$ et des morphismes
$U_{i, Z} \to V$ tels que $\{U_{i, Z} \to V\}$
soit un recouvrement de Zariski de $V$.
\end{lemma}

\begin{proof}
Puisqu'il nous suffit de trouver un recouvrement de Zariski de $V$ form\'e de sch\'emas
de la forme $U_Z$, avec $U$ \'etale sur $X$, nous pouvons raisonner localement pour la topologie de Zariski sur $X$
et sur $V$. Nous pouvons donc supposer $X$ et $V$ affines. Dans le cas affine, cela r\'esulte de
Algèbre, Lemme \ref{algebra-lemma-lift-etale}.
\end{proof}

\noindent
Si $\overline{x} : \Spec(k) \to X$ est un point g\'eom\'etrique de $X$, alors
ou bien $\overline{x}$ se factorise (de fa\c{c}on unique) par le sous-sch\'ema ferm\'e $Z$, ou bien
$Z_{\overline{x}} = \emptyset$. Si $\overline{x}$ se factorise par $Z$,
nous disons que $\overline{x}$ est un point g\'eom\'etrique de $Z$ (ce qu'il est) et
nous employons la notation ``$\overline{x} \in Z$'' pour l'indiquer.

\begin{lemma}
\label{lemma-stalk-pushforward-closed-immersion}
Soit $i : Z \to X$ une immersion ferm\'ee de sch\'emas.
Soit $\mathcal{G}$ un faisceau d'ensembles sur $Z_\etale$.
Soit $\overline{x}$ un point g\'eom\'etrique de $X$.
Alors
$$
(i_{small, *}\mathcal{G})_{\overline{x}} =
\left\{
\begin{matrix}
* & \text{si} & \overline{x} \not \in Z \\
\mathcal{G}_{\overline{x}} & \text{si} & \overline{x} \in Z
\end{matrix}
\right.
$$
o\`u $*$ d\'esigne un ensemble r\'eduit \`a un \'el\'ement.
\end{lemma}

\begin{proof}
Posons $U = X \setminus Z$.
Notons que $i_{small, *}\mathcal{G}|_{U_\etale} = *$ est l'objet final
de la cat\'egorie des faisceaux \'etales sur $U$, c'est-\`a-dire le faisceau
qui associe un ensemble r\'eduit \`a un \'el\'ement \`a chaque sch\'ema \'etale sur $U$.
Cela explique la valeur de $(i_{small, *}\mathcal{G})_{\overline{x}}$
lorsque $\overline{x} \not \in Z$.

\medskip\noindent
Supposons maintenant que $\overline{x} \in Z$. Notons que
$$
(i_{small, *}\mathcal{G})_{\overline{x}}
=
\colim_{(U, \overline{u})} \mathcal{G}(U_Z)
$$
et, d'autre part,
$$
\mathcal{G}_{\overline{x}}
=
\colim_{(V, \overline{v})} \mathcal{G}(V).
$$
Soit $\mathcal{C}_1 = \{(U, \overline{u})\}^{opp}$ la cat\'egorie oppos\'ee \`a la
cat\'egorie des voisinages \'etales de $\overline{x}$ dans $X$, et soit
$\mathcal{C}_2 = \{(V, \overline{v})\}^{opp}$ la cat\'egorie oppos\'ee \`a la
cat\'egorie des voisinages \'etales de $\overline{x}$ dans $Z$. L'application canonique
$$
\mathcal{G}_{\overline{x}}
\longrightarrow
(i_{small, *}\mathcal{G})_{\overline{x}}
$$
correspond au foncteur $F : \mathcal{C}_1 \to \mathcal{C}_2$,
$F(U, \overline{u}) = (U_Z, \overline{x})$. Or les
lemmes \ref{lemma-closed-immersion-almost-essentially-surjective} et
\ref{lemma-closed-immersion-almost-full}
impliquent que $\mathcal{C}_1$ est cofinale dans $\mathcal{C}_2$, voir
Catégories, définition \ref{categories-definition-cofinal}.
Il en r\'esulte que la fl\`eche ci-dessus est un isomorphisme, voir
Catégories, lemme \ref{categories-lemma-cofinal}.
\end{proof}

\begin{proposition}
\label{proposition-closed-immersion-pushforward}
Soit $i : Z \to X$ une immersion ferm\'ee de sch\'emas.
\begin{enumerate}
\item Le foncteur
$$
i_{small, *} :
\Sh(Z_\etale)
\longrightarrow
\Sh(X_\etale)
$$
est pleinement fid\`ele et son image essentielle est form\'ee des faisceaux d'ensembles
$\mathcal{F}$ sur $X_\etale$ dont la restriction \`a $X \setminus Z$ est
isomorphe \`a $*$, et
\item le foncteur
$$
i_{small, *} :
\textit{Ab}(Z_\etale)
\longrightarrow
\textit{Ab}(X_\etale)
$$
est pleinement fid\`ele et son image essentielle est form\'ee des faisceaux ab\'eliens sur
$X_\etale$ dont le support est contenu dans $Z$.
\end{enumerate}
Dans les deux cas, $i_{small}^{-1}$ est un inverse \`a gauche du foncteur
$i_{small, *}$.
\end{proposition}

\begin{proof}
Traitons le cas des faisceaux d'ensembles.
Pour tout faisceau $\mathcal{G}$ sur $Z$, le morphisme
$i_{small}^{-1}i_{small, *}\mathcal{G} \to \mathcal{G}$
est un isomorphisme d'apr\`es
le lemme \ref{lemma-stalk-pushforward-closed-immersion}
(et
le th\'eor\`eme \ref{theorem-exactness-stalks}).
Il en r\'esulte formellement que $i_{small, *}$ est pleinement fid\`ele, voir
Sites, lemme \ref{sites-lemma-exactness-properties}.
Il est clair que $i_{small, *}\mathcal{G}|_{U_\etale} \cong *$
o\`u $U = X \setminus Z$. R\'eciproquement, supposons que $\mathcal{F}$
soit un faisceau d'ensembles sur $X$ tel que $\mathcal{F}|_{U_\etale} \cong *$.
Consid\'erons le morphisme d'adjonction
$$
\mathcal{F} \longrightarrow i_{small, *}i_{small}^{-1}\mathcal{F}
$$
En combinant
les lemmes \ref{lemma-stalk-pushforward-closed-immersion} et
\ref{lemma-stalk-pullback},
nous voyons que c'est un isomorphisme. Cela ach\`eve la d\'emonstration de (1).
La d\'emonstration de (2) est identique.
\end{proof}





\section{Morphismes entiers universellement injectifs}
\label{section-integral-universally-injective}

\noindent
Voici la version g\'en\'erale de la
proposition \ref{proposition-closed-immersion-pushforward}.

\begin{proposition}
\label{proposition-integral-universally-injective-pushforward}
Soit $f : X \to Y$ un morphisme entier de sch\'emas
et universellement injectif.
\begin{enumerate}
\item Le foncteur
$$
f_{small, *} :
\Sh(X_\etale)
\longrightarrow
\Sh(Y_\etale)
$$
est pleinement fid\`ele et son image essentielle est form\'ee des faisceaux d'ensembles
$\mathcal{F}$ sur $Y_\etale$ dont la restriction \`a $Y \setminus f(X)$ est
isomorphe \`a $*$, et
\item le foncteur
$$
f_{small, *} :
\textit{Ab}(X_\etale)
\longrightarrow
\textit{Ab}(Y_\etale)
$$
est pleinement fid\`ele et son image essentielle est form\'ee des faisceaux ab\'eliens sur
$Y_\etale$ dont le support est contenu dans $f(X)$.
\end{enumerate}
Dans les deux cas, $f_{small}^{-1}$ est un inverse \`a gauche du foncteur
$f_{small, *}$.
\end{proposition}

\begin{proof}
Nous pouvons factoriser $f$ sous la forme
$$
\xymatrix{
X \ar[r]^h & Z \ar[r]^i & Y
}
$$
o\`u $h$ est entier, universellement injectif et surjectif
et $i : Z \to Y$ est une immersion ferm\'ee.
Appliquons la
proposition \ref{proposition-closed-immersion-pushforward}
\`a $i$ et le
th\'eor\`eme \ref{theorem-topological-invariance}
\`a $h$.
\end{proof}









\section{Grands sites et image directe}
\label{section-big}

\noindent
Dans cette section, nous d\'emontrons quelques r\'esultats techniques sur $f_{big, *}$ pour
certains types de morphismes de sch\'emas.

\begin{lemma}
\label{lemma-monomorphism-big-push-pull}
Soit $\tau \in \{Zariski, \etale, smooth, syntomic, fppf\}$.
Soit $f : X \to Y$ un monomorphisme de sch\'emas.
Alors le morphisme canonique
$f_{big}^{-1}f_{big, *}\mathcal{F} \to \mathcal{F}$
est un isomorphisme pour tout faisceau $\mathcal{F}$ sur
$(\Sch/X)_\tau$.
\end{lemma}

\begin{proof}
Dans ce cas, le foncteur $(\Sch/X)_\tau \to (\Sch/Y)_\tau$
est continu, cocontinu et pleinement fid\`ele. Le r\'esultat d\'ecoule donc de
Sites, lemme \ref{sites-lemma-back-and-forth}.
\end{proof}

\begin{remark}
\label{remark-push-pull-shriek}
Dans la situation du
lemme \ref{lemma-monomorphism-big-push-pull},
le morphisme canonique
$\mathcal{F} \to f_{big}^{-1}f_{big!}\mathcal{F}$
est un isomorphisme pour tout faisceau d'ensembles $\mathcal{F}$ sur
$(\Sch/X)_\tau$. La d\'emonstration est la m\^eme. Cela
vaut aussi pour les faisceaux ab\'eliens. Cependant, notons
que le foncteur $f_{big!}$ pour les faisceaux ab\'eliens est d\'efini dans
Modules sur les sites, section \ref{sites-modules-section-exactness-lower-shriek}
et est en g\'en\'eral diff\'erent de $f_{big!}$ sur les faisceaux d'ensembles.
Le r\'esultat pour les faisceaux ab\'eliens d\'ecoule de
Modules sur les sites, lemme \ref{sites-modules-lemma-back-and-forth}.
\end{remark}

\begin{lemma}
\label{lemma-closed-immersion-cover-from-below}
Soit $f : X \to Y$ une immersion ferm\'ee de sch\'emas.
Soit $U \to X$ un morphisme syntomique (resp.\ lisse, resp.\ \'etale).
Il existe alors des morphismes syntomiques (resp.\ lisses, resp.\ \'etales)
$V_i \to Y$ et des morphismes $V_i \times_Y X \to U$ tels que
$\{V_i \times_Y X \to U\}$ forme un recouvrement de Zariski de $U$.
\end{lemma}

\begin{proof}
D\'emontrons le lemme dans le cas o\`u $\tau = syntomic$.
La question est locale sur $U$. Nous pouvons donc supposer que $U$ est
un sch\'ema affine dont l'image est contenue dans un ouvert affine de $Y$.
Nous sommes donc ramen\'es \`a d\'emontrer le cas suivant :
$Y = \Spec(A)$, $X = \Spec(A/I)$, et
$U = \Spec(\overline{B})$, o\`u
$A/I \to \overline{B}$ est un morphisme d'anneaux syntomique.
D'apr\`es Algèbre, lemme \ref{algebra-lemma-lift-syntomic}
nous pouvons trouver des \'el\'ements $\overline{g}_i \in \overline{B}$
tels que
$\overline{B}_{\overline{g}_i} = A_i/IA_i$ pour certains morphismes d'anneaux syntomiques
$A \to A_i$.
Cela d\'emontre le lemme dans le cas syntomique.
La d\'emonstration du cas lisse est identique, \`a ceci pr\`es qu'elle utilise
Algèbre, lemme \ref{algebra-lemma-lift-smooth}.
Dans le cas \'etale, on utilise
Algèbre, lemme \ref{algebra-lemma-lift-etale}.
\end{proof}

\begin{lemma}
\label{lemma-prepare-closed-immersion-almost-cocontinuous}
Soit $f : X \to Y$ une immersion ferm\'ee de sch\'emas.
Soit $\{U_i \to X\}$ un recouvrement syntomique (resp.\ lisse, resp.\ \'etale).
Il existe un recouvrement syntomique (resp.\ lisse, resp.\ \'etale) $\{V_j \to Y\}$
tel que, pour chaque $j$, soit $V_j \times_Y X = \emptyset$, soit le
morphisme $V_j \times_Y X \to X$ se factorise par $U_i$ pour un certain $i$.
\end{lemma}

\begin{proof}
Pour chaque $i$, nous pouvons choisir des morphismes syntomiques (resp.\ lisses, resp.\ \'etales)
$g_{ij} : V_{ij} \to Y$ et des morphismes $V_{ij} \times_Y X \to U_i$ au-dessus de $X$,
de sorte que $\{V_{ij} \times_Y X \to U_i\}$ forme un recouvrement de Zariski, voir
le lemme \ref{lemma-closed-immersion-cover-from-below}.
Cela implique en particulier que
$\bigcup_{ij} g_{ij}(V_{ij})$ contient le sous-ensemble ferm\'e $f(X)$.
Par cons\'equent, la famille de morphismes syntomiques (resp.\ lisses, resp.\ \'etales) $g_{ij}$
ainsi que l'immersion ouverte $Y \setminus f(X) \to Y$ forment le recouvrement
syntomique (resp.\ lisse, resp.\ \'etale) voulu de $Y$.
\end{proof}

\begin{lemma}
\label{lemma-closed-immersion-almost-cocontinuous}
Soit $f : X \to Y$ une immersion ferm\'ee de sch\'emas.
Soit $\tau \in \{syntomic, smooth, \etale\}$.
Le foncteur $V \mapsto X \times_Y V$ d\'efinit un foncteur presque
cocontinu (voir
Sites, D\'efinition \ref{sites-definition-almost-cocontinuous})
$(\Sch/Y)_\tau \to (\Sch/X)_\tau$ entre les
grands $\tau$-sites.
\end{lemma}

\begin{proof}
Il faut montrer ce qui suit : \'etant donn\'e un morphisme $V \to Y$
et un recouvrement syntomique quelconque (resp.\ lisse, resp.\ \'etale)
$\{U_i \to X \times_Y V\}$, il existe un
recouvrement lisse (resp.\ lisse, resp.\ \'etale) $\{V_j \to V\}$
tel que, pour tout $j$, soit $X \times_Y V_j$ est vide, soit
$X \times_Y V_j \to Z \times_Y V$ se factorise par l'un des
$U_i$. Cela r\'esulte de l'application du
Lemme \ref{lemma-prepare-closed-immersion-almost-cocontinuous}
ci-dessus \`a l'immersion ferm\'ee $X \times_Y V \to V$.
\end{proof}

\begin{lemma}
\label{lemma-closed-immersion-pushforward-exact}
Soit $f : X \to Y$ une immersion ferm\'ee de sch\'emas.
Soit $\tau \in \{syntomic, smooth, \etale\}$.
\begin{enumerate}
\item Le foncteur d'image directe
$f_{big, *} :
\Sh((\Sch/X)_\tau)
\to
\Sh((\Sch/Y)_\tau)$
commute aux co\'egalisateurs et aux sommes amalgam\'ees.
\item Le foncteur d'image directe
$f_{big, *} :
\textit{Ab}((\Sch/X)_\tau)
\to
\textit{Ab}((\Sch/Y)_\tau)$
est exact.
\end{enumerate}
\end{lemma}

\begin{proof}
Cela r\'esulte de
Sites, Lemme \ref{sites-lemma-morphism-of-sites-almost-cocontinuous},
Modules sur les sites,
Lemme \ref{sites-modules-lemma-morphism-ringed-sites-almost-cocontinuous},
et du
Lemme \ref{lemma-closed-immersion-almost-cocontinuous}
ci-dessus.
\end{proof}

\begin{remark}
\label{remark-fppf-closed-immersion-not-closed}
Dans le Lemme \ref{lemma-closed-immersion-pushforward-exact}, le cas $\tau = fppf$
n'est pas trait\'e. La raison en est qu'\'etant donn\'es un anneau $A$, un id\'eal $I$ et un
homomorphisme d'anneaux $A/I \to \overline{B}$ fid\`element plat et de pr\'esentation finie, il
n'y a aucune raison de penser que l'on puisse trouver {\it un quelconque} homomorphisme d'anneaux
$A \to B$, plat et de pr\'esentation finie, avec $B/IB \not = 0$, tel que $A/I \to B/IB$ se factorise par
$\overline{B}$. Ainsi, la preuve du
Lemme \ref{lemma-closed-immersion-almost-cocontinuous}
ne fonctionne pas pour la topologie fppf.
En fait, il est probablement faux que
$f_{big, *} : \textit{Ab}((\Sch/X)_{fppf})
\to \textit{Ab}((\Sch/Y)_{fppf})$
soit exact lorsque $f$ est une immersion ferm\'ee.
Si vous connaissez un exemple, veuillez envoyer un courriel \`a
\href{mailto:stacks.project@gmail.com}{stacks.project@gmail.com}.
\end{remark}












\section{Exactitude du grand foncteur image directe \`a support propre}
\sectionmark{Image directe à support propre}
\label{section-exactness-lower-shriek}

\noindent
Il s'agit simplement du r\'esultat technique suivant. Notons que le foncteur $f_{big!}$
n'a absolument rien \`a voir avec la cohomologie \`a support compact en
g\'en\'eral.

\begin{lemma}
\label{lemma-exactness-lower-shriek}
Soit $\tau \in \{Zariski, \etale, smooth, syntomic, fppf\}$.
Soit $f : X \to Y$ un morphisme de sch\'emas. Soit
$$
f_{big} :
\Sh((\Sch/X)_\tau)
\longrightarrow
\Sh((\Sch/Y)_\tau)
$$
le morphisme de topos correspondant comme dans
Topologies, Lemme
\ref{topologies-lemma-morphism-big},
\ref{topologies-lemma-morphism-big-etale},
\ref{topologies-lemma-morphism-big-smooth},
\ref{topologies-lemma-morphism-big-syntomic}, ou
\ref{topologies-lemma-morphism-big-fppf}.
\begin{enumerate}
\item Le foncteur
$f_{big}^{-1} : \textit{Ab}((\Sch/Y)_\tau) \to \textit{Ab}((\Sch/X)_\tau)$
admet un adjoint \`a gauche
$$
f_{big!} : \textit{Ab}((\Sch/X)_\tau) \to \textit{Ab}((\Sch/Y)_\tau)
$$
qui est exact.
\item Le foncteur
$f_{big}^* :
\textit{Mod}((\Sch/Y)_\tau, \mathcal{O})
\to
\textit{Mod}((\Sch/X)_\tau, \mathcal{O})$
admet un adjoint \`a gauche
$$
f_{big!} :
\textit{Mod}((\Sch/X)_\tau, \mathcal{O})
\to
\textit{Mod}((\Sch/Y)_\tau, \mathcal{O})
$$
qui est exact.
\end{enumerate}
De plus, les deux foncteurs $f_{big!}$ co\"incident sur les faisceaux sous-jacents
de groupes ab\'eliens.
\end{lemma}

\begin{proof}
Rappelons que $f_{big}$ est le morphisme de topos associ\'e au
foncteur continu et cocontinu
$u : (\Sch/X)_\tau \to (\Sch/Y)_\tau$, $U/X \mapsto U/Y$.
De plus, on a $f_{big}^{-1}\mathcal{O} = \mathcal{O}$.
Ainsi, l'existence de $f_{big!}$ r\'esulte de
Modules sur les sites, Lemme \ref{sites-modules-lemma-g-shriek-adjoint},
et, respectivement, de
Modules sur les sites, Lemme \ref{sites-modules-lemma-lower-shriek-modules}.
Notons que, si $U$ est un objet de $(\Sch/X)_\tau$, alors le foncteur
$u$ induit une \'equivalence de cat\'egories
$$
u' :
(\Sch/X)_\tau/U
\longrightarrow
(\Sch/Y)_\tau/U
$$
car les deux membres de la fl\`eche sont \'egaux \`a $(\Sch/U)_\tau$.
Ainsi, la co\"incidence de $f_{big!}$ sur les faisceaux ab\'eliens sous-jacents
r\'esulte de la discussion dans
Modules sur les sites, Remarque \ref{sites-modules-remark-when-shriek-equal}.
L'exactitude de $f_{big!}$ r\'esulte de
Modules sur les sites, Lemme \ref{sites-modules-lemma-exactness-lower-shriek}
puisque le foncteur $u$ ci-dessus commute aux produits fibr\'es et aux \'egalisateurs.
\end{proof}

\noindent
D\'emontrons ensuite un lemme technique qui sera utile plus loin pour comparer
les faisceaux de modules sur les diff\'erents sites associ\'es aux champs alg\'ebriques.

\begin{lemma}
\label{lemma-compare-structure-sheaves}
Soit $X$ un sch\'ema. Soit
$\tau \in \{Zariski, \etale, smooth, syntomic, fppf\}$.
Soient $\mathcal{C}_1 \subset \mathcal{C}_2 \subset (\Sch/X)_\tau$ des sous-cat\'egories pleines
ayant les propri\'et\'es suivantes :
\begin{enumerate}
\item Pour tout objet $U/X$ de $\mathcal{C}_t$,
\begin{enumerate}
\item si $\{U_i \to U\}$ est un recouvrement de $(\Sch/X)_\tau$, alors
$U_i/X$ est un objet de $\mathcal{C}_t$,
\item $U \times \mathbf{A}^1/X$ est un objet de $\mathcal{C}_t$.
\end{enumerate}
\item $X/X$ est un objet de $\mathcal{C}_t$.
\end{enumerate}
Munissons $\mathcal{C}_t$ de la structure de site dont les recouvrements sont
exactement les recouvrements $\{U_i \to U\}$ de $(\Sch/X)_\tau$ tels que
$U \in \Ob(\mathcal{C}_t)$. Alors
\begin{enumerate}
\item[(a)] Le foncteur $\mathcal{C}_1 \to \mathcal{C}_2$
est pleinement fid\`ele, continu et cocontinu.
\end{enumerate}
Notons $g : \Sh(\mathcal{C}_1) \to \Sh(\mathcal{C}_2)$ le
morphisme de topos correspondant. Notons $\mathcal{O}_t$ la restriction de $\mathcal{O}$
\`a $\mathcal{C}_t$. Notons $g_!$ le foncteur de
Modules sur les sites, D\'efinition \ref{sites-modules-definition-g-shriek}.
\begin{enumerate}
\item[(b)] Le morphisme canonique $g_!\mathcal{O}_1 \to \mathcal{O}_2$
est un isomorphisme.
\end{enumerate}
\end{lemma}

\begin{proof}
L'assertion (a) est imm\'ediate d'apr\`es les d\'efinitions.
Dans cette preuve, tous les sch\'emas sont des sch\'emas sur $X$ et tous les morphismes de
sch\'emas sont des morphismes de sch\'emas sur $X$. Notons que $g^{-1}$ est
donn\'e par restriction, de sorte que, pour tout objet $U$ de $\mathcal{C}_1$
on a $\mathcal{O}_1(U) = \mathcal{O}_2(U) = \mathcal{O}(U)$.
Rappelons que $g_!\mathcal{O}_1$ est le faisceau associ\'e au pr\'efaisceau
$g_{p!}\mathcal{O}_1$ qui associe \`a $V$ dans $\mathcal{C}_2$ le groupe
$$
\colim_{V \to U} \mathcal{O}(U)
$$
o\`u $U$ parcourt les objets de $\mathcal{C}_1$ et o\`u la colimite est
prise dans la cat\'egorie des groupes ab\'eliens. Nous utiliserons souvent ci-dessous
le fait que, si
$$
V \to U \to U'
$$
sont des morphismes avec $U, U' \in \Ob(\mathcal{C}_1)$
et si $f' \in \mathcal{O}(U')$ se restreint en $f \in \mathcal{O}(U)$,
alors $(V \to U, f)$ et $(V \to U', f')$ d\'efinissent le m\^eme \'el\'ement de la
colimite. De plus, $g_!\mathcal{O}_1 \to \mathcal{O}_2$ envoie l'\'el\'ement
$(V \to U, f)$ simplement sur la restriction de $f$ \`a $V$.

\medskip\noindent
Surjectivit\'e. Soit $V$ un sch\'ema et soit $h \in \mathcal{O}(V)$.
On obtient alors un morphisme $V \to X \times \mathbf{A}^1$ induit par $h$
et le morphisme structural $V \to X$. En \'ecrivant
$\mathbf{A}^1 = \Spec(\mathbf{Z}[x])$, on voit que l'\'el\'ement
$x \in \mathcal{O}(X \times \mathbf{A}^1)$ a pour image inverse
$h$. Puisque $X \times \mathbf{A}^1$ est un objet de $\mathcal{C}_1$
d'apr\`es les hypoth\`eses (1)(b) et (2), on obtient la surjectivit\'e voulue.

\medskip\noindent
Injectivit\'e. Soit $V$ un sch\'ema. Soit
$s = \sum_{i = 1, \ldots, n} (V \to U_i, f_i)$ un \'el\'ement de la colimite
affich\'ee ci-dessus. Pour tout $i$, on peut utiliser le morphisme
$f_i : U_i \to X \times \mathbf{A}^1$
pour voir que $(V \to U_i, f_i)$ d\'efinit le m\^eme \'el\'ement de la colimite que
$(f_i : V \to X \times \mathbf{A}^1, x)$. Nous pouvons alors consid\'erer
$$
f_1 \times \ldots \times f_n : V \to X \times \mathbf{A}^n
$$
et l'on voit que $s$ est \'equivalent dans la colimite \`a
$$
\sum\nolimits_{i = 1, \ldots, n}
(f_1 \times \ldots \times f_n : V \to X \times \mathbf{A}^n, x_i) =
(f_1 \times \ldots \times f_n : V \to X \times \mathbf{A}^n,
x_1 + \ldots + x_n)
$$
Maintenant, si $x_1 + \ldots + x_n$ se restreint \`a z\'ero sur $V$, alors on voit
que $f_1 \times \ldots \times f_n$ se factorise par
$X \times \mathbf{A}^{n - 1} = V(x_1 + \ldots + x_n)$. Ainsi, nous voyons
que $s$ est \'equivalent \`a z\'ero dans la limite inductive.
\end{proof}











%9.29.09
\section{Cohomologie \'etale}
\label{section-etale-cohomology}

\noindent
Dans les sections suivantes, nous d\'emontrons quelques r\'esultats \`el\'ementaires sur la cohomologie \'etale.
Voici une propri\'et\'e connue de la cohomologie des espaces
topologiques et qui vaut \`egalement pour la cohomologie \'etale.

\begin{lemma}[Mayer-Vietoris en cohomologie \'etale]
\label{lemma-mayer-vietoris}
Soit $X$ un sch\'ema. Supposons que $X = U \cup V$ soit la
r\'eunion de deux ouverts. Pour tout faisceau ab\'elien $\mathcal{F}$ sur $X_\etale$
il existe une suite exacte longue de cohomologie
$$
\begin{matrix}
0 \to
H^0_\etale(X, \mathcal{F}) \to
H^0_\etale(U, \mathcal{F}) \oplus H^0_\etale(V, \mathcal{F}) \to
H^0_\etale(U \cap V, \mathcal{F}) \phantom{\to \ldots} \\
\phantom{0} \to H^1_\etale(X, \mathcal{F}) \to
H^1_\etale(U, \mathcal{F}) \oplus H^1_\etale(V, \mathcal{F}) \to
H^1_\etale(U \cap V, \mathcal{F}) \to \ldots
\end{matrix}
$$
Cette suite exacte longue est fonctorielle en $\mathcal{F}$.
\end{lemma}

\begin{proof}
Observons que si $\mathcal{I}$ est un faisceau ab\'elien injectif, alors
$$
0 \to \mathcal{I}(X) \to \mathcal{I}(U) \oplus \mathcal{I}(V) \to
\mathcal{I}(U \cap V) \to 0
$$
est exacte. L'exactitude aux premier et second termes vient de ce que $\mathcal{I}$
est un faisceau. L'exactitude \`a droite vient de ce que
$\mathcal{I}(U) \to \mathcal{I}(U \cap V)$ est surjective d'apr\`es
Cohomologie sur les sites, Lemme
\ref{sites-cohomology-lemma-restriction-along-monomorphism-surjective}.
Une autre mani\`ere de le d\'emontrer serait de montrer que le conoyau
de l'application $\mathcal{I}(U) \oplus \mathcal{I}(V) \to
\mathcal{I}(U \cap V)$ est le premier groupe de cohomologie de {\v C}ech
de $\mathcal{I}$ relativement au recouvrement
$X = U \cup V$, lequel s'annule d'apr\`es les
lemmes \ref{lemma-hom-injective} et \ref{lemma-forget-injectives}.
Ainsi, si $\mathcal{F} \to \mathcal{I}^\bullet$ est une
r\'esolution injective, alors
$$
0 \to \mathcal{I}^\bullet(X) \to
\mathcal{I}^\bullet(U) \oplus \mathcal{I}^\bullet(V) \to
\mathcal{I}^\bullet(U \cap V) \to 0
$$
est une suite exacte courte de complexes et la suite exacte longue
de cohomologie associ\'ee est celle de l'\'enonc\'e du lemme.
\end{proof}

\begin{lemma}[Mayer-Vietoris relatif]
\label{lemma-relative-mayer-vietoris}
Soit $f : X \to Y$ un morphisme de sch\'emas. Supposons que $X = U \cup V$
soit une r\'eunion de deux sous-sch\'emas ouverts. Notons
$a = f|_U : U \to Y$, $b = f|_V : V \to Y$, et
$c = f|_{U \cap V} : U \cap V \to Y$.
Pour tout faisceau ab\'elien $\mathcal{F}$ sur $X_\etale$
il existe une suite exacte longue
$$
0 \to
f_*\mathcal{F} \to
a_*(\mathcal{F}|_U) \oplus b_*(\mathcal{F}|_V) \to
c_*(\mathcal{F}|_{U \cap V}) \to
R^1f_*\mathcal{F} \to \ldots
$$
sur $Y_\etale$.
Cette suite exacte longue est fonctorielle en $\mathcal{F}$.
\end{lemma}

\begin{proof}
Soit $\mathcal{F} \to \mathcal{I}^\bullet$ une r\'esolution injective
de $\mathcal{F}$ sur $X_\etale$. Nous affirmons qu'il existe
une suite exacte courte de complexes
$$
0 \to
f_*\mathcal{I}^\bullet \to
a_*\mathcal{I}^\bullet|_U \oplus b_*\mathcal{I}^\bullet|_V \to
c_*\mathcal{I}^\bullet|_{U \cap V} \to
0.
$$
En effet, pour tout $W$ de $Y_\etale$ et tout $n \geq 0$, la
suite correspondante des groupes de sections sur $W$
$$
0 \to
\mathcal{I}^n(W \times_Y X) \to
\mathcal{I}^n(W \times_Y U)
\oplus \mathcal{I}^n(W \times_Y V) \to
\mathcal{I}^n(W \times_Y (U \cap V)) \to
0
$$
est une suite exacte courte, comme on l'a montr\'e dans la d\'emonstration du lemme \ref{lemma-mayer-vietoris}.
Le lemme r\'esulte alors du passage aux faisceaux de cohomologie et du fait que
$\mathcal{I}^\bullet|_U$ est une r\'esolution injective de $\mathcal{F}|_U$
et de m\^eme pour $\mathcal{I}^\bullet|_V$, $\mathcal{I}^\bullet|_{U \cap V}$.
\end{proof}







\section{Limites inductives}
\label{section-colimit}

\noindent
Rappelons que si $(\mathcal{F}_i, \varphi_{ii'})$ est un diagramme de faisceaux sur
un site $\mathcal{C}$, sa limite inductive (dans la cat\'egorie des faisceaux) est le
faisceau associ\'e au pr\'efaisceau $U \mapsto \colim_i \mathcal{F}_i(U)$. Voir
Sites, Lemme \ref{sites-lemma-colimit-sheaves}.
Si le syst\`eme est filtrant, si $U$ est un objet quasi-compact de
$\mathcal{C}$ qui poss\`ede un syst\`eme cofinal de recouvrements par des objets
quasi-compacts, alors $\mathcal{F}(U) = \colim \mathcal{F}_i(U)$, voir
Sites, Lemme \ref{sites-lemma-directed-colimits-sections}.
Voir Cohomologie sur les sites, Lemme
\ref{sites-cohomology-lemma-colim-works-over-collection}
pour un r\'esultat concernant les groupes de cohomologie sup\'erieure des limites inductives
de faisceaux ab\'eliens.

\medskip\noindent
Dans Cohomologie sur les sites, Lemme \ref{sites-cohomology-lemma-colimit},
nous g\'en\'eralisons ce r\'esultat \`a un syst\`eme de faisceaux
sur un syst\`eme projectif de sites. Voici la notion correspondante
dans le cas d'un syst\`eme de faisceaux \'etales d\'efinis
sur un syst\`eme projectif de sch\'emas.

\begin{definition}
\label{definition-inverse-system-sheaves}
Soit $I$ un ensemble pr\'eordonn\'e. Soit $(X_i, f_{i'i})$ un syst\`eme
projectif de sch\'emas index\'e par $I$.
Un {\it syst\`eme $(\mathcal{F}_i, \varphi_{i'i})$ de faisceaux
sur $(X_i, f_{i'i})$} est la donn\'ee de
\begin{enumerate}
\item un faisceau $\mathcal{F}_i$ sur $(X_i)_\etale$ pour tout $i \in I$,
\item pour $i' \geq i$, un morphisme
$\varphi_{i'i} : f_{i'i}^{-1}\mathcal{F}_i \to \mathcal{F}_{i'}$
de faisceaux sur $(X_{i'})_\etale$
\end{enumerate}
de telle sorte que $\varphi_{i''i} = \varphi_{i''i'} \circ f_{i'' i'}^{-1}\varphi_{i'i}$
lorsque $i'' \geq i' \geq i$.
\end{definition}

\noindent
Dans la situation de la d\'efinition \ref{definition-inverse-system-sheaves},
supposons que $I$ soit filtrant et que les morphismes de transition $f_{i'i}$ soient affines.
Soit $X = \lim X_i$ la limite dans la cat\'egorie des sch\'emas, voir
Limites, section \ref{limits-section-limits}.
Notons $f_i : X \to X_i$ les morphismes de projection et consid\'erons les morphismes
$$
f_i^{-1}\mathcal{F}_i = f_{i'}^{-1}f_{i'i}^{-1}\mathcal{F}_i
\xrightarrow{f_{i'}^{-1}\varphi_{i'i}}
f_{i'}^{-1}\mathcal{F}_{i'}
$$
Cela fait des $f_i^{-1}\mathcal{F}_i$ un syst\`eme de faisceaux sur $X_\etale$
index\'e par $I$ (c'est un bon exercice de le v\'erifier).
Nous voulons souvent savoir si l'on a un isomorphisme
$$
H^q_\etale(X, \colim f_i^{-1}\mathcal{F}_i) =
\colim H^q_\etale(X_i, \mathcal{F}_i)
$$
Il s'av\'erera que cela est vrai si $X_i$ est quasi-compact et quasi-s\'epar\'e
pour tout $i$, voir le th\'eor\`eme \ref{theorem-colimit}.

\begin{lemma}
\label{lemma-colimit-affine-sites}
Soit $I$ un ensemble pr\'eordonn\'e filtrant. Soit $(X_i, f_{i'i})$ un syst\`eme
projectif de sch\'emas index\'e par $I$ dont les morphismes de transition sont affines.
Soit $X = \lim_{i \in I} X_i$. Avec les
notations de Topologies, Lemme \ref{topologies-lemma-alternative}, on a
$$
X_{affine, \etale} = \colim (X_i)_{affine, \etale}
$$
comme sites au sens de
Sites, Lemme \ref{sites-lemma-colimit-sites}.
\end{lemma}

\begin{proof}
D\'emontrons d'abord ceci lorsque $X$ et les $X_i$ sont quasi-compacts et
quasi-s\'epar\'es pour tout $i$ (c'est le cas dans toutes les situations qui
nous int\'eressent). Dans ce cas, tout objet de
$X_{affine, \etale}$, resp.\ de $(X_i)_{affine, \etale}$, est
de pr\'esentation finie sur $X$. De plus, la
cat\'egorie des sch\'emas de pr\'esentation finie sur $X$ est la
limite inductive des cat\'egories de sch\'emas de pr\'esentation finie
sur les $X_i$, voir
Limites, Lemme \ref{limits-lemma-descend-finite-presentation}.
Il en va de m\^eme des sous-cat\'egories d'objets affines et \'etales
sur $X$, d'apr\`es les lemmes de Limites
\ref{limits-lemma-limit-affine} et \ref{limits-lemma-descend-etale}.
Enfin, si $\{U^j \to U\}$ est un recouvrement de $X_{affine, \etale}$
et si $U_i^j \to U_i$ est un morphisme de sch\'emas affines \'etales sur
$X_i$ dont le changement de base \`a $X$ est $U^j \to U$, alors nous voyons que
le changement de base de $\{U^j_i \to U_i\}$ \`a un certain $X_{i'}$ est
un recouvrement pour $i'$ assez grand, voir
Limites, Lemme \ref{limits-lemma-descend-surjective}.

\medskip\noindent
Dans le cas g\'en\'eral, soit $U$ un objet de $X_{affine, \etale}$.
Alors $U \to X$ est \'etale et s\'epar\'e (puisque $U$ est s\'epar\'e),
mais n'est en g\'en\'eral pas quasi-compact. Toutefois, $U \to X$ est localement
de pr\'esentation finie et donc, d'apr\`es le
lemme \ref{limits-lemma-descend-finite-presentation-variant} de Limites,
il existe un $i$, un sch\'ema $U_i$ quasi-compact et quasi-s\'epar\'e, et
un morphisme $U_i \to X_i$ qui est localement de pr\'esentation finie
dont le changement de base \`a $X$ est $U \to X$. Alors $U = \lim_{i' \geq i} U_{i'}$
o\`u $U_{i'} = U_i \times_{X_i} X_{i'}$.
Quitte \`a augmenter $i$, nous pouvons supposer $U_i$ affine, voir
Limites, Lemme \ref{limits-lemma-limit-affine}.
Pour v\'erifier que $U_i \to X_i$ est \'etale pour $i$ assez grand,
choisissons un recouvrement ouvert affine fini $U_i = U_{i, 1} \cup \ldots \cup U_{i, m}$
tel que $U_{i, j} \to U_i \to X_i$ se factorise par un ouvert affine
$W_{i, j} \subset X_i$. On peut alors appliquer
Limites, Lemme \ref{limits-lemma-descend-etale},
pour voir que $U_{i, j} \to W_{i, j}$ est \'etale
apr\`es avoir, si n\'ecessaire, augment\'e $i$.
Nous voyons ainsi que le foncteur
$\colim (X_i)_{affine, \etale} \to X_{affine, \etale}$
est essentiellement surjectif. Sa pleine fid\'elit\'e r\'esulte
directement du r\'esultat d\'ej\`a utilis\'e :
Limites, Lemme \ref{limits-lemma-descend-finite-presentation-variant}.
L'assertion concernant les recouvrements se d\'emontre exactement de la m\^eme
mani\`ere que dans le premier paragraphe de la d\'emonstration.
\end{proof}

\noindent
Ce qui pr\'ec\`ede donne le r\'esultat g\'en\'eral suivant sur les limites inductives et la cohomologie.

\begin{theorem}
\label{theorem-colimit}
Soit $X = \lim_{i \in I} X_i$ la limite d'un syst\`eme projectif de sch\'emas
dont les morphismes de transition $f_{i'i} : X_{i'} \to X_i$ sont affines. Nous supposons
que $X_i$ est quasi-compact et quasi-s\'epar\'e pour tout $i \in I$.
Soit $(\mathcal{F}_i, \varphi_{i'i})$ un syst\`eme de faisceaux ab\'eliens
sur $(X_i, f_{i'i})$. Notons $f_i : X \to X_i$ la projection et posons
$\mathcal{F} = \colim f_i^{-1}\mathcal{F}_i$. Alors
$$
\colim_{i\in I} H_\etale^p(X_i, \mathcal{F}_i) = H_\etale^p(X, \mathcal{F}).
$$
pour tout $p \geq 0$.
\end{theorem}

\begin{proof}
D'apr\`es Topologies, lemme \ref{topologies-lemma-alternative},
nous pouvons calculer la cohomologie
de $\mathcal{F}$ sur $X_{affine, \etale}$.
Le r\'esultat d\'ecoule donc de la combinaison du
lemme \ref{lemma-colimit-affine-sites}
et du
lemme \ref{sites-cohomology-lemma-colimit} de Cohomologie sur les sites.
\end{proof}

\noindent
Les deux r\'esultats suivants sont des cas particuliers du th\'eor\`eme ci-dessus.

\begin{lemma}
\label{lemma-colimit}
Soit $X$ un sch\'ema quasi-compact et quasi-s\'epar\'e. Soit $I$
un ensemble ordonn\'e filtrant. Soit $(\mathcal{F}_i, \varphi_{ij})$ un syst\`eme
de faisceaux ab\'eliens sur $X_\etale$ index\'e par $I$. Alors
$$
\colim_{i\in I} H_\etale^p(X, \mathcal{F}_i) = H_\etale^p(X,
\colim_{i\in I} \mathcal{F}_i).
$$
\end{lemma}

\begin{proof}
C'est un cas particulier du th\'eor\`eme \ref{theorem-colimit}.
Nous esquissons \'egalement une d\'emonstration directe.
Nous le d\'emontrons simultan\'ement pour tout $X$, par r\'ecurrence sur $p$.
\begin{enumerate}
\item
Pour tout sch\'ema quasi-compact et quasi-s\'epar\'e $X$ et tout recouvrement \'etale
$\mathcal{U}$ de $X$, montrer qu'il existe un raffinement
$\mathcal{V} = \{V_j \to X\}_{j\in J}$ avec $J$ fini et chaque $V_j$
quasi-compact et quasi-s\'epar\'e, de telle sorte que tous les
$V_{j_0} \times_X \ldots \times_X V_{j_p}$ soient eux aussi
quasi-compacts et quasi-s\'epar\'es.
\item
En utilisant l'\'etape pr\'ec\'edente et la d\'efinition des limites inductives dans la cat\'egorie des
faisceaux, montrer que le th\'eor\`eme vaut pour $p = 0$ et tout $X$.
\item
En utilisant la localit\'e de la cohomologie
(lemme \ref{lemma-locality-cohomology}),
la suite spectrale de {\v C}ech vers la cohomologie
(th\'eor\`eme \ref{theorem-cech-ss}) et le fait que l'hypoth\`ese de r\'ecurrence
s'applique \`a tous les
$V_{j_0} \times_X \ldots \times_X V_{j_p}$
dans la situation ci-dessus, d\'emontrer l'\'etape de r\'ecurrence $p \to p + 1$.
\end{enumerate}
\end{proof}

\begin{lemma}
\label{lemma-directed-colimit-cohomology}
Soient $A$ un anneau, $(I, \leq)$ un ensemble ordonn\'e filtrant et $(B_i, \varphi_{ij})$ un
syst\`eme de $A$-alg\`ebres. Posons $B = \colim_{i\in I} B_i$. Soit $X \to \Spec(A)$
un morphisme de sch\'emas quasi-compact et quasi-s\'epar\'e. Soit
$\mathcal{F}$ un faisceau ab\'elien sur $X_\etale$.
Notons $Y_i = X \times_{\Spec(A)} \Spec(B_i)$,
$Y = X \times_{\Spec(A)} \Spec(B)$,
$\mathcal{G}_i = (Y_i \to X)^{-1}\mathcal{F}$ et
$\mathcal{G} = (Y \to X)^{-1}\mathcal{F}$. Alors
$$
H_\etale^p(Y, \mathcal{G}) =
\colim_{i\in I} H_\etale^p (Y_i, \mathcal{G}_i).
$$
\end{lemma}

\begin{proof}
C'est un cas particulier du th\'eor\`eme \ref{theorem-colimit}.
Nous indiquons \'egalement une d\'emonstration directe comme suit.
\begin{enumerate}
\item \'Etant donn\'e $V \to Y$ \'etale, avec $V$ quasi-compact et
quasi-s\'epar\'e, il existe $i\in I$ et $V_i \to Y_i$ tels que
$V = V_i \times_{Y_i} Y$.
Si tous les sch\'emas consid\'er\'es \'etaient affines, cela correspondrait \`a
l'\'enonc\'e alg\'ebrique suivant : si $B = \colim B_i$ et si $B \to C$ est \'etale,
alors il existe $i \in I$ et $B_i \to C_i$ \'etale tels que
$C \cong B \otimes_{B_i} C_i$.
Ceci est d\'emontr\'e dans Alg\`ebre, lemme \ref{algebra-lemma-etale}.
\item Dans la situation de (1), montrer que
$\mathcal{G}(V) = \colim_{i' \geq i} \mathcal{G}_{i'}(V_{i'})$
o\`u $V_{i'}$ est le changement de base de $V_i$ sur $Y_{i'}$.
\item D'apr\`es (1), on voit que pour tout recouvrement \'etale
$\mathcal{V} = \{V_j \to Y\}_{j\in J}$ avec $J$ fini et les
$V_j$ quasi-compacts et quasi-s\'epar\'es, il existe $i \in I$ et
un recouvrement \'etale $\mathcal{V}_i = \{V_{ij} \to Y_i\}_{j \in J}$
tel que $\mathcal{V} \cong \mathcal{V}_i \times_{Y_i} Y$.
\item Montrer que (2) et (3) impliquent
$$
\check H^*(\mathcal{V}, \mathcal{G})=
\colim_{i\in I} \check H^*(\mathcal{V}_i, \mathcal{G}_i).
$$
\item Utiliser judicieusement la suite spectrale de {\v C}ech vers la cohomologie
(th\'eor\`eme \ref{theorem-cech-ss}).
\end{enumerate}
\end{proof}

\begin{lemma}
\label{lemma-higher-direct-images}
Soit $f: X\to Y$ un morphisme de sch\'emas et soit $\mathcal{F}\in
\textit{Ab}(X_\etale)$. Alors $R^pf_*\mathcal{F}$ est le faisceau
associ\'e au pr\'efaisceau
$$
(V \to Y) \longmapsto H_\etale^p(X \times_Y V, \mathcal{F}|_{X \times_Y V}).
$$
Plus g\'en\'eralement, pour $K \in D(X_\etale)$, $R^pf_*K$ est le
faisceau associ\'e au pr\'efaisceau
$$
(V \to Y) \longmapsto H_\etale^p(X \times_Y V, K|_{X \times_Y V}).
$$
\end{lemma}

\begin{proof}
Ce lemme est valable pour les espaces topologiques, et la d\'emonstration dans ce cas est la
m\^eme. Voir Cohomologie sur les sites, lemme
\ref{sites-cohomology-lemma-higher-direct-images}
pour le cas d'un faisceau, et voir
Cohomologie sur les sites, lemme
\ref{sites-cohomology-lemma-sheafification-cohomology}
pour le cas d'un complexe de faisceaux ab\'eliens.
\end{proof}

\begin{lemma}
\label{lemma-relative-colimit}
Soit $S$ un sch\'ema. Soit $X = \lim_{i \in I} X_i$ une limite d'un
syst\`eme projectif de sch\'emas sur $S$ avec morphismes de transition affines
$f_{i'i} : X_{i'} \to X_i$. Nous supposons que les morphismes structuraux
$g_i : X_i \to S$ et $g : X \to S$ sont quasi-compacts et quasi-s\'epar\'es.
Soit $(\mathcal{F}_i, \varphi_{i'i})$ un syst\`eme de faisceaux ab\'eliens
sur $(X_i, f_{i'i})$. Notons $f_i : X \to X_i$ la projection et posons
$\mathcal{F} = \colim f_i^{-1}\mathcal{F}_i$. Alors
$$
\colim_{i\in I} R^p g_{i, *} \mathcal{F}_i = R^p g_* \mathcal{F}
$$
pour tout $p \geq 0$.
\end{lemma}

\begin{proof}
Rappelons (lemme \ref{lemma-higher-direct-images})
que $R^p g_{i, *} \mathcal{F}_i$ est le faisceau associ\'e au
pr\'efaisceau $U \mapsto H^p_\etale(U \times_S X_i, \mathcal{F}_i)$
et de m\^eme pour $R^pg_*\mathcal{F}$. De plus, la limite inductive d'un
syst\`eme de faisceaux est le faisceau associ\'e \`a la limite inductive calcul\'ee dans la cat\'egorie
des pr\'efaisceaux. Remarquons que tout objet de $S_\etale$ admet un recouvrement
par des objets quasi-compacts et quasi-s\'epar\'es (par exemple, des sch\'emas affines).
De plus, si $U$ est un objet quasi-compact et quasi-s\'epar\'e,
alors on a
$$
\colim H^p_\etale(U \times_S X_i, \mathcal{F}_i) =
H^p_\etale(U \times_S X, \mathcal{F})
$$
d'apr\`es le th\'eor\`eme \ref{theorem-colimit}. Le lemme en r\'esulte.
\end{proof}

\begin{lemma}
\label{lemma-relative-colimit-general}
Soit $I$ un ensemble ordonn\'e filtrant. Soit $g_i : X_i \to S_i$ un syst\`eme projectif de
morphismes de sch\'emas index\'e par $I$. Supposons que $g_i$ soit quasi-compact et
quasi-s\'epar\'e et que, pour $i' \geq i$, les morphismes de transition
$f_{i'i} : X_{i'} \to X_i$ et $h_{i'i} : S_{i'} \to S_i$ soient affines.
Soit $g : X \to S$ la limite des morphismes $g_i$, voir
Limites, section \ref{limits-section-limits}.
Notons $f_i : X \to X_i$ et $h_i : S \to S_i$ les projections.
Soit $(\mathcal{F}_i, \varphi_{i'i})$ un syst\`eme de faisceaux
sur $(X_i, f_{i'i})$. Posons $\mathcal{F} = \colim f_i^{-1}\mathcal{F}_i$. Alors
$$
R^p g_* \mathcal{F} =
\colim_{i \in I} h_i^{-1}R^p g_{i, *} \mathcal{F}_i
$$
pour tout $p \geq 0$.
\end{lemma}

\begin{proof}
Comment construit-on l'application du lemme ?
Pour $i' \geq i$, nous avons un diagramme commutatif
$$
\xymatrix{
X \ar[r]_{f_{i'}} \ar[d]_g &
X_{i'} \ar[r]_{f_{i'i}} \ar[d]_{g_{i'}} &
X_i \ar[d]^{g_i} \\
S \ar[r]^{h_{i'}} &
S_{i'} \ar[r]^{h_{i'i}} &
S_i
}
$$
Si nous composons l'application de changement de base
$h_{i'i}^{-1}Rg_{i, *}\mathcal{F}_i \to Rg_{i', *}f_{i'i}^{-1}\mathcal{F}_i$
(Cohomologie sur les sites, lemme
\ref{sites-cohomology-lemma-base-change-map-flat-case} ou
Remarque \ref{sites-cohomology-remark-base-change})
avec l'application $Rg_{i', *}\varphi_{i'i}$, nous obtenons alors
$\psi_{i'i} : h_{i' i}^{-1} R^p g_{i, *} \mathcal{F}_i \to
R^pg_{i', *} \mathcal{F}_{i'}$. De m\^eme, en utilisant le carr\'e de gauche
du diagramme, nous obtenons les applications
$\psi_i : h_i^{-1}R^pg_{i, *}\mathcal{F}_i \to R^pg_*\mathcal{F}$.
Les applications $h_{i'}^{-1}\psi_{i'i}$ et $\psi_i$ sont celles utilis\'ees dans
l'\'enonc\'e du lemme. Pour que cela ait un sens, nous devons v\'erifier que
$\psi_{i''i} = \psi_{i''i'} \circ h_{i''i'}^{-1}\psi_{i'i}$ et
$\psi_{i'} \circ h_{i'}^{-1}\psi_{i'i} = \psi_i$ ; cela r\'esulte
de Cohomologie sur les sites, remarque
\ref{sites-cohomology-remark-compose-base-change-horizontal}.

\medskip\noindent
D\'emonstration de l'\'egalit\'e. Premi\`ere d\'emonstration par
d\'ecalage de dimension\footnote{On peut aussi utiliser cette m\'ethode
pour construire les applications du lemme.}. Pour tout
$U$ affine et \'etale sur $X$, d'apr\`es le Th\'eor\`eme \ref{theorem-colimit},
nous avons
$$
g_*\mathcal{F}(U) =
H^0(U \times_S X, \mathcal{F}) =
\colim H^0(U_i \times_{S_i} X_i, \mathcal{F}_i) =
\colim g_{i, *}\mathcal{F}_i(U_i)
$$
o\`u la limite inductive porte sur les $i$ assez grands tels qu'
il existe un $i$ et un $U_i$ affine et \'etale sur $S_i$
dont le changement de base est $U$ sur $S$ (voir
Lemme \ref{lemma-colimit-affine-sites}).
Le membre de droite est \'egal \`a
$(\colim h_i^{-1}g_{i, *}\mathcal{F}_i)(U)$ d'apr\`es
Sites, Lemme \ref{sites-lemma-colimit}.
Cela d\'emontre le lemme pour $p = 0$.
Si $(\mathcal{G}_i, \varphi_{i'i})$ est un syst\`eme tel que
$\mathcal{G} = \colim f_i^{-1}\mathcal{G}_i$
et tel que $\mathcal{G}_i$ soit un faisceau ab\'elien injectif sur $X_i$
pour tout $i$, alors, pour tout $U$ affine et \'etale sur $X$, d'apr\`es
le Th\'eor\`eme \ref{theorem-colimit}, nous avons
$$
H^p(U \times_S X, \mathcal{G}) =
\colim H^p(U_i \times_{S_i} X_i, \mathcal{G}_i) = 0
$$
pour $p > 0$ (la m\^eme limite inductive que pr\'ec\'edemment). Ainsi $R^pg_*\mathcal{G} = 0$
et nous obtenons le r\'esultat pour $p > 0$ pour un tel syst\`eme.
En g\'en\'eral, nous pouvons choisir une suite exacte courte de syst\`emes
$$
0 \to (\mathcal{F}_i, \varphi_{i'i}) \to
(\mathcal{G}_i, \varphi_{i'i}) \to
(\mathcal{Q}_i, \varphi_{i'i}) \to 0
$$
o\`u $(\mathcal{G}_i, \varphi_{i'i})$ est comme ci-dessus, voir
Cohomologie sur les sites, Lemme
\ref{sites-cohomology-lemma-colim-sites-injective}.
Par r\'ecurrence, le lemme vaut pour $p - 1$ et, d'apr\`es ce qui pr\'ec\`ede, on a
l'annulation en degr\'e $p$ pour $(\mathcal{G}_i, \varphi_{i'i})$.
On en d\'eduit le r\'esultat pour $p$
et $(\mathcal{F}_i, \varphi_{i'i})$ par la suite exacte longue
de cohomologie.

\medskip\noindent
Seconde d\'emonstration.
Rappelons que $S_{affine, \etale} = \colim (S_i)_{affine, \etale}$, voir
le lemme \ref{lemma-colimit-affine-sites}. Ainsi, si $U$ est un objet de
$S_{affine, \etale}$, alors on peut \'ecrire $U = U_i \times_{S_i} S$
pour un certain $i$ et un certain $U_i$ de $(S_i)_{affine, \etale}$, et
$$
(\colim_{i \in I} h_i^{-1}R^p g_{i, *} \mathcal{F}_i)(U) =
\colim_{i' \geq i} (R^p g_{i', *}\mathcal{F}_{i'})(U_i \times_{S_i} S_{i'})
$$
d'apr\`es Sites, Lemme \ref{sites-lemma-colimit}, et la construction des
morphismes de transition du syst\`eme d\'ecrit ci-dessus. Puisque
$R^pg_{i', *}\mathcal{F}_{i'}$ est le faisceau associ\'e au pr\'efaisceau
$U_{i'} \mapsto H^p(U_{i'} \times_{S_{i'}} X_{i'}, \mathcal{F}_{i'})$
et puisque $R^pg_*\mathcal{F}$ est le faisceau associ\'e au pr\'efaisceau
$U \mapsto H^p(U \times_S X, \mathcal{F})$
(Lemme \ref{lemma-higher-direct-images})
on obtient un diagramme commutatif canonique
$$
\xymatrix{
\colim_{i' \geq i}
H^p(U_i \times_{S_i} X_{i'}, \mathcal{F}_{i'}) \ar[r] \ar[d] &
\colim_{i' \geq i}
(R^p g_{i', *}\mathcal{F}_{i'})(U_i \times_{S_i} S_{i'}) \ar[d] \\
H^p(U \times_S X, \mathcal{F}) \ar[r] &
R^pg_*\mathcal{F}(U)
}
$$
Remarquons que la fl\`eche verticale de gauche est un isomorphisme
d'apr\`es le th\'eor\`eme \ref{theorem-colimit}. Nous cherchons \`a montrer que
la fl\`eche verticale de droite est un isomorphisme. Cependant, nous
savons d\'ej\`a que la source et le but de cette fl\`eche
sont des faisceaux sur $S_{affine, \etale}$. Il suffit donc de
montrer : (1) tout \'el\'ement du but provient localement d'un
\'el\'ement de la source, et (2) tout \'el\'ement de la source
qui s'envoie sur z\'ero dans le but s'annule localement.
L'assertion (1) d\'ecoule imm\'ediatement de ce qui pr\'ec\`ede et du fait que
la fl\`eche horizontale inf\'erieure provient d'un morphisme de pr\'efaisceaux
qui devient un isomorphisme apr\`es faisceautisation.
Pour (2), soit $\xi \in  \colim_{i' \geq i}
(R^p g_{i', *}\mathcal{F}_{i'})(U_i \times_{S_i} S_{i'})$
dans le noyau. Choisissons $i' \geq i$ et
$\xi_{i'} \in (R^p g_{i', *}\mathcal{F}_{i'})(U_i \times_{S_i} S_{i'})$
repr\'esentant $\xi$.
Choisissons un recouvrement \'etale standard
$\{U_{i', k} \to U_i \times_{S_i} S_{i'}\}_{k = 1, \ldots, m}$
tel que $\xi_{i'}|_{U_{i', k}}$ provienne de
$\xi_{i', k} \in H^p(U_{i', k} \times_{S_{i'}} X_{i'}, \mathcal{F}_{i'})$.
Puisqu'il suffit de prouver que $\xi$ s'annule localement, nous
pouvons remplacer $U$ par les membres du recouvrement \'etale
$\{U_{i', k} \times_{S_{i'}} S \to U = U_i \times_{S_i} S\}$.
Apr\`es ce remplacement, on voit que $\xi$ est l'image
d'un \'el\'ement $\xi'$ du groupe
$\colim_{i' \geq i} H^p(U_i \times_{S_i} X_{i'}, \mathcal{F}_{i'})$
dans le diagramme ci-dessus. Puisque $\xi'$ s'envoie sur z\'ero dans $R^pg_*\mathcal{F}(U)$
nous pouvons effectuer un autre remplacement et supposer que $\xi'$ s'envoie
sur z\'ero dans $H^p(U \times_S X, \mathcal{F})$.
Toutefois, puisque la fl\`eche verticale de gauche est un isomorphisme,
on conclut alors que $\xi' = 0$, d'o\`u $\xi = 0$, comme voulu.
\end{proof}

\begin{lemma}
\label{lemma-linus-hamann}
Soit $X = \lim_{i \in I} X_i$ la limite d'un syst\`eme projectif de sch\'emas
dont les morphismes de transition $f_{i'i}$ sont affines et dont les morphismes de projection
sont $f_i : X \to X_i$. Soit $\mathcal{F}$ un faisceau sur $X_\etale$. Alors
\begin{enumerate}
\item il existe des morphismes canoniques
$\varphi_{i'i} : f_{i'i}^{-1}f_{i, *}\mathcal{F} \to f_{i', *}\mathcal{F}$
tels que $(f_{i, *}\mathcal{F}, \varphi_{i'i})$ soit un syst\`eme de
faisceaux sur $(X_i, f_{i'i})$ au sens de la
d\'efinition \ref{definition-inverse-system-sheaves}, et
\item $\mathcal{F} = \colim f_i^{-1}f_{i, *}\mathcal{F}$.
\end{enumerate}
\end{lemma}

\begin{proof}
Par Topologies, Lemme \ref{topologies-lemma-alternative}, et le
lemme \ref{lemma-colimit-affine-sites}, ceci est un cas particulier de
Sites, Lemme \ref{sites-lemma-colimit-push-pull}.
\end{proof}

\begin{lemma}
\label{lemma-compute-strangely}
Soit $I$ un ensemble ordonn\'e filtrant. Soit $g_i : X_i \to S_i$ un syst\`eme projectif de
morphismes de sch\'emas index\'e par $I$. Supposons que $g_i$ soit quasi-compact et
quasi-s\'epar\'e et que, pour $i' \geq i$, les morphismes de transition
$X_{i'} \to X_i$ et $S_{i'} \to S_i$ soient affines.
Soit $g : X \to S$ la limite des morphismes $g_i$, voir
Limites, section \ref{limits-section-limits}.
Notons $f_i : X \to X_i$ et $h_i : S \to S_i$ les projections.
Soit $\mathcal{F}$ un faisceau ab\'elien sur $X$. Alors on a
$$
R^pg_*\mathcal{F} = \colim_{i \in I} h_i^{-1}R^pg_{i, *}(f_{i, *}\mathcal{F})
$$
\end{lemma}

\begin{proof}
Combinaison formelle des lemmes \ref{lemma-relative-colimit-general}
et \ref{lemma-linus-hamann}.
\end{proof}











\section{Limites inductives et complexes}
\label{section-colimit-complexes}

\noindent
Dans cette section, nous \'etudions la cohomologie des syst\`emes de complexes
dans divers contextes, en poursuivant l'\'etude des faisceaux commenc\'ee
\`a la section \ref{section-colimit}.
Nous conseillons vivement au lecteur de ne pas lire cette section, sauf si cela est absolument
n\'ecessaire.

\begin{lemma}
\label{lemma-colimit-variant-complexes}
Soit $X = \lim_{i \in I} X_i$ la limite d'un syst\`eme projectif de sch\'emas
dont les morphismes de transition $f_{i'i} : X_{i'} \to X_i$ sont affines. Supposons
que $X_i$ soit quasi-compact et quasi-s\'epar\'e pour tout $i \in I$.
Soit $\mathcal{F}_i^\bullet$ un complexe de faisceaux ab\'eliens sur
$X_{i, \etale}$. Soit $\varphi_{i'i} : f_{i'i}^{-1}\mathcal{F}_i^\bullet \to
\mathcal{F}_{i'}^\bullet$ un morphisme de complexes sur $X_{i, \etale}$
tel que $\varphi_{i''i} = \varphi_{i''i'} \circ f_{i'' i'}^{-1}\varphi_{i'i}$
lorsque $i'' \geq i' \geq i$. Supposons qu'il existe un entier $a$ tel que
$\mathcal{F}_i^n = 0$ pour $n < a$ et tout $i \in I$.
Alors on a
$$
H^p_\etale(X, \colim f_i^{-1}\mathcal{F}_i^\bullet) =
\colim H^p_\etale(X_i, \mathcal{F}^\bullet_i)
$$
o\`u $f_i : X \to X_i$ est la projection.
\end{lemma}

\begin{proof}
Cela r\'esulte du th\'eor\`eme \ref{theorem-colimit}. Posons
$\mathcal{F}^\bullet = \colim f_i^{-1}\mathcal{F}_i^\bullet$.
Le th\'eor\`eme affirme que
$$
\colim_{i \in I} H_\etale^p(X_i, \mathcal{F}_i^n) =
H_\etale^p(X, \mathcal{F}^n)
$$
pour tous $n, p \in \mathbf{Z}$. Utilisons les suites spectrales
$$
E_{1, i}^{s, t} = H_\etale^t(X_i, \mathcal{F}_i^s) \Rightarrow
H_\etale^{s + t}(X_i, \mathcal{F}_i^\bullet)
$$
et
$$
E_1^{s, t} = H_\etale^t(X, \mathcal{F}^s) \Rightarrow
H_\etale^{s + t}(X, \mathcal{F}^\bullet)
$$
de Cat\'egories d\'eriv\'ees, Lemme \ref{derived-lemma-two-ss-complex-functor}.
Puisque $\mathcal{F}_i^n = 0$ pour $n < a$ (avec $a$ ind\'ependant de $i$),
on voit que seul un nombre fini de termes $E_{1, i}^{s, t}$
(ind\'ependant de $i$) et $E_1^{s, t}$ contribuent \`a
$H^q_\etale(X_i, \mathcal{F}_i^\bullet)$ et
$H^q_\etale(X, \mathcal{F}^\bullet)$ et $E_1^{s, t} = \colim E_{i, i}^{s, t}$.
Cela d\'emontre l'assertion. Certains d\'etails sont omis.
(Il existe un autre argument utilisant les troncatures ``b\^etes''
des complexes, qui \'evite d'utiliser les suites spectrales.)
\end{proof}

\begin{lemma}
\label{lemma-direct-sum-bounded-below-cohomology}
Soit $X$ un sch\'ema quasi-compact et quasi-s\'epar\'e. Soit
$K_i \in D(X_\etale)$, $i \in I$, une famille d'objets.
Supposons donn\'e $a \in \mathbf{Z}$ tel que $H^n(K_i) = 0$ pour $n < a$
et $i \in I$. Alors $R\Gamma(X, \bigoplus_i K_i) =
\bigoplus_i R\Gamma(X, K_i)$.
\end{lemma}

\begin{proof}
Il faut montrer que $H^p(X, \bigoplus_i K_i) =
\bigoplus_i H^p(X, K_i)$ pour tout $p \in \mathbf{Z}$.
Choisissons des complexes $\mathcal{F}_i^\bullet$ repr\'esentant $K_i$
tels que $\mathcal{F}_i^n = 0$ pour $n < a$.
La somme directe des complexes $\mathcal{F}_i^\bullet$
repr\'esente l'objet $\bigoplus K_i$ d'apr\`es
Injectifs, Lemme \ref{injectives-lemma-derived-products}.
Puisque $\bigoplus \mathcal{F}^\bullet$ est la limite inductive filtrante
des sommes directes finies, le r\'esultat d\'ecoule
du lemme \ref{lemma-colimit-variant-complexes}.
\end{proof}

\begin{lemma}
\label{lemma-colimit-complexes}
Soit $S$ un sch\'ema. Soit $X = \lim_{i \in I} X_i$ la limite d'un
syst\`eme projectif de sch\'emas sur $S$ avec pour morphismes de transition affines
$f_{i'i} : X_{i'} \to X_i$. Supposons que $X_i$ soit quasi-compact et
quasi-s\'epar\'e pour tout $i \in I$. Soit $K \in D^+(S_\etale)$. Alors
$$
\colim_{i \in I} H_\etale^p(X_i, K|_{X_i}) = H_\etale^p(X, K|_X).
$$
pour tout $p \in \mathbf{Z}$, o\`u $K|_{X_i}$ et $K|_X$
sont les images inverses de $K$ sur $X_i$ et $X$.
\end{lemma}

\begin{proof}
On peut repr\'esenter $K$ par un complexe born\'e inf\'erieurement $\mathcal{G}^\bullet$
de faisceaux ab\'eliens sur $S_\etale$. Disons que $\mathcal{G}^n = 0$ pour $n < a$.
Notons $\mathcal{F}^\bullet_i$ et $\mathcal{F}^\bullet$ les images inverses
de ce complexe sur $X_i$ et $X$. Ces complexes repr\'esentent les
objets $K|_{X_i}$ et $K|_X$, et l'on a
$\mathcal{F}^\bullet = \colim f_i^{-1}\mathcal{F}_i^\bullet$
terme \`a terme. Le lemme d\'ecoule donc du
lemme \ref{lemma-colimit-variant-complexes}.
\end{proof}

\begin{lemma}
\label{lemma-relative-colimit-general-complexes}
Soient $I$, $g_i : X_i \to S_i$, $g : X \to S$, $f_i$, $g_i$, $h_i$ comme dans le
lemme \ref{lemma-relative-colimit-general}.
Soient $0 \in I$ et $K_0 \in D^+(X_{0, \etale})$.
Pour $i \geq 0$, notons $K_i$ l'image inverse de $K_0$ sur $X_i$.
Notons $K$ l'image inverse de $K_0$ sur $X$. Alors
$$
R^pg_*K = \colim_{i \geq 0} h_i^{-1}R^pg_{i, *}K_i
$$
pour tout $p \in \mathbf{Z}$.
\end{lemma}

\begin{proof}
Fixons un entier $p_0 \in \mathbf{Z}$.
Soit $a$ un entier tel que $H^j(K_0) = 0$ pour $j < a$.
Nous allons d\'emontrer que la formule vaut pour tout $p \leq p_0$
par r\'ecurrence descendante sur $a$. Si
$a > p_0$, alors on voit que les membres de gauche et de droite
de la formule sont nuls pour $p \leq p_0$ par annulation triviale, voir
Cat\'egories d\'eriv\'ees, Lemme \ref{derived-lemma-negative-vanishing}.
Supposons $a \leq p_0$. Consid\'erons le triangle distingu\'e
$$
H^a(K_0)[-a] \to K_0 \to \tau_{\geq a + 1}K_0
$$
Le tir\'e en arri\`ere de ce triangle distingu\'e sur $X_i$ et $X$
donne des triangles distingu\'es compatibles pour $K_i$ et $K$.
Pour $p \leq p_0$, consid\'erons le diagramme commutatif
$$
\xymatrix{
\colim_{i \geq 0} h_i^{-1}R^{p - 1}g_{i, *}(\tau_{\geq a + 1}K_i)
\ar[r]_-\alpha \ar[d] &
R^{p - 1}g_*(\tau_{\geq a + 1}K) \ar[d] \\
\colim_{i \geq 0} h_i^{-1}R^pg_{i, *}(H^a(K_i)[-a])
\ar[r]_-\beta \ar[d] &
R^pg_*(H^a(K)[-a]) \ar[d] \\
\colim_{i \geq 0} h_i^{-1}R^pg_{i, *}K_i
\ar[r]_-\gamma \ar[d] &
R^pg_*K \ar[d] \\
\colim_{i \geq 0} R^pg_{i, *}\tau_{\geq a + 1}K_i
\ar[r]_-\delta \ar[d] &
R^pg_*\tau_{\geq a + 1}K \ar[d] \\
\colim_{i \geq 0} R^{p + 1}g_{i, *}(H^a(K_i)[-a])
\ar[r]^-\epsilon &
R^{p + 1}g_*(H^a(K)[-a])
}
$$
\`a colonnes exactes. Les fl\`eches $\beta$ et $\epsilon$ sont des isomorphismes d'apr\`es le
lemme \ref{lemma-relative-colimit-general}.
Les fl\`eches $\alpha$ et $\delta$ sont des isomorphismes
par l'hypoth\`ese de r\'ecurrence.
Ainsi, $\gamma$ est un isomorphisme, comme voulu.
\end{proof}

\begin{lemma}
\label{lemma-relative-colimit-general-really-complexes}
Soient $I$, $g_i : X_i \to S_i$, $g : X \to S$, $f_{ii'}$, $f_i$, $g_i$, $h_i$
comme dans le lemme \ref{lemma-relative-colimit-general}.
Soit $\mathcal{F}_i^\bullet$ un complexe de faisceaux ab\'eliens sur
$X_{i, \etale}$. Soit $\varphi_{i'i} : f_{i'i}^{-1}\mathcal{F}_i^\bullet \to
\mathcal{F}_{i'}^\bullet$ un morphisme de complexes sur $X_{i, \etale}$
tel que $\varphi_{i''i} = \varphi_{i''i'} \circ f_{i'' i'}^{-1}\varphi_{i'i}$
lorsque $i'' \geq i' \geq i$. Supposons qu'il existe un entier $a$ tel que
$\mathcal{F}_i^n = 0$ pour $n < a$ et tout $i \in I$. Alors
$$
R^pg_*(\colim f_i^{-1}\mathcal{F}_i^\bullet) =
\colim_{i \geq 0} h_i^{-1}R^pg_{i, *}\mathcal{F}_i^\bullet
$$
pour tout $p \in \mathbf{Z}$.
\end{lemma}

\begin{proof}
Ceci r\'esulte du lemme \ref{lemma-relative-colimit-general}. Posons
$\mathcal{F}^\bullet = \colim f_i^{-1}\mathcal{F}_i^\bullet$.
Le lemme affirme que
$$
\colim_{i \in I} h_i^{-1}R^pg_{i, *}\mathcal{F}_i^n = R^pg_*\mathcal{F}^n
$$
pour tous $n, p \in \mathbf{Z}$. Utilisons les suites spectrales
$$
E_{1, i}^{s, t} = R^tg_{i, *}\mathcal{F}_i^s \Rightarrow
R^{s + t}g_{i, *}\mathcal{F}_i^\bullet
$$
et
$$
E_1^{s, t} = R^tg_*\mathcal{F}^s \Rightarrow
R^{s + t}g_*\mathcal{F}^\bullet
$$
de Cat\'egories d\'eriv\'ees, Lemme \ref{derived-lemma-two-ss-complex-functor}.
Puisque $\mathcal{F}_i^n = 0$ pour $n < a$ (avec $a$ ind\'ependant de $i$),
on voit que seul un nombre fini fix\'e de termes $E_{1, i}^{s, t}$
(ind\'ependant de $i$) et $E_1^{s, t}$ contribuent et que
$E_1^{s, t} = \colim E_{i, i}^{s, t}$.
Cela d\'emontre l'assertion. Certains d\'etails sont omis.
(Il existe un autre argument utilisant les troncatures ``b\^etes''
des complexes, qui \'evite d'utiliser les suites spectrales.)
\end{proof}

\begin{lemma}
\label{lemma-direct-sum-bounded-below-pushforward}
Soit $f : X \to Y$ un morphisme de sch\'emas quasi-compact et quasi-s\'epar\'e.
Soit $K_i \in D(X_\etale)$, $i \in I$, une famille d'objets.
Supposons donn\'e $a \in \mathbf{Z}$ tel que $H^n(K_i) = 0$ pour $n < a$
et $i \in I$. Alors $Rf_*(\bigoplus_i K_i) = \bigoplus_i Rf_*K_i$.
\end{lemma}

\begin{proof}
Il faut montrer que $R^pf_*(\bigoplus_i K_i) =
\bigoplus_i R^pf_*K_i$ pour tout $p \in \mathbf{Z}$.
Choisissons des complexes $\mathcal{F}_i^\bullet$ repr\'esentant $K_i$
tels que $\mathcal{F}_i^n = 0$ pour $n < a$.
La somme directe des complexes $\mathcal{F}_i^\bullet$
repr\'esente l'objet $\bigoplus K_i$ d'apr\`es
Injectifs, Lemme \ref{injectives-lemma-derived-products}.
Puisque $\bigoplus \mathcal{F}^\bullet$ est la limite inductive filtrante
des sommes directes finies, le r\'esultat d\'ecoule
du lemme \ref{lemma-relative-colimit-general-really-complexes}.
\end{proof}














\section{Fibres des images directes sup\'erieures}
\label{section-stalks-direct-image}

\noindent
Les fibres des images directes sup\'erieures se calculent souvent comme suit.

\begin{theorem}
\label{theorem-higher-direct-images}
Soit $f: X \to S$ un morphisme de sch\'emas quasi-compact et quasi-s\'epar\'e,
$\mathcal{F}$ un faisceau ab\'elien sur $X_\etale$, et $\overline{s}$ un
point g\'eom\'etrique de $S$ au-dessus de $s \in S$. Alors
$$
\left(R^nf_* \mathcal{F}\right)_{\overline{s}} =
H_\etale^n( X \times_S \Spec(\mathcal{O}_{S, \overline{s}}^{sh}),
p^{-1}\mathcal{F})
$$
o\`u $p : X \times_S \Spec(\mathcal{O}_{S, \overline{s}}^{sh}) \to X$
est la projection. Pour $K \in D^+(X_\etale)$ et $n \in \mathbf{Z}$,
on a
$$
\left(R^nf_*K\right)_{\overline{s}} =
H_\etale^n(X \times_S \Spec(\mathcal{O}_{S, \overline{s}}^{sh}), p^{-1}K)
$$
En fait, on a
$$
\left(Rf_*K\right)_{\overline{s}}
=
R\Gamma_\etale(X \times_S \Spec(\mathcal{O}_{S, \overline{s}}^{sh}), p^{-1}K)
$$
dans $D^+(\textit{Ab})$.
\end{theorem}

\begin{proof}
Soit $\mathcal{I}$ la cat\'egorie des voisinages \'etales de $\overline{s}$
sur $S$. D'apr\`es le lemme \ref{lemma-higher-direct-images},
on a
$$
(R^nf_*\mathcal{F})_{\overline{s}} =
\colim_{(V, \overline{v}) \in \mathcal{I}^{opp}}
H_\etale^n(X \times_S V, \mathcal{F}|_{X \times_S V}).
$$
Nous pouvons remplacer $\mathcal{I}$ par la sous-cat\'egorie initiale form\'ee
des voisinages \'etales affines de $\overline{s}$. Remarquons que
$$
\Spec(\mathcal{O}_{S, \overline{s}}^{sh}) =
\lim_{(V, \overline{v}) \in \mathcal{I}} V
$$
d'apr\`es le lemme \ref{lemma-describe-etale-local-ring} et
Limites, Lemme
\ref{limits-lemma-directed-inverse-system-affine-schemes-has-limit}.
Puisque les produits fibr\'es commutent aux limites, on obtient aussi
$$
X \times_S \Spec(\mathcal{O}_{S, \overline{s}}^{sh}) =
\lim_{(V, \overline{v}) \in \mathcal{I}} X \times_S V
$$
On conclut par le lemme \ref{lemma-directed-colimit-cohomology}.
Pour la seconde variante, on emploie le m\^eme argument avec
le lemme \ref{lemma-colimit-complexes} au lieu du
lemme \ref{lemma-directed-colimit-cohomology}.

\medskip\noindent
Pour voir que la derni\`ere assertion est vraie, il suffit de construire un morphisme
$\left(Rf_*K\right)_{\overline{s}} \to
R\Gamma_\etale(X \times_S \Spec(\mathcal{O}_{S, \overline{s}}^{sh}), p^{-1}K)$
dans $D^+(\textit{Ab})$ qui r\'ealise les isomorphismes ci-dessus sur les groupes de cohomologie
en degr\'e $n$ pour tout $n$. Pour cela, choisissons un
complexe born\'e inf\'erieurement $\mathcal{J}^\bullet$ de faisceaux ab\'eliens injectifs
sur $X_\etale$ repr\'esentant $K$. Le complexe
$f_*\mathcal{J}^\bullet$ repr\'esente $Rf_*K$. Ainsi, le complexe
$$
(f_*\mathcal{J}^\bullet)_{\overline{s}} =
\colim_{(V, \overline{v}) \in \mathcal{I}^{opp}}
(f_*\mathcal{J}^\bullet)(V)
$$
repr\'esente $(Rf_*K)_{\overline{s}}$. Pour chaque $V$, on a des morphismes
$$
(f_*\mathcal{J}^\bullet)(V) =
\Gamma(X \times_S V, \mathcal{J}^\bullet)
\longrightarrow
\Gamma(X \times_S \Spec(\mathcal{O}_{S, \overline{s}}^{sh}),
p^{-1}\mathcal{J}^\bullet)
$$
et le complexe cible repr\'esente
$R\Gamma_\etale(X \times_S \Spec(\mathcal{O}_{S, \overline{s}}^{sh}), p^{-1}K)$
dans $D^+(\textit{Ab})$.
En prenant la limite inductive de ces morphismes, on obtient le r\'esultat.
\end{proof}

\begin{remark}
\label{remark-stalk-fibre}
Soit $f : X \to S$ un morphisme de sch\'emas.
Soit $K \in D(X_\etale)$.
Soit $\overline{s}$ un point g\'eom\'etrique de $S$.
Il existe toujours des morphismes canoniques
$$
(Rf_*K)_{\overline{s}}
\longrightarrow
R\Gamma(X \times_S \Spec(\mathcal{O}_{S, \overline{s}}^{sh}), p^{-1}K)
\longrightarrow
R\Gamma(X_{\overline{s}}, K|_{X_{\overline{s}}})
$$
o\`u $p : X \times_S \Spec(\mathcal{O}_{S, \overline{s}}^{sh}) \to X$
est la projection. Consid\'erons en effet le diagramme commutatif
$$
\xymatrix{
X_{\overline{s}} \ar[r] \ar[d]^{f_{\overline{s}}} &
X \times_S \Spec(\mathcal{O}_{S, \overline{s}}^{sh}) \ar[r]_-p \ar[d]^{f'} &
X \ar[d]^f \\
\overline{s} \ar[r]^-i &
\Spec(\mathcal{O}_{S, \overline{s}}^{sh}) \ar[r]^-j &
S
}
$$
On dispose des morphismes de changement de base
$$
i^{-1}Rf'_*(p^{-1}K) \to Rf_{\overline{s}, *}(K|_{X_{\overline{s}}})
\quad\text{et}\quad
j^{-1}Rf_*K \to Rf'_*(p^{-1}K)
$$
(Cohomologie sur les sites, Remarque \ref{sites-cohomology-remark-base-change})
associ\'es aux deux carr\'es de ce diagramme.
En prenant les sections globales, on obtient les morphismes voulus.
D'apr\`es Cohomologie sur les sites, Remarque
\ref{sites-cohomology-remark-compose-base-change-horizontal}
la compos\'ee de ces deux morphismes est le
morphisme usuel de changement de base
$(Rf_*K)_{\overline{s}} \to R\Gamma(X_{\overline{s}}, K|_{X_{\overline{s}}})$.
\end{remark}






\section{La suite spectrale de Leray}
\label{section-leray}

\begin{lemma}
\label{lemma-prepare-leray}
Soient $f: X \to Y$ un morphisme et $\mathcal{I}$ un objet injectif de
$\textit{Ab}(X_\etale)$. Soit $V \in \Ob(Y_\etale)$. Alors
\begin{enumerate}
\item pour tout recouvrement $\mathcal{V} = \{V_j\to V\}_{j \in J}$, on a
$\check H^p(\mathcal{V}, f_*\mathcal{I}) = 0$ pour tout $p > 0$ ;
\item $f_*\mathcal{I}$ est acyclique pour le foncteur $\Gamma(V, -)$ ; et
\item si $g : Y \to Z$, alors $f_*\mathcal{I}$ est acyclique pour $g_*$. 
\end{enumerate}
\end{lemma}

\begin{proof}
Remarquons que $\check{\mathcal{C}}^\bullet(\mathcal{V}, f_*\mathcal{I}) =
\check{\mathcal{C}}^\bullet(\mathcal{V} \times_Y X, \mathcal{I})$
dont les groupes de cohomologie sup\'erieure sont nuls d'apr\`es le lemme \ref{lemma-hom-injective}.
Cela d\'emontre (1). La deuxi\`eme assertion d\'ecoule du fait qu'un faisceau dont
la cohomologie sup\'erieure de {\v C}ech s'annule pour tout recouvrement a des groupes de
cohomologie sup\'erieure nuls. C'est un excellent exercice d'utilisation de la
suite spectrale de {\v C}ech vers la cohomologie, mais voir
Cohomologie sur les sites, Lemme \ref{sites-cohomology-lemma-cech-vanish-collection}
pour les d\'etails et pour un \'enonc\'e plus pr\'ecis et plus g\'en\'eral.
L'assertion (3) r\'esulte de (2) et de la description de
$R^pg_*$ dans le lemme \ref{lemma-higher-direct-images}.
\end{proof}

\noindent
Le formalisme des suites spectrales de Grothendieck donne alors le
r\'esultat suivant.

\begin{proposition}[Suite spectrale de Leray]
\label{proposition-leray}
Soit $f: X \to Y$ un morphisme de sch\'emas et soit $\mathcal{F}$ un faisceau \'etale sur
$X$. Il existe alors une suite spectrale
$$
E_2^{p, q} = H_\etale^p(Y, R^qf_*\mathcal{F}) \Rightarrow
H_\etale^{p+q}(X, \mathcal{F}).
$$
\end{proposition}

\begin{proof}
Voir le lemme \ref{lemma-prepare-leray} et
Cat\'egories d\'eriv\'ees, section
\ref{derived-section-composition-right-derived-functors}.
\end{proof}









\section{Annulation des images directes sup\'erieures finies}
\label{section-vanishing-finite-morphism}

\noindent
Le but suivant est de d\'emontrer que les images directes sup\'erieures d'un morphisme fini de
sch\'emas s'annulent.

\begin{lemma}
\label{lemma-vanishing-etale-cohomology-strictly-henselian}
Soit $R$ un anneau local strictement hens\'elien. Posons $S = \Spec(R)$ et soit
$\overline{s}$ son point ferm\'e. Alors le foncteur des
sections globales
$\Gamma(S, -) : \textit{Ab}(S_\etale) \to \textit{Ab}$
est exact. En fait, $\Gamma(S, \mathcal{F}) = \mathcal{F}_{\overline{s}}$
pour tout faisceau d'ensembles $\mathcal{F}$. En particulier,
$$
\forall p\geq 1, \quad H_\etale^p(S, \mathcal{F})=0
$$
pour tout $\mathcal{F}\in \textit{Ab}(S_\etale)$.
\end{lemma}

\begin{proof}
Si l'on montre que $\Gamma(S, \mathcal{F}) = \mathcal{F}_{\overline{s}}$,
alors $\Gamma(S, -)$ est exact, puisque le foncteur fibre est exact.
Soit $(U, \overline{u})$ un voisinage \'etale de $\overline{s}$.
Choisissons un voisinage ouvert affine $\Spec(A)$ de $\overline{u}$ dans $U$.
Alors $R \to A$ est \'etale et $\kappa(\overline{s}) = \kappa(\overline{u})$.
D'apr\`es le th\'eor\`eme \ref{theorem-henselian}, on voit que $A \cong R \times A'$
comme $R$-alg\`ebre, de mani\`ere compatible avec les morphismes vers
$\kappa(\overline{s}) = \kappa(\overline{u})$.
On obtient donc une section
$$
\xymatrix{
\Spec(A) \ar[r] & U \ar[d]\\
& S \ar[ul]
}
$$
Il s'ensuit que, dans le syst\`eme des voisinages \'etales de $\overline{s}$,
le morphisme identit\'e $(S, \overline{s}) \to (S, \overline{s})$ est cofinal.
Ainsi $\Gamma(S, \mathcal{F}) = \mathcal{F}_{\overline{s}}$.
La derni\`ere assertion du lemme r\'esulte du fait que les foncteurs d\'eriv\'es sup\'erieurs
d'un foncteur exact sont nuls, voir
Cat\'egories d\'eriv\'ees, Lemme \ref{derived-lemma-right-derived-exact-functor}.
\end{proof}

\begin{proposition}
\label{proposition-finite-higher-direct-image-zero}
Soit $f : X \to Y$ un morphisme fini de sch\'emas.
\begin{enumerate}
\item Pour tout point g\'eom\'etrique $\overline{y} : \Spec(k) \to Y$, on a
$$
(f_*\mathcal{F})_{\overline{y}} =
\prod\nolimits_{\overline{x} : \Spec(k) \to X,\ f(\overline{x}) =
\overline{y}} \mathcal{F}_{\overline{x}}.
$$
pour $\mathcal{F}$ dans $\Sh(X_\etale)$ et
$$
(f_*\mathcal{F})_{\overline{y}} =
\bigoplus\nolimits_{\overline{x} : \Spec(k) \to X,\ f(\overline{x}) =
\overline{y}} \mathcal{F}_{\overline{x}}.
$$
pour $\mathcal{F}$ dans $\textit{Ab}(X_\etale)$.
\item Pour tout $q \geq 1$, on a $R^q f_*\mathcal{F} = 0$
pour $\mathcal{F}$ dans $\textit{Ab}(X_\etale)$.
\end{enumerate}
\end{proposition}

\begin{proof}
Notons $X_{\overline{y}}^{sh}$ le produit fibr\'e
$X \times_Y \Spec(\mathcal{O}_{Y, \overline{y}}^{sh})$.
D'apr\`es le th\'eor\`eme \ref{theorem-higher-direct-images},
la fibre de $R^qf_*\mathcal{F}$ en $\overline{y}$ se calcule par
$H_\etale^q(X_{\overline{y}}^{sh}, \mathcal{F})$.
Puisque $f$ est fini, $X_{\bar y}^{sh}$ est fini sur
$\Spec(\mathcal{O}_{Y, \overline{y}}^{sh})$ ; ainsi,
$X_{\bar y}^{sh} = \Spec(A)$ pour un anneau $A$
fini sur $\mathcal{O}_{Y, \bar y}^{sh}$.
Comme ce dernier est strictement hens\'elien,
le lemme \ref{lemma-finite-over-henselian}
implique que $A$ est un produit fini d'anneaux locaux hens\'eliens
$A = A_1 \times \ldots \times A_r$. Puisque le corps r\'esiduel de
$\mathcal{O}_{Y, \overline{y}}^{sh}$ est s\'eparablement clos, il en va de
m\^eme pour chaque $A_i$. Ainsi $A_i$ est strictement hens\'elien.
Il s'ensuit que $X_{\overline{y}}^{sh} = \coprod_{i = 1}^r \Spec(A_i)$.
L'annulation du
lemme \ref{lemma-vanishing-etale-cohomology-strictly-henselian}
implique que $(R^qf_*\mathcal{F})_{\overline{y}} = 0$ pour $q > 0$ ;
cela implique (2) d'apr\`es le th\'eor\`eme \ref{theorem-exactness-stalks}.
L'assertion (1) r\'esulte de l'assertion correspondante du
lemme \ref{lemma-vanishing-etale-cohomology-strictly-henselian}.
\end{proof}

\begin{lemma}
\label{lemma-finite-pushforward-commutes-with-base-change}
Consid\'erons un carr\'e cart\'esien
$$
\xymatrix{
X' \ar[r]_{g'} \ar[d]_{f'} & X \ar[d]^f \\
Y' \ar[r]^g & Y
}
$$
de sch\'emas, avec $f$ un morphisme fini.
Pour tout faisceau d'ensembles $\mathcal{F}$ sur $X_\etale$, on a
$f'_*(g')^{-1}\mathcal{F} = g^{-1}f_*\mathcal{F}$.
\end{lemma}

\begin{proof}
Dans une grande g\'en\'eralit\'e, il existe un morphisme de changement de base
$g^{-1}f_*\mathcal{F} \to f'_*(g')^{-1}\mathcal{F}$, voir
Sites, section \ref{sites-section-pullback}.
Il suffit de v\'erifier l'assertion sur les fibres (th\'eor\`eme \ref{theorem-exactness-stalks}).
Soit $\overline{y}' : \Spec(k) \to Y'$ un point g\'eom\'etrique.
On a
\begin{align*}
(f'_*(g')^{-1}\mathcal{F})_{\overline{y}'}
& =
\prod\nolimits_{\overline{x}' : \Spec(k) \to X',\ f' \circ \overline{x}' =
\overline{y}'}
((g')^{-1}\mathcal{F})_{\overline{x}'} \\
& =
\prod\nolimits_{\overline{x}' : \Spec(k) \to X',\ f' \circ \overline{x}' =
\overline{y}'} \mathcal{F}_{g' \circ \overline{x}'} \\
& =
\prod\nolimits_{\overline{x} : \Spec(k) \to X,\ f \circ \overline{x} =
g \circ \overline{y}'} \mathcal{F}_{\overline{x}} \\
& =
(f_*\mathcal{F})_{g \circ \overline{y}'} \\
& =
(g^{-1}f_*\mathcal{F})_{\overline{y}'}
\end{align*}
La premi\`ere \'egalit\'e r\'esulte de la
proposition \ref{proposition-finite-higher-direct-image-zero}.
La deuxi\`eme \'egalit\'e r\'esulte du
lemme \ref{lemma-stalk-pullback}.
La troisi\`eme \'egalit\'e vaut parce que le diagramme est un carr\'e cart\'esien
et que, par cons\'equent, l'application
$$
\{\overline{x}' : \Spec(k) \to X',\ f' \circ \overline{x}' =
\overline{y}'\}
\longrightarrow
\{\overline{x} : \Spec(k) \to X,\ f \circ \overline{x} =
g \circ \overline{y}'\}
$$
qui envoie $\overline{x}'$ sur $g' \circ \overline{x}'$ est bijective.
La quatri\`eme \'egalit\'e r\'esulte de la
proposition \ref{proposition-finite-higher-direct-image-zero}.
La cinqui\`eme \'egalit\'e r\'esulte du
lemme \ref{lemma-stalk-pullback}.
\end{proof}

\begin{lemma}
\label{lemma-integral-pushforward-commutes-with-base-change}
Consid\'erons un carr\'e cart\'esien
$$
\xymatrix{
X' \ar[r]_{g'} \ar[d]_{f'} & X \ar[d]^f \\
Y' \ar[r]^g & Y
}
$$
de sch\'emas, avec $f$ un morphisme entier.
Pour tout faisceau d'ensembles $\mathcal{F}$ sur $X_\etale$, on a
$f'_*(g')^{-1}\mathcal{F} = g^{-1}f_*\mathcal{F}$.
\end{lemma}

\begin{proof}
La question est locale sur $Y$ ; nous pouvons donc supposer $Y$ affine.
On peut alors \'ecrire $X = \lim X_i$, o\`u $f_i : X_i \to Y$ est fini
(c'est imm\'ediat dans le cas affine, mais voir
Limites, Lemme \ref{limits-lemma-integral-limit-finite-and-finite-presentation}
pour une r\'ef\'erence). Notons $p_{i'i} : X_{i'} \to X_i$
les morphismes de transition et $p_i : X \to X_i$ les projections.
En posant $\mathcal{F}_i = p_{i, *}\mathcal{F}$, on obtient
du lemme \ref{lemma-linus-hamann}
un syst\`eme $(\mathcal{F}_i, \varphi_{i'i})$
tel que $\mathcal{F} = \colim p_i^{-1}\mathcal{F}_i$.
On obtient $f_*\mathcal{F} = \colim f_{i, *}\mathcal{F}_i$
par le lemme \ref{lemma-relative-colimit}.
Posons $X'_i = Y' \times_Y X_i$, de projections $f'_i$ et $g'_i$.
Alors $X' = \lim X'_i$, puisque les limites commutent aux limites.
Notons $p'_i : X' \to X'_i$ les projections. On a
\begin{align*}
g^{-1}f_*\mathcal{F}
& =
g^{-1} \colim f_{i, *}\mathcal{F}_i \\
& =
\colim g^{-1}f_{i, *}\mathcal{F}_i \\
& =
\colim f'_{i, *}(g'_i)^{-1}\mathcal{F}_i \\
& =
f'_*(\colim (p'_i)^{-1}(g'_i)^{-1}\mathcal{F}_i) \\
& =
f'_*(\colim (g')^{-1}p_i^{-1}\mathcal{F}_i) \\
& =
f'_*(g')^{-1} \colim p_i^{-1}\mathcal{F}_i \\
& =
f'_*(g')^{-1}\mathcal{F}
\end{align*}
comme voulu. Pour la premi\`ere \'egalit\'e, voir ci-dessus.
Pour la deuxi\`eme, on utilise le fait que l'image inverse commute aux limites inductives.
Pour la troisi\`eme, on utilise le cas fini, voir
le lemme \ref{lemma-finite-pushforward-commutes-with-base-change}.
Pour la quatri\`eme, on utilise le lemme \ref{lemma-relative-colimit}.
Pour la cinqui\`eme, on utilise $g'_i \circ p'_i = p_i \circ g'$.
Pour la sixi\`eme, on utilise le fait que l'image inverse commute aux limites inductives.
Pour la septi\`eme, on utilise $\mathcal{F} = \colim p_i^{-1}\mathcal{F}_i$.
\end{proof}

\noindent
Le lemme suivant est un cas de descente cohomologique portant sur
les faisceaux \'etales et les morphismes finis surjectifs. Nous g\'en\'eraliserons
consid\'erablement ce r\'esultat une fois d\'emontr\'e le th\'eor\`eme de changement de base propre.

\begin{lemma}
\label{lemma-cohomological-descent-finite}
Soit $f : X \to Y$ un morphisme fini surjectif de sch\'emas.
Posons $f_n : X_n \to Y$ \'egal au produit fibr\'e it\'er\'e $(n + 1)$ fois
de $X$ au-dessus de $Y$. Pour $\mathcal{F} \in \textit{Ab}(Y_\etale)$, posons
$\mathcal{F}_n = f_{n, *}f_n^{-1}\mathcal{F}$. Il existe une suite
exacte
$$
0 \to \mathcal{F} \to \mathcal{F}_0 \to \mathcal{F}_1 \to
\mathcal{F}_2 \to \ldots
$$
sur $Y_\etale$. De plus, il existe une suite spectrale
$$
E_1^{p, q} = H^q_\etale(X_p, f_p^{-1}\mathcal{F})
$$
convergeant vers $H^{p + q}(Y_\etale, \mathcal{F})$.
Cette suite spectrale est fonctorielle en $\mathcal{F}$.
\end{lemma}

\begin{proof}
Si nous d\'emontrons la premi\`ere assertion du lemme, nous obtenons une suite
spectrale de terme $E_1^{p, q} = H^q_\etale(Y, \mathcal{F}_p)$ convergeant
vers $H^{p + q}(Y_\etale, \mathcal{F})$, voir
Cat\'egories d\'eriv\'ees, lemme \ref{derived-lemma-two-ss-complex-functor}.
D'autre part, puisque
$R^if_{p, *}f_p^{-1}\mathcal{F} = 0$ pour $i > 0$
(proposition \ref{proposition-finite-higher-direct-image-zero})
on obtient
$$
H^q_\etale(X_p, f_p^{-1}\mathcal{F}) =
H^q_\etale(Y, f_{p, *}f_p^{-1} \mathcal{F}) =
H^q_\etale(Y, \mathcal{F}_p)
$$
d'apr\`es la proposition \ref{proposition-leray},
ce qui donne la suite spectrale du lemme.

\medskip\noindent
Pour d\'emontrer la premi\`ere assertion du lemme, observons que
les $X_n$ forment un sch\'ema simplicial au-dessus de $Y$, voir
Simplicial, exemple \ref{simplicial-example-fibre-products-simplicial-object}.
Observons de plus que, pour chacune des projections
$d_j : X_{n + 1} \to X_n$, il existe un morphisme
$d_j^{-1} f_n^{-1}\mathcal{F} \to f_{n + 1}^{-1}\mathcal{F}$.
Ces morphismes induisent des morphismes
$$
\delta_j : \mathcal{F}_n \to \mathcal{F}_{n + 1}
$$
pour $j = 0, \ldots, n + 1$. Nous utilisons la somme altern\'ee de ces morphismes
pour d\'efinir les diff\'erentielles $\mathcal{F}_n \to \mathcal{F}_{n + 1}$.
De m\^eme, il existe une augmentation canonique $\mathcal{F} \to \mathcal{F}_0$,
qui est simplement le morphisme canonique $\mathcal{F} \to f_*f^{-1}\mathcal{F}$.
Pour v\'erifier que cette suite de faisceaux est un complexe exact, il suffit
de le v\'erifier sur les fibres aux points g\'eom\'etriques (th\'eor\`eme \ref{theorem-exactness-stalks}).
Soit donc $\overline{y} : \Spec(k) \to Y$ un point g\'eom\'etrique. Posons
$E = \{\overline{x} : \Spec(k) \to X \mid f(\overline{x}) = \overline{y}\}$.
Alors $E$ est un ensemble fini non vide et l'on voit que
$$
(\mathcal{F}_n)_{\overline{y}} =
\bigoplus\nolimits_{e \in E^{n + 1}} \mathcal{F}_{\overline{y}}
$$
d'apr\`es la proposition \ref{proposition-finite-higher-direct-image-zero}
et le lemme \ref{lemma-stalk-pullback}.
Il reste donc \`a voir que, pour tout groupe ab\'elien $M$, la suite
$$
0 \to M \to \bigoplus\nolimits_{e \in E} M \to
\bigoplus\nolimits_{e \in E^2} M \to \ldots
$$
est exacte. Ici, le premier morphisme est le morphisme diagonal et le morphisme
$\bigoplus_{e \in E^{n + 1}} M  \to \bigoplus_{e \in E^{n + 2}} M$
est la somme altern\'ee des morphismes induits par les $(n + 2)$
projections $E^{n + 2} \to E^{n + 1}$. Cela se d\'emontre directement
ou se d\'eduit en appliquant Simplicial, lemme
\ref{simplicial-lemma-fibre-products-simplicial-object-w-section}
au morphisme $E \to \{*\}$.
\end{proof}

\begin{remark}
\label{remark-cohomological-descent-finite}
Dans la situation du lemme \ref{lemma-cohomological-descent-finite},
si $\mathcal{G}$ est un faisceau d'ensembles sur $Y_\etale$, alors on a
$$
\Gamma(Y, \mathcal{G}) =
\text{\'Egaliseur}(
\xymatrix{
\Gamma(X_0, f_0^{-1}\mathcal{G})
\ar@<1ex>[r] \ar@<-1ex>[r] &
\Gamma(X_1, f_1^{-1}\mathcal{G})
}
)
$$
Cela se d\'emontre exactement de la m\^eme mani\`ere, en montrant que
le faisceau $\mathcal{G}$ est l'\'egalisateur des deux morphismes
$f_{0, *}f_0^{-1}\mathcal{G} \to f_{1, *}f_1^{-1}\mathcal{G}$.
\end{remark}









%10.06.09
\section{Action galoisienne sur les fibres}
\label{section-galois-action-stalks}

\noindent
Dans cette section, nous d\'efinissons une action du groupe de Galois absolu du corps
r\'esiduel d'un point $s$ de $S$ sur le foncteur fibre en tout point g\'eom\'etrique situ\'e
au-dessus de $s$.

\medskip\noindent
Action galoisienne sur les fibres.
Soit $S$ un sch\'ema.
Soit $\overline{s}$ un point g\'eom\'etrique de $S$.
Soit $\sigma \in \text{Aut}(\kappa(\overline{s})/\kappa(s))$.
D\'efinissons comme suit une action de $\sigma$ sur la fibre $\mathcal{F}_{\overline{s}}$
d'un faisceau $\mathcal{F}$
\begin{equation}
\label{equation-galois-action}
\begin{matrix}
\mathcal{F}_{\overline{s}} &
\longrightarrow &
\mathcal{F}_{\overline{s}} \\
(U, \overline{u}, t) &
\longmapsto &
(U, \overline{u} \circ \Spec(\sigma), t).
\end{matrix}
\end{equation}
o\`u nous utilisons la description des \'el\'ements de la fibre au moyen de triplets,
comme dans la discussion qui suit la
d\'efinition \ref{definition-stalk}.
Il s'agit d'une action \`a gauche, car si
$\sigma_i \in \text{Aut}(\kappa(\overline{s})/\kappa(s))$
alors
\begin{align*}
\sigma_1 \cdot (\sigma_2 \cdot (U, \overline{u}, t))
& =
\sigma_1 \cdot (U, \overline{u} \circ \Spec(\sigma_2), t) \\
& =
(U, \overline{u} \circ \Spec(\sigma_2) \circ \Spec(\sigma_1), t) \\
& =
(U, \overline{u} \circ \Spec(\sigma_1 \circ \sigma_2), t) \\
& =
(\sigma_1 \circ \sigma_2) \cdot (U, \overline{u}, t)
\end{align*}
Il est clair que cette action est fonctorielle par rapport au faisceau $\mathcal{F}$.
Notons que nous aurions pu d\'efinir cette action en renvoyant directement \`a la
remarque \ref{remark-map-stalks}.

\begin{definition}
\label{definition-algebraic-geometric-point}
Soit $S$ un sch\'ema.
Soit $\overline{s}$ un point g\'eom\'etrique situ\'e au-dessus du point $s$ de $S$.
Notons $\kappa(s) \subset \kappa(s)^{sep} \subset \kappa(\overline{s})$
la cl\^oture s\'eparable de $\kappa(s)$ dans le corps alg\'ebriquement
clos $\kappa(\overline{s})$.
\begin{enumerate}
\item Dans cette situation, le {\it groupe de Galois absolu} de $\kappa(s)$
est $\text{Gal}(\kappa(s)^{sep}/\kappa(s))$. On le note parfois
$\text{Gal}_{\kappa(s)}$.
\item Le point g\'eom\'etrique $\overline{s}$ est dit
{\it alg\'ebrique} si $\kappa(s) \subset \kappa(\overline{s})$ est
une cl\^oture alg\'ebrique de $\kappa(s)$.
\end{enumerate}
\end{definition}

\begin{example}
\label{example-stupid}
Le point g\'eom\'etrique
$\Spec(\mathbf{C}) \to \Spec(\mathbf{Q})$
n'est pas alg\'ebrique.
\end{example}

\noindent
Soit $\kappa(s) \subset \kappa(s)^{sep} \subset \kappa(\overline{s})$
comme dans la d\'efinition. Puisque $\kappa(\overline{s})$ est alg\'ebriquement
clos, le morphisme
$$
\text{Aut}(\kappa(\overline{s})/\kappa(s))
\longrightarrow
\text{Gal}(\kappa(s)^{sep}/\kappa(s)) = \text{Gal}_{\kappa(s)}
$$
est surjectif. Supposons que $(U, \overline{u})$ soit un
voisinage \'etale de $\overline{s}$ et que $\overline{u}$ soit situ\'e au-dessus
du point $u$ de $U$. Puisque $U \to S$ est \'etale, l'extension de corps r\'esiduels
$\kappa(u)/\kappa(s)$ est finie s\'eparable.
Il s'ensuit ce qui suit
\begin{enumerate}
\item Si $\sigma \in \text{Aut}(\kappa(\overline{s})/\kappa(s)^{sep})$,
alors $\sigma$ agit trivialement sur $\mathcal{F}_{\overline{s}}$.
\item Plus pr\'ecis\'ement, l'action de
$\text{Aut}(\kappa(\overline{s})/\kappa(s))$
d\'etermine une action du groupe de Galois absolu
$\text{Gal}_{\kappa(s)}$ sur $\mathcal{F}_{\overline{s}}$, et est d\'etermin\'ee par elle.
\item Si $(U, \overline{u}, t)$ repr\'esente un \'el\'ement $\xi$ de
$\mathcal{F}_{\overline{s}}$, tout \'el\'ement de
$\text{Gal}(\kappa(s)^{sep}/K)$ agit trivialement, o\`u
$\kappa(s) \subset K \subset \kappa(s)^{sep}$ est l'image de
$\overline{u}^\sharp : \kappa(u) \to \kappa(\overline{s})$.
\end{enumerate}
Au total, nous voyons que $\mathcal{F}_{\overline{s}}$ devient un
$\text{Gal}_{\kappa(s)}$-ensemble (voir
Groupes fondamentaux, d\'efinition \ref{pione-definition-G-set-continuous}).
Nous pouvons donc consid\'erer le foncteur fibre comme un foncteur
$$
\Sh(S_\etale) \longrightarrow
\text{Gal}_{\kappa(s)}\textit{-Ensembles},
\quad
\mathcal{F} \longmapsto \mathcal{F}_{\overline{s}}
$$
et, dor\'enavant, c'est g\'en\'eralement ainsi que nous consid\'ererons le foncteur fibre.

\begin{theorem}
\label{theorem-equivalence-sheaves-point}
Soit $S = \Spec(K)$, o\`u $K$ est un corps.
Soit $\overline{s}$ un point g\'eom\'etrique de $S$.
Notons $G = \text{Gal}_{\kappa(s)}$ le groupe de Galois absolu.
Le passage aux fibres induit une \'equivalence de cat\'egories
$$
\Sh(S_\etale) \longrightarrow G\textit{-Ensembles},
\quad
\mathcal{F} \longmapsto \mathcal{F}_{\overline{s}}.
$$
\end{theorem}

\begin{proof}
Construisons un quasi-inverse de ce foncteur. Dans
Groupes fondamentaux, lemme \ref{pione-lemma-sheaves-point},
nous avons vu que, pour tout $G$-ensemble $M$, il existe un morphisme \'etale
$X \to \Spec(K)$
tel que $\Mor_K(\Spec(K^{sep}), X)$ soit
isomorphe \`a $M$ comme $G$-ensemble. Consid\'erons le faisceau
$\mathcal{F}$ sur $\Spec(K)_\etale$ d\'efini par
la r\`egle $U \mapsto \Mor_K(U, X)$. C'est un faisceau, car la topologie
\'etale est sous-canonique. On voit alors que
$\mathcal{F}_{\overline{s}} = \Mor_K(\Spec(K^{sep}), X) = M$
comme $G$-ensembles (d\'etails omis). Cela fournit un quasi-inverse du foncteur et
ach\`eve la d\'emonstration.
\end{proof}

\begin{remark}
\label{remark-every-sheaf-representable}
Une autre mani\`ere d'\'enoncer la conclusion du
th\'eor\`eme \ref{theorem-equivalence-sheaves-point} et de
Groupes fondamentaux, lemme \ref{pione-lemma-sheaves-point},
consiste \`a dire que tout faisceau sur $\Spec(K)_\etale$ est repr\'esentable
par un sch\'ema $X$ \'etale sur $\Spec(K)$.
Cela ne signifie pas que tout faisceau soit repr\'esentable au sens de
Sites, d\'efinition \ref{sites-definition-representable-sheaf}.
La raison est que, dans notre construction de $\Spec(K)_\etale$,
nous avons choisi un ensemble suffisamment grand de sch\'emas \'etales sur $\Spec(K)$,
tandis que les faisceaux sur $\Spec(K)_\etale$ forment une classe propre.
\end{remark}

\begin{lemma}
\label{lemma-global-sections-point}
Hypoth\`eses et notations comme dans le
th\'eor\`eme \ref{theorem-equivalence-sheaves-point}.
Il existe une bijection fonctorielle
$$
\Gamma(S, \mathcal{F}) = (\mathcal{F}_{\overline{s}})^G
$$
\end{lemma}

\begin{proof}
Nous pouvons le d\'emontrer par des arguments formels et \`a l'aide du
th\'eor\`eme \ref{theorem-equivalence-sheaves-point},
comme suit. Pour un faisceau $\mathcal{F}$ correspondant au
$G$-ensemble $M = \mathcal{F}_{\overline{s}}$, on a
\begin{eqnarray*}
\Gamma(S, \mathcal{F}) & = &
\Mor_{\Sh(S_\etale)}(h_{\Spec(K)},  \mathcal{F})
\\
& = & \Mor_{G\textit{-Ensembles}}(\{*\}, M) \\
& = & M^G
\end{eqnarray*}
Ici, la premi\`ere identification est expliqu\'ee dans
Sites, sections \ref{sites-section-presheaves} et
\ref{sites-section-representable-sheaves},
la deuxi\`eme r\'esulte du
th\'eor\`eme \ref{theorem-equivalence-sheaves-point}
et la troisi\`eme est claire. Nous donnerons aussi une preuve directe\footnote{Pour
les saint Thomas parmi nous.}.

\medskip\noindent
Supposons que $t \in \Gamma(S, \mathcal{F})$ soit une section globale.
Alors le triplet $(S, \overline{s}, t)$ d\'efinit un \'el\'ement de
$\mathcal{F}_{\overline{s}}$ qui est clairement invariant sous
l'action de $G$. R\'eciproquement, supposons que $(U, \overline{u}, t)$
d\'efinisse un \'el\'ement invariant de $\mathcal{F}_{\overline{s}}$.
Nous pouvons alors r\'etr\'ecir $U$ et supposer que $U = \Spec(L)$ pour une
extension de corps finie s\'eparable $L$ de $K$, voir la
proposition \ref{proposition-etale-morphisms}.
Dans ce cas, le morphisme $\mathcal{F}(U) \to \mathcal{F}_{\overline{s}}$
est injectif, car pour tout morphisme de voisinages \'etales
$(U', \overline{u}') \to (U, \overline{u})$, le morphisme de restriction
$\mathcal{F}(U) \to \mathcal{F}(U')$ est injectif puisque $U' \to U$
est un recouvrement de $S_\etale$.
Apr\`es avoir un peu agrandi $L$, nous pouvons supposer que $K \subset L$ est une extension
galoisienne finie. Nous utilisons alors l'\'egalit\'e
$$
\Spec(L) \times_{\Spec(K)} \Spec(L)
=
\coprod\nolimits_{\sigma \in \text{Gal}(L/K)} \Spec(L)
$$
o\`u les morphismes $\Spec(L) \to \Spec(L \otimes_K L)$
proviennent des morphismes d'anneaux $a \otimes b \mapsto a\sigma(b)$. Nous
voyons donc que l'invariance de $(U, \overline{u}, t)$
sous tout $G$ implique que $t \in \mathcal{F}(\Spec(L))$
a la m\^eme image dans
$\mathcal{F}(\Spec(L) \times_{\Spec(K)} \Spec(L))$
par restriction suivant chacune des deux projections (on utilise ici l'injectivit\'e mentionn\'ee
ci-dessus; d\'etails omis). La condition de faisceau pour $\mathcal{F}$
et le recouvrement \'etale $\{\Spec(L) \to \Spec(K)\}$ s'applique
donc, et nous concluons que $t$ provient d'une unique section de $\mathcal{F}$
sur $\Spec(K)$.
\end{proof}

\begin{remark}
\label{remark-stalk-pullback}
Soit $S$ un sch\'ema et soit $\overline{s} : \Spec(k) \to S$
un point g\'eom\'etrique de $S$. Par d\'efinition, cela signifie que $k$
est alg\'ebriquement clos. En particulier, le groupe de Galois absolu de $k$
est trivial. Par cons\'equent, d'apr\`es le
th\'eor\`eme \ref{theorem-equivalence-sheaves-point},
la cat\'egorie des faisceaux sur $\Spec(k)_\etale$ est \'equivalente
\`a la cat\'egorie des ensembles. L'\'equivalence est donn\'ee par la prise des
sections sur $\Spec(k)$. Cela nous fournit enfin une
autre d\'efinition du foncteur fibre. Plus pr\'ecis\'ement, le foncteur
$$
\Sh(S_\etale) \longrightarrow \textit{Ensembles}, \quad
\mathcal{F} \longmapsto \mathcal{F}_{\overline{s}}
$$
est isomorphe au foncteur
$$
\Sh(S_\etale)
\longrightarrow
\Sh(\Spec(k)_\etale) = \textit{Ensembles},
\quad
\mathcal{F} \longmapsto \overline{s}^*\mathcal{F}
$$
Pour le d\'emontrer rigoureusement, on peut utiliser
la partie (3) du lemme \ref{lemma-stalk-pullback}
avec $f = \overline{s}$. De plus, ceci \'etant dit, le cas g\'en\'eral de la
partie (3) du lemme \ref{lemma-stalk-pullback}
r\'esulte de la fonctorialit\'e des images inverses.
\end{remark}






\section{Cohomologie des groupes}
\label{section-group-cohomology}

\noindent
Dans la suite, lorsque nous \'ecrivons $H^i(G, M)$, nous entendons que $G$
est un groupe topologique, que $M$ est un $G$-module discret muni d'une
action continue de $G$, et que $H^i(G, -)$ est le $i$-i\`eme foncteur d\'eriv\'e
droit sur la cat\'egorie $\text{Mod}_G$ de ces $G$-modules, voir les
d\'efinitions \ref{definition-G-module-continuous} et
\ref{definition-galois-cohomology}. Cela comprend
le cas d'un groupe abstrait $G$, ce qui signifie simplement que l'on consid\`ere $G$
comme un groupe topologique muni de la topologie discr\`ete.

\medskip\noindent
Lorsque le module est muni d'une topologie non discr\`ete, nous emploierons la notation
$H^i_{cont}(G, M)$ pour d\'esigner les groupes de cohomologie continue
introduits dans \cite{Tate}, voir la
section \ref{section-continuous-group-cohomology}.

\begin{definition}
\label{definition-G-module-continuous}
Soit $G$ un groupe topologique.
\begin{enumerate}
\item Un {\it $G$-module}, parfois appel\'e {\it $G$-module discret},
est un groupe ab\'elien $M$ muni d'une action \`a gauche $a : G \times M \to M$
par homomorphismes de groupes, telle que $a$ soit continue lorsque $M$ est muni de la
topologie discr\`ete.
\item Un {\it morphisme de $G$-modules} $f : M \to N$ est un
homomorphisme $G$-\'equivariant de $M$ dans $N$.
\item La cat\'egorie des $G$-modules est not\'ee $\text{Mod}_G$.
\end{enumerate}
Soit $R$ un anneau.
\begin{enumerate}
\item Un {\it $R\text{-}G$-module} est un $R$-module $M$ muni
d'une action \`a gauche $a : G \times M \to M$ par applications $R$-lin\'eaires, telle que $a$
soit continue lorsque $M$ est muni de la topologie discr\`ete.
\item Un {\it morphisme de $R\text{-}G$-modules} $f : M \to N$ est une
application $R$-lin\'eaire $G$-\'equivariante de $M$ dans $N$.
\item La cat\'egorie des $R\text{-}G$-modules est not\'ee $\text{Mod}_{R, G}$.
\end{enumerate}
\end{definition}

\noindent
La continuit\'e de $a : G \times M \to M$ \'equivaut
\`a demander que le stabilisateur de tout $x \in M$ soit ouvert dans $G$.
Si $G$ est un groupe abstrait, cela correspond \`a la notion d'un
groupe ab\'elien muni d'une action de $G$, pourvu que l'on munisse $G$ de la
topologie discr\`ete. Observons que $\text{Mod}_{\mathbf{Z}, G} = \text{Mod}_G$.

\medskip\noindent
La cat\'egorie $\text{Mod}_G$ a assez d'injectifs, voir
Injectifs, lemme \ref{injectives-lemma-G-modules}.
Consid\'erons le foncteur exact \`a gauche
$$
\text{Mod}_G \longrightarrow \textit{Ab},
\quad
M \longmapsto M^G =
\{x \in M \mid g \cdot x = x\ \forall g \in G\}
$$
Nous notons parfois $M^G = H^0(G, M)$ et parfois
$M^G = \Gamma_G(M)$. Ce foncteur admet un foncteur d\'eriv\'e droit total
$R\Gamma_G(M)$ et, pour tout $i \geq 0$, un $i$-i\`eme foncteur d\'eriv\'e droit
$R^i\Gamma_G(M) = H^i(G, M)$.

\medskip\noindent
La m\^eme construction s'applique \`a
$H^0(G, -) : \text{Mod}_{R, G} \to \text{Mod}_R$. Nous verrons au
lemme \ref{lemma-modules-abelian} qu'elle co\"\i ncide avec la cohomologie
du $G$-module sous-jacent.

\begin{definition}
\label{definition-galois-cohomology}
Soit $G$ un groupe topologique. Soit $M$ un $G$-module discret
muni d'une action continue de $G$. Autrement dit, $M$ est un objet
de la cat\'egorie $\text{Mod}_G$ introduite dans la
d\'efinition \ref{definition-G-module-continuous}.
\begin{enumerate}
\item Les foncteurs d\'eriv\'es droits $H^i(G, M)$ de $H^0(G, M)$ sur la
cat\'egorie $\text{Mod}_G$ sont appel\'es les
{\it groupes de cohomologie continue de $G$ \`a coefficients dans $M$}.
\item Si $G$ est un groupe abstrait muni de la topologie discr\`ete,
les $H^i(G, M)$ sont appel\'es les {\it groupes de cohomologie de $G$ \`a coefficients dans $M$}.
\item Si $G$ est un groupe de Galois, les groupes $H^i(G, M)$ sont appel\'es
les {\it groupes de cohomologie galoisienne \`a coefficients dans $M$}.
\item Si $G$ est le groupe de Galois absolu d'un corps $K$, les groupes
$H^i(G, M)$ sont parfois appel\'es les {\it groupes de cohomologie galoisienne de $K$
\`a coefficients dans $M$}. Dans ce cas, nous \'ecrivons parfois
$H^i(K, M)$ au lieu de $H^i(G, M)$.
\end{enumerate}
\end{definition}

\begin{lemma}
\label{lemma-modules-abelian}
Soit $G$ un groupe topologique. Soit $R$ un anneau.
Pour tout $i \geq 0$, le diagramme
$$
\xymatrix{
\text{Mod}_{R, G} \ar[rr]_{H^i(G, -)} \ar[d] & &
\text{Mod}_R \ar[d] \\
\text{Mod}_G \ar[rr]^{H^i(G, -)} & &
\textit{Ab}
}
$$
dont les fl\`eches verticales sont les foncteurs d'oubli est commutatif.
\end{lemma}

\begin{proof}
Notons $F : \text{Mod}_{R, G} \to \text{Mod}_G$ le foncteur d'oubli.
Alors $F$ admet un adjoint \`a gauche $H : \text{Mod}_G \to \text{Mod}_{R, G}$
donn\'e par $H(M) = M \otimes_\mathbf{Z} R$. Observons que tout objet de
$\text{Mod}_G$ est un quotient d'une somme directe de modules de la forme
$\mathbf{Z}[G/U]$, o\`u $U \subset G$ est un sous-groupe ouvert.
Ici, $\mathbf{Z}[G/U]$ d\'esigne le $G$-module form\'e des
combinaisons $\mathbf{Z}$-lin\'eaires finies
des classes \`a droite modulo $U$ dans $G$, muni de l'action \`a gauche de $G$.
Ainsi, tout complexe born\'e sup\'erieurement dans $\text{Mod}_G$ est quasi-isomorphe
\`a un complexe born\'e sup\'erieurement dans $\text{Mod}_G$ dont les
termes sous-jacents sont des $\mathbf{Z}$-modules plats
(Cat\'egories d\'eriv\'ees, lemme \ref{derived-lemma-subcategory-left-resolution}).
Il est donc clair que $LH$ existe sur $D^-(\text{Mod}_G)$ et se calcule en
appliquant $H$ \`a tout complexe dont les termes sont des $\mathbf{Z}$-modules plats;
cela r\'esulte de
Cat\'egories d\'eriv\'ees, lemme \ref{derived-lemma-subcategory-left-acyclics} et
proposition \ref{derived-proposition-enough-acyclics}.
Nous concluons de Cat\'egories d\'eriv\'ees, lemme
\ref{derived-lemma-pre-derived-adjoint-functors}
que
$$
\text{Ext}^i(\mathbf{Z}, F(M)) = \text{Ext}^i(R, M)
$$
pour $M$ dans $\text{Mod}_{R, G}$.
Observons que $H^0(G, -) = \Hom(\mathbf{Z}, -)$ sur
$\text{Mod}_G$, o\`u $\mathbf{Z}$ d\'esigne le $G$-module
\`a action triviale. Par cons\'equent,
$H^i(G, -) = \text{Ext}^i(\mathbf{Z}, -)$ sur $\text{Mod}_G$.
De m\^eme, on a $H^i(G, -) = \text{Ext}^i(R, -)$ sur
$\text{Mod}_{R, G}$. En combinant ces faits, on obtient le lemme.
\end{proof}

\begin{lemma}
\label{lemma-ext-modules-hom}
Soit $G$ un groupe topologique. Soit $R$ un anneau.
Soient $M$, $N$ des $R\text{-}G$-modules. Si $M$ est projectif de type fini
comme $R$-module, alors
$\text{Ext}^i(M, N) = H^i(G, M^\vee \otimes_R N)$ (pour les notations,
voir la preuve).
\end{lemma}

\begin{proof}
Munissons le module $M^\vee = \Hom_R(M, R)$ de l'action contragr\'ediente
de $G$. Ainsi, $(g \cdot \lambda)(m) = \lambda(g^{-1} \cdot m)$
pour $g \in G$, $\lambda \in M^\vee$, $m \in M$. L'action de $G$ sur
$M^\vee \otimes_R N$ est l'action diagonale, c'est-\`a-dire qu'elle est donn\'ee par
$g \cdot (\lambda \otimes n) = g \cdot \lambda \otimes g \cdot n$.
Notons que, pour un troisi\`eme $R\text{-}G$-module $E$, on a
$\Hom(E, M^\vee \otimes_R N) = \Hom(M \otimes_R E, N)$.
En effet, cela est vrai au niveau des $R$-modules d'apr\`es
Alg\`ebre, lemmes \ref{algebra-lemma-hom-from-tensor-product} et
\ref{algebra-lemma-evaluation-map-iso-finite-projective}
et les d\'efinitions des actions de $G$ sont choisies de sorte que cela
reste vrai pour les $R\text{-}G$-modules. Il en r\'esulte que
$M^\vee \otimes_R N$ est un $R\text{-}G$-module injectif
si $N$ est un $R\text{-}G$-module injectif. Ainsi, si
$N \to N^\bullet$ est une r\'esolution injective, alors
$M^\vee \otimes_R N \to M^\vee \otimes_R N^\bullet$
est une r\'esolution injective. On a alors
$$
\Hom(M, N^\bullet) = \Hom(R, M^\vee \otimes_R N^\bullet) =
(M^\vee \otimes_R N^\bullet)^G
$$
Puisque le membre de gauche calcule $\text{Ext}^i(M, N)$ et celui de droite
calcule $H^i(G, M^\vee \otimes_R N)$, la preuve est termin\'ee.
\end{proof}

\begin{lemma}
\label{lemma-finite-dim-group-cohomology}
Soit $G$ un groupe topologique. Soit $k$ un corps.
Soit $V$ un $k\text{-}G$-module.
Si $G$ est topologiquement engendr\'e par un nombre fini d'\'el\'ements et si
$\dim_k(V) < \infty$, alors $\dim_k H^1(G, V) < \infty$.
\end{lemma}

\begin{proof}
Soient $g_1, \ldots, g_r \in G$ des \'el\'ements qui engendrent topologiquement $G$,
c'est-\`a-dire que le sous-groupe engendr\'e par $g_1, \ldots, g_r$ est dense.
D'apr\`es le lemme \ref{lemma-ext-modules-hom},
on voit que $H^1(G, V)$ est l'espace vectoriel sur $k$ des extensions
$$
0 \to V \to E \to k \to 0
$$
de $k\text{-}G$-modules. Choisissons $e \in E$ d'image $1 \in k$.
\'Ecrivons
$$
g_i \cdot e = v_i + e
$$
pour un $v_i \in V$. Cela est possible puisque $g_i \cdot 1 = 1$.
Nous affirmons que la liste des \'el\'ements $v_1, \ldots, v_r \in V$
d\'etermine la classe d'isomorphisme de l'extension $E$.
Une fois ceci d\'emontr\'e, le lemme en r\'esulte, car notre espace vectoriel
Ext est alors isomorphe \`a un sous-quotient de l'espace vectoriel sur $k$
$V^{\oplus r}$; quelques d\'etails sont omis.
Puisque $E$ est un objet de la cat\'egorie d\'efinie dans la
d\'efinition \ref{definition-G-module-continuous},
il existe un sous-groupe ouvert $U$ tel que
$u \cdot e = e$ pour tout $u \in U$.
Prenons maintenant un $g \in G$. Alors $gU$ contient un mot $w$ en
les \'el\'ements $g_1, \ldots, g_r$.
\'Ecrivons $gu = w$. Puisque l'\'el\'ement $w \cdot e$ est d\'etermin\'e par
$v_1, \ldots, v_r$, on voit que $g \cdot e = (gu) \cdot e = w \cdot e$
l'est lui aussi.
\end{proof}

\begin{lemma}
\label{lemma-profinite-group-cohomology-torsion}
Soit $G$ un groupe topologique profini.
Alors
\begin{enumerate}
\item $H^i(G, M)$ est de torsion pour $i > 0$ et tout $G$-module $M$, et
\item $H^i(G, M) = 0$ si $M$ est un $\mathbf{Q}$-espace vectoriel.
\end{enumerate}
\end{lemma}

\begin{proof}
Preuve de (1). Par d\'ecalage de dimension, il suffit
de montrer que $H^1(G, M)$ est de torsion pour tout $G$-module $M$.
Choisissons une suite exacte $0 \to M \to I \to N \to 0$, o\`u $I$
est un objet injectif de la cat\'egorie des $G$-modules.
Tout \'el\'ement de $H^1(G, M)$ est alors l'image d'un \'el\'ement
$y \in N^G$. Choisissons $x \in I$ d'image $y$.
Le stabilisateur $U \subset G$ de $x$ est ouvert, donc
d'indice fini $r$. Soit $g_1, \ldots, g_r \in G$ un syst\`eme
de repr\'esentants de $G/U$. Alors $\sum g_i(x)$ est un \'el\'ement
invariant de $I$ dont l'image est $ry$. Ainsi, $r$ annule l'\'el\'ement
de $H^1(G, M)$ dont nous \'etions partis. La partie (2) en r\'esulte, puisque
$H^i(G, M)$ est alors \`a la fois un $\mathbf{Q}$-espace vectoriel et de torsion.
\end{proof}





\section{Cohomologie continue de Tate}
\label{section-continuous-group-cohomology}

\noindent
La cohomologie continue de Tate (\cite{Tate}) est d\'efinie par le complexe des
cocha\^ines inhomog\`enes continues. On peut la d\'efinir lorsque $M$ est un
groupe ab\'elien topologique quelconque muni d'une action continue de $G$.
Plus pr\'ecis\'ement, consid\'erons le complexe
$$
C^\bullet_{cont}(G, M) :
M \to \text{Maps}_{cont}(G, M) \to
\text{Maps}_{cont}(G \times G, M) \to \ldots
$$
o\`u l'application de bord est d\'efinie, pour $n \geq 1$, par la formule
\begin{align*}
\text{d}(f)(g_1, \ldots, g_{n + 1})
& = g_1(f(g_2, \ldots, g_{n + 1})) \\
&
+ \sum\nolimits_{j = 1, \ldots, n}
(-1)^jf(g_1, \ldots, g_jg_{j + 1}, \ldots, g_{n + 1}) \\
&
+ (-1)^{n + 1}f(g_1, \ldots, g_n)
\end{align*}
et, pour $n = 0$, envoie $m \in M$ sur l'application $g \mapsto g(m) - m$. Posons
$$
H^i_{cont}(G, M) = H^i(C^\bullet_{cont}(G, M))
$$
Comme les termes du complexe font intervenir les applications continues de $G$ et
de ses puissances cart\'esiennes dans le module topologique $M$, il n'est pas clair
qu'une suite exacte courte de modules topologiques soit ainsi transform\'ee en
une suite exacte longue de cohomologie. Une autre difficult\'e tient \`a ce que la cat\'egorie
des groupes ab\'eliens topologiques n'est pas une cat\'egorie ab\'elienne !

\medskip\noindent
Cependant, une suite exacte courte de $G$-modules discrets fournit bien
une suite exacte courte de complexes de cocha\^ines continues
et donc une suite exacte longue de groupes de cohomologie
continue $H^i_{cont}(G, -)$.
Par cons\'equent, dans la cat\'egorie $\text{Mod}_G$ de la
d\'efinition \ref{definition-G-module-continuous}, les foncteurs
$H^i_{cont}(G, M)$ forment un $\delta$-foncteur cohomologique au sens de
Homologie, section \ref{homology-section-cohomological-delta-functor}.
Puisque la cohomologie $H^i(G, M)$
de la d\'efinition \ref{definition-galois-cohomology}
est un $\delta$-foncteur universel
(Cat\'egories d\'eriv\'ees, lemme \ref{derived-lemma-right-derived-delta-functor}),
on obtient des morphismes canoniques
$$
H^i(G, M) \longrightarrow H^i_{cont}(G, M)
$$
pour $M \in \text{Mod}_G$. On sait que ces morphismes sont des
isomorphismes lorsque $G$ est un groupe abstrait (c'est-\`a-dire muni de
la topologie discr\`ete) ou lorsque $G$ est un groupe profini
(ins\'erer ici une r\'ef\'erence ult\'erieure).
Si vous connaissez un exemple o\`u ce morphisme n'est pas un isomorphisme
pour un groupe topologique $G$ et $M \in \Ob(\text{Mod}_G)$,
veuillez l'envoyer par courriel \`a
\href{mailto:stacks.project@gmail.com}{stacks.project@gmail.com}.




\section{Cohomologie d'un point}
\label{section-cohomology-point}

\noindent
Comme cons\'equence de la discussion des sections pr\'ec\'edentes,
nous obtenons l'\'equivalence entre la cohomologie \'etale du spectre d'un
corps et la cohomologie galoisienne.

\begin{lemma}
\label{lemma-equivalence-abelian-sheaves-point}
Soit $S = \Spec(K)$, o\`u $K$ est un corps.
Soit $\overline{s}$ un point g\'eom\'etrique de $S$.
Notons $G = \text{Gal}_{\kappa(s)}$ le groupe de Galois absolu.
Le foncteur fibre induit une \'equivalence de cat\'egories
$$
\textit{Ab}(S_\etale) \longrightarrow \text{Mod}_G,
\quad
\mathcal{F} \longmapsto \mathcal{F}_{\overline{s}}.
$$
\end{lemma}

\begin{proof}
Dans le
th\'eor\`eme \ref{theorem-equivalence-sheaves-point},
nous avons vu l'\'equivalence entre les faisceaux d'ensembles et les $G$-ensembles.
Le pr\'esent lemme s'en d\'eduit formellement, car un faisceau ab\'elien n'est autre
qu'un faisceau d'ensembles muni d'une loi de groupe commutative, et un $G$-module
n'est autre qu'un $G$-ensemble muni d'une loi de groupe commutative.
\end{proof}

\begin{lemma}
\label{lemma-compare-cohomology-point}
Les notations et hypoth\`eses sont celles du
lemme \ref{lemma-equivalence-abelian-sheaves-point}.
Soit $\mathcal{F}$ un faisceau ab\'elien sur $\Spec(K)_\etale$
qui correspond au $G$-module $M$.
Alors
\begin{enumerate}
\item dans $D(\textit{Ab})$, on a un isomorphisme canonique
$R\Gamma(S, \mathcal{F}) = R\Gamma_G(M)$,
\item $H_\etale^0(S, \mathcal{F}) = M^G$, et
\item $H_\etale^q(S, \mathcal{F}) = H^q(G, M)$.
\end{enumerate}
\end{lemma}

\begin{proof}
Combiner le
lemme \ref{lemma-equivalence-abelian-sheaves-point}
avec le
lemme \ref{lemma-global-sections-point}.
\end{proof}

\begin{example}
\label{example-sheaves-point}
Faisceaux sur $\Spec(K)_\etale$.
Soit $G = \text{Gal}(K^{sep}/K)$ le groupe de Galois absolu de $K$.
\begin{enumerate}
\item Le faisceau constant $\underline{\mathbf{Z}/n\mathbf{Z}}$ correspond au
module $\mathbf{Z}/n\mathbf{Z}$ muni de l'action triviale de $G$,
\item le faisceau $\mathbf{G}_m|_{\Spec(K)_\etale}$ correspond \`a
$(K^{sep})^*$ muni de son action de $G$,
\item le faisceau $\mathbf{G}_a|_{\Spec(K^{sep})}$ correspond \`a
$(K^{sep}, +)$ muni de son action de $G$, et
\item le faisceau $\mu_n|_{\Spec(K^{sep})}$ correspond \`a
$\mu_n(K^{sep})$ muni de son action de $G$.
\end{enumerate}
D'apr\`es la
remarque \ref{remark-special-case-fpqc-cohomology-quasi-coherent}
et le
th\'eor\`eme \ref{theorem-picard-group},
on a les identifications suivantes pour les groupes de cohomologie :
\begin{align*}
H_\etale^0(S_\etale, \mathbf{G}_m) & =
\Gamma(S, \mathcal{O}_S^*) \\
H_\etale^1(S_\etale, \mathbf{G}_m) & =
H_{Zar}^1(S, \mathcal{O}_S^*) = \Pic(S) \\
H_\etale^i(S_\etale, \mathbf{G}_a) & =
H_{Zar}^i(S, \mathcal{O}_S)
\end{align*}
De plus, pour tout faisceau quasi-coh\'erent $\mathcal{F}$ sur $S_\etale$, on a
$$
H^i(S_\etale, \mathcal{F}) = H_{Zar}^i(S, \mathcal{F}),
$$
voir le
th\'eor\`eme \ref{theorem-zariski-fpqc-quasi-coherent}.
En particulier, on obtient la suite d'\'egalit\'es suivante
$$
0 =
\Pic(\Spec(K)) =
H_\etale^1(\Spec(K)_\etale, \mathbf{G}_m) =
H^1(G, (K^{sep})^*)
$$
qui n'est autre que le th\'eor\`eme 90 de Hilbert. De m\^eme, pour $i \geq 1$,
$$
0 = H^i(\Spec(K), \mathcal{O})
= H_\etale^i(\Spec(K)_\etale, \mathbf{G}_a)
= H^i(G, K^{sep})
$$
o\`u $K^{sep}$ est ici consid\'er\'e comme module galoisien, l'addition \'etant
sa loi de groupe. De la sorte, nous pouvons consid\'erer le travail accompli jusqu'ici comme
une mani\`ere compliqu\'ee de calculer des groupes de cohomologie galoisienne.
\end{example}

\noindent
Le r\'esultat suivant est une curiosit\'e et devrait \^etre omis lors d'une
premi\`ere lecture.

\begin{lemma}
\label{lemma-all-modules-quasi-coherent}
Soit $R$ un anneau local de dimension $0$. Soit $S = \Spec(R)$.
Alors tout $\mathcal{O}_S$-module sur $S_\etale$ est quasi-coh\'erent.
\end{lemma}

\begin{proof}
Soit $\mathcal{F}$ un $\mathcal{O}_S$-module sur $S_\etale$.
Nous devons montrer que $\mathcal{F}$ est d\'etermin\'e par le
$R$-module $M = \Gamma(S, \mathcal{F})$.
Plus pr\'ecis\'ement, si $\pi : X \to S$ est \'etale, nous devons
montrer que $\Gamma(X, \mathcal{F}) = \Gamma(X, \pi^*\widetilde{M})$.

\medskip\noindent
Soit $\mathfrak m \subset R$ l'id\'eal maximal et soit
$\kappa$ le corps r\'esiduel. D'apr\`es
Alg\`ebre, lemme \ref{algebra-lemma-local-dimension-zero-henselian},
l'anneau local $R$ est hens\'elien. Si $X \to S$ est \'etale,
alors l'espace topologique sous-jacent \`a $X$ est discret
d'apr\`es Morphismes, lemme \ref{morphisms-lemma-etale-over-field},
et $X$ est donc une r\'eunion disjointe de sch\'emas affines
ayant chacun un seul point. De plus, si $X = \Spec(A)$ est affine et
n'a qu'un point, alors $A$ est une $R$-alg\`ebre finie \'etale d'apr\`es
Alg\`ebre, lemme \ref{algebra-lemma-mop-up}.
Nous devons montrer que $\Gamma(X, \mathcal{F}) = M \otimes_R A$
dans ce cas.

\medskip\noindent
Le foncteur $A \mapsto A/\mathfrak m A$ d\'efinit une \'equivalence entre
la cat\'egorie des $R$-alg\`ebres finies \'etales
et celle des $\kappa$-alg\`ebres finies s\'eparables, d'apr\`es
Alg\`ebre, lemme \ref{algebra-lemma-henselian-cat-finite-etale}.
Consid\'erons d'abord le cas o\`u $A/\mathfrak m A$
est une extension galoisienne de $\kappa$, de groupe de Galois $G$.
Pour chaque $\sigma \in G$, notons $\sigma : A \to A$
l'automorphisme correspondant de $A$ au-dessus de $R$.
Posons $N = \Gamma(X, \mathcal{F})$.
Alors $\Spec(\sigma) : X \to X$ est un automorphisme au-dessus de $S$
et l'image inverse par celui-ci d\'efinit donc une application $\sigma : N \to N$
qui est $\sigma$-lin\'eaire : $\sigma(an) = \sigma(a) \sigma(n)$
pour $a \in A$ et $n \in N$.
Nous allons appliquer la descente galoisienne au module quasi-coh\'erent
$\widetilde{N}$ sur $X$, muni des isomorphismes
provenant de l'action de $\sigma$ sur $N$. Voir Descente, lemme
\ref{descent-lemma-galois-descent-more-general}.
Ce lemme fournit un isomorphisme $N = N^G \otimes_R A$.
D'autre part, il est clair que $N^G = M$ par la propri\'et\'e de faisceau
de $\mathcal{F}$. L'isomorphisme voulu est donc \'etabli.

\medskip\noindent
Le cas g\'en\'eral (o\`u $A$ est une $R$-alg\`ebre locale finie \'etale)
se d\'eduit du cas galoisien comme suit. Choisissons $A \to B$
finie \'etale, avec $B$ local et son corps r\'esiduel
galoisien sur $\kappa$. Posons $G = \text{Aut}(B/R) = \text{Gal}(\kappa_B/\kappa)$.
Soit $H \subset G$ le groupe de Galois correspondant \`a
l'extension galoisienne $\kappa_B/\kappa_A$. Alors, comme ci-dessus, on
montre que $\Gamma(X, \mathcal{F}) = \Gamma(\Spec(B), \mathcal{F})^H$.
En appliquant deux fois le r\'esultat relatif aux extensions galoisiennes, on obtient
$$
\Gamma(X, \mathcal{F}) = (M \otimes_R B)^H = M \otimes_R A
$$
comme voulu.
\end{proof}








\section{Cohomologie des courbes}
\label{section-cohomology-curves}

\noindent
La t\^ache suivante consiste \`a calculer la cohomologie \'etale d'une courbe lisse
sur un corps alg\'ebriquement clos, \`a coefficients de torsion, et en
particulier \`a montrer qu'elle s'annule en degr\'e au moins $3$. Pour cela, nous
calculerons la cohomologie au point g\'en\'erique, ce qui
revient \`a calculer de la cohomologie galoisienne.





\section{Groupes de Brauer}
\label{section-brauer-groups}

\noindent
Les groupes de Brauer des corps sont d\'efinis au moyen d'alg\`ebres centrales simples de dimension finie.
Dans cette section, nous rappelons les faits pertinents sur les groupes de Brauer, dont la plupart
sont expos\'es dans le chapitre
Groupes de Brauer, section \ref{brauer-section-introduction}.
Pour d'autres r\'ef\'erences, voir \cite{SerreCorpsLocaux},
\cite{SerreGaloisCohomology} ou \cite{Weil}.

\begin{theorem}
\label{theorem-central-simple-algebra}
Soit $K$ un corps. Pour une $K$-alg\`ebre $A$ unitaire et associative (non n\'ecessairement commutative),
les conditions suivantes sont \'equivalentes
\begin{enumerate}
\item $A$ est une $K$-alg\`ebre centrale simple de dimension finie,
\item $A$ est un $K$-espace vectoriel de dimension finie, $K$ est le centre de $A$,
et $A$ n'a pas d'id\'eal bilat\`ere non trivial,
\item il existe $d \geq 1$ tel que
$A \otimes_K \bar K \cong \text{Mat}(d \times d, \bar K)$,
\item il existe $d \geq 1$ tel que
$A \otimes_K K^{sep} \cong \text{Mat}(d \times d, K^{sep})$,
\item il existe $d \geq 1$ et une extension galoisienne finie $K'/K$
tels que
$A \otimes_K K' \cong \text{Mat}(d \times d, K')$,
\item il existe $n \geq 1$ et un corps gauche central $D$ de dimension finie
sur $K$ tel que $A \cong \text{Mat}(n \times n, D)$.
\end{enumerate}
L'entier $d$ est appel\'e le {\it degr\'e} de $A$.
\end{theorem}

\begin{proof}
Il s'agit d'une copie de
Groupes de Brauer, lemme \ref{brauer-lemma-finite-central-simple-algebra}.
\end{proof}

\begin{lemma}
\label{lemma-brauer-inverse}
Soit $A$ une alg\`ebre centrale simple de dimension finie sur $K$. Alors
$$
\begin{matrix}
A \otimes_K A^{opp} & \longrightarrow & \text{End}_K(A) \\
\ a \otimes a' & \longmapsto & (x \mapsto a x a')
\end{matrix}
$$
est un isomorphisme de $K$-alg\`ebres.
\end{lemma}

\begin{proof}
Voir
Groupes de Brauer, lemme \ref{brauer-lemma-inverse}.
\end{proof}

\begin{definition}
\label{definition-brauer-equivalent}
Deux alg\`ebres centrales simples de dimension finie $A_1$ et $A_2$ sur $K$ sont dites
{\it semblables}, ou {\it \'equivalentes}, s'il existe $m, n \geq 1$
tels que $\text{Mat}(n \times n, A_1)
\cong \text{Mat}(m \times m, A_2)$. On \'ecrit $A_1 \sim A_2$.
\end{definition}

\noindent
D'apr\`es Groupes de Brauer, lemme \ref{brauer-lemma-similar}, c'est une
relation d'\'equivalence.

\begin{definition}
\label{definition-brauer-group}
Soit $K$ un corps. Le {\it groupe de Brauer} de $K$ est l'ensemble $\text{Br} (K)$
des classes de similitude d'alg\`ebres centrales simples de dimension finie sur $K$, muni de
la loi de groupe induite par le produit tensoriel (sur $K$). La classe de $A$ dans
$\text{Br}(K)$ est not\'ee $[A]$. L'\'el\'ement neutre est
$[K] = [\text{Mat}(d \times d, K)]$ pour tout $d \geq 1$.
\end{definition}

\noindent
Le lemme pr\'ec\'edent entra\^ine l'existence des inverses et l'\'egalit\'e $-[A] = [A^{opp}]$.
Le groupe de Brauer d'un corps est toujours de torsion.
En fait, nous verrons que $[A]$ est d'ordre divisant $\deg(A)$
pour toute alg\`ebre centrale simple de dimension finie $A$ (voir le
lemme \ref{lemma-annihilated-by-degree}).
En g\'en\'eral, le groupe de Brauer n'est pas de type fini; par exemple,
le groupe de Brauer d'un corps local non archim\'edien est $\mathbf{Q}/\mathbf{Z}$.
Le groupe de Brauer de $\mathbf{C}(x, y)$ n'est pas d\'enombrable.

\begin{lemma}
\label{lemma-central-simple-algebra-pgln}
\begin{slogan}
Les alg\`ebres centrales simples sont class\'ees par la cohomologie galoisienne de PGL.
\end{slogan}
Soient $K$ un corps et $K^{sep}$ une cl\^oture s\'eparable.
Alors l'ensemble des classes d'isomorphisme d'alg\`ebres centrales simples de degr\'e
$d$ sur $K$ est en bijection avec la cohomologie non ab\'elienne
$H^1(\text{Gal}(K^{sep}/K), \text{PGL}_d(K^{sep}))$.
\end{lemma}

\begin{proof}[Esquisse de preuve.]
Le th\'eor\`eme de Skolem-Noether (voir
Groupes de Brauer, th\'eor\`eme \ref{brauer-theorem-skolem-noether})
entra\^ine que, pour tout corps $L$, le groupe
$\text{Aut}_{L\text{-Algebras}}(\text{Mat}_d(L))$
est \`egal \`a $\text{PGL}_d(L)$. D'apr\`es le
th\'eor\`eme \ref{theorem-central-simple-algebra}, on voit que les
alg\`ebres centrales simples de degr\'e $d$ correspondent
aux formes de la $K$-alg\`ebre $\text{Mat}_d(K)$.
En combinant ces faits, on voit que les classes d'isomorphisme d'alg\`ebres
centrales simples de degr\'e $d$ correspondent aux \'el\'ements de
$H^1(\text{Gal}(K^{sep}/K), \text{PGL}_d(K^{sep}))$.
Pour plus de d\'etails sur les formes tordues, voir par exemple
\cite{SilvermanEllipticCurves}.
\end{proof}

\noindent
Si $A$ est une alg\`ebre centrale simple de dimension finie et de degr\'e $d$ sur un corps $K$,
nous notons $\xi_A$ la classe de cohomologie correspondante dans
$H^1(\text{Gal}(K^{sep}/K), \text{PGL}_d(K^{sep}))$.
Consid\'erons la suite exacte courte
$$
1 \to (K^{sep})^* \to \text{GL}_d(K^{sep}) \to \text{PGL}_d(K^{sep}) \to 1,
$$
qui donne naissance \`a une suite exacte longue de cohomologie (jusqu'au degr\'e $2$), avec
pour homomorphisme bord
$$
\delta_d :
H ^1(\text{Gal}(K^{sep}/K), \text{PGL}_d(K^{sep}))
\longrightarrow
H^2(\text{Gal}(K^{sep}/K), (K^{sep})^*).
$$
Plus pr\'ecis\'ement, il est d\'efini comme suit : si $\xi$ est une classe de cohomologie
repr\'esent\'ee par le $1$-cocycle $(g_\sigma)$, alors $\delta_d(\xi)$ est la
classe du $2$-cocycle
\begin{equation}
\label{equation-two-cocycle}
(\sigma, \tau)
\longmapsto
\tilde g_\sigma^{-1} \tilde g_{\sigma \tau} \sigma(\tilde g_\tau^{-1})
\in (K^{sep})^*
\end{equation}
o\`u $\tilde g_\sigma \in \text{GL}_d(K^{sep})$ est un rel\`evement de $g_\sigma$.
Cela permet d'expliciter l'application
$$
\delta : \text{Br}(K) \longrightarrow H^2(\text{Gal}(K^{sep}/K), (K^{sep})^*),
\quad
[A] \longmapsto \delta_{\deg A} (\xi_A)
$$
comme suit. Supposons que $A$ soit de degr\'e $d$ sur $K$. Choisissons un isomorphisme
$\varphi : \text{Mat}_d(K^{sep}) \to A \otimes_K K^{sep}$. Pour
$\sigma \in \text{Gal}(K^{sep}/K)$, choisissons un \'el\'ement
$\tilde g_\sigma \in \text{GL}_d(K^{sep})$ tel que
$\varphi^{-1} \circ \sigma(\varphi)$ soit \'egal \`a l'application
$x \mapsto \tilde g_\sigma x \tilde g_\sigma^{-1}$. La classe dans $H^2$
est d\'efinie par le $2$-cocycle (\ref{equation-two-cocycle}).

\begin{theorem}
\label{theorem-brauer-delta}
Soit $K$ un corps de cl\^oture s\'eparable $K^{sep}$. L'application
$\delta : \text{Br}(K) \to H^2(\text{Gal}(K^{sep}/K), (K^{sep})^*)$
d\'efinie ci-dessus est un isomorphisme de groupes.
\end{theorem}

\begin{proof}[Esquisse de preuve]
Pour montrer que $\delta$ d\'efinit un homomorphisme de groupes, c'est-\`a-dire que
$\delta(A \otimes_K B) = \delta(A) + \delta(B)$, on calcule
directement sur les cocycles.

\medskip\noindent
Injectivit\'e de $\delta$. Dans le cas ab\'elien ($d = 1$), on a
l'identification
$$
H^1(\text{Gal}(K^{sep}/K), \text{GL}_d(K^{sep})) =
H_\etale^1(\Spec(K), \text{GL}_d(\mathcal{O}))
$$
dont le second membre est trivial par descente fpqc. Si cela \'etait vrai dans le
cas non ab\'elien, l'injectivit\'e de $\delta$ en r\'esulterait aussit\^ot. (Voir
\cite{SGA4.5}.) En fait, pour la d\'emontrer, on peut r\'einterpr\'eter $\delta([A])$ comme
l'obstruction \`a l'existence d'un $K$-espace vectoriel $V$ muni d'une structure de $A$-module
\`a gauche et tel que $\dim_K V = \deg A$. Lorsque $V$ existe, on
a $A \cong \text{End}_K(V)$.

\medskip\noindent
Pour la surjectivit\'e, choisissons une
classe de cohomologie $\xi \in H^2(\text{Gal}(K^{sep}/K), (K^{sep})^*)$ ;
il existe alors une extension galoisienne finie $K^{sep}/K'/K$
telle que $\xi$ soit l'image d'une classe
$\xi' \in H^2(\text{Gal}(K'|K), (K')^*)$. On construit alors
explicitement une alg\`ebre centrale simple sur $K$ \`a partir des donn\'ees $K', \xi'$.
\end{proof}










\section{Le groupe de Brauer d'un sch\'ema}
\label{section-brauer-scheme}

\noindent
Soit $S$ un sch\'ema. Une $\mathcal{O}_S$-alg\`ebre
$\mathcal{A}$ est dite {\it d'Azumaya} si elle est, localement pour la topologie \'etale, une alg\`ebre
de matrices, c'est-\`a-dire s'il existe un recouvrement \'etale
$\mathcal{U} = \{ \varphi_i : U_i \to S\}_{i \in I}$ tel que
$\varphi_i^*\mathcal{A} \cong \text{Mat}_{d_i}(\mathcal{O}_{U_i})$
pour un $d_i \geq 1$. Deux telles alg\`ebres
$\mathcal{A}$ et $\mathcal{B}$ sont dites {\it \'equivalentes} s'il existe
des $\mathcal{O}_S$-modules localement libres de type fini $\mathcal{F}$ et $\mathcal{G}$
de rang strictement positif en tout $s \in S$, tels que
$$
\mathcal{A} \otimes_{\mathcal{O}_S}
\SheafHom_{\mathcal{O}_S}(\mathcal{F}, \mathcal{F})
\cong
\mathcal{B} \otimes_{\mathcal{O}_S}
\SheafHom_{\mathcal{O}_S}(\mathcal{G}, \mathcal{G})
$$
comme $\mathcal{O}_S$-alg\`ebres. Le {\it groupe de Brauer} de
$S$ est l'ensemble $\text{Br}(S)$ des classes d'\'equivalence d'alg\`ebres
d'Azumaya sur $\mathcal{O}_S$, muni de l'op\'eration induite par le produit tensoriel (sur
$\mathcal{O}_S$).

\begin{lemma}
\label{lemma-end-unique-up-to-invertible}
Soit $S$ un sch\'ema. Soient $\mathcal{F}$ et $\mathcal{G}$ des faisceaux de
$\mathcal{O}_S$-modules localement libres de type fini et de rang strictement positif. S'il
existe un isomorphisme
$\SheafHom_{\mathcal{O}_S}(\mathcal{F}, \mathcal{F}) \cong
\SheafHom_{\mathcal{O}_S}(\mathcal{G}, \mathcal{G})$
de $\mathcal{O}_S$-alg\`ebres, alors il existe un faisceau inversible
$\mathcal{L}$ sur $S$ tel que
$\mathcal{F} \otimes_{\mathcal{O}_S} \mathcal{L} \cong \mathcal{G}$
et tel que cet isomorphisme induise l'isomorphisme donn\'e des
alg\`ebres d'endomorphismes.
\end{lemma}

\begin{proof}
Fixons un isomorphisme
$\SheafHom_{\mathcal{O}_S}(\mathcal{F}, \mathcal{F}) \to
\SheafHom_{\mathcal{O}_S}(\mathcal{G}, \mathcal{G})$.
Consid\'erons le faisceau $\mathcal{L} \subset \SheafHom(\mathcal{F}, \mathcal{G})$
engendr\'e, comme $\mathcal{O}_S$-module, par les isomorphismes locaux
$\varphi : \mathcal{F} \to \mathcal{G}$ tels que la conjugaison par
$\varphi$ soit l'isomorphisme donn\'e des alg\`ebres d'endomorphismes.
Un calcul local (qui se ram\`ene au cas o\`u $\mathcal{F}$ et $\mathcal{G}$
sont libres de type fini et $S$ est affine) montre que $\mathcal{L}$ est inversible.
Un autre calcul local montre que le morphisme d'\'evaluation
$$
\mathcal{F} \otimes_{\mathcal{O}_S} \mathcal{L} \longrightarrow \mathcal{G}
$$
est un isomorphisme.
\end{proof}

\noindent
L'argument donn\'e dans la preuve du lemme suivant se trouve dans
\cite{Saltman-torsion}.

\begin{lemma}
\label{lemma-annihilated-by-degree}
\begin{reference}
Argument tir\'e de \cite{Saltman-torsion}.
\end{reference}
Soit $S$ un sch\'ema. Soit $\mathcal{A}$ une alg\`ebre d'Azumaya qui est
localement libre de rang $d^2$ sur $S$. Alors la classe
de $\mathcal{A}$ dans le groupe de Brauer de $S$ est annul\'ee par $d$.
\end{lemma}

\begin{proof}
Choisissons un recouvrement \'etale $\{U_i \to S\}$ et des isomorphismes
$\mathcal{A}|_{U_i} \to \SheafHom(\mathcal{F}_i, \mathcal{F}_i)$
pour des $\mathcal{O}_{U_i}$-modules localement libres $\mathcal{F}_i$
de rang $d$. (On peut supposer $\mathcal{F}_i$ libre.) Consid\'erons la
compos\'ee
$$
p_i : \mathcal{F}_i^{\otimes d} \to
\wedge^d(\mathcal{F}_i) \to \mathcal{F}_i^{\otimes d}
$$
La premi\`ere fl\`eche est la projection usuelle et la seconde est
l'isomorphisme de la puissance ext\'erieure maximale de $\mathcal{F}_i$ sur
le sous-module des sections de $\mathcal{F}_i^{\otimes d}$ qui se transforment
suivant le caract\`ere signature sous l'action du groupe sym\'etrique
sur $d$ lettres. Alors $p_i^2 = d! p_i$ et le rang de $p_i$ vaut $1$.
En utilisant l'isomorphisme donn\'e
$\mathcal{A}|_{U_i} \to \SheafHom(\mathcal{F}_i, \mathcal{F}_i)$
et l'isomorphisme canonique
$$
\SheafHom(\mathcal{F}_i, \mathcal{F}_i)^{\otimes d} =
\SheafHom(\mathcal{F}_i^{\otimes d}, \mathcal{F}_i^{\otimes d})
$$
nous pouvons consid\'erer $p_i$ comme une section de $\mathcal{A}^{\otimes d}$
sur $U_i$. Nous affirmons que $p_i|_{U_i \times_S U_j} = p_j|_{U_i \times_S U_j}$
comme sections de $\mathcal{A}^{\otimes d}$. En effet, en appliquant le
lemme \ref{lemma-end-unique-up-to-invertible},
nous obtenons un faisceau inversible $\mathcal{L}_{ij}$ et un isomorphisme canonique
$$
\mathcal{F}_i|_{U_i \times_S U_j} \otimes \mathcal{L}_{ij}
\longrightarrow
\mathcal{F}_j|_{U_i \times_S U_j}.
$$
Cet isomorphisme montre que $p_i$ s'envoie sur $p_j$.
Comme $\mathcal{A}^{\otimes d}$ est un faisceau sur $S_\etale$
(proposition \ref{proposition-quasi-coherent-sheaf-fpqc}), on obtient une
section globale canonique $p \in \Gamma(S, \mathcal{A}^{\otimes d})$. Un calcul local
montre que
$$
\mathcal{H} =
\Im(\mathcal{A}^{\otimes d} \to \mathcal{A}^{\otimes d}, f \mapsto fp)
$$
est un module localement libre de rang $d^d$ et que la multiplication (\`a gauche)
par $\mathcal{A}^{\otimes d}$ induit un isomorphisme
$\mathcal{A}^{\otimes d} \to \SheafHom(\mathcal{H}, \mathcal{H})$.
Autrement dit, $\mathcal{A}^{\otimes d}$ est l'\'el\'ement neutre
du groupe de Brauer de $S$, comme voulu.
\end{proof}

\noindent
Dans ce cadre, l'analogue de l'isomorphisme $\delta$ du
th\'eor\`eme \ref{theorem-brauer-delta}
est une application
$$
\delta_S: \text{Br}(S) \to H_\etale^2(S, \mathbf{G}_m).
$$
L'application $\delta_S$ est injective. Si $S$ est quasi-compact ou
connexe, alors $\text{Br}(S)$ est un groupe de torsion; dans ce cas, l'image
de $\delta_S$ est contenue dans le {\it groupe de Brauer cohomologique} de $S$
$$
\text{Br}'(S) := H_\etale^2(S, \mathbf{G}_m)_\text{torsion}.
$$
Ainsi, si $S$ est quasi-compact ou connexe, on a une inclusion $\text{Br}(S)
\subset \text{Br}'(S)$. Ce n'est pas toujours une \'egalit\'e : il existe une
surface singuli\`ere non s\'epar\'ee $S$ pour laquelle $\text{Br}(S) \subset
\text{Br}'(S)$ est une inclusion stricte. Si $S$ est quasi-projectif, alors
$\text{Br}(S) = \text{Br}'(S)$. Cependant, on ignore si cela reste vrai pour
une vari\'et\'e lisse propre sur $\mathbf{C}$, par exemple.




\section{La suite d'Artin-Schreier}
\label{section-artin-schreier}

\noindent
Soit $p$ un nombre premier. Soit $S$ un sch\'ema de caract\'eristique $p$.
La suite {\it d'Artin-Schreier} est la suite exacte courte
$$
0 \longrightarrow \underline{\mathbf{Z}/p\mathbf{Z}}_S \longrightarrow
\mathbf{G}_{a, S} \xrightarrow{F-1} \mathbf{G}_{a, S} \longrightarrow 0
$$
o\`u $F - 1$ est l'application $x \mapsto x^p - x$.

\begin{lemma}
\label{lemma-vanishing-affine-char-p-p}
Soit $p$ un nombre premier. Soit $S$ un sch\'ema de caract\'eristique $p$.
\begin{enumerate}
\item Si $S$ est affine, alors
$H_\etale^q(S, \underline{\mathbf{Z}/p\mathbf{Z}}) = 0$ pour tout
$q \geq 2$.
\item Si $S$ est un sch\'ema quasi-compact et quasi-s\'epar\'e de
dimension $d$, alors $H_\etale^q(S, \underline{\mathbf{Z}/p\mathbf{Z}}) = 0$
pour tout $q \geq 2 + d$.
\end{enumerate}
\end{lemma}

\begin{proof}
Rappelons que la cohomologie \'etale du faisceau structural est \'egale
\`a sa cohomologie sur l'espace topologique sous-jacent
(théorème \ref{theorem-zariski-fpqc-quasi-coherent}).
La premi\`ere assertion r\'esulte de la suite exacte d'Artin-Schreier
et de l'annulation de la cohomologie du faisceau structural sur un sch\'ema
affine (Cohomologie des schémas, lemme
\ref{coherent-lemma-quasi-coherent-affine-cohomology-zero}).
La seconde assertion r\'esulte par le m\^eme argument de l'annulation
donn\'ee par Cohomologie, Proposition
\ref{cohomology-proposition-cohomological-dimension-spectral}
et du fait que $S$ est un espace spectral
(Properties, lemme
\ref{properties-lemma-quasi-compact-quasi-separated-spectral}).
\end{proof}

\begin{lemma}
\label{lemma-F-1}
Soit $k$ un corps alg\'ebriquement clos de caract\'eristique $p > 0$.
Soit $V$ un $k$-espace vectoriel de dimension finie. Soit $F : V \to V$
une application semi-lin\'eaire relativement au Frobenius, c'est-\`a-dire une application additive telle que
$F(\lambda v) = \lambda^p F(v)$ pour tous $\lambda \in k$ et $v \in V$.
Alors $F - 1 : V \to V$ est surjective, et son noyau est un
$\mathbf{F}_p$-espace vectoriel de dimension finie, celle-ci \'etant $\leq \dim_k(V)$.
\end{lemma}

\begin{proof}
Si $F = 0$, l'\'enonc\'e est vrai. Si $V$ est muni d'une filtration par des
sous-espaces vectoriels stables par $F$ telle que l'\'enonc\'e soit vrai pour chaque
quotient gradu\'e, alors il est vrai pour $(V, F)$. En combinant ces deux remarques,
nous pouvons supposer que le noyau de $F$ est nul.

\medskip\noindent
Choisissons une base $v_1, \ldots, v_n$ de $V$ et \'ecrivons
$F(v_i) = \sum a_{ij} v_j$. Remarquons que $v = \sum \lambda_i v_i$
appartient au noyau si et seulement si $\sum \lambda_i^p a_{ij} v_j = 0$.
Comme $k$ est alg\'ebriquement clos, cela implique que la matrice $(a_{ij})$
est inversible. Notons $(b_{ij})$ son inverse. Pour voir que $F - 1$
est surjective, choisissons $w = \sum \mu_i v_i \in V$ et cherchons \`a r\'esoudre
$$
(F - 1)(\sum \lambda_iv_i) =
\sum \lambda_i^p a_{ij} v_j - \sum \lambda_j v_j = \sum \mu_j v_j
$$
Cela \'equivaut \`a
$$
\sum \lambda_j^p v_j - \sum b_{ij} \lambda_i v_j = \sum b_{ij} \mu_i v_j
$$
autrement dit
$$
\lambda_j^p - \sum b_{ij} \lambda_i = \sum b_{ij} \mu_i,
\quad j = 1, \ldots, \dim(V).
$$
L'alg\`ebre
$$
A = k[x_1, \ldots, x_n]/
(x_j^p - \sum b_{ij} x_i - \sum b_{ij} \mu_i)
$$
est lisse standard sur $k$
(Algèbre, définition \ref{algebra-definition-standard-smooth})
car la matrice $(b_{ij})$ est inversible et les d\'eriv\'ees partielles
des $x_j^p$ sont nulles. Une base de $A$ sur $k$ est l'ensemble des mon\^omes
$x_1^{e_1} \ldots x_n^{e_n}$ avec $e_i < p$, donc $\dim_k(A) = p^n$.
Comme $k$ est alg\'ebriquement clos, on voit que $\Spec(A)$ poss\`ede exactement
$p^n$ points. Il s'ensuit que $F - 1$ est surjective et que chaque fibre
poss\`ede $p^n$ points, c'est-\`a-dire que le noyau de $F - 1$ est un groupe d'ordre $p^n$.
\end{proof}

\begin{lemma}
\label{lemma-F-2}
Dans la situation du lemme \ref{lemma-F-1}, soit $k'/k$ une extension
de corps alg\'ebriquement clos. Posons $V' = V \otimes_k k'$ et soit
$F' : V' \to V'$ l'application semi-lin\'eaire relativement au Frobenius induite par $F$.
Alors $\Ker(F - 1)$ s'envoie isomorphiquement sur $\Ker(F' - 1)$.
\end{lemma}

\begin{proof}
Soit $v_1, \ldots, v_n$ une base de $V$ et \'ecrivons $F(v_i) = \sum a_{ij}v_j$.
Alors $v = \sum \lambda_i v_i$ appartient \`a $\Ker(F - 1)$ si et seulement si
$\sum \lambda_i^pa_{ij} = \lambda_j$ pour $j = 1, \ldots, n$.
D'apr\`es le lemme \ref{lemma-F-1}, ces \'equations
d\'efinissent un sous-sch\'ema ferm\'e $Z$ de
$\mathbf{A}^n_k = \Spec(k[\lambda_1, \ldots, \lambda_n])$
fini sur $\Spec(k)$. Alors $Z(k) = Z(k')$, ce qui permet de conclure.
\end{proof}

\begin{lemma}
\label{lemma-top-cohomology-coherent}
Soit $X$ un sch\'ema s\'epar\'e de type fini sur un corps $k$.
Soit $\mathcal{F}$ un faisceau coh\'erent de $\mathcal{O}_X$-modules.
Alors $\dim_k H^d(X, \mathcal{F}) < \infty$, o\`u $d = \dim(X)$.
\end{lemma}

\begin{proof}
Nous allons le d\'emontrer par r\'ecurrence sur $d$. Le cas $d = 0$ est vrai car
dans ce cas $X$ est le spectre d'une $k$-alg\`ebre $A$ de dimension finie
(Variétés, lemme \ref{varieties-lemma-algebraic-scheme-dim-0})
et tout faisceau coh\'erent $\mathcal{F}$ correspond \`a un $A$-module de type fini
$M = H^0(X, \mathcal{F})$ qui v\'erifie $\dim_k M < \infty$.

\medskip\noindent
Supposons $d > 0$ et le r\'esultat d\'emontr\'e pour les sch\'emas s\'epar\'es
de type fini de dimension $< d$. Le sch\'ema $X$ est noeth\'erien. Consid\'erons
la propri\'et\'e $\mathcal{P}$ des faisceaux coh\'erents sur $X$ d\'efinie par la r\`egle
$$
\mathcal{P}(\mathcal{F}) \Leftrightarrow
\dim_k H^d(X, \mathcal{F}) < \infty
$$
Nous allons utiliser le r\'esultat de
Cohomologie des schémas, lemme \ref{coherent-lemma-property-initial}
pour montrer que $\mathcal{P}$ est satisfaite par tout faisceau coh\'erent sur $X$.

\medskip\noindent
Soit
$$
0 \to \mathcal{F}_1 \to \mathcal{F} \to \mathcal{F}_2 \to 0
$$
une suite exacte courte de faisceaux coh\'erents sur $X$.
Consid\'erons la suite exacte longue de cohomologie
$$
H^d(X, \mathcal{F}_1) \to
H^d(X, \mathcal{F}) \to
H^d(X, \mathcal{F}_2)
$$
Ainsi, si $\mathcal{P}$ est satisfaite par $\mathcal{F}_1$ et $\mathcal{F}_2$,
elle l'est par $\mathcal{F}$.

\medskip\noindent
Soit $Z \subset X$ un sous-sch\'ema ferm\'e int\`egre. Soit $\mathcal{I}$
un faisceau coh\'erent d'id\'eaux sur $Z$. Pour achever la d\'emonstration, il faut montrer
que $H^d(X, i_*\mathcal{I}) = H^d(Z, \mathcal{I})$ est de dimension finie.
Si $\dim(Z) < d$, le r\'esultat est vrai car ce groupe de cohomologie
est nul (Cohomologie, Proposition
\ref{cohomology-proposition-vanishing-Noetherian}).
Nous sommes ainsi ramen\'es \`a la situation examin\'ee au paragraphe suivant.

\medskip\noindent
Supposons que $X$ soit une vari\'et\'e de dimension $d$ et que
$\mathcal{F} = \mathcal{I}$ soit un faisceau coh\'erent d'id\'eaux. Dans ce
cas, nous avons une suite exacte courte
$$
0 \to \mathcal{I} \to \mathcal{O}_X \to i_*\mathcal{O}_Z \to 0
$$
o\`u $i : Z \to X$ est l'immersion ferm\'ee d\'efinie par $\mathcal{I}$.
L'hypoth\`ese de r\'ecurrence montre que
$H^{d - 1}(Z, \mathcal{O}_Z) = H^{d - 1}(X, i_*\mathcal{O}_Z)$ est
de dimension finie. Il suffit donc de d\'emontrer le r\'esultat
pour le faisceau structural.

\medskip\noindent
Nous pouvons appliquer le lemme de Chow
(Cohomologie des schémas, lemme \ref{coherent-lemma-chow-Noetherian})
au morphisme $X \to \Spec(k)$. Nous obtenons ainsi un diagramme
$$
\xymatrix{
X \ar[rd]_g & X' \ar[d]^-{g'} \ar[l]^\pi \ar[r]_i & \mathbf{P}^n_k \ar[dl] \\
& \Spec(k) &
}
$$
comme dans l'\'enonc\'e du lemme de Chow. Soit en outre $U \subset X$
le sous-sch\'ema ouvert dense tel que $\pi^{-1}(U) \to U$ soit un isomorphisme.
Nous pouvons aussi supposer que $X'$ est une vari\'et\'e, voir
Cohomologie des schémas, remarque \ref{coherent-remark-chow-Noetherian}.
Le morphisme $i' = (i, \pi) : X' \to \mathbf{P}^n_X$ est
une immersion ferm\'ee (loc. cit.). Par cons\'equent
$$
\mathcal{L} = i^*\mathcal{O}_{\mathbf{P}^n_k}(1) \cong
(i')^*\mathcal{O}_{\mathbf{P}^n_X}(1)
$$
est ample relativement \`a $\pi$ (par exemple d'apr\`es
Morphismes, lemme \ref{morphisms-lemma-characterize-ample-on-finite-type}).
Ainsi, d'apr\`es Cohomologie des schémas, lemme \ref{coherent-lemma-kill-by-twisting},
il existe un $n \geq 0$ tel que
$R^p\pi_*\mathcal{L}^{\otimes n} = 0$ pour tout $p > 0$.
Posons $\mathcal{G} = \pi_*\mathcal{L}^{\otimes n}$.
Choisissons une section globale non nulle $s$ de $\mathcal{L}^{\otimes n}$.
Puisque $\mathcal{G} = \pi_*\mathcal{L}^{\otimes n}$, la section $s$
correspond \`a une section de $\mathcal{G}$, c'est-\`a-dire \`a un morphisme
$\mathcal{O}_X \to \mathcal{G}$.
Comme $s|_U \not = 0$, puisque $X'$ est une vari\'et\'e et $\mathcal{L}$ est
inversible, on voit que $\mathcal{O}_X|_U \to \mathcal{G}|_U$
est non nul. Comme $\mathcal{G}|_U = \mathcal{L}^{\otimes n}|_{\pi^{-1}(U)}$
est inversible, nous obtenons une suite exacte courte
$$
0 \to \mathcal{O}_X \to \mathcal{G} \to \mathcal{Q} \to 0
$$
o\`u $\mathcal{Q}$ est coh\'erent et \`a support dans un sous-sch\'ema
ferm\'e strict de $X$. En raisonnant comme ci-dessus avec notre hypoth\`ese de
r\'ecurrence, on voit qu'il
suffit de d\'emontrer $\dim H^d(X, \mathcal{G}) < \infty$.

\medskip\noindent
La suite spectrale de Leray
(Cohomologie, lemme \ref{cohomology-lemma-apply-Leray})
montre que $H^d(X, \mathcal{G}) = H^d(X', \mathcal{L}^{\otimes n})$.
Soit $\overline{X}' \subset \mathbf{P}^n_k$ l'adh\'erence
de $X'$. Alors $\overline{X}'$ est une vari\'et\'e projective de dimension $d$
sur $k$ et $X' \subset \overline{X}'$ est un ouvert dense.
Le faisceau inversible $\mathcal{L}^{\otimes n}$ est la restriction de
$\mathcal{O}_{\overline{X}'}(n)$ \`a $X'$. D'apr\`es
Cohomologie, Proposition
\ref{cohomology-proposition-cohomological-dimension-spectral}
l'application
$$
H^d(\overline{X}', \mathcal{O}_{\overline{X}'}(n))
\longrightarrow
H^d(X', \mathcal{L}^{\otimes n})
$$
est surjective. Comme le groupe de cohomologie de gauche est
de dimension finie d'apr\`es
Cohomologie des schémas, lemme \ref{coherent-lemma-coherent-projective}
la d\'emonstration est achev\'ee.
\end{proof}

\begin{lemma}
\label{lemma-vanishing-variety-char-p-p}
Soit $X$ un sch\'ema s\'epar\'e de type fini sur un corps alg\'ebriquement clos
$k$ de caract\'eristique $p > 0$. Alors
$H_\etale^q(X, \underline{\mathbf{Z}/p\mathbf{Z}}) = 0$ pour
$q \geq dim(X) + 1$.
\end{lemma}

\begin{proof}
Posons $d = \dim(X)$. D'apr\`es l'annulation \'etablie dans le
lemme \ref{lemma-vanishing-affine-char-p-p},
il suffit de montrer que
$H_\etale^{d + 1}(X, \underline{\mathbf{Z}/p\mathbf{Z}}) = 0$.
D'apr\`es le lemme \ref{lemma-top-cohomology-coherent}, on voit que
$H^d(X, \mathcal{O}_X)$ est un espace vectoriel de dimension finie sur $k$.
Par cons\'equent, la suite exacte longue de cohomologie associ\'ee \`a la
suite d'Artin-Schreier se termine par
$$
H^d(X, \mathcal{O}_X) \xrightarrow{F - 1}
H^d(X, \mathcal{O}_X) \to H^{d + 1}_\etale(X, \mathbf{Z}/p\mathbf{Z}) \to 0
$$
D'apr\`es le lemme \ref{lemma-F-1}, l'application $F - 1$ de cette suite est surjective.
Cela d\'emontre le lemme.
\end{proof}

\begin{lemma}
\label{lemma-finiteness-proper-variety-char-p-p}
Soit $X$ un sch\'ema propre sur un corps alg\'ebriquement clos
$k$ de caract\'eristique $p > 0$. Alors
\begin{enumerate}
\item $H_\etale^q(X, \underline{\mathbf{Z}/p\mathbf{Z}})$
est un $\mathbf{Z}/p\mathbf{Z}$-module de type fini pour tout $q$, et
\item $H^q_\etale(X, \underline{\mathbf{Z}/p\mathbf{Z}}) \to
H^q_\etale(X_{k'}, \underline{\mathbf{Z}/p\mathbf{Z}})$
est un isomorphisme si $k'/k$ est une extension de corps alg\'ebriquement
clos.
\end{enumerate}
\end{lemma}

\begin{proof}
D'apr\`es Cohomologie des schémas, lemme
\ref{coherent-lemma-proper-over-affine-cohomology-finite}
et la comparaison de la cohomologie donn\'ee par le
th\'eor\`eme \ref{theorem-zariski-fpqc-quasi-coherent}, les groupes de cohomologie
$H^q_\etale(X, \mathbf{G}_a) = H^q(X, \mathcal{O}_X)$ sont
des espaces vectoriels de dimension finie sur $k$. Par cons\'equent, d'apr\`es le
lemme \ref{lemma-F-1}, la suite exacte longue de cohomologie
associ\'ee \`a la suite d'Artin-Schreier se d\'ecompose en
suites exactes courtes
$$
0 \to H_\etale^q(X, \underline{\mathbf{Z}/p\mathbf{Z}}) \to
H^q(X, \mathcal{O}_X) \xrightarrow{F - 1} H^q(X, \mathcal{O}_X) \to 0
$$
et, de plus, la dimension sur $\mathbf{F}_p$ des groupes de cohomologie
$H_\etale^q(X, \underline{\mathbf{Z}/p\mathbf{Z}})$ est inf\'erieure ou \'egale \`a la
dimension sur $k$ de l'espace vectoriel $H^q(X, \mathcal{O}_X)$.
Cela d\'emontre la premi\`ere assertion. La seconde r\'esulte
du lemme \ref{lemma-F-2}
car $H^q(X, \mathcal{O}_X) \otimes_k k' \to H^q(X_{k'}, \mathcal{O}_{X_{k'}})$
est un isomorphisme par changement de base plat
(Cohomologie des schémas,
lemme \ref{coherent-lemma-flat-base-change-cohomology}).
\end{proof}








\section{Faisceaux localement constants}
\label{section-locally-constant}

\noindent
Cette section est l'analogue de
Modules sur les sites, section \ref{sites-modules-section-locally-constant}
pour le site \'etale.

\begin{definition}
\label{definition-finite-locally-constant}
Soit $X$ un sch\'ema.
Soit $\mathcal{F}$ un faisceau d'ensembles sur $X_\etale$.
\begin{enumerate}
\item Soit $E$ un ensemble. Nous disons que $\mathcal{F}$ est le
{\it faisceau constant de valeur $E$} si $\mathcal{F}$ est le
faisceau associ\'e au pr\'efaisceau $U \mapsto E$.
Notation : $\underline{E}_X$ ou $\underline{E}$.
\item Nous disons que $\mathcal{F}$ est un {\it faisceau constant} s'il est
isomorphe \`a un faisceau comme en (1).
\item Nous disons que $\mathcal{F}$ est {\it localement constant} s'il existe un
recouvrement $\{U_i \to X\}$ tel que $\mathcal{F}|_{U_i}$ soit un faisceau constant.
\item Nous disons que $\mathcal{F}$ est {\it localement constant fini} s'il
est localement constant et si ses valeurs locales sont des ensembles finis.
\end{enumerate}
Soit $\mathcal{F}$ un faisceau de groupes ab\'eliens sur $X_\etale$.
\begin{enumerate}
\item Soit $A$ un groupe ab\'elien.
Nous disons que $\mathcal{F}$ est le {\it faisceau constant de valeur $A$} si
$\mathcal{F}$ est le faisceau associ\'e au pr\'efaisceau $U \mapsto A$.
Notation : $\underline{A}_X$ ou $\underline{A}$.
\item Nous disons que $\mathcal{F}$ est un {\it faisceau constant} s'il est isomorphe
comme faisceau ab\'elien \`a un faisceau comme en (1).
\item Nous disons que $\mathcal{F}$ est {\it localement constant} s'il existe un
recouvrement $\{U_i \to X\}$ tel que $\mathcal{F}|_{U_i}$ soit un faisceau constant.
\item Nous disons que $\mathcal{F}$ est {\it localement constant fini} s'il
est localement constant et si ses valeurs locales sont des groupes ab\'eliens finis.
\end{enumerate}
Soit $\Lambda$ un anneau. Soit $\mathcal{F}$ un faisceau de $\Lambda$-modules
sur $X_\etale$.
\begin{enumerate}
\item Soit $M$ un $\Lambda$-module.
Nous disons que $\mathcal{F}$ est le {\it faisceau constant de valeur $M$} si
$\mathcal{F}$ est le faisceau associ\'e au pr\'efaisceau $U \mapsto M$.
Notation : $\underline{M}_X$ ou $\underline{M}$.
\item Nous disons que $\mathcal{F}$ est un {\it faisceau constant} s'il est isomorphe
comme faisceau de $\Lambda$-modules \`a un faisceau comme en (1).
\item Nous disons que $\mathcal{F}$ est {\it localement constant} s'il existe un
recouvrement $\{U_i \to X\}$ tel que $\mathcal{F}|_{U_i}$ soit un faisceau constant.
\end{enumerate}
\end{definition}

\begin{lemma}
\label{lemma-pullback-locally-constant}
Soit $f : X \to Y$ un morphisme de sch\'emas. Si $\mathcal{G}$ est un
faisceau localement constant d'ensembles, de groupes ab\'eliens ou de $\Lambda$-modules
sur $Y_\etale$, alors $f^{-1}\mathcal{G}$ l'est aussi
sur $X_\etale$.
\end{lemma}

\begin{proof}
Cela vaut pour tout morphisme de topos, voir
Modules sur les sites, lemme \ref{sites-modules-lemma-pullback-locally-constant}.
\end{proof}

\begin{lemma}
\label{lemma-pushforward-locally-constant}
Soit $f : X \to Y$ un morphisme fini \'etale de sch\'emas.
Si $\mathcal{F}$ est un faisceau d'ensembles localement constant (fini),
un faisceau de groupes ab\'eliens localement constant (fini), ou
un faisceau de $\Lambda$-modules localement constant (de type fini)
sur $X_\etale$, alors $f_*\mathcal{F}$ l'est aussi
sur $Y_\etale$.
\end{lemma}

\begin{proof}
La construction de $f_*$ commute \`a la localisation \'etale.
Un morphisme fini \'etale est localement isomorphe \`a une somme disjointe
d'isomorphismes, voir
Morphismes étales, lemme \ref{etale-lemma-finite-etale-etale-local}.
Ainsi le lemme affirme que, si $\mathcal{F}_i$, $i = 1, \ldots, n$
sont des faisceaux d'ensembles localement constants (finis), alors
$\prod_{i = 1, \ldots, n} \mathcal{F}_i$ l'est aussi.
C'est clair. Il en va de m\^eme pour les faisceaux de groupes ab\'eliens et de modules.
\end{proof}

\begin{lemma}
\label{lemma-characterize-finite-locally-constant}
Soient $X$ un sch\'ema et $\mathcal{F}$ un faisceau d'ensembles sur $X_\etale$.
Alors les assertions suivantes sont \`equivalentes
\begin{enumerate}
\item $\mathcal{F}$ est localement constant fini, et
\item $\mathcal{F} = h_U$ pour un morphisme fini \'etale $U \to X$.
\end{enumerate}
\end{lemma}

\begin{proof}
Un morphisme fini \'etale est localement isomorphe \`a une somme disjointe
d'isomorphismes, voir
Morphismes étales, lemme \ref{etale-lemma-finite-etale-etale-local}.
Ainsi (2) implique (1). R\'eciproquement, si $\mathcal{F}$ est localement
constant fini, il existe un recouvrement \'etale $\{X_i \to X\}$ tel que
$\mathcal{F}|_{X_i}$ soit repr\'esentable par un morphisme fini \'etale $U_i \to X_i$.
En raisonnant exactement comme dans la d\'emonstration du
Descente, lemme \ref{descent-lemma-descent-data-sheaves}
nous obtenons une donn\'ee de descente de sch\'emas $(U_i, \varphi_{ij})$ relativement \`a
$\{X_i \to X\}$ (d\'etails omis). Cette donn\'ee de descente est effective, par
exemple d'apr\`es Descente, lemme \ref{descent-lemma-affine},
et le morphisme de sch\'emas obtenu $U \to X$ est fini \'etale
d'apr\`es Descente, lemmes \ref{descent-lemma-descending-property-finite} et
\ref{descent-lemma-descending-property-etale}.
\end{proof}

\begin{lemma}
\label{lemma-morphism-locally-constant}
Soit $X$ un sch\'ema.
\begin{enumerate}
\item Soit $\varphi : \mathcal{F} \to \mathcal{G}$ un morphisme
de faisceaux d'ensembles localement constants sur $X_\etale$.
Si $\mathcal{F}$ est localement constant fini, il existe un
recouvrement \'etale $\{U_i \to X\}$ tel que
$\varphi|_{U_i}$ soit le morphisme de faisceaux constants associ\'e \`a
une application d'ensembles.
\item Soit $\varphi : \mathcal{F} \to \mathcal{G}$ un morphisme
de faisceaux de groupes ab\'eliens localement constants sur $X_\etale$.
Si $\mathcal{F}$ est localement constant fini, il existe un recouvrement \'etale
$\{U_i \to X\}$ tel que $\varphi|_{U_i}$ soit le morphisme de
faisceaux ab\'eliens constants associ\'e \`a un homomorphisme de groupes ab\'eliens.
\item Soit $\Lambda$ un anneau.
Soit $\varphi : \mathcal{F} \to \mathcal{G}$ un morphisme
de faisceaux de $\Lambda$-modules localement constants sur $X_\etale$.
Si $\mathcal{F}$ est de type fini, alors il existe un recouvrement \'etale
$\{U_i \to X\}$ tel que $\varphi|_{U_i}$ soit le morphisme de
faisceaux constants de $\Lambda$-modules associ\'e \`a un homomorphisme de $\Lambda$-modules.
\end{enumerate}
\end{lemma}

\begin{proof}
Cela vaut sur tout site, voir
Modules sur les sites, lemme \ref{sites-modules-lemma-morphism-locally-constant}.
\end{proof}

\begin{lemma}
\label{lemma-kernel-finite-locally-constant}
Soit $X$ un sch\'ema.
\begin{enumerate}
\item La cat\'egorie des faisceaux d'ensembles localement constants finis
est stable par limites projectives et inductives finies dans $\Sh(X_\etale)$.
\item La cat\'egorie des faisceaux ab\'eliens localement constants finis est une
sous-cat\'egorie de Serre faible de $\textit{Ab}(X_\etale)$.
\item Soit $\Lambda$ un anneau noeth\'erien. La cat\'egorie des
faisceaux de $\Lambda$-modules localement constants de type fini sur
$X_\etale$ est une sous-cat\'egorie de Serre faible de
$\textit{Mod}(X_\etale, \Lambda)$.
\end{enumerate}
\end{lemma}

\begin{proof}
Cela vaut sur tout site, voir
Modules sur les sites, lemme
\ref{sites-modules-lemma-kernel-finite-locally-constant}.
\end{proof}

\begin{lemma}
\label{lemma-tensor-product-locally-constant}
Soit $X$ un sch\'ema. Soit $\Lambda$ un anneau.
Le produit tensoriel de deux faisceaux de $\Lambda$-modules localement constants
sur $X_\etale$ est un faisceau de $\Lambda$-modules localement constant.
\end{lemma}

\begin{proof}
Cela vaut sur tout site, voir
Modules sur les sites, lemme
\ref{sites-modules-lemma-tensor-product-locally-constant}.
\end{proof}

\begin{lemma}
\label{lemma-connected-locally-constant}
Soit $X$ un sch\'ema connexe. Soit $\Lambda$ un anneau et soit
$\mathcal{F}$ un faisceau de $\Lambda$-modules localement constant.
Alors il existe un $\Lambda$-module $M$ et un recouvrement \'etale
$\{U_i \to X\}$ tels que $\mathcal{F}|_{U_i} \cong \underline{M}|_{U_i}$.
\end{lemma}

\begin{proof}
Choisissons un recouvrement \'etale
$\{U_i \to X\}$ tel que $\mathcal{F}|_{U_i}$ soit constant, disons
$\mathcal{F}|_{U_i} \cong \underline{M_i}_{U_i}$.
Remarquons que $U_i \times_X U_j$ est vide si $M_i$ n'est pas isomorphe
\`a $M_j$.
Pour chaque classe d'isomorphisme de $\Lambda$-modules repr\'esent\'ee par $M$, posons $I_M = \{i \in I \mid M_i \cong M\}$.
Comme les morphismes \'etales sont ouverts, nous voyons que
$U_M = \bigcup_{i \in I_M} \Im(U_i \to X)$
est un ouvert de $X$. Alors $X = \coprod U_M$ est un recouvrement
ouvert disjoint de $X$. Comme $X$ est connexe, un seul $U_M$ est non vide
et le lemme en r\'esulte.
\end{proof}






\section{Faisceaux localement constants et groupe fondamental}
\label{section-pione}

\noindent
Dans certains cas, nous pouvons relier les faisceaux localement constants
au groupe fondamental d'un sch\'ema.

\begin{lemma}
\label{lemma-locally-constant-on-connected}
Soit $X$ un sch\'ema connexe. Soit $\overline{x}$ un point g\'eom\'etrique de $X$.
\begin{enumerate}
\item Il existe une \`equivalence de cat\'egories
$$
\left\{
\begin{matrix}
\text{faisceaux d'ensembles localement constants finis}\\
\text{sur }X_\etale
\end{matrix}
\right\}
\longleftrightarrow
\left\{
\begin{matrix}
\text{ensembles finis munis d'une action de }\pi_1(X, \overline{x})
\end{matrix}
\right\}
$$
\item Il existe une \`equivalence de cat\'egories
$$
\left\{
\begin{matrix}
\text{faisceaux de groupes ab\'eliens localement constants finis}\\
\text{sur }X_\etale
\end{matrix}
\right\}
\longleftrightarrow
\left\{
\begin{matrix}
\pi_1(X, \overline{x})\text{-modules finis}
\end{matrix}
\right\}
$$
\item Soit $\Lambda$ un anneau fini. Il existe une \`equivalence de cat\'egories
$$
\left\{
\begin{matrix}
\text{faisceaux localement constants de type fini}\\
\text{de }\Lambda\text{-modules sur }X_\etale
\end{matrix}
\right\}
\longleftrightarrow
\left\{
\begin{matrix}
\pi_1(X, \overline{x})\text{-modules finis munis}\\
\text{d'une structure commutante de }\Lambda\text{-module}
\end{matrix}
\right\}
$$
\end{enumerate}
\end{lemma}

\begin{proof}
Remarquons que $\pi_1(X, \overline{x})$ est un groupe topologique
profini, voir Fundamental Groups, définition
\ref{pione-definition-fundamental-group}.
Les cat\'egories des membres de gauche sont d\'efinies dans la
section \ref{section-locally-constant}.
La notation employ\'ee dans les membres de droite vient de
Fundamental Groups, définition \ref{pione-definition-G-set-continuous}
pour les ensembles et de la
d\'efinition \ref{definition-G-module-continuous} pour les groupes ab\'eliens.
Cela explique la notation.

\medskip\noindent
L'assertion (1) r\'esulte du
lemme \ref{lemma-characterize-finite-locally-constant}
et de Fundamental Groups, théorème \ref{pione-theorem-fundamental-group}.
Les parties (2) et (3) s'en d\'eduisent imm\'ediatement en munissant les
(faisceaux d')ensembles sous-jacents d'une structure suppl\'ementaire. Par exemple, un
faisceau de groupes ab\'eliens localement constant fini sur $X_\etale$ revient
\`a un faisceau d'ensembles localement constant fini $\mathcal{F}$
muni d'un morphisme $+ : \mathcal{F} \times \mathcal{F} \to \mathcal{F}$
satisfaisant aux axiomes usuels. L'\'equivalence de (1) pr\'eserve les produits
et transforme donc $+$ en une addition sur l'ensemble fini correspondant
muni d'une action de $\pi_1(X, \overline{x})$. Puisque les $\pi_1(X, \overline{x})$-modules
reviennent aux ensembles munis d'une action de $\pi_1(X, \overline{x})$ et d'une
structure compatible de groupe ab\'elien, nous obtenons (2). La partie (3) se d\'emontre
exactement de la m\^eme mani\`ere.
\end{proof}

\begin{lemma}
\label{lemma-locally-constant-on-connected-geom-unibranch}
Soit $X$ un sch\'ema irr\'eductible et g\'eom\'etriquement unibranche.
Soit $\overline{x}$ un point g\'eom\'etrique de $X$.
Soit $\Lambda$ un anneau. Il existe une \`equivalence de cat\'egories
$$
\left\{
\begin{matrix}
\text{faisceaux localement constants de type fini}\\
\text{de }\Lambda\text{-modules sur }X_\etale
\end{matrix}
\right\}
\longleftrightarrow
\left\{
\begin{matrix}
\Lambda\text{-modules de type fini }M\text{ munis}\\
\text{d'une action continue de }\pi_1(X, \overline{x})
\end{matrix}
\right\}
$$
\end{lemma}

\begin{proof}
La d\'emonstration du lemme \ref{lemma-locally-constant-on-connected}
ne s'applique pas, car un $\Lambda$-module de type fini $M$ peut ne pas avoir
un ensemble sous-jacent fini.

\medskip\noindent
Soit $\nu : X^\nu \to X$ le morphisme de normalisation. D'apr\`es Morphismes, lemme
\ref{morphisms-lemma-normalization-geom-unibranch-univ-homeo}
il s'agit d'un hom\'eomorphisme universel. D'apr\`es
Fundamental Groups, Proposition \ref{pione-proposition-universal-homeomorphism}
il induit un isomorphisme
$\pi_1(X^\nu, \overline{x}) \to \pi_1(X, \overline{x})$
et, d'apr\`es le th\'eor\`eme \ref{theorem-topological-invariance}, nous obtenons
une \`equivalence de cat\'egories entre les faisceaux localement constants de type fini
de $\Lambda$-modules sur $X_\etale$ et sur $X^\nu_\etale$.
Nous sommes ainsi ramen\'es au cas o\`u $X$ est un sch\'ema int\`egre normal.

\medskip\noindent
Supposons que $X$ soit un sch\'ema int\`egre normal. Soit $\eta \in X$ le point g\'en\'erique.
Soit $\overline{\eta}$ un point g\'eom\'etrique au-dessus de $\eta$.
D'apr\`es Fundamental Groups, Proposition \ref{pione-proposition-normal},
nous avons une surjection continue
$$
\text{Gal}(\kappa(\eta)^{sep}/\kappa(\eta)) = \pi_1(\eta, \overline{\eta})
\longrightarrow
\pi_1(X, \overline{\eta})
$$
dont le noyau est d\'ecrit dans Fundamental Groups, lemme
\ref{pione-lemma-normal-pione-quotient-inertia}.
Soit $\mathcal{F}$ un faisceau localement constant de type fini de $\Lambda$-modules
sur $X_\etale$. Soit $M = \mathcal{F}_{\overline{\eta}}$
la fibre de $\mathcal{F}$ en $\overline{\eta}$.
La section \ref{section-galois-action-stalks} nous donne une action continue de
$\text{Gal}(\kappa(\eta)^{sep}/\kappa(\eta))$ sur $M$.
Notre but est de montrer que cette action se factorise par la
surjection affich\'ee.
Puisque $\mathcal{F}$ est de type fini, $M$ est un $\Lambda$-module de type fini.
Puisque $\mathcal{F}$ est localement constant, pour tout $x \in X$,
la restriction de $\mathcal{F}$ \`a $\Spec(\mathcal{O}_{X, x}^{sh})$
est constante. L'action de $\text{Gal}(K^{sep}/K_x^{sh})$
(avec les notations de Fundamental Groups, lemme
\ref{pione-lemma-normal-pione-quotient-inertia})
sur $M$ est donc triviale. Nous obtenons la factorisation voulue.

\medskip\noindent
R\'eciproquement, supposons donn\'e un $\Lambda$-module de type fini $M$
muni d'une action continue de $\pi_1(X, \overline{\eta})$.
Nous allons construire un $\mathcal{F}$ tel que
$M \cong \mathcal{F}_{\overline{\eta}}$ comme
$\Lambda[\pi_1(X, \overline{\eta})]$-modules.
Choisissons des g\'en\'erateurs $m_1, \ldots, m_r \in M$.
Comme l'action de $\pi_1(X, \overline{\eta})$ sur $M$ est
continue, il existe pour tout $i$ un sous-groupe ouvert
$H_i$ du groupe profini $\pi_1(X, \overline{\eta})$
tel que tout $\gamma \in H_i$ fixe $m_i$. Nous en d\'eduisons que
tout \`el\'ement du sous-groupe ouvert $H = \bigcap_{i = 1, \ldots, r} H_i$
fixe tout \`el\'ement de $M$. En rempla\c cant $H$ par un sous-groupe plus petit, nous pouvons supposer que $H$
est un sous-groupe ouvert normal de $\pi_1(X, \overline{\eta})$.
Posons $G = \pi_1(X, \overline{\eta})/H$. Soit $f : Y \to X$ le
rev\^etement galoisien fini \'etale de groupe $G$ correspondant.
Nous pouvons voir $f_*\underline{\mathbf{Z}}$ comme un faisceau de
$\mathbf{Z}[G]$-modules sur $X_\etale$. Il suffit alors de poser
$$
\mathcal{F} =
f_*\underline{\mathbf{Z}} \otimes_{\underline{\mathbf{Z}[G]}} \underline{M}
$$
Nous laissons au lecteur le soin de calculer
$\mathcal{F}_{\overline{\eta}}$.
Nous omettons aussi de v\'erifier que cette construction
est inverse de celle du paragraphe pr\'ec\'edent.
\end{proof}

\begin{remark}
\label{remark-functorial-locally-constant-on-connected}
Les \`equivalences des lemmes \ref{lemma-locally-constant-on-connected} et
\ref{lemma-locally-constant-on-connected-geom-unibranch}
sont compatibles aux images inverses. Par exemple, soit $f : Y \to X$
un morphisme de sch\'emas connexes. Soit $\overline{y}$ un point g\'eom\'etrique
de $Y$ et posons $\overline{x} = f(\overline{y})$.
Alors le diagramme
$$
\xymatrix{
\substack{\text{faisceaux d'ensembles finis} \\
\text{localement constants sur }Y_\etale}
\ar[r] &
\substack{\text{ensembles finis} \\
\text{munis d'une action de }\pi_1(Y, \overline{y})} \\
\substack{\text{faisceaux d'ensembles finis} \\
\text{localement constants sur }X_\etale}
\ar[r] \ar[u]_{f^{-1}} &
\substack{\text{ensembles finis} \\
\text{munis d'une action de }\pi_1(X, \overline{x})} \ar[u]
}
$$
est commutatif, o\`u la fl\`eche verticale de droite provient de
l'homomorphisme continu
$\pi_1(Y, \overline{y}) \to \pi_1(X, \overline{x})$
induit par $f$. Cela r\'esulte imm\'ediatement du
diagramme commutatif de
Fundamental Groups, théorème \ref{pione-theorem-fundamental-group}.
Un r\'esultat analogue vaut dans les autres cas.
\end{remark}










\section{M\'ethode de la trace}
\label{section-trace-method}

\noindent
Une r\'ef\'erence pour cette section est \cite[Expos\'e IX, \S 5]{SGA4}.
Les r\'esultats qui suivent seront utilis\'es dans la d\'emonstration du
lemme \ref{lemma-vanishing-easier} ci-dessous.

\medskip\noindent
Soit $f : Y \to X$ un morphisme \'etale de sch\'emas. Il existe
une suite
$$
f_!, f^{-1}, f_*
$$
de foncteurs adjoints entre
$\textit{Ab}(X_\etale)$ et $\textit{Ab}(Y_\etale)$.
Le foncteur $f_!$ est \'etudi\'e dans la section \ref{section-extension-by-zero}.
Le morphisme d'adjonction $\text{id} \to f_* f^{-1}$ est appel\'e {\it restriction}.
Le morphisme d'adjonction $f_! f^{-1} \to \text{id}$ est souvent
appel\'e le {\it morphisme trace}. Si $f$ est fini \'etale, alors $f_* = f_!$
(lemme \ref{lemma-shriek-equals-star-finite-etale}) et
nous pouvons le voir comme un morphisme $f_*f^{-1} \to \text{id}$.

\begin{definition}
\label{definition-trace-map}
Soit $f : Y \to X$ un morphisme fini \'etale de sch\'emas.
Le morphisme $f_* f^{-1} \to \text{id}$ d\'ecrit ci-dessus et explicit\'e ci-dessous
est appel\'e la {\it trace}.
\end{definition}

\noindent
Soit $f : Y \to X$ un morphisme fini \'etale de sch\'emas. Le morphisme trace est
caract\'eris\'e par les deux propri\'et\'es suivantes :
\begin{enumerate}
\item il commute \`a la localisation \'etale sur $X$, et
\item si $Y = \coprod_{i = 1}^d X$, alors le morphisme trace est
l'application somme $f_*f^{-1} \mathcal{F} = \mathcal{F}^{\oplus d} \to \mathcal{F}$.
\end{enumerate}
D'apr\`es Morphismes étales, lemme \ref{etale-lemma-finite-etale-etale-local},
tout morphisme fini \'etale $f : Y \to X$ est, localement pour la topologie \'etale sur $X$,
de la forme donn\'ee en (2) pour un entier $d \geq 0$. Nous pouvons donc
d\'efinir le morphisme trace par cette caract\'erisation; en particulier,
nous n'avons besoin ni de l'existence de $f_!$, ni de l'\'egalit\'e
de $f_!$ et $f_*$ pour construire le morphisme trace.
Cette description montre que, si $f$ est de degr\'e constant $d$, alors
la compos\'ee
$$
\mathcal{F} \xrightarrow{res}
f_* f^{-1} \mathcal{F} \xrightarrow{trace}
\mathcal{F}
$$
est la multiplication par $d$. La « m\'ethode de la trace »
est l'observation suivante : si $\mathcal{F}$
est un faisceau ab\'elien sur $X_\etale$ tel que la multiplication par $d$
sur $\mathcal{F}$ soit un isomorphisme, alors l'application
$$
H^n_\etale(X, \mathcal{F}) \longrightarrow H^n_\etale(Y, f^{-1}\mathcal{F})
$$
est injective. En effet, nous avons
$$
H^n_\etale(Y, f^{-1}\mathcal{F}) = H^n_\etale(X, f_*f^{-1}\mathcal{F})
$$
par l'annulation des images directes sup\'erieures
(proposition \ref{proposition-finite-higher-direct-image-zero})
et la suite spectrale de Leray
(proposition \ref{proposition-leray}).
Nous pouvons donc consid\'erer les applications
$$
H^n_\etale(X, \mathcal{F}) \to
H^n_\etale(Y, f^{-1}\mathcal{F})= H^n_\etale(X, f_*f^{-1}\mathcal{F})
\xrightarrow{trace}
H^n_\etale(X, \mathcal{F})
$$
et leur compos\'ee est un isomorphisme (sous notre hypoth\`ese sur $\mathcal{F}$
et $f$). En particulier, si
$H_\etale^q(Y, f^{-1}\mathcal{F}) = 0$, alors
$H_\etale^q(X, \mathcal{F}) = 0$ aussi.
En effet, la multiplication par $d$ induit un
isomorphisme de $H_\etale^q(X, \mathcal{F})$ qui se factorise par
$H_\etale^q(Y, f^{-1}\mathcal{F})= 0$.

\medskip\noindent
On combine souvent ceci avec le lemme suivant.

\begin{lemma}
\label{lemma-pullback-filtered}
Soit $S$ un sch\'ema connexe. Soit $\ell$ un nombre premier. Soit
$\mathcal{F}$ un faisceau localement constant de type fini de
$\mathbf{F}_\ell$-espaces vectoriels sur $S_\etale$.
Alors il existe un morphisme fini \'etale
$f : T \to S$ de degr\'e premier \`a $\ell$ tel que $f^{-1}\mathcal{F}$
admette une filtration finie dont les quotients successifs sont
$\underline{\mathbf{Z}/\ell\mathbf{Z}}_T$.
\end{lemma}

\begin{proof}
Choisissons un point g\'eom\'etrique $\overline{s}$ de $S$.
Par l'\'equivalence du lemme \ref{lemma-locally-constant-on-connected},
le faisceau $\mathcal{F}$ correspond \`a un espace vectoriel
$V$ de dimension finie sur $\mathbf{F}_\ell$, muni d'une action continue de
$\pi_1(S, \overline{s})$.
Soit $G \subset \text{Aut}(V)$ l'image de l'homomorphisme
$\rho : \pi_1(S, \overline{s}) \to \text{Aut}(V)$ qui donne l'action.
Remarquons que $G$ est fini.
L'homomorphisme continu surjectif
$\overline{\rho} : \pi_1(S, \overline{s}) \to G$
correspond \`a un objet galoisien $Y \to S$ de
$\textit{F\'Et}_S$, de groupe d'automorphismes $G = \text{Aut}(Y/S)$, voir
Fundamental Groups, section \ref{pione-section-finite-etale-under-galois}.
Soit $H \subset G$ un $\ell$-sous-groupe de Sylow.
Nous affirmons que $T = Y/H \to S$ convient. En effet, soit $\overline{t} \in T$
un point g\'eom\'etrique au-dessus de $\overline{s}$. L'image de
$\pi_1(T, \overline{t}) \to \pi_1(S, \overline{s})$
est $(\overline{\rho})^{-1}(H)$ par fonctorialit\'e
des groupes fondamentaux. Par cons\'equent, l'action de $\pi_1(T, \overline{t})$
sur $V$ correspondant \`a $f^{-1}\mathcal{F}$ se fait par
l'application $\pi_1(T, \overline{t}) \to H$, voir la
remarque \ref{remark-functorial-locally-constant-on-connected}. Comme
$H$ est un $\ell$-groupe fini, les constituants irr\'eductibles de la
repr\'esentation $\rho|_{\pi_1(T, \overline{t})}$
sont tous triviaux de rang $1$ (c'est un lemme \`el\'ementaire de
th\'eorie des repr\'esentations des groupes finis).
Par l'\'equivalence du
lemme \ref{lemma-locally-constant-on-connected},
cela signifie que $f^{-1}\mathcal{F}$ est une extension successive de
faisceaux constants $\underline{\mathbf{Z}/\ell\mathbf{Z}}_T$.
De plus, le degr\'e de $T = Y/H \to S$ est premier \`a $\ell$,
car il est \'egal \`a l'indice de $H$ dans $G$.
\end{proof}

\begin{lemma}
\label{lemma-action-l-group}
Soit $\Lambda$ un anneau noeth\'erien.
Soit $\ell$ un nombre premier et $n \geq 1$.
Soit $H$ un $\ell$-groupe fini.
Soit $M$ un $\Lambda[H]$-module de type fini annul\'e par $\ell^n$.
Alors il existe une filtration finie
$0 = M_0 \subset M_1 \subset \ldots \subset M_t = M$
par des sous-$\Lambda[H]$-modules
telle que $H$ agisse trivialement sur $M_{i + 1}/M_i$ pour tout
$i = 0, \ldots, t - 1$.
\end{lemma}

\begin{proof}
Omis. Indication : montrer que l'id\'eal d'augmentation $\mathfrak m$
de l'anneau non commutatif $\mathbf{Z}/\ell^n\mathbf{Z}[H]$
est nilpotent.
\end{proof}

\begin{lemma}
\label{lemma-pullback-filtered-modules}
Soit $S$ un sch\'ema irr\'eductible et g\'eom\'etriquement unibranche.
Soit $\ell$ un nombre premier et $n \geq 1$. Soit $\Lambda$
un anneau noeth\'erien. Soit $\mathcal{F}$ un faisceau localement
constant de type fini de $\Lambda$-modules sur $S_\etale$
qui est annul\'e par $\ell^n$. Alors il existe un
morphisme fini \'etale $f : T \to S$ de degr\'e premier \`a $\ell$
tel que $f^{-1}\mathcal{F}$ admette une filtration finie dont les
quotients successifs sont de la forme $\underline{M}_T$
pour des $\Lambda$-modules de type fini $M$.
\end{lemma}

\begin{proof}
Choisissons un point g\'eom\'etrique $\overline{s}$ de $S$.
Par l'\'equivalence du
lemme \ref{lemma-locally-constant-on-connected-geom-unibranch},
le faisceau $\mathcal{F}$ correspond \`a un $\Lambda$-module de type fini
$M$ muni d'une action continue de $\pi_1(S, \overline{s})$.
Soit $G \subset \text{Aut}(M)$ l'image de l'homomorphisme
$\rho : \pi_1(S, \overline{s}) \to \text{Aut}(M)$ qui donne l'action.
Remarquons que $G$ est fini par continuit\'e et parce que $M$ est engendr\'e
par un nombre fini d'\'el\'ements comme $\Lambda$-module (voir la preuve du lemme \ref{lemma-locally-constant-on-connected-geom-unibranch}).
L'homomorphisme continu surjectif
$\overline{\rho} : \pi_1(S, \overline{s}) \to G$
correspond \`a un objet galoisien $Y \to S$ de
$\textit{F\'Et}_S$, de groupe d'automorphismes $G = \text{Aut}(Y/S)$, voir
Fundamental Groups, section \ref{pione-section-finite-etale-under-galois}.
Soit $H \subset G$ un $\ell$-sous-groupe de Sylow.
Nous affirmons que $T = Y/H \to S$ convient. En effet, soit $\overline{t} \in T$
un point g\'eom\'etrique au-dessus de $\overline{s}$. L'image de
$\pi_1(T, \overline{t}) \to \pi_1(S, \overline{s})$
est $(\overline{\rho})^{-1}(H)$ par fonctorialit\'e
des groupes fondamentaux. Par cons\'equent, l'action de $\pi_1(T, \overline{t})$
sur $M$ correspondant \`a $f^{-1}\mathcal{F}$ se fait par
l'application $\pi_1(T, \overline{t}) \to H$, voir la
remarque \ref{remark-functorial-locally-constant-on-connected}.
Soit $0 = M_0 \subset M_1 \subset \ldots \subset M_t = M$
comme dans le lemme \ref{lemma-action-l-group}.
Cela induit une filtration
$0 = \mathcal{F}_0 \subset \mathcal{F}_1 \subset \ldots
\subset \mathcal{F}_t = f^{-1}\mathcal{F}$
telle que les quotients successifs soient constants de
valeur $M_{i + 1}/M_i$.
Enfin, le degr\'e de $T = Y/H \to S$ est premier \`a $\ell$,
car il est \'egal \`a l'indice de $H$ dans $G$.
\end{proof}






\section{Cohomologie galoisienne}
\label{section-galois-cohomology}

\noindent
Dans cette section, nous d\'emontrons un r\'esultat de cohomologie galoisienne
(proposition \ref{proposition-serre-galois})
au moyen de la cohomologie \'etale et de l'argument de la
section \ref{section-trace-method}.
Cela nous permettra de d\'emontrer l'annulation des groupes de cohomologie \'etale sup\'erieure
sur le spectre d'un corps.

\begin{lemma}
\label{lemma-nonvanishing-inherited}
Soit $\ell$ un nombre premier et $n$ un entier $> 0$.
Soit $S$ un sch\'ema quasi-compact et quasi-s\'epar\'e.
Soit $X = \lim_{i \in I} X_i$ la limite projective d'un
syst\`eme projectif filtrant de $S$-sch\'emas, chaque $X_i \to S$
\'etant fini \'etale de degr\'e constant premier \`a $\ell$.
Les assertions suivantes sont \`equivalentes :
\begin{enumerate}
\item il existe un faisceau de torsion $\ell$-primaire
$\mathcal{G}$ sur $S$ tel que $H_\etale^n(S, \mathcal{G}) \neq 0$, et
\item il existe un faisceau de torsion $\ell$-primaire $\mathcal{F}$ sur $X$
tel que $H_\etale^n(X, \mathcal{F}) \neq 0$.
\end{enumerate}
En fait, \`a partir de
$\mathcal{G}$ nous pouvons prendre $\mathcal{F} = g^{-1}\mathcal{G}$,
et \`a partir de
$\mathcal{F}$ nous pouvons prendre $\mathcal{G} = g_*\mathcal{F}$.
\end{lemma}

\begin{proof}
Notons $g : X \to S$ et $g_i : X_i \to S$ les morphismes structuraux.
Fixons un faisceau de torsion $\ell$-primaire $\mathcal{G}$ sur $S$
tel que $H^n_\etale(S, \mathcal{G}) \not = 0$.
Le syst\`eme form\'e par $\mathcal{G}_i = g_i^{-1}\mathcal{G}$
satisfait aux conditions du th\'eor\`eme \ref{theorem-colimit}
et son faisceau limite inductive est $g^{-1}\mathcal{G}$. Nous en
d\'eduisons que :
$$
\colim_{i\in I} H^n_\etale(X_i, g_i^{-1}\mathcal{G}) =
H^n_\etale(X, g^{-1}\mathcal{G})
$$
Comme les $g_i$ sont des morphismes finis \'etales de degr\'e premier
\`a $\ell$, nous pouvons appliquer la « m\'ethode de la trace » et obtenons que
les applications
$$
H^n_\etale(S, \mathcal{G}) \to H^n_\etale(X_i, g_i^{-1}\mathcal{G})
$$
sont toutes injectives (et compatibles aux applications de transition).
Voir la section \ref{section-trace-method}. La limite inductive est donc non nulle, c'est-\`a-dire
$H^n(X,g^{-1}\mathcal{G}) \neq 0$, ce qui donne le r\'esultat voulu avec
$\mathcal{F} = g^{-1}\mathcal{G}$.

\medskip\noindent
R\'eciproquement, soit $\mathcal{F}$ un faisceau de torsion $\ell$-primaire sur $X$
tel que $H^n_\etale(X, \mathcal{F}) \not = 0$. Comme les $g_i$
sont des morphismes finis, leurs images directes sup\'erieures sont nulles
(proposition \ref{proposition-finite-higher-direct-image-zero}).
En appliquant alors le lemme \ref{lemma-relative-colimit},
nous pouvons conclure qu'il en va de m\^eme pour $g$.
L'annulation des images directes sup\'erieures donne
$H^n_\etale(X, \mathcal{F}) = H^n(S, g_*\mathcal{F}) \neq 0$
par Leray (proposition \ref{proposition-leray}),
ce qui donne le r\'esultat voulu avec $\mathcal{G} = g_*\mathcal{F}$.
\end{proof}

\begin{lemma}
\label{lemma-reduce-to-l-group}
Soit $\ell$ un nombre premier et $n$ un entier $> 0$.
Soit $K$ un corps, posons $G = \text{Gal}(K^{sep}/K)$, et soit
$H \subset G$ un $\ell$-sous-groupe de Sylow, l'extension de corps
correspondante \'etant not\'ee $L/K$. Alors
$H^n_\etale(\Spec(K), \mathcal{F}) = 0$ pour tout faisceau
de torsion $\ell$-primaire $\mathcal{F}$ si et seulement si
$H^n_\etale(\Spec(L), \underline{\mathbf{Z}/\ell\mathbf{Z}}) = 0$.
\end{lemma}

\begin{proof}
\'Ecrivons $L = \bigcup L_i$ comme r\'eunion de ses sous-extensions finies sur $K$.
Le choix de $H$ implique que $[L_i : K]$ est premier \`a $\ell$.
Ainsi $\Spec(L) = \lim_{i \in I} \Spec(L_i)$ comme dans le
lemme \ref{lemma-nonvanishing-inherited}.
Nous pouvons donc remplacer $K$ par $L$ et supposer que
le groupe de Galois absolu $G$ de $K$ est un
pro-$\ell$-groupe profini.

\medskip\noindent
Supposons $H^n(\Spec(K), \underline{\mathbf{Z}/\ell\mathbf{Z}}) = 0$.
Soit $\mathcal{F}$ un faisceau de torsion $\ell$-primaire sur $\Spec(K)_\etale$.
Nous allons montrer que $H^n_\etale(\Spec(K), \mathcal{F}) = 0$.
Par la correspondance du
lemme \ref{lemma-equivalence-abelian-sheaves-point},
notre faisceau $\mathcal{F}$ correspond \`a un $G$-module de torsion
$\ell$-primaire $M$. Par continuit\'e, tout ensemble fini d'\'el\'ements $x_1, \ldots, x_m \in M$
est fix\'e par un sous-groupe ouvert $U$.
Soit $M'$ le sous-module engendr\'e par les orbites de $x_1, \ldots, x_m$.
C'est un $\ell$-groupe ab\'elien fini,
car chaque $x_i$ est annul\'e par une puissance de $\ell$
et les orbites sont finies. Comme $M$ est la limite inductive filtrante de
ces sous-modules $M'$, nous voyons que $\mathcal{F}$ est la limite
inductive filtrante des sous-faisceaux correspondants $\mathcal{F}' \subset \mathcal{F}$.
En appliquant le th\'eor\`eme \ref{theorem-colimit} \`a cette limite inductive, nous nous ramenons
au cas o\`u $\mathcal{F}$ est un faisceau localement constant fini.

\medskip\noindent
Soit $M$ un $\ell$-groupe ab\'elien fini muni d'une action continue
du pro-$\ell$-groupe profini $G$. Il existe alors une filtration
stable par $G$
$$
0 = M_0 \subset M_1 \subset \ldots \subset M_r = M
$$
telle que $M_{i + 1}/M_i \cong \mathbf{Z}/\ell \mathbf{Z}$, avec
action triviale de $G$ (c'est un lemme \'el\'ementaire de th\'eorie des repr\'esentations
des groupes finis).
Le faisceau correspondant $\mathcal{F}$ admet donc une filtration
$$
0 = \mathcal{F}_0 \subset \mathcal{F}_1 \subset \ldots \subset
\mathcal{F}_r = \mathcal{F}
$$
dont les quotients successifs sont isomorphes \`a
$\underline{\mathbf{Z}/\ell \mathbf{Z}}$.
Nous concluons alors par r\'ecurrence et par la suite exacte longue de
cohomologie.
\end{proof}

\begin{lemma}
\label{lemma-reduce-to-l-group-higher}
Soit $\ell$ un nombre premier et $n$ un entier $> 0$.
Soit $K$ un corps, posons $G = \text{Gal}(K^{sep}/K)$, et soit
$H \subset G$ un $\ell$-sous-groupe de Sylow,
l'extension de corps correspondante \'etant not\'ee $L/K$.
Alors $H^q_\etale(\Spec(K),\mathcal{F}) = 0$ pour $q \geq n$ et tout
faisceau de torsion $\ell$-primaire $\mathcal{F}$ si et seulement si
$H^n_\etale(\Spec(L), \underline{\mathbf{Z}/\ell\mathbf{Z}}) = 0$.
\end{lemma}

\begin{proof}
L'implication directe est imm\'ediate; il suffit donc de d\'emontrer la r\'eciproque.
Nous raisonnons par r\'ecurrence sur $q$. Le cas $q = n$ est le
lemme \ref{lemma-reduce-to-l-group}. Soit maintenant
$\mathcal{F}$ un faisceau de torsion $\ell$-primaire sur $\Spec(K)$.
Soit $f : \Spec(K^{sep}) \rightarrow \Spec(K)$
le morphisme d'inclusion d'un point g\'eom\'etrique.
Consid\'erons alors la suite exacte :
$$
0 \rightarrow \mathcal{F} \xrightarrow{res}
f_* f^{-1} \mathcal{F} \rightarrow f_* f^{-1} \mathcal{F}/\mathcal{F} 
\rightarrow 0
$$
Remarquons que $K^{sep}$ s'\'ecrit comme limite inductive filtrante d'extensions
s\'eparables finies. Ainsi $f$
est la limite projective d'un syst\`eme projectif de morphismes finis \'etales.
Comme nous l'avons vu dans la d\'emonstration du
lemme \ref{lemma-nonvanishing-inherited}, nous pouvons conclure que les
images directes sup\'erieures de $f$ sont nulles. La cohomologie sup\'erieure de
$f_* f^{-1} \mathcal{F}$ s'identifie donc \`a celle du point g\'eom\'etrique, qui
est manifestement nulle. Comme tout reste de torsion $\ell$-primaire, nous pouvons appliquer
l'hypoth\`ese de r\'ecurrence conjointement avec la suite exacte longue de cohomologie
pour conclure au rang $q + 1$.
\end{proof}

\begin{proposition}
\label{proposition-serre-galois}
\begin{reference}
\cite[Chapter II, Section 3, Proposition 5]{SerreGaloisCohomology}
\end{reference}
Soit $K$ un corps muni d'une clôture séparable $K^{sep}$.
Supposons que, pour toute extension finie $K'$ de $K$, on ait
$\text{Br}(K') = 0$. Alors
\begin{enumerate}
\item $H^q(\text{Gal}(K^{sep}/K), (K^{sep})^*) = 0$
pour tout $q \geq 1$, et
\item $H^q(\text{Gal}(K^{sep}/K), M) = 0$
pour tout $\text{Gal}(K^{sep}/K)$-module de torsion $M$ et tout $q \geq 2$,
\end{enumerate}
\end{proposition}

\begin{proof}
Posons $p = \text{char}(K)$.
D'après le lemme \ref{lemma-compare-cohomology-point},
le théorème \ref{theorem-brauer-delta}
et l'exemple \ref{example-sheaves-point},
la proposition revient à montrer que, si
$H^2(\Spec(K'),\mathbf{G}_m|_{\Spec(K')_\etale}) = 0$
pour toute extension finie $K'/K$, alors :
\begin{itemize}
\item $H^q(\Spec(K),\mathbf{G}_m|_{\Spec(K)_\etale}) = 0$
pour tout $q \geq 1$, et
\item $H^q(\Spec(K),\mathcal{F}) = 0$
pour tout faisceau de torsion $\mathcal{F}$ et tout $q \geq 2$.
\end{itemize}
Commençons par démontrer le second point.
Comme $\mathcal{F}$ est un faisceau de torsion, la décomposition en
composantes $\ell$-primaires et la compatibilité de la cohomologie aux
limites inductives (ici, aux sommes directes; voir le théorème \ref{theorem-colimit})
nous ramènent à établir $H^q(\Spec(K),\mathcal{F}) = 0$, $q \geq 2$,
pour tout faisceau de torsion $\ell$-primaire et tout nombre premier $\ell$. Nous
pouvons ainsi étudier séparément chaque nombre premier.

\medskip\noindent
Supposons que $\ell \neq p$. Pour toute extension $K'/K$,
considérons la suite exacte de Kummer (lemme \ref{lemma-kummer-sequence})
$$
0 \to
\mu_{\ell, \Spec{K'}} \to
\mathbf{G}_{m, \Spec{K'}} \xrightarrow{(\cdot)^{\ell}}
\mathbf{G}_{m, \Spec{K'}} \to 0
$$
L'hypothèse donne $H^q(\Spec{K'},\mathbf{G}_m|_{\Spec(K')_\etale}) = 0$
pour $q = 2$; pour $q = 1$, la même annulation résulte du
théorème \ref{theorem-picard-group} conjointement à
$\Pic(K') = (0)$. La suite exacte longue de cohomologie
permet donc de conclure que $H^2(\Spec{K'}, \mu_\ell) = 0$ pour
toute extension finie $K'/K$. Soit maintenant $H$ un $\ell$-sous-groupe de Sylow
du groupe de Galois absolu de $K$, et soit $L$ l'extension
correspondante. On peut écrire $L$ comme limite inductive d'extensions finies;
en appliquant le théorème \ref{theorem-colimit} à cette limite, on voit que
$H^2(\Spec(L), \mu_\ell) = 0$.
Le faisceau $\mu_\ell$ est alors nécessairement constant. S'il ne l'était pas, il
existerait une extension galoisienne de degré premier à
$\ell$ de $L$, contrairement à la définition de $L$ (à savoir l'extension
obtenue en adjoignant les racines $\ell$-ièmes de l'unité à $L$).
Le lemme \ref{lemma-reduce-to-l-group-higher}
donne donc le résultat pour $\ell \neq p$.

\medskip\noindent
Supposons maintenant que $\ell = p$. Considérons la
suite exacte d'Artin-Schreier (section \ref{section-artin-schreier})
$$
0 \longrightarrow \underline{\mathbf{Z}/p\mathbf{Z}}_{\Spec{K}} \longrightarrow
\mathbf{G}_{a, \Spec{K}} \xrightarrow{F-1} \mathbf{G}_{a, \Spec{K}} 
\longrightarrow 0
$$
où $F - 1$ est le morphisme $x \mapsto x^p - x$. La cohomologie supérieure
de $\mathbf{G}_{a, \Spec{K}}$ est nulle, d'après la remarque
\ref{remark-special-case-fpqc-cohomology-quasi-coherent} et l'annulation de la
cohomologie supérieure du faisceau structural d'un schéma affine
(Cohomologie des schémas, lemme
\ref{coherent-lemma-quasi-coherent-affine-cohomology-zero}).
Cette observation vaut sur tout corps de
caractéristique $p$. En particulier, elle s'applique à l'extension $L$
correspondant à un $p$-sous-groupe de Sylow $H$. On en déduit
$H^n(\Spec{L}, \underline{\mathbf{Z}/p\mathbf{Z}}_{\Spec{L}}) = 0$
pour $n \geq 2$, d'où le résultat pour $\ell = p$, par le
lemme \ref{lemma-reduce-to-l-group-higher}.

\medskip\noindent
Pour achever la démonstration, il reste à montrer que
$H^q(\text{Gal}(K^{sep}/K), (K^{sep})^*) = 0$ pour tout $q \geq 1$.
Posons $G = \text{Gal}(K^{sep}/K)$ et $M = (K^{sep})^*$,
considéré comme un $G$-module. Nous avons déjà montré ci-dessus que
$H^1(G, M) = 0$ et $H^2(G, M) = 0$. Considérons la suite exacte
$$
0 \to A \to M \to M \otimes \mathbf{Q} \to B \to 0
$$
de $G$-modules. Ce qui précède donne $H^i(G, A) = 0$
et $H^i(G, B) = 0$ pour $i > 1$, car $A$ et $B$ sont des
$G$-modules de torsion. D'après le
lemme \ref{lemma-profinite-group-cohomology-torsion},
on a $H^i(G, M \otimes \mathbf{Q}) = 0$ pour $i > 0$.
Il est facile de vérifier que cela entraîne également
$H^i(G, M) = 0$ pour $i \geq 3$.
\end{proof}

%10.08.09

\begin{definition}
\label{definition-Cr}
Un corps $K$ est dit {\it $C_r$}
si, pour tout $0 < d^r < n$ et tout $f \in K[T_1,
\ldots, T_n]$ homogène de degré $d$, il existe $\alpha = (\alpha_1,
\ldots, \alpha_n)$, les $\alpha_i \in K$ n'étant pas tous nuls, tel que $f(\alpha) = 0$.
Un tel $\alpha$ est appelé une {\it solution non triviale} de $f$.
\end{definition}

\begin{example}
\label{example-algebraically-closed-field-Cr}
Tout corps algébriquement clos est $C_r$.
\end{example}

\noindent
En fait, on a le lemme élémentaire suivant.

\begin{lemma}
\label{lemma-algebraically-closed-find-solutions}
Soit $k$ un corps algébriquement clos. Soient
$f_1, \ldots, f_s \in k[T_1, \ldots, T_n]$
des polynômes homogènes de degrés $d_1, \ldots, d_s$, avec $d_i
> 0$. Si $s < n$, les équations $f_1 = \ldots = f_s = 0$ ont une solution
commune non triviale.
\end{lemma}

\begin{proof}
Cela résulte de la théorie de la dimension, par exemple sous la forme du
lemme \ref{varieties-lemma-intersection-in-affine-space} de Variétés,
appliqué $s - 1$ fois.
\end{proof}

\noindent
Le résultat suivant calcule le groupe de Brauer des corps $C_1$.

\begin{theorem}
\label{theorem-C1-brauer-group-zero}
Soit $K$ un corps $C_1$. Alors $\text{Br}(K) = 0$.
\end{theorem}

\begin{proof}
Soit $D$ un corps gauche de dimension finie sur $K$, de centre $K$. Nous
avons vu que
$$
D \otimes_K K^{sep} \cong \text{Mat}_d(K^{sep})
$$
à automorphisme intérieur près. Le déterminant $\det :
\text{Mat}_d(K^{sep}) \to K^{sep}$ est donc invariant par Galois et se descend en une
application homogène de degré $d$
$$
\det = N_\text{red} : D \longrightarrow K
$$
appelée la {\it norme réduite}. Comme $K$ est $C_1$, si $d > 1$, il existe
un élément non nul $x \in D$ tel que $N_\text{red}(x) = 0$. Cela implique clairement
que $x$ n'est pas inversible, contradiction. Par conséquent, $\text{Br}(K) = 0$.
\end{proof}

\begin{definition}
\label{definition-variety}
Soit $k$ un corps. Une {\it variété} est un schéma intègre, séparé et de
type fini sur $k$. Une {\it courbe} est une variété de dimension $1$.
\end{definition}

\begin{theorem}[Théorème de Tsen]
\label{theorem-tsen}
Le corps des fonctions d'une variété de dimension $r$ sur un corps algébriquement clos
$k$ est $C_r$.
\end{theorem}

\begin{proof}
Pour l'espace projectif, on peut montrer directement que le corps
$k(x_1, \ldots, x_r)$ est $C_r$ (exercice).

\medskip\noindent
Cas général. Sans perte de généralité, on peut supposer $X$ projectif.
Soit $f \in k(X)[T_1, \ldots, T_n]_d$ tel que $0 < d^r < n$.
Supposons que les coefficients de $f$ appartiennent à $\Gamma(X, \mathcal{O}_X(H))$
pour un diviseur ample $H \subset X$. Posons
$\mathbf{\alpha} = (\alpha_1, \ldots, \alpha_n)$, où $\alpha_i \in \Gamma(X,
\mathcal{O}_X(eH))$. Alors $f(\mathbf{\alpha}) \in \Gamma(X,
\mathcal{O}_X((de + 1)H))$. Considérons le système d'équations $f(\mathbf{\alpha})
=0$. D'après Riemann--Roch asymptotique
(Variétés, proposition \ref{varieties-proposition-asymptotic-riemann-roch}),
il existe un $c > 0$ tel que
\begin{itemize}
\item le nombre de variables est
$n\dim_k \Gamma(X, \mathcal{O}_X(eH)) \sim n e^r c$, et
\item le nombre d'équations est
$\dim_k \Gamma(X, \mathcal{O}_X((de + 1)H)) \sim (de + 1)^r c$.
\end{itemize}
Comme $n > d^r$, il y a plus de variables que d'équations. Celles-ci sont
homogènes; il existe donc une solution d'après le
lemme \ref{lemma-algebraically-closed-find-solutions}.
\end{proof}

\begin{lemma}
\label{lemma-curve-brauer-zero}
Soit $C$ une courbe sur un corps algébriquement clos $k$. Alors
le groupe de Brauer du corps des fonctions de $C$ est nul :
$\text{Br}(k(C)) = 0$.
\end{lemma}

\begin{proof}
Cela résulte immédiatement du théorème de Tsen
\ref{theorem-tsen} et du
théorème \ref{theorem-C1-brauer-group-zero}.
\end{proof}

\begin{lemma}
\label{lemma-cohomology-Gm-function-field-curve}
Soient $k$ un corps algébriquement clos et $K/k$ une extension de corps
de degré de transcendance 1. Alors, pour tout $q \geq 1$,
$H_\etale^q(\Spec(K), \mathbf{G}_m) = 0$.
\end{lemma}

\begin{proof}
Rappelons que
$H_\etale^q(\Spec(K), \mathbf{G}_m) = H^q(\text{Gal}(K^{sep}/K), (K^{sep})^*)$
d'après le lemme \ref{lemma-compare-cohomology-point}.
La proposition \ref{proposition-serre-galois} montre donc
qu'il suffit d'établir que, si $K'/K$ est une extension finie de corps, alors
$\text{Br}(K') = 0$. Or $K' = \colim K''$, où $K''$ parcourt
les extensions de type fini de $k$ contenues comme sous-corps dans $K'$ et de
degré de transcendance $1$.
On a $\text{Br}(K') = \colim \text{Br}(K'')$, ce qui nous ramène
à une extension de corps de type fini $K''/k$ de degré de transcendance
$1$. Un tel corps est le corps des fonctions d'une courbe sur $k$,
dont le groupe de Brauer est donc nul d'après le
lemme \ref{lemma-curve-brauer-zero}.
\end{proof}






\section{Annulation en degrés supérieurs pour le groupe multiplicatif}
\sectionmark{Annulation pour le groupe multiplicatif}
\label{section-higher-Gm}

\noindent
Dans cette section, on fixe un corps algébriquement clos $k$ et une courbe lisse
$X$ sur $k$. On note $i_x : x \hookrightarrow X$ l'inclusion d'un point fermé
de $X$ et $j : \eta \hookrightarrow X$ l'inclusion du point générique.
On note aussi $X_0$ l'ensemble des points fermés de $X$.

\begin{theorem}[La suite exacte fondamentale]
\label{theorem-fundamental-exact-sequence}
Il existe une suite exacte courte de faisceaux étales sur $X$
$$
0 \longrightarrow
\mathbf{G}_{m, X} \longrightarrow
j_* \mathbf{G}_{m, \eta} \longrightarrow
\bigoplus\nolimits_{x \in X_0} {i_x}_* \underline{\mathbf{Z}}
\longrightarrow 0.
$$
\end{theorem}

\begin{proof}
Soit $\varphi : U \to X$ un morphisme étale. D'après les propriétés des
morphismes étales (proposition \ref{proposition-etale-morphisms}),
$U = \coprod_i U_i$, où chaque $U_i$ est une courbe lisse munie d'un morphisme vers $X$.
La suite ci-dessus pour $U$ est le produit des suites correspondantes
pour les $U_i$; il suffit donc de traiter le cas où $U$ est connexe,
donc irréductible. Dans ce cas, on a la suite exacte bien connue
$$
1 \longrightarrow
\Gamma(U, \mathcal{O}_U^*) \longrightarrow
k(U)^* \longrightarrow \bigoplus\nolimits_{y \in U_0} \mathbf{Z}_y.
$$
Cela revient à une suite
$$
0 \longrightarrow \Gamma(U, \mathcal{O}_U^*) \longrightarrow
\Gamma(\eta \times_X U, \mathcal{O}_{\eta \times_X U}^*) \longrightarrow
\bigoplus\nolimits_{x \in X_0}
\Gamma(x \times_X U, \underline{\mathbf{Z}})
$$
qui, en développant les définitions, n'est autre que la suite
$$
0 \longrightarrow \mathbf{G}_m(U) \longrightarrow
j_* \mathbf{G}_{m, \eta}(U) \longrightarrow
\left(\bigoplus\nolimits_{x \in X_0} {i_x}_* \underline{\mathbf{Z}}\right)(U).
$$
Cette construction définit les applications de la suite exacte fondamentale et montre
qu'elle est exacte, sauf peut-être au dernier terme. Pour voir la surjectivité, rappelons que si
$U$ est une courbe non singulière et $D$ un diviseur sur $U$,
il existe un recouvrement ouvert de Zariski $\{U_j \to U\}$ de $U$
tel que $D|_{U_j} = \text{div}(f_j)$ pour un $f_j \in k(U)^*$.
\end{proof}

\begin{lemma}
\label{lemma-higher-direct-jstar-Gm}
Pour tout $q \geq 1$, $R^q j_*\mathbf{G}_{m, \eta} = 0$.
\end{lemma}

\begin{proof}
Il faut montrer que $(R^q j_*\mathbf{G}_{m, \eta})_{\bar x} = 0$ pour tout
point géométrique $\bar x$ de $X$.

\medskip\noindent
Supposons que $\bar x$ soit au-dessus d'un point fermé $x$ de $X$. Soit $\Spec(A)$
un voisinage ouvert affine de $x$ dans $X$, et soit $K$ le corps des fractions
de $A$. Alors
$$
\Spec(\mathcal{O}^{sh}_{X, \bar x}) \times_X \eta =
\Spec(\mathcal{O}^{sh}_{X, \bar x} \otimes_A K).
$$
L'anneau $\mathcal{O}^{sh}_{X, \bar x} \otimes_A K$ est un localisé de
l'anneau de valuation discrète $\mathcal{O}^{sh}_{X, \bar x}$; c'est donc soit
encore $\mathcal{O}^{sh}_{X, \bar x}$, soit son corps des fractions
$K^{sh}_{\bar x}$. Mais comme une uniformisante locale y devient inversible, c'est
nécessairement le second. Par conséquent,
$$
(R^q j_*\mathbf{G}_{m, \eta})_{(X, \bar x)} = H_\etale^q(\Spec
K^{sh}_{\bar x}, \mathbf{G}_m).
$$
Rappelons que $\mathcal{O}^{sh}_{X, \bar x} =
\colim_{(U, \bar u) \to \bar x} \mathcal{O}(U) = \colim_{A \subset B} B$
où $A \to B$ est étale; ainsi $K^{sh}_{\bar x}$ est une extension algébrique
de $K = k(X)$, et on peut appliquer le
lemme \ref{lemma-cohomology-Gm-function-field-curve}
pour obtenir l'annulation.

\medskip\noindent
Supposons que $\bar x = \bar \eta$ soit au-dessus du point générique $\eta$ de $X$ (en
fait, ce cas est superflu). Alors
$\mathcal{O}^{sh}_{X, \bar \eta} = \kappa(\eta)^{sep}$ et donc
\begin{eqnarray*}
(R^q j_*\mathbf{G}_{m, \eta})_{\bar \eta}
& = &
H_\etale^q(\Spec(\kappa(\eta)^{sep}) \times_X \eta, \mathbf{G}_m) \\
& = & H_\etale^q (\Spec(\kappa(\eta)^{sep}), \mathbf{G}_m) \\
& = & 0 \ \ \text{ pour } q \geq 1
\end{eqnarray*}
puisque le groupe de Galois correspondant est trivial.
\end{proof}

\begin{lemma}
\label{lemma-cohomology-jstar-Gm}
Pour tout $p \geq 1$, $H_\etale^p(X, j_*\mathbf{G}_{m, \eta}) = 0$.
\end{lemma}

\begin{proof}
La suite spectrale de Leray s'écrit
$$
E_2^{p, q} = H_\etale^p(X, R^qj_*\mathbf{G}_{m, \eta}) \Rightarrow
H_\etale^{p+q}(\eta, \mathbf{G}_{m, \eta}),
$$
Le membre de droite est nul pour $p + q \geq 1$ d'après le
lemme \ref{lemma-cohomology-Gm-function-field-curve}.
Le membre de gauche est nul pour $q \geq 1$ d'après le
lemme \ref{lemma-higher-direct-jstar-Gm}. On obtient donc l'annulation voulue.
\end{proof}

\begin{lemma}
\label{lemma-cohomology-istar-Z}
Pour tout $q \geq 1$, $H_\etale^q(X, \bigoplus_{x \in X_0} {i_x}_*
\underline{\mathbf{Z}}) = 0$.
\end{lemma}

\begin{proof}
Comme $X$ est quasi-compact et quasi-séparé, la cohomologie commute aux limites inductives;
il suffit donc de montrer l'annulation de $H_\etale^q(X, {i_x}_*
\underline{\mathbf{Z}})$. Mais l'inclusion $i_x$ d'un point fermé est
finie, donc $R^p {i_x}_* \underline{\mathbf{Z}} = 0$ pour tout $p \geq 1$ d'après la
proposition \ref{proposition-finite-higher-direct-image-zero}.
La suite spectrale de Leray donne alors
$H_\etale^q(X, {i_x}_* \underline{\mathbf{Z}}) =
H_\etale^q(x, \underline{\mathbf{Z}})$.
Enfin, comme $x$ est le spectre d'un
corps algébriquement clos, toute cohomologie supérieure de $x$ est nulle.
\end{proof}

\noindent
Cette série de lemmes donne le résultat suivant.

\begin{theorem}
\label{theorem-vanishing-cohomology-Gm-curve}
Soit $X$ une courbe lisse sur un corps algébriquement clos. Alors
$$
H_\etale^q(X, \mathbf{G}_m) = 0 \ \ \text{ pour tout } q \geq 2.
$$
\end{theorem}

\begin{proof}
Voir la discussion ci-dessus.
\end{proof}

\noindent
On obtient également la suite exacte longue de cohomologie
$$
0 \to
H_\etale^0(X, \mathbf{G}_m) \to
H_\etale^0(X, j_*\mathbf{G}_{m, \eta}) \to
H_\etale^0(X, \bigoplus\nolimits_{x \in X_0} {i_x}_*\underline{\mathbf{Z}}) \to
H_\etale^1(X, \mathbf{G}_m) \to 0
$$
bien qu'il s'agisse de la suite familière
$$
0 \to H_{Zar}^0(X, \mathcal{O}_X^*) \to k(X)^* \to \text{Div}(X)
\to \Pic(X) \to 0.
$$





\section{Groupes de Picard des courbes}
\label{section-pic-curves}

\noindent
L'étape suivante consiste à utiliser la suite de Kummer pour déduire des renseignements
sur le groupe de cohomologie d'une courbe à coefficients finis. Afin
d'obtenir une annulation dans la suite exacte longue, rappelons quelques faits sur les
groupes de Picard.

\medskip\noindent
Soit $X$ une courbe projective lisse sur un corps algébriquement clos $k$.
Soit $g = \dim_k H^1(X, \mathcal{O}_X)$ le genre de $X$.
Il existe une suite exacte courte
$$
0 \to \Pic^0(X) \to \Pic(X) \xrightarrow{\deg} \mathbf{Z} \to 0.
$$
Le groupe abélien $\Pic^0(X)$ s'identifie à
$\Pic^0(X) = \underline{\Picardfunctor}^0_{X/k}(k)$, c'est-à-dire
au groupe des points rationnels sur $k$ d'une variété abélienne
$\underline{\Picardfunctor}^0_{X/k}$ sur $k$ de dimension $g$.
Par conséquent, si $n \in k^*$, alors
$\Pic^0(X)[n] \cong (\mathbf{Z}/n\mathbf{Z})^{2g}$
en tant que groupes abéliens. Voir
Schémas de Picard des courbes, section \ref{pic-section-picard-curve},
et
Groupoïdes, section \ref{groupoids-section-abelian-varieties}.
Ce fait essentiel, à savoir la description de la torsion dans le
groupe de Picard d'une courbe projective lisse sur un corps
algébriquement clos, ne semble pas admettre de démonstration élémentaire.

\begin{lemma}
\label{lemma-cohomology-smooth-projective-curve}
Soit $X$ une courbe projective lisse de genre $g$ sur un
corps algébriquement clos $k$, et soit $n \geq 1$ inversible dans $k$.
Alors on a les identifications canoniques
$$
H_\etale^q(X, \mu_n) =
\left\{
\begin{matrix}
\mu_n(k) & \text{ si }q = 0, \\
\Pic^0(X)[n] & \text{ si }q = 1, \\
\mathbf{Z}/n\mathbf{Z} & \text{ si }q = 2, \\
0 & \text{ si }q \geq 3.
\end{matrix}
\right.
$$
Comme $\mu_n \cong \underline{\mathbf{Z}/n\mathbf{Z}}$, on en déduit des
identifications (non canoniques)
$$
H_\etale^q(X, \underline{\mathbf{Z}/n\mathbf{Z}}) \cong
\left\{
\begin{matrix}
\mathbf{Z}/n\mathbf{Z} & \text{ si }q = 0, \\
(\mathbf{Z}/n\mathbf{Z})^{2g} & \text{ si }q = 1, \\
\mathbf{Z}/n\mathbf{Z} & \text{ si }q = 2, \\
0 & \text{ si }q \geq 3.
\end{matrix}
\right.
$$
\end{lemma}

\begin{proof}
Les théorèmes \ref{theorem-picard-group} et
\ref{theorem-vanishing-cohomology-Gm-curve}
déterminent la cohomologie étale de $\mathbf{G}_m$ sur $X$
en fonction du groupe de Picard de $X$.
La suite exacte de Kummer $0\to \mu_{n, X} \to \mathbf{G}_{m, X}
\to \mathbf{G}_{m, X}\to 0$ (lemme \ref{lemma-kummer-sequence})
donne alors la suite exacte longue de cohomologie
$$
\xymatrix{
0 \ar[r] & \mu_n(k) \ar[r] &
k^* \ar[r]^{(\cdot)^n} &
k^* \ar@(rd, ul)[rdllllr] \\
& H_\etale^1(X, \mu_n) \ar[r] &
\Pic(X) \ar[r]^{(\cdot)^n} &
\Pic(X) \ar@(rd, ul)[rdllllr] \\
& H_\etale^2(X, \mu_n) \ar[r] & 0 \ar[r] & 0 \ldots
}
$$
L'application d'élévation à la puissance $n$, $k^* \to k^*$, est surjective puisque $k$ est
algébriquement clos. Il reste donc à calculer le noyau et le conoyau de l'application
$n : \Pic(X) \to \Pic(X)$. Considérons le diagramme commutatif à
lignes exactes
$$
\xymatrix{
0 \ar[r] &
\Pic^0(X) \ar[r] \ar@{>>}[d]^{(\cdot)^n} &
\Pic(X) \ar[r]_-\deg \ar[d]^{(\cdot)^n} &
\mathbf{Z} \ar[r] \ar@{^{(}->}[d]^n & 0 \\
0 \ar[r] &
\Pic^0(X) \ar[r] &
\Pic(X) \ar[r]^-\deg &
\mathbf{Z} \ar[r] & 0
}
$$
Le groupe $\Pic^0(X)$ est le groupe des points rationnels sur $k$ du
schéma en groupes $\underline{\Picardfunctor}^0_{X/k}$; voir
Schémas de Picard des courbes, lemme \ref{pic-lemma-picard-pieces}.
Le même lemme affirme que $\underline{\Picardfunctor}^0_{X/k}$
est une variété abélienne de dimension $g$ sur $k$ au sens de
Groupoïdes, définition \ref{groupoids-definition-abelian-variety}.
L'application verticale de gauche est donc surjective d'après
Groupoïdes, proposition \ref{groupoids-proposition-review-abelian-varieties}.
Le lemme du serpent donne alors les identifications canoniques énoncées
dans le lemme.

\medskip\noindent
Pour obtenir les identifications non canoniques du lemme, il faut
montrer que le noyau de $n : \Pic^0(X) \to \Pic^0(X)$
est isomorphe à $(\mathbf{Z}/n\mathbf{Z})^{\oplus 2g}$.
Cela fait aussi partie de Groupoïdes, proposition
\ref{groupoids-proposition-review-abelian-varieties}.
\end{proof}

\begin{lemma}
\label{lemma-pullback-on-h2-curve}
Soit $\pi : X \to Y$ un morphisme non constant de courbes projectives lisses
sur un corps algébriquement clos $k$, et soit $n \geq 1$ inversible dans $k$.
L'application
$$
\pi^* : H^2_\etale(Y, \mu_n) \longrightarrow H^2_\etale(X, \mu_n)
$$
est la multiplication par le degré de $\pi$.
\end{lemma}

\begin{proof}
L'énoncé a un sens puisque nous avons identifié les deux
groupes de cohomologie $ H^2_\etale(Y, \mu_n)$ et $H^2_\etale(X, \mu_n)$
à $\mathbf{Z}/n\mathbf{Z}$ dans le
lemme \ref{lemma-cohomology-smooth-projective-curve}.
En effet, si $\mathcal{L}$ est un faisceau inversible de degré $1$
sur $Y$, de classe $[\mathcal{L}] \in H^1_\etale(Y, \mathbf{G}_m)$,
alors l'image de $[\mathcal{L}]$ par l'homomorphisme bord est un générateur de
$H^2_\etale(Y, \mu_n)$. Il s'agit ici de l'homomorphisme bord
de la suite exacte longue de cohomologie associée à la suite exacte
de Kummer. Le résultat du lemme découle donc du fait que
le degré du faisceau inversible $\pi^*\mathcal{L}$ sur $X$ vaut $\deg(\pi)$.
Nous omettons quelques détails.
\end{proof}

\begin{lemma}
\label{lemma-vanishing-cohomology-mu-smooth-curve}
Soit $X$ une courbe affine lisse sur un corps algébriquement clos $k$, et
soit $n \in k^*$. Soit $X \subset \overline{X}$ une
compactification projective lisse
(Variétés, remarque \ref{varieties-remark-smooth-projective-compactification}).
Soit $g$ le genre de $\overline{X}$ et soit $r$ le nombre de
points de $\overline{X} \setminus X$. Alors
\begin{enumerate}
\item $H_\etale^0(X, \mu_n) = \mu_n(k)$;
\item $H_\etale^1(X, \mu_n) \cong (\mathbf{Z}/n\mathbf{Z})^{2g+r-1}$; et
\item $H_\etale^q(X, \mu_n) = 0$ pour tout $q \geq 2$.
\end{enumerate}
\end{lemma}

\begin{proof}
\'Ecrivons $X = \overline{X} - \{ x_1, \ldots, x_r\}$. Alors
$\Pic(X) = \Pic(\overline{X})/ R$, où $R$ est le sous-groupe engendré par
$\mathcal{O}_{\overline{X}}(x_i)$, $1 \leq i \leq r$.
Comme $r \geq 1$, on voit que $\Pic^0(\overline{X}) \to \Pic(X)$
est surjective; le groupe $\Pic(X)$ est donc divisible (voir la discussion dans la preuve
du lemme \ref{lemma-cohomology-smooth-projective-curve}).
La suite exacte de Kummer donne (1) et (3). Pour (2), on utilise
\begin{align*}
H_\etale^1(X, \mu_n)
& =
\{(\mathcal L, \alpha)\}/\cong \\
& =
(\{(\bar{\mathcal L}, D, \bar \alpha)\}/\cong)/\tilde{R}
\end{align*}
Pour le premier signe d'égalité et les notations, voir le lemme \ref{lemma-describe-h1-mun}.
Les triplets $(\bar{\mathcal L}, D, \bar \alpha)$ de la seconde ligne
sont formés d'un module inversible $\bar{\mathcal L}$ sur $\overline{X}$
de degré $0$, d'un diviseur $D$ sur $\overline{X}$ de degré $0$ supporté par
$\left\{x_1, \ldots, x_r\right\}$ et d'un isomorphisme
$\bar{\alpha}: \bar{\mathcal L}^{\otimes n} \cong \mathcal{O}_{\bar{X}}(D)$.
Les isomorphismes de triplets sont définis comme dans la discussion qui précède le
lemme \ref{lemma-describe-h1-mun}.
De plus, $\tilde{R}$ est le sous-groupe engendré par les triplets de la forme
$(\mathcal{O}_{\overline{X}}(D'), n D', 1^{\otimes n})$
où $D'$ est supporté par $\left\{x_1, \ldots, x_r\right\}$ et de degré 0.
La discussion ci-dessus montre que le second signe d'égalité est donné par
l'application qui envoie le triplet $(\bar{\mathcal L},\ D,\ \bar \alpha)$
sur la paire $(\bar{\mathcal L}|_X, \bar \alpha|_X)$. On obtient donc
une suite exacte
$$
0 \longrightarrow
H_\etale^1(\overline{X}, \mu_n) \longrightarrow
H_\etale^1(X, \mu_n) \longrightarrow
\bigoplus_{i = 1}^r \mathbf{Z}/n\mathbf{Z}
\xrightarrow{\ \sum\ }
\mathbf{Z}/n\mathbf{Z} \longrightarrow 0
$$
où l'application du milieu envoie la classe d'un triplet
$(\bar{ \mathcal L}, D, \bar \alpha)$ avec $D = \sum_{i = 1}^r a_i (x_i)$
sur le $r$-uplet $(a_i)_{i = 1}^r$. Il suffit alors d'utiliser le
lemme \ref{lemma-cohomology-smooth-projective-curve}
pour compter les rangs.
\end{proof}

\begin{remark}
\label{remark-natural-proof}
La manière ``naturelle'' de démontrer le lemme précédent consiste à appliquer l'excision à l'ouvert $X$ dans $\bar
X$. Cela est possible, mais nous n'avons pas développé cette théorie.
\end{remark}

\begin{remark}
\label{remark-normalize-H1-Gm}
Soit $k$ un corps algébriquement clos. Soit $n$ un entier premier à
la caractéristique de $k$. Rappelons que
$$
\mathbf{G}_{m, k} = \mathbf{A}^1_k \setminus \{0\} =
\mathbf{P}^1_k \setminus \{0, \infty\}
$$
Nous affirmons qu'il existe un isomorphisme canonique
$$
H^1_\etale(\mathbf{G}_{m, k}, \mu_n) = \mathbf{Z}/n\mathbf{Z}
$$
Que signifie cette affirmation? Elle signifie qu'il existe un élément $1_k$ de
$H^1_\etale(\mathbf{G}_{m, k}, \mu_n)$ tel que, pour
tout morphisme $\Spec(k') \to \Spec(k)$, l'application d'image inverse en
cohomologie étale associée à $\mathbf{G}_{m, k'} \to \mathbf{G}_{m, k}$
envoie $1_k$ sur $1_{k'}$. (En particulier, cet élément est
fixé par tout automorphisme de $k$.) Pour le voir, considérons le
torseur sous $\mu_{n, \mathbf{Z}}$
$\mathbf{G}_{m, \mathbf{Z}} \to \mathbf{G}_{m, \mathbf{Z}}$,
$x \mapsto x^n$. Par l'identification des classes de torseurs avec la
cohomologie de degré un, son image inverse fournit nos éléments canoniques $1_k$.
En revenant au faisceau constant, on obtient des identifications canoniques
$$
H^1_\etale(\mathbf{G}_{m, k}, \mathbf{Z}/n\mathbf{Z}) =
\Hom(\mu_n(k), \mathbf{Z}/n\mathbf{Z}),
$$
c'est-à-dire que ces isomorphismes sont compatibles avec les morphismes de
corps algébriquement clos, en particulier aux
automorphismes de $k$.
\end{remark}










\section{Prolongement par zéro}
\label{section-extension-by-zero}

\noindent
Les résultats généraux de
Modules sur les sites, section \ref{sites-modules-section-localize},
permettent de poser la définition suivante.

\begin{definition}
\label{definition-extension-zero}
Soit $j : U \to X$ un morphisme étale de schémas.
\begin{enumerate}
\item Le foncteur de restriction
$j^{-1} : \Sh(X_\etale) \to \Sh(U_\etale)$
admet un adjoint à gauche
$j_!^{Sh} : \Sh(U_\etale) \to \Sh(X_\etale)$.
\item Le foncteur de restriction
$j^{-1} : \textit{Ab}(X_\etale) \to \textit{Ab}(U_\etale)$
admet un adjoint à gauche, noté
$j_! : \textit{Ab}(U_\etale) \to \textit{Ab}(X_\etale)$
et appelé {\it prolongement par zéro}.
\item Soit $\Lambda$ un anneau. Le foncteur de restriction
$j^{-1} : \textit{Mod}(X_\etale, \Lambda) \to
\textit{Mod}(U_\etale, \Lambda)$
admet un adjoint à gauche, noté
$j_! : \textit{Mod}(U_\etale, \Lambda) \to
\textit{Mod}(X_\etale, \Lambda)$
et appelé {\it prolongement par zéro}.
\end{enumerate}
\end{definition}

\noindent
Si $\mathcal{F}$ est un faisceau abélien sur $U_\etale$, alors
$j_!\mathcal{F} \not = j_!^{Sh}\mathcal{F}$ en général. En revanche,
le $j_!$ des faisceaux de $\Lambda$-modules coïncide avec le $j_!$ des
faisceaux abéliens sous-jacents
(Modules sur les sites, remarque \ref{sites-modules-remark-localize-shriek-equal}).
Le foncteur $j_!$ est caractérisé par l'isomorphisme fonctoriel
$$
\Hom_X(j_!\mathcal{F}, \mathcal{G}) = \Hom_U(\mathcal{F}, j^{-1}\mathcal{G})
$$
pour tout $\mathcal{F} \in \textit{Ab}(U_\etale)$ et
$\mathcal{G} \in \textit{Ab}(X_\etale)$. Il en va de même pour les
faisceaux de $\Lambda$-modules.

\medskip\noindent
Pour décrire plus explicitement les foncteurs de la définition \ref{definition-extension-zero},
rappelons que $j^{-1}$ n'est autre que la restriction
le long du foncteur $U_\etale \to X_\etale$. Autrement dit,
$j^{-1}\mathcal{G}(U') = \mathcal{G}(U')$ pour $U'$ étale sur $U$.
Pour $\mathcal{F} \in \textit{Ab}(U_\etale)$, considérons le
préfaisceau
\begin{equation}
\label{equation-j-p-shriek}
j_{p!}\mathcal{F} : X_\etale \longrightarrow \textit{Ab},
\quad
V \longmapsto \bigoplus\nolimits_{V \to U} \mathcal{F}(V \to U)
\end{equation}
Alors $j_!\mathcal{F}$ est le faisceau associé au préfaisceau $j_{p!}\mathcal{F}$.
C'est démontré dans Modules sur les sites, lemme
\ref{sites-modules-lemma-extension-by-zero};
plus généralement, voir la discussion dans
Modules sur les sites, sections \ref{sites-modules-section-localize} et
\ref{sites-modules-section-exactness-lower-shriek}.

\begin{exercise}
\label{exercise-jshriek-direct}
Montrer directement que le foncteur $j_!$, défini comme le faisceau associé au
préfaisceau $j_{p!}$ donné dans
(\ref{equation-j-p-shriek}), est adjoint à gauche de $j^{-1}$.
\end{exercise}

\begin{proposition}
\label{proposition-describe-jshriek}
Soit $j : U \to X$ un morphisme étale de schémas.
Soit $\mathcal{F}$ un objet de $\textit{Ab}(U_\etale)$.
Si $\overline{x} : \Spec(k) \to X$ est un point géométrique de $X$, alors
$$
(j_!\mathcal{F})_{\overline{x}} =
\bigoplus\nolimits_{\overline{u} : \Spec(k) \to U,\ j(\overline{u}) =
\overline{x}} \mathcal{F}_{\bar{u}}.
$$
En particulier, $j_!$ est un foncteur exact.
\end{proposition}

\begin{proof}
L'exactitude de $j_!$ est très générale; voir Modules sur les sites,
lemme \ref{sites-modules-lemma-extension-by-zero-exact}.
Elle résulte bien sûr aussi de la description des fibres.
La formule de la fibre résulte de
Modules sur les sites, lemme \ref{sites-modules-lemma-stalk-j-shriek},
et de la description des points du petit site étale
en termes de points géométriques; voir le lemme \ref{lemma-points-small-etale-site}.

\medskip\noindent
Pour un usage ultérieur, notons que l'isomorphisme
\begin{align*}
(j_!\mathcal{F})_{\overline{x}}
& =
(j_{p!}\mathcal{F})_{\overline{x}} \\
& =
\colim_{(V, \overline{v})} j_{p!}\mathcal{F}(V) \\
& =
\colim_{(V, \overline{v})}
\bigoplus\nolimits_{\varphi : V \to U}
\mathcal{F}(V \xrightarrow{\varphi} U) \\
& \to
\bigoplus\nolimits_{\overline{u} : \Spec(k) \to U,\ j(\overline{u}) =
\overline{x}} \mathcal{F}_{\bar{u}}.
\end{align*}
construit dans Modules sur les sites, lemme \ref{sites-modules-lemma-stalk-j-shriek},
envoie $(V, \overline{v}, \varphi, s)$ sur la classe de $s$ dans la fibre
de $\mathcal{F}$ en $\overline{u} = \varphi(\overline{v})$.
\end{proof}

\begin{lemma}
\label{lemma-jshriek-open}
Soit $j : U \to X$ une immersion ouverte de schémas. Pour tout
faisceau abélien $\mathcal{F}$ sur $U_\etale$, les morphismes d'adjonction
$j^{-1}j_*\mathcal{F} \to \mathcal{F}$ et
$\mathcal{F} \to j^{-1}j_!\mathcal{F}$ sont des isomorphismes.
En fait, $j_!\mathcal{F}$ est l'unique faisceau abélien sur $X_\etale$
dont la restriction à $U$ est $\mathcal{F}$ et dont les fibres aux
points géométriques de $X \setminus U$ sont nulles.
\end{lemma}

\begin{proof}
Nous invitons le lecteur à démontrer la première assertion directement à partir des
définitions; nous utilisons ici seulement le fait qu'elle est un cas particulier du très général
lemme \ref{sites-modules-lemma-restrict-back} de Modules sur les sites.
Pour la seconde assertion, observons que si $\mathcal{G}$ est un faisceau abélien
sur $X_\etale$ dont la restriction à $U$ est $\mathcal{F}$, l'adjonction fournit
un morphisme $j_!\mathcal{F} \to \mathcal{G}$. Ce morphisme
induit alors un isomorphisme sur les fibres aux points géométriques de $U$ d'après la
proposition \ref{proposition-describe-jshriek}.
Ainsi, si les fibres de $\mathcal{G}$ aux points géométriques
de $X \setminus U$ sont nulles, alors $j_!\mathcal{F} \to \mathcal{G}$ est un
isomorphisme d'après le théorème \ref{theorem-exactness-stalks}.
\end{proof}

\begin{lemma}[Le prolongement par zéro commute au changement de base]
\label{lemma-shriek-base-change}
Soit $f: Y \to X$ un morphisme de schémas. Soit $j: V \to X$ un morphisme
étale. Considérons le produit fibré
$$
\xymatrix{
V' = Y \times_X V \ar[d]_{f'} \ar[r]_-{j'} & Y \ar[d]^f \\
V \ar[r]^j & X
}
$$
Alors on a $j'_! f'^{-1} = f^{-1} j_!$ sur les faisceaux abéliens et les
faisceaux de modules.
\end{lemma}

\begin{proof}
Cela résulte du fait que $j'_! f'^{-1}$ est adjoint à gauche de
$f'_* (j')^{-1}$, tandis que $f^{-1} j_!$ est adjoint à gauche de $j^{-1}f_*$.
De plus, $f'_* (j')^{-1} = j^{-1}f_*$, car $f_*$ commute à la
localisation étale (par construction). En fait, le lemme vaut très généralement
dans le cadre d'un morphisme de sites; voir
Modules sur les sites, lemme
\ref{sites-modules-lemma-localize-morphism-ringed-sites}.
\end{proof}

\begin{lemma}
\label{lemma-shriek-into-star-separated-etale}
Soit $j : U \to X$ séparé et étale. Il existe alors un morphisme
injectif fonctoriel $j_!\mathcal{F} \to j_*\mathcal{F}$
sur les faisceaux abéliens et les faisceaux de $\Lambda$-modules.
\end{lemma}

\begin{proof}
Démontrons-le dans le cas des faisceaux abéliens. Construisons un morphisme canonique
$$
j_{p!}\mathcal{F} \to j_*\mathcal{F}
$$
de préfaisceaux abéliens sur $X_\etale$, pour tout faisceau abélien
$\mathcal{F}$ sur $U_\etale$, où $j_{p!}$ est celui de
(\ref{equation-j-p-shriek}). Le faisceau associé à ce morphisme sera
le morphisme cherché $j_!\mathcal{F} \to j_*\mathcal{F}$.
En évaluant les deux membres sur $V \to X$ étale, on obtient
$$
j_{p!}\mathcal{F}(V) =
\bigoplus\nolimits_{\varphi : V \to U} \mathcal{F}(V \xrightarrow{\varphi} U)
\quad\text{et}\quad
j_*\mathcal{F}(V) = \mathcal{F}(V \times_X U)
$$
Pour chaque $\varphi$, on a une immersion ouverte et fermée
$$
\Gamma_\varphi = (1, \varphi) : V \longrightarrow V \times_X U
$$
au-dessus de $U$. Elle est ouverte, car c'est un morphisme entre schémas étales
sur $U$, et elle est fermée, car c'est une section d'un schéma séparé
sur $V$ (Schémas, lemme \ref{schemes-lemma-section-immersion}).
Ainsi, pour une section $s_\varphi \in \mathcal{F}(V \xrightarrow{\varphi} U)$,
il existe une unique section $s'_\varphi$ de $\mathcal{F}(V \times_X U)$
dont l'image inverse est $s_\varphi$ par $\Gamma_\varphi$
et dont la restriction au complémentaire de l'image de $\Gamma_\varphi$ est nulle.

\medskip\noindent
Pour montrer que notre morphisme est injectif, supposons que
$\sum_{i = 1}^n s_{\varphi_i}$ soit un élément de
$j_{p!}\mathcal{F}(V)$ dont l'image par la formule ci-dessus est nulle
dans $j_*\mathcal{F}(V)$. Il s'agit de montrer que la restriction de
$\sum_{i = 1}^n s_{\varphi_i}$ est nulle
sur les objets d'un recouvrement étale de $V$.
En considérant tous les égalisateurs deux à deux
(qui sont ouverts et fermés dans $V$) des morphismes
$\varphi_i : V \to U$ et en raisonnant localement sur $V$, on
peut supposer que les images des morphismes
$\Gamma_{\varphi_1}, \ldots, \Gamma_{\varphi_n}$
sont deux à deux disjointes. Comme notre hypothèse est que
$\sum_{i = 1}^n s'_{\varphi_i} = 0$,
on en déduit immédiatement que $s'_{\varphi_i} = 0$
pour tout $i$ (par disjonction des supports de ces
sections), d'où $s_{\varphi_i} = 0$ pour tout $i$, comme voulu.
\end{proof}

\begin{lemma}
\label{lemma-shriek-equals-star-finite-etale}
Soit $j : U \to X$ fini et étale. Alors le morphisme
$j_! \to j_*$ du lemme \ref{lemma-shriek-into-star-separated-etale}
est un isomorphisme
sur les faisceaux abéliens et les faisceaux de $\Lambda$-modules.
\end{lemma}

\begin{proof}
Il suffit de vérifier que $j_!\mathcal{F} \to j_*\mathcal{F}$
est un isomorphisme étale-localement sur $X$.
On peut donc supposer que $U \to X$ est une union disjointe finie
d'isomorphismes; voir
Morphismes étales, lemme \ref{etale-lemma-finite-etale-etale-local}.
Nous omettons la démonstration dans ce cas.
\end{proof}

\begin{lemma}
\label{lemma-ses-associated-to-open}
Soit $X$ un schéma. Soit $Z \subset X$ un sous-schéma fermé, et soit
$U \subset X$ son complémentaire. Notons $i : Z \to X$ et $j : U \to X$
les morphismes d'inclusion. Pour tout faisceau abélien $\mathcal{F}$ sur $X_\etale$,
il existe une suite exacte courte canonique
$$
0 \to j_!j^{-1}\mathcal{F} \to \mathcal{F} \to i_*i^{-1}\mathcal{F} \to 0
$$
sur $X_\etale$.
\end{lemma}

\begin{proof}
Les propriétés d'adjonction des foncteurs en jeu fournissent les
morphismes. Pour un point géométrique $\overline{x}$ de $X$, on a soit
$\overline{x} \in U$, auquel cas le morphisme de gauche induit
un isomorphisme sur les fibres et la fibre de $i_*i^{-1}\mathcal{F}$
est nulle, soit $\overline{x} \in Z$, auquel cas le morphisme de droite induit
un isomorphisme sur les fibres et la fibre de $j_!j^{-1}\mathcal{F}$
est nulle. Nous avons utilisé ici la description des fibres du
lemme \ref{lemma-stalk-pushforward-closed-immersion} et de la
proposition \ref{proposition-describe-jshriek}.
\end{proof}

\begin{lemma}
\label{lemma-compatible-shriek-push-finite}
Considérons un diagramme cartésien de schémas
$$
\xymatrix{
U \ar[d]_g \ar[r]_{j'} & X \ar[d]^f \\
V \ar[r]^j & Y
}
$$
où $f$ est fini et $j$ est une immersion ouverte.
Alors $f_* \circ j'_! = j_! \circ g_*$ comme foncteurs
$\textit{Ab}(U_\etale) \to \textit{Ab}(Y_\etale)$.
\end{lemma}

\begin{proof}
Soit $\mathcal{F}$ un objet de $\textit{Ab}(U_\etale)$.
Soit $\overline{y}$ un point géométrique de $Y$ qui n'appartient pas
à l'ouvert $V$. Alors
$$
(f_*j'_!\mathcal{F})_{\overline{y}} =
\bigoplus\nolimits_{\overline{x},\ f(\overline{x}) = \overline{y}}
(j'_!\mathcal{F})_{\overline{x}} = 0
$$
d'après la proposition \ref{proposition-finite-higher-direct-image-zero}, et
parce que la fibre de $j'_!\mathcal{F}$ en $\overline{x} \not \in U$
est nulle d'après le lemme \ref{lemma-jshriek-open}. D'autre part, on a
$$
j^{-1}f_*j'_!\mathcal{F} =
g_*(j')^{-1}j'_!\mathcal{F} =
g_*\mathcal{F}
$$
d'après les lemmes \ref{lemma-finite-pushforward-commutes-with-base-change} et
\ref{lemma-jshriek-open}.
La caractérisation de
$j_!$ dans le lemme \ref{lemma-jshriek-open} montre donc que
$f_*j'_!\mathcal{F} = j_!g_*\mathcal{F}$.
Nous omettons la vérification que cette identification
est fonctorielle en $\mathcal{F}$.
\end{proof}







\section{Faisceaux constructibles}
\label{section-constructible}

\noindent
Soit $X$ un schéma. Un {\it sous-schéma constructible localement fermé} de $X$
est un sous-schéma localement fermé $T \subset X$ dont l'espace
topologique sous-jacent à $T$ est une partie constructible de $X$.
Si $T, T' \subset X$ sont des sous-schémas localement
fermés ayant le même espace topologique sous-jacent, alors
$T_\etale \cong T'_\etale$ par l'invariance topologique
du site étale (théorème \ref{theorem-topological-invariance}).
Ainsi, dans la définition suivante, nous pouvons supposer réduits nos
sous-schémas localement fermés.

\begin{definition}
\label{definition-constructible}
Soit $X$ un schéma.
\begin{enumerate}
\item Un faisceau d'ensembles sur $X_\etale$ est {\it constructible}
si, pour tout ouvert affine $U \subset X$, il existe une décomposition finie
de $U$ en sous-schémas constructibles localement fermés $U = \coprod_i U_i$
telle que $\mathcal{F}|_{U_i}$ soit localement constant, de valeur finie, pour tout $i$.
\item Un faisceau de groupes abéliens sur $X_\etale$ est {\it constructible}
si, pour tout ouvert affine $U \subset X$, il existe une décomposition finie
de $U$ en sous-schémas constructibles localement fermés $U = \coprod_i U_i$
telle que $\mathcal{F}|_{U_i}$ soit localement constant, de valeur finie, pour tout $i$.
\item Soit $\Lambda$ un anneau noethérien. Un faisceau de $\Lambda$-modules
sur $X_\etale$ est {\it constructible} si, pour tout ouvert affine
$U \subset X$, il existe une décomposition finie
de $U$ en sous-schémas constructibles localement fermés
$U = \coprod_i U_i$ telle que
$\mathcal{F}|_{U_i}$ soit de type fini et localement constant pour tout $i$.
\end{enumerate}
\end{definition}

\noindent
Il semble que ce soit la définition admise. Une autre définition, qui se prête
plus aisément à des généralisations au-delà du site étale d'un schéma,
aurait consisté à définir les faisceaux constructibles en partant de
$h_U$, $j_{U!}\mathbf{Z}/n\mathbf{Z}$ et $j_{U!}\underline{\Lambda}$
où $U$ parcourt tous les objets quasi-compacts et quasi-séparés
de $X_\etale$, puis à prendre la plus petite sous-catégorie pleine de
$\Sh(X_\etale)$, $\textit{Ab}(X_\etale)$ et
$\textit{Mod}(X_\etale, \underline{\Lambda})$ qui les contienne
et soit stable par limites et colimites finies. Il résulte du
lemme \ref{lemma-constructible-abelian}
et des
lemmes \ref{lemma-category-constructible-sets},
\ref{lemma-category-constructible-abelian} et
\ref{lemma-category-constructible-modules}
que cette construction donne la même catégorie si $X$ est quasi-compact et
quasi-séparé. En général, elle ne donne cependant pas la même
catégorie.

\medskip\noindent
Une décomposition en union disjointe $U = \coprod U_i$ d'un schéma en
sous-schémas localement fermés sera appelée une {\it partition} de $U$
(comparer à Topologie, section \ref{topology-section-stratifications}).

\begin{lemma}
\label{lemma-constructible-quasi-compact-quasi-separated}
Soit $X$ un schéma quasi-compact et quasi-séparé. Soit $\mathcal{F}$
un faisceau d'ensembles sur $X_\etale$. Les conditions suivantes sont équivalentes :
\begin{enumerate}
\item $\mathcal{F}$ est constructible,
\item il existe un recouvrement ouvert $X = \bigcup U_i$ tel que
$\mathcal{F}|_{U_i}$ soit constructible, et
\item il existe une partition $X = \bigcup X_i$ en sous-schémas
constructibles localement fermés telle que $\mathcal{F}|_{X_i}$ soit
localement constant, de valeur finie.
\end{enumerate}
Un énoncé analogue vaut pour les faisceaux abéliens et les faisceaux de
$\Lambda$-modules si $\Lambda$ est noethérien.
\end{lemma}

\begin{proof}
Il est clair que (1) implique (2).

\medskip\noindent
Supposons (2). Pour tout $x \in X$, il existe un $i$ et un voisinage ouvert
affine $V_x \subset U_i$ de $x$. On peut donc trouver un recouvrement
ouvert affine fini $X = \bigcup V_j$ tel que, pour tout $j$, il
existe une décomposition finie $V_j = \coprod V_{j, k}$ en parties localement fermées
constructibles telle que $\mathcal{F}|_{V_{j, k}}$ soit localement constant,
de valeur finie. D'après Topologie, lemme
\ref{topology-lemma-quasi-compact-open-immersion-constructible-image}
chaque $V_{j, k}$ est une partie constructible de $X$.
D'après Topologie, lemme
\ref{topology-lemma-constructible-partition-refined-by-stratification}
il existe une stratification finie $X = \coprod X_l$ en strates constructibles
localement fermées telle que chacun des
$V_{j, k}$ soit une réunion de $X_l$. Ainsi, (3) est satisfaite.

\medskip\noindent
Supposons (3) satisfaite. Soit $U \subset X$ un ouvert affine.
Alors $U \cap X_i$ est une partie constructible localement fermée de $U$
(par exemple d'après Propriétés, lemme \ref{properties-lemma-locally-constructible})
et $U = \coprod U \cap X_i$ est une partition de $U$ comme dans la
définition \ref{definition-constructible}. Ainsi, (1) est satisfaite.
\end{proof}

\begin{lemma}
\label{lemma-constructible-constructible}
Soit $X$ un schéma quasi-compact et quasi-séparé.
Soit $\mathcal{F}$ un faisceau d'ensembles, de groupes abéliens ou de
$\Lambda$-modules (avec $\Lambda$ noethérien) sur $X_\etale$.
S'il existe des sous-schémas constructibles localement fermés $T_i \subset X$
tels que (a) $X = \bigcup T_j$ et que (b) $\mathcal{F}|_{T_j}$ soit
constructible, alors $\mathcal{F}$ est constructible.
\end{lemma}

\begin{proof}
Tout d'abord, on peut supposer le recouvrement fini puisque $X$ est quasi-compact
pour la topologie spectrale
(Topologie, lemme \ref{topology-lemma-constructible-hausdorff-quasi-compact}
et
Propriétés, lemme
\ref{properties-lemma-quasi-compact-quasi-separated-spectral}).
Remarquons que chaque $T_i$ est lui-même un schéma quasi-compact et quasi-séparé
(puisqu'il est constructible dans $X$; détails omis).
Il existe donc une partition finie $T_i = \coprod T_{i, j}$ en
parties localement fermées constructibles de $T_i$ telle que
$\mathcal{F}|_{T_{i, j}}$ soit localement constant, de valeur finie
(lemme \ref{lemma-constructible-quasi-compact-quasi-separated}).
D'après Topologie, lemme \ref{topology-lemma-constructible-in-constructible},
$T_{i, j}$ est un sous-schéma constructible localement fermé de $X$.
On peut alors appliquer Topologie, lemme
\ref{topology-lemma-constructible-partition-refined-by-stratification}
à $X = \bigcup T_{i, j}$ pour obtenir la partition cherchée de $X$.
\end{proof}

\begin{lemma}
\label{lemma-constructible-local}
Soit $X$ un schéma. La constructibilité d'un faisceau
d'ensembles, de groupes abéliens ou de $\Lambda$-modules (avec $\Lambda$ noethérien)
se vérifie localement pour la topologie de Zariski sur $X$.
\end{lemma}

\begin{proof}
L'énoncé signifie que, si $X = \bigcup U_i$ est un recouvrement ouvert
tel que $\mathcal{F}|_{U_i}$ soit constructible, alors $\mathcal{F}$
est constructible. Si $U \subset X$ est un ouvert affine, alors
$U = \bigcup U \cap U_i$ et $\mathcal{F}|_{U \cap U_i}$ est constructible
(il est évident que la restriction d'un faisceau constructible à
un ouvert est constructible). Il résulte du
lemme \ref{lemma-constructible-quasi-compact-quasi-separated}
que $\mathcal{F}|_U$ est constructible, c'est-à-dire qu'il existe une partition convenable
de $U$.
\end{proof}

\begin{lemma}
\label{lemma-pullback-constructible}
Soit $f : X \to Y$ un morphisme de schémas. Si $\mathcal{F}$ est un
faisceau constructible d'ensembles, de groupes abéliens ou de $\Lambda$-modules
(avec $\Lambda$ noethérien) sur $Y_\etale$, il en est de même
de $f^{-1}\mathcal{F}$ sur $X_\etale$.
\end{lemma}

\begin{proof}
D'après le lemme \ref{lemma-constructible-local}, on se ramène au cas
où $X$ et $Y$ sont affines. D'après le
lemme \ref{lemma-constructible-quasi-compact-quasi-separated},
il suffit de trouver une partition finie de $X$ en sous-schémas constructibles
localement fermés sur chacun desquels $f^{-1}\mathcal{F}$ soit localement
constant, de valeur finie.
Pour la trouver, il suffit de prendre l'image inverse de la partition de $Y$ adaptée à
$\mathcal{F}$ et d'utiliser le
lemme \ref{lemma-pullback-locally-constant}.
\end{proof}

\begin{lemma}
\label{lemma-constructible-abelian}
Soit $X$ un schéma.
\begin{enumerate}
\item La catégorie des faisceaux constructibles d'ensembles
est stable par limites et colimites finies dans $\Sh(X_\etale)$.
\item La catégorie des faisceaux abéliens constructibles est une
sous-catégorie de Serre faible de $\textit{Ab}(X_\etale)$.
\item Soit $\Lambda$ un anneau noethérien. La catégorie des
faisceaux constructibles de $\Lambda$-modules sur
$X_\etale$ est une sous-catégorie de Serre faible de
$\textit{Mod}(X_\etale, \Lambda)$.
\end{enumerate}
\end{lemma}

\begin{proof}
Nous démontrons (3). Nous utiliserons le critère donné dans
Homologie, lemme \ref{homology-lemma-characterize-weak-serre-subcategory}.
Supposons que $\varphi : \mathcal{F} \to \mathcal{G}$
soit un morphisme de faisceaux constructibles de $\Lambda$-modules.
Il faut montrer que $\mathcal{K} = \Ker(\varphi)$ et
$\mathcal{Q} = \Coker(\varphi)$ sont constructibles.
De même, supposons que
$0 \to \mathcal{F} \to \mathcal{E} \to \mathcal{G} \to 0$
soit une suite exacte courte de faisceaux de $\Lambda$-modules
avec $\mathcal{F}$ et $\mathcal{G}$ constructibles. Il faut montrer
que $\mathcal{E}$ est constructible.
Dans les deux cas, on peut remplacer $X$ par les membres d'un
recouvrement ouvert affine. On peut donc supposer $X$ affine.
On peut ensuite remplacer $X$ par les membres d'une
partition finie de $X$ en sous-schémas constructibles localement fermés
sur lesquels $\mathcal{F}$ et $\mathcal{G}$ sont de type fini et
localement constants. On conclut alors par le
lemme \ref{lemma-kernel-finite-locally-constant}.

\medskip\noindent
Les démonstrations de (1) et (2), très semblables, sont omises.
\end{proof}

\begin{lemma}
\label{lemma-support-constructible}
Soit $X$ un schéma quasi-compact et quasi-séparé.
\begin{enumerate}
\item Soit $\mathcal{F} \to \mathcal{G}$ un morphisme de faisceaux constructibles
d'ensembles sur $X_\etale$. Alors l'ensemble des points $x \in X$
où $\mathcal{F}_{\overline{x}} \to \mathcal{G}_{\overline{x}}$
est surjectif, resp.\ injectif, resp.\ isomorphe à une application donnée
d'ensembles, est constructible dans $X$.
\item Soit $\mathcal{F}$ un faisceau abélien constructible sur $X_\etale$.
Le support de $\mathcal{F}$ est constructible.
\item Soit $\Lambda$ un anneau noethérien.
Soit $\mathcal{F}$ un faisceau constructible de $\Lambda$-modules sur $X_\etale$.
Le support de $\mathcal{F}$ est constructible.
\end{enumerate}
\end{lemma}

\begin{proof}
Démontrons (1).
Soit $X = \coprod X_i$ une partition de $X$ en sous-schémas localement fermés
constructibles telle que $\mathcal{F}$ et $\mathcal{G}$ soient tous deux
localement constants, de valeur finie, sur les parties (appliquer le
lemme \ref{lemma-constructible-quasi-compact-quasi-separated}
à $\mathcal{F}$ et à $\mathcal{G}$, puis choisir un
raffinement commun). On applique alors le lemme \ref{lemma-morphism-locally-constant}
à la restriction du morphisme à chaque partie.

\medskip\noindent
Les démonstrations de (2) et (3) sont omises.
\end{proof}

\noindent
Le lemme suivant se révélera très utile par la suite.
Il affirme en substance que la catégorie des faisceaux constructibles
possède une forme faible de propriété ``noethérienne''.

\begin{lemma}
\label{lemma-colimit-constructible}
Soit $X$ un schéma quasi-compact et quasi-séparé. Soit
$\mathcal{F} = \colim_{i \in I} \mathcal{F}_i$ une limite inductive filtrante
de faisceaux d'ensembles, de faisceaux abéliens ou de faisceaux de modules.
\begin{enumerate}
\item Si $\mathcal{F}$ et les $\mathcal{F}_i$ sont des faisceaux constructibles
d'ensembles, alors l'ind-objet $\mathcal{F}_i$ est essentiellement constant de
valeur $\mathcal{F}$.
\item Si $\mathcal{F}$ et les $\mathcal{F}_i$ sont des faisceaux constructibles de
groupes abéliens, alors l'ind-objet $\mathcal{F}_i$ est essentiellement constant
de valeur $\mathcal{F}$.
\item Soit $\Lambda$ un anneau noethérien.
Si $\mathcal{F}$ et les $\mathcal{F}_i$ sont des faisceaux constructibles de
$\Lambda$-modules, alors l'ind-objet $\mathcal{F}_i$ est essentiellement constant
de valeur $\mathcal{F}$.
\end{enumerate}
\end{lemma}

\begin{proof}
Démontrons (1). Nous utiliserons sans plus le mentionner que les limites
et colimites finies de faisceaux constructibles sont constructibles
(lemme \ref{lemma-constructible-abelian}).
Pour tout $i$, soit $T_i \subset X$ l'ensemble des points $x \in X$
où $\mathcal{F}_{i, \overline{x}} \to \mathcal{F}_{\overline{x}}$
n'est pas surjectif. Puisque $\mathcal{F}_i$ et $\mathcal{F}$ sont
constructibles, $T_i$ est une partie constructible de $X$
(lemme \ref{lemma-support-constructible}).
Comme les fibres de $\mathcal{F}$ sont finies
et que $\mathcal{F} = \colim_{i \in I} \mathcal{F}_i$, on voit
que, pour tout $x \in X$, on a $x \not \in T_i$ pour $i$ assez grand.
Comme $X$ est un espace spectral d'après Propriétés, lemme
\ref{properties-lemma-quasi-compact-quasi-separated-spectral}
la topologie constructible sur $X$ est quasi-compact d'après
Topologie, lemme \ref{topology-lemma-constructible-hausdorff-quasi-compact}.
Ainsi, $T_i = \emptyset$ pour $i$ assez grand. Par conséquent,
$\mathcal{F}_i \to \mathcal{F}$ est surjectif pour $i$ assez grand.
Supposons maintenant que $\mathcal{F}_i \to \mathcal{F}$ soit surjectif pour tout $i$.
Fixons $i \in I$. Pour $i' \geq i$, notons $S_{i'} \subset X$ l'ensemble des
points $x$ pour lesquels le nombre d'éléments de
$\Im(\mathcal{F}_{i, \overline{x}} \to \mathcal{F}_{\overline{x}})$
n'est pas égal au nombre d'éléments de
$\Im(\mathcal{F}_{i, \overline{x}} \to \mathcal{F}_{i', \overline{x}})$.
Puisque $\mathcal{F}_i$, $\mathcal{F}_{i'}$ et $\mathcal{F}$ sont
constructibles, $S_{i'}$ est une partie constructible de $X$
(détails omis ; indication : utiliser le lemme \ref{lemma-support-constructible}).
Comme les fibres de $\mathcal{F}_i$ et de $\mathcal{F}$
sont finies et que $\mathcal{F} = \colim_{i' \geq i} \mathcal{F}_{i'}$,
on voit que, pour tout $x \in X$, on a $x \not \in S_{i'}$ pour $i'$
assez grand. Par le même argument que ci-dessus, il existe un $i'$ assez grand
tel que $S_{i'} = \emptyset$. Ainsi, $\mathcal{F}_i \to \mathcal{F}_{i'}$
se factorise par $\mathcal{F}$, comme souhaité.

\medskip\noindent
Démontrons (2). Remarquons qu'un faisceau abélien constructible est un faisceau
d'ensembles constructible. Le cas (2) résulte donc de (1).

\medskip\noindent
Démontrons (3). Nous utiliserons sans plus le mentionner que la catégorie des
faisceaux constructibles de $\Lambda$-modules est abélienne
(lemme \ref{lemma-constructible-abelian}).
Pour tout $i$, soit $\mathcal{Q}_i$ le conoyau du morphisme
$\mathcal{F}_i \to \mathcal{F}$. Le support $T_i$ de $\mathcal{Q}_i$
est une partie constructible de $X$ puisque $\mathcal{Q}_i$ est constructible
(lemme \ref{lemma-support-constructible}).
Comme les fibres de $\mathcal{F}$ sont des $\Lambda$-modules de type fini
et que $\mathcal{F} = \colim_{i \in I} \mathcal{F}_i$, on voit
que, pour tout $x \in X$, on a $x \not \in T_i$ pour $i$ assez grand.
Comme $X$ est un espace spectral d'après Propriétés, lemme
\ref{properties-lemma-quasi-compact-quasi-separated-spectral}
la topologie constructible sur $X$ est quasi-compact d'après
Topologie, lemme \ref{topology-lemma-constructible-hausdorff-quasi-compact}.
Ainsi, $T_i = \emptyset$ pour $i$ assez grand. Cela démontre la première
assertion. Pour la seconde, supposons maintenant que
$\mathcal{F}_i \to \mathcal{F}$ soit surjectif pour tout $i$.
Fixons $i \in I$. Pour $i' \geq i$, désignons par $\mathcal{K}_{i'}$
l'image de $\Ker(\mathcal{F}_i \to \mathcal{F})$ dans $\mathcal{F}_{i'}$.
Le support $S_{i'}$ de $\mathcal{K}_{i'}$
est une partie constructible de $X$ puisque $\mathcal{K}_{i'}$ est constructible.
Comme les fibres de $\Ker(\mathcal{F}_i \to \mathcal{F})$
sont des $\Lambda$-modules de type fini et que
$\mathcal{F} = \colim_{i' \geq i} \mathcal{F}_{i'}$, on voit
que, pour tout $x \in X$, on a $x \not \in S_{i'}$ pour $i'$ assez grand.
Par le même argument que ci-dessus, il existe un $i'$ assez grand
tel que $S_{i'} = \emptyset$. Ainsi, $\mathcal{F}_i \to \mathcal{F}_{i'}$
se factorise par $\mathcal{F}$, comme souhaité.
\end{proof}

\begin{lemma}
\label{lemma-tensor-product-constructible}
Soit $X$ un schéma. Soit $\Lambda$ un anneau noethérien.
Le produit tensoriel de deux faisceaux constructibles de $\Lambda$-modules
sur $X_\etale$ est un faisceau constructible de $\Lambda$-modules.
\end{lemma}

\begin{proof}
L'assertion se ramène immédiatement au cas où $X$ est affine.
Comme deux partitions quelconques de $X$ en strates constructibles localement
fermées admettent un raffinement commun du même type et
que les images inverses commutent au produit tensoriel, on se ramène au
lemme \ref{lemma-tensor-product-locally-constant}.
\end{proof}

\begin{lemma}
\label{lemma-tensor-constructible}
Soit $\Lambda \to \Lambda'$ un homomorphisme d'anneaux noethériens.
Soit $X$ un schéma. Soit $\mathcal{F}$ un faisceau constructible
de $\Lambda$-modules sur $X_\etale$. Alors
$\mathcal{F} \otimes_{\underline{\Lambda}} \underline{\Lambda'}$
est un faisceau constructible de $\Lambda'$-modules.
\end{lemma}

\begin{proof}
Omis. Indication : localement sur un ouvert affine, on peut utiliser la même stratification.
\end{proof}










\section{Lemmes auxiliaires sur les morphismes}
\label{section-stratify-morphisms}

\noindent
Voici quelques lemmes utiles pour d\'emontrer des propri\'et\'es de fonctorialit\'e
des faisceaux constructibles.

\begin{lemma}
\label{lemma-etale-stratified-finite}
Soit $U \to X$ un morphisme \'etale de sch\'emas quasi-compacts et quasi-s\'epar\'es
(par exemple un morphisme \'etale de sch\'emas noeth\'eriens). Alors il
existe une partition $X = \coprod_i X_i$ en sous-sch\'emas constructibles
localement ferm\'es telle que $X_i \times_X U \to X_i$ soit fini \'etale pour tout $i$.
\end{lemma}

\begin{proof}
Si $U \to X$ est s\'epar\'e, cela r\'esulte de
Compléments sur les morphismes, lemme \ref{more-morphisms-lemma-stratify-flat-fp-qf}.
En g\'en\'eral, on peut supposer $X$ affine. Choisissons un recouvrement ouvert
affine fini $U = \bigcup U_j$. Appliquons le cas pr\'ec\'edent \`a tous les morphismes
$U_j \to X$ et $U_j \cap U_{j'} \to X$, puis choisissons un raffinement
commun $X = \coprod X_i$ des partitions obtenues.
En raffinant encore la partition, on peut aussi supposer $X_i$ affine.
Fixons $i$ et posons $V = U \times_X X_i$. Les morphismes
$V_j = U_j \times_X X_i \to X_i$ et
$V_{jj'} = (U_j \cap U_{j'}) \times_X X_i \to X_i$ sont finis \'etales.
Ainsi, $V_j$ et $V_{jj'}$ sont des sch\'emas affines et $V_{jj'} \subset V_j$
est \`a la fois ferm\'e et ouvert (puisque $V_{jj'} \to X_i$ est propre, si bien que
Morphismes, lemme \ref{morphisms-lemma-image-proper-scheme-closed}
s'applique). Alors $V = \bigcup V_j$ est s\'epar\'e car
$\mathcal{O}(V_j) \to \mathcal{O}(V_{jj'})$ est surjectif, voir
Schémas, lemme \ref{schemes-lemma-characterize-separated}.
Le cas pr\'ec\'edent s'applique donc \`a $V \to X_i$, et l'on peut encore
raffiner la partition si n\'ecessaire (en fait, cela ne l'est pas, mais nous n'en
avons pas besoin).
\end{proof}

\noindent
Dans le cas noeth\'erien, on peut d\'emontrer le lemme pr\'ec\'edent par
r\'ecurrence noeth\'erienne \`a l'aide du petit lemme amusant suivant.

\begin{lemma}
\label{lemma-generically-finite}
Soit $f: X \to Y$ un morphisme de sch\'emas quasi-compact,
quasi-s\'epar\'e et localement de type fini. Si $\eta$ est un point g\'en\'erique
d'une composante irr\'eductible de $Y$ tel que $f^{-1}(\eta)$ soit fini, alors
il existe un ouvert $V \subset Y$ contenant $\eta$ tel que
$f^{-1}(V) \to V$ soit fini.
\end{lemma}

\begin{proof}
C'est Morphismes, lemme \ref{morphisms-lemma-generically-finite}.
\end{proof}

\noindent
L'\'enonc\'e du lemme suivant peut \^etre l\'eg\`erement renforc\'e.

\begin{lemma}
\label{lemma-decompose-quasi-finite-morphism}
Soit $f : Y \to X$ un morphisme quasi-fini et de pr\'esentation finie
de sch\'emas affines.
\begin{enumerate}
\item Il existe un morphisme surjectif de sch\'emas affines $X' \to X$ et un
sous-sch\'ema ferm\'e $Z' \subset Y' = X' \times_X Y$ tels que
\begin{enumerate}
\item l'immersion $Z' \subset Y'$ est un \'epaississement, et
\item $Z' \to X'$ est un morphisme fini \'etale.
\end{enumerate}
\item Il existe une partition finie $X = \coprod X_i$ en
strates affines, constructibles et localement ferm\'ees, et des morphismes surjectifs
finis localement libres $X'_i \to X_i$ tels que le morphisme r\'eduit associ\'e \`a
$Y'_i = X'_i \times_X Y \to X'_i$ soit isomorphe \`a
$\coprod_{j = 1}^{n_i} (X'_i)_{red} \to (X'_i)_{red}$ pour un certain $n_i$.
\end{enumerate}
\end{lemma}

\begin{proof}
En posant $X' = \coprod X'_i$, on voit que (2) implique (1).
\'Ecrivons $X = \Spec(A)$ et $Y = \Spec(B)$. Exprimons $A$ comme limite inductive filtrante
de $\mathbf{Z}$-alg\`ebres de type fini $A_i$. Comme $B$ est une $A$-alg\`ebre de
pr\'esentation finie, on voit qu'il existe $0 \in I$ et un
homomorphisme d'anneaux de type fini $A_0 \to B_0$ tels que $B = \colim B_i$, o\`u
$B_i = A_i \otimes_{A_0} B_0$, voir
Algèbre, lemme \ref{algebra-lemma-colimit-category-fp-algebras}.
Pour $i$ suffisamment grand, on voit que $A_i \to B_i$ est
quasi-fini, voir Limits, lemme \ref{limits-lemma-descend-quasi-finite}.
Nous sommes donc ramen\'es au cas des alg\`ebres de type fini sur $\mathbf{Z}$,
et en particulier au cas noeth\'erien. (D\'etails omis.)

\medskip\noindent
Supposons $X$ et $Y$ noeth\'eriens. Dans ce cas, tout sous-ensemble localement ferm\'e
de $X$ est constructible. Le lemme \ref{lemma-generically-finite}
et la r\'ecurrence noeth\'erienne montrent qu'il existe
une partition finie $X = \coprod X_i$ de $X$
en strates localement ferm\'ees telle que $Y \times_X X_i \to X_i$ soit fini.
On peut raffiner cette partition de sorte que les strates soient affines.
Ainsi, apr\`es avoir remplac\'e $X$ par $X' = \coprod X_i$, on peut supposer que
$Y \to X$ est fini.

\medskip\noindent
Supposons $X$ et $Y$ noeth\'eriens et $Y \to X$ fini.
Supposons pouvoir d\'emontrer (2) apr\`es changement de base par un morphisme
surjectif, plat et quasi-fini $U \to X$. On dispose donc d'une partition
$U = \coprod U_i$ et de morphismes finis localement libres $U'_i \to U_i$
tels que $U'_i \times_X Y \to U'_i$ soit isomorphe \`a
$\coprod_{j = 1}^{n_i} (U'_i)_{red} \to (U'_i)_{red}$ pour un certain $n_i$.
Alors, par l'argument du paragraphe pr\'ec\'edent, on peut trouver une
partition $X = \coprod X_j$ en strates affines localement ferm\'ees telle que
$X_j \times_X U_i \to X_j$ soit fini pour tous $i, j$. D'apr\`es
Morphismes, lemme \ref{morphisms-lemma-finite-flat},
chaque $X_j \times_X U_i \to X_j$ est fini localement libre.
Par cons\'equent, $X_j \times_X U'_i \to X_j$ est fini localement libre
(Morphismes, lemme \ref{morphisms-lemma-composition-finite-locally-free}).
Il s'ensuit que $X = \coprod X_j$ et $X_j' = \coprod_i X_j \times_X U'_i$
fournissent une solution pour $Y \to X$. Il suffit donc de d\'emontrer
le r\'esultat (dans le cas noeth\'erien) apr\`es un changement de base surjectif,
plat et quasi-fini.

\medskip\noindent
En appliquant Morphismes, lemme \ref{morphisms-lemma-massage-finite},
on voit qu'on peut supposer que $Y$ est un sous-sch\'ema ferm\'e d'un
sch\'ema affine $Z$ qui est, ensemblistement, une r\'eunion finie
$Z = \bigcup_{i \in I} Z_i$ de sous-sch\'emas ferm\'es s'envoyant isomorphiquement
sur $X$. Dans ce cas, on trouvera une partition finie $X = \coprod X_j$
en strates affines localement ferm\'ees qui convient (autrement dit $X'_j = X_j$).
Posons $T_i = Y \cap Z_i$. C'est un sous-sch\'ema ferm\'e de $X$.
Comme $X$ est noeth\'erien, on peut trouver une partition finie $X = \coprod X_j$
en sous-sch\'emas affines localement ferm\'es, telle que chaque
$X_j \times_X T_i$ soit vide ou soit, ensemblistement, \`egal \`a $X_j \times_X Z_i$.
En rempla\c cant $X$ par $X_j$, on voit qu'on peut supposer $I = I_1 \amalg I_2$,
avec $Z_i \subset Y$ pour $i \in I_1$ et $Z_i \cap Y = \emptyset$ pour
$i \in I_2$. En rempla\c cant $Z$ par $\bigcup_{i \in I_1} Z_i$, on voit qu'on
peut supposer $Y = Z$.
Enfin, on peut remplacer de nouveau $X$ par les membres d'une partition
comme ci-dessus telle que, pour tous $i, i' \in I$, l'intersection
$Z_i \cap Z_{i'}$ soit vide ou soit, ensemblistement, \`egale
\`a $Z_i$ et \`a $Z_{i'}$. Cela signifie clairement que $Y$ est, ensemblistement,
\'egal \`a une r\'eunion disjointe des $Z_i$, ce qu'il fallait d\'emontrer.
\end{proof}








\section{Compl\'ements sur les faisceaux constructibles}
\label{section-more-constructible}

\noindent
Soit $\Lambda$ un anneau noeth\'erien. Soit $X$ un sch\'ema.
Nous consid\'erons souvent $X_\etale$ comme un site annel\'e dont le
faisceau d'anneaux est $\underline{\Lambda}$. Pour les faisceaux ab\'eliens,
nous prenons souvent $\Lambda = \mathbf{Z}/n\mathbf{Z}$ pour un
entier $n$ convenable.

\begin{lemma}
\label{lemma-jshriek-constructible}
Soit $j : U \to X$ un morphisme \'etale de sch\'emas quasi-compacts et
quasi-s\'epar\'es.
\begin{enumerate}
\item Le faisceau $h_U$ est un faisceau d'ensembles constructible.
\item Le faisceau $j_!\underline{M}$ est un faisceau ab\'elien constructible
pour tout groupe ab\'elien fini $M$.
\item Si $\Lambda$ est un anneau noeth\'erien et $M$ un $\Lambda$-module de type fini,
alors $j_!\underline{M}$ est un faisceau constructible de $\Lambda$-modules
sur $X_\etale$.
\end{enumerate}
\end{lemma}

\begin{proof}
D'apr\`es le lemme \ref{lemma-etale-stratified-finite}, il existe une partition
$\coprod_i X_i$ telle que $\pi_i : j^{-1}(X_i) \to X_i$ soit fini \'etale.
La restriction de $h_U$ \`a $X_i$ est $h_{j^{-1}(X_i)}$, qui est fini et
localement constant d'apr\`es le lemme \ref{lemma-characterize-finite-locally-constant}.
Pour les cas (2) et (3), observons que
$$
j_!(\underline{M})|_{X_i} =
\pi_{i!}(\underline{M}) =
\pi_{i*}(\underline{M})
$$
par les lemmes \ref{lemma-shriek-base-change} et
\ref{lemma-shriek-equals-star-finite-etale}.
Il suffit donc de d\'emontrer le lemme lorsque $\pi : Y \to X$ est fini \'etale.
C'est le lemme \ref{lemma-pushforward-locally-constant}.
\end{proof}

\begin{lemma}
\label{lemma-torsion-colimit-constructible}
Soit $X$ un sch\'ema quasi-compact et quasi-s\'epar\'e.
\begin{enumerate}
\item Soit $\mathcal{F}$ un faisceau d'ensembles sur $X_\etale$.
Alors $\mathcal{F}$ est une limite inductive filtrante de
faisceaux d'ensembles constructibles.
\item Soit $\mathcal{F}$ un faisceau ab\'elien de torsion sur $X_\etale$.
Alors $\mathcal{F}$ est une limite inductive filtrante de faisceaux ab\'eliens constructibles.
\item Soient $\Lambda$ un anneau noeth\'erien et $\mathcal{F}$ un faisceau
de $\Lambda$-modules sur $X_\etale$. Alors
$\mathcal{F}$ est une limite inductive filtrante de faisceaux constructibles de
$\Lambda$-modules.
\end{enumerate}
\end{lemma}

\begin{proof}
Soit $\mathcal{B}$ la collection des objets quasi-compacts et quasi-s\'epar\'es
de $X_\etale$. D'apr\`es Modules sur les sites,
lemme \ref{sites-modules-lemma-module-filtered-colimit-constructibles},
tout faisceau d'ensembles est une limite inductive filtrante de faisceaux de la forme
$$
\text{Conoyau de la double fl\`eche}\left(
\xymatrix{
\coprod\nolimits_{j = 1}^m h_{V_j}
\ar@<1ex>[r] \ar@<-1ex>[r] &
\coprod\nolimits_{i = 1}^n h_{U_i}
}
\right)
$$
o\`u $V_j$ et $U_i$ sont des objets quasi-compacts et quasi-s\'epar\'es
de $X_\etale$. D'apr\`es les
lemmes \ref{lemma-jshriek-constructible} et \ref{lemma-constructible-abelian},
ces conoyaux de doubles fl\`eches sont constructibles. Cela d\'emontre (1).

\medskip\noindent
Soit $\Lambda$ un anneau noeth\'erien.
D'apr\`es Modules sur les sites,
lemme \ref{sites-modules-lemma-module-filtered-colimit-constructibles},
le faisceau de $\Lambda$-modules $\mathcal{F}$ est une limite inductive filtrante
de modules de la forme
$$
\Coker\left(
\bigoplus\nolimits_{j = 1}^m j_{V_j!}\underline{\Lambda}_{V_j}
\longrightarrow
\bigoplus\nolimits_{i = 1}^n j_{U_i!}\underline{\Lambda}_{U_i}
\right)
$$
o\`u $V_j$ et $U_i$ sont des objets quasi-compacts et quasi-s\'epar\'es
de $X_\etale$. D'apr\`es les
lemmes \ref{lemma-jshriek-constructible} et \ref{lemma-constructible-abelian},
ces conoyaux sont constructibles. Cela d\'emontre (3).

\medskip\noindent
D\'emontrons (2). \'Ecrivons d'abord $\mathcal{F} = \bigcup \mathcal{F}[n]$, o\`u
$\mathcal{F}[n]$ est le sous-faisceau de $n$-torsion. On peut alors consid\'erer
$\mathcal{F}[n]$ comme un faisceau de $\mathbf{Z}/n\mathbf{Z}$-modules
et appliquer (3).
\end{proof}

\begin{lemma}
\label{lemma-check-constructible}
Soit $f : X \to Y$ un morphisme surjectif de sch\'emas quasi-compacts et
quasi-s\'epar\'es.
\begin{enumerate}
\item Soit $\mathcal{F}$ un faisceau d'ensembles sur $Y_\etale$. Alors $\mathcal{F}$
est constructible si et seulement si $f^{-1}\mathcal{F}$ est constructible.
\item Soit $\mathcal{F}$ un faisceau ab\'elien sur $Y_\etale$. Alors $\mathcal{F}$
est constructible si et seulement si $f^{-1}\mathcal{F}$ est constructible.
\item Soit $\Lambda$ un anneau noeth\'erien.
Soit $\mathcal{F}$ un faisceau de $\Lambda$-modules sur $Y_\etale$.
Alors $\mathcal{F}$ est constructible si et seulement si $f^{-1}\mathcal{F}$
est constructible.
\end{enumerate}
\end{lemma}

\begin{proof}
Une implication r\'esulte du lemme \ref{lemma-pullback-constructible}.
Pour la r\'eciproque, supposons $f^{-1}\mathcal{F}$ constructible. Dans le cas ab\'elien,
la surjectivit\'e montre en outre que le faisceau initial est de torsion. \'Ecrivons $\mathcal{F} = \colim \mathcal{F}_i$ comme
limite inductive filtrante de faisceaux constructibles (d'ensembles, de groupes ab\'eliens ou de modules)
\`a l'aide du lemme \ref{lemma-torsion-colimit-constructible}.
Comme $f^{-1}$ est un adjoint \`a gauche, il commute aux limites inductives
(Catégories, lemme \ref{categories-lemma-adjoint-exact}), et l'on obtient
$f^{-1}\mathcal{F} = \colim f^{-1}\mathcal{F}_i$.
D'apr\`es le lemme \ref{lemma-colimit-constructible}, on voit que
$f^{-1}\mathcal{F}_i \to f^{-1}\mathcal{F}$
est surjectif pour tout $i$ suffisamment grand.
Comme $f$ est surjectif, on conclut (en examinant les fibres \`a l'aide du
lemme \ref{lemma-stalk-pullback} et du
th\'eor\`eme \ref{theorem-exactness-stalks})
que $\mathcal{F}_i \to \mathcal{F}$ est surjectif pour tout $i$ suffisamment grand.
Ainsi, $\mathcal{F}$ est le quotient d'un faisceau constructible $\mathcal{G}$.
En appliquant encore une fois l'argument \`a
$\mathcal{G} \times_\mathcal{F} \mathcal{G}$ ou
au noyau de $\mathcal{G} \to \mathcal{F}$,
on conclut en utilisant l'exactitude de $f^{-1}$ et le fait que la cat\'egorie des
faisceaux constructibles (d'ensembles, de groupes ab\'eliens ou de modules) est
stable par limites et colimites finies, ou par noyaux et conoyaux dans
$\Sh(Y_\etale)$, $\Sh(X_\etale)$, $\textit{Ab}(Y_\etale)$,
$\textit{Ab}(X_\etale)$, $\textit{Mod}(Y_\etale, \Lambda)$ et
$\textit{Mod}(X_\etale, \Lambda)$, voir le
lemme \ref{lemma-constructible-abelian}.
\end{proof}

\begin{lemma}
\label{lemma-pushforward-constructible}
Soit $f : X \to Y$ un morphisme fini \'etale de sch\'emas. Soit $\Lambda$ un
anneau noeth\'erien. Si $\mathcal{F}$ est un faisceau d'ensembles constructible,
un faisceau de groupes ab\'eliens constructible ou un faisceau constructible de
$\Lambda$-modules sur $X_\etale$, il en va de m\^eme de
$f_*\mathcal{F}$ sur $Y_\etale$.
\end{lemma}

\begin{proof}
D'apr\`es le lemme \ref{lemma-constructible-local}, il suffit de le v\'erifier
localement pour la topologie de Zariski sur $Y$ et, d'apr\`es le lemme \ref{lemma-check-constructible},
on peut remplacer $Y$ par un recouvrement \'etale (la construction de $f_*$
commute \`a la localisation \'etale). Un morphisme fini \'etale est
\'etale-localement isomorphe \`a une union disjointe d'isomorphismes, voir
Morphismes étales, lemme \ref{etale-lemma-finite-etale-etale-local}.
Ainsi, dans le cas des faisceaux d'ensembles, le lemme affirme que si
$\mathcal{F}_i$, $i = 1, \ldots, n$ sont des faisceaux d'ensembles constructibles, alors
$\prod_{i = 1}^n \mathcal{F}_i$ l'est aussi.
C'est clair. Il en va de m\^eme pour les faisceaux de groupes ab\'eliens et de modules.
\end{proof}

\begin{lemma}
\label{lemma-category-constructible-sets}
Soit $X$ un sch\'ema quasi-compact et quasi-s\'epar\'e. La cat\'egorie des
faisceaux d'ensembles constructibles est la sous-cat\'egorie pleine de $\Sh(X_\etale)$
form\'ee des faisceaux $\mathcal{F}$ qui sont les conoyaux de doubles fl\`eches
$$
\xymatrix{
\mathcal{F}_1
\ar@<1ex>[r] \ar@<-1ex>[r]
&
\mathcal{F}_0 \ar[r]
&
\mathcal{F}}
$$
pour lesquelles $\mathcal{F}_i$, pour $i = 0, 1$, est un coproduit fini de faisceaux de
la forme $h_U$, o\`u $U$ est un objet quasi-compact et quasi-s\'epar\'e
de $X_\etale$.
\end{lemma}

\begin{proof}
Dans la d\'emonstration du lemme \ref{lemma-torsion-colimit-constructible},
nous avons vu que les faisceaux de cette forme sont constructibles.
R\'eciproquement, supposons que, pour tout faisceau d'ensembles constructible
$\mathcal{F}$, on puisse trouver une surjection $\mathcal{F}_0 \to \mathcal{F}$
avec $\mathcal{F}_0$ comme dans le lemme. On trouve alors la surjection voulue
$\mathcal{F}_1 \to \mathcal{F}_0 \times_\mathcal{F} \mathcal{F}_0$,
car le dernier faisceau est constructible d'apr\`es le lemme \ref{lemma-constructible-abelian}.

\medskip\noindent
D'apr\`es Topologie, lemme
\ref{topology-lemma-constructible-partition-refined-by-stratification}
on peut choisir une stratification finie
$X = \coprod_{i \in I} X_i$ telle que $\mathcal{F}$ soit finie et localement
constante sur chaque strate. D\'emontrons le r\'esultat par r\'ecurrence sur
le cardinal de $I$. Soit $i \in I$ un \'el\'ement minimal pour l'ordre
partiel de $I$. Alors $X_i \subset X$ est ferm\'e.
Par r\'ecurrence, il existe un nombre fini d'objets quasi-compacts et quasi-s\'epar\'es
$U_\alpha$ de $(X \setminus X_i)_\etale$ et un morphisme surjectif
$\coprod h_{U_\alpha} \to \mathcal{F}|_{X \setminus X_i}$.
Ils d\'eterminent un morphisme
$$
\coprod h_{U_\alpha} \to \mathcal{F}
$$
qui est surjectif apr\`es restriction \`a $X \setminus X_i$. D'apr\`es le
lemme \ref{lemma-characterize-finite-locally-constant},
on a $\mathcal{F}|_{X_i} = h_V$ pour un sch\'ema $V$ fini \'etale
sur $X_i$. Soit $\overline{v}$ un point g\'eom\'etrique de $V$ au-dessus
de $\overline{x} \in X_i$. On peut consid\'erer $\overline{v}$ comme un
\'el\'ement de la fibre $\mathcal{F}_{\overline{x}} = V_{\overline{x}}$.
On peut donc trouver un voisinage \'etale $(U, \overline{u})$
de $\overline{x}$ et une section $s \in \mathcal{F}(U)$ dont la fibre en
$\overline{x}$ donne $\overline{v}$. En consid\'erant $s$ comme un morphisme
$s : h_U \to \mathcal{F}$ et en le restreignant \`a $X_i$, on obtient un morphisme
$s|_{X_i} : U \times_X X_i \to V$ au-dessus de $X_i$ qui envoie $\overline{u}$
sur $\overline{v}$. Comme $V$ est quasi-compact (il est fini sur le sous-sch\'ema
ferm\'e $X_i$ du sch\'ema quasi-compact $X$), un nombre fini
$s^{(1)}, \ldots, s^{(m)}$ de ces sections de $\mathcal{F}$ sur
$U^{(1)}, \ldots, U^{(m)}$ d\'eterminent conjointement le morphisme
surjectif
$$
\coprod s^{(j)}|_{X_i} : \coprod U^{(j)} \times_X X_i \longrightarrow V
$$
On obtient alors la surjection
$$
\coprod h_{U_\alpha} \amalg \coprod h_{U^{(j)}} \to \mathcal{F}
$$
voulue.
\end{proof}

\begin{lemma}
\label{lemma-category-constructible-modules}
Soit $X$ un schéma quasi-compact et quasi-séparé. Soit $\Lambda$
un anneau noethérien. La catégorie des faisceaux constructibles de
$\Lambda$-modules est exactement la catégorie des modules de la forme
$$
\Coker\left(
\bigoplus\nolimits_{j = 1}^m j_{V_j!}\underline{\Lambda}_{V_j}
\longrightarrow
\bigoplus\nolimits_{i = 1}^n j_{U_i!}\underline{\Lambda}_{U_i}
\right)
$$
où $V_j$ et $U_i$ sont des objets quasi-compacts et quasi-séparés de
$X_\etale$. En fait, on peut même supposer $U_i$ et $V_j$ affines.
\end{lemma}

\begin{proof}
Dans la démonstration du lemme \ref{lemma-torsion-colimit-constructible},
nous avons vu que les modules de cette forme sont constructibles. Comme la
catégorie des modules constructibles est abélienne
(lemme \ref{lemma-constructible-abelian}),
il suffit de montrer que, pour tout module constructible $\mathcal{F}$,
il existe une surjection
$$
\bigoplus\nolimits_{i = 1}^n j_{U_i!}\underline{\Lambda}_{U_i}
\longrightarrow \mathcal{F}
$$
pour certains objets affines $U_i$ de $X_\etale$. D'après
Modules sur les sites, lemme
\ref{sites-modules-lemma-module-filtered-colimit-constructibles}
il existe une surjection
$$
\Psi :
\bigoplus\nolimits_{i \in I} j_{U_i!}\underline{\Lambda}_{U_i}
\longrightarrow
\mathcal{F}
$$
où les $U_i$ sont affines et où la somme directe est prise sur un
ensemble d'indices $I$ éventuellement infini. Pour toute partie finie $I' \subset I$, posons
$$
T_{I'} = \text{Supp}(\Coker(
\bigoplus\nolimits_{i \in I'} j_{U_i!}\underline{\Lambda}_{U_i}
\longrightarrow \mathcal{F}))
$$
Par la définition même des faisceaux constructibles, l'ensemble $T_{I'}$
est une partie constructible de $X$. Nous voulons montrer que $T_{I'} = \emptyset$
pour un certain $I'$. Comme toute fibre $\mathcal{F}_{\overline{x}}$ est
un $\Lambda$-module de type fini et que $\Psi$ est surjective, pour
tout $x \in X$ il existe un $I'$ tel que $x \not \in T_{I'}$.
Autrement dit, on a
$\emptyset = \bigcap_{I' \subset I\text{ fini}} T_{I'}$. Comme
$X$ est un espace spectral d'après Properties, lemme
\ref{properties-lemma-quasi-compact-quasi-separated-spectral}
la topologie constructible sur $X$ est quasi-compacte d'après
Topologie, lemme \ref{topology-lemma-constructible-hausdorff-quasi-compact}.
Ainsi, $T_{I'} = \emptyset$ pour une partie finie $I' \subset I$,
ce qu'il fallait démontrer.
\end{proof}

\begin{lemma}
\label{lemma-category-constructible-abelian}
Soit $X$ un schéma quasi-compact et quasi-séparé. La catégorie des
faisceaux abéliens constructibles est exactement la catégorie des faisceaux
abéliens de la forme
$$
\Coker\left(
\bigoplus\nolimits_{j = 1}^m
j_{V_j!}\underline{\mathbf{Z}/m_j\mathbf{Z}}_{V_j}
\longrightarrow
\bigoplus\nolimits_{i = 1}^n
j_{U_i!}\underline{\mathbf{Z}/n_i\mathbf{Z}}_{U_i}
\right)
$$
où $V_j$ et $U_i$ sont des objets quasi-compacts et quasi-séparés de
$X_\etale$ et où $m_j$, $n_i$ sont des entiers positifs.
En fait, on peut même supposer $U_i$ et $V_j$ affines.
\end{lemma}

\begin{proof}
Cela résulte du lemme \ref{lemma-category-constructible-modules}
appliqué avec $\Lambda = \mathbf{Z}/n\mathbf{Z}$
et du fait que, puisque $X$ est quasi-compact, tout faisceau
abélien constructible est annulé par un entier positif $n$ (détails omis).
\end{proof}

\begin{lemma}
\label{lemma-constructible-is-compact}
Soit $X$ un schéma quasi-compact et quasi-séparé. Soit $\Lambda$ un
anneau noethérien. Soit $\mathcal{F}$ un faisceau constructible d'ensembles, de groupes
abéliens ou de $\Lambda$-modules sur $X_\etale$. Soit
$\mathcal{G} = \colim \mathcal{G}_i$ une limite inductive filtrante de faisceaux
d'ensembles, de groupes abéliens ou de $\Lambda$-modules. Alors
$$
\Mor(\mathcal{F}, \mathcal{G}) = \colim \Mor(\mathcal{F}, \mathcal{G}_i)
$$
dans la catégorie des faisceaux d'ensembles, de groupes abéliens ou de $\Lambda$-modules sur
$X_\etale$.
\end{lemma}

\begin{proof}
Traitons le cas des faisceaux d'ensembles. D'après le lemme \ref{lemma-category-constructible-sets},
il suffit de démontrer le lemme pour $h_U$, où $U$ est un objet quasi-compact
et quasi-séparé de $X_\etale$. Rappelons que
$\Mor(h_U, \mathcal{G}) = \mathcal{G}(U)$. Le résultat
découle donc de Sites, lemme \ref{sites-lemma-directed-colimits-sections}.

\medskip\noindent
Dans le cas des faisceaux abéliens ou des faisceaux de modules, le résultat se démontre
de la même manière à l'aide des
lemmes \ref{lemma-category-constructible-abelian} et
\ref{lemma-category-constructible-modules}.
Pour les faisceaux abéliens, ajoutons que
$\Mor(j_{U!}\underline{\mathbf{Z}/n\mathbf{Z}}, \mathcal{G})$
est égal à l'ensemble des éléments de $n$-torsion de $\mathcal{G}(U)$.
\end{proof}

\begin{lemma}
\label{lemma-finite-pushforward-constructible}
Soit $f : X \to Y$ un morphisme fini et de présentation finie de schémas.
Soit $\Lambda$ un anneau noethérien. Si $\mathcal{F}$ est un faisceau constructible
d'ensembles, de groupes abéliens ou de $\Lambda$-modules sur $X_\etale$,
alors $f_*\mathcal{F}$ l'est aussi.
\end{lemma}

\begin{proof}
Il suffit de le démontrer lorsque $X$ et $Y$ sont affines, d'après le
lemme \ref{lemma-constructible-local}.
D'après les
lemmes \ref{lemma-finite-pushforward-commutes-with-base-change} et
\ref{lemma-check-constructible}, on peut effectuer un changement de base par n'importe quel
schéma affine surjectif sur $Y$. Le
lemme \ref{lemma-decompose-quasi-finite-morphism}
nous ramène alors au cas d'un morphisme fini \'etale
(car un épaississement induit une équivalence des topos \'etales
et même des petits sites \'etales, voir le
théorème \ref{theorem-topological-invariance}).
Le cas fini \'etale est le
lemme \ref{lemma-pushforward-constructible}.
\end{proof}

\begin{lemma}
\label{lemma-category-is-colimit}
Soit $X = \lim_{i \in I} X_i$ la limite projective d'un système
projectif filtrant de schémas à morphismes de transition affines.
Nous supposons que $X_i$ est quasi-compact et quasi-séparé
pour tout $i \in I$.
\begin{enumerate}
\item La catégorie des faisceaux d'ensembles constructibles sur $X_\etale$
est la limite inductive des catégories des faisceaux d'ensembles constructibles
sur $(X_i)_\etale$.
\item La catégorie des faisceaux abéliens constructibles sur $X_\etale$
est la limite inductive des catégories des faisceaux abéliens constructibles
sur $(X_i)_\etale$.
\item Soit $\Lambda$ un anneau noethérien. La catégorie des faisceaux constructibles
de $\Lambda$-modules sur $X_\etale$ est la limite inductive des
catégories des faisceaux constructibles de $\Lambda$-modules sur $(X_i)_\etale$.
\end{enumerate}
\end{lemma}

\begin{proof}
Démontrons (1). Notons $f_i : X \to X_i$ les morphismes de projection.
La démonstration comporte trois parties, correspondant à ``fidèle'',
``pleinement fidèle'' et ``essentiellement surjectif''.

\medskip\noindent
Fidélité. Choisissons $0 \in I$ et soient $\mathcal{F}_0$, $\mathcal{G}_0$ des
faisceaux constructibles sur $X_0$. Supposons que
$a, b : \mathcal{F}_0 \to \mathcal{G}_0$ soient des morphismes tels que
$f_0^{-1}a = f_0^{-1}b$. Soit $E \subset X_0$ l'ensemble
des points $x \in X_0$ tels que $a_{\overline{x}} = b_{\overline{x}}$.
D'après le lemme \ref{lemma-support-constructible}, la partie
$E \subset X_0$ est constructible. Par hypothèse, $X \to X_0$ est à valeurs dans $E$.
D'après Limits, lemme \ref{limits-lemma-limit-contained-in-constructible},
il existe $i \geq 0$ tel que $X_i \to X_0$ soit à valeurs dans $E$.
Par conséquent, $f_{i0}^{-1}a = f_{i0}^{-1}b$.

\medskip\noindent
Pleine fidélité. Choisissons $0 \in I$ et soient $\mathcal{F}_0$, $\mathcal{G}_0$ des
faisceaux constructibles sur $X_0$. Supposons donné un morphisme
$a : f_0^{-1}\mathcal{F}_0 \to f_0^{-1}\mathcal{G}_0$.
Nous affirmons qu'il existe un $i$ et un morphisme
$a_i : f_{i0}^{-1}\mathcal{F}_0 \to f_{i0}^{-1}\mathcal{G}_0$
dont l'image inverse est $a$ sur $X$.
D'après le lemme \ref{lemma-category-constructible-sets},
on peut remplacer $\mathcal{F}_0$
par un coproduit fini de faisceaux représentés par des objets quasi-compacts
et quasi-séparés de $(X_0)_\etale$.
Il suffit donc d'établir ceci : si $U_0 \to X_0$ est un tel objet
de $(X_0)_\etale$, alors
$$
(f_0^{-1}\mathcal{G}_0)(U) = \colim_{i \geq 0} (f_{i0}^{-1}\mathcal{G}_0)(U_i)
$$
où $U = X \times_{X_0} U_0$ et $U_i = X_i \times_{X_0} U_0$.
C'est un cas particulier du théorème \ref{theorem-colimit}.

\medskip\noindent
Surjectivité essentielle. Il faut montrer que tout faisceau constructible $\mathcal{F}$
sur $X$ est isomorphe à $f_i^{-1}\mathcal{F}_i$ pour un faisceau constructible
$\mathcal{F}_i$ sur $X_i$. En appliquant le
lemme \ref{lemma-category-constructible-sets}
et en utilisant les résultats des deux paragraphes précédents, on voit ceci :
il suffit de le démontrer pour $h_U$, pour un certain objet quasi-compact
et quasi-séparé $U$ de $X_\etale$.
Dans ce cas, il faut montrer que $U$ s'obtient par changement de base à partir d'un
schéma quasi-compact et quasi-séparé \'etale sur $X_i$ pour un certain $i$.
Cela résulte des
lemmes \ref{limits-lemma-descend-finite-presentation} et
\ref{limits-lemma-descend-etale}.

\medskip\noindent
Démontrons (3). L'argument est très semblable à celui du cas des
faisceaux d'ensembles, mais on utilise le
lemme \ref{lemma-category-constructible-modules}
au lieu du
lemme \ref{lemma-category-constructible-sets}. Les détails sont omis.
L'assertion (2) résulte de (3), car tout faisceau abélien constructible
sur un schéma quasi-compact est un faisceau constructible de
$\mathbf{Z}/n\mathbf{Z}$-modules pour un certain $n$.
\end{proof}

\begin{lemma}
\label{lemma-category-loc-cst-is-colimit}
Soit $X = \lim_{i \in I} X_i$ la limite projective d'un système
projectif filtrant de schémas à morphismes de transition affines.
Nous supposons que $X_i$ est quasi-compact et quasi-séparé
pour tout $i \in I$.
\begin{enumerate}
\item La catégorie des faisceaux localement constants finis sur $X_\etale$
est la limite inductive des catégories des faisceaux localement constants finis
sur $(X_i)_\etale$.
\item La catégorie des faisceaux abéliens localement constants finis sur $X_\etale$
est la limite inductive des catégories des faisceaux abéliens localement constants finis
sur $(X_i)_\etale$.
\item Soit $\Lambda$ un anneau noethérien. La catégorie des faisceaux
localement constants de type fini de $\Lambda$-modules sur $X_\etale$
est la limite inductive des catégories des faisceaux localement constants
de type fini de $\Lambda$-modules sur $(X_i)_\etale$.
\end{enumerate}
\end{lemma}

\begin{proof}
D'après le lemme \ref{lemma-category-is-colimit}, le foncteur est dans chaque cas
pleinement fidèle. D'après le même lemme, il suffit, pour achever
la démonstration dans le cas (1), d'établir ceci : étant donné un faisceau constructible
$\mathcal{F}_i$ sur $X_i$ dont l'image inverse $\mathcal{F}$ sur $X$ est
localement constant fini, il existe $i' \geq i$ tel que
l'image inverse $\mathcal{F}_{i'}$ de $\mathcal{F}_i$ sur $X_{i'}$
soit localement constant fini. Par hypothèse, il existe un recouvrement \'etale
$\mathcal{U} = \{U_j \to X\}_{j \in J}$
tel que $\mathcal{F}|_{U_j} \cong \underline{S_j}$ pour un ensemble fini $S_j$.
On peut supposer $U_j$ affine pour tout $j \in J$.
Comme $X$ est quasi-compact, on peut supposer $J$ fini.
D'après le lemme \ref{lemma-colimit-affine-sites}, il existe $i' \geq i$
et un recouvrement \'etale $\mathcal{U}_{i'} = \{U_{i', j} \to X_{i'}\}_{j \in J}$
dont le changement de base à $X$ est $\mathcal{U}$. Alors
$\mathcal{F}_{i'}|_{U_{i', j}}$ et $\underline{S_j}$ sont des
faisceaux constructibles sur $(U_{i', j})_\etale$ dont les images inverses sur
$U_j$ sont isomorphes. Ainsi, quitte à augmenter $i'$, on obtient
que $\mathcal{F}_{i'}|_{U_{i', j}}$ et $\underline{S_j}$
sont isomorphes. Par conséquent, $\mathcal{F}_{i'}$ est localement constant fini.
La démonstration des cas (2) et (3) est exactement la même.
\end{proof}

\begin{lemma}
\label{lemma-irreducible-subsheaf-constant-zero}
Soit $X$ un schéma irréductible de point générique $\eta$.
\begin{enumerate}
\item Soit $S' \subset S$ une inclusion d'ensembles. Si l'on a
$\underline{S'} \subset \mathcal{G} \subset \underline{S}$
dans $\Sh(X_\etale)$ et si $S' = \mathcal{G}_{\overline{\eta}}$, alors
$\mathcal{G} = \underline{S'}$.
\item Soit $A' \subset A$ une inclusion de groupes abéliens. Si l'on a
$\underline{A'} \subset \mathcal{G} \subset \underline{A}$
dans $\textit{Ab}(X_\etale)$ et si $A' = \mathcal{G}_{\overline{\eta}}$, alors
$\mathcal{G} = \underline{A'}$.
\item Soit $M' \subset M$ une inclusion de modules sur un anneau $\Lambda$.
Si l'on a $\underline{M'} \subset \mathcal{G} \subset \underline{M}$
dans $\textit{Mod}(X_\etale, \underline{\Lambda})$
et si $M' = \mathcal{G}_{\overline{\eta}}$, alors
$\mathcal{G} = \underline{M'}$.
\end{enumerate}
\end{lemma}

\begin{proof}
En effet, pour tout morphisme \'etale $U \to X$
tel que $U \not = \emptyset$, le point $\eta$ appartient à l'image.
\end{proof}

\begin{lemma}
\label{lemma-push-constant-sheaf-from-open}
Soit $X$ un schéma intègre normal de corps des fonctions $K$.
Soit $E$ un ensemble.
\begin{enumerate}
\item Soit $g : \Spec(K) \to X$ l'inclusion du point générique.
Alors $g_*\underline{E} = \underline{E}$.
\item Soit $j : U \to X$ l'inclusion d'un ouvert non vide. Alors
$j_*\underline{E} = \underline{E}$.
\end{enumerate}
\end{lemma}

\begin{proof}
Démontrons (1). Soit $x \in X$ un point. Soit
$\mathcal{O}^{sh}_{X, \overline{x}}$
un hensélisé strict de $\mathcal{O}_{X, x}$. 
D'après Compléments d'algèbre, lemme \ref{more-algebra-lemma-henselization-normal},
on voit que $\mathcal{O}^{sh}_{X, \overline{x}}$ est un anneau intègre normal.
Ainsi, $\Spec(K) \times_X \Spec(\mathcal{O}^{sh}_{X, \overline{x}})$
est irréductible. Il s'ensuit que
la fibre $(g_*\underline{E})_{\overline{x}}$ est égale à $E$,
voir le théorème \ref{theorem-higher-direct-images}.

\medskip\noindent
Démontrons (2). Comme $g$ se factorise par $j$, il existe un morphisme
$j_*\underline{E} \to g_*\underline{E}$. Ce morphisme est injectif, car
pour tout schéma $V$ \'etale sur $X$, l'ensemble $\Spec(K) \times_X V$
est dense dans $U \times_X V$. D'autre part, il existe un morphisme
$\underline{E} \to j_*\underline{E}$, et on conclut.
\end{proof}

\begin{lemma}
\label{lemma-zero-in-generic-point}
Soit $X$ un schéma quasi-compact et quasi-séparé. Soit
$\eta \in X$ un point générique d'une composante irréductible de $X$.
\begin{enumerate}
\item Soit $\mathcal{F}$ un faisceau abélien de torsion sur $X_\etale$
dont la fibre $\mathcal{F}_{\overline{\eta}}$ est nulle.
Alors $\mathcal{F} = \colim \mathcal{F}_i$ est une limite inductive filtrante de
faisceaux abéliens constructibles $\mathcal{F}_i$ tels que, pour tout $i$,
le support de $\mathcal{F}_i$ soit contenu
dans un sous-schéma fermé ne contenant pas $\eta$.
\item Soit $\Lambda$ un anneau noethérien et $\mathcal{F}$ un faisceau
de $\Lambda$-modules sur $X_\etale$ dont la fibre
$\mathcal{F}_{\overline{\eta}}$ est nulle. Alors
$\mathcal{F} = \colim \mathcal{F}_i$
est une limite inductive filtrante de faisceaux constructibles de
$\Lambda$-modules $\mathcal{F}_i$ tels que, pour tout $i$,
le support de $\mathcal{F}_i$ soit contenu dans un sous-schéma fermé
ne contenant pas $\eta$.
\end{enumerate}
\end{lemma}

\begin{proof}
Démontrons (1). On peut écrire $\mathcal{F} = \colim_{i \in I} \mathcal{F}_i$,
où $\mathcal{F}_i$ est un faisceau abélien constructible, d'après le
lemme \ref{lemma-torsion-colimit-constructible}.
Choisissons $i \in I$. Comme $\mathcal{F}|_\eta$ est nul par hypothèse, on
voit qu'il existe $i'(i) \geq i$ tel que
$\mathcal{F}_i|_\eta \to \mathcal{F}_{i'(i)}|_\eta$ soit nul, voir le
lemme \ref{lemma-colimit-constructible}.
Alors $\mathcal{G}_i = \Im(\mathcal{F}_i \to \mathcal{F}_{i'(i)})$
est un faisceau abélien constructible (lemme \ref{lemma-constructible-abelian})
dont la fibre en $\eta$ est nulle.
Par conséquent, le support $E_i$ de $\mathcal{G}_i$ est une partie constructible
de $X$ ne contenant pas $\eta$. Comme
$\eta$ est un point générique d'une composante irréductible de
$X$, on voit que $\eta \not \in Z_i = \overline{E_i}$ d'après
Topologie, lemme \ref{topology-lemma-generic-point-in-constructible}.
Définissons un nouvel ensemble filtrant $I'$ ayant pour ensemble sous-jacent $I$ et
pour relation d'ordre la règle suivante :
$i_1$ est supérieur ou égal à $i_2$ si et seulement si $i_1 = i_2\ \text{ou}\ i_1 \geq i'(i_2)$.
Alors les faisceaux $\mathcal{G}_i$ forment un système indexé par $I'$
de limite inductive $\mathcal{F}$, ce qui achève la démonstration.

\medskip\noindent
La démonstration du cas (2) est exactement la même et nous l'omettons.
\end{proof}








\section{Faisceaux constructibles sur les schémas noethériens}
\label{section-noetherian-constructible}

\noindent
Si $X$ est un schéma noethérien, toute partie localement fermée est une
partie localement fermée constructible
(Topologie, lemme \ref{topology-lemma-constructible-Noetherian-space}).
Par conséquent, un faisceau abélien $\mathcal{F}$ sur $X_\etale$ est constructible
si et seulement s'il existe une partition finie $X = \coprod X_i$
telle que $\mathcal{F}|_{X_i}$ soit localement constant fini.
(Par convention, les parties d'une partition d'un espace topologique sont
localement fermées, voir Topologie, section \ref{topology-section-stratifications}.)
Autrement dit, on peut omettre l'adjectif ``constructible'' dans la
définition \ref{definition-constructible}. En fait, la catégorie
des faisceaux constructibles sur les
schémas noethériens possède quelques propriétés supplémentaires que nous
répertorions dans cette section.

\begin{proposition}
\label{proposition-constructible-over-noetherian}
Soit $X$ un schéma noethérien. Soit $\Lambda$ un anneau noethérien.
\begin{enumerate}
\item Tout sous-faisceau ou faisceau quotient d'un faisceau constructible
d'ensembles est constructible.
\item La catégorie des faisceaux abéliens constructibles sur $X_\etale$ est une
sous-catégorie de Serre (forte) de $\textit{Ab}(X_\etale)$. En particulier,
tout sous-faisceau et tout faisceau quotient d'un faisceau abélien constructible
sur $X_\etale$ sont constructibles.
\item La catégorie des faisceaux constructibles de $\Lambda$-modules
sur $X_\etale$ est une sous-catégorie de Serre (forte) de
$\textit{Mod}(X_\etale, \Lambda)$. En particulier, tout sous-module
et tout module quotient d'un faisceau constructible de $\Lambda$-modules
sur $X_\etale$ sont constructibles.
\end{enumerate}
\end{proposition}

\begin{proof}
Démontrons (1). Soit $\mathcal{G} \subset \mathcal{F}$, où $\mathcal{F}$ est
un faisceau constructible d'ensembles sur $X_\etale$. Soit $\eta \in X$ un point
générique d'une composante irréductible de $X$. Par récurrence noethérienne,
il suffit de trouver un voisinage ouvert $U$ de $\eta$ tel que
$\mathcal{G}|_U$ soit localement constant. Pour cela, on peut remplacer $X$
par un voisinage \'etale de $\eta$.
On peut donc supposer que $\mathcal{F}$ est constant et que $X$ est irréductible.

\medskip\noindent
Écrivons $\mathcal{F} = \underline{S}$ pour un ensemble fini $S$.
Alors $S' = \mathcal{G}_{\overline{\eta}} \subset S$;
écrivons $S' = \{s_1, \ldots, s_t\}$.
Choisissons un voisinage \'etale $(U, \overline{u})$ de $\overline{\eta}$
et des sections $\sigma_1, \ldots, \sigma_t \in \mathcal{G}(U)$ qui ont pour images
les $s_i$ dans $\mathcal{G}_{\overline{\eta}} \subset S$.
Comme $\sigma_i$ a pour image l'élément
$s_i \in S' \subset S = \Gamma(X, \mathcal{F})$,
on voit que les deux images inverses de $\sigma_i$ sur $U \times_X U$
sont égales en tant que sections de $\mathcal{G}$. Par la condition de faisceau
pour $\mathcal{G}$, on en déduit que $\sigma_i$ provient d'une section
de $\mathcal{G}$ sur l'ouvert $\Im(U \to X)$ de $X$.
Quitte à rétrécir $X$, on peut supposer que
$\underline{S'} \subset \mathcal{G} \subset \underline{S}$.
Alors $\underline{S'} = \mathcal{G}$ d'après le
lemme \ref{lemma-irreducible-subsheaf-constant-zero}.

\medskip\noindent
Soit $\mathcal{F} \to \mathcal{Q}$ un morphisme surjectif, où $\mathcal{F}$ est
un faisceau constructible d'ensembles sur $X_\etale$. Posons
$\mathcal{G} = \mathcal{F} \times_\mathcal{Q} \mathcal{F}$.
La première partie de la démonstration montre que $\mathcal{G}$ est
constructible comme sous-faisceau de $\mathcal{F} \times \mathcal{F}$.
Il en résulte que $\mathcal{Q}$ est constructible, voir le
lemme \ref{lemma-constructible-abelian}.

\medskip\noindent
Démontrons (3). Nous savons déjà que les faisceaux constructibles de modules
forment une sous-catégorie de Serre faible, voir le lemme \ref{lemma-constructible-abelian}.
Il suffit donc d'établir l'assertion concernant les sous-modules.

\medskip\noindent
Soit $\mathcal{G} \subset \mathcal{F}$ un sous-module d'un
faisceau constructible de $\Lambda$-modules sur $X_\etale$. Soit $\eta \in X$
un point générique d'une composante irréductible de $X$. Par récurrence noethérienne,
il suffit de trouver un voisinage ouvert $U$ de $\eta$ tel que
$\mathcal{G}|_U$ soit localement constant. Pour cela, on peut remplacer $X$
par un voisinage \'etale de $\eta$. On peut donc supposer que $\mathcal{F}$
est constant et que $X$ est irréductible.

\medskip\noindent
Écrivons $\mathcal{F} = \underline{M}$ pour un $\Lambda$-module de type fini $M$.
Alors $M' = \mathcal{G}_{\overline{\eta}} \subset M$. Choisissons un nombre fini
d'éléments $s_1, \ldots, s_t$ qui engendrent $M'$ comme $\Lambda$-module.
(C'est possible puisque $\Lambda$ est noethérien et que $M$ est de type fini.)
Choisissons un voisinage \'etale $(U, \overline{u})$ de $\overline{\eta}$
et des sections $\sigma_1, \ldots, \sigma_t \in \mathcal{G}(U)$ qui ont pour images
les $s_i$ dans $\mathcal{G}_{\overline{\eta}} \subset M$.
Comme $\sigma_i$ a pour image l'élément
$s_i \in M' \subset M = \Gamma(X, \mathcal{F})$,
on voit que les deux images inverses de $\sigma_i$ sur $U \times_X U$
sont égales en tant que sections de $\mathcal{G}$. Par la condition de faisceau
pour $\mathcal{G}$, on en déduit que $\sigma_i$ provient d'une section
de $\mathcal{G}$ sur l'ouvert $\Im(U \to X)$ de $X$.
Quitte à rétrécir $X$, on peut supposer que
$\underline{M'} \subset \mathcal{G} \subset \underline{M}$.
Alors $\underline{M'} = \mathcal{G}$ d'après le
lemme \ref{lemma-irreducible-subsheaf-constant-zero}.

\medskip\noindent
Démontrons (2). Cela se déduit de (3) de la manière habituelle. Nous omettons
les détails.
\end{proof}

\noindent
Le lemme suivant affirme que tout objet de la catégorie abélienne des
faisceaux constructibles sur $X$ est ``noethérien'', c'est-à-dire satisfait
à la condition de chaîne ascendante pour les sous-objets.

\begin{lemma}
\label{lemma-constructible-over-noetherian-noetherian}
Soit $X$ un schéma noethérien. Soit $\Lambda$ un anneau noethérien.
Considérons des inclusions
$$
\mathcal{F}_1 \subset \mathcal{F}_2 \subset \mathcal{F}_3 \subset \ldots
\subset \mathcal{F}
$$
dans la catégorie des faisceaux d'ensembles, de groupes abéliens ou de $\Lambda$-modules.
Si $\mathcal{F}$ est constructible, alors il existe $n$ tel que
l'on ait $\mathcal{F}_n = \mathcal{F}_{n + 1} = \mathcal{F}_{n + 2} = \ldots$.
\end{lemma}

\begin{proof}
D'après la proposition \ref{proposition-constructible-over-noetherian},
les $\mathcal{F}_i$ et $\colim \mathcal{F}_i$ sont constructibles.
Le lemme résulte alors du
lemme \ref{lemma-colimit-constructible}.
\end{proof}

\begin{lemma}
\label{lemma-constructible-maps-into-constant}
Soit $X$ un schéma noethérien.
\begin{enumerate}
\item Soit $\mathcal{F}$ un faisceau constructible d'ensembles sur $X_\etale$.
Il existe un monomorphisme de faisceaux
$$
\mathcal{F} \longrightarrow
\prod\nolimits_{i = 1}^n f_{i, *}\underline{E_i}
$$
où $f_i : Y_i \to X$ est un morphisme fini et $E_i$ un ensemble fini.
\item Soit $\mathcal{F}$ un faisceau abélien constructible sur $X_\etale$.
Il existe un monomorphisme de faisceaux abéliens
$$
\mathcal{F} \longrightarrow
\bigoplus\nolimits_{i = 1}^n f_{i, *}\underline{M_i}
$$
où $f_i : Y_i \to X$ est un morphisme fini et où
$M_i$ est un groupe abélien fini.
\item Soit $\Lambda$ un anneau noethérien.
Soit $\mathcal{F}$ un faisceau constructible de $\Lambda$-modules sur $X_\etale$.
Il existe un monomorphisme de faisceaux de modules
$$
\mathcal{F} \longrightarrow
\bigoplus\nolimits_{i = 1}^n f_{i, *}\underline{M_i}
$$
où $f_i : Y_i \to X$ est un morphisme fini et où
$M_i$ est un $\Lambda$-module de type fini.
\end{enumerate}
De plus, on peut supposer que chaque $Y_i$ est irréductible et réduit, et s'envoie
surjectivement sur un sous-schéma fermé irréductible et réduit $Z_i \subset X$ tel que
$Y_i \to Z_i$ soit fini \'etale au-dessus d'un ouvert non vide de $Z_i$.
\end{lemma}

\begin{proof}
Démontrons (1). Comme on dispose de la condition de chaîne ascendante pour les
sous-faisceaux de $\mathcal{F}$
(lemme \ref{lemma-constructible-over-noetherian-noetherian}), il
suffit de montrer que, pour tout point $x \in X$, on
peut trouver un morphisme $\varphi : \mathcal{F} \to f_*\underline{E}$, où
$f : Y \to X$ est fini et $E$ est un ensemble fini, tel que
$\varphi_{\overline{x}} : \mathcal{F}_{\overline{x}} \to
(f_*\underline{E})_{\overline{x}}$ soit injectif.
(On peut éviter cet argument en choisissant une partition de $X$ comme dans le
lemme \ref{lemma-constructible-quasi-compact-quasi-separated}
et en construisant un $Y_i \to X$ pour chaque composante irréductible
de chaque partie.)
Soit $Z \subset X$ le sous-schéma réduit induit
(Schémas, définition \ref{schemes-definition-reduced-induced-scheme})
sur $\overline{\{x\}}$.
Comme $\mathcal{F}$ est constructible, il existe une extension finie séparable
$K/\kappa(x)$ telle que
$\mathcal{F}|_{\Spec(K)}$ soit le faisceau constant de valeur $E$
pour un ensemble fini $E$. Soit $h : Y \to Z$ la normalisation
de $Z$ dans $\Spec(K)$, et notons $f : Y \to X$ le composé de $h$ avec $Z \to X$.
D'après Morphismes, lemme \ref{morphisms-lemma-normal-normalization},
on voit que $Y$ est un schéma intègre normal.
Comme $K/\kappa(x)$ est une extension finie, il est clair que $K$ est le corps
des fonctions de $Y$. Notons $g : \Spec(K) \to Y$ l'inclusion.
Le morphisme $\mathcal{F}|_{\Spec(K)} \to \underline{E}$ correspond par adjonction
à un morphisme $\mathcal{F}|_Y \to g_*\underline{E} = \underline{E}$
(lemme \ref{lemma-push-constant-sheaf-from-open}).
Celui-ci correspond à son tour par adjonction à un morphisme
$\varphi : \mathcal{F} \to f_*\underline{E}$.
Observons que le morphisme induit par $\varphi$ sur la fibre en un point géométrique
$\overline{x}$ est injectif : on peut choisir un relèvement $\overline{y} \in Y$
de $\overline{x}$, et le diagramme commutatif
$$
\xymatrix{
\mathcal{F}_{\overline{x}} \ar@{=}[r] \ar[d] &
(\mathcal{F}|_Y)_{\overline{y}} \ar@{=}[d] \\
(f_*\underline{E})_{\overline{x}} \ar[r] &
\underline{E}_{\overline{y}}
}
$$
prouve l'injectivité. Cependant, la démonstration n'est pas encore achevée, car le
morphisme $h : Y \to Z$ est entier, mais n'est pas fini en
général\footnote{Si $X$ est un schéma de Nagata, par exemple de type fini
sur un corps, alors $Y \to Z$ est fini.}.

\medskip\noindent
Pour résoudre le problème signalé dans la dernière phrase du paragraphe précédent,
écrivons $Y = \lim_{i \in I} Y_i$, où $Y_i$ est irréductible, intègre et
fini sur $Z$; notons $f_i : Y_i \to X$ les morphismes composés. Appliquons Properties, lemme
\ref{properties-lemma-integral-algebra-directed-colimit-finite},
à $h_*\mathcal{O}_Y$, considéré comme un faisceau de $\mathcal{O}_Z$-algèbres,
puis appliquons le foncteur $\underline{\Spec}_Z$.
Alors $f_*\underline{E} = \colim f_{i, *}\underline{E}$
d'après le lemme \ref{lemma-relative-colimit}.
D'après le lemme \ref{lemma-constructible-is-compact}, le morphisme
$\mathcal{F} \to f_*\underline{E}$
se factorise par $f_{i, *}\underline{E}$ pour un certain $i$.
Comme $Y_i \to Z$ est un morphisme fini entre schémas intègres
et que l'extension de corps de fonctions
induite par ce morphisme est finie séparable, on voit que le
morphisme est fini \'etale au-dessus d'un ouvert non vide de $Z$ (utiliser
Algèbre, lemme \ref{algebra-lemma-smooth-at-generic-point}; détails omis).
Ceci achève la démonstration de (1).

\medskip\noindent
Les démonstrations de (2) et (3) sont identiques à celle de (1).
\end{proof}

\noindent
Dans le lemme suivant, nous employons un procédé standard pour ramener une
assertion très générale au cas noethérien.

\begin{lemma}
\label{lemma-constructible-maps-into-constant-general}
\begin{reference}
\cite[Expos\'e IX, Proposition 2.14]{SGA4}
\end{reference}
Soit $X$ un schéma quasi-compact et quasi-séparé.
\begin{enumerate}
\item Soit $\mathcal{F}$ un faisceau constructible d'ensembles sur $X_\etale$.
Il existe un monomorphisme de faisceaux
$$
\mathcal{F} \longrightarrow
\prod\nolimits_{i = 1}^n f_{i, *}\underline{E_i}
$$
où $f_i : Y_i \to X$ est un morphisme fini et de présentation finie et où
$E_i$ est un ensemble fini.
\item Soit $\mathcal{F}$ un faisceau abélien constructible sur $X_\etale$.
Il existe un monomorphisme de faisceaux abéliens
$$
\mathcal{F} \longrightarrow
\bigoplus\nolimits_{i = 1}^n f_{i, *}\underline{M_i}
$$
où $f_i : Y_i \to X$ est un morphisme fini et de présentation finie et où
$M_i$ est un groupe abélien fini.
\item Soit $\Lambda$ un anneau noethérien.
Soit $\mathcal{F}$ un faisceau constructible de $\Lambda$-modules sur $X_\etale$.
Il existe un monomorphisme de faisceaux de modules
$$
\mathcal{F} \longrightarrow
\bigoplus\nolimits_{i = 1}^n f_{i, *}\underline{M_i}
$$
où $f_i : Y_i \to X$ est un morphisme fini et de présentation finie et où
$M_i$ est un $\Lambda$-module de type fini.
\end{enumerate}
\end{lemma}

\begin{proof}
Nous allons ramener ce lemme au cas noethérien par approximation noethérienne
absolue. Plus précisément, d'après la
proposition \ref{limits-proposition-approximate} de Limits,
on peut écrire $X = \lim_{t \in T} X_t$, où chaque $X_t$ est de type fini sur
$\Spec(\mathbf{Z})$ et les morphismes de transition sont affines. D'après le
lemme \ref{lemma-category-is-colimit},
la catégorie des faisceaux constructibles (d'ensembles, de groupes abéliens ou de
$\Lambda$-modules) sur $X_\etale$ est la limite inductive des
catégories correspondantes sur $X_t$. Ainsi, notre faisceau constructible $\mathcal{F}$
est l'image inverse d'un faisceau constructible du même type $\mathcal{F}_t$
sur $X_t$ pour un certain $t$. Appliquons alors le cas noethérien
(lemme \ref{lemma-constructible-maps-into-constant})
pour trouver un monomorphisme
$$
\mathcal{F}_t \longrightarrow
\prod\nolimits_{i = 1}^n f_{i, *}\underline{E_i}
\quad\text{ou}\quad
\mathcal{F}_t \longrightarrow
\bigoplus\nolimits_{i = 1}^n f_{i, *}\underline{M_i}
$$
sur $X_t$, pour certains morphismes finis $f_i : Y_i \to X_t$.
Comme $X_t$ est noethérien, les morphismes $f_i$ sont de présentation finie.
Comme l'image inverse est exacte et que la formation de $f_{i, *}$ commute
au changement de base
(lemme \ref{lemma-finite-pushforward-commutes-with-base-change}), on conclut.
\end{proof}

\begin{lemma}
\label{lemma-support-in-subset}
Soit $X$ un schéma noethérien. Soit $E \subset X$ une partie stable
par spécialisation.
\begin{enumerate}
\item Soit $\mathcal{F}$ un faisceau abélien de torsion sur $X_\etale$
dont le support est contenu dans $E$. Alors $\mathcal{F} = \colim \mathcal{F}_i$
est une limite inductive filtrante de faisceaux abéliens constructibles $\mathcal{F}_i$
tels que, pour tout $i$, le support de $\mathcal{F}_i$ soit contenu dans
une partie fermée elle-même contenue dans $E$.
\item Soient $\Lambda$ un anneau noethérien et $\mathcal{F}$ un faisceau
de $\Lambda$-modules sur $X_\etale$ dont le support est contenu dans $E$.
Alors $\mathcal{F} = \colim \mathcal{F}_i$
est une limite inductive filtrante de faisceaux constructibles de
$\Lambda$-modules $\mathcal{F}_i$ tels que, pour tout $i$,
le support de $\mathcal{F}_i$ soit contenu dans une partie fermée
elle-même contenue dans $E$.
\end{enumerate}
\end{lemma}

\begin{proof}
Démontrons (1). On peut écrire $\mathcal{F} = \colim_{i \in I} \mathcal{F}_i$,
où $\mathcal{F}_i$ est un faisceau abélien constructible, d'après le
lemme \ref{lemma-torsion-colimit-constructible}.
D'après la proposition \ref{proposition-constructible-over-noetherian},
l'image $\mathcal{F}'_i \subset \mathcal{F}$
du morphisme $\mathcal{F}_i \to \mathcal{F}$ est constructible.
Ainsi, $\mathcal{F} = \colim \mathcal{F}'_i$, et le support
de $\mathcal{F}'_i$ est contenu dans $E$.
Comme le support de $\mathcal{F}'_i$ est constructible
(par notre définition des faisceaux constructibles), on voit
que son adhérence est aussi contenue dans $E$, voir par exemple
Topologie, lemme
\ref{topology-lemma-constructible-stable-specialization-closed}.

\medskip\noindent
La démonstration du cas (2) est exactement la même et nous l'omettons.
\end{proof}








\section{Spécialisations et faisceaux étales}
\label{section-specialization}

\noindent
Image topologique : soit $X$ un espace topologique et soit $x' \leadsto x$
une spécialisation de points de $X$. Alors tout voisinage ouvert de $x$
contient $x'$. Par conséquent, pour tout faisceau $\mathcal{F}$ sur $X$, il existe un
{\it homomorphisme de spécialisation}
$$
sp : \mathcal{F}_x \longrightarrow \mathcal{F}_{x'}
$$
entre les fibres, qui envoie la classe d'équivalence du couple $(U, s)$ dans
$\mathcal{F}_x$ sur la classe d'équivalence du couple $(U, s)$ dans
$\mathcal{F}_{x'}$; voir Faisceaux, section \ref{sheaves-section-stalks},
pour la description des fibres en termes de classes d'équivalence de couples.
Cette application est bien entendu fonctorielle en $\mathcal{F}$, autrement dit $sp$
est une transformation de foncteurs.

\medskip\noindent
Pour les faisceaux de la topologie étale, on peut imiter cette construction; voir
\cite[Expos\'e VIII, 7.7, page 397]{SGA4}. Supposons donnés
un schéma $S$, un point géométrique $\overline{s}$ de $S$ et un point
géométrique $\overline{t}$ de $\Spec(\mathcal{O}^{sh}_{S, \overline{s}})$.
Pour tout faisceau $\mathcal{F}$ sur $S_\etale$, nous allons construire
l'{\it homomorphisme de spécialisation}
$$
sp : \mathcal{F}_{\overline{s}} \longrightarrow \mathcal{F}_{\overline{t}}
$$
Nous avons ici commis un abus de langage : au lieu d'écrire
$\mathcal{F}_{\overline{t}}$, il faudrait écrire $\mathcal{F}_{p(\overline{t})}$,
où $p : \Spec(\mathcal{O}^{sh}_{S, \overline{s}}) \to S$ est le morphisme
canonique. Rappelons que
$$
\mathcal{F}_{\overline{s}} = \colim_{(U, \overline{u})} \mathcal{F}(U)
$$
où la limite inductive porte sur tous les voisinages étales $(U, \overline{u})$
de $(S, \overline{s})$; voir section \ref{section-stalks}. Comme
$\mathcal{O}^{sh}_{S, \overline{s}}$ est la fibre du
faisceau structural, tout voisinage étale
$(U, \overline{u})$ de $(S, \overline{s})$ fournit un homomorphisme canonique
$\mathcal{O}_{U, u} \to \mathcal{O}^{sh}_{S, \overline{s}}$.
On obtient donc une unique factorisation
$$
\Spec(\mathcal{O}^{sh}_{S, \overline{s}}) \to U \to S
$$
Si $\overline{v}$ désigne l'image de $\overline{t}$ dans $U$, alors
$(U, \overline{v})$ est un voisinage étale de $(S, \overline{t})$.
Cette construction définit un foncteur de la catégorie des voisinages étales
de $(S, \overline{s})$ vers la catégorie des voisinages étales de
$(S, \overline{t})$. On peut donc définir l'application
$sp : \mathcal{F}_{\overline{s}} \to \mathcal{F}_{\overline{t}}$
en envoyant la classe d'équivalence de
$(U, \overline{u}, \sigma)$, où $\sigma \in \mathcal{F}(U)$,
sur la classe d'équivalence de $(U, \overline{v}, \sigma)$.

\medskip\noindent
Soit $K \in D(S_\etale)$. Pour $\overline{s}$ et $\overline{t}$
comme ci-dessus, on dispose de l'{\it homomorphisme de spécialisation}
$$
sp : K_{\overline{s}} \longrightarrow K_{\overline{t}}
\quad\text{dans}\quad D(\textit{Ab})
$$
En effet, si $K$ est représenté par le complexe $\mathcal{F}^\bullet$
de faisceaux abéliens, on prend simplement l'application
$$
K_{\overline{s}} = \mathcal{F}^\bullet_{\overline{s}}
\longrightarrow
\mathcal{F}^\bullet_{\overline{t}} = K_{\overline{t}}
$$
qui est donnée terme à terme par les homomorphismes de spécialisation des faisceaux
construits ci-dessus. Elle ne dépend pas du choix du complexe
représentant $K$, par exactitude des foncteurs fibres (autrement dit,
prendre les fibres d'un complexe est bien défini sur la catégorie dérivée).

\medskip\noindent
La construction est manifestement fonctorielle en le faisceau $\mathcal{F}$ sur
$S_\etale$. Si l'on considère les foncteurs fibres comme des morphismes de topos
$\overline{s}, \overline{t} : \textit{Ensembles} \to \Sh(S_\etale)$, alors
on peut considérer $sp$ comme un $2$-morphisme
$$
\xymatrix{
\textit{Ensembles}
\rrtwocell^{\overline{t}}_{\overline{s}}{\ sp}
&
&
\Sh(S_\etale)
}
$$
de topos.

\begin{remark}[Autre description de sp]
\label{remark-alternative-sp}
Soient $S$, $\overline{s}$ et $\overline{t}$ comme ci-dessus.
Une autre manière de décrire l'homomorphisme de spécialisation consiste à utiliser
$$
\mathcal{F}_{\overline{s}} =
\Gamma(\Spec(\mathcal{O}^{sh}_{S, \overline{s}}), p^{-1}\mathcal{F})
\quad\text{et}\quad
\mathcal{F}_{\overline{t}} = \Gamma(\overline{t},
\overline{t}^{-1}p^{-1}\mathcal{F})
$$
La première égalité résulte du Théorème \ref{theorem-higher-direct-images}
appliqué à $\text{id}_S : S \to S$, et la seconde résulte
du Lemme \ref{lemma-stalk-pullback}. On peut alors considérer $sp$
comme l'application
$$
sp :
\mathcal{F}_{\overline{s}} =
\Gamma(\Spec(\mathcal{O}^{sh}_{S, \overline{s}}), p^{-1}\mathcal{F})
\xrightarrow{\text{image inverse par }\overline{t}}
\Gamma(\overline{t}, \overline{t}^{-1}p^{-1}\mathcal{F}) =
\mathcal{F}_{\overline{t}}
$$
\end{remark}

\begin{remark}[Encore une description de sp]
\label{remark-another-sp}
Soient $S$, $\overline{s}$ et $\overline{t}$ comme ci-dessus.
Une autre possibilité consiste à utiliser l'unique morphisme
$$
c :
\Spec(\mathcal{O}^{sh}_{S, \overline{t}})
\longrightarrow
\Spec(\mathcal{O}^{sh}_{S, \overline{s}})
$$
au-dessus de $S$ qui est compatible avec le morphisme donné
$\overline{t} \to \Spec(\mathcal{O}^{sh}_{S, \overline{s}})$
et le morphisme
$\overline{t} \to \Spec(\mathcal{O}^{sh}_{S, \overline{t}})$.
L'existence et l'unicité de la flèche affichée
résultent d'Algèbre, Lemme \ref{algebra-lemma-map-into-henselian-colimit},
appliqué à $\mathcal{O}_{S, s}$, $\mathcal{O}^{sh}_{S, \overline{t}}$ et
$\mathcal{O}^{sh}_{S, \overline{s}} \to \kappa(\overline{t})$.
On obtient
$$
sp :
\mathcal{F}_{\overline{s}} =
\Gamma(\Spec(\mathcal{O}^{sh}_{S, \overline{s}}), \mathcal{F})
\xrightarrow{\text{image inverse par }c}
\Gamma(\Spec(\mathcal{O}^{sh}_{S, \overline{t}}), \mathcal{F}) =
\mathcal{F}_{\overline{t}}
$$
(avec les conventions de notation évidentes).
En fait, ce procédé fonctionne également pour les objets $K$ de $D(S_\etale)$ :
l'homomorphisme de spécialisation de $K$ est égal à l'application
$$
sp :
K_{\overline{s}} =
R\Gamma(\Spec(\mathcal{O}^{sh}_{S, \overline{s}}), K)
\xrightarrow{\text{image inverse par }c}
R\Gamma(\Spec(\mathcal{O}^{sh}_{S, \overline{t}}), K) =
K_{\overline{t}}
$$
Les signes d'égalité sont valables, car la prise des sections globales sur
les schémas strictement henséliens
$\Spec(\mathcal{O}^{sh}_{S, \overline{s}})$ et
$\Spec(\mathcal{O}^{sh}_{S, \overline{t}})$
est exacte (et revient à prendre les fibres en $\overline{s}$ et $\overline{t}$);
aucune difficulté liée au fait que $K$ puisse être non borné
n'apparaît donc.
\end{remark}

\begin{remark}[Relèvement des spécialisations]
\label{remark-can-lift}
Soit $S$ un schéma et soit $t \leadsto s$ une spécialisation de
points de $S$. Choisissons des points géométriques $\overline{t}$ et $\overline{s}$
au-dessus de $t$ et de $s$. Comme $t$ correspond à un point de
$\Spec(\mathcal{O}_{S, s})$ d'après
Schémas, Lemme \ref{schemes-lemma-specialize-points},
et comme $\mathcal{O}_{S, s} \to \mathcal{O}^{sh}_{S, \overline{s}}$
est fidèlement plat, on peut trouver un point
$t' \in \Spec(\mathcal{O}^{sh}_{S, \overline{s}})$ qui s'envoie sur $t$.
Puisque $\Spec(\mathcal{O}^{sh}_{S, \overline{s}})$ est une limite de
schémas étales sur $S$, l'extension $\kappa(t')/\kappa(t)$
est algébrique séparable (mais généralement pas finie, bien entendu).
Comme $\kappa(\overline{t})$ est algébriquement clos, on peut
choisir un plongement $\kappa(t') \to \kappa(\overline{t})$
d'extensions de $\kappa(t)$. Ce choix fournit un diagramme
commutatif
$$
\xymatrix{
\overline{t} \ar[d] \ar[r] &
\Spec(\mathcal{O}^{sh}_{S, \overline{s}}) \ar[d] &
\overline{s} \ar[l] \ar[d] \\
t \ar[r] & S & s \ar[l]
}
$$
de points et de points géométriques. Ainsi, si $t \leadsto s$,
on peut toujours « relever » $\overline{t}$ en un point géométrique
du hensélisé strict de $S$ en $\overline{s}$
et obtenir les homomorphismes de spécialisation ci-dessus.
\end{remark}

\begin{lemma}
\label{lemma-specialization-map-pullback}
Soit $g : S' \to S$ un morphisme de schémas. Soit $\mathcal{F}$ un faisceau
sur $S_\etale$. Soit $\overline{s}'$ un point géométrique de $S'$, et soit
$\overline{t}'$ un point géométrique de
$\Spec(\mathcal{O}^{sh}_{S', \overline{s}'})$. Posons
$\overline{s} = g(\overline{s}')$ et $\overline{t} = h(\overline{t}')$,
où $h : \Spec(\mathcal{O}^{sh}_{S', \overline{s}'}) \to
\Spec(\mathcal{O}^{sh}_{S, \overline{s}})$ est le morphisme canonique.
Pour tout faisceau $\mathcal{F}$ sur $S_\etale$, l'homomorphisme de spécialisation
$$
sp :
(g^{-1}\mathcal{F})_{\overline{s}'}
\longrightarrow
(g^{-1}\mathcal{F})_{\overline{t}'}
$$
est égal à l'homomorphisme de spécialisation
$sp : \mathcal{F}_{\overline{s}} \to \mathcal{F}_{\overline{t}}$
au moyen des identifications
$(g^{-1}\mathcal{F})_{\overline{s}'} = \mathcal{F}_{\overline{s}}$ et
$(g^{-1}\mathcal{F})_{\overline{t}'} = \mathcal{F}_{\overline{t}}$
du Lemme \ref{lemma-stalk-pullback}.
\end{lemma}

\begin{proof}
Omis.
\end{proof}

\begin{lemma}
\label{lemma-characterize-locally-constant}
Soit $S$ un schéma tel que tout ouvert quasi-compact de $S$ possède
un nombre fini de composantes irréductibles (par exemple, si $S$ a un
espace topologique sous-jacent noethérien ou si $S$ est localement noethérien).
Soit $\mathcal{F}$ un faisceau d'ensembles sur $S_\etale$.
Les conditions suivantes sont équivalentes :
\begin{enumerate}
\item $\mathcal{F}$ est localement constant fini;
\item toutes les fibres de $\mathcal{F}$ sont des ensembles finis et tous les homomorphismes
de spécialisation $sp : \mathcal{F}_{\overline{s}} \to \mathcal{F}_{\overline{t}}$
sont bijectifs.
\end{enumerate}
\end{lemma}

\begin{proof}
Supposons (2). Soit $\overline{s}$ un point géométrique de $S$ au-dessus de
$s \in S$. Pour démontrer (1), il faut trouver un voisinage étale
$(U, \overline{u})$ de $(S, \overline{s})$ tel que $\mathcal{F}|_U$
soit constant. On peut et on va supposer $S$ affine.

\medskip\noindent
Comme $\mathcal{F}_{\overline{s}}$ est fini, on peut choisir
$(U, \overline{u})$, un entier $n \geq 0$ et des éléments deux à deux distincts
$\sigma_1, \ldots, \sigma_n \in \mathcal{F}(U)$
tels que $\{\sigma_1, \ldots, \sigma_n\} \subset \mathcal{F}(U)$ s'envoie
bijectivement sur $\mathcal{F}_{\overline{s}}$ par l'application
$\mathcal{F}(U) \to \mathcal{F}_{\overline{s}}$.
Considérons le morphisme
$$
\varphi : \underline{\{1, \ldots, n\}} \longrightarrow \mathcal{F}|_U
$$
sur $U_\etale$ défini par $\sigma_1, \ldots, \sigma_n$.
Ce morphisme induit une bijection sur les fibres en $\overline{u}$
par construction. Considérons la partie
$$
E = \{u' \in U \mid \varphi_{\overline{u}'}\text{ est bijectif}\} \subset U
$$
Ici $\overline{u}'$ est un point géométrique quelconque de $U$ au-dessus de $u'$
(la condition ne dépend pas du choix, par Remarque \ref{remark-map-stalks}).
L'image $u \in U$ de $\overline{u}$ appartient à $E$.
Notre hypothèse sur les homomorphismes de spécialisation de $\mathcal{F}$,
la Remarque \ref{remark-can-lift} et le
Lemme \ref{lemma-specialization-map-pullback}
montrent que $E$ est stable par spécialisation
et par générisation dans l'espace topologique $U$.

\medskip\noindent
Après avoir rétréci $U$, on peut supposer $U$ affine lui aussi. Par
Descente, Lemme \ref{descent-lemma-locally-finite-nr-irred-local-fppf},
$U$ possède un nombre fini de composantes irréductibles.
Après avoir retiré celles qui ne passent pas
par $u$, on peut supposer que toute composante irréductible de $U$
passe par $u$. Puisque $U$ est un espace topologique sobre,
il s'ensuit que $E = U$ et que $\varphi$ est un isomorphisme par le
Théorème \ref{theorem-exactness-stalks}. Ainsi, (1) est vérifiée.

\medskip\noindent
Nous omettons la démonstration du fait que (1) implique (2).
\end{proof}

\begin{lemma}
\label{lemma-characterize-locally-constant-module}
Soit $S$ un schéma tel que tout ouvert quasi-compact de $S$ possède
un nombre fini de composantes irréductibles (par exemple, si $S$ a un
espace topologique sous-jacent noethérien ou si $S$ est localement noethérien).
Soit $\Lambda$ un anneau noethérien.
Soit $\mathcal{F}$ un faisceau de $\Lambda$-modules sur $S_\etale$.
Les conditions suivantes sont équivalentes :
\begin{enumerate}
\item $\mathcal{F}$ est un faisceau localement constant de type fini
de $\Lambda$-modules;
\item toutes les fibres de $\mathcal{F}$ sont des $\Lambda$-modules de type fini et
tous les homomorphismes de spécialisation
$sp : \mathcal{F}_{\overline{s}} \to \mathcal{F}_{\overline{t}}$
sont bijectifs.
\end{enumerate}
\end{lemma}

\begin{proof}
La démonstration est la même que celle du
Lemme \ref{lemma-characterize-locally-constant}.
Supposons (2). Soit $\overline{s}$ un point géométrique de $S$ au-dessus de
$s \in S$. Pour démontrer (1), il faut trouver un voisinage étale
$(U, \overline{u})$ de $(S, \overline{s})$ tel que $\mathcal{F}|_U$
soit constant. On peut et on va supposer $S$ affine.

\medskip\noindent
Comme $M = \mathcal{F}_{\overline{s}}$ est un $\Lambda$-module de type fini
et que $\Lambda$ est noethérien, on peut choisir une présentation
$$
\Lambda^{\oplus m} \xrightarrow{A} \Lambda^{\oplus n} \to M \to 0
$$
pour une matrice $A = (a_{ji})$ à coefficients dans $\Lambda$.
On peut choisir $(U, \overline{u})$ et des éléments
$\sigma_1, \ldots, \sigma_n \in \mathcal{F}(U)$
tels que $\sum a_{ji}\sigma_i = 0$ dans $\mathcal{F}(U)$ et
que les images des $\sigma_i$ dans $\mathcal{F}_{\overline{s}} = M$
soient celles des éléments de la base canonique de
$\Lambda^n$ dans la présentation de $M$ donnée ci-dessus.
Considérons le morphisme
$$
\varphi : \underline{M} \longrightarrow \mathcal{F}|_U
$$
sur $U_\etale$ défini par $\sigma_1, \ldots, \sigma_n$.
Ce morphisme induit une bijection sur les fibres en $\overline{u}$
par construction. Considérons la partie
$$
E = \{u' \in U \mid \varphi_{\overline{u}'}\text{ est bijectif}\} \subset U
$$
Ici $\overline{u}'$ est un point géométrique quelconque de $U$ au-dessus de $u'$
(la condition ne dépend pas du choix, par Remarque \ref{remark-map-stalks}).
L'image $u \in U$ de $\overline{u}$ appartient à $E$.
Notre hypothèse sur les homomorphismes de spécialisation de $\mathcal{F}$,
la Remarque \ref{remark-can-lift} et le
Lemme \ref{lemma-specialization-map-pullback}
montrent que $E$ est stable par spécialisation
et par générisation dans l'espace topologique $U$.

\medskip\noindent
Après avoir rétréci $U$, on peut supposer $U$ affine lui aussi. Par
Descente, Lemme \ref{descent-lemma-locally-finite-nr-irred-local-fppf},
$U$ possède un nombre fini de composantes irréductibles.
Après avoir retiré celles qui ne passent pas
par $u$, on peut supposer que toute composante irréductible de $U$
passe par $u$. Puisque $U$ est un espace topologique sobre,
il s'ensuit que $E = U$ et que $\varphi$ est un isomorphisme par le
Théorème \ref{theorem-exactness-stalks}. Ainsi, (1) est vérifiée.

\medskip\noindent
Nous omettons la démonstration du fait que (1) implique (2).
\end{proof}

\begin{lemma}
\label{lemma-specialization-map-pushforward}
Soit $f : X \to S$ un morphisme de schémas quasi-compact et quasi-séparé.
Soit $K \in D^+(X_\etale)$. Soit $\overline{s}$ un point géométrique de $S$
et soit $\overline{t}$ un point géométrique de
$\Spec(\mathcal{O}^{sh}_{S, \overline{s}})$. On a un diagramme
commutatif
$$
\xymatrix{
(Rf_*K)_{\overline{s}} \ar[r]_{sp} \ar@{=}[d] &
(Rf_*K)_{\overline{t}} \ar@{=}[d] \\
R\Gamma(X \times_S \Spec(\mathcal{O}^{sh}_{S, \overline{s}}), K)
\ar[r] &
R\Gamma(X \times_S \Spec(\mathcal{O}^{sh}_{S, \overline{t}}), K)
}
$$
où la flèche horizontale inférieure est l'image inverse par le morphisme
$\text{id}_X \times c$, où
$c : \Spec(\mathcal{O}^{sh}_{S, \overline{t}}) \to
\Spec(\mathcal{O}^{sh}_{S, \overline{s}})$
est le morphisme introduit dans la Remarque \ref{remark-another-sp}.
Les flèches verticales sont données par le
Théorème \ref{theorem-higher-direct-images}.
\end{lemma}

\begin{proof}
Cela résulte immédiatement de la description
de $sp$ donnée dans la Remarque \ref{remark-another-sp}.
\end{proof}

\begin{remark}
\label{remark-specialization-map-and-fibres}
Soit $f : X \to S$ un morphisme de schémas. Soit $K \in D(X_\etale)$.
Soit $\overline{s}$ un point géométrique de $S$ et soit $\overline{t}$
un point géométrique de $\Spec(\mathcal{O}^{sh}_{S, \overline{s}})$.
Soit $c$ comme dans la Remarque \ref{remark-another-sp}.
On peut toujours former un diagramme commutatif
$$
\xymatrix{
(Rf_*K)_{\overline{s}} \ar[r] \ar[d]_{sp} &
R\Gamma(X \times_S \Spec(\mathcal{O}_{S, \overline{s}}^{sh}), K)
\ar[r] \ar[d]_{(\text{id}_X \times c)^{-1}} &
R\Gamma(X_{\overline{s}}, K) \\
(Rf_*K)_{\overline{t}} \ar[r] &
R\Gamma(X \times_S \Spec(\mathcal{O}_{S, \overline{t}}^{sh}), K) \ar[r] &
R\Gamma(X_{\overline{t}}, K)
}
$$
où les flèches horizontales sont celles de la Remarque \ref{remark-stalk-fibre}.
En général, il n'existe pas de flèche verticale à droite
entre les cohomologies de $K$ sur les fibres qui complète ce
diagramme, même dans le cas du Lemme \ref{lemma-specialization-map-pushforward}.
\end{remark}















\section{Complexes à cohomologie constructible}
\label{section-Dc}

\noindent
Soit $\Lambda$ un anneau. Notons $D(X_\etale, \Lambda)$ la catégorie dérivée
des faisceaux de $\Lambda$-modules sur $X_\etale$.
On note $D^b(X_\etale, \Lambda)$ (respectivement $D^+$, $D^-$)
la sous-catégorie pleine des complexes bornés (respectivement bornés inférieurement,
supérieurement) de $D(X_\etale, \Lambda)$.

\begin{definition}
\label{definition-c}
Soit $X$ un schéma. Soit $\Lambda$ un anneau noethérien.
On note {\it $D_c(X_\etale, \Lambda)$} la sous-catégorie pleine
de $D(X_\etale, \Lambda)$ formée des complexes dont les faisceaux de cohomologie
sont des faisceaux constructibles de $\Lambda$-modules.
\end{definition}

\noindent
Cette définition a un sens par le Lemme \ref{lemma-constructible-abelian} et
Catégories dérivées, section \ref{derived-section-triangulated-sub}.
Ainsi, $D_c(X_\etale, \Lambda)$ est une sous-catégorie triangulée strictement
pleine et saturée de $D(X_\etale, \Lambda)$.

\begin{lemma}
\label{lemma-restrict-and-shriek-from-etale-c}
Soit $\Lambda$ un anneau noethérien.
Si $j : U \to X$ est un morphisme étale de schémas, alors
\begin{enumerate}
\item $K|_U \in D_c(U_\etale, \Lambda)$ si $K \in D_c(X_\etale, \Lambda)$;
\item $j_!M \in D_c(X_\etale, \Lambda)$ si $M \in D_c(U_\etale, \Lambda)$ et si
le morphisme $j$ est quasi-compact et quasi-séparé.
\end{enumerate}
\end{lemma}

\begin{proof}
La première assertion est claire. La seconde résulte de l'exactitude
de $j_!$ et du Lemme \ref{lemma-jshriek-constructible}.
\end{proof}

\begin{lemma}
\label{lemma-pullback-c}
Soit $\Lambda$ un anneau noethérien.
Soit $f : X \to Y$ un morphisme de schémas. Si $K \in D_c(Y_\etale, \Lambda)$,
alors $Lf^*K \in D_c(X_\etale, \Lambda)$.
\end{lemma}

\begin{proof}
Cela résulte de l'exactitude de $f^{-1} = f^*$ et du
Lemme \ref{lemma-pullback-constructible}.
\end{proof}

\begin{lemma}
\label{lemma-one-constructible}
Soit $X$ un schéma quasi-compact et quasi-séparé.
Soit $\Lambda$ un anneau noethérien. Soient $K \in D(X_\etale, \Lambda)$
et $b \in \mathbf{Z}$ tels que $H^b(K)$ soit constructible.
Il existe alors un faisceau $\mathcal{F}$ qui est une somme directe finie
de faisceaux $j_{U!}\underline{\Lambda}$, où $U \in \Ob(X_\etale)$ est affine, et
un morphisme $\mathcal{F}[-b] \to K$ dans $D(X_\etale, \Lambda)$
qui induit une surjection $\mathcal{F} \to H^b(K)$.
\end{lemma}

\begin{proof}
Représentons $K$ par un complexe $\mathcal{K}^\bullet$ de faisceaux de
$\Lambda$-modules. Considérons la surjection
$$
\Ker(\mathcal{K}^b \to \mathcal{K}^{b + 1})
\longrightarrow
H^b(K)
$$
Par Modules sur les sites, Lemme
\ref{sites-modules-lemma-module-quotient-direct-sum}
on peut choisir une surjection
$\bigoplus_{i \in I} j_{U_i!} \underline{\Lambda} \to
\Ker(\mathcal{K}^b \to \mathcal{K}^{b + 1})$
avec $U_i$ affine. Pour toute partie finie $I' \subset I$, notons
$\mathcal{H}_{I'} \subset H^b(K)$ l'image de
$\bigoplus_{i \in I'} j_{U_i!} \underline{\Lambda}$. Par le
Lemme \ref{lemma-colimit-constructible},
on a $\mathcal{H}_{I'} = H^b(K)$ pour une partie finie $I' \subset I$.
Le lemme s'ensuit en prenant
$\mathcal{F} = \bigoplus_{i \in I'} j_{U_i!} \underline{\Lambda}$.
\end{proof}

\begin{lemma}
\label{lemma-bounded-above-c}
Soit $X$ un schéma quasi-compact et quasi-séparé.
Soit $\Lambda$ un anneau noethérien. Soit $K \in D^-(X_\etale, \Lambda)$. Alors
les conditions suivantes sont équivalentes :
\begin{enumerate}
\item $K$ appartient à $D_c(X_\etale, \Lambda)$;
\item $K$ peut être représenté par un complexe borné supérieurement
dont les termes sont des sommes directes finies de $j_{U!}\underline{\Lambda}$,
où $U \in \Ob(X_\etale)$ est affine;
\item $K$ peut être représenté par un complexe borné supérieurement
de faisceaux constructibles plats de $\Lambda$-modules.
\end{enumerate}
\end{lemma}

\begin{proof}
Il est clair que (2) implique (3) et que (3) implique (1).
Supposons que $K$ appartienne à $D_c^-(X_\etale, \Lambda)$.
Supposons que $H^i(K) = 0$ pour $i > b$. Par récurrence sur $a$,
nous allons construire un complexe $\mathcal{F}^a \to \ldots \to \mathcal{F}^b$
tel que chaque $\mathcal{F}^i$ soit une somme directe finie
de $j_{U!}\underline{\Lambda}$, avec $U \in \Ob(X_\etale)$ affine,
et un morphisme $\mathcal{F}^\bullet \to K$ qui induise un isomorphisme
$H^i(\mathcal{F}^\bullet) \to H^i(K)$ pour $i > a$ et une surjection
$H^a(\mathcal{F}^\bullet) \to H^a(K)$.
Pour $a = b$, cela résulte du Lemme \ref{lemma-one-constructible}.
Une telle donnée étant construite, choisissons un triangle distingué
$$
\mathcal{F}^\bullet \to K \to L \to \mathcal{F}^\bullet[1]
$$
On voit alors que $H^i(L) = 0$ pour $i \geq a$. Choisissons
$\mathcal{F}^{a - 1}[-a +1] \to L$ comme dans le
Lemme \ref{lemma-one-constructible}. Le composé
$\mathcal{F}^{a - 1}[-a +1] \to L \to \mathcal{F}^\bullet$
correspond à un morphisme $\mathcal{F}^{a - 1} \to \mathcal{F}^a$
dont le composé avec $\mathcal{F}^a \to \mathcal{F}^{a + 1}$
est nul. Par TR4, on obtient un morphisme
$$
(\mathcal{F}^{a - 1} \to \ldots \to \mathcal{F}^b) \to K
$$
dans $D(X_\etale, \Lambda)$. Cela achève l'étape de récurrence et la
démonstration du lemme.
\end{proof}

\begin{lemma}
\label{lemma-tensor-c}
Soit $X$ un schéma. Soit $\Lambda$ un anneau noethérien.
Soient $K, L \in D_c^-(X_\etale, \Lambda)$. Alors
$K \otimes_\Lambda^\mathbf{L} L$ appartient à $D_c^-(X_\etale, \Lambda)$.
\end{lemma}

\begin{proof}
Cela résulte des Lemmes \ref{lemma-bounded-above-c} et
\ref{lemma-tensor-product-constructible}.
\end{proof}




\section{Tor-dimension finie et cohomologie constructible}
\label{section-ctf}

\noindent
Soit $X$ un schéma et soit $\Lambda$ un anneau noethérien. Une sous-catégorie
fréquemment utilisée de la catégorie dérivée $D_c(X_\etale, \Lambda)$ définie
à la section \ref{section-Dc} est la sous-catégorie pleine formée des objets
de tor-dimension (localement) finie. En voici la définition formelle.

\begin{definition}
\label{definition-ctf}
Soit $X$ un schéma. Soit $\Lambda$ un anneau noethérien. On note
{\it $D_{ctf}(X_\etale, \Lambda)$} la sous-catégorie pleine de
$D_c(X_\etale, \Lambda)$
formée des objets de tor-dimension localement finie.
\end{definition}

\noindent
C'est une sous-catégorie triangulée strictement pleine et saturée de
$D_c(X_\etale, \Lambda)$ et de $D(X_\etale, \Lambda)$. D'après nos conventions, voir
Cohomologie sur les sites, Définition \ref{sites-cohomology-definition-tor-amplitude},
on a
$$
D_{ctf}(X_\etale, \Lambda) \subset D^b_c(X_\etale, \Lambda)
\subset D^b(X_\etale, \Lambda)
$$
si $X$ est quasi-compact. Le Lemme \ref{lemma-when-ctf} donne une bonne
description des objets de $D_{ctf}(X_\etale, \Lambda)$.

\begin{remark}
\label{remark-different}
Les objets de la catégorie dérivée $D_{ctf}(X_\etale, \Lambda)$ possèdent, en un
certain sens, de meilleures propriétés globales que les objets parfaits de $D(\mathcal{O}_X)$.
En effet, un complexe de $\mathcal{O}_X$-modules peut être
localement quasi-isomorphe à un complexe fini dont les termes sont
des $\mathcal{O}_X$-modules localement libres de type fini, sans être
globalement quasi-isomorphe à un complexe borné de
$\mathcal{O}_X$-modules localement libres. Le lemme suivant montre que cela ne se produit pas pour
$D_{ctf}$ sur un schéma noethérien.
\end{remark}

\begin{lemma}
\label{lemma-when-ctf}
Soit $\Lambda$ un anneau noethérien. Soit $X$ un schéma quasi-compact
et quasi-séparé. Soit $K \in D(X_\etale, \Lambda)$. Les conditions suivantes
sont équivalentes :
\begin{enumerate}
\item $K \in D_{ctf}(X_\etale, \Lambda)$;
\item $K$ peut être représenté par un complexe fini de faisceaux
constructibles plats de $\Lambda$-modules.
\end{enumerate}
En fait, si $K$ est d'amplitude de Tor contenue dans $[a, b]$, on peut représenter
$K$ par un complexe $\mathcal{F}^a \to \ldots \to \mathcal{F}^b$ où
$\mathcal{F}^p$ est un faisceau constructible plat de $\Lambda$-modules.
\end{lemma}

\begin{proof}
Il est clair qu'un complexe fini de faisceaux constructibles
plats de $\Lambda$-modules est de tor-dimension finie.
Il est également clair qu'il appartient à $D_c(X_\etale, \Lambda)$.
Ainsi, (2) implique (1).

\medskip\noindent
Supposons (1). Choisissons $a, b \in \mathbf{Z}$ tels que
$H^i(K \otimes_\Lambda^\mathbf{L} \mathcal{G}) = 0$ si
$i \not \in [a, b]$, pour tout faisceau de $\Lambda$-modules $\mathcal{G}$.
Nous démontrerons la dernière assertion par récurrence sur $b - a$. Si
$a = b$, alors $K = H^a(K)[-a]$ est un faisceau constructible plat
et le résultat est acquis. Supposons maintenant $b > a$. Représentons $K$
par un complexe $\mathcal{K}^\bullet$ de faisceaux de $\Lambda$-modules.
Considérons la surjection
$$
\Ker(\mathcal{K}^b \to \mathcal{K}^{b + 1})
\longrightarrow
H^b(K)
$$
Par le Lemme \ref{lemma-category-constructible-modules},
on peut trouver un nombre fini de schémas affines $U_i$ étales sur $X$ et une
surjection $\bigoplus j_{U_i!}\underline{\Lambda}_{U_i} \to H^b(K)$.
Après avoir remplacé les $U_i$ par des recouvrements étales standard $\{U_{ij} \to U_i\}$,
on peut supposer que cette surjection se relève en un morphisme
$\mathcal{F} = \bigoplus j_{U_i!}\underline{\Lambda}_{U_i} \to
\Ker(\mathcal{K}^b \to \mathcal{K}^{b + 1})$.
Ce morphisme détermine un triangle distingué
$$
\mathcal{F}[-b] \to K \to L \to \mathcal{F}[-b + 1]
$$
dans $D(X_\etale, \Lambda)$. Comme $D_{ctf}(X_\etale, \Lambda)$ est une
sous-catégorie triangulée, $L$ lui appartient également. En fait, $L$ est
d'amplitude de Tor contenue dans $[a, b - 1]$, puisque $\mathcal{F}$ se surjecte sur
$H^b(K)$ (nous omettons les détails). Par l'hypothèse de récurrence, il existe
un complexe fini $\mathcal{F}^a \to \ldots \to \mathcal{F}^{b - 1}$
de faisceaux constructibles plats de $\Lambda$-modules qui représente $L$.
Le morphisme $L \to \mathcal{F}[-b + 1]$ correspond à un morphisme
$\mathcal{F}^{b-1} \to \mathcal{F}$ qui annule l'image
de $\mathcal{F}^{b-2} \to \mathcal{F}^{b-1}$. Posons $\mathcal{F}^b=\mathcal{F}$; il résulte alors
de l'axiome TR3 que $K$ est représenté par le complexe
$$
\mathcal{F}^a \to \ldots \to \mathcal{F}^{b - 1} \to \mathcal{F}^b
$$
ce qui achève la démonstration.
\end{proof}

\begin{remark}
\label{remark-projective-each-degree}
Soit $\Lambda$ un anneau noethérien. Soit $X$ un schéma.
Pour un complexe borné $K^\bullet$ de faisceaux constructibles plats de $\Lambda$-modules
sur $X_\etale$,
chaque fibre $K^p_{\overline{x}}$ est un $\Lambda$-module projectif de type fini.
Les fibres du complexe sont donc des complexes parfaits de $\Lambda$-modules.
\end{remark}

\begin{lemma}
\label{lemma-restrict-and-shriek-from-etale-ctf}
Soit $\Lambda$ un anneau noethérien.
Si $j : U \to X$ est un morphisme étale de schémas, alors
\begin{enumerate}
\item $K|_U \in D_{ctf}(U_\etale, \Lambda)$ si
$K \in D_{ctf}(X_\etale, \Lambda)$;
\item $j_!M \in D_{ctf}(X_\etale, \Lambda)$ si
$M \in D_{ctf}(U_\etale, \Lambda)$ et si
le morphisme $j$ est quasi-compact et quasi-séparé.
\end{enumerate}
\end{lemma}

\begin{proof}
Le moyen le plus simple de démontrer ce lemme est peut-être de se ramener au
cas où $X$ est affine, puis d'appliquer le Lemme \ref{lemma-when-ctf}
pour le reformuler comme un énoncé sur les complexes finis
de faisceaux constructibles plats de $\Lambda$-modules;
le résultat découle alors du Lemme \ref{lemma-jshriek-constructible}.
\end{proof}

\begin{lemma}
\label{lemma-pullback-ctf}
Soit $\Lambda$ un anneau noethérien. Soit $f : X \to Y$ un morphisme de schémas.
Si $K \in D_{ctf}(Y_\etale, \Lambda)$, alors
$Lf^*K \in D_{ctf}(X_\etale, \Lambda)$.
\end{lemma}

\begin{proof}
Appliquons le Lemme \ref{lemma-when-ctf} pour nous ramener à une question
sur les complexes finis de faisceaux constructibles plats de $\Lambda$-modules.
L'énoncé résulte alors de l'exactitude de $f^{-1} = f^*$ et du
Lemme \ref{lemma-pullback-constructible}.
\end{proof}

\begin{lemma}
\label{lemma-connected-ctf-locally-constant}
Soit $X$ un schéma connexe. Soit $\Lambda$ un anneau noethérien.
Soit $K \in D_{ctf}(X_\etale, \Lambda)$ à faisceaux de cohomologie localement constants.
Il existe alors un complexe fini de $\Lambda$-modules projectifs de type fini
$M^\bullet$ et un recouvrement étale $\{U_i \to X\}$ tels que
$K|_{U_i} \cong \underline{M^\bullet}|_{U_i}$ dans $D(U_{i, \etale}, \Lambda)$.
\end{lemma}

\begin{proof}
Choisissons un recouvrement étale $\{U_i \to X\}$ tel que $K|_{U_i}$
soit constant, disons $K|_{U_i} \cong \underline{M_i^\bullet}_{U_i}$
pour un complexe fini de $\Lambda$-modules $M_i^\bullet$ de type fini.
Voir Cohomologie sur les sites, Lemme
\ref{sites-cohomology-lemma-locally-constant}.
Remarquons que $U_i \times_X U_j$ est vide si $M_i^\bullet$
n'est pas isomorphe à $M_j^\bullet$ dans $D(\Lambda)$.
Pour chaque complexe de $\Lambda$-modules $M^\bullet$, posons
$I_{M^\bullet} =
\{i \in I \mid M_i^\bullet \cong M^\bullet\text{ dans }D(\Lambda)\}$.
Comme les morphismes étales sont ouverts,
$U_{M^\bullet} = \bigcup_{i \in I_{M^\bullet}} \Im(U_i \to X)$
est une partie ouverte de $X$. Alors $X = \coprod U_{M^\bullet}$ est un
recouvrement ouvert disjoint de $X$. Comme $X$ est connexe, un seul $U_{M^\bullet}$
est non vide. Puisque $K$ appartient à $D_{ctf}(X_\etale, \Lambda)$, $M^\bullet$
est un complexe parfait de $\Lambda$-modules; voir
Compléments d'algèbre, Lemme \ref{more-algebra-lemma-perfect}.
On peut donc supposer que $M^\bullet$ est un complexe fini de
$\Lambda$-modules projectifs de type fini.
\end{proof}








\section{Faisceaux de torsion}
\label{section-torsion}

\noindent
Voici une brève section sur les faisceaux abéliens de torsion et leur cohomologie
étale. Soit $\mathcal{C}$ un site. Nous avons montré dans
Cohomologie sur les sites, Lemme \ref{sites-cohomology-lemma-torsion},
que tout objet de $D(\mathcal{C})$ dont les faisceaux de cohomologie sont
de torsion peut être représenté par un complexe dont tous les termes
sont de torsion.

\begin{lemma}
\label{lemma-torsion-cohomology}
Soit $X$ un schéma quasi-compact et quasi-séparé.
\begin{enumerate}
\item Si $\mathcal{F}$ est un faisceau abélien de torsion sur $X_\etale$, alors
$H^n_\etale(X, \mathcal{F})$ est un groupe abélien de torsion pour tout $n$.
\item Si $K$ dans $D^+(X_\etale)$ a des faisceaux de cohomologie de torsion, alors
$H^n_\etale(X, K)$ est un groupe abélien de torsion pour tout $n$.
\end{enumerate}
\end{lemma}

\begin{proof}
Pour démontrer (1), écrivons $\mathcal{F} = \bigcup_d \mathcal{F}[d]$, où
$\mathcal{F}[d]$ est le sous-faisceau de $d$-torsion. Par le
Lemme \ref{lemma-colimit}, on a
$H^n_\etale(X, \mathcal{F}) = \colim H^n_\etale(X, \mathcal{F}[d])$.
Cela démontre (1), car $H^n_\etale(X, \mathcal{F}[d])$ est annulé par $d$.

\medskip\noindent
Pour démontrer (2), on peut utiliser la suite spectrale
$E_2^{p, q} = H^p_\etale(X, H^q(K))$ qui aboutit à $H^n_\etale(X, K)$
(Catégories dérivées, Lemme \ref{derived-lemma-two-ss-complex-functor})
et le résultat pour les faisceaux.
\end{proof}

\begin{lemma}
\label{lemma-torsion-direct-image}
Soit $f : X \to Y$ un morphisme de schémas quasi-compact et
quasi-séparé.
\begin{enumerate}
\item Si $\mathcal{F}$ est un faisceau abélien de torsion sur $X_\etale$, alors
$R^nf_*\mathcal{F}$ est un faisceau abélien de torsion sur $Y_\etale$ pour tout $n$.
\item Si $K$ dans $D^+(X_\etale)$ a des faisceaux de cohomologie de torsion, alors
$Rf_*K$ est un objet de $D^+(Y_\etale)$ dont les faisceaux de cohomologie sont
des faisceaux abéliens de torsion.
\end{enumerate}
\end{lemma}

\begin{proof}
Démontrons (1). Rappelons que $R^nf_*\mathcal{F}$ est le faisceau associé
au préfaisceau $V \mapsto H^n_\etale(X \times_Y V, \mathcal{F})$
sur $Y_\etale$. Voir Cohomologie sur les sites,
Lemme \ref{sites-cohomology-lemma-higher-direct-images}.
Si l'on choisit $V$ affine, alors $X \times_Y V$ est
quasi-compact et quasi-séparé puisque $f$ l'est; on peut donc
appliquer le Lemme \ref{lemma-torsion-cohomology} pour voir que
$H^n_\etale(X \times_Y V, \mathcal{F})$ est de torsion.

\medskip\noindent
Démontrons (2). Rappelons que $R^nf_*K$ est le faisceau associé
au préfaisceau $V \mapsto H^n_\etale(X \times_Y V, K)$
sur $Y_\etale$. Voir Cohomologie sur les sites,
Lemme \ref{sites-cohomology-lemma-unbounded-describe-higher-direct-images}.
Si l'on choisit $V$ affine, alors $X \times_Y V$ est
quasi-compact et quasi-séparé puisque $f$ l'est; on peut donc
appliquer le Lemme \ref{lemma-torsion-cohomology} pour voir que
$H^n_\etale(X \times_Y V, K)$ est de torsion.
\end{proof}









\section{Cohomologie à supports dans un sous-schéma fermé}
\label{section-cohomology-support}

\noindent
Soit $X$ un schéma et soit $Z \subset X$ un sous-schéma fermé.
Soit $\mathcal{F}$ un faisceau abélien sur $X_\etale$. Posons
$$
\Gamma_Z(X, \mathcal{F}) =
\{s \in \mathcal{F}(X) \mid \text{Supp}(s) \subset Z\}
$$
pour le groupe des sections à support dans $Z$ (Définition \ref{definition-support}).
C'est un foncteur exact à gauche, qui n'est pas exact en général.
On obtient donc un foncteur dérivé
$$
R\Gamma_Z(X, -) : D(X_\etale) \longrightarrow D(\textit{Ab})
$$
et les groupes de cohomologie à support dans $Z$ définis par
$H^q_Z(X, \mathcal{F}) = R^q\Gamma_Z(X, \mathcal{F})$.

\medskip\noindent
Soit $\mathcal{I}$ un faisceau abélien injectif sur $X_\etale$.
Posons $U = X \setminus Z$.
Alors l'application de restriction $\mathcal{I}(X) \to \mathcal{I}(U)$ est surjective
(Cohomologie sur les sites, Lemme
\ref{sites-cohomology-lemma-restriction-along-monomorphism-surjective})
et son noyau est $\Gamma_Z(X, \mathcal{I})$. Il s'ensuit immédiatement que,
pour $K \in D(X_\etale)$, il existe un triangle distingué
$$
R\Gamma_Z(X, K) \to R\Gamma(X, K) \to R\Gamma(U, K) \to R\Gamma_Z(X, K)[1]
$$
dans $D(\textit{Ab})$. On en déduit une longue suite exacte de
cohomologie
$$
\ldots \to H^i_Z(X, K) \to H^i(X, K) \to H^i(U, K) \to
H^{i + 1}_Z(X, K) \to \ldots
$$
pour tout $K$ de $D(X_\etale)$.

\medskip\noindent
Pour un faisceau abélien $\mathcal{F}$ sur $X_\etale$, on peut considérer
le {\it sous-faisceau des sections à support dans $Z$}, noté
$\mathcal{H}_Z(\mathcal{F})$, défini par la règle
$$
\mathcal{H}_Z(\mathcal{F})(U) =
\{s \in \mathcal{F}(U) \mid \text{Supp}(s) \subset U \times_X Z\}
$$
Nous utilisons ici le support d'une section défini dans la
Définition \ref{definition-support}. Grâce à l'équivalence de la
Proposition \ref{proposition-closed-immersion-pushforward},
on peut considérer $\mathcal{H}_Z(\mathcal{F})$ comme un faisceau abélien sur
$Z_\etale$. On obtient ainsi un foncteur
$$
\textit{Ab}(X_\etale) \longrightarrow \textit{Ab}(Z_\etale),\quad
\mathcal{F} \longmapsto \mathcal{H}_Z(\mathcal{F})
$$
qui est exact à gauche, mais qui n'est pas exact en général.

\begin{lemma}
\label{lemma-sections-with-support-acyclic}
Soit $i : Z \to X$ une immersion fermée de schémas.
Soit $\mathcal{I}$ un faisceau abélien injectif sur $X_\etale$.
Alors $\mathcal{H}_Z(\mathcal{I})$ est un faisceau abélien injectif
sur $Z_\etale$.
\end{lemma}

\begin{proof}
Remarquons que, pour tout faisceau abélien $\mathcal{G}$ sur $Z_\etale$,
on a
$$
\Hom_Z(\mathcal{G}, \mathcal{H}_Z(\mathcal{F})) =
\Hom_X(i_*\mathcal{G}, \mathcal{F})
$$
car toute section de $i_*\mathcal{G}$ a son support dans $Z$.
Comme $i_*$ est exact (section \ref{section-closed-immersions}) et que
$\mathcal{I}$ est injectif sur $X_\etale$, on en conclut que
$\mathcal{H}_Z(\mathcal{I})$ est injectif sur $Z_\etale$.
\end{proof}

\noindent
Notons
$$
R\mathcal{H}_Z : D(X_\etale) \longrightarrow D(Z_\etale)
$$
le foncteur dérivé. Posons
$\mathcal{H}^q_Z(\mathcal{F}) = R^q\mathcal{H}_Z(\mathcal{F})$, de sorte que
$\mathcal{H}^0_Z(\mathcal{F}) = \mathcal{H}_Z(\mathcal{F})$.
Par le lemme précédent, on dispose d'une suite spectrale de Grothendieck
$$
E_2^{p, q} = H^p(Z, \mathcal{H}^q_Z(\mathcal{F}))
\Rightarrow H^{p + q}_Z(X, \mathcal{F})
$$

\begin{lemma}
\label{lemma-cohomology-with-support-sheaf-on-support}
Soit $i : Z \to X$ une immersion fermée de schémas.
Soit $\mathcal{G}$ un faisceau abélien injectif sur $Z_\etale$.
Alors $\mathcal{H}^p_Z(i_*\mathcal{G}) = 0$ pour $p > 0$.
\end{lemma}

\begin{proof}
Cela est vrai parce que le foncteur $i_*$ est exact et transforme
les faisceaux abéliens injectifs en faisceaux abéliens injectifs
(Cohomologie sur les sites, Lemme
\ref{sites-cohomology-lemma-pushforward-injective-flat}).
\end{proof}

\begin{lemma}
\label{lemma-cohomology-with-support-triangle}
Soit $i : Z \to X$ une immersion fermée de schémas.
Soit $j : U \to X$ l'inclusion du complémentaire de $Z$.
Soit $\mathcal{F}$ un faisceau abélien sur $X_\etale$.
Il existe un triangle distingué
$$
i_*R\mathcal{H}_Z(\mathcal{F}) \to \mathcal{F} \to Rj_*(\mathcal{F}|_U) \to
i_*R\mathcal{H}_Z(\mathcal{F})[1]
$$
dans $D(X_\etale)$. Il donne une suite exacte
$$
0 \to i_*\mathcal{H}_Z(\mathcal{F}) \to \mathcal{F} \to
j_*(\mathcal{F}|_U) \to i_*\mathcal{H}^1_Z(\mathcal{F}) \to 0
$$
et des isomorphismes
$R^pj_*(\mathcal{F}|_U) \cong i_*\mathcal{H}^{p + 1}_Z(\mathcal{F})$
pour $p \geq 1$.
\end{lemma}

\begin{proof}
Pour obtenir le triangle distingué, choisissons une résolution injective
$\mathcal{F} \to \mathcal{I}^\bullet$. On obtient alors une suite exacte courte
de complexes
$$
0 \to
i_*\mathcal{H}_Z(\mathcal{I}^\bullet) \to \mathcal{I}^\bullet
\to j_*(\mathcal{I}^\bullet|_U) \to 0
$$
d'après la discussion précédente. Cela fournit le triangle distingué, par
Catégories dérivées, section \ref{derived-section-canonical-delta-functor}.
\end{proof}

\noindent
Soit $X$ un schéma et soit $Z \subset X$ un sous-schéma fermé.
On note $D_Z(X_\etale)$ la sous-catégorie triangulée strictement pleine et
saturée de $D(X_\etale)$ formée des complexes dont les faisceaux de cohomologie
sont à support dans $Z$. Remarquons que $D_Z(X_\etale)$ ne dépend que
de la partie fermée sous-jacente de $X$.

\begin{lemma}
\label{lemma-complexes-with-support-on-closed}
Soit $i : Z \to X$ une immersion fermée de schémas.
Le foncteur $Ri_{small, *} = i_{small, *} : D(Z_\etale) \to D(X_\etale)$
induit une équivalence $D(Z_\etale) \to D_Z(X_\etale)$ dont un quasi-inverse est
$$
i_{small}^{-1}|_{D_Z(X_\etale)} = R\mathcal{H}_Z|_{D_Z(X_\etale)}
$$
\end{lemma}

\begin{proof}
Rappelons que $i_{small}^{-1}$ et $i_{small, *}$ forment une paire de foncteurs
adjoints exacts telle que $i_{small}^{-1}i_{small, *}$ soit isomorphe au
foncteur identité sur les faisceaux abéliens. Voir la
Proposition \ref{proposition-closed-immersion-pushforward} et le
Lemme \ref{lemma-stalk-pullback}. Ainsi,
$i_{small, *} : D(Z_\etale) \to D_Z(X_\etale)$ est pleinement fidèle et
$i_{small}^{-1}$ en fournit
un inverse à gauche. D'autre part, supposons que $K$ soit un objet de
$D_Z(X_\etale)$ et considérons le morphisme d'adjonction
$K \to i_{small, *}i_{small}^{-1}K$.
Par exactitude de $i_{small, *}$ et de $i_{small}^{-1}$,
il induit sur les faisceaux de cohomologie les morphismes d'adjonction
$H^n(K) \to i_{small, *}i_{small}^{-1}H^n(K)$.
Comme ces faisceaux de cohomologie
sont à support dans $Z$, ces morphismes d'adjonction sont des isomorphismes;
on en conclut que $D(Z_\etale) \to D_Z(X_\etale)$ est une équivalence.

\medskip\noindent
Pour achever la démonstration, il faut montrer que $R\mathcal{H}_Z(K) = i_{small}^{-1}K$
si $K$ est un objet de $D_Z(X_\etale)$. Pour cela, on peut utiliser l'égalité
$K = i_{small, *}i_{small}^{-1}K$
que l'on vient de démontrer. On peut alors
choisir un représentant K-injectif $\mathcal{I}^\bullet$ de
$i_{small}^{-1}K$.
Comme $i_{small, *}$ est l'adjoint à droite du foncteur exact
$i_{small}^{-1}$, le
complexe $i_{small, *}\mathcal{I}^\bullet$ est K-injectif
(Catégories dérivées, Lemme \ref{derived-lemma-adjoint-preserve-K-injectives}).
On voit que $R\mathcal{H}_Z(K)$ est calculé par
$\mathcal{H}_Z(i_{small, *}\mathcal{I}^\bullet) = \mathcal{I}^\bullet$
comme voulu.
\end{proof}

\begin{lemma}
\label{lemma-cohomology-with-support-quasi-coherent}
Soit $X$ un schéma. Soit $Z \subset X$ un sous-schéma fermé, et soit $i : Z \to X$ l'immersion canonique.
Soit $\mathcal{F}$ un $\mathcal{O}_X$-module quasi-cohérent
et notons $\mathcal{F}^a$ le faisceau quasi-cohérent associé
sur le petit site étale de $X$
(Proposition \ref{proposition-quasi-coherent-sheaf-fpqc}). Alors
\begin{enumerate}
\item $H^q_Z(X, \mathcal{F})$ s'identifie à $H^q_Z(X_\etale, \mathcal{F}^a)$;
\item si le complémentaire de $Z$ est rétrocompact dans $X$, alors
$i_*\mathcal{H}^q_Z(\mathcal{F}^a)$ est un faisceau quasi-cohérent de
$\mathcal{O}_X$-modules égal à $(i_*\mathcal{H}^q_Z(\mathcal{F}))^a$.
\end{enumerate}
\end{lemma}

\begin{proof}
Soit $j : U \to X$ l'inclusion du complémentaire de $Z$.
L'assertion (1) sur les groupes de cohomologie résulte des longues
suites exactes de cohomologie à supports et des identifications
$H^q(X_\etale, \mathcal{F}^a) = H^q(X, \mathcal{F})$ et
$H^q(U_\etale, \mathcal{F}^a) = H^q(U, \mathcal{F})$; voir le
Théorème \ref{theorem-zariski-fpqc-quasi-coherent}.
Si $j : U \to X$ est un morphisme quasi-compact, c'est-à-dire si $U \subset X$
est rétrocompact, alors $R^qj_*$ transforme les faisceaux quasi-cohérents
en faisceaux quasi-cohérents
(Cohomologie des schémas, Lemme
\ref{coherent-lemma-quasi-coherence-higher-direct-images})
et commute à la formation du
faisceau associé sur les sites étales
(Descente, Lemme \ref{descent-lemma-higher-direct-images-small-etale}).
On conclut en appliquant le
Lemme \ref{lemma-cohomology-with-support-triangle}.
\end{proof}







\section{Schémas dont les anneaux locaux sont strictement henséliens}
\sectionmark{Anneaux locaux strictement henséliens}
\label{section-strictly-henselian-local-rings}

\noindent
Dans cette section, nous rassemblons quelques résultats sur la cohomologie étale
des schémas dont les anneaux locaux sont strictement henséliens. Voici par exemple
une généralisation utile du
Lemme \ref{lemma-vanishing-etale-cohomology-strictly-henselian}.

\begin{lemma}
\label{lemma-local-rings-strictly-henselian}
Soit $S$ un schéma dont tous les anneaux locaux sont strictement henséliens.
Alors, pour tout faisceau abélien $\mathcal{F}$ sur $S_\etale$, on a
$H^i(S_\etale, \mathcal{F}) = H^i(S_{Zar}, \mathcal{F})$.
\end{lemma}

\begin{proof}
Soit $\epsilon : S_\etale \to S_{Zar}$ le morphisme de sites défini
par le foncteur d'inclusion. Le faisceau de Zariski $R^p\epsilon_*\mathcal{F}$
est le faisceau associé au préfaisceau $U \mapsto H^p_\etale(U, \mathcal{F})$.
Sa fibre en $x \in S$ est donc
$\colim H^p_\etale(U, \mathcal{F}) =
H^p_\etale(\Spec(\mathcal{O}_{S, x}), \mathcal{G}_x)$,
où $\mathcal{G}_x$ désigne l'image inverse de $\mathcal{F}$
sur $\Spec(\mathcal{O}_{S, x})$; voir le
Lemme \ref{lemma-directed-colimit-cohomology}.
Ainsi, les faisceaux d'images directes supérieures $R^p\epsilon_*\mathcal{F}$ sont
nuls d'après le
Lemme \ref{lemma-vanishing-etale-cohomology-strictly-henselian},
et l'on conclut par la suite spectrale de Leray.
\end{proof}

\begin{lemma}
\label{lemma-gabber-criterion}
Soit $R$ un anneau dont tous les anneaux locaux sont strictement henséliens.
Soit $\mathcal{F}$ un faisceau sur $\Spec(R)_\etale$.
Supposons que, pour tous $f,g\in R$, le noyau de
$$
H^1_\etale(D(f + g), \mathcal{F})
\longrightarrow
H^1_\etale(D(f(f + g)), \mathcal{F}) \oplus
H^1_\etale(D(g(f + g)), \mathcal{F})
$$
soit nul. Alors $H^q_\etale(\Spec(R), \mathcal{F}) = 0$ pour $q > 0$.
\end{lemma}

\begin{proof}
D'après le Lemme \ref{lemma-local-rings-strictly-henselian}, la
cohomologie étale de $\mathcal{F}$ coïncide avec sa cohomologie de Zariski
sur tout ouvert de $\Spec(R)$. Nous démontrerons par récurrence sur $i$
l'assertion suivante : pour $h \in R$, on a $H^q_\etale(D(h), \mathcal{F}) = 0$
si $1 \leq q \leq i$. Le cas initial $i = 0$ est trivial. Supposons $i \geq 1$.

\medskip\noindent
Soient $\xi \in H^q_\etale(D(h), \mathcal{F})$, avec $1 \leq q \leq i$,
et $h \in R$. Si $q < i$, la récurrence conclut; supposons donc $q = i$.
Après avoir remplacé $R$ par $R_h$, on peut supposer que
$\xi \in H^i_\etale(\Spec(R), \mathcal{F})$; nous omettons quelques détails.
Soit $I \subset R$ l'ensemble des éléments $f \in R$ tels que
$\xi|_{D(f)} = 0$. Comme $\xi$ est localement nulle pour la topologie de Zariski,
pour tout idéal premier $\mathfrak p$ de $R$, il existe $f \in I$
tel que $f \not \in \mathfrak p$. Si l'on montre que $I$ est un
idéal, alors $1 \in I$ et la preuve sera achevée. Il est clair que
$f \in I$ et $r \in R$ impliquent $rf \in I$. Supposons donc $f,g \in I$
et montrons que $f + g \in I$. Si $q = i = 1$, c'est exactement
l'hypothèse du lemme, d'où le résultat pour $i = 1$. Si
$q = i > 1$, remarquons que
$$
D(f + g) = D(f(f + g)) \cup D(g(f + g))
$$
Par Mayer--Vietoris (Cohomologie, Lemme \ref{cohomology-lemma-mayer-vietoris},
qui s'applique puisque, sur les sous-schémas ouverts de $\Spec(R)$, la cohomologie
étale coïncide avec la cohomologie de Zariski), on dispose d'une suite exacte
$$
\xymatrix{
H^{i - 1}_\etale(D(fg(f + g)), \mathcal{F}) \ar[d] \\
H^i_\etale(D(f + g), \mathcal{F}) \ar[d] \\
H^i_\etale(D(f(f + g)), \mathcal{F}) \oplus
H^i_\etale(D(g(f + g)), \mathcal{F})
}
$$
et le résultat s'ensuit, puisque le premier groupe est nul par récurrence.
\end{proof}

\begin{lemma}
\label{lemma-affine-only-closed-points}
Soit $S$ un schéma affine tel que
(1) tous ses points soient fermés et (2) tous ses corps résiduels soient
séparablement clos. Alors,
pour tout faisceau abélien $\mathcal{F}$ sur $S_\etale$, on a
$H^i(S_\etale, \mathcal{F}) = 0$ pour $i > 0$.
\end{lemma}

\begin{proof}
La condition (1) implique que l'espace topologique sous-jacent
à $S$ est profini; voir
Algèbre, Lemme \ref{algebra-lemma-ring-with-only-minimal-primes}.
Les groupes de cohomologie supérieure d'un faisceau abélien sur l'espace
topologique $S$ (c'est-à-dire la cohomologie de Zariski) sont donc nuls; voir
Cohomologie, Lemme \ref{cohomology-lemma-vanishing-for-profinite}.
Les anneaux locaux sont strictement henséliens d'après
Algèbre, Lemme \ref{algebra-lemma-local-dimension-zero-henselian}.
La cohomologie étale de $S$ se calcule donc par la cohomologie de Zariski
d'après le Lemme \ref{lemma-local-rings-strictly-henselian},
ce qui achève la démonstration.
\end{proof}

\noindent
Le spectre d'un anneau absolument intégralement clos
est un exemple de schéma dont tous les anneaux locaux sont
strictement henséliens; voir Compléments d'algèbre, Lemme
\ref{more-algebra-lemma-absolutely-integrally-closed-strictly-henselian}.
Les domaines normaux à corps des fractions séparablement clos
possèdent même une propriété plus forte, comme l'explique le
lemme suivant.

\begin{lemma}
\label{lemma-normal-scheme-with-alg-closed-function-field}
Soit $X$ un schéma intègre normal à corps des fonctions
séparablement clos.
\begin{enumerate}
\item Tout morphisme étale séparé $U \to X$ est une
réunion disjointe d'immersions ouvertes.
\item Tous les anneaux locaux de $X$ sont strictement henséliens.
\end{enumerate}
\end{lemma}

\begin{proof}
Soit $R$ un domaine normal dont le corps des fractions est séparablement
clos. Soit $R \to A$ un homomorphisme d'anneaux étale. Alors
$A \otimes_R K$ est, comme $K$-algèbre, un produit fini
$\prod_{i=1}^n K$ de copies de $K$. Soient $e_i$, $i = 1, \ldots, n$,
les idempotents correspondants de $A \otimes_R K$. Comme $A$ est normal
(Algèbre, Lemme \ref{algebra-lemma-normal-goes-up}),
les idempotents $e_i$ appartiennent à $A$
(Algèbre, Lemme \ref{algebra-lemma-normal-ring-integrally-closed}).
Ainsi $A = \prod Ae_i$, et l'on peut supposer $A \otimes_R K = K$.
Puisque $A \subset A \otimes_R K = K$ (par platitude de $R \to A$ et
parce que $R \subset K$), on conclut que $A$ est un domaine.
Le même argument donne
$A \otimes_R A \subset (A \otimes_R A) \otimes_R K = K$.
L'application $A \otimes_R A \to A$ est donc
injective aussi bien que surjective. Ainsi, $R \to A$ définit une
immersion ouverte d'après
Morphismes, Lemme \ref{morphisms-lemma-universally-injective},
et
Morphismes étales, Théorème \ref{etale-theorem-etale-radicial-open}.

\medskip\noindent
Soit $f : U \to X$ un morphisme étale séparé. Soit $\eta \in X$
le point générique, et posons $f^{-1}(\{\eta\}) = \{\xi_i\}_{i \in I}$.
Le résultat du paragraphe précédent montre ceci :
pour tout ouvert affine $U' \subset U$ dont l'image dans $X$ est contenue dans
un ouvert affine, on a $U' = \coprod_{i \in I} U'_i$, où $U'_i$
est l'ensemble des points de $U'$ qui sont des spécialisations de $\xi_i$.
De plus, le morphisme $U'_i \to X$ est une immersion ouverte.
Il s'ensuit que $U_i = \overline{\{\xi_i\}}$ est un sous-schéma ouvert et fermé
de $U$, et que $U_i \to X$ est localement, sur la source,
un isomorphisme. D'après Morphismes,
Lemme \ref{morphisms-lemma-distinct-local-rings},
le fait que $U_i \to X$ soit séparé implique que
$U_i \to X$ est injectif; on conclut que $U_i \to X$
est une immersion ouverte, ce qui prouve (1).

\medskip\noindent
La partie (2) résulte de la partie (1) et de la description de l'hensélisé
strict de $\mathcal{O}_{X, x}$ comme anneau local en $\overline{x}$
sur le site étale de $X$ (Lemme \ref{lemma-describe-etale-local-ring}).
On peut aussi la démontrer directement; voir
Groupes fondamentaux, Lemme
\ref{pione-lemma-normal-local-domain-separablly-closed-fraction-field}.
\end{proof}

\begin{lemma}
\label{lemma-Rf-star-zero-normal-with-alg-closed-function-field}
Soit $f : X \to Y$ un morphisme de schémas, où $X$ est un schéma intègre
normal à corps des fonctions séparablement clos. Alors
$R^qf_*\underline{M} = 0$ pour $q > 0$ et tout groupe abélien $M$.
\end{lemma}

\begin{proof}
Rappelons que $R^qf_*\underline{M}$ est le faisceau associé
au préfaisceau $V \mapsto H^q_\etale(V \times_Y X, M)$ sur $Y_\etale$; voir le
Lemme \ref{lemma-higher-direct-images}.
Si $V$ est affine, alors $V \times_Y X \to X$ est séparé et étale.
Par conséquent, $V \times_Y X = \coprod U_i$ est une réunion disjointe de
sous-schémas ouverts $U_i$ de $X$; voir le
Lemme \ref{lemma-normal-scheme-with-alg-closed-function-field}.
D'après le Lemme \ref{lemma-local-rings-strictly-henselian},
$H^q_\etale(U_i, M)$ est égal à
$H^q_{Zar}(U_i, M)$. Ce dernier est nul d'après
Cohomologie, Lemme \ref{cohomology-lemma-irreducible-constant-cohomology-zero}.
\end{proof}

\begin{lemma}
\label{lemma-closed-of-affine-normal-scheme-with-alg-closed-function-field}
Soit $X$ un schéma affine intègre normal à corps des fonctions
séparablement clos. Soit $Z \subset X$ un sous-schéma fermé. Soit
$V \to Z$ un morphisme étale, avec $V$ affine. Alors $V$ est une réunion
disjointe finie de sous-schémas ouverts de $Z$. Si $V \to Z$ est
surjectif et fini étale, alors $V \to Z$ admet une section.
\end{lemma}

\begin{proof}
D'après Algèbre, Lemme \ref{algebra-lemma-lift-etale},
on peut relever $V$ en un schéma affine $U$ étale sur $X$.
Appliquons le Lemme \ref{lemma-normal-scheme-with-alg-closed-function-field}
à $U \to X$ pour obtenir la première assertion.

\medskip\noindent
La dernière assertion est une conséquence de la première.
Soit $V = \coprod_{i=1}^n V_i$ une décomposition finie
en sous-schémas ouverts et
fermés telle que $V_i \to Z$ soit une immersion ouverte.
Comme $V \to Z$ est fini, $V_i \to Z$ est aussi fermé.
Soit $U_i \subset Z$ son image. On a alors une décomposition
en sous-schémas ouverts et fermés
$$
Z =
\coprod\nolimits_{(A, B)}
\bigcap\nolimits_{i \in A} U_i \cap
\bigcap\nolimits_{i \in B} U_i^c
$$
où la réunion disjointe porte sur les partitions $\{1, \ldots, n\} = A \amalg B$
pour lesquelles $A$ est non vide.
Chacune des strates est contenue dans l'un des $U_i$, ce qui
fournit la section cherchée.
\end{proof}

\begin{lemma}
\label{lemma-gabber-for-h1-absolutely-algebraically-closed}
Soit $X$ un schéma affine intègre normal à corps des fonctions
séparablement clos. Soit $Z \subset X$ un sous-schéma fermé.
Pour tout groupe abélien fini $M$, on a $H^1_\etale(Z, \underline{M}) = 0$.
\end{lemma}

\begin{proof}
D'après Cohomologie sur les sites, Lemme \ref{sites-cohomology-lemma-torsors-h1},
un élément de $H^1_\etale(Z, \underline{M})$ correspond à un
$\underline{M}$-torseur $\mathcal{F}$ sur $Z_\etale$.
Un tel torseur est manifestement un faisceau localement constant fini.
Ainsi, $\mathcal{F}$ est représentable par un schéma $V$ fini
étale sur $Z$; voir le Lemme \ref{lemma-characterize-finite-locally-constant}.
Le morphisme $V \to Z$ est surjectif, puisque tout torseur est localement trivial.
Comme $V \to Z$ admet une section d'après le
Lemme \ref{lemma-closed-of-affine-normal-scheme-with-alg-closed-function-field},
la preuve est achevée.
\end{proof}

\begin{lemma}
\label{lemma-gabber-for-absolutely-algebraically-closed}
Soit $X$ un schéma affine intègre normal à corps des fonctions
séparablement clos. Soit $Z \subset X$ un sous-schéma fermé.
Pour tout groupe abélien fini $M$, on a
$H^q_\etale(Z, \underline{M}) = 0$ pour $q \geq 1$.
\end{lemma}

\begin{proof}
Écrivons $X = \Spec(R)$ et $Z = \Spec(R')$, de sorte que l'on dispose d'une surjection
d'anneaux $R \to R'$. Tous les anneaux locaux de $R'$ sont strictement henséliens d'après le
Lemme \ref{lemma-normal-scheme-with-alg-closed-function-field} et
Algèbre, Lemme \ref{algebra-lemma-quotient-strict-henselization}.
De plus, pour tout $f' \in R'$, il existe une surjection
$R_f \to R'_{f'}$, où $f \in R$ est un relèvement de $f'$.
Comme $R_f$ est un domaine normal à corps des fractions séparablement clos,
on a $H^1_\etale(D(f'), \underline{M}) = 0$ d'après le
Lemme \ref{lemma-gabber-for-h1-absolutely-algebraically-closed}.
On peut donc appliquer le Lemme \ref{lemma-gabber-criterion} à $Z = \Spec(R')$
pour conclure.
\end{proof}

\begin{lemma}
\label{lemma-integral-cover-trivial-cohomology}
Soit $X$ un schéma affine.
\begin{enumerate}
\item Il existe un morphisme entier surjectif $X' \to X$ tel que, pour
tout sous-schéma fermé $Z' \subset X'$, tout groupe abélien fini $M$ et
tout $q \geq 1$, on ait $H^q_\etale(Z', \underline{M}) = 0$.
\item Pour tout sous-schéma fermé $Z \subset X$, tout groupe abélien fini $M$,
tout $q \geq 1$ et tout $\xi \in H^q_\etale(Z, \underline{M})$, il existe un
morphisme fini surjectif $X' \to X$ de présentation finie tel que
l'image inverse de $\xi$ soit nulle dans $H^q_\etale(X' \times_X Z, \underline{M})$.
\end{enumerate}
\end{lemma}

\begin{proof}
Écrivons $X = \Spec(A)$. Écrivons $A = \mathbf{Z}[x_i]/J$ pour un idéal $J$.
Soit $R$ la clôture intégrale de $\mathbf{Z}[x_i]$ dans une clôture
algébrique du corps des fractions de $\mathbf{Z}[x_i]$. Posons
$A' = R/JR$ et $X' = \Spec(A')$. Le Lemme
\ref{lemma-gabber-for-absolutely-algebraically-closed} montre que cet exemple satisfait (1).

\medskip\noindent
Démontrons (2). Soit $X' \to X$ le morphisme entier surjectif construit
ci-dessus. L'image de $\xi$ dans
$H^q_\etale(X' \times_X Z, \underline{M})$ est bien nulle. On peut écrire $X'$ comme une
limite projective $X' = \lim X'_i$ de schémas finis et de présentation finie
sur $X$; cela se fait aisément dans le présent cas affine, mais c'est
un cas particulier du lemme plus général Limites, Lemme
\ref{limits-lemma-integral-limit-finite-and-finite-presentation}.
D'après le Lemme \ref{lemma-directed-colimit-cohomology},
l'image de $\xi$ dans $H^q_\etale(X'_i \times_X Z, \underline{M})$ est nulle
pour un indice $i$ assez grand.
\end{proof}








\section{Annulation pour les anneaux absolument intégralement clos}
\sectionmark{Anneaux absolument intégralement clos}
\label{section-aic-vanishing}

\noindent
Rappelons qu'un anneau $R$ est dit absolument intégralement clos
si tout polynôme unitaire à coefficients dans $R$ a une racine dans $R$
(Compléments d'algèbre, définition
\ref{more-algebra-definition-absolutely-integrally-closed}).
Dans cette section, nous démontrons que la cohomologie étale de
$\Spec(R)$ à coefficients dans un groupe de torsion fini
est nulle en degrés strictement positifs (proposition \ref{proposition-aic-vanishing}),
ce qui améliore légèrement le
lemme antérieur \ref{lemma-gabber-for-absolutely-algebraically-closed}.
Nous suggérons au lecteur de sauter cette section.

\begin{lemma}
\label{lemma-find-extension}
Soit $A$ un anneau. Soient $a, b \in A$ tels que
$aA + bA = A$ et que $a \bmod bA$ soit une racine de l'unité.
Il existe alors une extension monogène $A \subset B$
et un élément $y \in B$ tels que $u = a - by$ soit inversible.
\end{lemma}

\begin{proof}
Écrivons $a^n \equiv 1 \bmod bA$. En particulier, $a^i$ est inversible modulo
$b^mA$ pour tous $i, m \geq 1$. Nous affirmons qu'il existe des éléments
$a_1, \ldots, a_n \in A$ tels que
$$
1 = a^n + a_1 a^{n - 1}b + a_2 a^{n - 2}b^2 + \ldots + a_n b^n
$$
En effet, puisque $1 - a^n \in bA$, nous pouvons trouver un élément $a_1 \in A$
tel que $1 - a^n - a_1 a^{n - 1} b \in b^2A$, grâce à l'inversibilité
de $a^{n - 1}$ modulo $bA$. Ensuite, nous pouvons trouver $a_2 \in A$
tel que $1 - a^n - a_1 a^{n - 1} b - a_2 a^{n - 2} b^2 \in b^3A$.
Et ainsi de suite. Nous finissons par trouver $a_1, \ldots, a_{n - 1} \in A$
tels que $1 - (a^n + a_1 a^{n - 1}b + a_2 a^{n - 2}b^2 +
\ldots + a_{n - 1} ab^{n - 1}) \in b^nA$. Nous pouvons alors trouver
$a_n \in A$ de sorte que l'égalité affichée soit satisfaite.

\medskip\noindent
Pour $a_1, \ldots, a_n$ comme ci-dessus, nous affirmons que le choix suivant convient :
$$
B = A[y]/(y^n + a_1 y^{n - 1} + a_2 y^{n - 2} + \ldots + a_n)
$$
En effet, supposons que $\mathfrak q \subset B$ soit un idéal premier
situé au-dessus de $\mathfrak p \subset A$. Pour obtenir une contradiction,
supposons que $u = a - by$ appartienne à $\mathfrak q$. Si $b \in \mathfrak p$,
alors $a \not \in \mathfrak p$ puisque $aA + bA = A$, et donc $u$
n'appartient pas à $\mathfrak q$. Nous pouvons ainsi supposer que $b \not \in \mathfrak p$, c'est-à-dire
que $b \not \in \mathfrak q$. Il s'ensuit que $y \bmod \mathfrak q$
est égal à $a/b \bmod \mathfrak q$. Mais nous obtenons alors
$$
0 = y^n + a_1 y^{n - 1} + a_2 y^{n - 2} + \ldots + a_n =
b^{-n}(a^n + a_1 a^{n - 1}b + a_2 a^{n - 2}b^2 + \ldots + a_nb^n) =
b^{-n}
$$
une contradiction. Cela achève la preuve.
\end{proof}

\noindent
Pour expliquer la preuve, nous devons introduire quelques schémas en groupes.
Fixons un nombre premier $\ell$. Posons
$$
A = \mathbf{Z}[\zeta] =
\mathbf{Z}[x]/(x^{\ell - 1} + x^{\ell - 2} + \ldots + 1)
$$
Autrement dit, $A$ est l'extension monogène de $\mathbf{Z}$ engendrée par
une racine primitive $\ell$-ième de l'unité $\zeta$. Posons
$$
\pi = \zeta - 1
$$
Un calcul (omis) montre que $\ell$ est divisible par $\pi^{\ell - 1}$
dans $A$. Notre premier schéma en groupes sur $A$ est
$$
G = \Spec(A[s, \frac{1}{\pi s + 1}])
$$
dont la loi de groupe est donnée par la comultiplication
$$
\mu :
A[s, \frac{1}{\pi s + 1}]
\longrightarrow
A[s, \frac{1}{\pi s + 1}] \otimes_A A[s, \frac{1}{\pi s + 1}],\quad
s \longmapsto \pi s \otimes s + s \otimes 1 + 1 \otimes s
$$
Avec ce choix, nous avons
$$
\mu(\pi s + 1) = (\pi s + 1) \otimes (\pi s + 1)
$$
et nous avons donc bien le morphisme de $A$-algèbres indiqué. Nous omettons
la vérification que cela définit bien une loi de groupe. Notre second schéma
en groupes sur $A$ est
$$
H = \Spec(A[t, \frac{1}{\pi^\ell t + 1}])
$$
dont la loi de groupe est donnée par la comultiplication
$$
\mu : A[t, \frac{1}{\pi^\ell t + 1}]
\longrightarrow
A[t, \frac{1}{\pi^\ell t + 1}] \otimes_A A[t, \frac{1}{\pi^\ell t + 1}],\quad
t \longmapsto \pi^\ell t \otimes t + t \otimes 1 + 1 \otimes t
$$
La même vérification que précédemment montre que cela définit une loi de groupe.
Observons ensuite que le polynôme
$$
\Phi(s) = \frac{(\pi s + 1)^\ell - 1}{\pi^\ell}
$$
appartient à $A[s]$, est de degré $\ell$ et est unitaire en $s$. En effet, le coefficient
de $s^i$, pour $0 < i < \ell$, est égal à ${\ell \choose i}\pi^{i - \ell}$;
comme $\pi^{\ell - 1}$ divise $\ell$ dans $A$, c'est un élément de $A$.
Nous obtenons un morphisme d'anneaux
$$
A[t, \frac{1}{\pi^\ell t + 1}]
\longrightarrow
A[s, \frac{1}{\pi s + 1}],\quad
t \longmapsto \Phi(s)
$$
dont le lecteur vérifiera facilement la compatibilité avec les comultiplications.
Nous obtenons ainsi un morphisme de schémas en groupes
$$
f : G \to H
$$
Le lemme suivant montre en particulier que ce morphisme est fidèlement
plat (nous verrons en fait qu'il est fini étale surjectif).

\begin{lemma}
\label{lemma-monogenic-one}
Nous avons
$$
A[s, \frac{1}{\pi s + 1}] =
\left(A[t, \frac{1}{\pi^\ell t + 1}]\right)[s]/(\Phi(s) - t)
$$
En particulier, l'algèbre de Hopf de $G$ est une extension monogène
de l'algèbre de Hopf de $H$.
\end{lemma}

\begin{proof}
Cela résulte de la discussion ci-dessus et de la forme de $\Phi(s)$.
Notons en particulier que, par la relation $\Phi(s) = t$, l'élément
$\frac{1}{\pi^\ell t + 1}$ devient l'élément
$\frac{1}{(\pi s + 1)^\ell}$.
\end{proof}

\noindent
Calculons maintenant le noyau de $f$. Comme l'origine de $H$ est
donnée par $t = 0$ dans $H$, nous voyons que le noyau de $f$ est défini
par $\Phi(s) = 0$. Observons que les points à valeurs dans $A$
$\sigma_0, \ldots, \sigma_{\ell - 1}$
de $G$ donnés par
$$
\sigma_i :
s = \frac{\zeta^i - 1}{\pi} = \frac{\zeta^i - 1}{\zeta - 1} =
\zeta^{i - 1} + \zeta^{i - 2} + \ldots + 1,\quad
i = 0, 1, \ldots, \ell - 1
$$
appartiennent certainement à $\Ker(f)$. De plus, ils sont deux à deux
distincts dans {\bf toutes} les fibres de $G \to \Spec(A)$. Le lecteur
calculera aussi que $\sigma_i +_G \sigma_j = \sigma_{i + j \bmod \ell}$.
Nous obtenons donc une immersion fermée de schémas en groupes
$$
\underline{\mathbf{Z}/\ell \mathbf{Z}}_A \longrightarrow \Ker(f)
$$
qui envoie $i$ sur $\sigma_i$. Or, par construction, $\Ker(f)$
est fini et plat sur $\Spec(A)$, de degré $\ell$. Nous en concluons
que ce morphisme est un isomorphisme. En définitive, nous obtenons
une suite exacte courte
\begin{equation}
\label{equation-ses}
0 \to
\underline{\mathbf{Z}/\ell \mathbf{Z}}_A
\to G
\to H
\to 0
\end{equation}
de schémas en groupes sur $A$.

\begin{lemma}
\label{lemma-lift-points-H-to-G}
Soit $R$ une $A$-algèbre absolument intégralement close. Alors
$G(R) \to H(R)$ est surjectif.
\end{lemma}

\begin{proof}
Soit $h \in H(R)$ correspondant au morphisme de $A$-algèbres
$A[t, \frac{1}{\pi^\ell t + 1}] \to R$ qui envoie $t$ sur $a \in R$.
Comme $\Phi(s)$ est unitaire, nous pouvons trouver $b \in R$
tel que $\Phi(b) = a$. Par le lemme \ref{lemma-monogenic-one},
en envoyant $s$ sur $b$, nous obtenons un unique morphisme de $A$-algèbres
$A[s, \frac{1}{\pi s + 1}] \to R$ compatible
avec le morphisme $A[t, \frac{1}{\pi^\ell t + 1}] \to R$ ci-dessus.
Celui-ci correspond à son tour à un élément $g \in G(R)$
qui s'envoie sur $h \in H(R)$.
\end{proof}

\begin{lemma}
\label{lemma-interpolate}
Soit $R$ une $A$-algèbre absolument intégralement close.
Soient $I, J \subset R$ des idéaux tels que $I + J = R$.
Il existe un élément $g \in G(R)$ tel
que $g \bmod I = \sigma_0$ et $g \bmod J = \sigma_1$.
\end{lemma}

\begin{proof}
Choisissons $x \in I$ tel que $x \equiv 1 \bmod J$.
Nous pouvons remplacer, et remplaçons, $I$ par $xR$ et $J$ par $(x - 1)R$.
Nous cherchons alors un élément $s \in R$ tel que
\begin{enumerate}
\item $1 + \pi s$ soit inversible,
\item $s \equiv 0 \bmod xR$, et
\item $s \equiv 1 \bmod (x - 1)R$.
\end{enumerate}
Les deux dernières conditions signifient que $s = x + x(x - 1)y$ pour un $y \in R$.
La première exige que $1 + \pi s = 1 + \pi x + \pi x (x - 1) y$
soit inversible dans $R$.
Or $1 + \pi x$ et $\pi x (x - 1)$ engendrent l'idéal
unité de $R$, et $1 + \pi x$ est une racine $\ell$-ième de $1$
modulo $\pi x (x - 1)$\footnote{En effet, $1 + \pi x$ est congru
à $1$ modulo $\pi$, à $1$ modulo $x$ et à
$1 + \pi = \zeta$ modulo $x - 1$, et nous avons
$(\pi) \cap (x) \cap (x - 1) = (\pi x (x - 1))$ dans $A[x]$.}.
Le lemme \ref{lemma-find-extension} et le fait que $R$ soit
absolument intégralement clos donnent donc le résultat.
\end{proof}

\begin{proposition}
\label{proposition-aic-vanishing}
Soit $R$ un anneau absolument intégralement clos.
Soit $M$ un groupe abélien fini.
Alors $H^i_\etale(\Spec(R), \underline{M}) = 0$ pour $i > 0$.
\end{proposition}

\begin{proof}
Comme tout groupe abélien fini admet une filtration finie dont les
sous-quotients sont cycliques d'ordre premier, nous pouvons supposer
que $M = \mathbf{Z}/\ell\mathbf{Z}$, où $\ell$ est un nombre premier.

\medskip\noindent
Observons que tous les anneaux locaux de $R$ sont strictement henséliens; voir
Compléments d'algèbre, lemme
\ref{more-algebra-lemma-absolutely-integrally-closed-strictly-henselian}.
De plus, tout localisé de $R$ est encore absolument intégralement clos
par Compléments d'algèbre, lemme
\ref{more-algebra-lemma-absolutely-integrally-closed-quotient-localization}.
Ainsi, le lemme \ref{lemma-gabber-criterion} nous dit qu'il suffit de
montrer que le noyau de
$$
H^1_\etale(D(f + g), \mathbf{Z}/\ell\mathbf{Z})
\longrightarrow
H^1_\etale(D(f(f + g)), \mathbf{Z}/\ell\mathbf{Z}) \oplus
H^1_\etale(D(g(f + g)), \mathbf{Z}/\ell\mathbf{Z})
$$
est nul pour tous $f, g \in R$. Après avoir remplacé $R$ par $R_{f + g}$, nous
nous ramenons à l'affirmation suivante : si
$\xi \in H^1_\etale(\Spec(R), \mathbf{Z}/\ell\mathbf{Z})$
et si $\Spec(R) = U \cup V$ est un recouvrement par deux ouverts affines
tel que $\xi|_U$ et $\xi|_V$ soient triviaux, alors $\xi = 0$.

\medskip\noindent
Soit $A = \mathbf{Z}[\zeta]$ comme ci-dessus.
Puisque $\mathbf{Z} \subset A$ est monogène, nous pouvons trouver un morphisme
d'anneaux $A \to R$. Dorénavant, nous considérons $R$ comme une $A$-algèbre
et $\Spec(R)$ comme un schéma sur $\Spec(A)$.
Après changement de base de la suite exacte courte (\ref{equation-ses})
à $\Spec(R)$ et passage à la cohomologie étale, nous obtenons
$$
G(R) \to H(R) \to
H^1_\etale(\Spec(R), \mathbf{Z}/\ell\mathbf{Z}) \to
H^1_\etale(\Spec(R), G)
$$
Nous utiliserons cela dans la suite de la preuve.

\medskip\noindent
Soit $\tau \in \Gamma(U \cap V, \mathbf{Z}/\ell\mathbf{Z})$
une section dont l'image par l'homomorphisme bord de la suite de Mayer--Vietoris
(lemme \ref{lemma-mayer-vietoris}) est $\xi$.
Pour $i = 0, 1, \ldots, \ell - 1$, soit $A_i \subset U \cap V$
la partie ouverte et fermée où $\tau$ prend la valeur $i \bmod \ell$.
Nous avons donc une décomposition en union disjointe finie
$$
U \cap V = A_0 \amalg \ldots \amalg A_{\ell - 1}
$$
telle que $\tau$ soit constante sur chaque $A_i$. Pour
$i = 0, 1, \ldots, \ell - 1$, notons
$\tau_i \in H^0(U \cap V, \mathbf{Z}/\ell\mathbf{Z})$
l'élément égal à $1$ sur $A_i$ et
à $0$ sur $A_j$ pour $j \not = i$.
Alors $\tau$ est une somme de multiples des $\tau_i$\footnote{À d'éventuelles
erreurs de calcul près, nous avons $\tau = \sum i \tau_i$.}.
Il suffit donc de montrer que la classe de cohomologie correspondant
à $\tau_i$ est triviale. Nous sommes ainsi ramenés au cas où
$\tau$ ne prend que deux valeurs distinctes, à savoir $1$ et $0$.

\medskip\noindent
Supposons que $\tau$ ne prenne que les valeurs $1$ et $0$. Écrivons
$$
U \cap V = A \amalg B
$$
où $A$ est le lieu où $\tau = 0$ et $B$ le lieu
où $\tau = 1$. Alors $A$ et $B$ sont des parties fermées disjointes.
Notons $\overline{A}$ et $\overline{B}$ les adhérences de $A$ et $B$ dans
$\Spec(R)$. Nous avons alors une ``banane'', c'est-à-dire
que nous avons
$$
\overline{A} \cap \overline{B} = Z_1 \amalg Z_2
$$
où $Z_1 \subset U$ et $Z_2 \subset V$ sont des parties fermées disjointes.
Posons $T_1 = \Spec(R) \setminus V$ et $T_2 = \Spec(R) \setminus U$.
Observons que $Z_1 \subset T_1 \subset U$, que $Z_2 \subset T_2 \subset V$ et que
$T_1 \cap T_2 = \emptyset$. Au niveau des espaces topologiques, nous pouvons écrire
$$
\Spec(R) = \overline{A} \cup \overline{B} \cup T_1 \cup T_2
$$
Nous suggérons de faire un dessin pour visualiser cette situation.
Pour prouver que $\xi$ est nul, nous pouvons remplacer, et remplaçons, $R$ par
son réduit (proposition \ref{proposition-topological-invariance}).
Dans la suite, nous considérons $A$, $\overline{A}$, $B$, $\overline{B}$, $T_1$, $T_2$
comme des sous-schémas fermés réduits de $\Spec(R)$. Munissons ensuite $Z_1$
et $Z_2$ des structures de schémas
$$
Z_1 = \overline{A} \cap (\overline{B} \cup T_1)
\quad\text{et}\quad
Z_2 = \overline{A} \cap (\overline{B} \cup T_2)
$$
(réunions et intersections schématiques comme dans Morphismes, définition
\ref{morphisms-definition-scheme-theoretic-intersection-union}).

\medskip\noindent
Notons $X$ le torseur sous $G$ sur $\Spec(R)$ correspondant à l'image
de $\xi$ dans $H^1(\Spec(R), G)$. Si $X$ est trivial, alors $\xi$
provient d'un élément $h \in H(R)$ (voir la suite exacte de cohomologie
ci-dessus). Mais le lemme \ref{lemma-lift-points-H-to-G} permet alors
de relever $h$ en un élément de $G(R)$, et nous concluons
que $\xi = 0$, comme souhaité. Notre but est donc de prouver que $X$ est trivial.

\medskip\noindent
Rappelons que le plongement $\mathbf{Z}/\ell \mathbf{Z} \to G(R)$
envoie $i \bmod \ell$ sur $\sigma_i \in G(R)$. Observons que $\overline{A}$
est le spectre d'un anneau absolument intégralement clos (à savoir
un quotient de $R$). Par le
lemme \ref{lemma-interpolate}, nous pouvons trouver $g \in G(\overline{A})$ tel que
$g|_{\overline{A} \cap Z_1} = \sigma_0$ et
$g|_{\overline{A} \cap Z_2} = \sigma_1$ (schématiquement).
Nous pouvons alors définir
\begin{enumerate}
\item $g_1 \in G(U)$ qui vaut $g$ sur $\overline{A} \cap U$,
$\sigma_0$ sur $\overline{B} \cap U$ et $\sigma_0$ sur $T_1$, et
\item $g_2 \in G(V)$ qui vaut $g$ sur $\overline{A} \cap V$,
$\sigma_1$ sur $\overline{B} \cap V$ et $\sigma_1$ sur $T_2$.
\end{enumerate}
En effet, pour trouver $g_1$ comme dans (1), nous recollons la section
$\sigma_0$ sur $\Omega = (\overline{B} \cup T_1) \cap U$
avec la restriction de la section $g$ à $\Omega' = \overline{A} \cap U$.
Notons que $U = \Omega \cup \Omega'$ (schématiquement),
car $U$ est réduit et $\Omega \cap \Omega' = Z_1$ (schématiquement)
par notre choix de $Z_1$. Ainsi, d'après
Morphismes, lemme \ref{morphisms-lemma-scheme-theoretic-union},
$U$ est la somme amalgamée de $\Omega$ et $\Omega'$
au-dessus de $Z_1$. Nous pouvons donc trouver $g_1$.
L'existence de $g_2$ dans (2) se démontre de même.
Nous avons alors
$$
\tau = g_2|_{A \cup B} - g_1|_{A \cup B} \quad(\text{addition selon la loi de groupe})
$$
et nous voyons que $X$ est trivial, ce qui achève la preuve.
\end{proof}




















\section{Analogue affine du changement de base propre}
\label{section-gabber-affine-proper}

\noindent
Dans cette section, nous étudions un résultat d'Ofer Gabber, voir
\cite{gabber-affine-proper}. Il a aussi été démontré par Roland Huber, voir
\cite{Huber-henselian}. Nous avons déjà effectué une partie du travail nécessaire
à la démonstration de Gabber dans la section \ref{section-strictly-henselian-local-rings}.

\begin{lemma}
\label{lemma-efface-cohomology-on-closed-by-finite-cover}
Soit $X$ un schéma affine. Soit $\mathcal{F}$ un faisceau abélien de torsion
sur $X_\etale$. Soit $Z \subset X$ un sous-schéma fermé. Soit
$\xi \in H^q_\etale(Z, \mathcal{F}|_Z)$ pour un $q > 0$.
Il existe alors un morphisme injectif $\mathcal{F} \to \mathcal{F}'$
de faisceaux abéliens de torsion sur $X_\etale$ tel que
l'image de $\xi$ dans $H^q_\etale(Z, \mathcal{F}'|_Z)$ soit nulle.
\end{lemma}

\begin{proof}
D'après les lemmes \ref{lemma-torsion-colimit-constructible} et \ref{lemma-colimit},
on peut trouver un morphisme $\mathcal{G} \to \mathcal{F}$ où $\mathcal{G}$
est un faisceau abélien constructible et où $\xi$ provient d'un élément $\zeta$ de
$H^q_\etale(Z, \mathcal{G}|_Z)$. Supposons qu'on puisse trouver un morphisme injectif
$\mathcal{G} \to \mathcal{G}'$ de faisceaux abéliens de torsion sur $X_\etale$
tel que l'image de $\zeta$ dans $H^q_\etale(Z, \mathcal{G}'|_Z)$ soit nulle.
On peut alors prendre pour $\mathcal{F}'$ la somme amalgamée
$$
\mathcal{F}' = \mathcal{G}' \amalg_{\mathcal{G}} \mathcal{F}
$$
et l'on obtient la conclusion du lemme. (Observons que la restriction
à $Z$ est exacte, donc commute aux limites et colimites finies, et qu'en outre
elle commute aux colimites quelconques comme adjoint à gauche de l'image directe.)
On peut donc supposer $\mathcal{F}$ constructible.

\medskip\noindent
Supposons $\mathcal{F}$ constructible. D'après le
lemme \ref{lemma-constructible-maps-into-constant-general},
il suffit de démontrer le résultat lorsque $\mathcal{F}$
est de la forme $f_*\underline{M}$, où $M$ est un groupe abélien fini
et $f : Y \to X$ un morphisme fini de présentation finie
(de tels faisceaux sont encore constructibles d'après le
lemme \ref{lemma-finite-pushforward-constructible},
mais nous n'en aurons pas besoin).
Comme la formation de $f_*$ commute à tout changement de base
(lemme \ref{lemma-finite-pushforward-commutes-with-base-change}),
la restriction de $f_*\underline{M}$ à $Z$ est
égale à l'image directe de $\underline{M}$ par
$Y \times_X Z \to Z$. Par la suite spectrale de Leray
(proposition \ref{proposition-leray})
et l'annulation des images directes supérieures
(proposition \ref{proposition-finite-higher-direct-image-zero}),
on obtient
$$
H^q_\etale(Z, f_*\underline{M}|_Z) = H^q_\etale(Y \times_X Z, \underline{M}).
$$
D'après le lemme \ref{lemma-integral-cover-trivial-cohomology},
on peut trouver un morphisme fini surjectif $Y' \to Y$ de présentation finie
tel que l'image de $\xi$ dans $H^q(Y' \times_X Z, \underline{M})$ soit nulle.
Notons $f' : Y' \to X$ le composé $Y' \to Y \to X$. Nous affirmons que
le morphisme
$$
f_*\underline{M} \longrightarrow f'_*\underline{M}
$$
est injectif, ce qui achève la démonstration d'après ce qui précède.
Pour établir cette injectivité, il suffit de regarder les fibres. En effet,
si $\overline{x} : \Spec(k) \to X$ est un point géométrique, alors
$$
(f_*\underline{M})_{\overline{x}} =
\bigoplus\nolimits_{f(\overline{y}) = \overline{x}} M
$$
d'après la proposition \ref{proposition-finite-higher-direct-image-zero},
et de même pour l'autre faisceau.
Comme $Y' \to Y$ est surjectif et fini, l'application induite
entre les points géométriques relevant $\overline{x}$ est
elle aussi surjective, d'où la conclusion.
\end{proof}

\noindent
Le lemme précédent traitera les groupes de cohomologie supérieure dans
le résultat de Gabber. Le lemme suivant servira à traiter les
sections globales.

\begin{lemma}
\label{lemma-gabber-h0}
Soit $X$ un schéma quasi-compact et quasi-séparé.
Soit $i : Z \to X$ une immersion fermée. Supposons que
\begin{enumerate}
\item pour tout faisceau $\mathcal{F}$ sur $X_{Zar}$, l'application
$\Gamma(X, \mathcal{F}) \to \Gamma(Z, i^{-1}\mathcal{F})$
soit bijective, et
\item pour tout morphisme fini $X' \to X$, l'hypothèse (1) soit satisfaite
pour $Z \times_X X' \to X'$.
\end{enumerate}
Alors, pour tout faisceau $\mathcal{F}$ sur $X_\etale$, on a
$\Gamma(X, \mathcal{F}) = \Gamma(Z, i^{-1}_{small}\mathcal{F})$.
\end{lemma}

\begin{proof}
Soit $\mathcal{F}$ un faisceau sur $X_\etale$. Il existe un morphisme canonique
(de changement de base)
$$
i^{-1}(\mathcal{F}|_{X_{Zar}})
\longrightarrow
(i_{small}^{-1}\mathcal{F})|_{Z_{Zar}}
$$
de faisceaux sur $Z_{Zar}$. Nous montrerons son injectivité en regardant
les fibres. La fibre du membre de gauche en $z \in Z$
est celle de $\mathcal{F}|_{X_{Zar}}$ en $z$. La fibre du membre de droite
est la limite inductive portant sur tous les voisinages étales élémentaires
$(U, u) \to (X, z)$ tels que $U \times_X Z \to Z$ admette une section sur
un voisinage de $z$. Comme les morphismes étales sont ouverts, l'image de
$U \to X$ est un voisinage ouvert $U_0$ de $z$ dans $X$. L'application
$\mathcal{F}(U_0) \to \mathcal{F}(U)$ est injective par la condition de faisceau
pour $\mathcal{F}$ relativement au recouvrement étale $U \to U_0$.
En passant à la limite inductive sur tous les $U$ et $U_0$, on obtient l'injectivité sur les fibres.

\medskip\noindent
Il en résulte, avec l'hypothèse (1), que l'application
$\Gamma(X, \mathcal{F}) \to \Gamma(Z, i^{-1}_{small}\mathcal{F})$
est injective. Par (2), il en va de même sur tout $X'$ fini sur $X$.

\medskip\noindent
Avant de démontrer la surjectivité, introduisons une notation. Pour
tout objet $U$ de $X_\etale$ et tout $s \in \mathcal{F}(U)$, nous
notons $i^{-1}s \in (i_{small}^{-1}\mathcal{F})(U \times_X Z)$
l'image inverse (plus précisément, c'est l'image de $s$ par le morphisme
de sites du lemme \ref{sites-lemma-recover}, composé avec le morphisme du
préfaisceau image inverse vers le faisceau image inverse). Cette construction est
compatible avec les images inverses par les morphismes de $X_\etale$ et par les morphismes (finis)
$X' \to X$, d'une manière dont nous laissons au lecteur le soin de préciser les détails.

\medskip\noindent
Soit $t \in \Gamma(Z, i^{-1}_{small}\mathcal{F})$. Par construction de
$i^{-1}_{small}\mathcal{F}$, il existe un recouvrement étale
$\{V_j \to Z\}_{j \in J}$, des morphismes étales $U_j \to X$, des sections
$s_j \in \mathcal{F}(U_j)$ et des morphismes $V_j \to U_j$ au-dessus de $X$
tels que $t|_{V_j}$ soit l'image inverse de $i^{-1}s_j$ par
$V_j \to U_j \times_X Z$. Comme $V_j \to U_j \times_X Z$ est étale,
son image est un sous-schéma ouvert. Comme $U_j \times_X Z \subset U_j$ est
fermé, on peut, après avoir remplacé $U_j$ par un sous-schéma ouvert, supposer
que $V_j \to U_j \times_X Z$ est surjectif. Alors $\{V_j \to U_j \times_X Z\}$
est un recouvrement étale, d'où $t|_{U_j \times_X Z} = i^{-1}s_j$.
Observons que tout sous-schéma fermé non vide $T \subset X$ rencontre $Z$
par l'hypothèse (1), appliquée par exemple au faisceau $(T \to X)_*\underline{\mathbf{Z}}$.
On voit donc que $\coprod_{j \in J} U_j \to X$ est surjectif.
D'après le lemme
\ref{more-morphisms-lemma-there-is-a-scheme-integral-over} de Compléments sur les morphismes,
on peut trouver un morphisme fini surjectif $X' \to X$
tel que le morphisme $X' \to X$ se factorise, localement pour la topologie de Zariski, par $\coprod_{j \in J} U_j \to X$.
Écrivons $X' = \bigcup_{\alpha \in A} W_\alpha$, où cette réunion est un recouvrement ouvert,
et soit $j : A \to J$ une application telle qu'il existe des morphismes
$g_\alpha : W_\alpha \to U_{j(\alpha)}$ au-dessus de $X$. Notons
$\mathcal{F}'$ l'image inverse de $\mathcal{F}$ sur $X'_\etale$.
Notons $s'_\alpha \in \mathcal{F}'(W_\alpha)$ l'image inverse
de $s_{j(\alpha)}$ par $g_\alpha$.
Posons $Z' = X' \times_X Z$. Notons $i' : Z' \to X'$
le changement de base de $i$. Alors $(i'_{small})^{-1}\mathcal{F}' =
(Z' \to Z)_{small}^{-1}i_{small}^{-1}\mathcal{F}$. Par cette identification,
on peut noter $t' \in
\Gamma(Z', (i')^{-1}_{small}\mathcal{F}')$
l'image inverse de $t$ sur $Z'$.
Par construction, $t'|_{W_\alpha \times_X Z}$ est égal à $(i')^{-1}s'_\alpha$.
On en déduit que $t'$ est une section du
sous-faisceau
$$
(i')^{-1}(\mathcal{F}'|_{X'_{Zar}})
\subset
(i'_{small})^{-1}\mathcal{F}')|_{Z'_{Zar}}
$$
Par l'hypothèse (2), la section $t'$ provient d'une section
$s' \in \Gamma(X', \mathcal{F}'|_{X'_{Zar}}) = \Gamma(X', \mathcal{F}')$.
Par l'injectivité démontrée au deuxième alinéa,
on en déduit que les deux images inverses de $s'$ sur
$X' \times_X X'$ sont égales (car cela est vrai pour
les images inverses de $t'$ sur $Z' \times_Z Z'$). On en conclut que
$s'$ provient d'une section de $\mathcal{F}$ sur $X$ grâce à la
remarque \ref{remark-cohomological-descent-finite}.
\end{proof}

\begin{lemma}
\label{lemma-connected-topological}
Soit $Z \subset X$ une partie fermée d'un espace topologique $X$.
Supposons que
\begin{enumerate}
\item $X$ soit un espace spectral
(Topologie, définition \ref{topology-definition-spectral-space}), et
\item pour $x \in X$, l'intersection $Z \cap \overline{\{x\}}$
soit connexe (en particulier non vide).
\end{enumerate}
Si $Z = Z_1 \amalg Z_2$, où les $Z_i$ sont fermés dans $Z$,
il existe alors une décomposition $X = X_1 \amalg X_2$ où les
$X_i$ sont fermés dans $X$ et $Z_i = Z \cap X_i$.
\end{lemma}

\begin{proof}
Observons que $Z_i$ est quasi-compact. L'ensemble $W_i$ des points
qui se spécialisent en un point de $Z_i$ est donc fermé dans la topologie constructible
d'après le lemme \ref{topology-lemma-make-spectral-space} de Topologie.
L'hypothèse (2) implique que $X = W_1 \amalg W_2$.
Soit $x \in \overline{W_1}$. D'après la partie (1) du lemme
\ref{topology-lemma-constructible-stable-specialization-closed} de Topologie,
il existe une spécialisation $x_1 \leadsto x$ avec $x_1 \in W_1$.
Ainsi $\overline{\{x\}} \subset \overline{\{x_1\}}$, et l'on voit
que $x \in W_1$. Autrement dit, le choix $X_i = W_i$ convient.
\end{proof}

\begin{lemma}
\label{lemma-h0-topological}
Soit $Z \subset X$ une partie fermée d'un espace topologique $X$.
Supposons que
\begin{enumerate}
\item $X$ soit un espace spectral
(Topologie, définition \ref{topology-definition-spectral-space}), et
\item pour $x \in X$, l'intersection $Z \cap \overline{\{x\}}$
soit connexe (en particulier non vide).
\end{enumerate}
Alors, pour tout faisceau $\mathcal{F}$ sur $X$, on a
$\Gamma(X, \mathcal{F}) = \Gamma(Z, \mathcal{F}|_Z)$.
\end{lemma}

\begin{proof}
Si $x \leadsto x'$ est une spécialisation de points, il existe un
morphisme canonique $\mathcal{F}_{x'} \to \mathcal{F}_x$, compatible avec les
sections sur les ouverts et fonctoriel en $\mathcal{F}$. Comme tout point
de $X$ se spécialise en un point de $Z$, il s'ensuit que
$\Gamma(X, \mathcal{F}) \to \Gamma(Z, \mathcal{F}|_Z)$ est injective.
La difficulté est d'en démontrer la surjectivité.

\medskip\noindent
Notons $\mathcal{B}$ l'ensemble de tous les ouverts quasi-compacts de $X$.
Écrivons $\mathcal{F}$ comme une limite inductive filtrante $\mathcal{F} = \colim \mathcal{F}_i$,
où chaque $\mathcal{F}_i$ est comme dans
Modules, équation (\ref{modules-equation-towards-constructible-sets}).
Voir Modules, lemme \ref{modules-lemma-filtered-colimit-constructibles}.
Alors $\mathcal{F}|_Z = \colim \mathcal{F}_i|_Z$, car la restriction à $Z$
est un adjoint à gauche (Catégories, lemme \ref{categories-lemma-adjoint-exact}, et
Faisceaux, lemme \ref{sheaves-lemma-f-map}).
D'après le lemme \ref{sheaves-lemma-directed-colimits-sections} de Faisceaux,
les foncteurs $\Gamma(X, -)$ et $\Gamma(Z, -)$ commutent aux limites inductives filtrantes.
On peut donc supposer que notre faisceau $\mathcal{F}$ est comme dans
Modules, équation (\ref{modules-equation-towards-constructible-sets}).

\medskip\noindent
Supposons disposer d'une injection $\mathcal{F} \subset \mathcal{G}$.
On a alors
$$
\Gamma(X, \mathcal{F}) =
\Gamma(Z, \mathcal{F}|_Z) \cap \Gamma(X, \mathcal{G})
$$
l'intersection étant prise dans $\Gamma(Z, \mathcal{G}|_Z)$.
Cela résulte de la première remarque de la démonstration, car on peut vérifier
si une section globale de $\mathcal{G}$ appartient à $\mathcal{F}$ en regardant
les fibres, et tout point de $X$ se spécialise en un point de $Z$.

\medskip\noindent
D'après le lemme \ref{modules-lemma-constructible-in-constant} de Modules,
il existe une injection $\mathcal{F} \to \prod (Z_i \to X)_*\underline{S_i}$,
où le produit est fini, $Z_i \subset X$ est fermé et $S_i$ est fini.
Il suffit donc de démontrer la surjectivité pour les faisceaux
$(Z_i \to X)_*\underline{S_i}$. Observons que
$$
\Gamma(X, (Z_i \to X)_*\underline{S_i}) = \Gamma(Z_i, \underline{S_i})
\quad\text{et}\quad
\Gamma(Z, (Z_i \to X)_*\underline{S_i}|_Z) =
\Gamma(Z \cap Z_i, \underline{S_i})
$$
De plus, les conditions (1) et (2) sont héritées par $Z_i$; c'est clair
pour (2), et cela résulte du
lemme \ref{topology-lemma-spectral-sub} de Topologie pour (1). Il
suffit donc de démontrer le lemme dans le cas d'un faisceau constant de valeur finie.
Ce cas est une reformulation du lemme \ref{lemma-connected-topological},
ce qui achève la démonstration.
\end{proof}

\begin{example}
\label{example-quinard}
Le lemme \ref{lemma-h0-topological} est faux si $X$ n'est pas spectral.
Voici un exemple : soit $Y$ un espace topologique $T_1$, et soit
$y \in Y$ un point non ouvert. Soit $X = Y \amalg \{ x \}$, muni de
la topologie dont les fermés sont $\emptyset$, $\{y\}$ et tous les
$F \amalg \{ x \}$, où $F$ est une partie fermée de $Y$. Alors
$Z = \{x, y\}$ est une partie fermée de $X$ qui satisfait à l'hypothèse (2)
du lemme \ref{lemma-h0-topological}. Mais $X$ est connexe, tandis que $Z$ ne l'est pas.
La conclusion du lemme échoue donc pour le faisceau constant
de valeur $\{0, 1\}$ sur $X$.
\end{example}

\begin{lemma}
\label{lemma-h0-henselian-pair}
Soit $(A, I)$ un couple hensélien. Posons $X = \Spec(A)$ et
$Z = \Spec(A/I)$. Pour tout faisceau $\mathcal{F}$ sur $X_\etale$,
on a $\Gamma(X, \mathcal{F}) = \Gamma(Z, \mathcal{F}|_Z)$.
\end{lemma}

\begin{proof}
Rappelons que le spectre de tout anneau est un espace spectral, voir
Algèbre, lemme \ref{algebra-lemma-spec-spectral}. D'après le
lemme
\ref{more-algebra-lemma-irreducible-henselian-pair-connected} de Compléments d'algèbre,
on voit que $\overline{\{x\}} \cap Z$ est connexe pour tout $x \in X$.
Le lemme \ref{lemma-h0-topological} montre que l'énoncé
est vrai pour les faisceaux sur $X_{Zar}$. Pour tout morphisme fini $X' \to X$,
on a $X' = \Spec(A')$ et $Z \times_X X' = \Spec(A'/IA')$,
où $(A', IA')$ est un couple hensélien, voir le lemme
\ref{more-algebra-lemma-integral-over-henselian-pair} de Compléments d'algèbre,
et l'on obtient le même énoncé pour les faisceaux sur $(X')_{Zar}$.
On peut donc conclure en appliquant le lemme \ref{lemma-gabber-h0}.
\end{proof}

\noindent
Enfin, nous pouvons énoncer et démontrer le théorème de Gabber.

\begin{theorem}[Gabber]
\label{theorem-gabber}
Soit $(A, I)$ un couple hensélien. Posons $X = \Spec(A)$ et
$Z = \Spec(A/I)$. Pour tout faisceau abélien de torsion $\mathcal{F}$ sur $X_\etale$,
on a $H^q_\etale(X, \mathcal{F}) = H^q_\etale(Z, \mathcal{F}|_Z)$.
\end{theorem}

\begin{proof}
Le résultat vaut pour $q = 0$ par le lemme \ref{lemma-h0-henselian-pair}.
Soit $q \geq 1$. Supposons le résultat démontré en tout degré $< q$.
Soit $\mathcal{F}$ un faisceau abélien de torsion. Soit
$\mathcal{F} \to \mathcal{F}'$
un morphisme injectif de faisceaux abéliens de torsion (qui sera choisi plus tard),
de conoyau $\mathcal{Q}$, de sorte que l'on ait la suite exacte courte
$$
0 \to \mathcal{F} \to \mathcal{F}' \to \mathcal{Q} \to 0
$$
de faisceaux abéliens de torsion sur $X_\etale$. Cela donne un morphisme de suites exactes longues
de cohomologie sur $X$ et $Z$, dont une partie est de la forme
$$
\xymatrix{
H^{q - 1}_\etale(X, \mathcal{F}') \ar[d] \ar[r] &
H^{q - 1}_\etale(X, \mathcal{Q}) \ar[d] \ar[r] &
H^q_\etale(X, \mathcal{F}) \ar[d] \ar[r] &
H^q_\etale(X, \mathcal{F}') \ar[d] \\
H^{q - 1}_\etale(Z, \mathcal{F}'|_Z) \ar[r] &
H^{q - 1}_\etale(Z, \mathcal{Q}|_Z) \ar[r] &
H^q_\etale(Z, \mathcal{F}|_Z) \ar[r] &
H^q_\etale(Z, \mathcal{F}'|_Z)
}
$$
Nous allons terminer la démonstration à l'aide de ce diagramme commutatif de groupes abéliens
à lignes exactes.

\medskip\noindent
Injectivité pour $\mathcal{F}$. Soit $\xi$ un élément non nul de
$H^q_\etale(X, \mathcal{F})$.
D'après le lemme \ref{lemma-efface-cohomology-on-closed-by-finite-cover}, appliqué avec
$Z = X$ (!), on peut trouver $\mathcal{F} \subset \mathcal{F}'$ de sorte que
l'image de $\xi$ dans le groupe de droite soit nulle. Alors $\xi$ est l'image d'un
élément de $H^{q - 1}_\etale(X, \mathcal{Q})$, et la bijectivité
en degré $q - 1$ implique que l'image de $\xi$ n'est pas nulle dans
$H^q_\etale(Z, \mathcal{F}|_Z)$.

\medskip\noindent
Surjectivité pour $\mathcal{F}$. Soit $\xi$ un élément de
$H^q_\etale(Z, \mathcal{F}|_Z)$.
D'après le lemme \ref{lemma-efface-cohomology-on-closed-by-finite-cover}, appliqué avec
$Z = Z$, on peut trouver $\mathcal{F} \subset \mathcal{F}'$ de sorte que
l'image de $\xi$ dans le groupe de droite soit nulle. Alors $\xi$ est l'image d'un
élément de $H^{q - 1}_\etale(Z, \mathcal{Q}|_Z)$, et la bijectivité
en degré $q - 1$ implique que $\xi$ appartient à l'image du morphisme vertical.
\end{proof}

\begin{lemma}
\label{lemma-vanishing-restriction-injective}
Soit $X$ un schéma à diagonale affine pouvant être recouvert par
$n + 1$ ouverts affines. Soit $Z \subset X$ un sous-schéma fermé.
Soit $\mathcal{A}$ un faisceau d'anneaux de torsion sur $X_\etale$
et soit $\mathcal{I}$ un faisceau injectif de $\mathcal{A}$-modules
sur $X_\etale$.
Alors $H^q_\etale(Z, \mathcal{I}|_Z) = 0$ pour $q > n$.
\end{lemma}

\begin{proof}
Nous procédons par récurrence sur $n$. Si $n = 0$, alors $X$ est affine.
Écrivons $X = \Spec(A)$ et $Z = \Spec(A/I)$. Soit $A^h$ la limite inductive filtrante
des $A$-algèbres étales $B$ telles que $A/I \to B/IB$ soit un isomorphisme.
Alors $(A^h, IA^h)$ est un couple hensélien et $A/I = A^h/IA^h$, voir
le lemme \ref{more-algebra-lemma-henselization} de Compléments d'algèbre
et sa démonstration. Posons $X^h = \Spec(A^h)$.
D'après le théorème \ref{theorem-gabber},
on a
$$
H^q_\etale(Z, \mathcal{I}|_Z) = H^q_\etale(X^h, \mathcal{I}|_{X^h})
$$
D'après le théorème \ref{theorem-colimit}, on a
$$
H^q_\etale(X^h, \mathcal{I}|_{X^h}) =
\colim_{A \to B} H^q_\etale(\Spec(B), \mathcal{I}|_{\Spec(B)})
$$
où la limite inductive porte sur les $A$-algèbres $B$ ci-dessus.
Comme les morphismes $\Spec(B) \to \Spec(A)$ sont étales,
la restriction $\mathcal{I}|_{\Spec(B)}$ est un faisceau injectif
de $\mathcal{A}|_{\Spec(B)}$-modules
(Cohomologie sur les sites, lemme \ref{sites-cohomology-lemma-cohomology-of-open}).
Les groupes de cohomologie du membre de droite sont donc nuls, ce qui donne le
résultat dans ce cas.

\medskip\noindent
Étape de récurrence. On peut utiliser Mayer--Vietoris.
Supposons en effet que $X = U \cup V$, où $U$ est une réunion de $n$ ouverts
affines et où $V$ est affine. Puisque la diagonale de $X$ est affine,
on voit que $U \cap V$ est la réunion de $n$ ouverts affines. Mayer--Vietoris
fournit une suite exacte
$$
H^{q - 1}_\etale(U \cap V \cap Z, \mathcal{I}|_Z) \to
H^q_\etale(Z, \mathcal{I}|_Z) \to
H^q_\etale(U \cap Z, \mathcal{I}|_Z) \oplus
H^q_\etale(V \cap Z, \mathcal{I}|_Z)
$$
et l'hypothèse de récurrence donne l'annulation pour $q > n$, comme voulu.
\end{proof}





\section{Cohomologie des faisceaux de torsion sur les courbes}
\label{section-vanishing-torsion}

\noindent
Le but de cette section est d'établir les résultats fondamentaux de finitude et d'annulation
pour la cohomologie des faisceaux de torsion sur les courbes, voir le
théorème \ref{theorem-vanishing-affine-curves}.
Dans la section \ref{section-vanishing-torsion-coefficients},
nous étudierons les faisceaux constructibles de modules de torsion
sur un anneau noethérien.

\begin{situation}
\label{situation-what-to-prove}
Ici, $k$ est un corps algébriquement clos, $X$ est un schéma séparé de type fini
de dimension $\leq 1$ sur $k$, et $\mathcal{F}$ est un faisceau abélien
de torsion sur $X_\etale$.
\end{situation}

\noindent
Dans la situation \ref{situation-what-to-prove},
nous voulons démontrer les assertions suivantes :
\begin{enumerate}
\item
\label{item-vanishing}
$H^q_\etale(X, \mathcal{F}) = 0$ pour $q > 2$,
\item
\label{item-vanishing-affine}
$H^q_\etale(X, \mathcal{F}) = 0$ pour $q > 1$ si $X$ est affine,
\item
\label{item-vanishing-p-p}
$H^q_\etale(X, \mathcal{F}) = 0$ pour $q > 1$ si $p = \text{char}(k) > 0$
et si $\mathcal{F}$ est de torsion $p$-primaire,
\item
\label{item-finite-prime-to-p}
$H^q_\etale(X, \mathcal{F})$ est fini si $\mathcal{F}$ est constructible
et de torsion première à $\text{char}(k)$,
\item
\label{item-finite-proper}
$H^q_\etale(X, \mathcal{F})$ est fini si $X$ est propre et si
$\mathcal{F}$ est constructible,
\item
\label{item-base-change-prime-to-p}
$H^q_\etale(X, \mathcal{F}) \to
H^q_\etale(X_{k'}, \mathcal{F}|_{X_{k'}})$ est un isomorphisme
pour toute extension $k'/k$ de corps algébriquement clos
si $\mathcal{F}$ est de torsion première à $\text{char}(k)$,
\item
\label{item-base-change-proper}
$H^q_\etale(X, \mathcal{F}) \to
H^q_\etale(X_{k'}, \mathcal{F}|_{X_{k'}})$ est un isomorphisme
pour toute extension $k'/k$ de corps algébriquement clos
si $X$ est propre,
\item
\label{item-surjective}
$H^2_\etale(X, \mathcal{F}) \to H^2_\etale(U, \mathcal{F})$
est surjectif pour tout ouvert $U \subset X$.
\end{enumerate}
Étant donnée une situation \ref{situation-what-to-prove},
nous dirons que
« les assertions (\ref{item-vanishing}) -- (\ref{item-surjective}) sont vérifiées »
si celles qui s'appliquent à la situation considérée sont vraies.
Nous commençons la démonstration par la conséquence suivante de notre calcul
de la cohomologie à coefficients constants.

\begin{lemma}
\label{lemma-constant-smooth-statements}
Dans la situation \ref{situation-what-to-prove},
supposons que $X$ soit lisse et que $\mathcal{F} = \underline{\mathbf{Z}/\ell\mathbf{Z}}$
pour un nombre premier $\ell$. Alors les assertions
(\ref{item-vanishing}) -- (\ref{item-surjective}) sont vérifiées
pour $\mathcal{F}$.
\end{lemma}

\begin{proof}
Puisque $X$ est lisse, nous voyons que $X$ est une réunion disjointe finie de
courbes lisses et de points. La cohomologie étale supérieure de tout faisceau
sur un point (le spectre de $k$) est nulle, par exemple d'après le
lemme \ref{lemma-compare-cohomology-point}
ou le lemme \ref{lemma-affine-only-closed-points}.
Nous pouvons donc supposer que $X$ est une courbe lisse.

\medskip\noindent
Cas I : $\ell$ est distinct de la caractéristique de $k$.
Ce cas résulte du
lemme \ref{lemma-cohomology-smooth-projective-curve}
(cas projectif) et du
lemme \ref{lemma-vanishing-cohomology-mu-smooth-curve}
(cas affine). L'assertion (\ref{item-base-change-prime-to-p})
sur la cohomologie et l'extension du corps de base algébriquement clos
résulte de ce que le genre $g$ et le nombre $r$
de « points manquants » ne changent pas quand on passe de $k$ à $k'$.
L'assertion (\ref{item-surjective}) résulte de ce que $H^2_\etale(U, \mathcal{F})$
est nul dès que $U \not = X$, car $U$ est alors affine
(Variétés, lemmes \ref{varieties-lemma-proper-minus-point} et
\ref{varieties-lemma-curve-affine-projective}).

\medskip\noindent
Cas II : $\ell$ est égal à la caractéristique de $k$.
L'annulation résulte du lemme \ref{lemma-vanishing-variety-char-p-p}.
Les assertions (\ref{item-finite-proper}) et (\ref{item-base-change-proper})
résultent du
lemme \ref{lemma-finiteness-proper-variety-char-p-p}.
\end{proof}

\begin{remark}[Invariance par extension du corps de base algébriquement clos]
\label{remark-invariance}
Soit $k$ un corps algébriquement clos de caractéristique $p > 0$.
Dans la section \ref{section-artin-schreier}, nous avons vu qu'il existe
une suite exacte
$$
k[x] \to k[x] \to
H^1_\etale(\mathbf{A}^1_k, \mathbf{Z}/p\mathbf{Z}) \to 0
$$
où la première flèche envoie $f(x)$ sur $f^p - f$. Un système de représentants
du conoyau est formé des polynômes
$$
\sum\nolimits_{p \not | n} \lambda_n x^n
$$
avec $\lambda_n \in k$. (Si $k$ n'est pas algébriquement clos,
il faut aussi y ajouter certaines constantes.) En particulier,
lorsque $k'/k$ est une extension de corps algébriquement clos,
l'application
$$
H^1_\etale(\mathbf{A}^1_k, \mathbf{Z}/p\mathbf{Z})
\to
H^1_\etale(\mathbf{A}^1_{k'}, \mathbf{Z}/p\mathbf{Z})
$$
n'est pas un isomorphisme en général. En particulier, l'application
$\pi_1(\mathbf{A}^1_{k'}) \to \pi_1(\mathbf{A}^1_k)$
entre groupes fondamentaux étales (référence à ajouter ultérieurement)
n'est pas non plus un isomorphisme. Ainsi, le type d'homotopie étale
de la droite affine dépend du corps de base algébriquement clos.
Le lemme \ref{lemma-constant-smooth-statements} ci-dessus montre
que ce phénomène ne se produit qu'en caractéristique $p$
avec des coefficients de torsion $p$-primaire.
\end{remark}

\begin{lemma}
\label{lemma-ses-statements}
Soit $k$ un corps algébriquement clos. Soit $X$ un schéma séparé de type
fini sur $k$ et de dimension $\leq 1$. Soit
$0 \to \mathcal{F}_1 \to \mathcal{F} \to \mathcal{F}_2 \to 0$
une suite exacte courte de faisceaux abéliens de torsion sur $X$.
Si les assertions (\ref{item-vanishing}) -- (\ref{item-surjective}) sont vérifiées
pour $\mathcal{F}_1$ et $\mathcal{F}_2$, elles le sont aussi
pour $\mathcal{F}$.
\end{lemma}

\begin{proof}
Cela résulte pour l'essentiel immédiatement des définitions et de la suite exacte longue
de cohomologie. Observons aussi que $\mathcal{F}$ est constructible
(resp.\ de torsion première à la caractéristique de $k$) si et seulement si
$\mathcal{F}_1$ et $\mathcal{F}_2$ sont tous deux constructibles
(resp.\ de torsion première à la caractéristique de $k$). Voir la
proposition \ref{proposition-constructible-over-noetherian}.
Nous omettons quelques détails.
\end{proof}

\begin{lemma}
\label{lemma-finite-pushforward-statements}
Soit $k$ un corps algébriquement clos. Soit $f : X \to Y$ un
morphisme fini de schémas séparés de type fini sur $k$ et de
dimension $\leq 1$. Soit $\mathcal{F}$ un faisceau abélien de torsion sur $X$.
Si les assertions (\ref{item-vanishing}) -- (\ref{item-surjective}) sont vérifiées
pour $\mathcal{F}$, elles le sont aussi pour $f_*\mathcal{F}$.
\end{lemma}

\begin{proof}
En effet, nous avons $H^q_\etale(X, \mathcal{F}) = H^q_\etale(Y, f_*\mathcal{F})$
par l'annulation de $R^qf_*$ pour $q > 0$
(proposition \ref{proposition-finite-higher-direct-image-zero}) et par
la suite spectrale de Leray
(Cohomologie sur les sites, lemme \ref{sites-cohomology-lemma-apply-Leray}).
Pour (\ref{item-surjective}), on utilise que la formation de $f_*$ commute à
tout changement de base
(lemme \ref{lemma-finite-pushforward-commutes-with-base-change}).
\end{proof}

\begin{lemma}
\label{lemma-restrict-to-open}
Dans la situation \ref{situation-what-to-prove}, supposons $\mathcal{F}$
constructible. Soit $j : X' \to X$ l'inclusion d'un sous-schéma ouvert dense.
Alors les assertions
(\ref{item-vanishing}) -- (\ref{item-surjective}) sont vérifiées pour $\mathcal{F}$
si et seulement si elles le sont pour $j_!j^{-1}\mathcal{F}$.
\end{lemma}

\begin{proof}
Puisque $X'$ est dense, nous voyons que $Z = X \setminus X'$ est de dimension $0$
et est donc un ensemble fini $Z = \{x_1, \ldots, x_n\}$ de points $k$-rationnels.
Considérons la suite exacte courte
$$
0 \to j_!j^{-1}\mathcal{F} \to \mathcal{F} \to i_*i^{-1}\mathcal{F} \to 0
$$
du lemme \ref{lemma-ses-associated-to-open}. Observons que
$H^q_\etale(X, i_*i^{-1}\mathcal{F}) = H^q_\etale(Z, i^{-1}\mathcal{F})$.
En effet, $i : Z \to X$ est une immersion fermée, donc finie,
et nous avons donc l'annulation de $R^qi_*$ pour $q > 0$ d'après la
proposition \ref{proposition-finite-higher-direct-image-zero};
l'égalité résulte alors de la suite spectrale de Leray
(Cohomologie sur les sites, lemme \ref{sites-cohomology-lemma-apply-Leray}).
Puisque $Z$ est une réunion disjointe de spectres de corps algébriquement
clos, nous en déduisons que $H^q_\etale(Z, i^*\mathcal{F}) = 0$
pour $q > 0$ et que
$$
H^0_\etale(Z, i^{-1}\mathcal{F}) =
\bigoplus\nolimits_{i=1}^n \mathcal{F}_{x_i}
$$
est fini, car $\mathcal{F}_{x_i}$ est fini en raison de
l'hypothèse de constructibilité de $\mathcal{F}$.
La suite exacte longue de cohomologie donne une suite exacte
$$
0 \to
H^0_\etale(X, j_!j^{-1}\mathcal{F}) \to
H^0_\etale(X, \mathcal{F}) \to
H^0_\etale(Z, i^{-1}\mathcal{F}) \to
H^1_\etale(X, j_!j^{-1}\mathcal{F}) \to
H^1_\etale(X, \mathcal{F}) \to 0
$$
et des isomorphismes $H^q_\etale(X, j_!j^{-1}\mathcal{F}) \to
H^q_\etale(X, \mathcal{F})$ pour $q > 1$.

\medskip\noindent
Il est maintenant facile de déduire que chacune des assertions 
(\ref{item-vanishing}) -- (\ref{item-surjective})
est vérifiée pour $\mathcal{F}$ si et seulement si elle l'est pour
$j_!j^{-1}\mathcal{F}$. Faisons quelques remarques pour aider le lecteur :
(a) si $\mathcal{F}$ est de torsion première à la caractéristique
de $k$, alors $j_!j^{-1}\mathcal{F}$ l'est aussi,
(b) le faisceau $j_!j^{-1}\mathcal{F}$ est constructible,
(c) nous avons $H^0_\etale(Z, i^{-1}\mathcal{F}) =
H^0_\etale(Z_{k'}, i^{-1}\mathcal{F}|_{Z_{k'}})$, et (d)
si $U \subset X$ est un ouvert, alors $U' = U \cap X'$ est dense dans $U$.
\end{proof}

\begin{lemma}
\label{lemma-even-easier}
Dans la situation \ref{situation-what-to-prove}, supposons $X$ lisse.
Soit $j : U \to X$ une immersion ouverte. Soit $\ell$ un nombre premier.
Posons $\mathcal{F} = j_!\underline{\mathbf{Z}/\ell\mathbf{Z}}$.
Alors les assertions (\ref{item-vanishing}) -- (\ref{item-surjective}) sont vérifiées
pour $\mathcal{F}$.
\end{lemma}

\begin{proof}
Puisque $X$ est lisse, il est une réunion disjointe de courbes lisses;
nous pouvons donc supposer que $X$ est une courbe (c'est-à-dire irréductible).
Alors soit $U = \emptyset$, auquel cas il n'y a rien à démontrer,
soit $U \subset X$ est dense. Dans ce cas, le lemme résulte des
lemmes \ref{lemma-constant-smooth-statements} et
\ref{lemma-restrict-to-open}.
\end{proof}

\begin{lemma}
\label{lemma-somewhat-easier}
Dans la situation \ref{situation-what-to-prove}, supposons $X$ réduit.
Soit $j : U \to X$ une immersion ouverte. Soit $\ell$ un nombre premier
et posons $\mathcal{F} = j_!\underline{\mathbf{Z}/\ell\mathbf{Z}}$.
Alors les assertions (\ref{item-vanishing}) -- (\ref{item-surjective}) sont vérifiées
pour $\mathcal{F}$.
\end{lemma}

\begin{proof}
La différence avec le lemme \ref{lemma-even-easier} est qu'ici nous ne
supposons pas $X$ lisse. Soit $\nu : X^\nu \to X$ le morphisme de
normalisation. Alors $\nu$ est fini
(Variétés, lemme \ref{varieties-lemma-normalization-locally-algebraic})
et $X^\nu$ est lisse
(Variétés, lemme \ref{varieties-lemma-regular-point-on-curve}).
Soit $j^\nu : U^\nu \to X^\nu$ l'image inverse de $U$.
D'après le lemme \ref{lemma-even-easier}, le résultat vaut pour
$j^\nu_!\underline{\mathbf{Z}/\ell\mathbf{Z}}$.
D'après le lemme \ref{lemma-finite-pushforward-statements},
le résultat vaut pour $\nu_*j^\nu_!\underline{\mathbf{Z}/\ell\mathbf{Z}}$.
En général, il n'est pas vrai que
$\nu_*j^\nu_!\underline{\mathbf{Z}/\ell\mathbf{Z}}$ soit égal à
$j_!\underline{\mathbf{Z}/\ell\mathbf{Z}}$
mais nous pouvons contourner cette difficulté comme suit.
Puisque $X$ est réduit, le morphisme $\nu : X^\nu \to X$
est un isomorphisme au-dessus d'un ouvert dense $j' : X' \to X$
(Variétés, lemme \ref{varieties-lemma-normalization-locally-algebraic}).
Au-dessus de cet ouvert, les deux faisceaux coïncident :
$$
(j')^{-1}(\nu_*j^\nu_!\underline{\mathbf{Z}/\ell\mathbf{Z}})
=
(j')^{-1}(j_!\underline{\mathbf{Z}/\ell\mathbf{Z}})
$$
En appliquant deux fois le lemme \ref{lemma-restrict-to-open}
à $j' : X' \to X$ et aux faisceaux ci-dessus,
nous concluons.
\end{proof}

\begin{lemma}
\label{lemma-vanishing-easier}
Dans la situation \ref{situation-what-to-prove}, supposons $X$ réduit.
Soit $j : U \to X$ une immersion ouverte avec $U$ connexe. Soit
$\ell$ un nombre premier. Soit $\mathcal{G}$ un faisceau localement constant
fini de $\mathbf{F}_\ell$-espaces vectoriels sur $U$. Posons
$\mathcal{F} = j_!\mathcal{G}$. Alors les assertions
(\ref{item-vanishing}) -- (\ref{item-surjective}) sont vérifiées pour $\mathcal{F}$.
\end{lemma}

\begin{proof}
Soit $f : V \to U$ un morphisme fini étale de degré premier à $\ell$
comme dans le lemme \ref{lemma-pullback-filtered}. La discussion de la
section \ref{section-trace-method} donne des applications
$$
\mathcal{G} \to f_*f^{-1}\mathcal{G} \to \mathcal{G}
$$
dont le composé est un isomorphisme. Il suffit donc de démontrer le
lemme avec $\mathcal{F} = j_!f_*f^{-1}\mathcal{G}$.
En vertu du théorème principal de Zariski
(Compléments sur les morphismes, lemme
\ref{more-morphisms-lemma-quasi-finite-separated-pass-through-finite})
nous pouvons choisir un diagramme
$$
\xymatrix{
V \ar[r]_{j'} \ar[d]_f & Y \ar[d]^{\overline{f}} \\
U \ar[r]^j & X
}
$$
où $\overline{f} : Y \to X$ est fini et $j'$ une immersion ouverte
d'image dense. Nous pouvons remplacer $Y$ par son réduit (cela ne
change pas $V$, puisque $V$ est réduit car étale sur $U$).
Comme $f$ est fini et $V$ dense dans $Y$, nous avons $V = U \times_X Y$. D'après le
lemme \ref{lemma-compatible-shriek-push-finite}, nous avons
$$
j_!f_*f^{-1}\mathcal{G} = \overline{f}_*j'_!f^{-1}\mathcal{G}
$$
D'après le lemme \ref{lemma-finite-pushforward-statements}, il suffit de
considérer $j'_!f^{-1}\mathcal{G}$.
L'existence de la filtration donnée par le
lemme \ref{lemma-pullback-filtered},
le fait que $j'_!$ est exact et le
lemme \ref{lemma-ses-statements}
nous ramènent au cas
$\mathcal{F} = j'_!\underline{\mathbf{Z}/\ell\mathbf{Z}}$
qui est le lemme \ref{lemma-somewhat-easier}.
\end{proof}


%10.20.09

\begin{theorem}
\label{theorem-vanishing-affine-curves}
Si $k$ est un corps algébriquement clos, $X$ un schéma séparé de type fini
et de dimension $\leq 1$ sur $k$, et $\mathcal{F}$ un faisceau abélien
de torsion sur $X_\etale$, alors
\begin{enumerate}
\item
$H^q_\etale(X, \mathcal{F}) = 0$ pour $q > 2$,
\item
$H^q_\etale(X, \mathcal{F}) = 0$ pour $q > 1$ si $X$ est affine,
\item
$H^q_\etale(X, \mathcal{F}) = 0$ pour $q > 1$ si $p = \text{char}(k) > 0$
et si $\mathcal{F}$ est de torsion $p$-primaire,
\item
$H^q_\etale(X, \mathcal{F})$ est fini si $\mathcal{F}$ est constructible
et de torsion première à $\text{char}(k)$,
\item
$H^q_\etale(X, \mathcal{F})$ est fini si $X$ est propre et si
$\mathcal{F}$ est constructible,
\item
$H^q_\etale(X, \mathcal{F}) \to
H^q_\etale(X_{k'}, \mathcal{F}|_{X_{k'}})$ est un isomorphisme
pour toute extension $k'/k$ de corps algébriquement clos
si $\mathcal{F}$ est de torsion première à $\text{char}(k)$,
\item
$H^q_\etale(X, \mathcal{F}) \to
H^q_\etale(X_{k'}, \mathcal{F}|_{X_{k'}})$ est un isomorphisme
pour toute extension $k'/k$ de corps algébriquement clos
si $X$ est propre,
\item
$H^2_\etale(X, \mathcal{F}) \to H^2_\etale(U, \mathcal{F})$
est surjectif pour tout ouvert $U \subset X$.
\end{enumerate}
\end{theorem}

\begin{proof}
Le théorème affirme que, dans la situation \ref{situation-what-to-prove},
les assertions (\ref{item-vanishing}) -- (\ref{item-surjective}) sont vérifiées.
Nous commençons par remplacer $X$ par son réduit, ce qui est permis
par la proposition \ref{proposition-topological-invariance}.
D'après le lemme \ref{lemma-torsion-colimit-constructible}, nous pouvons écrire
$\mathcal{F}$ comme limite inductive filtrante de faisceaux abéliens constructibles.
La cohomologie commute aux limites inductives, voir le lemme \ref{lemma-colimit}.
De plus, l'image inverse par $X_{k'} \to X$ commute aux limites inductives en tant
qu'adjoint à gauche. Il suffit donc de démontrer les assertions pour un faisceau constructible.

\medskip\noindent
Dans ce paragraphe, nous utilisons le lemme \ref{lemma-ses-statements} sans le
mentionner davantage. Écrivons
$\mathcal{F} = \mathcal{F}_1 \oplus \ldots \oplus \mathcal{F}_r$
où $\mathcal{F}_i$ est $\ell_i$-primaire pour un nombre premier $\ell_i$; nous pouvons
supposer que $\ell^n$ annule $\mathcal{F}$ pour un nombre premier $\ell$. Considérons la
suite exacte
$$
0 \to \mathcal{F}[\ell] \to \mathcal{F} \to \mathcal{F}/\mathcal{F}[\ell] \to 0.
$$
Nous voyons ainsi qu'il suffit de supposer $\mathcal{F}$ de $\ell$-torsion.
Cela signifie que $\mathcal{F}$ est un faisceau constructible de
$\mathbf{F}_\ell$-espaces vectoriels pour un nombre premier $\ell$.

\medskip\noindent
Par définition, cela signifie qu'il existe un ouvert dense $U \subset X$
tel que $\mathcal{F}|_U$ soit un faisceau localement constant fini de
$\mathbf{F}_\ell$-espaces vectoriels. Comme $\dim(X) \leq 1$, nous pouvons
supposer, après avoir rétréci $U$, que $U = U_1 \amalg \ldots \amalg U_n$
est une réunion disjointe de schémas irréductibles (il suffit de supprimer les points
fermés appartenant aux intersections d'au moins $2$ composantes de $U$).
D'après le lemme \ref{lemma-restrict-to-open}, nous nous ramenons au cas
$\mathcal{F} = j_!\mathcal{G}$ où $\mathcal{G}$ est un faisceau
localement constant fini de $\mathbf{F}_\ell$-espaces vectoriels sur $U$.

\medskip\noindent
Comme nous avons choisi $U = U_1 \amalg \ldots \amalg U_n$ avec les $U_i$ irréductibles,
nous avons
$$
j_!\mathcal{G} =
j_{1!}(\mathcal{G}|_{U_1}) \oplus \ldots \oplus
j_{n!}(\mathcal{G}|_{U_n})
$$
où $j_i : U_i \to X$ est le morphisme d'inclusion.
Le cas de $j_{i!}(\mathcal{G}|_{U_i})$ est traité dans le
lemme \ref{lemma-vanishing-easier}.
\end{proof}

\begin{theorem}
\label{theorem-vanishing-curves}
Soit $X$ un schéma de type fini et de dimension $1$ sur un
corps algébriquement clos $k$. Soit $\mathcal{F}$ un faisceau de torsion
sur $X_\etale$. Alors
$$
H_\etale^q(X, \mathcal{F}) = 0, \quad \forall q \geq 3.
$$
Si $X$ est affine, alors $H_\etale^2(X, \mathcal{F}) = 0$ également.
\end{theorem}

\begin{proof}
Si $X$ est séparé, cela résulte immédiatement du théorème plus précis
\ref{theorem-vanishing-affine-curves}.
Si $X$ n'est pas séparé, choisissons un recouvrement ouvert affine
$X = X_1 \cup \ldots \cup X_n$. Par récurrence sur $n$, nous pouvons supposer
l'annulation vérifiée sur $U = X_1 \cup \ldots \cup X_{n - 1}$.
Alors Mayer--Vietoris (lemme \ref{lemma-mayer-vietoris}) donne
$$
H^2_\etale(U, \mathcal{F}) \oplus H^2_\etale(X_n, \mathcal{F}) \to
H^2_\etale(U \cap X_n, \mathcal{F}) \to
H^3_\etale(X, \mathcal{F}) \to 0
$$
Cependant, puisque $U \cap X_n$ est un ouvert d'un schéma affine,
il est affine en vertu de notre hypothèse de dimension. Le groupe
$H^2_\etale(U \cap X_n, \mathcal{F})$ s'annule donc
d'après le théorème \ref{theorem-vanishing-affine-curves}.
\end{proof}

\begin{lemma}
\label{lemma-base-change-dim-1-separably-closed}
Soit $k'/k$ une extension de corps séparablement clos.
Soit $X$ un schéma propre sur $k$ de dimension $\leq 1$.
Soit $\mathcal{F}$ un faisceau abélien de torsion sur $X$.
Alors l'application $H^q_\etale(X, \mathcal{F}) \to
H^q_\etale(X_{k'}, \mathcal{F}|_{X_{k'}})$ est un isomorphisme
pour $q \geq 0$.
\end{lemma}

\begin{proof}
Nous l'avons vu pour les corps algébriquement clos dans le
théorème \ref{theorem-vanishing-affine-curves}.
Étant donnés $k \subset k'$ comme dans l'énoncé du lemme, nous pouvons
choisir un diagramme
$$
\xymatrix{
k' \ar[r] & \overline{k}' \\
k \ar[u] \ar[r] & \overline{k} \ar[u]
}
$$
où $k \subset \overline{k}$ et $k' \subset \overline{k}'$ sont
les clôtures algébriques. Puisque $k$ et $k'$ sont séparablement clos,
les extensions de corps
$\overline{k}/k$ et $\overline{k}'/k'$
sont algébriques et purement inséparables. Dans ce cas, les morphismes
$X_{\overline{k}} \to X$ et $X_{\overline{k}'} \to X_{k'}$
sont des homéomorphismes universels. Ainsi, la cohomologie de $\mathcal{F}$
peut se calculer sur $X_{\overline{k}}$ et la cohomologie
de $\mathcal{F}|_{X_{k'}}$ peut se calculer sur $X_{\overline{k}'}$,
voir la proposition \ref{proposition-topological-invariance}.
Nous déduisons donc le cas général du cas des corps algébriquement
clos.
\end{proof}






\section{Cohomologie des modules de torsion sur les courbes}
\label{section-vanishing-torsion-coefficients}

\noindent
Dans cette section, nous reprenons les arguments de la
section \ref{section-vanishing-torsion}
pour les faisceaux constructibles de modules de torsion
sur un anneau noethérien. Nous commençons par l'étape la plus intéressante.

\begin{lemma}
\label{lemma-constant-statements-coefficients}
Soit $\Lambda$ un anneau noethérien, soit $M$ un $\Lambda$-module de type fini annulé
par un entier $n > 0$, soit $k$ un corps algébriquement clos et soit $X$ un
schéma séparé de type fini et de dimension $\leq 1$ sur $k$.
Alors
\begin{enumerate}
\item $H^q_\etale(X, \underline{M})$ est un $\Lambda$-module de type fini
si $n$ est premier à $\text{char}(k)$,
\item $H^q_\etale(X, \underline{M})$ est un $\Lambda$-module de type fini
si $X$ est propre.
\end{enumerate}
\end{lemma}

\begin{proof}
Si $n = \ell n'$ pour un nombre premier $\ell$, nous obtenons une suite
exacte courte $0 \to M[\ell] \to M \to M' \to 0$ de $\Lambda$-modules
finis, et $M'$ est annulé par $n'$.
Cela fournit une suite exacte courte correspondante de faisceaux constants,
laquelle donne à son tour une suite
exacte de modules de cohomologie
$$
H^q_\etale(X, \underline{M[\ell]}) \to
H^q_\etale(X, \underline{M}) \to
H^q_\etale(X, \underline{M'})
$$
Ainsi, si nous pouvons démontrer le résultat lorsque $M$ est annulé par
un nombre premier, une récurrence sur $n$ permet de conclure.

\medskip\noindent
Soit $\ell$ un nombre premier qui annule $M$.
Nous pouvons alors remplacer $\Lambda$ par la $\mathbf{F}_\ell$-algèbre
$\Lambda/\ell \Lambda$. En effet,
la cohomologie de $\underline{M}$ considéré comme faisceau de
$\Lambda$-modules est la même que sa cohomologie comme faisceau de
$\Lambda/\ell \Lambda$-modules, par exemple d'après
Cohomologie sur les sites, lemme
\ref{sites-cohomology-lemma-cohomology-modules-abelian-agree}.

\medskip\noindent
Supposons que $\ell$ soit un nombre premier qui annule $M$
et $\Lambda$.
Ramenons-nous au cas où $M$ est un $\Lambda$-module libre de type fini.
Choisissons pour cela une suite exacte courte
$$
0 \to N \to \Lambda^{\oplus m} \to M \to 0
$$
Elle détermine une suite exacte
$$
H^q_\etale(X, \underline{\Lambda^{\oplus m}}) \to
H^q_\etale(X, \underline{M}) \to
H^{q + 1}_\etale(X, \underline{N})
$$
Une récurrence descendante sur $q$ donne le résultat pour $M$
si nous le connaissons pour $\Lambda^{\oplus m}$.
Nous utilisons ici le fait que nos groupes de cohomologie
s'annulent en degrés $> 2$
d'après le théorème \ref{theorem-vanishing-affine-curves}.

\medskip\noindent
Soit $\ell$ un nombre premier et supposons que $\ell$ annule $\Lambda$.
Il reste à montrer que les groupes de cohomologie
$H^q_\etale(X, \underline{\Lambda})$ sont des $\Lambda$-modules de type fini.
Nous allons employer une astuce; l'argument « correct » utilise
un théorème sur les coefficients que nous démontrerons plus loin.
Choisissons une base $\Lambda = \bigoplus_{i \in I} \mathbf{F}_\ell e_i$
telle que $e_0 = 1$ pour un certain $0 \in I$.
Le choix de cette base détermine un isomorphisme
$$
\underline{\Lambda} = \bigoplus_{i \in I} \underline{\mathbf{F}_\ell} e_i
$$
de faisceaux sur $X_\etale$. Nous voyons donc que
$$
H^q_\etale(X, \underline{\Lambda}) =
H^q_\etale(X, \bigoplus_{i \in I} \underline{\mathbf{F}_\ell} e_i) =
\bigoplus_{i \in I} H^q_\etale(X, \underline{\mathbf{F}_\ell})e_i
$$
car la cohomologie sur $X$ commute aux sommes directes d'après le
théorème \ref{theorem-colimit} (ou le lemme \ref{lemma-colimit}, ou le
lemme \ref{lemma-direct-sum-bounded-below-cohomology}).
Comme nous savons déjà que $H^q_\etale(X, \underline{\mathbf{F}_\ell})$
est un $\mathbf{F}_\ell$-espace vectoriel de dimension finie
(théorème \ref{theorem-vanishing-affine-curves}),
nous voyons que $H^q_\etale(X, \underline{\Lambda})$
est libre sur $\Lambda$ et de même rang.
En effet, si $\xi_1, \ldots, \xi_m$ est une base de
$H^q_\etale(X, \underline{\mathbf{F}_\ell})$, alors
$\xi_1 e_0, \ldots, \xi_m e_0$ forment une $\Lambda$-base de
$H^q_\etale(X, \underline{\Lambda})$.
\end{proof}

\begin{lemma}
\label{lemma-finite-pushforward-coefficients}
Soit $\Lambda$ un anneau noethérien, soit $k$ un corps algébriquement clos,
soit $f : X \to Y$ un morphisme fini de schémas séparés de type fini
sur $k$ et de dimension $\leq 1$, et soit $\mathcal{F}$ un faisceau
de $\Lambda$-modules sur $X_\etale$. Si
$H^q_\etale(X, \mathcal{F})$ est un $\Lambda$-module de type fini, alors
$H^q_\etale(Y, f_*\mathcal{F})$ l'est aussi.
\end{lemma}

\begin{proof}
En effet, nous avons $H^q_\etale(X, \mathcal{F}) = H^q_\etale(Y, f_*\mathcal{F})$
par l'annulation de $R^qf_*$ pour $q > 0$
(proposition \ref{proposition-finite-higher-direct-image-zero}) et par
la suite spectrale de Leray
(Cohomologie sur les sites, lemme \ref{sites-cohomology-lemma-apply-Leray}).
\end{proof}

\begin{lemma}
\label{lemma-restrict-to-open-coefficients}
Soit $\Lambda$ un anneau noethérien, soit $k$ un corps algébriquement clos,
soit $X$ un schéma séparé de type fini sur $k$ et de dimension $\leq 1$,
soit $\mathcal{F}$ un faisceau constructible de $\Lambda$-modules sur $X_\etale$,
et soit $j : X' \to X$ l'inclusion d'un sous-schéma ouvert dense. Alors
$H^q_\etale(X, \mathcal{F})$ est un $\Lambda$-module de type fini si et seulement si
$H^q_\etale(X, j_!j^{-1}\mathcal{F})$ est un $\Lambda$-module de type fini.
\end{lemma}

\begin{proof}
Puisque $X'$ est dense, nous voyons que $Z = X \setminus X'$ est de dimension $0$
et est donc un ensemble fini $Z = \{x_1, \ldots, x_n\}$ de points $k$-rationnels.
Considérons la suite exacte courte
$$
0 \to j_!j^{-1}\mathcal{F} \to \mathcal{F} \to i_*i^{-1}\mathcal{F} \to 0
$$
du lemme \ref{lemma-ses-associated-to-open}. Observons que
$H^q_\etale(X, i_*i^{-1}\mathcal{F}) = H^q_\etale(Z, i^{-1}\mathcal{F})$.
En effet, $i : Z \to X$ est une immersion fermée, donc finie,
et nous avons donc l'annulation de $R^qi_*$ pour $q > 0$ d'après la
proposition \ref{proposition-finite-higher-direct-image-zero};
l'égalité résulte alors de la suite spectrale de Leray
(Cohomologie sur les sites, lemme \ref{sites-cohomology-lemma-apply-Leray}).
Puisque $Z$ est une réunion disjointe de spectres de corps algébriquement
clos, nous en déduisons que $H^q_\etale(Z, i^{-1}\mathcal{F}) = 0$
pour $q > 0$ et que
$$
H^0_\etale(Z, i^{-1}\mathcal{F}) =
\bigoplus\nolimits_{i=1}^n \mathcal{F}_{x_i}
$$
qui est un $\Lambda$-module de type fini, car $\mathcal{F}_{x_i}$ est de type fini en vertu de
l'hypothèse que $\mathcal{F}$ est un faisceau constructible de $\Lambda$-modules.
La suite exacte longue de cohomologie donne une suite exacte
$$
0 \to
H^0_\etale(X, j_!j^{-1}\mathcal{F}) \to
H^0_\etale(X, \mathcal{F}) \to
H^0_\etale(Z, i^{-1}\mathcal{F}) \to
H^1_\etale(X, j_!j^{-1}\mathcal{F}) \to
H^1_\etale(X, \mathcal{F}) \to 0
$$
et des isomorphismes $H^q_\etale(X, j_!j^{-1}\mathcal{F}) \to
H^q_\etale(X, \mathcal{F})$ pour $q > 1$. Le lemme en résulte aisément.
\end{proof}

\begin{lemma}
\label{lemma-somewhat-easier-coefficients}
Soit $\Lambda$ un anneau noethérien, soit $M$ un $\Lambda$-module de type fini annulé
par un entier $n > 0$, soit $k$ un corps algébriquement clos, soit $X$ un
schéma séparé de type fini et de dimension $\leq 1$ sur $k$, et soit
$j : U \to X$ une immersion ouverte. Alors
\begin{enumerate}
\item $H^q_\etale(X, j_!\underline{M})$ est un $\Lambda$-module de type fini
si $n$ est premier à $\text{char}(k)$,
\item $H^q_\etale(X, j_!\underline{M})$ est un $\Lambda$-module de type fini
si $X$ est propre.
\end{enumerate}
\end{lemma}

\begin{proof}
Comme $\dim(X) \leq 1$, il existe un ouvert $V \subset X$ disjoint de
$U$ tel que $X' = U \cup V$ soit ouvert dense dans $X$ (détails omis).
Si $j' : X' \to X$ désigne le morphisme d'inclusion, nous voyons que
$j_!\underline{M}$ est un facteur direct de $j'_!\underline{M}$.
Il suffit donc de démontrer le lemme lorsque $U$ est ouvert et dense dans $X$.
Ce cas résulte des lemmes \ref{lemma-restrict-to-open-coefficients}
et \ref{lemma-constant-statements-coefficients}.
\end{proof}

\begin{lemma}
\label{lemma-ses-statements-coefficients}
Soit $\Lambda$ un anneau noethérien, soit $k$ un corps algébriquement clos,
soit $X$ un schéma séparé de type fini sur $k$ et de dimension $\leq 1$, et soit
$0 \to \mathcal{F}_1 \to \mathcal{F} \to \mathcal{F}_2 \to 0$
une suite exacte courte de faisceaux de $\Lambda$-modules sur $X_\etale$. Si
les $H^q_\etale(X, \mathcal{F}_i)$, $i = 1, 2$, sont des $\Lambda$-modules de type fini,
alors $H^q_\etale(X, \mathcal{F})$ est un $\Lambda$-module de type fini.
\end{lemma}

\begin{proof}
Cela résulte immédiatement de la suite exacte longue de cohomologie.
\end{proof}

\begin{lemma}
\label{lemma-vanishing-easier-coefficients}
Soit $\Lambda$ un anneau noethérien, soit $k$ un corps algébriquement clos,
soit $X$ un schéma séparé de type fini et de dimension $\leq 1$ sur $k$,
soit $j : U \to X$ une immersion ouverte avec $U$ connexe, soit
$\ell$ un nombre premier, soit $n > 0$, et soit $\mathcal{G}$ un faisceau
localement constant de type fini de $\Lambda$-modules sur $U_\etale$, annulé par
$\ell^n$. Alors
\begin{enumerate}
\item $H^q_\etale(X, j_!\mathcal{G})$ est un $\Lambda$-module de type fini
si $\ell$ est premier à $\text{char}(k)$,
\item $H^q_\etale(X, j_!\mathcal{G})$ est un $\Lambda$-module de type fini
si $X$ est propre.
\end{enumerate}
\end{lemma}

\begin{proof}
Soit $f : V \to U$ un morphisme fini étale de degré premier à $\ell$
comme dans le lemme \ref{lemma-pullback-filtered-modules}. La discussion de la
section \ref{section-trace-method} donne des applications
$$
\mathcal{G} \to f_*f^{-1}\mathcal{G} \to \mathcal{G}
$$
dont le composé est un isomorphisme. Il suffit donc de démontrer la
finitude de $H^q_\etale(X, j_!f_*f^{-1}\mathcal{G})$.
En vertu du théorème principal de Zariski
(Compléments sur les morphismes, lemme
\ref{more-morphisms-lemma-quasi-finite-separated-pass-through-finite})
nous pouvons choisir un diagramme
$$
\xymatrix{
V \ar[r]_{j'} \ar[d]_f & Y \ar[d]^{\overline{f}} \\
U \ar[r]^j & X
}
$$
où $\overline{f} : Y \to X$ est fini et $j'$ une immersion ouverte d'image
dense. Comme $f$ est fini et $V$ dense dans $Y$, nous avons
$V = U \times_X Y$. D'après le lemme \ref{lemma-compatible-shriek-push-finite}, nous avons
$$
j_!f_*f^{-1}\mathcal{G} = \overline{f}_*j'_!f^{-1}\mathcal{G}
$$
D'après le lemme \ref{lemma-finite-pushforward-coefficients}, il suffit de
considérer $j'_!f^{-1}\mathcal{G}$. L'existence de la filtration donnée par le
lemme \ref{lemma-pullback-filtered-modules},
le fait que $j'_!$ est exact et le
lemme \ref{lemma-ses-statements-coefficients}
nous ramènent au cas
$\mathcal{F} = j'_!\underline{M}$ pour un $\Lambda$-module de type fini $M$,
qui est le lemme \ref{lemma-somewhat-easier-coefficients}.
\end{proof}

\begin{theorem}
\label{theorem-vanishing-affine-curves-coefficients}
Soit $\Lambda$ un anneau noethérien, soit $k$ un corps algébriquement clos,
soit $X$ un schéma séparé de type fini et de dimension $\leq 1$ sur $k$,
et soit $\mathcal{F}$ un faisceau constructible de $\Lambda$-modules
de torsion sur $X_\etale$. Alors
\begin{enumerate}
\item
\label{item-finite-prime-to-p-coefficients}
$H^q_\etale(X, \mathcal{F})$ est un $\Lambda$-module de type fini
si $\mathcal{F}$ est de torsion première à $\text{char}(k)$,
\item
\label{item-finite-proper-coefficients}
$H^q_\etale(X, \mathcal{F})$ est un $\Lambda$-module de type fini si $X$ est propre.
\end{enumerate}
\end{theorem}

\begin{proof}
Nous appliquerons le dévissage par suites exactes sans autre mention.
Écrivons $\mathcal{F} = \mathcal{F}_1 \oplus \ldots \oplus \mathcal{F}_r$
où $\mathcal{F}_i$ est annulé par $\ell_i^{n_i}$
pour un nombre premier $\ell_i$ et un entier $n_i > 0$.
D'après le lemme \ref{lemma-ses-statements-coefficients}, il suffit
de démontrer le théorème pour $\mathcal{F}_i$. Nous pouvons donc supposer,
et supposons, que $\ell^n$ annule $\mathcal{F}$ pour un nombre premier $\ell$
et un entier $n > 0$.

\medskip\noindent
Puisque $\mathcal{F}$ est constructible comme faisceau de $\Lambda$-modules,
il existe un ouvert dense $U \subset X$ tel que $\mathcal{F}|_U$ soit un
faisceau localement constant de type fini de $\Lambda$-modules.
Comme $\dim(X) \leq 1$, nous pouvons supposer, après avoir rétréci $U$, que
$U = U_1 \amalg \ldots \amalg U_n$ est une réunion disjointe de schémas
irréductibles (il suffit de supprimer les points fermés qui appartiennent aux
intersections d'au moins $2$ composantes de $U$). D'après le
lemme \ref{lemma-restrict-to-open-coefficients},
nous nous ramenons au cas
$\mathcal{F} = j_!\mathcal{G}$ où $\mathcal{G}$ est un faisceau localement
constant de type fini de $\Lambda$-modules sur $U$ (et annulé
par $\ell^n$).

\medskip\noindent
Comme nous avons choisi $U = U_1 \amalg \ldots \amalg U_n$ avec les $U_i$
irréductibles, nous avons
$$
j_!\mathcal{G} =
j_{1!}(\mathcal{G}|_{U_1}) \oplus \ldots \oplus
j_{n!}(\mathcal{G}|_{U_n})
$$
où $j_i : U_i \to X$ est le morphisme d'inclusion.
Le cas de $j_{i!}(\mathcal{G}|_{U_i})$ est traité dans le
lemme \ref{lemma-vanishing-easier-coefficients}.
\end{proof}








\section{Première cohomologie des schémas propres}
\label{section-finite-etale-over-proper}

\noindent
Dans Fundamental Groups, section \ref{pione-section-finite-etale-over-proper},
nous avons vu, en un certain sens, que la formation de
$R^1f_*\underline{G}$ commute au changement de base si $f : X \to Y$
est un morphisme propre et si $G$ est un groupe fini (non nécessairement
commutatif). Dans cette section,
nous déduisons une conséquence utile de ces résultats.

\begin{lemma}
\label{lemma-proper-over-henselian-and-h1}
Soit $A$ un anneau local hensélien. Soit $X$ un schéma propre sur $A$
de fibre fermée $X_0$. Soit $M$ un groupe abélien fini.
Alors $H^1_\etale(X, \underline{M}) = H^1_\etale(X_0, \underline{M})$.
\end{lemma}

\begin{proof}
D'après Cohomologie sur les sites, lemme \ref{sites-cohomology-lemma-torsors-h1},
un élément de $H^1_\etale(X, \underline{M})$ correspond à un
$\underline{M}$-torseur $\mathcal{F}$ sur $X_\etale$.
Un tel torseur est manifestement un faisceau localement constant fini.
Ainsi, $\mathcal{F}$ est représentable par un schéma $V$ fini
étale sur $X$, lemme \ref{lemma-characterize-finite-locally-constant}.
Réciproquement, un schéma $V$ fini étale sur $X$, muni d'une action de $M$
qui en fait un $M$-torseur sur $X$, donne une classe de cohomologie.
La même correspondance entre les classes de cohomologie sur $X_0$ et les
torseurs finis étales sur $X_0$ vaut.
Le lemme résulte donc de
l'équivalence de catégories de Fundamental Groups, lemme
\ref{pione-lemma-finite-etale-on-proper-over-henselian}.
\end{proof}

\noindent
Le lemme technique suivant est un ingrédient essentiel de la démonstration du
théorème de changement de base propre. L'argument vaut mot pour mot
pour tout schéma propre sur $A$ dont la fibre spéciale est de dimension
$\leq 1$, mais la conclusion sera en fait une conséquence du
théorème de changement de base propre, et nous n'avons besoin que de cette
version particulière pour le démontrer.

\begin{lemma}
\label{lemma-efface-cohomology-on-fibre-by-finite-cover}
Soit $A$ un anneau local hensélien. Soit $X = \mathbf{P}^1_A$.
Soit $X_0 \subset X$ la fibre fermée. Soit $\ell$ un nombre
premier. Soit $\mathcal{I}$ un faisceau injectif de
$\mathbf{Z}/\ell\mathbf{Z}$-modules sur $X_\etale$. Alors
$H^q_\etale(X_0, \mathcal{I}|_{X_0}) = 0$ pour $q > 0$.
\end{lemma}

\begin{proof}
Observons que $X$ est un schéma séparé que l'on peut recouvrir par $2$
ouverts affines. Pour $q > 1$, le résultat découle donc de la variante affine
du théorème de changement de base propre due à Gabber, voir le
lemme \ref{lemma-vanishing-restriction-injective}.
Nous pouvons donc supposer $q = 1$. Soit
$\xi \in H^1_\etale(X_0, \mathcal{I}|_{X_0})$.
Il s'agit de montrer que $\xi$ est nul.
D'après les lemmes \ref{lemma-torsion-colimit-constructible} et
\ref{lemma-colimit}, nous pouvons trouver un morphisme $\mathcal{F} \to \mathcal{I}$
où $\mathcal{F}$ est un faisceau constructible de
$\mathbf{Z}/\ell\mathbf{Z}$-modules
et où $\xi$ provient d'un élément $\zeta$ de
$H^1_\etale(X_0, \mathcal{F}|_{X_0})$. Supposons donné un morphisme injectif
$\mathcal{F} \to \mathcal{F}'$ de faisceaux de
$\mathbf{Z}/\ell\mathbf{Z}$-modules sur $X_\etale$.
Comme $\mathcal{I}$ est injectif, nous pouvons prolonger
le morphisme donné $\mathcal{F} \to \mathcal{I}$ en un morphisme $\mathcal{F}' \to \mathcal{I}$.
Dans cette situation, nous pouvons remplacer $\mathcal{F}$ par $\mathcal{F}'$
et $\zeta$ par son image dans $H^1_\etale(X_0, \mathcal{F}'|_{X_0})$.
De même, si $\mathcal{F} = \mathcal{F}_1 \oplus \mathcal{F}_2$ est une somme directe,
nous pouvons remplacer $\mathcal{F}$ par $\mathcal{F}_i$
et $\zeta$ par son image dans $H^1_\etale(X_0, \mathcal{F}_i|_{X_0})$.

\medskip\noindent
D'après le lemme \ref{lemma-constructible-maps-into-constant-general}
et les remarques ci-dessus, nous pouvons supposer que $\mathcal{F}$
est de la forme $f_*\underline{M}$, où $M$ est un
$\mathbf{Z}/\ell\mathbf{Z}$-module de type fini
et $f : Y \to X$ est un morphisme fini de présentation finie
(de tels faisceaux restent constructibles d'après le
lemme \ref{lemma-finite-pushforward-constructible},
mais nous n'en aurons pas besoin).
Puisque la formation de $f_*$ commute à tout changement de base
(lemme \ref{lemma-finite-pushforward-commutes-with-base-change}),
nous voyons que la restriction de $f_*\underline{M}$ à $X_0$ est
égale à l'image directe de $\underline{M}$ par le morphisme induit
$Y_0 \to X_0$ entre les fibres spéciales. La suite spectrale de Leray
(proposition \ref{proposition-leray})
et l'annulation des images directes supérieures
(proposition \ref{proposition-finite-higher-direct-image-zero})
donnent
$$
H^1_\etale(X_0, f_*\underline{M}|_{X_0}) = H^1_\etale(Y_0, \underline{M}).
$$
Puisque $Y \to \Spec(A)$ est propre, le
lemme \ref{lemma-proper-over-henselian-and-h1} montre que
$H^1_\etale(Y_0, \underline{M})$ est égal à
$H^1_\etale(Y, \underline{M})$. Ainsi, notre classe de cohomologie
$\zeta$ se relève en une classe de cohomologie
$$
\tilde\zeta \in H^1_\etale(Y, \underline{M}) = H^1_\etale(X, f_*\underline{M})
$$
Cependant, $\tilde \zeta$ a une image nulle dans
$H^1_\etale(X, \mathcal{I})$ puisque $\mathcal{I}$ est injectif;
par commutativité de
$$
\xymatrix{
H^1_\etale(X, f_*\underline{M}) \ar[r] \ar[d] &
H^1_\etale(X, \mathcal{I}) \ar[d] \\
H^1_\etale(X_0, (f_*\underline{M})|_{X_0}) \ar[r] &
H^1_\etale(X_0, \mathcal{I}|_{X_0})
}
$$
nous en concluons que l'image $\xi$ de $\zeta$ est nulle elle aussi.
\end{proof}











\section{Préliminaires sur le changement de base}
\label{section-base-change-preliminaries}

\noindent
Si seul le théorème de changement de base lisse ou le théorème de changement
de base propre vous intéresse, vous pouvez passer directement aux
sections correspondantes.
Dans cette section et les quelques suivantes, nous considérons des diagrammes commutatifs
$$
\xymatrix{
X \ar[d]_f & Y \ar[l]^h \ar[d]^e \\
S & T \ar[l]_g
}
$$
de schémas; nous supposons généralement ce diagramme cartésien,
c'est-à-dire $Y = X \times_S T$. Un diagramme commutatif comme ci-dessus
donne lieu à un diagramme commutatif
$$
\xymatrix{
X_\etale \ar[d]_{f_{small}} & Y_\etale \ar[d]^{e_{small}} \ar[l]^{h_{small}} \\
S_\etale & T_\etale \ar[l]_{g_{small}}
}
$$
de petits sites étales. Employons les notations
$$
f^{-1} = f_{small}^{-1}, \quad
g_* = g_{small, *}, \quad
e^{-1} = e_{small}^{-1}, \text{ et}\quad
h_* = h_{small, *}.
$$
D'après Sites, section \ref{sites-section-pullback},
nous obtenons un morphisme de changement de base ou d'image inverse
$$
f^{-1}g_*\mathcal{F}
\longrightarrow
h_*e^{-1}\mathcal{F}
$$
pour un faisceau $\mathcal{F}$ sur $T_\etale$. Si $\mathcal{F}$ est un faisceau
abélien sur $T_\etale$, nous obtenons alors un morphisme de changement de base dérivé
$$
f^{-1}Rg_*\mathcal{F}
\longrightarrow
Rh_*e^{-1}\mathcal{F}
$$
voir Cohomologie sur les sites, lemme
\ref{sites-cohomology-lemma-base-change-map-flat-case}.
Enfin, si $K$ est un objet arbitraire de $D(T_\etale)$,
il existe un morphisme de changement de base
$$
f^{-1}Rg_*K
\longrightarrow
Rh_*e^{-1}K
$$
voir
Cohomologie sur les sites, remarque \ref{sites-cohomology-remark-base-change}.

\begin{lemma}
\label{lemma-base-change-local}
Considérons un diagramme cartésien de schémas
$$
\xymatrix{
X \ar[d]_f & Y \ar[l]^h \ar[d]^e \\
S & T \ar[l]_g
}
$$
Soit $\{U_i \to X\}$ un recouvrement étale tel que $U_i \to S$
se factorise en $U_i \to V_i \to S$, avec $V_i \to S$ étale,
et considérons les diagrammes cartésiens
$$
\xymatrix{
U_i \ar[d]_{f_i} & U_i \times_X Y \ar[l]^{h_i} \ar[d]^{e_i} \\
V_i & V_i \times_S T \ar[l]_{g_i}
}
$$
Soit $\mathcal{F}$ un faisceau sur $T_\etale$. Soit $K \in D(T_\etale)$.
Posons $K_i = K|_{V_i \times_S T}$ et
$\mathcal{F}_i = \mathcal{F}|_{V_i \times_S T}$.
\begin{enumerate}
\item Si $f_i^{-1}g_{i, *}\mathcal{F}_i = h_{i, *}e_i^{-1}\mathcal{F}_i$
pour tout $i$, alors $f^{-1}g_*\mathcal{F} = h_*e^{-1}\mathcal{F}$.
\item Si $f_i^{-1}Rg_{i, *}K_i = Rh_{i, *}e_i^{-1}K_i$
pour tout $i$, alors $f^{-1}Rg_*K = Rh_*e^{-1}K$.
\item Si $\mathcal{F}$ est un faisceau abélien et si
$f_i^{-1}R^qg_{i, *}\mathcal{F}_i = R^qh_{i, *}e_i^{-1}\mathcal{F}_i$
pour tout $i$, alors
$f^{-1}R^qg_*\mathcal{F} = R^qh_*e^{-1}\mathcal{F}$.
\end{enumerate}
\end{lemma}

\begin{proof}
Démonstration de (1). Observons d'abord que
$$
(f^{-1}g_*\mathcal{F})|_{U_i} =
f_i^{-1}(g_*\mathcal{F}|_{V_i}) =
f_i^{-1}g_{i, *}\mathcal{F}_i
$$
La première égalité vient de ce que $U_i \to X \to S$ est égal à
$U_i \to V_i \to S$, et la seconde de ce que
$g_*\mathcal{F}|_{V_i} = g_{i, *}\mathcal{F}_i$ par le
lemme \ref{sites-lemma-localize-morphism-strong} de Sites.
De même, nous avons
$$
(h_*e^{-1}\mathcal{F})|_{U_i} =
h_{i, *}(e^{-1}\mathcal{F}|_{U_i \times_X Y}) =
h_{i, *}e_i^{-1}\mathcal{F}_i
$$
Ainsi, si les morphismes de changement de base
$f_i^{-1}g_{i, *}\mathcal{F}_i \to h_{i, *}e_i^{-1}\mathcal{F}_i$
sont des isomorphismes pour tout $i$, alors le morphisme de changement de base
$f^{-1}g_*\mathcal{F} \to h_*e^{-1}\mathcal{F}$
se restreint en un isomorphisme sur $U_i$ pour tout $i$,
et nous concluons qu'il est un isomorphisme puisque $\{U_i \to X\}$
est un recouvrement étale.

\medskip\noindent
Pour les deux autres assertions, nous remplaçons l'appel au
lemme \ref{sites-lemma-localize-morphism-strong} de Sites
par un appel au
lemme de Cohomologie sur les sites
\ref{sites-cohomology-lemma-restrict-direct-image-open}.
\end{proof}

\begin{lemma}
\label{lemma-base-change-compose}
Considérons une tour de diagrammes cartésiens de schémas
$$
\xymatrix{
W \ar[d]_i & Z \ar[l]^j \ar[d]^k \\
X \ar[d]_f & Y \ar[l]^h \ar[d]^e \\
S & T \ar[l]_g
}
$$
Soit $K \in D(T_\etale)$. Si
$$
f^{-1}Rg_*K \to Rh_*e^{-1}K
\quad\text{et}\quad
i^{-1}Rh_*e^{-1}K \to Rj_*k^{-1}e^{-1}K
$$
sont des isomorphismes, alors
$(f \circ i)^{-1}Rg_*K \to Rj_*(e \circ k)^{-1}K$
est un isomorphisme.
De même, si $\mathcal{F}$ est un faisceau abélien sur $T_\etale$ et si
$$
f^{-1}R^qg_*\mathcal{F} \to R^qh_*e^{-1}\mathcal{F}
\quad\text{et}\quad
i^{-1}R^qh_*e^{-1}\mathcal{F} \to R^qj_*k^{-1}e^{-1}\mathcal{F}
$$
sont des isomorphismes, alors
$(f \circ i)^{-1}R^qg_*\mathcal{F} \to R^qj_*(e \circ k)^{-1}\mathcal{F}$
est un isomorphisme.
\end{lemma}

\begin{proof}
Cela est formel, pourvu que l'on vérifie que la composée de ces
morphismes de changement de base est le morphisme de changement de base
du rectangle extérieur; voir la remarque de Cohomologie sur les sites
\ref{sites-cohomology-remark-compose-base-change-horizontal}.
\end{proof}

\begin{lemma}
\label{lemma-base-change-Rf-star-colim}
Soit $I$ un ensemble ordonné filtrant. Considérons un système inverse de
diagrammes cartésiens de schémas
$$
\xymatrix{
X_i \ar[d]_{f_i} & Y_i \ar[l]^{h_i} \ar[d]^{e_i} \\
S_i & T_i \ar[l]_{g_i}
}
$$
dont les morphismes de transition sont affines et où $g_i$ est quasi-compact et
quasi-séparé. Posons $X = \lim X_i$,
$S = \lim S_i$, $T = \lim T_i$ et $Y = \lim Y_i$ afin
d'obtenir le diagramme cartésien
$$
\xymatrix{
X \ar[d]_f & Y \ar[l]^h \ar[d]^e \\
S & T \ar[l]_g
}
$$
Soit $(\mathcal{F}_i, \varphi_{i'i})$ un système de faisceaux sur
$(T_i)$ comme dans la définition \ref{definition-inverse-system-sheaves}. Posons
$\mathcal{F} = \colim p_i^{-1}\mathcal{F}_i$ sur $T$,
où $p_i : T \to T_i$ est la projection.
Alors les assertions suivantes sont vraies.
\begin{enumerate}
\item Si $f_i^{-1}g_{i, *}\mathcal{F}_i = h_{i, *}e_i^{-1}\mathcal{F}_i$
pour tout $i$, alors
$f^{-1}g_*\mathcal{F} = h_*e^{-1}\mathcal{F}$.
\item Si $\mathcal{F}_i$ est un faisceau abélien pour tout $i$ et si
$f_i^{-1}R^qg_{i, *}\mathcal{F}_i = R^qh_{i, *}e_i^{-1}\mathcal{F}_i$
pour tout $i$, alors
$f^{-1}R^qg_*\mathcal{F} = R^qh_*e^{-1}\mathcal{F}$.
\end{enumerate}
\end{lemma}

\begin{proof}
Nous démontrons (2) et omettons la démonstration de (1). Nous utiliserons sans
autre mention le fait que l'image inverse des faisceaux commute aux limites inductives,
car c'est un adjoint à gauche. Remarquons que $h_i$ est quasi-compact et quasi-séparé en tant que
changement de base de $g_i$.
En notant $q_i : Y \to Y_i$ les projections, remarquons que
$e^{-1}\mathcal{F} = \colim e^{-1}p_i^{-1}\mathcal{F}_i =
\colim q_i^{-1}e_i^{-1}\mathcal{F}_i$.
Par le lemme \ref{lemma-relative-colimit-general},
cela donne
$$
R^qh_*e^{-1}\mathcal{F} = \colim r_i^{-1}R^qh_{i, *}e_i^{-1}\mathcal{F}_i
$$
où $r_i : X \to X_i$ est la projection.
De même, nous avons
$$
f^{-1}R^qg_*\mathcal{F} =
f^{-1}\colim s_i^{-1}R^qg_{i, *}\mathcal{F}_i =
\colim r_i^{-1}f_i^{-1}R^qg_{i, *}\mathcal{F}_i
$$
où $s_i : S \to S_i$ est la projection. Le lemme en résulte.
\end{proof}

\begin{lemma}
\label{lemma-base-change-Rf-star-colim-complexes}
Soient $I$, $X_i$, $Y_i$, $S_i$, $T_i$, $f_i$, $h_i$, $e_i$, $g_i$,
$X$, $Y$, $S$, $T$, $f$, $h$, $e$, $g$ comme dans l'énoncé
du lemme \ref{lemma-base-change-Rf-star-colim}.
Soient $0 \in I$ et $K_0 \in D^+(T_{0, \etale})$.
Pour $i \in I$, $i \geq 0$, notons $K_i$ l'image inverse de
$K_0$ sur $T_i$. Notons $K$ l'image inverse de $K_0$ sur $T$.
Si $f_i^{-1}Rg_{i, *}K_i = Rh_{i, *}e_i^{-1}K_i$
pour tout $i \geq 0$, alors $f^{-1}Rg_*K = Rh_*e^{-1}K$.
\end{lemma}

\begin{proof}
Il suffit de montrer que le morphisme de changement de base
$f^{-1}Rg_*K \to Rh_*e^{-1}K$ induit un isomorphisme
sur les faisceaux de cohomologie. Autrement dit, nous devons montrer
que $f^{-1}R^pg_*K \to R^ph_*e^{-1}K$ est un isomorphisme
pour tout $p \in \mathbf{Z}$, sachant que
$f_i^{-1}R^pg_{i, *}K_i \to R^ph_{i, *}e_i^{-1}K_i$
est un isomorphisme pour tous $i \geq 0$ et $p \in \mathbf{Z}$.
À ce point, nous pouvons raisonner exactement comme dans la démonstration du
lemme \ref{lemma-base-change-Rf-star-colim},
en remplaçant le renvoi au
lemme \ref{lemma-relative-colimit-general}
par un renvoi au
lemme \ref{lemma-relative-colimit-general-complexes}.
\end{proof}

\begin{lemma}
\label{lemma-base-change-f-star-general-stalks}
Considérons un diagramme cartésien de schémas
$$
\xymatrix{
X \ar[d]_f & Y \ar[l]^h \ar[d]^e \\
S & T \ar[l]_g
}
$$
où $g : T \to S$ est quasi-compact et quasi-séparé.
Soit $\mathcal{F}$ un
faisceau abélien sur $T_\etale$. Soit $q \geq 0$. Les assertions suivantes sont équivalentes.
\begin{enumerate}
\item Pour tout point géométrique $\overline{x}$ de $X$, d'image
$\overline{s} = f(\overline{x})$, nous avons
$$
H^q(\Spec(\mathcal{O}^{sh}_{X, \overline{x}}) \times_S T, \mathcal{F})
=
H^q(\Spec(\mathcal{O}^{sh}_{S, \overline{s}}) \times_S T, \mathcal{F})
$$
\item $f^{-1}R^qg_*\mathcal{F} \to R^qh_*e^{-1}\mathcal{F}$
est un isomorphisme.
\end{enumerate}
\end{lemma}

\begin{proof}
Puisque $Y = X \times_S T$, nous avons
$\Spec(\mathcal{O}^{sh}_{X, \overline{x}}) \times_X Y =
\Spec(\mathcal{O}^{sh}_{X, \overline{x}}) \times_S T$. Ainsi,
le morphisme de (1) est le morphisme sur les fibres en $\overline{x}$ associé au morphisme
de (2), par le théorème \ref{theorem-higher-direct-images} (et le
lemme \ref{lemma-stalk-pullback}).
Le résultat découle donc du théorème \ref{theorem-exactness-stalks}.
\end{proof}

\begin{lemma}
\label{lemma-check-stalks-better}
Soit $f : X \to S$ un morphisme de schémas.
Soit $\overline{x}$ un point géométrique de $X$, d'image $\overline{s}$ dans $S$.
Soit $\Spec(K) \to \Spec(\mathcal{O}^{sh}_{S, \overline{s}})$
un morphisme, où $K$ est un corps séparablement clos. Soit $\mathcal{F}$ un
faisceau abélien sur $\Spec(K)_\etale$. Soit $q \geq 0$. Les assertions suivantes sont
équivalentes.
\begin{enumerate}
\item
$H^q(\Spec(\mathcal{O}^{sh}_{X, \overline{x}}) \times_S \Spec(K), \mathcal{F}) =
H^q(\Spec(\mathcal{O}^{sh}_{S, \overline{s}}) \times_S \Spec(K), \mathcal{F})$
\item
$H^q(\Spec(\mathcal{O}^{sh}_{X, \overline{x}})
\times_{\Spec(\mathcal{O}^{sh}_{S, \overline{s}})} \Spec(K), \mathcal{F}) =
H^q(\Spec(K), \mathcal{F})$
\end{enumerate}
\end{lemma}

\begin{proof}
Remarquons que $\Spec(K) \times_S \Spec(\mathcal{O}^{sh}_{S, \overline{s}})$
est le spectre d'une limite inductive filtrante d'algèbres étales sur $K$.
Puisque $K$ est séparablement clos, toute $K$-algèbre étale
est un produit fini de copies de $K$. Nous pouvons donc écrire
$$
\Spec(K) \times_S \Spec(\mathcal{O}^{sh}_{S, \overline{s}}) =
\lim_{i \in I} \coprod\nolimits_{a \in A_i} \Spec(K)
$$
comme une limite cofiltrante dont chaque terme est une réunion disjointe
de copies de $\Spec(K)$ indexées par un ensemble fini $A_i$.
Remarquons que $A_i$ est non vide, puisque nous disposons de
$\Spec(K) \to \Spec(\mathcal{O}^{sh}_{S, \overline{s}})$.
Il s'ensuit que
\begin{align*}
\Spec(\mathcal{O}^{sh}_{X, \overline{x}}) \times_S \Spec(K)
& =
\Spec(\mathcal{O}^{sh}_{X, \overline{x}})
\times_{\Spec(\mathcal{O}^{sh}_{S, \overline{s}})}
\left(
\Spec(\mathcal{O}^{sh}_{S, \overline{s}})
\times_S \Spec(K)\right) \\
& =
\lim_{i \in I} \coprod\nolimits_{a \in A_i}
\Spec(\mathcal{O}^{sh}_{X, \overline{x}})
\times_{\Spec(\mathcal{O}^{sh}_{S, \overline{s}})} \Spec(K)
\end{align*}
Comme, dans notre situation, la cohomologie commute aux limites
de schémas (théorème \ref{theorem-colimit}), nous concluons.
\end{proof}













\section{Changement de base pour l'image directe}
\label{section-base-change-f-star}

\noindent
Cette section est préliminaire et devrait être omise en première lecture.
Nous y étudions pour quels morphismes $f : X \to S$ on a
$f^{-1}g_* = h_*e^{-1}$ sur tous les faisceaux (d'ensembles), pour tout
diagramme cartésien
$$
\xymatrix{
X \ar[d]_f & Y \ar[l]^h \ar[d]^e \\
S & T \ar[l]_g
}
$$
où $g$ est quasi-compact et quasi-séparé.

\begin{lemma}
\label{lemma-base-change-f-star-general}
Considérons le diagramme cartésien de schémas
$$
\xymatrix{
X \ar[d]_f & Y \ar[l]^h \ar[d]^e \\
S & T \ar[l]_g
}
$$
Supposons que $f$ soit plat et que tout objet $U$ de $X_\etale$ possède
un recouvrement $\{U_i \to U\}$ tel que $U_i \to S$
se factorise en $U_i \to V_i \to S$, avec $V_i \to S$
étale et $U_i \to V_i$ quasi-compact à
fibres géométriquement connexes.
Alors, pour tout faisceau d'ensembles $\mathcal{F}$ sur $T_\etale$, nous avons
$f^{-1}g_*\mathcal{F} = h_*e^{-1}\mathcal{F}$.
\end{lemma}

\begin{proof}
Soit $U \to X$ un morphisme étale tel que $U \to S$ se factorise en
$U \to V \to S$, avec $V \to S$ étale et $U \to V$ quasi-compact
à fibres géométriquement connexes. Remarquons que $U \to V$ est plat
(lemme \ref{flat-lemma-etale-flat-up-down} de Compléments sur la platitude).
Nous affirmons que
\begin{align*}
f^{-1}g_*\mathcal{F}(U)
& = g_*\mathcal{F}(V) \\
& = \mathcal{F}(V \times_S T) \\
& = e^{-1}\mathcal{F}(U \times_X Y) \\
& = h_*e^{-1}\mathcal{F}(U)
\end{align*}
En effet, en considérant $U$ comme un objet de $X_\etale$ et
$V$ comme un objet de $S_\etale$, nous voyons que la première égalité
découle du lemme \ref{lemma-sections-upstairs}\footnote{Plus
précisément, nous utilisons aussi le fait que la restriction de $f^{-1}g_*\mathcal{F}$
à $U_\etale$ est l'image inverse par $U \to V$ de la restriction de
$g_*\mathcal{F}$ à $V_\etale$. Voir le
lemme \ref{sites-lemma-localize-morphism-strong} de Sites.}.
En considérant $V \times_S T$ comme un objet de $T_\etale$,
la seconde égalité découle de la définition de $g_*$.
Remarquons que $U \times_X Y = U \times_S T$ (car $Y = X \times_S T$)
et que, par conséquent, $U \times_X Y \to V \times_S T$
est à fibres géométriquement connexes comme changement de base de $U \to V$.
En considérant $U \times_X Y$ comme un objet de $Y_\etale$, nous voyons que
la troisième égalité découle du lemme \ref{lemma-sections-upstairs},
comme précédemment. Enfin, la quatrième égalité découle de la définition
de $h_*$.

\medskip\noindent
Puisque, par hypothèse, tout objet de $X_\etale$ possède un recouvrement
étale auquel s'applique l'argument du paragraphe précédent,
nous voyons que le lemme est vrai.
\end{proof}

\begin{lemma}
\label{lemma-fppf-reduced-fibres-base-change-f-star}
Considérons un diagramme cartésien de schémas
$$
\xymatrix{
X \ar[d]_f & Y \ar[l]^h \ar[d]^e \\
S & T \ar[l]_g
}
$$
où $f$ est plat et localement de présentation finie,
à fibres géométriquement réduites.
Alors $f^{-1}g_*\mathcal{F} = h_*e^{-1}\mathcal{F}$
pour tout faisceau $\mathcal{F}$ sur $T_\etale$.
\end{lemma}

\begin{proof}
Combiner le lemme \ref{lemma-base-change-f-star-general} avec le
lemme de Compléments sur les morphismes
\ref{more-morphisms-lemma-cover-by-geometrically-connected}.
\end{proof}

\begin{lemma}
\label{lemma-base-change-f-star-field}
Considérons le diagramme cartésien de schémas
$$
\xymatrix{
X \ar[d]_f & Y \ar[l]^h \ar[d]^e \\
S & T \ar[l]_g
}
$$
Supposons que $S$ soit le spectre d'un corps séparablement clos.
Alors $f^{-1}g_*\mathcal{F} = h_*e^{-1}\mathcal{F}$
pour tout faisceau $\mathcal{F}$ sur $T_\etale$.
\end{lemma}

\begin{proof}
Nous pouvons travailler localement sur $X$. Nous pouvons donc supposer $X$ affine.
Nous pouvons alors écrire $X$ comme une limite cofiltrante de schémas affines
de type fini sur $S$. Par le lemme \ref{lemma-base-change-Rf-star-colim},
nous pouvons supposer que $X$ est de type fini sur $S$.
Le lemme \ref{lemma-base-change-f-star-general}
s'applique alors, car tout schéma de type
fini sur un corps séparablement clos est une réunion disjointe finie
de schémas connexes et géométriquement connexes
(voir le lemme de Variétés
\ref{varieties-lemma-separably-closed-field-connected-components}).
\end{proof}

\begin{lemma}
\label{lemma-base-change-f-star-valuation}
Considérons un diagramme cartésien de schémas
$$
\xymatrix{
X \ar[d]_f & Y \ar[l]^h \ar[d]^e \\
S & T \ar[l]_g
}
$$
Supposons que
\begin{enumerate}
\item $f$ soit plat et ouvert,
\item les corps résiduels de $S$ soient séparablement algébriquement clos,
\item étant donné un morphisme étale $U \to X$ avec $U$ affine,
nous puissions écrire $U$ comme une réunion disjointe finie de sous-schémas ouverts
de $X$ (par exemple si $X$ est un schéma intègre normal
dont le corps des fonctions est séparablement clos),
\item tout ouvert non vide d'une fibre $X_s$ de $f$ soit connexe
(par exemple si $X_s$ est irréductible ou vide).
\end{enumerate}
Alors, pour tout faisceau d'ensembles $\mathcal{F}$ sur $T_\etale$, nous avons
$f^{-1}g_*\mathcal{F} = h_*e^{-1}\mathcal{F}$.
\end{lemma}

\begin{proof}
Omis. Indication : les hypothèses impliquent presque trivialement
la condition du lemme \ref{lemma-base-change-f-star-general}.
L'exemple de la partie (3) découle du
lemme \ref{lemma-normal-scheme-with-alg-closed-function-field}.
\end{proof}

\noindent
Le lemme suivant n'a pas vraiment sa place ici, mais il ne semble
pas y avoir ailleurs de meilleur endroit où le placer.

\begin{lemma}
\label{lemma-fppf-reduced-fibres-pullback-products}
Soit $f : X \to S$ un morphisme de schémas plat et
localement de présentation finie, à fibres géométriquement réduites.
Alors $f^{-1} : \Sh(S_\etale) \to \Sh(X_\etale)$ commute
aux produits.
\end{lemma}

\begin{proof}
Soit $I$ un ensemble et soit $\mathcal{G}_i$ un faisceau sur $S_\etale$
pour $i \in I$.
Soit $U \to X$ un morphisme étale tel que $U \to S$ se factorise en
$U \to V \to S$, avec $V \to S$ étale et $U \to V$ plat de
présentation finie à fibres géométriquement connexes. Nous avons alors
\begin{align*}
f^{-1}(\prod \mathcal{G}_i)(U)
& =
(\prod \mathcal{G}_i)(V) \\
& =
\prod \mathcal{G}_i(V) \\
& =
\prod f^{-1}\mathcal{G}_i(U) \\
& =
(\prod f^{-1}\mathcal{G}_i)(U)
\end{align*}
où nous avons utilisé le lemme \ref{lemma-sections-upstairs}
dans les première et troisième égalités
(nous utilisons aussi le fait que la restriction de $f^{-1}\mathcal{G}$
à $U_\etale$ est l'image inverse par $U \to V$ de la restriction de
$\mathcal{G}$ à $V_\etale$; voir le
lemme \ref{sites-lemma-localize-morphism-strong} de Sites).
Par le lemme de Compléments sur les morphismes
\ref{more-morphisms-lemma-cover-by-geometrically-connected}
tout objet $U$ de $X_\etale$ possède un recouvrement étale
$\{U_i \to U\}$ tel que la discussion du paragraphe
précédent s'applique à $U_i$. Le lemme en résulte.
\end{proof}

\begin{lemma}
\label{lemma-base-change-f-star}
Soit $f : X \to S$ un morphisme plat de schémas tel
que, pour tout point géométrique $\overline{x}$ de $X$, le morphisme
$$
\mathcal{O}_{S, f(\overline{x})}^{sh}
\longrightarrow
\mathcal{O}_{X, \overline{x}}^{sh}
$$
soit à fibres géométriquement connexes. Alors, pour tout
diagramme cartésien de schémas
$$
\xymatrix{
X \ar[d]_f & Y \ar[l]^h \ar[d]^e \\
S & T \ar[l]_g
}
$$
où $g$ est quasi-compact et quasi-séparé, nous avons
$f^{-1}g_*\mathcal{F} = h_*e^{-1}\mathcal{F}$
pour tout faisceau d'ensembles $\mathcal{F}$ sur $T_\etale$.
\end{lemma}

\begin{proof}
Il suffit de vérifier l'égalité sur les fibres; voir le
théorème \ref{theorem-exactness-stalks}.
Par le théorème \ref{theorem-higher-direct-images}, nous avons
$$
(h_*e^{-1}\mathcal{F})_{\overline{x}} =
\Gamma(\Spec(\mathcal{O}_{X, \overline{x}}^{sh}) \times_X Y, e^{-1}\mathcal{F})
$$
et, de même, nous avons
$$
(f^{-1}g_*\mathcal{F})_{\overline{x}} =
(g_*\mathcal{F})_{f(\overline{x})} =
\Gamma(\Spec(\mathcal{O}_{S, f(\overline{x})}^{sh}) \times_S T, \mathcal{F})
$$
Ces ensembles sont égaux par application du lemme \ref{lemma-sections-upstairs}
au morphisme
$$
\Spec(\mathcal{O}_{X, \overline{x}}^{sh}) \times_X Y
\longrightarrow
\Spec(\mathcal{O}_{S, f(\overline{x})}^{sh}) \times_S T
$$
qui est un changement de base de
$\Spec(\mathcal{O}_{X, \overline{x}}^{sh}) \to
\Spec(\mathcal{O}_{S, f(\overline{x})}^{sh})$
car $Y = X \times_S T$.
\end{proof}







\section{Changement de base pour les images directes supérieures}
\label{section-base-change}

\noindent
Cette section est l'analogue de la section \ref{section-base-change-f-star}
pour les images directes supérieures.
Elle est préliminaire et devrait être omise en première lecture.

\begin{remark}
\label{remark-base-change-holds}
Soit $f : X \to S$ un morphisme de schémas. Soit $n$ un entier.
Nous dirons que $BC(f, n, q_0)$ est vraie si, pour tout diagramme commutatif
$$
\xymatrix{
X \ar[d]_f & X' \ar[l] \ar[d]_{f'} & Y \ar[l]^h \ar[d]^e \\
S & S' \ar[l] & T \ar[l]_g
}
$$
où $X' = X \times_S S'$ et $Y = X' \times_{S'} T$, où
$g$ est quasi-compact et quasi-séparé, et pour tout faisceau abélien
$\mathcal{F}$ sur $T_\etale$ annulé par $n$, le morphisme de changement de base
$$
(f')^{-1}R^qg_*\mathcal{F}
\longrightarrow
R^qh_*e^{-1}\mathcal{F}
$$
est un isomorphisme pour $q \leq q_0$.
\end{remark}

\begin{lemma}
\label{lemma-base-change-q-injective}
Avec $f : X \to S$ et $n$ comme dans la remarque \ref{remark-base-change-holds},
supposons que, pour un certain $q \geq 1$, on ait $BC(f, n, q - 1)$. Alors,
pour tout diagramme commutatif
$$
\xymatrix{
X \ar[d]_f & X' \ar[l] \ar[d]_{f'} & Y \ar[l]^h \ar[d]^e \\
S & S' \ar[l] & T \ar[l]_g
}
$$
où $X' = X \times_S S'$ et $Y = X' \times_{S'} T$, où
$g$ est quasi-compact et quasi-séparé, et pour tout faisceau abélien
$\mathcal{F}$ sur $T_\etale$ annulé par $n$, on a :
\begin{enumerate}
\item le morphisme de changement de base
$(f')^{-1}R^qg_*\mathcal{F}\to R^qh_*e^{-1}\mathcal{F}$
est injectif;
\item si $\mathcal{F} \subset \mathcal{G}$, où $\mathcal{G}$
sur $T_\etale$ est annulé par $n$, alors
$$
\Coker\left(
(f')^{-1}R^qg_*\mathcal{F}\to R^qh_*e^{-1}\mathcal{F}
\right)
\subset
\Coker\left(
(f')^{-1}R^qg_*\mathcal{G}\to R^qh_*e^{-1}\mathcal{G}
\right)
$$
\item si, dans (2), le faisceau $\mathcal{G}$ est un faisceau injectif
de $\mathbf{Z}/n\mathbf{Z}$-modules, alors
$$
\Coker\left((f')^{-1}R^qg_*\mathcal{F}\to R^qh_*e^{-1}\mathcal{F} \right)
\subset R^qh_*e^{-1}\mathcal{G}
$$
\end{enumerate}
\end{lemma}

\begin{proof}
Choisissons une suite exacte courte
$0 \to \mathcal{F} \to \mathcal{I} \to \mathcal{Q} \to 0$
où $\mathcal{I}$ est un faisceau injectif de $\mathbf{Z}/n\mathbf{Z}$-modules.
Considérons le diagramme induit
$$
\xymatrix{
(f')^{-1}R^{q - 1}g_*\mathcal{I} \ar[d]_{\cong} \ar[r] &
(f')^{-1}R^{q - 1}g_*\mathcal{Q} \ar[d]_{\cong} \ar[r] &
(f')^{-1}R^qg_*\mathcal{F} \ar[d] \ar[r] &
0 \ar[d] \\
R^{q - 1}h_*e^{-1}\mathcal{I} \ar[r] &
R^{q - 1}h_*e^{-1}\mathcal{Q} \ar[r] &
R^qh_*e^{-1}\mathcal{F} \ar[r] &
R^qh_*e^{-1}\mathcal{I}
}
$$
à lignes exactes. Le zéro figure dans le coin supérieur droit
car $\mathcal{I}$ est injectif. Les deux flèches verticales de gauche sont
des isomorphismes par $BC(f, n, q - 1)$. Nous concluons que (1) est vraie.
Ce qui précède montre aussi que
$$
\Coker\left(
(f')^{-1}R^qg_*\mathcal{F}\to R^qh_*e^{-1}\mathcal{F}
\right)
\subset
R^qh_*e^{-1}\mathcal{I}
$$
donc (3) est vraie. Pour démontrer (2), choisissons
$\mathcal{F} \subset \mathcal{G} \subset \mathcal{I}$.
\end{proof}

\begin{lemma}
\label{lemma-base-change-q-integral-top}
Avec $f : X \to S$ et $n$ comme dans la remarque \ref{remark-base-change-holds},
supposons que, pour un certain $q \geq 1$, on ait $BC(f, n, q - 1)$. Considérons
les diagrammes commutatifs
$$
\vcenter{
\xymatrix{
X \ar[d]_f &
X' \ar[d]_{f'} \ar[l] &
Y \ar[l]^h \ar[d]^e &
Y' \ar[l]^{\pi'} \ar[d]^{e'} \\
S &
S' \ar[l] &
T \ar[l]_g &
T' \ar[l]_\pi
}
}
\quad\text{et}\quad
\vcenter{
\xymatrix{
X' \ar[d]_{f'} & & Y' \ar[ll]^{h' = h \circ \pi'} \ar[d]^{e'} \\
S' & & T' \ar[ll]_{g' = g \circ \pi}
}
}
$$
où tous les carrés sont cartésiens, $g$ est quasi-compact et quasi-séparé, et
$\pi$ est entier surjectif. Soit $\mathcal{F}$ un faisceau abélien
sur $T_\etale$ annulé par $n$, et posons $\mathcal{F}' = \pi^{-1}\mathcal{F}$.
Si le morphisme de changement de base
$$
(f')^{-1}R^qg'_*\mathcal{F}' \longrightarrow R^qh'_*(e')^{-1}\mathcal{F}'
$$
est un isomorphisme, alors le morphisme de changement de base
$(f')^{-1}R^qg_*\mathcal{F} \to R^qh_*e^{-1}\mathcal{F}$
est un isomorphisme.
\end{lemma}

\begin{proof}
Remarquons que $\mathcal{F} \to \pi_*\mathcal{F}'$ est injectif
car $\pi$ est surjectif (cela se vérifie sur les fibres). Ainsi, par le
lemme \ref{lemma-base-change-q-injective},
nous voyons qu'il suffit de montrer que le morphisme de changement de base
$$
(f')^{-1}R^qg_*\pi_*\mathcal{F}'
\longrightarrow
R^qh_*e^{-1}\pi_*\mathcal{F}'
$$
est un isomorphisme. Cela découle de l'hypothèse, car
nous avons $R^qg_*\pi_*\mathcal{F}' = R^qg'_*\mathcal{F}'$,
nous avons $e^{-1}\pi_*\mathcal{F}' =\pi'_*(e')^{-1}\mathcal{F}'$, et
nous avons $R^qh_*\pi'_*(e')^{-1}\mathcal{F}' = R^qh'_*(e')^{-1}\mathcal{F}'$.
Ces égalités découlent des lemmes
\ref{lemma-integral-pushforward-commutes-with-base-change} et
\ref{lemma-what-integral} et de la suite spectrale de Leray relative
(lemme \ref{sites-cohomology-lemma-relative-Leray} de Cohomologie sur les sites).
\end{proof}

\begin{lemma}
\label{lemma-base-change-q-integral-bottom}
Avec $f : X \to S$ et $n$ comme dans la remarque \ref{remark-base-change-holds},
supposons que, pour un certain $q \geq 1$, on ait $BC(f, n, q - 1)$. Considérons
les diagrammes commutatifs
$$
\vcenter{
\xymatrix{
X \ar[d]_f &
X' \ar[d]_{f'} \ar[l] &
X'' \ar[l]^{\pi'} \ar[d]_{f''} &
Y \ar[l]^{h'} \ar[d]^e \\
S &
S' \ar[l] &
S'' \ar[l]_\pi &
T \ar[l]_{g'}
}
}
\quad\text{et}\quad
\vcenter{
\xymatrix{
X' \ar[d]_{f'} & & Y \ar[ll]^{h = h' \circ \pi'} \ar[d]^e \\
S' & & T \ar[ll]_{g = g' \circ \pi}
}
}
$$
où tous les carrés sont cartésiens, $g'$ est quasi-compact et quasi-séparé, et
$\pi$ est entier. Soit $\mathcal{F}$ un faisceau abélien
sur $T_\etale$ annulé par $n$. Si le morphisme de changement de base
$$
(f')^{-1}R^qg_*\mathcal{F} \longrightarrow R^qh_*e^{-1}\mathcal{F}
$$
est un isomorphisme, alors le morphisme de changement de base
$(f'')^{-1}R^qg'_*\mathcal{F} \to R^qh'_*e^{-1}\mathcal{F}$
est un isomorphisme.
\end{lemma}

\begin{proof}
Puisque $\pi$ et $\pi'$ sont entiers, nous avons $R\pi_* = \pi_*$ et
$R\pi'_* = \pi'_*$; voir le lemme \ref{lemma-what-integral}.
Nous avons aussi $(f')^{-1}\pi_* = \pi'_*(f'')^{-1}$. Nous voyons donc que
$\pi'_*(f'')^{-1}R^qg'_*\mathcal{F} = (f')^{-1}R^qg_*\mathcal{F}$
et
$\pi'_*R^qh'_*e^{-1}\mathcal{F} = R^qh_*e^{-1}\mathcal{F}$.
Ainsi, l'hypothèse signifie que notre morphisme devient un
isomorphisme après application du foncteur $\pi'_*$.
Nous voyons donc que c'est un isomorphisme par le lemme \ref{lemma-what-integral}.
\end{proof}

\begin{lemma}
\label{lemma-formal-argument}
Soit $T$ un schéma quasi-compact et quasi-séparé.
Soit $P$ une propriété des schémas quasi-compacts et quasi-séparés
sur $T$. Supposons que :
\begin{enumerate}
\item si $T'' \to T'$ est un épaississement de schémas quasi-compacts et
quasi-séparés sur $T$, alors $P(T'')$ si et seulement si $P(T')$;
\item si $T' = \lim T_i$ est la limite d'un système inverse de
schémas quasi-compacts et quasi-séparés sur $T$, à morphismes de
transition affines, et si $P(T_i)$ est vraie pour tout $i$, alors
$P(T')$ est vraie;
\item si $Z \subset T'$ est un sous-schéma fermé de complément
quasi-compact $V \subset T'$ et si $P(T')$ est vraie,
alors $P(V)$ ou $P(Z)$ est vraie.
\end{enumerate}
Alors $P(T)$ implique $P(\Spec(K))$ pour un certain morphisme $\Spec(K) \to T$,
où $K$ est un corps.
\end{lemma}

\begin{proof}
Considérons l'ensemble $\mathfrak T$ des sous-schémas fermés $T' \subset T$
tels que $P(T')$ soit vraie. Par l'hypothèse (2), cet ensemble possède un élément minimal,
noté $T'$. Par l'hypothèse (1), nous voyons que $T'$ est réduit.
Soit $\eta \in T'$ le point générique d'une composante
irréductible de $T'$. Alors $\eta = \Spec(K)$ pour un certain corps
$K$, et $\eta = \lim V$, où la limite porte sur les sous-schémas
ouverts affines $V \subset T'$ contenant $\eta$.
Par l'hypothèse (3) et la minimalité de $T'$, nous voyons
que $P(V)$ est vraie pour tous ces $V$. Ainsi $P(\eta)$ est vraie
par (2), ce qui achève la démonstration.
\end{proof}

\begin{lemma}
\label{lemma-base-change-does-not-hold-pre}
Avec $f : X \to S$ et $n$ comme dans la remarque \ref{remark-base-change-holds},
supposons que, pour un certain $q \geq 1$, $BC(f, n, q - 1)$ soit vraie, mais que
$BC(f, n, q)$ ne le soit pas. Il existe alors un diagramme commutatif
$$
\xymatrix{
X \ar[d]_f & X' \ar[d]_{f'} \ar[l] & Y \ar[l]^h \ar[d]^e \\
S & S' \ar[l] & \Spec(K) \ar[l]_g
}
$$
où $X' = X \times_S S'$, $Y = X' \times_{S'} \Spec(K)$,
$K$ est un corps et $\mathcal{F}$ est un faisceau abélien
sur $\Spec(K)$ annulé par $n$, tel que
$(f')^{-1}R^qg_*\mathcal{F} \to R^qh_*e^{-1}\mathcal{F}$
n'est pas un isomorphisme.
\end{lemma}

\begin{proof}
Choisissons un diagramme commutatif
$$
\xymatrix{
X \ar[d]_f & X' \ar[l] \ar[d]_{f'} & Y \ar[l]^h \ar[d]^e \\
S & S' \ar[l] & T \ar[l]_g
}
$$
où $X' = X \times_S S'$ et $Y = X' \times_{S'} T$, où
$g$ est quasi-compact et quasi-séparé, ainsi qu'un faisceau abélien
$\mathcal{F}$ sur $T_\etale$ annulé par $n$, tels que
le morphisme de changement de base
$(f')^{-1}R^qg_*\mathcal{F} \to R^qh_*e^{-1}\mathcal{F}$
n'est pas un isomorphisme. Nous pouvons bien entendu remplacer $S'$
par un ouvert affine de $S'$, et nous le faisons; cela implique que $T$ est quasi-compact
et quasi-séparé. Par le lemme \ref{lemma-base-change-q-injective}, nous voyons que
$(f')^{-1}R^qg_*\mathcal{F} \to R^qh_*e^{-1}\mathcal{F}$
est injectif. Choisissons un point géométrique $\overline{x}$ de $X'$
et un élément $\xi$ de $(R^qh_*e^{-1}\mathcal{F})_{\overline{x}}$
qui n'appartient pas à l'image du morphisme
$((f')^{-1}R^qg_*\mathcal{F})_{\overline{x}} \to
(R^qh_*e^{-1}\mathcal{F})_{\overline{x}}$.

\medskip\noindent
Considérons un morphisme $\pi : T' \to T$, avec $T'$ quasi-compact
et quasi-séparé, et notons $\mathcal{F}' = \pi^{-1}\mathcal{F}$.
Notons $\pi' : Y' = Y \times_T T' \to Y$ le changement de base de $\pi$
et $e' : Y' \to T'$ le changement de base de $e$. Nous avons les diagrammes
$$
\vcenter{
\xymatrix{
X' \ar[d]_{f'} & Y \ar[l]^h \ar[d]^e & Y' \ar[l]^{\pi'} \ar[d]^{e'} \\
S' & T \ar[l]_g & T' \ar[l]_\pi
}
}
\quad\text{et}\quad
\vcenter{
\xymatrix{
X' \ar[d]_{f'} & & Y' \ar[ll]^{h' = h \circ \pi'} \ar[d]^{e'} \\
S' & & T' \ar[ll]_{g' = g \circ \pi}
}
}
$$
À l'aide des morphismes d'image inverse, nous obtenons un diagramme commutatif canonique
$$
\xymatrix{
(f')^{-1}R^qg_*\mathcal{F} \ar[r] \ar[d] &
(f')^{-1}R^qg'_*\mathcal{F}' \ar[d] \\
R^qh_*e^{-1}\mathcal{F} \ar[r] &
R^qh'_*(e')^{-1}\mathcal{F}'
}
$$
de faisceaux abéliens sur $X'$. Soit $P(T')$ la propriété suivante :
\begin{itemize}
\item L'image $\xi'$ de $\xi$ dans
$(R^qh'_*(e')^{-1}\mathcal{F}')_{\overline{x}}$ n'appartient
pas à l'image du morphisme
$((f')^{-1}R^qg'_*\mathcal{F}')_{\overline{x}} \to
(R^qh'_*(e')^{-1}\mathcal{F}')_{\overline{x}}$.
\end{itemize}
Nous affirmons que les hypothèses (1), (2) et (3) du
lemme \ref{lemma-formal-argument} sont satisfaites par $P$,
ce qui démontre notre lemme.

\medskip\noindent
La condition (1) du lemme \ref{lemma-formal-argument}
est satisfaite par $P$, car la topologie étale d'un schéma
et celle d'un épaississement de ce schéma sont les mêmes. Voir la
proposition \ref{proposition-topological-invariance}.

\medskip\noindent
Supposons que $I$ soit un ensemble ordonné filtrant et que $(T_i)$
soit un système inverse indexé par $I$ de schémas quasi-compacts et quasi-séparés
sur $T$, à morphismes de transition affines.
Posons $T' = \lim T_i$. Notons $\mathcal{F}'$ et $\mathcal{F}_i$
les images inverses de $\mathcal{F}$ sur $T'$ et sur $T_i$, respectivement. Considérons
les diagrammes
$$
\vcenter{
\xymatrix{
X' \ar[d]_{f'} & Y \ar[l]^h \ar[d]^e & Y_i \ar[l]^{\pi_i'} \ar[d]^{e_i} \\
S' & T \ar[l]_g & T_i \ar[l]_{\pi_i}
}
}
\quad\text{et}\quad
\vcenter{
\xymatrix{
X' \ar[d]_{f'} & & Y_i \ar[ll]^{h_i = h \circ \pi_i'} \ar[d]^{e_i} \\
S' & & T_i \ar[ll]_{g_i = g \circ \pi_i}
}
}
$$
comme au paragraphe précédent. Il est clair que $\mathcal{F}'$ sur
$T'$ est la limite inductive des images inverses des $\mathcal{F}_i$ sur $T'$
et que $(e')^{-1}\mathcal{F}'$ est la limite inductive des images inverses
des $e_i^{-1}\mathcal{F}_i$ sur $Y'$.
Par le lemme \ref{lemma-relative-colimit-general},
nous avons
$$
R^qh'_*(e')^{-1}\mathcal{F}' = \colim R^qh_{i, *}e_i^{-1}\mathcal{F}_i
\quad\text{et}\quad
(f')^{-1}R^qg'_*\mathcal{F}' = \colim (f')^{-1}R^qg_{i, *}\mathcal{F}_i
$$
Il s'ensuit que, si $P(T_i)$ est vraie pour tout $i$, alors
$P(T')$ est vraie. Ainsi, la condition (2) du lemme \ref{lemma-formal-argument}
est satisfaite par $P$.

\medskip\noindent
La plus intéressante est la condition (3) du lemme \ref{lemma-formal-argument}.
Supposons que $T'$ soit un schéma quasi-compact et quasi-séparé sur $T$
tel que $P(T')$ soit vraie.
Soit $Z \subset T'$ un sous-schéma fermé de complément quasi-compact
$V \subset T'$. Considérons le diagramme
$$
\xymatrix{
Y' \times_{T'} Z \ar[d]_{e_Z} \ar[r]_{i'} &
Y' \ar[d]_{e'} &
Y' \times_{T'} V \ar[l]^{j'} \ar[d]^{e_V} \\
Z \ar[r]^i &
T' &
V \ar[l]_j
}
$$
Choisissons un morphisme injectif $j^{-1}\mathcal{F}' \to \mathcal{J}$,
où $\mathcal{J}$ est un faisceau injectif de $\mathbf{Z}/n\mathbf{Z}$-modules
sur $V$. En regardant les fibres, nous voyons que le morphisme
$$
\mathcal{F}' \to \mathcal{G} = j_*\mathcal{J} \oplus i_*i^{-1}\mathcal{F}'
$$
est injectif. Ainsi, $\xi'$ a pour image un élément non nul de
\begin{align*}
& \Coker\left(
((f')^{-1}R^qg'_*\mathcal{G})_{\overline{x}}
\to
(R^qh'_*(e')^{-1}\mathcal{G})_{\overline{x}}
\right) = \\
&
\Coker\left(
((f')^{-1}R^qg'_*j_*\mathcal{J})_{\overline{x}}
\to
(R^qh'_*(e')^{-1}j_*\mathcal{J})_{\overline{x}}
\right) \oplus \\
& \Coker\left(
((f')^{-1}R^qg'_*i_*i^{-1}\mathcal{F}')_{\overline{x}}
\to
(R^qh'_*(e')^{-1}i_*i^{-1}\mathcal{F}')_{\overline{x}}
\right)
\end{align*}
par la partie (2) du lemme \ref{lemma-base-change-q-injective}.
Si $\xi'$ ne s'envoie pas sur zéro dans le second facteur,
nous utilisons
$$
(f')^{-1}R^qg'_*i_*i^{-1}\mathcal{F}' =
(f')^{-1}R^q(g' \circ i)_*i^{-1}\mathcal{F}'
$$
(car $Ri_* = i_*$ par la
proposition \ref{proposition-finite-higher-direct-image-zero}) et
$$
R^qh'_*(e')^{-1}i_*i^{-1}\mathcal{F}' =
R^qh'_*i'_*e_Z^{-1}i^{-1}\mathcal{F}' =
R^q(h' \circ i')_*e_Z^{-1}i^{-1}\mathcal{F}'
$$
(la première égalité par le
lemme \ref{lemma-finite-pushforward-commutes-with-base-change}
et la seconde parce que
$Ri'_* = i'_*$ par la
proposition \ref{proposition-finite-higher-direct-image-zero}),
pour voir que $P(Z)$ est vraie.
Enfin, supposons que $\xi'$ ne s'envoie pas sur zéro dans le premier facteur.
Nous avons
$$
(e')^{-1}j_*\mathcal{J} = j'_*e_V^{-1}\mathcal{J}
\quad\text{et}\quad
R^aj'_*e_V^{-1}\mathcal{J} = 0, \quad a = 1, \ldots, q - 1
$$
par $BC(f, n, q - 1)$ appliquée au diagramme
$$
\xymatrix{
X' \ar[d]_{f'} & Y' \ar[l]^{h'} \ar[d]_{e'} & Y' \times_{T'} V \ar[l]^{j'} \ar[d]^{e_V} \\
S' & T' \ar[l]_{g'} & V \ar[l]_j
}
$$
et au fait que $\mathcal{J}$ est injectif.
Par la suite spectrale de Leray relative pour $h' \circ j'$
(lemme \ref{sites-cohomology-lemma-relative-Leray} de Cohomologie sur les sites),
nous en déduisons que
$$
R^qh'_*(e')^{-1}j_*\mathcal{J} =
R^qh'_*j'_*e_V^{-1}\mathcal{J}
\longrightarrow
R^q(h' \circ j')_* e_V^{-1}\mathcal{J}
$$
est injectif. Ainsi, $\xi'$ a pour image un élément non nul de
$(R^q(h' \circ j')_* e_V^{-1}\mathcal{J})_{\overline{x}}$.
En appliquant la partie (3) du lemme \ref{lemma-base-change-q-injective}
à l'injection $j^{-1}\mathcal{F}' \to \mathcal{J}$,
nous concluons que $P(V)$ est vraie.
\end{proof}

\begin{lemma}
\label{lemma-base-change-does-not-hold}
Avec $f : X \to S$ et $n$ comme dans la remarque \ref{remark-base-change-holds},
supposons que, pour un certain $q \geq 1$,
$BC(f, n, q - 1)$ soit vraie, mais que $BC(f, n, q)$ ne le soit pas.
Il existe alors un diagramme commutatif
$$
\xymatrix{
X \ar[d]_f & X' \ar[d] \ar[l] & Y \ar[l]^h \ar[d] \\
S & S' \ar[l] & \Spec(K) \ar[l]
}
$$
dont les deux carrés sont cartésiens, où :
\begin{enumerate}
\item $S'$ est affine, intègre et normal, à corps des fonctions
algébriquement clos;
\item $K$ est algébriquement clos et $\Spec(K) \to S'$
est dominant (autrement dit, $K$ est une extension du
corps des fonctions de $S'$).
\end{enumerate}
De plus, il existe un entier $d | n$
tel que $R^qh_*(\mathbf{Z}/d\mathbf{Z})$ soit non nul.
\end{lemma}

\noindent
Réciproquement, la non-annulation de $R^qh_*(\mathbf{Z}/d\mathbf{Z})$
dans le lemme implique que $BC(f, n, q)$ n'est pas vraie, puisque le
lemme \ref{lemma-Rf-star-zero-normal-with-alg-closed-function-field}
montre que $R^q(\Spec(K) \to S')_*\mathbf{Z}/d\mathbf{Z} = 0$.

\begin{proof}
Choisissons d'abord un diagramme et $\mathcal{F}$ comme dans le
lemme \ref{lemma-base-change-does-not-hold-pre}.
Nous pouvons supposer $S'$ affine, et nous le faisons (c'est évident, mais
voir la démonstration du lemme en cas de doute).
Par le lemme \ref{lemma-base-change-q-integral-top},
nous pouvons supposer $K$ algébriquement clos.
Alors $\mathcal{F}$ correspond à un $\mathbf{Z}/n\mathbf{Z}$-module.
Un tel module est une somme directe de copies de $\mathbf{Z}/d\mathbf{Z}$
pour divers $d | n$; nous pouvons donc supposer $\mathcal{F}$
constant, de valeur $\mathbf{Z}/d\mathbf{Z}$.
Par le lemme \ref{lemma-base-change-q-integral-bottom},
nous pouvons remplacer $S'$ par sa normalisation
dans $\Spec(K)$, ce qui achève la démonstration.
\end{proof}










\section{Changement de base lisse}
\label{section-smooth-base-change}

\noindent
Dans cette section, nous démontrons le théorème de changement de base lisse.

\begin{lemma}
\label{lemma-smooth-base-change-fields}
Soit $K/k$ une extension de corps. Soit $X$ une courbe affine lisse
sur $k$, munie d'un point rationnel $x \in X(k)$. Soit $\mathcal{F}$ un faisceau
abélien sur $\Spec(K)$ annulé par un entier $n$ inversible dans $k$.
Soit $q > 0$ et soit
$$
\xi \in H^q(X_K, (X_K \to \Spec(K))^{-1}\mathcal{F})
$$
Il existe :
\begin{enumerate}
\item des extensions finies $K'/K$ et $k'/k$, avec $k' \subset K'$;
\item un revêtement galoisien fini étale $Z \to X_{k'}$, de groupe $G$.
\end{enumerate}
L'ordre de $G$ divise une puissance de $n$, le morphisme
$Z \to X_{k'}$ est scindé au-dessus de $x_{k'}$, et
$\xi$ s'annule dans $H^q(Z_{K'}, (Z_{K'} \to \Spec(K))^{-1}\mathcal{F})$.
\end{lemma}

\begin{proof}
Pour $q > 1$, nous savons que $\xi$ s'annule dans
$H^q(X_{\overline{K}}, (X_{\overline{K}} \to \Spec(K))^{-1}\mathcal{F})$
(théorème \ref{theorem-vanishing-affine-curves}).
Par le lemme \ref{lemma-directed-colimit-cohomology}, nous voyons que
cela signifie qu'il existe une extension finie $K'/K$ telle que
$\xi$ s'annule dans $H^q(X_{K'}, (X_{K'} \to \Spec(K))^{-1}\mathcal{F})$.
Dans ce cas, nous pouvons donc prendre $k' = k$ et $Z = X$.

\medskip\noindent
Supposons $q = 1$. Rappelons que $\mathcal{F}$ correspond à un
module discret $M$ muni d'une action continue de $\text{Gal}_K$; voir le
lemme \ref{lemma-equivalence-abelian-sheaves-point}.
Puisque $M$ est de $n$-torsion, il est la réunion de ses sous-groupes finis
stables par $\text{Gal}_K$. Nous nous ramenons donc
au cas où $M$ est un groupe abélien fini annulé par $n$;
voir le lemme \ref{lemma-colimit}. Après avoir remplacé $K$ par une extension finie,
nous pouvons supposer que l'action de $\text{Gal}_K$ sur $M$ est triviale.
Nous pouvons donc supposer que $\mathcal{F} = \underline{M}$ est le faisceau
constant de valeur le groupe abélien fini $M$, annulé par $n$.

\medskip\noindent
Nous pouvons écrire $M$ comme une somme directe de groupes cycliques.
Deux revêtements galoisiens finis étales quelconques, dont les groupes de Galois
sont d'ordre inversible dans $k$, sont dominés par un troisième dont le
groupe de Galois est d'ordre inversible dans $k$
(section \ref{pione-section-finite-etale-under-galois} de Fundamental Groups).
Il suffit donc de démontrer le lemme lorsque
$M = \mathbf{Z}/d\mathbf{Z}$, où $d | n$.

\medskip\noindent
Supposons $M = \mathbf{Z}/d\mathbf{Z}$, où $d | n$.
Dans ce cas, $\overline{\xi} = \xi|_{X_{\overline{K}}}$ est un élément de
$$
H^1(X_{\overline{k}}, \mathbf{Z}/d\mathbf{Z}) =
H^1(X_{\overline{K}}, \mathbf{Z}/d\mathbf{Z})
$$
Voir le théorème \ref{theorem-vanishing-affine-curves}.
Ce groupe classifie les $\mathbf{Z}/d\mathbf{Z}$-torseurs; voir le
lemme \ref{sites-cohomology-lemma-torsors-h1} de Cohomologie sur les sites.
Le torseur correspondant à $\overline{\xi}$ (considéré comme un faisceau sur
$X_{\overline{k}, \etale}$) donne à son tour un morphisme fini étale
$T \to X_{\overline{k}}$ muni d'une action de $\mathbf{Z}/d\mathbf{Z}$,
transitive sur la fibre de $T$ au-dessus de $x_{\overline{k}}$; voir le
lemme \ref{lemma-characterize-finite-locally-constant}.
Choisissons une composante connexe $T' \subset T$ (si $\overline{\xi}$ est
d'ordre $d$, alors $T$ est déjà connexe).
Alors $T' \to X_{\overline{k}}$ est un revêtement galoisien
fini étale dont le groupe de Galois est un sous-groupe $G \subset \mathbf{Z}/d\mathbf{Z}$
(un petit détail est omis). En outre, l'élément $\overline{\xi}$ s'envoie sur zéro
par le morphisme $H^1(X_{\overline{k}}, \mathbf{Z}/d\mathbf{Z}) \to
H^1(T', \mathbf{Z}/d\mathbf{Z})$, car c'est l'une
des propriétés qui définissent $T$.

\medskip\noindent
Nous utilisons ensuite un argument de limite pour choisir une extension finie $k'/k$
contenue dans $\overline{k}$ telle que $T' \to X_{\overline{k}}$
descende en un revêtement galoisien fini étale $Z \to X_{k'}$, de groupe $G$.
Voir les lemmes de Limits \ref{limits-lemma-descend-finite-presentation},
\ref{limits-lemma-descend-finite-finite-presentation} et
\ref{limits-lemma-descend-etale}.
Après avoir agrandi $k'$, nous pouvons supposer que $Z$ est scindé au-dessus de $x_{k'}$.
Par construction, l'image de $\xi$ dans
$H^1(Z_{\overline{K}}, \mathbf{Z}/d\mathbf{Z})$ est nulle.
Ainsi, par le lemme \ref{lemma-directed-colimit-cohomology},
nous pouvons trouver une sous-extension finie $\overline{K}/K'/K$
contenant $k'$ telle que $\xi$ s'annule dans $H^1(Z_{K'}, \mathbf{Z}/d\mathbf{Z})$,
ce qui achève la démonstration.
\end{proof}

\begin{theorem}[Changement de base lisse]
\label{theorem-smooth-base-change}
Considérons un diagramme cartésien de schémas
$$
\xymatrix{
X \ar[d]_f & Y \ar[l]^h \ar[d]^e \\
S & T \ar[l]_g
}
$$
où $f$ est lisse et $g$ quasi-compact et quasi-séparé. Alors
$$
f^{-1}R^qg_*\mathcal{F} = R^qh_*e^{-1}\mathcal{F}
$$
pour tout $q$ et tout faisceau abélien $\mathcal{F}$
sur $T_\etale$ dont toutes les fibres aux points géométriques sont de torsion
d'ordres inversibles sur $S$.
\end{theorem}

\begin{proof}[Première démonstration du changement de base lisse]
Cette démonstration est très longue, mais plus directe (elle utilise moins de théorie générale)
que la seconde démonstration donnée ci-dessous.

\medskip\noindent
Le théorème est local sur $X_\etale$. Plus précisément, supposons donné
$U \to X$ étale tel que $U \to S$ se factorise en $U \to V \to S$,
avec $V \to S$ étale. Nous pouvons alors considérer le carré cartésien
$$
\xymatrix{
U \ar[d]_{f'} & U \times_X Y \ar[l]^{h'} \ar[d]^{e'} \\
V & V \times_S T \ar[l]_{g'}
}
$$
et, en posant $\mathcal{F}' = \mathcal{F}|_{V \times_S T}$,
nous avons $f^{-1}R^qg_*\mathcal{F}|_U = (f')^{-1}R^qg'_*\mathcal{F}'$
et $R^qh_*e^{-1}\mathcal{F}|_U = R^qh'_*(e')^{-1}\mathcal{F}'$
(cela découle de la compatibilité de la localisation avec les morphismes de sites; voir le
lemme \ref{sites-lemma-localize-morphism-strong} de Sites,
ainsi que le
lemme de Cohomologie sur les sites
\ref{sites-cohomology-lemma-restrict-direct-image-open}).
Il suffit donc de produire un recouvrement étale de $X$ par des
$U \to X$ et des factorisations $U \to V \to S$
comme ci-dessus, tels que le théorème vaille pour le diagramme formé de
$f'$, $h'$, $g'$ et $e'$.

\medskip\noindent
Par la structure locale des morphismes lisses, voir le
lemme \ref{morphisms-lemma-smooth-etale-over-affine-space} de Morphismes,
nous pouvons supposer $X$ et $S$ affines et que $X \to S$
se factorise par un morphisme étale $X \to \mathbf{A}^d_S$.
Si nous avons une tour de diagrammes cartésiens
$$
\xymatrix{
W \ar[d]_i & Z \ar[l]^j \ar[d]^k \\
X \ar[d]_f & Y \ar[l]^h \ar[d]^e \\
S & T \ar[l]_g
}
$$
et si le théorème vaut pour les carrés inférieur et supérieur, alors
il vaut pour le rectangle extérieur; cela est formel.
En écrivant $X \to S$ comme la composée
$$
X \to \mathbf{A}^{d - 1}_S \to \mathbf{A}^{d - 2}_S \to \ldots \to
\mathbf{A}^1_S \to S
$$
nous concluons qu'il suffit de démontrer le théorème lorsque $X$ et $S$
sont affines et que $X \to S$ est de dimension relative $1$.

\medskip\noindent
Pour tout $n \geq 1$ inversible sur $S$, soit $\mathcal{F}[n]$
le sous-faisceau des sections de $\mathcal{F}$ annulées par $n$. Alors
$\mathcal{F} = \colim \mathcal{F}[n]$ par notre hypothèse sur
les fibres de $\mathcal{F}$. Les foncteurs $e^{-1}$ et $f^{-1}$
commutent aux limites inductives car ce sont des adjoints à gauche. Les foncteurs
$R^qh_*$ et $R^qg_*$ commutent aux limites inductives filtrantes par le
lemme \ref{lemma-relative-colimit}.
Il suffit donc de démontrer le théorème pour $\mathcal{F}[n]$.
Désormais, nous fixons un entier $n$, travaillons avec des
faisceaux de $\mathbf{Z}/n\mathbf{Z}$-modules et
supposons que $S$ est un schéma sur $\Spec(\mathbf{Z}[1/n])$.

\medskip\noindent
Nous nous ramenons ensuite au cas où $T$ est affine. Puisque $g$ est quasi-compact
et quasi-séparé et que $S$ est affine, le schéma $T$ est quasi-compact et
quasi-séparé. Nous pouvons donc utiliser le principe de récurrence du
lemme \ref{coherent-lemma-induction-principle} de Cohomologie des schémas.
Par conséquent, il suffit de montrer que, si $T = W \cup W'$
est un recouvrement ouvert et si le théorème vaut pour les carrés
$$
\xymatrix{
X \ar[d] & e^{-1}(W) \ar[l]^i \ar[d] \\
S & W \ar[l]_a
}
\quad
\xymatrix{
X \ar[d] & e^{-1}(W') \ar[l]^j \ar[d] \\
S & W' \ar[l]_b
}
\quad
\xymatrix{
X \ar[d] & e^{-1}(W \cap W') \ar[l]^-k \ar[d] \\
S & W \cap W' \ar[l]_c
}
$$
alors le théorème vaut pour le diagramme initial. Pour le voir, considérons le
diagramme
$$
\xymatrix{
f^{-1}R^{q - 1}c_*\mathcal{F}|_{W \cap W'} \ar[d]^{\cong} \ar[r] &
f^{-1}R^qg_*\mathcal{F} \ar[d] \ar[r] &
f^{-1}R^qa_*\mathcal{F}|_W \oplus
f^{-1}R^qb_*\mathcal{F}|_{W'} \ar[d]_{\cong} \\
R^{q - 1}k_*e^{-1}\mathcal{F}|_{e^{-1}(W \cap W')} \ar[r] &
R^qh_*e^{-1}\mathcal{F} \ar[r] &
R^qi_*e^{-1}\mathcal{F}|_{e^{-1}(W)} \oplus
R^qj_*e^{-1}\mathcal{F}|_{e^{-1}(W')}
}
$$
dont les lignes sont les suites exactes longues du
lemme \ref{lemma-relative-mayer-vietoris}.
Le lemme des cinq donne donc la conclusion voulue.

\medskip\noindent
En résumé, nous pouvons supposer $S$, $X$, $T$ et $Y$ affines,
$\mathcal{F}$ de $n$-torsion, $X \to S$ lisse de dimension relative $1$,
et $S$ schéma sur $\mathbf{Z}[1/n]$.
Nous allons démontrer le théorème par récurrence sur $q$. Le cas initial $q = 0$
est traité par le lemme \ref{lemma-fppf-reduced-fibres-base-change-f-star}.
Supposons $q > 0$ et le théorème vrai en tout degré inférieur.
Choisissons une suite exacte courte
$0 \to \mathcal{F} \to \mathcal{I} \to \mathcal{Q} \to 0$
où $\mathcal{I}$ est un faisceau injectif de $\mathbf{Z}/n\mathbf{Z}$-modules.
Considérons le diagramme induit
$$
\xymatrix{
f^{-1}R^{q - 1}g_*\mathcal{I} \ar[d]_{\cong} \ar[r] &
f^{-1}R^{q - 1}g_*\mathcal{Q} \ar[d]_{\cong} \ar[r] &
f^{-1}R^qg_*\mathcal{F} \ar[d] \ar[r] &
0 \ar[d] \\
R^{q - 1}h_*e^{-1}\mathcal{I} \ar[r] &
R^{q - 1}h_*e^{-1}\mathcal{Q} \ar[r] &
R^qh_*e^{-1}\mathcal{F} \ar[r] &
R^qh_*e^{-1}\mathcal{I}
}
$$
à lignes exactes. Le zéro figure dans le coin supérieur droit
car $\mathcal{I}$ est injectif. Les deux flèches verticales de gauche sont
des isomorphismes par l'hypothèse de récurrence. Il suffit donc de démontrer
que $R^qh_*e^{-1}\mathcal{I} = 0$.

\medskip\noindent
Écrivons $S = \Spec(A)$ et $T = \Spec(B)$, et supposons que le morphisme
$T \to S$ soit donné par le morphisme d'anneaux $A \to B$. Nous pouvons écrire
$A \to B = \colim_{i \in I} (A_i \to B_i)$ comme une limite inductive
filtrante de morphismes d'anneaux de type fini sur $\mathbf{Z}[1/n]$
(voir le lemme \ref{algebra-lemma-limit-no-condition} d'Algèbre).
Pour $i \in I$, posons $S_i = \Spec(A_i)$ et $T_i = \Spec(B_i)$.
Pour $i$ assez grand, nous pouvons trouver un morphisme lisse $X_i \to S_i$
de dimension relative $1$ tel que $X = X_i \times_{S_i} S$; voir les
lemmes de Limits \ref{limits-lemma-descend-finite-presentation},
\ref{limits-lemma-descend-smooth} et
\ref{limits-lemma-descend-dimension-d}. Posons $Y_i = X_i \times_{S_i} T_i$
pour obtenir les carrés
$$
\xymatrix{
X_i \ar[d]_{f_i} & Y_i \ar[l]^{h_i} \ar[d]^{e_i} \\
S_i & T_i \ar[l]_{g_i}
}
$$
Remarquons que $\mathcal{I}_i = (T \to T_i)_*\mathcal{I}$
est un faisceau injectif de $\mathbf{Z}/n\mathbf{Z}$-modules sur
$T_i$; voir le lemme de Cohomologie sur les sites
\ref{sites-cohomology-lemma-pushforward-injective-flat}.
Nous avons $\mathcal{I} = \colim (T \to T_i)^{-1}\mathcal{I}_i$
par le lemme \ref{lemma-linus-hamann}. En prenant l'image inverse par $e$, nous obtenons
$e^{-1}\mathcal{I} = \colim (Y \to Y_i)^{-1}e_i^{-1}\mathcal{I}_i$.
Par le lemme \ref{lemma-relative-colimit-general}, appliqué au système
de morphismes $Y_i \to X_i$ de limite $Y \to X$, nous avons
$$
R^qh_*e^{-1}\mathcal{I} =
\colim (X \to X_i)^{-1} R^qh_{i, *} e_i^{-1}\mathcal{I}_i
$$
Cela nous ramène au cas où $T$ et $S$ sont affines de type
fini sur $\mathbf{Z}[1/n]$.

\medskip\noindent
En résumé, nous avons un entier $q \geq 1$ tel que le théorème vaille
en degrés $< q$, les schémas $S$ et $T$ sont affines de type fini
sur $\mathbf{Z}[1/n]$, le morphisme
$X \to S$ est lisse de dimension relative $1$, avec $X$ affine, et
$\mathcal{I}$ est un faisceau injectif de $\mathbf{Z}/n\mathbf{Z}$-modules;
nous devons montrer que $R^qh_*e^{-1}\mathcal{I} = 0$.
Nous allons le faire par récurrence sur $\dim(T)$.

\medskip\noindent
Le cas initial est $T = \emptyset$, c'est-à-dire $\dim(T) < 0$.
Si l'on préfère, on peut prendre comme cas initial
le cas $\dim(T) = 0$. Dans ce cas, $T \to S$ est fini
(en fait, même $T \to \Spec(\mathbf{Z}[1/n])$ est fini puisque la
cible est jacobsonienne; les détails sont omis), donc
$h$ est lui aussi fini et ses
images directes supérieures sont nulles (voir les références ci-dessous).

\medskip\noindent
Supposons $\dim(T) = d \geq 0$ et le résultat connu dans toutes les situations où $T$
est de dimension inférieure. Choisissons $U$ affine et étale sur $X$, ainsi qu'une section $\xi$
de $R^qh_*e^{-1}\mathcal{I}$ sur $U$. Nous devons montrer
que $\xi$ est nulle. Nous pouvons bien entendu remplacer $X$ par $U$
(et, en conséquence, $Y$ par $U \times_X Y$)
et supposer $\xi \in H^0(X, R^qh_*e^{-1}\mathcal{I})$.
De plus, puisque $R^qh_*e^{-1}\mathcal{I}$ est un faisceau,
il suffit de démontrer que $\xi$ est nulle localement sur $X$.
Nous pouvons donc remplacer $X$ par les membres d'un recouvrement étale.
En particulier, à l'aide du lemme \ref{lemma-higher-direct-images},
nous pouvons supposer que $\xi$ est l'image d'un élément
$\tilde \xi \in H^q(Y, e^{-1}\mathcal{I})$.
En termes de $\tilde \xi$, notre tâche est de montrer que
$\tilde \xi$ s'annule dans $H^q(U_i \times_X Y, e^{-1}\mathcal{I})$
pour un certain recouvrement étale $\{U_i \to X\}$.

\medskip\noindent
Par le lemme \ref{more-morphisms-lemma-cover-smooth-by-special} de Compléments sur les morphismes,
nous pouvons supposer que $X \to S$ se factorise en $X \to V \to S$,
où $V \to S$ est étale et $X \to V$ est un morphisme lisse
de schémas affines, de dimension relative $1$, muni d'une section et
à fibres géométriquement connexes. Remarquons que
$\dim(V \times_S T) \leq \dim(T) = d$
par exemple par le lemme de Compléments d'algèbre
\ref{more-algebra-lemma-dimension-etale-extension}.
Nous pouvons donc remplacer $S$ par $V$ et $T$ par $V \times_S T$
(exactement comme dans la discussion du premier paragraphe de la démonstration).
Ainsi, nous pouvons supposer $X \to S$ lisse de dimension relative $1$,
à fibres géométriquement connexes et muni d'une section $\sigma : S \to X$.

\medskip\noindent
Soit $\pi : T' \to T$ un morphisme fini surjectif.
Nous utiliserons ci-dessous que $\dim(T') \leq \dim(T) = d$; voir le
lemme \ref{algebra-lemma-integral-dim-up} d'Algèbre.
Choisissons un morphisme injectif $\pi^{-1}\mathcal{I} \to \mathcal{I}'$
vers un faisceau injectif de $\mathbf{Z}/n\mathbf{Z}$-modules.
Alors $\mathcal{I} \to \pi_*\mathcal{I}'$ est injectif
et possède donc une rétraction (car $\mathcal{I}$ est un faisceau injectif
de $\mathbf{Z}/n\mathbf{Z}$-modules).
Notons $\pi' : Y' = Y \times_T T' \to Y$ le changement de base de $\pi$
et $e' : Y' \to T'$ le changement de base de $e$. Nous avons le diagramme
$$
\xymatrix{
X \ar[d]_f & Y \ar[l]^h \ar[d]^e & Y' \ar[l]^{\pi'} \ar[d]^{e'} \\
S & T \ar[l]_g & T' \ar[l]_\pi
}
$$
Par la proposition \ref{proposition-finite-higher-direct-image-zero} et le
lemme \ref{lemma-finite-pushforward-commutes-with-base-change}, nous avons
$R\pi'_*(e')^{-1}\mathcal{I}' = e^{-1}\pi_*\mathcal{I}'$.
Ainsi, par la suite spectrale de Leray
(lemme \ref{sites-cohomology-lemma-Leray} de Cohomologie sur les sites),
nous avons
$$
H^q(Y', (e')^{-1}\mathcal{I}') =
H^q(Y, e^{-1}\pi_*\mathcal{I}') \supset H^q(Y, e^{-1}\mathcal{I})
$$
et cela reste vrai après changement de base par tout $U \to X$ étale.
Nous pouvons donc remplacer $T$ par $T'$, $\mathcal{I}$ par $\mathcal{I}'$
et $\tilde \xi$ par son image dans $H^q(Y', (e')^{-1}\mathcal{I}')$.

\medskip\noindent
Supposons donnée une factorisation $T \to S' \to S$, où
$\pi : S' \to S$ est fini.
En posant $X' = S' \times_S X$, nous pouvons considérer le diagramme induit
$$
\xymatrix{
X \ar[d]_f & X' \ar[l]^{\pi'} \ar[d]^{f'} & Y \ar[l]^{h'} \ar[d]^e \\
S & S' \ar[l]_\pi & T \ar[l]_g
}
$$
Puisque les images directes supérieures de $\pi'$ sont nulles, nous voyons que
$R^qh_*e^{-1}\mathcal{I} = \pi'_*R^qh'_*e^{-1}\mathcal{I}$
par la suite spectrale de Leray. Par conséquent,
$H^0(X, R^qh_*e^{-1}\mathcal{I}) = H^0(X', R^qh'_*e^{-1}\mathcal{I})$.
Ainsi, $\xi$ est nulle si et seulement si la section correspondante de
$R^qh'_*e^{-1}\mathcal{I}$ est nulle\footnote{Cette
étape peut aussi se voir autrement. En effet, nous devons montrer
qu'il existe un recouvrement étale $\{U_i \to X\}$ tel que
$\tilde \xi$ s'annule dans $H^q(U_i \times_X Y, e^{-1}\mathcal{I})$.
Cependant, si nous montrons qu'il existe un recouvrement étale
$\{U'_j \to X'\}$ tel que $\tilde \xi$ s'annule dans
$H^q(U'_j \times_{X'} Y, e^{-1}\mathcal{I})$, alors, par la propriété (B)
pour $X' \to X$ (lemme \ref{lemma-finite-B}), il existe un recouvrement étale
$\{U_i \to X\}$ tel que $U_i \times_X X'$ soit une réunion disjointe
de schémas sur $X'$, chacun se factorisant par un $U'_j$.
Nous voyons donc que $\tilde \xi$ s'annule dans $H^q(U_i \times_X Y, e^{-1}\mathcal{I})$,
comme voulu.}.
Nous pouvons donc remplacer $S$ par $S'$ et $X$ par $X'$.
Remarquons que le changement de base de $\sigma : S \to X$ est $\sigma' : S' \to X'$;
ainsi, après ce remplacement, il reste vrai que $X \to S$
possède une section $\sigma$ et des fibres géométriquement connexes.

\medskip\noindent
Nous utiliserons le fait que $S$ et $T$ sont des schémas de Nagata; voir la
proposition \ref{algebra-proposition-ubiquity-nagata} d'Algèbre.
Cela garantira
que diverses normalisations sont finies; voir les
lemmes de Morphismes \ref{morphisms-lemma-nagata-normalization-finite} et
\ref{morphisms-lemma-nagata-normalization}.
En particulier, nous pouvons d'abord remplacer $T$ par sa normalisation,
puis remplacer $S$ par la normalisation de $S$ dans $T$.
Alors $T \to S$ est une réunion disjointe de morphismes
dominants de schémas intègres normaux; voir le
lemme \ref{morphisms-lemma-normal-normalization} de Morphismes.
Nous pouvons clairement raisonner composante connexe
par composante connexe; nous pouvons donc supposer que $T \to S$ est un
morphisme dominant de schémas intègres normaux.

\medskip\noindent
Soient $s \in S$ et $t \in T$ les points génériques. Par le
lemme \ref{lemma-smooth-base-change-fields}, il existe des extensions finies
de corps $K/\kappa(t)$ et $k/\kappa(s)$ telles que $k$ soit contenu dans $K$,
ainsi qu'un revêtement galoisien fini étale $Z \to X_k$ de groupe de Galois $G$,
d'ordre divisant une puissance de $n$, scindé au-dessus de $\sigma(\Spec(k))$,
tel que $\tilde \xi$ s'envoie sur zéro dans $H^q(Z_K, e^{-1}\mathcal{I}|_{Z_K})$.
Soit $T' \to T$ la normalisation de $T$ dans $\Spec(K)$,
et soit $S' \to S$ la normalisation de $S$ dans $\Spec(k)$.
Nous obtenons alors un diagramme commutatif
$$
\xymatrix{
S' \ar[d] & T' \ar[l] \ar[d] \\
S & T \ar[l]
}
$$
dont les flèches verticales sont finies. Par les arguments donnés ci-dessus, nous
pouvons remplacer $S$ et $T$ par $S'$ et $T'$, et nous le faisons (et, en conséquence,
$X$ par $X \times_S S'$ et $Y$ par $Y \times_T T'$). Après ce remplacement,
nous concluons que nous avons un revêtement galoisien fini étale
$Z \to X_s$ de la fibre générique de $X \to S$,
de groupe de Galois $G$ d'ordre divisant une puissance de $n$,
scindé au-dessus de $\sigma(s)$, tel que $\tilde \xi$ s'envoie sur zéro dans
$H^q(Z_t, (Z_t \to Y)^{-1}e^{-1}\mathcal{I})$.
Ici, $Z_t = Z \times_S t = Z \times_s t = Z \times_{X_s} Y_t$.
Comme $n$ est inversible sur $S$,
le lemme \ref{pione-lemma-extend-covering} de Fundamental Groups
permet de trouver un morphisme fini étale $U \to X$ dont la restriction à $X_s$
est $Z$.

\medskip\noindent
À ce point, nous remplaçons $X$ par $U$ et $Y$ par $U \times_X Y$.
Après ce remplacement, il peut arriver que
les fibres de $X \to S$ ne soient plus géométriquement connexes
(il subsiste une section, mais nous ne l'utiliserons pas);
en revanche, ce remplacement assure que $\tilde \xi$ s'envoie sur zéro
dans $H^q(Y_t, e^{-1}\mathcal{I})$, c'est-à-dire que $\tilde \xi$ se restreint à
zéro sur la fibre générique de $Y \to T$.

\medskip\noindent
Rappelons que $t$ est le spectre du corps des fonctions de $T$, c'est-à-dire
que, comme schéma, $t$ est la limite des sous-schémas ouverts affines non vides de $T$.
Par le lemme \ref{lemma-directed-colimit-cohomology}, nous concluons qu'il existe
un sous-schéma ouvert non vide $V \subset T$ tel que $\tilde \xi$ s'envoie sur zéro dans
$H^q(Y \times_T V, e^{-1}\mathcal{I}|_{Y \times_T V})$.

\medskip\noindent
Notons $Z = T \setminus V$. Considérons le diagramme
$$
\xymatrix{
Y \times_T Z \ar[d]_{e_Z} \ar[r]_{i'} &
Y \ar[d]_e &
Y \times_T V \ar[l]^{j'} \ar[d]^{e_V} \\
Z \ar[r]^i &
T &
V \ar[l]_j
}
$$
Choisissons une injection $i^{-1}\mathcal{I} \to \mathcal{I}'$
dans un faisceau injectif de $\mathbf{Z}/n\mathbf{Z}$-modules sur $Z$.
En regardant les fibres, nous voyons que le morphisme
$$
\mathcal{I} \to j_*\mathcal{I}|_V \oplus i_*\mathcal{I}'
$$
est injectif et possède donc une rétraction, puisque $\mathcal{I}$ est un faisceau injectif
de $\mathbf{Z}/n\mathbf{Z}$-modules. Il suffit donc de montrer que
$\tilde \xi$ s'envoie sur zéro dans
$$
H^q(Y, e^{-1}j_*\mathcal{I}|_V) \oplus
H^q(Y, e^{-1}i_*\mathcal{I}')
$$
au moins après avoir remplacé $X$ par les membres d'un recouvrement étale.
Remarquons que
$$
e^{-1}j_*\mathcal{I}|_V = j'_*e_V^{-1}\mathcal{I}|_V,\quad
e^{-1}i_*\mathcal{I}' = i'_*e_Z^{-1}\mathcal{I}'
$$
Par l'hypothèse de récurrence sur $q$, nous voyons que
$$
R^aj'_*e_V^{-1}\mathcal{I}|_V = 0, \quad a = 1, \ldots, q - 1
$$
Par la suite spectrale de Leray pour $j'$ et l'annulation
ci-dessus, il s'ensuit que
$$
H^q(Y, j'_*(e_V^{-1}\mathcal{I}|_V))
\longrightarrow
H^q(Y \times_T V, e_V^{-1}(\mathcal{I}|_V)) =
H^q(Y \times_T V, e^{-1}\mathcal{I}|_{Y \times_T V})
$$
est injectif. Ainsi, l'image de $\tilde \xi$ dans le premier facteur
ci-dessus s'annule, car nous savons que $\tilde \xi$ s'annule
dans $H^q(Y \times_T V, e^{-1}\mathcal{I}|_{Y \times_T V})$.
Puisque $\dim(Z) < \dim(T) = d$, l'hypothèse de récurrence montre que l'image de $\tilde \xi$
dans le second facteur
$$
H^q(Y, e^{-1}i_*\mathcal{I}') =
H^q(Y, i'_*e_Z^{-1}\mathcal{I}') =
H^q(Y \times_T Z, e_Z^{-1}\mathcal{I}')
$$
s'annule après avoir remplacé $X$ par les membres d'un recouvrement étale convenable.
Cela achève la démonstration du théorème de changement de base lisse.
\end{proof}

\begin{proof}[Seconde démonstration du changement de base lisse]
Cette démonstration est la même que la première, plus longue;
elle n'est plus courte que parce que nous avons isolé dans plusieurs
lemmes les arguments utilisés.

\medskip\noindent
Le cas $q = 0$ est le
lemme \ref{lemma-fppf-reduced-fibres-base-change-f-star}.
Nous pouvons donc supposer $q > 0$ et le résultat vrai
en tout degré inférieur.

\medskip\noindent
Pour tout $n \geq 1$ inversible sur $S$, soit $\mathcal{F}[n]$
le sous-faisceau des sections de $\mathcal{F}$ annulées par $n$. Alors
$\mathcal{F} = \colim \mathcal{F}[n]$ par notre hypothèse sur
les fibres de $\mathcal{F}$. Les foncteurs $e^{-1}$ et $f^{-1}$
commutent aux limites inductives car ce sont des adjoints à gauche. Les foncteurs
$R^qh_*$ et $R^qg_*$ commutent aux limites inductives filtrantes par le
lemme \ref{lemma-relative-colimit}.
Il suffit donc de démontrer le théorème pour $\mathcal{F}[n]$.
Désormais, nous fixons un entier $n$ inversible sur $S$ et travaillons avec des
faisceaux de $\mathbf{Z}/n\mathbf{Z}$-modules.

\medskip\noindent
Par le lemme \ref{lemma-base-change-local}, la question est locale pour la topologie étale
sur $X$ et $S$. Par la structure locale des morphismes lisses, voir le
lemme \ref{morphisms-lemma-smooth-etale-over-affine-space} de Morphismes,
nous pouvons supposer $X$ et $S$ affines et que $X \to S$
se factorise par un morphisme étale $X \to \mathbf{A}^d_S$.
En écrivant $X \to S$ comme la composée
$$
X \to \mathbf{A}^{d - 1}_S \to \mathbf{A}^{d - 2}_S \to \ldots \to
\mathbf{A}^1_S \to S
$$
nous concluons du lemme \ref{lemma-base-change-compose}
qu'il suffit de démontrer le théorème lorsque $X$ et $S$
sont affines et que $X \to S$ est de dimension relative $1$.

\medskip\noindent
Par le lemme \ref{lemma-base-change-does-not-hold}, il suffit de montrer que
$R^qh_*\mathbf{Z}/d\mathbf{Z} = 0$ pour tout $d | n$, chaque fois que nous avons un
diagramme cartésien
$$
\xymatrix{
X \ar[d] & Y \ar[d] \ar[l]^h \\
S & \Spec(K) \ar[l]
}
$$
où $X \to S$ est affine et lisse de dimension relative $1$,
$S$ est le spectre d'un anneau intègre
normal $A$, de corps des fractions algébriquement clos $L$,
et $K/L$ est une extension de corps algébriquement clos.

\medskip\noindent
Rappelons que $R^qh_*\mathbf{Z}/d\mathbf{Z}$ est le faisceau associé
au préfaisceau
$$
U \longmapsto
H^q(U \times_X Y, \mathbf{Z}/d\mathbf{Z}) =
H^q(U \times_S \Spec(K), \mathbf{Z}/d\mathbf{Z})
$$
sur $X_\etale$ (lemme \ref{lemma-higher-direct-images}).
Il suffit donc de montrer ceci : étant donnés $U$ et
$\xi \in H^q(U \times_S \Spec(K), \mathbf{Z}/d\mathbf{Z})$,
il existe un recouvrement étale $\{U_i \to U\}$
tel que $\xi$ s'annule dans $H^q(U_i \times_S \Spec(K), \mathbf{Z}/d\mathbf{Z})$.

\medskip\noindent
Nous pouvons bien entendu prendre $U$ affine. Alors $U \times_S \Spec(K)$
est une courbe affine (lisse) sur $K$, et nous avons donc l'annulation
pour $q > 1$ par le théorème \ref{theorem-vanishing-affine-curves}.

\medskip\noindent
Dernier cas : $q = 1$. Nous pouvons remplacer $U$ par les membres d'un
recouvrement étale comme dans le lemme de Compléments sur les morphismes
\ref{more-morphisms-lemma-cover-smooth-by-special}.
Alors $U \to S$ se factorise en $U \to V \to S$, où
$U \to V$ est à fibres géométriquement connexes,
$U$ et $V$ sont affines, $V \to S$ est étale et il existe
une section $\sigma : V \to U$. Par le
lemme \ref{lemma-normal-scheme-with-alg-closed-function-field},
nous voyons que $V$ est isomorphe à une réunion disjointe (finie) de
sous-schémas ouverts (affines) de $S$.
Nous pouvons clairement remplacer $S$ par l'un d'eux
et $X$ par la composante correspondante de $U$.
Nous pouvons donc supposer $X \to S$ à fibres géométriquement connexes,
muni d'une section $\sigma$, et
$\xi \in H^1(Y, \mathbf{Z}/d\mathbf{Z})$.
Puisque $K$ et $L$ sont algébriquement clos, nous avons
$$
H^1(X_L, \mathbf{Z}/d\mathbf{Z}) =
H^1(Y, \mathbf{Z}/d\mathbf{Z})
$$
Voir le lemme \ref{lemma-base-change-dim-1-separably-closed}.
Il existe donc un revêtement galoisien fini étale
$Z \to X_L$, de groupe de Galois $G \subset \mathbf{Z}/d\mathbf{Z}$,
qui annule $\xi$.
On peut le voir soit dans l'énoncé, soit dans la
démonstration du lemme \ref{lemma-smooth-base-change-fields},
soit en utilisant directement que $\xi$ correspond à un
$\mathbf{Z}/d\mathbf{Z}$-torseur sur $X_L$.
Enfin, par le
lemme \ref{pione-lemma-extend-covering-general} de Fundamental Groups,
nous trouvons un morphisme fini étale
$X' \to X$ (nécessairement surjectif)
dont la restriction à $X_L$ est $Z \to X_L$.
Puisque $\xi$ s'annule sur $X'_K$, cela achève la démonstration.
\end{proof}

\noindent
La conséquence immédiate suivante du théorème de changement de base
lisse est celle que l'on utilise souvent en pratique.

\begin{lemma}
\label{lemma-smooth-base-change-general}
Soit $S$ un schéma. Soit $S' = \lim S_i$ une limite inverse
filtrante de schémas $S_i$ lisses sur $S$, à morphismes de transition
affines. Soit $f : X \to S$ quasi-compact et quasi-séparé,
et formons le carré cartésien
$$
\xymatrix{
X' \ar[d]_{f'} \ar[r]_{g'} & X \ar[d]^f \\
S' \ar[r]^g & S
}
$$
Alors
$$
g^{-1}Rf_*E = R(f')_*(g')^{-1}E
$$
pour tout $E \in D^+(X_\etale)$ dont les faisceaux de cohomologie $H^q(E)$
ont des fibres de torsion, d'ordres inversibles sur $S$.
\end{lemma}

\begin{proof}
Considérons les suites spectrales
$$
E_2^{p, q} = R^pf_*H^q(E)
\quad\text{et}\quad
{E'}_2^{p, q} = R^pf'_*H^q((g')^{-1}E) = R^pf'_*(g')^{-1}H^q(E)
$$
convergeant vers $R^nf_*E$ et $R^nf'_*(g')^{-1}E$.
Ces suites spectrales sont construites dans le
lemme \ref{derived-lemma-two-ss-complex-functor} de Catégories dérivées.
En combinant le théorème de changement de base lisse
(théorème \ref{theorem-smooth-base-change})
avec le lemme \ref{lemma-base-change-Rf-star-colim}, nous voyons que
$$
g^{-1}R^pf_*H^q(E) = R^p(f')_*(g')^{-1}H^q(E)
$$
En combinant tout ce qui précède, nous obtenons le lemme.
\end{proof}












\section{Applications du changement de base lisse}
\label{section-applications-smooth-base-change}

\noindent
Dans cette section, nous étudions quelques conséquences plus ou moins immédiates
du théorème de changement de base lisse.

\begin{lemma}
\label{lemma-base-change-field-extension}
Soit $L/K$ une extension de corps. Soit $g : T \to S$ un morphisme quasi-compact
et quasi-séparé de schémas sur $K$. Notons
$g_L : T_L \to S_L$ le changement de base de $g$ à $\Spec(L)$.
Soit $E \in D^+(T_\etale)$ dont les faisceaux de cohomologie ont des fibres
de torsion, d'ordres inversibles dans $K$. Soit $E_L$ l'image
inverse de $E$ sur $(T_L)_\etale$. Alors
$Rg_{L, *}E_L$ est l'image inverse de $Rg_*E$ sur $S_L$.
\end{lemma}

\begin{proof}
Si $L/K$ est séparable, alors $L$ est une limite inductive filtrante de
$K$-algèbres lisses; voir le
lemme \ref{algebra-lemma-colimit-syntomic} d'Algèbre.
Dans ce cas, le lemme découle donc immédiatement du
lemme \ref{lemma-smooth-base-change-general}.
Dans le cas général, soient $K'$ et $L'$ les clôtures parfaites
(définition \ref{algebra-definition-perfection} d'Algèbre)
de $K$ et $L$. Alors $\Spec(K') \to \Spec(K)$ et
$\Spec(L') \to \Spec(L)$ sont des homéomorphismes universels,
car $K'/K$ et $L'/L$ sont purement inséparables
(voir le lemme \ref{algebra-lemma-p-ring-map} d'Algèbre).
Nous avons donc $(T_{K'})_\etale = T_\etale$,
$(S_{K'})_\etale = S_\etale$,
$(T_{L'})_\etale = (T_L)_\etale$ et
$(S_{L'})_\etale = (S_L)_\etale$, par
l'invariance topologique de la cohomologie étale; voir la
proposition \ref{proposition-topological-invariance}.
Cela ramène le lemme au cas de l'extension de corps
$L'/K'$, qui est séparable (par définition des
corps parfaits; voir la définition \ref{algebra-definition-perfect} d'Algèbre).
\end{proof}

\begin{lemma}
\label{lemma-smooth-base-change-separably-closed}
Soit $K/k$ une extension de corps séparablement clos. Soit $X$
un schéma quasi-compact et quasi-séparé sur $k$.
Soit $E \in D^+(X_\etale)$ dont les faisceaux de cohomologie ont des fibres
de torsion, d'ordres inversibles dans $k$. Alors
\begin{enumerate}
\item les morphismes $H^q_\etale(X, E) \to H^q_\etale(X_K, E|_{X_K})$
sont des isomorphismes;
\item $E \to R(X_K \to X)_*E|_{X_K}$ est un isomorphisme.
\end{enumerate}
\end{lemma}

\begin{proof}
Démonstration de (1). Soient d'abord $\overline{k}$ et $\overline{K}$
les clôtures algébriques
de $k$ et $K$. Les morphismes $\Spec(\overline{k}) \to \Spec(k)$ et
$\Spec(\overline{K}) \to \Spec(K)$ sont des homéomorphismes universels,
car $\overline{k}/k$ et $\overline{K}/K$ sont purement inséparables
(voir le lemme \ref{algebra-lemma-p-ring-map} d'Algèbre).
Ainsi $H^q_\etale(X, E) =
H^q_\etale(X_{\overline{k}}, E|_{X_{\overline{k}}})$ par
l'invariance topologique de la cohomologie étale; voir la
proposition \ref{proposition-topological-invariance}.
De même pour $X_K$ et $X_{\overline{K}}$.
Nous pouvons donc supposer $k$ et $K$ algébriquement clos.
Dans ce cas, $K$ est une limite de $k$-algèbres lisses; voir le
lemme \ref{algebra-lemma-colimit-syntomic} d'Algèbre.
Nous concluons que notre lemme est un cas particulier du
théorème \ref{theorem-smooth-base-change}, sous la forme reformulée dans le
lemme \ref{lemma-smooth-base-change-general}.

\medskip\noindent
Démonstration de (2). Pour tout $U$ quasi-compact et quasi-séparé dans $X_\etale$,
ce qui précède montre que la restriction du morphisme
$E \to R(X_K \to X)_*E|_{X_K}$ détermine un isomorphisme en cohomologie.
Puisque tout objet de $X_\etale$ possède un recouvrement étale par de tels $U$,
cela démontre l'assertion voulue.
\end{proof}

\begin{lemma}
\label{lemma-base-change-does-not-hold-post}
Avec $f : X \to S$ et $n$ comme dans la remarque \ref{remark-base-change-holds},
supposons $n$ inversible sur $S$ et, pour un certain $q \geq 1$,
$BC(f, n, q - 1)$ vraie, mais $BC(f, n, q)$ fausse.
Il existe alors un diagramme commutatif
$$
\xymatrix{
X \ar[d]_f & X' \ar[d] \ar[l] & Y \ar[l]^h \ar[d] \\
S & S' \ar[l] & \Spec(K) \ar[l]
}
$$
dont les deux carrés sont cartésiens, où $S'$ est affine, intègre et normal,
de corps des fonctions algébriquement clos $K$, ainsi qu'un entier
$d | n$ tel que $R^qh_*(\mathbf{Z}/d\mathbf{Z})$ soit non nul.
\end{lemma}

\begin{proof}
Choisissons d'abord un diagramme et $\mathcal{F}$ comme dans le
lemme \ref{lemma-base-change-does-not-hold}.
Nous pouvons supposer $S'$ affine, et nous le faisons (c'est évident, mais
voir la démonstration du lemme en cas de doute).
Soit $K'$ le corps des fonctions de $S'$,
et soit $Y' = X' \times_{S'} \Spec(K')$; nous obtenons le diagramme
$$
\xymatrix{
X \ar[d]_f &
X' \ar[d] \ar[l] &
Y' \ar[l]^{h'} \ar[d] &
Y \ar[l] \ar[d] \\
S &
S' \ar[l] &
\Spec(K') \ar[l] &
\Spec(K) \ar[l]
}
$$
Par le lemme \ref{lemma-smooth-base-change-separably-closed},
l'image directe totale $R(Y \to Y')_*\mathbf{Z}/d\mathbf{Z}$
est isomorphe à $\mathbf{Z}/d\mathbf{Z}$ dans $D(Y'_\etale)$; nous
utilisons ici que $n$ est inversible sur $S$.
Ainsi $Rh'_*\mathbf{Z}/d\mathbf{Z} = Rh_*\mathbf{Z}/d\mathbf{Z}$
par la suite spectrale de Leray relative. Cela achève la démonstration.
\end{proof}










\section{Le théorème de changement de base propre}
\label{section-proper-base-change}

\noindent
Le théorème de changement de base propre est énoncé et démontré dans cette section.
Notre approche suit approximativement la démonstration de \cite[XII, théorème 5.1]{SGA4},
en utilisant les idées de Gabber (issues du cas affine) pour simplifier
légèrement les arguments.

\begin{lemma}
\label{lemma-zariski-h0-proper-over-henselian-pair}
Soit $(A, I)$ un couple hensélien. Soit $f : X \to \Spec(A)$ un morphisme propre
de schémas. Soit $Z = X \times_{\Spec(A)} \Spec(A/I)$. Pour tout
faisceau $\mathcal{F}$ sur l'espace topologique associé à $X$, nous
avons $\Gamma(X, \mathcal{F}) = \Gamma(Z, \mathcal{F}|_Z)$.
\end{lemma}

\begin{proof}
Nous utiliserons le lemme \ref{lemma-h0-topological} pour le démontrer. Remarquons d'abord
que l'espace topologique sous-jacent à $X$ est spectral par le lemme de Properties
\ref{properties-lemma-quasi-compact-quasi-separated-spectral}.
Soit $Y \subset X$ un sous-schéma fermé irréductible. Pour achever la démonstration,
montrons que $Y \cap Z = Y \times_{\Spec(A)} \Spec(A/I)$ est connexe.
En remplaçant $X$ par $Y$,
nous pouvons supposer $X$ irréductible, et nous devons montrer que $Z$
est connexe. Soit $X \to \Spec(B) \to \Spec(A)$ la factorisation de Stein
de $f$ (théorème de Compléments sur les morphismes
\ref{more-morphisms-theorem-stein-factorization-general}).
Alors $A \to B$ est entier et $(B, IB)$ est un couple hensélien
(lemme \ref{more-algebra-lemma-integral-over-henselian-pair} de Compléments d'algèbre).
Nous pouvons donc supposer que les fibres de $X \to \Spec(A)$ sont géométriquement
connexes. D'autre part, l'image $T \subset \Spec(A)$ de $f$
est irréductible et fermée puisque $X$ est propre sur $A$. Ainsi $T \cap V(I)$
est connexe par le lemme de Compléments d'algèbre
\ref{more-algebra-lemma-irreducible-henselian-pair-connected}.
Or $Y \times_{\Spec(A)} \Spec(A/I) \to T \cap V(I)$
est une application fermée surjective à fibres connexes.
Le résultat découle alors du lemme de Topologie
\ref{topology-lemma-connected-fibres-quotient-topology-connected-components}.
\end{proof}

\begin{lemma}
\label{lemma-h0-proper-over-henselian-pair}
Soit $(A, I)$ un couple hensélien. Soit $f : X \to \Spec(A)$ un morphisme propre
de schémas. Soit $i : Z \to X$ l'immersion fermée de
$X \times_{\Spec(A)} \Spec(A/I)$ dans $X$. Pour tout
faisceau $\mathcal{F}$ sur $X_\etale$, nous
avons $\Gamma(X, \mathcal{F}) = \Gamma(Z, i_{small}^{-1}\mathcal{F})$.
\end{lemma}

\begin{proof}
Cela découle des lemmes \ref{lemma-gabber-h0} et
\ref{lemma-zariski-h0-proper-over-henselian-pair}
ainsi que du fait que tout schéma fini sur $X$ est propre sur $\Spec(A)$.
\end{proof}

\begin{lemma}
\label{lemma-h0-proper-over-henselian-local}
Soit $A$ un anneau local hensélien. Soit $f : X \to \Spec(A)$
un morphisme propre de schémas. Soit $X_0 \subset X$ la fibre de
$f$ au-dessus du point fermé. Pour tout faisceau $\mathcal{F}$ sur $X_\etale$, nous
avons $\Gamma(X, \mathcal{F}) = \Gamma(X_0, \mathcal{F}|_{X_0})$.
\end{lemma}

\begin{proof}
C'est un cas particulier du lemme \ref{lemma-h0-proper-over-henselian-pair}.
\end{proof}

\noindent
Soit $f : X \to S$ un morphisme de schémas. Soit
$\overline{s} : \Spec(k) \to S$ un point géométrique. La fibre
de $f$ en $\overline{s}$ est le schéma
$X_{\overline{s}} = \Spec(k) \times_{\overline{s}, S} X$, considéré
comme un schéma sur $\Spec(k)$. Si $\mathcal{F}$ est un faisceau sur
$X_\etale$, notons
$\mathcal{F}_{\overline{s}} = p_{small}^{-1}\mathcal{F}$
l'image inverse de $\mathcal{F}$ sur $(X_{\overline{s}})_\etale$.
Dans la suite, nous considérerons l'ensemble
$$
\Gamma(X_{\overline{s}}, \mathcal{F}_{\overline{s}})
$$
Soit $s \in S$ le point image de $\overline{s}$. Soit $\kappa(s)^{sep}$
la clôture séparable de $\kappa(s)$ dans $k$, comme dans la
définition \ref{definition-algebraic-geometric-point}.
Par le lemme \ref{lemma-sections-base-field-extension},
l'image inverse définit une bijection
$$
\Gamma(X_{\kappa(s)^{sep}},
p_{sep}^{-1} \mathcal{F})
\longrightarrow
\Gamma(X_{\overline{s}}, \mathcal{F}_{\overline{s}})
$$
où $p_{sep} : X_{\kappa(s)^{sep}} = \Spec(\kappa(s)^{sep}) \times_S X \to X$
est la projection.

\begin{lemma}
\label{lemma-proper-pushforward-stalk}
Soit $f : X \to S$ un morphisme propre de schémas. Soit
$\overline{s} \to S$ un point géométrique.
Pour tout faisceau $\mathcal{F}$ sur $X_\etale$,
le morphisme canonique
$$
(f_*\mathcal{F})_{\overline{s}} \longrightarrow
\Gamma(X_{\overline{s}}, \mathcal{F}_{\overline{s}})
$$
est bijectif.
\end{lemma}

\begin{proof}
Par le théorème \ref{theorem-higher-direct-images} (pour les faisceaux d'ensembles),
nous avons
$$
(f_*\mathcal{F})_{\overline{s}} =
\Gamma(X \times_S \Spec(\mathcal{O}_{S, \overline{s}}^{sh}),
p_{small}^{-1}\mathcal{F})
$$
où $p : X \times_S \Spec(\mathcal{O}_{S, \overline{s}}^{sh}) \to X$
est la projection. Puisque le corps résiduel de l'anneau local strictement
hensélien $\mathcal{O}_{S, \overline{s}}^{sh}$ est $\kappa(s)^{sep}$,
le lemme résulte de la discussion précédente et du
lemme \ref{lemma-h0-proper-over-henselian-local}.
\end{proof}

\begin{lemma}
\label{lemma-proper-base-change-f-star}
Soit $f : X \to Y$ un morphisme propre de schémas. Soit $g : Y' \to Y$
un morphisme de schémas. Posons $X' = Y' \times_Y X$, de projections
$f' : X' \to Y'$ et $g' : X' \to X$. Soit $\mathcal{F}$ un faisceau quelconque sur
$X_\etale$. Alors $g^{-1}f_*\mathcal{F} = f'_*(g')^{-1}\mathcal{F}$.
\end{lemma}

\begin{proof}
Il existe un morphisme canonique $g^{-1}f_*\mathcal{F} \to f'_*(g')^{-1}\mathcal{F}$.
En effet, il est adjoint au morphisme
$$
f_*\mathcal{F} \longrightarrow
g_*f'_*(g')^{-1}\mathcal{F} = f_*g'_*(g')^{-1}\mathcal{F}
$$
qui est obtenu en appliquant $f_*$ au morphisme canonique
$\mathcal{F} \to g'_*(g')^{-1}\mathcal{F}$. Pour vérifier que ce morphisme est un
isomorphisme, nous pouvons calculer ce qui se passe sur les fibres.
Soit $y' : \Spec(k) \to Y'$ un point géométrique d'image $y$ dans $Y$.
Par le lemme \ref{lemma-proper-pushforward-stalk}, les fibres sont respectivement
$\Gamma(X'_{y'}, \mathcal{F}_{y'})$ et $\Gamma(X_y, \mathcal{F}_y)$.
Ici, les faisceaux $\mathcal{F}_y$ et $\mathcal{F}_{y'}$
sont les images inverses de $\mathcal{F}$ par les projections $X_y \to X$
et $X'_{y'} \to X$. Puisque $X_y = X'_{y'}$ comme schémas,
nous voyons que les fibres coïncident. Nous omettons la
vérification de la compatibilité de cet isomorphisme avec notre morphisme.
\end{proof}

\noindent
Nous commençons maintenant l'étude du théorème de changement de base propre.
Pour cela, introduisons quelques notations. Considérons un diagramme commutatif
\begin{equation}
\label{equation-base-change-diagram}
\vcenter{
\xymatrix{
X' \ar[r]_{g'} \ar[d]_{f'} & X \ar[d]^f \\
Y' \ar[r]^g & Y
}
}
\end{equation}
de morphismes de schémas. Nous obtenons alors un diagramme commutatif de sites
$$
\xymatrix{
X'_\etale \ar[r]_{g'_{small}} \ar[d]_{f'_{small}} &
X_\etale \ar[d]^{f_{small}} \\
Y'_\etale \ar[r]^{g_{small}} &
Y_\etale
}
$$
Pour tout objet $E$ de $D(X_\etale)$, nous obtenons un morphisme canonique de changement de base
\begin{equation}
\label{equation-base-change}
g_{small}^{-1}Rf_{small, *}E \longrightarrow Rf'_{small, *}(g'_{small})^{-1}E
\end{equation}
dans $D(Y'_\etale)$. Voir la remarque de Cohomologie sur les sites
\ref{sites-cohomology-remark-base-change}, où nous employons le faisceau
constant $\mathbf{Z}$ comme faisceau d'anneaux.
Nous omettrons généralement les indices ${}_{small}$ dans cette formule.
Par exemple, si $E = \mathcal{F}[0]$, où $\mathcal{F}$ est un faisceau
abélien sur $X_\etale$, le morphisme de changement de base est un morphisme
\begin{equation}
\label{equation-base-change-sheaf}
g^{-1}Rf_*\mathcal{F} \longrightarrow Rf'_*(g')^{-1}\mathcal{F}
\end{equation}
dans $D(Y'_\etale)$.

\medskip\noindent
Le morphisme (\ref{equation-base-change}) n'a aucune chance d'être un isomorphisme
dans la généralité ci-dessus. Le but est de montrer
qu'il est un isomorphisme si le diagramme (\ref{equation-base-change-diagram})
est cartésien, si $f : X \to Y$ est propre, si les faisceaux de cohomologie de $E$
sont de torsion et si $E$ est borné inférieurement.
Pour étudier cette question, introduisons la terminologie suivante.
Nous dirons que {\it la cohomologie commute au changement de base
pour $f : X \to Y$} si (\ref{equation-base-change-sheaf})
est un isomorphisme pour tout diagramme (\ref{equation-base-change-diagram})
où $X' = Y' \times_Y X$, et pour tout faisceau abélien de torsion $\mathcal{F}$.

\begin{lemma}
\label{lemma-proper-base-change-in-terms-of-injectives}
Soit $f : X \to Y$ un morphisme propre de schémas.
Les assertions suivantes sont équivalentes :
\begin{enumerate}
\item la cohomologie commute au changement de base pour $f$ (voir ci-dessus);
\item pour tout nombre premier $\ell$, tout faisceau injectif
de $\mathbf{Z}/\ell\mathbf{Z}$-modules $\mathcal{I}$
sur $X_\etale$ et tout diagramme (\ref{equation-base-change-diagram})
où $X' = Y' \times_Y X$, les faisceaux
$R^qf'_*(g')^{-1}\mathcal{I}$ sont nuls pour $q > 0$.
\end{enumerate}
\end{lemma}

\begin{proof}
Il est clair que (1) implique (2). Réciproquement, supposons (2) et soit
$\mathcal{F}$ un faisceau abélien de torsion sur $X_\etale$. Soit $Y' \to Y$
un morphisme de schémas et soit $X' = Y' \times_Y X$,
de projections $g' : X' \to X$ et $f' : X' \to Y'$, comme dans le
diagramme (\ref{equation-base-change-diagram}).
Nous voulons montrer que les morphismes de faisceaux
$$
g^{-1}R^qf_*\mathcal{F} \longrightarrow R^qf'_*(g')^{-1}\mathcal{F}
$$
sont des isomorphismes pour tout $q \geq 0$.

\medskip\noindent
Pour tout $n \geq 1$, soit $\mathcal{F}[n]$ le sous-faisceau des sections
de $\mathcal{F}$ annulées par $n$. Alors
$\mathcal{F} = \colim \mathcal{F}[n]$.
Les foncteurs $g^{-1}$ et $(g')^{-1}$ commutent aux limites inductives quelconques
(en tant qu'adjoints à gauche). La prise des images directes supérieures par $f$ ou $f'$
commute aux limites inductives filtrantes par le lemme \ref{lemma-relative-colimit}.
Nous voyons donc que
$$
g^{-1}R^qf_*\mathcal{F} = \colim g^{-1}R^qf_*\mathcal{F}[n]
\quad\text{et}\quad
R^qf'_*(g')^{-1}\mathcal{F} =
\colim R^qf'_*(g')^{-1}\mathcal{F}[n]
$$
Il suffit donc de démontrer le résultat lorsque $\mathcal{F}$ est
annulé par un entier positif $n$.

\medskip\noindent
Si $n = \ell n'$ pour un certain nombre premier $\ell$, nous obtenons une suite
exacte courte
$$
0 \to \mathcal{F}[\ell] \to \mathcal{F} \to
\mathcal{F}/\mathcal{F}[\ell] \to 0
$$
Remarquons que $\mathcal{F}/\mathcal{F}[\ell]$ est annulé par $n'$.
De plus, si le résultat vaut à la fois pour $\mathcal{F}[\ell]$ et
$\mathcal{F}/\mathcal{F}[\ell]$, alors le résultat découle de
la suite exacte longue des images directes supérieures (et du lemme des $5$).
Nous nous ramenons ainsi au cas où $\mathcal{F}$ est annulé
par un nombre premier $\ell$.

\medskip\noindent
Supposons que $\mathcal{F}$ soit annulé par un nombre premier $\ell$.
Choisissons une résolution injective $\mathcal{F} \to \mathcal{I}^\bullet$
dans $D(X_\etale, \mathbf{Z}/\ell\mathbf{Z})$. En appliquant l'hypothèse
(2) et le lemme d'acyclicité de Leray
(lemme \ref{derived-lemma-leray-acyclicity} de Catégories dérivées),
nous voyons que
$$
f'_*(g')^{-1}\mathcal{I}^\bullet
$$
calcule $Rf'_*(g')^{-1}\mathcal{F}$. Nous concluons en appliquant le
lemme \ref{lemma-proper-base-change-f-star}.
\end{proof}

\begin{lemma}
\label{lemma-sandwich}
Soient $f : X \to Y$ et $g : Y \to Z$ des morphismes propres de schémas. Supposons que :
\begin{enumerate}
\item la cohomologie commute au changement de base pour $f$;
\item la cohomologie commute au changement de base pour $g \circ f$;
\item $f$ est surjectif.
\end{enumerate}
Alors la cohomologie commute au changement de base pour $g$.
\end{lemma}

\begin{proof}
Nous utiliserons sans autre mention l'équivalence du
lemme \ref{lemma-proper-base-change-in-terms-of-injectives}.
Soit $\ell$ un nombre premier.
Soit $\mathcal{I}$ un faisceau injectif de
$\mathbf{Z}/\ell\mathbf{Z}$-modules sur $Y_\etale$.
Choisissons un morphisme injectif de faisceaux $f^{-1}\mathcal{I} \to \mathcal{J}$,
où $\mathcal{J}$ est un faisceau injectif de
$\mathbf{Z}/\ell\mathbf{Z}$-modules sur $X_\etale$.
Puisque $f$ est surjectif, le morphisme $\mathcal{I} \to f_*\mathcal{J}$
est injectif (il suffit de le vérifier sur les fibres aux points géométriques).
Puisque $\mathcal{I}$ est injectif, nous voyons que $\mathcal{I}$
est un facteur direct de $f_*\mathcal{J}$. Il suffit donc
de démontrer l'annulation voulue pour $f_*\mathcal{J}$.

\medskip\noindent
Soit $Z' \to Z$ un morphisme de schémas et posons
$Y' = Z' \times_Z Y$ et $X' = Z' \times_Z X = Y' \times_ Y X$.
Notons $a : X' \to X$, $b : Y' \to Y$ et $c : Z' \to Z$ les
projections. De même pour $f' : X' \to Y'$ et $g' : Y' \to Z'$.
Par le lemme \ref{lemma-proper-base-change-f-star}, nous avons
$b^{-1}f_*\mathcal{J} = f'_*a^{-1}\mathcal{J}$.
D'autre part, nous savons que $R^qf'_*a^{-1}\mathcal{J}$ et
$R^q(g' \circ f')_*a^{-1}\mathcal{J}$ sont nuls pour $q > 0$.
En utilisant la suite spectrale
(lemme \ref{sites-cohomology-lemma-relative-Leray} de Cohomologie sur les sites)
$$
R^pg'_* R^qf'_* a^{-1}\mathcal{J} \Rightarrow
R^{p + q}(g' \circ f')_* a^{-1}\mathcal{J}
$$
nous concluons que
$ R^pg'_*(b^{-1}f_*\mathcal{J}) = R^pg'_*(f'_*a^{-1}\mathcal{J}) = 0$
pour $p > 0$, comme voulu.
\end{proof}

\begin{lemma}
\label{lemma-composition}
Soient $f : X \to Y$ et $g : Y \to Z$ des morphismes propres de schémas. Supposons que :
\begin{enumerate}
\item la cohomologie commute au changement de base pour $f$;
\item la cohomologie commute au changement de base pour $g$.
\end{enumerate}
Alors la cohomologie commute au changement de base pour $g \circ f$.
\end{lemma}

\begin{proof}
Nous utiliserons sans autre mention l'équivalence du
lemme \ref{lemma-proper-base-change-in-terms-of-injectives}.
Soit $\ell$ un nombre premier.
Soit $\mathcal{I}$ un faisceau injectif de
$\mathbf{Z}/\ell\mathbf{Z}$-modules sur $X_\etale$.
Alors $f_*\mathcal{I}$ est un faisceau injectif de
$\mathbf{Z}/\ell\mathbf{Z}$-modules sur $Y_\etale$
(lemme de Cohomologie sur les sites
\ref{sites-cohomology-lemma-pushforward-injective-flat}).
Le résultat en découle formellement, mais nous allons aussi
l'expliciter.

\medskip\noindent
Soit $Z' \to Z$ un morphisme de schémas et posons
$Y' = Z' \times_Z Y$ et $X' = Z' \times_Z X = Y' \times_ Y X$.
Notons $a : X' \to X$, $b : Y' \to Y$ et $c : Z' \to Z$ les
projections. De même pour $f' : X' \to Y'$ et $g' : Y' \to Z'$.
Par le lemme \ref{lemma-proper-base-change-f-star}, nous avons
$b^{-1}f_*\mathcal{I} = f'_*a^{-1}\mathcal{I}$.
D'autre part, nous savons que $R^qf'_*a^{-1}\mathcal{I}$ et
$R^q(g')_*b^{-1}f_*\mathcal{I}$ sont nuls pour $q > 0$.
En utilisant la suite spectrale
(lemme \ref{sites-cohomology-lemma-relative-Leray} de Cohomologie sur les sites)
$$
R^pg'_* R^qf'_* a^{-1}\mathcal{I} \Rightarrow
R^{p + q}(g' \circ f')_* a^{-1}\mathcal{I}
$$
nous concluons que $R^p(g' \circ f')_*a^{-1}\mathcal{I} = 0$ pour
$p > 0$, comme voulu.
\end{proof}

\begin{lemma}
\label{lemma-finite}
\begin{slogan}
Le changement de base propre en cohomologie étale est valable pour les morphismes finis.
\end{slogan}
Soit $f : X \to Y$ un morphisme fini de schémas.
Alors la cohomologie commute au changement de base pour $f$.
\end{lemma}

\begin{proof}
Remarquons qu'un morphisme fini est propre; voir le
lemme \ref{morphisms-lemma-finite-proper} de Morphismes.
De plus, le changement de base d'un morphisme fini est fini; voir le
lemme \ref{morphisms-lemma-base-change-finite} de Morphismes.
Le résultat découle donc du
lemme \ref{lemma-proper-base-change-in-terms-of-injectives},
combiné à la
proposition \ref{proposition-finite-higher-direct-image-zero}.
\end{proof}

\begin{lemma}
\label{lemma-reduce-to-P1}
Pour démontrer que la cohomologie commute au changement de base pour
tout morphisme propre de schémas, il suffit de démontrer qu'elle
commute pour le morphisme $\mathbf{P}^1_S \to S$, pour tout schéma $S$.
\end{lemma}

\begin{proof}
Soit $f : X \to Y$ un morphisme propre de schémas.
Soit $Y = \bigcup Y_i$ un recouvrement ouvert affine,
et posons $X_i = f^{-1}(Y_i)$. Si nous pouvons démontrer
que la cohomologie commute au changement de base pour $X_i \to Y_i$,
alors elle commute au changement de base pour $f$.
En effet, la formation des images directes supérieures
commute à la localisation de Zariski (et même \'etale)
sur la base; voir le
lemme \ref{lemma-higher-direct-images}.
Nous pouvons donc supposer $Y$ affine.

\medskip\noindent
Soit $Y$ un schéma affine et soit $X \to Y$ un morphisme propre.
Par le lemme de Chow, il existe un diagramme commutatif
$$
\xymatrix{
X \ar[rd] & X' \ar[d] \ar[l]^\pi \ar[r] & \mathbf{P}^n_Y \ar[dl] \\
& Y &
}
$$
où $X' \to \mathbf{P}^n_Y$ est une immersion et
$\pi : X' \to X$ est propre et surjectif; voir le
lemme \ref{limits-lemma-chow-finite-type} de Limits.
Puisque $X \to Y$ est propre, on voit que $X' \to Y$ est propre
(lemme \ref{morphisms-lemma-composition-proper} de Morphismes).
Par conséquent, $X' \to \mathbf{P}^n_Y$ est une immersion fermée
(lemme \ref{morphisms-lemma-image-proper-scheme-closed} de Morphismes).
Il s'ensuit que $X' \to X \times_Y \mathbf{P}^n_Y = \mathbf{P}^n_X$
est une immersion fermée (en tant qu'immersion d'image fermée).

\medskip\noindent
Par le lemme \ref{lemma-sandwich},
il suffit de démontrer que la cohomologie commute au changement de base pour
$\pi$ et $X' \to Y$. Ces deux morphismes se factorisent en une immersion
fermée suivie d'une projection $\mathbf{P}^n_S \to S$ (pour un certain $S$).
Par le lemme \ref{lemma-finite}, le résultat vaut pour les immersions
fermées (car les immersions fermées sont finies).
Par le lemme \ref{lemma-composition}, il suffit de démontrer le
résultat pour les projections $\mathbf{P}^n_S \to S$.

\medskip\noindent
Pour tout $n \geq 1$, il existe un morphisme fini surjectif
$$
\mathbf{P}^1_S \times_S \ldots \times_S \mathbf{P}^1_S
\longrightarrow
\mathbf{P}^n_S
$$
donné en coordonnées par
$$
((x_1 : y_1), (x_2 : y_2), \ldots, (x_n : y_n))
\longmapsto
(F_0 : \ldots : F_n)
$$
où $F_0, \ldots, F_n$ sont les polynômes en $x_1, \ldots, y_n$
à coefficients entiers tels que
$$
\prod (x_i t + y_i) = F_0 t^n + F_1 t^{n - 1} + \ldots + F_n
$$
En appliquant encore une fois
les lemmes \ref{lemma-sandwich}, \ref{lemma-finite} et \ref{lemma-composition},
nous concluons que le lemme est vrai.
\end{proof}

\begin{theorem}[Changement de base propre]
\label{theorem-proper-base-change}
Soit $f : X \to Y$ un morphisme propre de schémas. Soit $g : Y' \to Y$
un morphisme de schémas. Posons $X' = Y' \times_Y X$
et considérons le diagramme cartésien
$$
\xymatrix{
X' \ar[r]_{g'} \ar[d]_{f'} & X \ar[d]^f \\
Y' \ar[r]^g & Y
}
$$
Soit $\mathcal{F}$ un faisceau abélien de torsion sur $X_\etale$.
Alors le morphisme de changement de base
$$
g^{-1}Rf_*\mathcal{F} \longrightarrow Rf'_*(g')^{-1}\mathcal{F}
$$
est un isomorphisme.
\end{theorem}

\begin{proof}
Dans la terminologie introduite ci-dessus, cela signifie que la cohomologie commute
au changement de base pour tout morphisme propre de schémas. Par le
lemme \ref{lemma-reduce-to-P1},
il suffit de démontrer que la cohomologie commute au changement de base
pour le morphisme $\mathbf{P}^1_S \to S$, pour tout schéma $S$.

\medskip\noindent
Soit $S$ le spectre d'un anneau local strictement hensélien, de point
fermé $s$. Posons $X = \mathbf{P}^1_S$ et $X_0 = X_s = \mathbf{P}^1_s$.
Soit $\mathcal{F}$ un faisceau de $\mathbf{Z}/\ell\mathbf{Z}$-modules
sur $X_\etale$. Le point clef de notre démonstration est l'égalité
$$
H^q_\etale(X, \mathcal{F}) = H^q_\etale(X_0, \mathcal{F}|_{X_0}).
$$
En effet, choisissons une résolution $\mathcal{F} \to \mathcal{I}^\bullet$
par des faisceaux injectifs de $\mathbf{Z}/\ell\mathbf{Z}$-modules.
Alors $\mathcal{I}^\bullet|_{X_0}$ est une résolution de $\mathcal{F}|_{X_0}$
par des objets acycliques à droite pour $H^0_\etale(X_0, -)$; voir le
lemme \ref{lemma-efface-cohomology-on-fibre-by-finite-cover}.
D'après le lemme d'acyclicité de Leray, le membre de droite est calculé par
le complexe $H^0_\etale(X_0, \mathcal{I}^\bullet|_{X_0})$,
qui est égal à $H^0_\etale(X, \mathcal{I}^\bullet)$ par le
lemme \ref{lemma-h0-proper-over-henselian-local}. Ce complexe
calcule le membre de gauche.

\medskip\noindent
Supposons maintenant $S$ quelconque et $\mathcal{F}$ un faisceau de
$\mathbf{Z}/\ell\mathbf{Z}$-modules sur $X_\etale$.
Soit $\overline{s} : \Spec(k) \to S$ un point géométrique
de $S$ au-dessus de $s \in S$. Nous avons
$$
(R^qf_*\mathcal{F})_{\overline{s}} =
H^q_\etale(\mathbf{P}^1_{\mathcal{O}_{S, \overline{s}}^{sh}},
\mathcal{F}|_{\mathbf{P}^1_{\mathcal{O}_{S, \overline{s}}^{sh}}}) =
H^q_\etale(\mathbf{P}^1_{\kappa(s)^{sep}},
\mathcal{F}|_{\mathbf{P}^1_{\kappa(s)^{sep}}})
$$
où $\kappa(s)^{sep}$ est le corps résiduel de
$\mathcal{O}_{S, \overline{s}}^{sh}$, c'est-à-dire la clôture algébrique
séparable de $\kappa(s)$ dans $k$.
La première égalité découle du théorème \ref{theorem-higher-direct-images},
et la seconde de la formule affichée au
paragraphe précédent.

\medskip\noindent
Enfin, considérons un morphisme quelconque de schémas $g : T \to S$, où
$S$ et $\mathcal{F}$ sont comme ci-dessus.
Soit $f' : \mathbf{P}^1_T \to T$ la projection et soit
$g' : \mathbf{P}^1_T \to \mathbf{P}^1_S$ le morphisme induit
par $g$. Considérons le morphisme de changement de base
$$
g^{-1}R^qf_*\mathcal{F}
\longrightarrow
R^qf'_*(g')^{-1}\mathcal{F}
$$
Soit $\overline{t}$ un point géométrique de $T$, d'image
$\overline{s} = g(\overline{t})$. D'après la discussion
ci-dessus, le morphisme induit sur les fibres en $\overline{t}$ est le morphisme
$$
H^q_\etale(\mathbf{P}^1_{\kappa(s)^{sep}},
\mathcal{F}|_{\mathbf{P}^1_{\kappa(s)^{sep}}})
\longrightarrow
H^q_\etale(\mathbf{P}^1_{\kappa(t)^{sep}},
\mathcal{F}|_{\mathbf{P}^1_{\kappa(t)^{sep}}})
$$
Puisque $\kappa(s)^{sep} \subset \kappa(t)^{sep}$, ce morphisme est un
isomorphisme par le lemme \ref{lemma-base-change-dim-1-separably-closed}.

\medskip\noindent
Cela démontre que la cohomologie commute au changement de base pour
$\mathbf{P}^1_S \to S$, avec des faisceaux de $\mathbf{Z}/\ell\mathbf{Z}$-modules.
En particulier, pour un faisceau injectif de $\mathbf{Z}/\ell\mathbf{Z}$-modules,
les images directes supérieures de tout changement de base sont nulles.
Autrement dit, la condition (2) du
lemme \ref{lemma-proper-base-change-in-terms-of-injectives}
est satisfaite, ce qui achève la démonstration.
\end{proof}

\begin{lemma}
\label{lemma-proper-base-change}
Soit $f : X \to Y$ un morphisme propre de schémas. Soit $g : Y' \to Y$
un morphisme de schémas. Posons $X' = Y' \times_Y X$ et notons
$f' : X' \to Y'$ et $g' : X' \to X$ les projections.
Soit $E \in D^+(X_\etale)$ à faisceaux de cohomologie de torsion.
Alors le morphisme de changement de base (\ref{equation-base-change})
$g^{-1}Rf_*E \to Rf'_*(g')^{-1}E$
est un isomorphisme.
\end{lemma}

\begin{proof}
C'est une conséquence simple du théorème de changement de base propre
(théorème \ref{theorem-proper-base-change}), en utilisant les suites
spectrales
$$
E_2^{p, q} = R^pf_*H^q(E)
\quad\text{et}\quad
{E'}_2^{p, q} = R^pf'_*(g')^{-1}H^q(E)
$$
convergeant vers $R^nf_*E$ et $R^nf'_*(g')^{-1}E$.
Les suites spectrales sont construites dans le
lemme \ref{derived-lemma-two-ss-complex-functor} de Catégories dérivées.
Certains détails sont omis.
\end{proof}

\begin{lemma}
\label{lemma-proper-base-change-stalk}
Soit $f : X \to Y$ un morphisme propre de schémas. Soit $\overline{y} \to Y$
un point géométrique.
\begin{enumerate}
\item Pour un faisceau abélien de torsion $\mathcal{F}$ sur $X_\etale$, nous avons
$(R^nf_*\mathcal{F})_{\overline{y}} =
H^n_\etale(X_{\overline{y}}, \mathcal{F}_{\overline{y}})$.
\item Pour $E \in D^+(X_\etale)$ à faisceaux de cohomologie de torsion, nous avons
$(R^nf_*E)_{\overline{y}} =
H^n_\etale(X_{\overline{y}}, E|_{X_{\overline{y}}})$.
\end{enumerate}
\end{lemma}

\begin{proof}
Dans l'énoncé, $\mathcal{F}_{\overline{y}}$ désigne l'image inverse
de $\mathcal{F}$ sur la fibre schématique
$X_{\overline{y}} = \overline{y} \times_Y X$.
Puisque l'image inverse par $\overline{y} \to Y$ fournit la
fibre de $\mathcal{F}$, la première assertion du lemme
est un cas particulier du théorème \ref{theorem-proper-base-change}.
La seconde est un cas particulier du lemme \ref{lemma-proper-base-change}.
\end{proof}






\section{Applications du changement de base propre}
\label{section-applications-proper-base-change}

\noindent
Dans cette section, nous étudions quelques conséquences plus ou moins immédiates
du théorème de changement de base propre.

\begin{lemma}
\label{lemma-base-change-separably-closed}
Soit $K/k$ une extension de corps séparablement clos.
Soit $X$ un schéma propre sur $k$.
Soit $\mathcal{F}$ un faisceau abélien de torsion sur $X_\etale$.
Alors le morphisme $H^q_\etale(X, \mathcal{F}) \to
H^q_\etale(X_K, \mathcal{F}|_{X_K})$ est un isomorphisme
pour $q \geq 0$.
\end{lemma}

\begin{proof}
L'examen des fibres montre qu'il s'agit
d'un cas particulier du
théorème \ref{theorem-proper-base-change}.
\end{proof}

\begin{lemma}
\label{lemma-cohomological-dimension-proper}
Soit $f : X \to Y$ un morphisme propre de schémas
dont toutes les fibres sont de dimension $\leq n$.
Alors, pour tout faisceau abélien de torsion $\mathcal{F}$ sur $X_\etale$,
nous avons $R^qf_*\mathcal{F} = 0$ pour $q > 2n$.
\end{lemma}

\begin{proof}
Nous le démontrons par récurrence sur $n$, pour tous les morphismes propres.

\medskip\noindent
Si $n = 0$, alors $f$ est un morphisme fini
(lemme \ref{more-morphisms-lemma-characterize-finite} de Compléments sur les morphismes),
et le
résultat est vrai par la proposition \ref{proposition-finite-higher-direct-image-zero}.

\medskip\noindent
Si $n > 0$, le lemme \ref{lemma-proper-base-change-stalk}
montre qu'il suffit de démontrer $H^i_\etale(X, \mathcal{F}) = 0$
pour $i > 2n$, lorsque $X$ est un schéma propre, $\dim(X) \leq n$,
sur un corps algébriquement clos $k$,
et que $\mathcal{F}$ est un faisceau abélien de torsion sur $X$.

\medskip\noindent
Si $n = 1$, cela découle du théorème \ref{theorem-vanishing-curves}.
Supposons $n > 1$. Par la proposition \ref{proposition-topological-invariance},
nous pouvons remplacer $X$ par sa réduction.
Soit $\nu : X^\nu \to X$ la normalisation.
C'est un morphisme fini birationnel surjectif
(voir le lemme \ref{varieties-lemma-normalization-locally-algebraic} de Variétés),
donc un isomorphisme au-dessus d'un ouvert dense $U \subset X$
(lemme \ref{morphisms-lemma-birational-birational} de Morphismes).
Nous voyons alors que
$c : \mathcal{F} \to \nu_*\nu^{-1}\mathcal{F}$ est injectif
(car $\nu$ est surjectif) et est un isomorphisme au-dessus de $U$.
Notons $i : Z \to X$ l'inclusion du complémentaire de $U$.
Puisque $U$ est dense dans $X$, nous avons $\dim(Z) < \dim(X) = n$.
Par la proposition \ref{proposition-closed-immersion-pushforward}, nous avons
$\Coker(c) = i_*\mathcal{G}$ pour un certain
faisceau abélien de torsion $\mathcal{G}$ sur $Z_\etale$.
Alors $H^q_\etale(X, \Coker(c)) = H^q_\etale(Z, \mathcal{G})$
(par la proposition \ref{proposition-finite-higher-direct-image-zero}
et la suite spectrale de Leray), et l'hypothèse de récurrence montre que
le conoyau de $c$ a sa cohomologie en degrés $\leq 2(n - 1)$.
Il suffit donc de démontrer le résultat pour $\nu_*\nu^{-1}\mathcal{F}$.
Comme $\nu$ est fini, cela nous ramène à montrer
que $H^i_\etale(X^\nu, \nu^{-1}\mathcal{F})$ est
nul pour $i > 2n$. Ce cas est traité au paragraphe suivant.

\medskip\noindent
Supposons que $X$ soit un schéma propre, intègre et normal sur $k$, de dimension $n$.
Choisissons une fonction rationnelle non constante $f$ sur $X$. Le graphe
$X' \subset X \times \mathbf{P}^1_k$ de $f$ s'insère dans un diagramme
$$
X \xleftarrow{b} X' \xrightarrow{f} \mathbf{P}^1_k
$$
Remarquons que $b$ est un isomorphisme au-dessus d'un sous-schéma ouvert
$U \subset X$ dont le complémentaire est un sous-schéma fermé
$Z \subset X$ de codimension $\geq 2$. En effet, $U$ est le
domaine de définition de $f$, qui contient tous les points de codimension $1$
de $X$; voir les lemmes de Morphismes
\ref{morphisms-lemma-rational-map-from-reduced-to-separated}
et \ref{morphisms-lemma-extend-across}
(combinés au critère de normalité de Serre; voir le
lemme \ref{properties-lemma-criterion-normal} de Properties).
De plus, les fibres de $b$ sont de dimension $\leq 1$ (comme sous-schémas fermés
de $\mathbf{P}^1$). Ainsi, par récurrence, $R^ib_*b^{-1}\mathcal{F}$ n'est non nul que
si $i \in \{0, 1, 2\}$. Choisissons un triangle distingué
$$
\mathcal{F} \to Rb_*b^{-1}\mathcal{F} \to Q \to \mathcal{F}[1]
$$
Comme précédemment, l'injectivité de $\mathcal{F} \to b_*b^{-1}\mathcal{F}$
et ce que nous venons de montrer impliquent que $Q$ a des
faisceaux de cohomologie non nuls seulement en degrés $0, 1, 2$, portés par $Z$.
De plus, ces faisceaux de cohomologie sont de torsion par le
lemme \ref{lemma-torsion-direct-image}.
Par récurrence et par la suite spectrale du
lemme \ref{derived-lemma-two-ss-complex-functor} de Catégories dérivées,
nous voyons que $H^i(X, Q)$ est nul pour
$i > 2 + 2\dim(Z) \leq 2 + 2(n - 2) = 2n - 2$. Il suffit donc
de démontrer que $H^i(X', b^{-1}\mathcal{F}) = 0$ pour
$i > 2n$. Nous utilisons alors le morphisme
$$
f : X' \to \mathbf{P}^1_k
$$
dont les fibres sont de dimension $< n$. Ainsi, par récurrence, nous voyons que
$R^if_*b^{-1}\mathcal{F} = 0$ pour $i > 2(n - 1)$.
Nous concluons par la suite spectrale de Leray
$$
H^i(\mathbf{P}^1_k, R^jf_*b^{-1}\mathcal{F})
\Rightarrow
H^{i + j}(X', b^{-1}\mathcal{F})
$$
et le fait que $\dim(\mathbf{P}^1_k) = 1$.
\end{proof}

\noindent
Avec des coefficients modulo $n$, nous pouvons appliquer le changement
de base propre aux complexes non bornés.

\begin{lemma}
\label{lemma-proper-base-change-mod-n}
Soit $f : X \to Y$ un morphisme propre de schémas. Soit $g : Y' \to Y$
un morphisme de schémas. Posons $X' = Y' \times_Y X$ et notons
$f' : X' \to Y'$ et $g' : X' \to X$ les projections.
Soit $n \geq 1$ un entier.
Soit $E \in D(X_\etale, \mathbf{Z}/n\mathbf{Z})$.
Alors le morphisme de changement de base (\ref{equation-base-change})
$g^{-1}Rf_*E \to Rf'_*(g')^{-1}E$
est un isomorphisme.
\end{lemma}

\begin{proof}
Il suffit de le démontrer lorsque $Y$ et $Y'$ sont quasi-compacts.
Par le lemme de Morphismes
\ref{morphisms-lemma-morphism-finite-type-bounded-dimension}
nous voyons que les dimensions des fibres de
$f : X \to Y$ et de $f' : X' \to Y'$ sont bornées. Ainsi, le
lemme \ref{lemma-cohomological-dimension-proper} implique que
$$
f_* :
\textit{Mod}(X_\etale, \mathbf{Z}/n\mathbf{Z})
\longrightarrow
\textit{Mod}(Y_\etale, \mathbf{Z}/n\mathbf{Z})
$$
et
$$
f'_* :
\textit{Mod}(X'_\etale, \mathbf{Z}/n\mathbf{Z})
\longrightarrow
\textit{Mod}(Y'_\etale, \mathbf{Z}/n\mathbf{Z})
$$
ont une dimension cohomologique finie au sens du
lemme \ref{derived-lemma-unbounded-right-derived} de Catégories dérivées.
Choisissons un complexe K-injectif $\mathcal{I}^\bullet$
de $\mathbf{Z}/n\mathbf{Z}$-modules dont chacun des
termes $\mathcal{I}^n$ est un faisceau injectif de
$\mathbf{Z}/n\mathbf{Z}$-modules; ce complexe représente $E$.
Voir le théorème d'Injectifs
\ref{injectives-theorem-K-injective-embedding-grothendieck}.
Par le théorème usuel de changement de base propre, on obtient
que $R^qf'_*(g')^{-1}\mathcal{I}^n = 0$ pour $q > 0$; voir le
théorème \ref{theorem-proper-base-change}.
Nous concluons donc, par le
lemme \ref{derived-lemma-unbounded-right-derived} de Catégories dérivées,
que nous pouvons calculer $Rf'_*(g')^{-1}E$ par le complexe
$f'_*(g')^{-1}\mathcal{I}^\bullet$. Une autre application
du théorème usuel de changement de base propre montre que
ce complexe est égal à $g^{-1}f_*\mathcal{I}^\bullet$, comme souhaité.
\end{proof}

\begin{lemma}
\label{lemma-pull-out-constant}
Soit $X$ un schéma quasi-compact et quasi-séparé.
Soient $E \in D^+(X_\etale)$ et $K \in D^+(\mathbf{Z})$.
Alors
$$
R\Gamma(X, E \otimes_\mathbf{Z}^\mathbf{L} \underline{K}) =
R\Gamma(X, E) \otimes_\mathbf{Z}^\mathbf{L} K
$$
\end{lemma}

\begin{proof}
Supposons $H^i(E) = 0$ pour $i \geq a$ et $H^j(K) = 0$ pour $j \geq b$.
Nous pouvons représenter $K$ par un complexe borné inférieurement
$K^\bullet$ de $\mathbf{Z}$-modules sans torsion.
(Choisissons un complexe K-plat $L^\bullet$ représentant $K$, puis
prenons $K^\bullet = \tau_{\geq b - 1}L^\bullet$. Cela fonctionne parce que
$\mathbf{Z}$ est de dimension globale $1$. Voir le
lemme \ref{more-algebra-lemma-last-one-flat} de Compléments d'algèbre.)
Nous pouvons représenter $E$ par un complexe borné inférieurement
$\mathcal{E}^\bullet$.
Alors $E \otimes_\mathbf{Z}^\mathbf{L} \underline{K}$ est
représenté par
$$
\text{Tot}(\mathcal{E}^\bullet \otimes_\mathbf{Z} \underline{K}^\bullet)
$$
En utilisant les triangles distingués
$$
\sigma_{\geq -b + n + 1}K^\bullet
\to K^\bullet \to
\sigma_{\leq -b + n}K^\bullet
$$
et l'annulation évidente
$$
H^n(X,
\text{Tot}(\mathcal{E}^\bullet \otimes_\mathbf{Z}
\sigma_{\geq -a + n + 1}\underline{K}^\bullet) = 0
$$
et
$$
H^n(R\Gamma(X, E)
\otimes_\mathbf{Z}^\mathbf{L} \sigma_{\geq -a + n + 1}K^\bullet) = 0
$$
nous nous ramenons au cas où $K^\bullet$ est un complexe borné de
$\mathbf{Z}$-modules plats. En répétant l'argument, nous nous ramenons au
cas où $K^\bullet$ est réduit à un unique $\mathbf{Z}$-module plat
placé en un certain degré. Ensuite, en utilisant les troncations bêtes
de $\mathcal{E}^\bullet$,
nous nous ramenons exactement de la même manière au cas où
$\mathcal{E}^\bullet$ est réduit à un unique faisceau abélien placé
en un certain degré. Il suffit donc de montrer que
$$
H^n(X, \mathcal{E} \otimes_\mathbf{Z} \underline{M}) =
H^n(X, \mathcal{E}) \otimes_\mathbf{Z} M
$$
Lorsque $M$ est un $\mathbf{Z}$-module plat et que $\mathcal{E}$ est
un faisceau abélien sur $X$,
on écrit $M$ comme une limite inductive filtrante de $\mathbf{Z}$-modules libres finis
(théorème de Lazard; voir le théorème \ref{algebra-theorem-lazard} d'Algèbre).
Par le théorème \ref{theorem-colimit}, cela nous ramène au cas d'un module libre fini
sur $\mathbf{Z}$, noté $M$, auquel cas le résultat est trivialement vrai.
\end{proof}

\begin{lemma}
\label{lemma-projection-formula-proper}
Soit $f : X \to Y$ un morphisme propre de schémas.
Soit $E \in D^+(X_\etale)$ à faisceaux de cohomologie de torsion.
Soit $K \in D^+(Y_\etale)$. Alors
$$
Rf_*E \otimes_\mathbf{Z}^\mathbf{L} K =
Rf_*(E \otimes_\mathbf{Z}^\mathbf{L} f^{-1}K)
$$
dans $D^+(Y_\etale)$.
\end{lemma}

\begin{proof}
Un morphisme canonique du membre de gauche vers celui de droite est fourni par la
section \ref{sites-cohomology-section-projection-formula} de Cohomologie sur les sites.
Nous allons vérifier l'égalité sur les fibres.
Rappelons que le calcul des produits tensoriels dérivés commute aux images inverses.
Voir le lemme de Cohomologie sur les sites
\ref{sites-cohomology-lemma-pullback-tensor-product}.
Nous avons donc
$$
(E \otimes_\mathbf{Z}^\mathbf{L} f^{-1}K)_{\overline{x}}
=
E_{\overline{x}} \otimes_\mathbf{Z}^\mathbf{L} K_{\overline{y}}
$$
où $\overline{y}$ est l'image de $\overline{x}$ dans $Y$.
Puisque $\mathbf{Z}$ est de dimension globale $1$, nous voyons que
la cohomologie de ce complexe est nulle en degrés $< - 1 + a + b$
si $H^i(E) = 0$ pour $i \geq a$ et $H^j(K) = 0$ pour $j \geq b$.
De plus, puisque $H^i(E)$ est un faisceau abélien
de torsion pour tout $i$, il en va de même des faisceaux de cohomologie
du complexe $E \otimes_\mathbf{Z}^\mathbf{L} K$.
En effet, nous avons
$$
(E \otimes_\mathbf{Z}^\mathbf{L} f^{-1}K)
\otimes_{\mathbf{Z}}^\mathbf{L} \mathbf{Q} =
(E \otimes_\mathbf{Z}^\mathbf{L} \mathbf{Q})
\otimes_{\mathbf{Q}}^\mathbf{L}
(f^{-1}K \otimes_{\mathbf{Z}}^\mathbf{L} \mathbf{Q})
$$
qui est nul dans la catégorie dérivée.
Le lemme \ref{lemma-proper-base-change-stalk} s'applique donc
aux deux membres, et il suffit de montrer
$$
R\Gamma(X_{\overline{y}},
E|_{X_{\overline{y}}} \otimes_\mathbf{Z}^\mathbf{L}
(X_{\overline{y}} \to \overline{y})^{-1}K_{\overline{y}}) =
R\Gamma(X_{\overline{y}},
E|_{X_{\overline{y}}}) \otimes_\mathbf{Z}^\mathbf{L} K_{\overline{y}}
$$
Cela résulte du lemme \ref{lemma-pull-out-constant}.
\end{proof}








\section{Acyclicité locale}
\label{section-local-acyclicity}

\noindent
Dans cette section, nous déduisons l'acyclicité locale des morphismes lisses
du théorème de changement de base lisse. Dans SGA 4 ou SGA 4.5, les auteurs
démontrent d'abord une version de l'acyclicité locale pour les morphismes lisses,
puis en déduisent le théorème de changement de base lisse.

\medskip\noindent
Nous utiliserons la formulation de l'acyclicité locale donnée par
Deligne \cite[définition 2.12, page 242]{SGA4.5}.
Soit $f : X \to S$ un morphisme de schémas.
Soit $\overline{x}$ un point géométrique de $X$, d'image
$\overline{s} = f(\overline{x})$ dans $S$.
Soit $\overline{t}$ un point géométrique de
$\Spec(\mathcal{O}^{sh}_{S, \overline{s}})$.
Nous obtenons un diagramme commutatif
$$
\xymatrix{
F_{\overline{x}, \overline{t}} =
\overline{t} \times_{\Spec(\mathcal{O}^{sh}_{S, \overline{s}})}
\Spec(\mathcal{O}^{sh}_{X, \overline{x}}) \ar[r] \ar[d] &
\Spec(\mathcal{O}^{sh}_{X, \overline{x}}) \ar[r] \ar[d] &
X \ar[d] \\
\overline{t} \ar[r] &
\Spec(\mathcal{O}^{sh}_{S, \overline{s}}) \ar[r] &
S
}
$$
Le schéma $F_{\overline{x}, \overline{t}}$ est appelé
une {\it variété de cycles évanescents de $f$ en $\overline{x}$}.
Soit $K$ un objet de $D(X_\etale)$. Pour tout morphisme de schémas
$g : Y\to X$, nous écrivons $R\Gamma(Y, K)$ au lieu de
$R\Gamma(Y_\etale, g_{small}^{-1}K)$. Puisque
$\mathcal{O}^{sh}_{X, \overline{x}}$ est strictement hensélien,
nous avons
$K_{\overline{x}} = R\Gamma(\Spec(\mathcal{O}^{sh}_{X, \overline{x}}), K)$.
Nous obtenons donc un morphisme canonique
\begin{equation}
\label{equation-alpha-K}
\alpha_{K, \overline{x}, \overline{t}} :
K_{\overline{x}}
\longrightarrow
R\Gamma(F_{\overline{x}, \overline{t}}, K)
\end{equation}
en prenant l'image inverse en cohomologie le long de
$F_{\overline{x}, \overline{t}} \to \Spec(\mathcal{O}^{sh}_{X, \overline{x}})$.

\begin{definition}
\label{definition-locally-acyclic}
\begin{reference}
\cite[définition 2.12, page 242]{SGA4.5} et
\cite[définition (1.3), page 54]{SGA4.5}
\end{reference}
Soit $f : X \to S$ un morphisme de schémas.
Soit $K$ un objet de $D(X_\etale)$.
\begin{enumerate}
\item
Soit $\overline{x}$ un point géométrique de $X$, d'image
$\overline{s} = f(\overline{x})$.
Nous dirons que $f$ est {\it localement acyclique en $\overline{x}$ relativement à $K$}
si, pour tout point géométrique $\overline{t}$ de
$\Spec(\mathcal{O}^{sh}_{S, \overline{s}})$, le
morphisme (\ref{equation-alpha-K}) est un isomorphisme\footnote{Nous ne
supposons pas que $\overline{t}$ soit un point géométrique algébrique de
$\Spec(\mathcal{O}^{sh}_{S, \overline{s}})$. Le
lemme \ref{lemma-smooth-base-change-separably-closed} permet souvent
de se ramener à ce cas.}.
\item Nous dirons que $f$ est {\it localement acyclique relativement à $K$}
si $f$ est localement acyclique en $\overline{x}$ relativement à $K$
pour tout point géométrique $\overline{x}$ de $X$.
\item Nous dirons que $f$ est {\it universellement localement acyclique relativement à $K$}
si, pour tout morphisme de schémas $S' \to S$, le changement de base $f' : X' \to S'$
est localement acyclique relativement à l'image inverse de $K$ sur $X'$.
\item Nous dirons que $f$ est {\it localement acyclique} si, pour tout point
géométrique $\overline{x}$ de $X$ et tout entier $n$ premier à la caractéristique
de $\kappa(\overline{x})$, le morphisme $f$ est localement acyclique
en $\overline{x}$ relativement au faisceau constant de valeur
$\mathbf{Z}/n\mathbf{Z}$.
\item Nous dirons que $f$ est {\it universellement localement acyclique} si,
pour tout morphisme de schémas $S' \to S$, le changement de base $f' : X' \to S'$
est localement acyclique.
\end{enumerate}
\end{definition}

\noindent
Soit $M$ un groupe abélien. Alors l'acyclicité locale de $f : X \to S$
relativement au faisceau constant $\underline{M}$ se ramène
à la condition
$$
H^q(F_{\overline{x}, \overline{t}}, \underline{M}) =
\left\{
\begin{matrix}
M & \text{si} & q = 0 \\
0 & \text{si} & q \not = 0
\end{matrix}
\right.
$$
pour tout point géométrique $\overline{x}$ de $X$ et tout point
géométrique $\overline{t}$ de $\Spec(\mathcal{O}^{sh}_{S, f(\overline{x})})$.
Nous voyons ainsi qu'être localement acyclique correspond à
l'annulation des groupes de cohomologie supérieure des fibres géométriques
$F_{\overline{x}, \overline{t}}$ des morphismes entre les
hensélisés stricts en $\overline{x}$ et $\overline{s}$.

\begin{proposition}
\label{proposition-smooth-locally-acyclic}
Soit $f : X \to S$ un morphisme lisse de schémas.
Alors $f$ est universellement localement acyclique.
\end{proposition}

\begin{proof}
Puisque le changement de base d'un morphisme lisse est lisse, il suffit
de montrer que les morphismes lisses sont localement acycliques.
Soit $\overline{x}$ un point géométrique de $X$, d'image
$\overline{s} = f(\overline{x})$. Soit
$\overline{t}$ un point géométrique de
$\Spec(\mathcal{O}^{sh}_{S, f(\overline{x})})$.
Puisque nous cherchons à démontrer une propriété du morphisme d'anneaux
$\mathcal{O}^{sh}_{S, \overline{s}} \to \mathcal{O}^{sh}_{X, \overline{x}}$
(voir la discussion qui suit la définition \ref{definition-locally-acyclic}),
nous pouvons remplacer $f : X \to S$ par le changement de base
$X \times_S \Spec(\mathcal{O}^{sh}_{S, \overline{s}}) \to
\Spec(\mathcal{O}^{sh}_{S, \overline{s}})$.
Nous pouvons donc supposer que $S$ est le spectre d'un anneau local
strictement hensélien et que $\overline{s}$ est au-dessus
du point fermé de $S$.

\medskip\noindent
Nous allons appliquer le lemme \ref{lemma-base-change-f-star-general-stalks}
au diagramme
$$
\xymatrix{
X \ar[d]_f & X_{\overline{t}} \ar[l]^h \ar[d]^e \\
S & \overline{t} \ar[l]_g
}
$$
et au faisceau $\mathcal{F} = \underline{M}$, où $M = \mathbf{Z}/n\mathbf{Z}$
pour un certain entier $n$ premier à la caractéristique du corps résiduel de
$\overline{x}$. Nous savons que le morphisme
$f^{-1}R^qg_*\mathcal{F} \to R^qh_*e^{-1}\mathcal{F}$
est un isomorphisme par changement de base lisse; voir le
théorème \ref{theorem-smooth-base-change} (l'hypothèse
sur la torsion est satisfaite par notre choix de $n$).
Le lemme \ref{lemma-base-change-f-star-general-stalks}
nous donne donc l'égalité médiane dans
$$
H^q(F_{\overline{x}, \overline{t}}, \underline{M}) =
H^q(\Spec(\mathcal{O}^{sh}_{X, \overline{x}}) \times_S \overline{t},
\underline{M})
=
H^q(\Spec(\mathcal{O}^{sh}_{S, \overline{s}}) \times_S \overline{t},
\underline{M}) =
H^q(\overline{t}, \underline{M})
$$
Pour les deux égalités extérieures, nous utilisons le fait que
$S = \Spec(\mathcal{O}^{sh}_{S, \overline{s}})$.
Puisque $\overline{t}$ est le spectre d'un corps séparablement
clos, nous concluons que
$$
H^q(F_{\overline{x}, \overline{t}}, \underline{M}) =
\left\{
\begin{matrix}
M & \text{si} & q = 0 \\
0 & \text{si} & q \not = 0
\end{matrix}
\right.
$$
ce qu'il fallait montrer (voir la discussion qui suit la
définition \ref{definition-locally-acyclic}).
\end{proof}

\begin{lemma}
\label{lemma-locally-acyclic-locally-constant}
Soit $f : X \to S$ un morphisme de schémas. Soit $\mathcal{F}$ un
faisceau abélien localement constant sur $X_\etale$ tel que, pour tout point
géométrique $\overline{x}$ de $X$, le groupe abélien
$\mathcal{F}_{\overline{x}}$ soit de torsion et que chacun de ses éléments soit
d'ordre premier à la caractéristique du
corps résiduel de $\overline{x}$. Si $f$ est localement acyclique, alors $f$ est
localement acyclique relativement à $\mathcal{F}$.
\end{lemma}

\begin{proof}
En effet, soit $\overline{x}$ un point géométrique de $X$.
Puisque $\mathcal{F}$ est localement constant, nous voyons que
la restriction de $\mathcal{F}$ à $\Spec(\mathcal{O}^{sh}_{X, \overline{x}})$
est isomorphe au faisceau constant $\underline{M}$, où
$M = \mathcal{F}_{\overline{x}}$. Par hypothèse, nous pouvons écrire
$M = \colim M_i$ comme une limite inductive filtrante de groupes abéliens finis
$M_i$ d'ordre premier à la caractéristique du corps résiduel de
$\overline{x}$. Considérons un point géométrique $\overline{t}$ de
$\Spec(\mathcal{O}^{sh}_{S, f(\overline{x})})$.
Puisque $F_{\overline{x}, \overline{t}}$ est affine, nous avons
$$
H^q(F_{\overline{x}, \overline{t}}, \underline{M}) =
\colim H^q(F_{\overline{x}, \overline{t}}, \underline{M_i})
$$
par le lemme \ref{lemma-colimit}.
Pour tout $i$, nous pouvons écrire $M_i = \bigoplus \mathbf{Z}/n_{i, j}\mathbf{Z}$
comme une somme directe finie, pour certains entiers $n_{i, j}$ premiers à la
caractéristique du corps résiduel de $\overline{x}$.
Puisque $f$ est localement acyclique, nous voyons que
$$
H^q(F_{\overline{x}, \overline{t}},
\underline{\mathbf{Z}/n_{i, j}\mathbf{Z}}) =
\left\{
\begin{matrix}
\mathbf{Z}/n_{i, j}\mathbf{Z} & \text{si} & q = 0 \\
0 & \text{si} & q \not = 0
\end{matrix}
\right.
$$
Voir la discussion qui suit la définition \ref{definition-locally-acyclic}.
En prenant les sommes directes puis la limite inductive, nous concluons que
$$
H^q(F_{\overline{x}, \overline{t}}, \underline{M}) =
\left\{
\begin{matrix}
M & \text{si} & q = 0 \\
0 & \text{si} & q \not = 0
\end{matrix}
\right.
$$
ce qui achève la démonstration.
\end{proof}

\begin{lemma}
\label{lemma-locally-acyclic-quasi-finite-base-change}
Soit
$$
\xymatrix{
X' \ar[r]_{g'} \ar[d]_{f'} & X \ar[d]^f \\
S' \ar[r]^g & S
}
$$
un diagramme cartésien de schémas. Soit $K$ un objet de $D(X_\etale)$.
Soit $\overline{x}'$ un point géométrique de $X'$, d'image $\overline{x}$
dans $X$. Si :
\begin{enumerate}
\item $f$ est localement acyclique en $\overline{x}$ relativement à $K$;
\item $g$ est localement quasi-fini, ou bien $S' = \lim S_i$
est une limite projective filtrante de schémas localement quasi-finis sur $S$
à morphismes de transition affines, ou bien $g : S' \to S$ est entier,
\end{enumerate}
alors $f'$ est localement acyclique en $\overline{x}'$ relativement à $(g')^{-1}K$.
\end{lemma}

\begin{proof}
Notons $\overline{s}'$ et $\overline{s}$ les images de
$\overline{x}'$ et $\overline{x}$ dans $S'$ et $S$.
Soit $\overline{t}'$ un point géométrique de
$\Spec(\mathcal{O}^{sh}_{S', \overline{s}'})$, et notons
$\overline{t}$ son image dans $\Spec(\mathcal{O}^{sh}_{S, \overline{s}})$.
Par le lemme d'Algèbre
\ref{algebra-lemma-base-change-strict-henselization-quasi-finite}
et nos hypothèses sur $g$, nous avons un isomorphisme
$$
\mathcal{O}^{sh}_{X, \overline{x}}
\otimes_{\mathcal{O}^{sh}_{S, \overline{s}}}
\mathcal{O}^{sh}_{S', \overline{s}'}
\longrightarrow
\mathcal{O}^{sh}_{X', \overline{x}'}
$$
Puisque, d'après nos conventions,
$\kappa(\overline{t}) = \kappa(\overline{t}')$,
nous concluons que
$$
F_{\overline{x}', \overline{t}'} =
\Spec\left(
\mathcal{O}^{sh}_{X', \overline{x}'}
\otimes_{\mathcal{O}^{sh}_{S', \overline{s}'}}
\kappa(\overline{t}')\right) =
\Spec\left(
\mathcal{O}^{sh}_{X, \overline{x}}
\otimes_{\mathcal{O}^{sh}_{S, \overline{s}}}
\kappa(\overline{t})\right) =
F_{\overline{x}, \overline{t}}
$$
Autrement dit, les variétés des cycles évanescents
de $f'$ en $\overline{x}'$ sont des exemples de variétés des cycles
évanescents de $f$ en $\overline{x}$. Le lemme découle
immédiatement de cela et des définitions.
\end{proof}










\section{Le morphisme de cospécialisation}
\label{section-cospecialization}

\noindent
Soit $f : X \to S$ un morphisme de schémas. Soit $\overline{x}$ un
point géométrique de $X$, d'image $\overline{s} = f(\overline{x})$ dans $S$.
Soit $\overline{t}$ un point géométrique de
$\Spec(\mathcal{O}^{sh}_{S, \overline{s}})$.
Soit $K \in D(X_\etale)$. Pour tout morphisme de schémas $g : Y \to X$,
nous écrivons $K|_Y$ au lieu de $g_{small}^{-1}K$ et $R\Gamma(Y, K)$
au lieu de $R\Gamma(Y_\etale, g_{small}^{-1}K)$.
Nous affirmons que si :
\begin{enumerate}
\item $K$ est borné inférieurement, c'est-à-dire $K \in D^+(X_\etale)$;
\item $f$ est localement acyclique relativement à $K$,
\end{enumerate}
alors il existe un {\it morphisme de cospécialisation}
$$
cosp :
R\Gamma(X_{\overline{t}}, K)
\longrightarrow
R\Gamma(X_{\overline{s}}, K)
$$
qui sera étroitement lié au morphisme de spécialisation
considéré dans la section \ref{section-specialization}, et en particulier
dans la remarque \ref{remark-specialization-map-and-fibres}.

\medskip\noindent
Pour construire ce morphisme, considérons les morphismes
$$
X_{\overline{t}}
\xrightarrow{h}
X \times_S \Spec(\mathcal{O}^{sh}_{S, \overline{s}})
\xleftarrow{i}
X_{\overline{s}}
$$
L'unité de l'adjonction entre $h^{-1}$ et $Rh_*$ donne un morphisme
$$
\beta_{K, \overline{s}, \overline{t}} :
K|_{X \times_S \Spec(\mathcal{O}^{sh}_{S, \overline{s}})}
\longrightarrow
Rh_*(K|_{X_{\overline{t}}})
$$
dans $D((X \times_S \Spec(\mathcal{O}^{sh}_{S, \overline{s}}))_\etale)$.
Le lemme \ref{lemma-beta-is-isomorphism} ci-dessous montre que l'image inverse
$i^{-1}\beta_{K, \overline{s}, \overline{t}}$
est un isomorphisme sous les hypothèses ci-dessus.
Nous pouvons donc définir le morphisme de cospécialisation comme la composée
\begin{align*}
R\Gamma(X_{\overline{t}}, K)
& =
R\Gamma(X \times_S \Spec(\mathcal{O}^{sh}_{S, \overline{s}}),
Rh_*(K|_{X_{\overline{t}}})) \\
&
\xrightarrow{i^{-1}}
R\Gamma(X_{\overline{s}}, i^{-1}Rh_*(K|_{X_{\overline{t}}})) \\
&
\xrightarrow{(i^{-1}\beta_{K, \overline{s}, \overline{t}})^{-1}}
R\Gamma(X_{\overline{s}},
i^{-1}(K|_{X \times_S \Spec(\mathcal{O}^{sh}_{S, \overline{s}})})) \\
& =
R\Gamma(X_{\overline{s}}, K)
\end{align*}

\begin{lemma}
\label{lemma-beta-is-isomorphism}
Le morphisme $i^{-1}\beta_{K, \overline{s}, \overline{t}}$ est un isomorphisme.
\end{lemma}

\begin{proof}
La construction des morphismes $h$, $i$ et $\beta_{K, \overline{s}, \overline{t}}$
ne dépend que du
changement de base de $X$ et $K$ à $\Spec(\mathcal{O}^{sh}_{S, \overline{s}})$.
Nous pouvons donc supposer que $S$ est un schéma strictement hensélien
de point fermé $\overline{s}$. Remarquons que l'acyclicité locale
de $f$ relativement à $K$ est préservée par ce changement de base (par exemple
par le lemme \ref{lemma-locally-acyclic-quasi-finite-base-change}, ou simplement
en comparant directement les anneaux strictement henséliens dans ce cas très particulier).

\medskip\noindent
Soit $\overline{x}$ un point géométrique de $X_{\overline{s}}$.
De manière équivalente, soit $\overline{x}$ un point géométrique dont l'image
par $f$ est $\overline{s}$. Calculons la fibre de
$i^{-1}\beta_{K, \overline{s}, \overline{t}}$
en $\overline{x}$. Tout d'abord, nous avons
$$
(i^{-1}\beta_{K, \overline{s}, \overline{t}})_{\overline{x}} =
(\beta_{K, \overline{s}, \overline{t}})_{\overline{x}}
$$
car l'image inverse préserve les fibres; voir le lemme \ref{lemma-stalk-pullback}.
Puisque nous sommes dans la situation $S = \Spec(\mathcal{O}^{sh}_{S, \overline{s}})$,
nous voyons que $h : X_{\overline{t}} \to X$ possède la propriété
$X_{\overline{t}} \times_X \Spec(\mathcal{O}^{sh}_{X, \overline{x}})
= F_{\overline{x}, \overline{t}}$. Nous voyons donc que
$$
(\beta_{K, \overline{s}, \overline{t}})_{\overline{x}} :
K_{\overline{x}}
\longrightarrow
Rh_*(K|_{X_{\overline{t}}})_{\overline{x}} =
R\Gamma(F_{\overline{x}, \overline{t}}, K)
$$
où l'égalité est donnée par le théorème \ref{theorem-higher-direct-images}.
Il s'ensuit que le morphisme $(\beta_{K, \overline{s}, \overline{t}})_{\overline{x}}$
n'est autre que le morphisme $\alpha_{K, \overline{x}, \overline{t}}$ employé
dans la définition \ref{definition-locally-acyclic}. Le résultat en découle, car nous pouvons
vérifier sur les fibres qu'un morphisme est un isomorphisme, par le
théorème \ref{theorem-exactness-stalks}.
\end{proof}

\noindent
Le morphisme de cospécialisation, lorsqu'il existe, est destiné à être l'inverse
du morphisme de spécialisation.

\begin{lemma}
\label{lemma-specialization-cospecialization}
Dans la situation ci-dessus, si de plus
$f$ est quasi-compact et quasi-séparé, alors le diagramme
$$
\xymatrix{
(Rf_*K)_{\overline{s}} \ar[r] \ar[d]_{sp} &
R\Gamma(X_{\overline{s}}, K) \\
(Rf_*K)_{\overline{t}} \ar[r] &
R\Gamma(X_{\overline{t}}, K) \ar[u]_{cosp}
}
$$
est commutatif.
\end{lemma}

\begin{proof}
Comme dans la démonstration du lemme \ref{lemma-beta-is-isomorphism},
nous pouvons remplacer $S$ par $\Spec(\mathcal{O}^{sh}_{S, \overline{s}})$.
Nos morphismes se simplifient alors en $h : X_{\overline{t}} \to X$,
$i : X_{\overline{s}} \to X$ et
$\beta_{K, \overline{s}, \overline{t}} : K \to Rh_*(K|_{X_{\overline{t}}})$.
En utilisant l'égalité $(Rf_*K)_{\overline{s}} = R\Gamma(X, K)$ donnée par le
théorème \ref{theorem-higher-direct-images},
la composée de $sp$ avec le morphisme de changement de base
$(Rf_*K)_{\overline{t}} \to R\Gamma(X_{\overline{t}}, K)$
n'est autre que l'image inverse en cohomologie le long de $h$.
C'est le même morphisme que
$$
R\Gamma(X, K) \xrightarrow{\beta_{K, \overline{s}, \overline{t}}}
R\Gamma(X, Rh_*(K|_{X_{\overline{t}}})) =
R\Gamma(X_{\overline{t}}, K)
$$
Or le morphisme $cosp$ inverse d'abord le signe $=$ dans cette
formule, puis prend l'image inverse le long de $i$, et enfin applique
l'inverse de $i^{-1}\beta_{K, \overline{s}, \overline{t}}$.
Nous obtenons donc la commutativité voulue.
\end{proof}

\begin{lemma}
\label{lemma-sp-isom-proper-torsion-loc-ac}
Soit $f : X \to S$ un morphisme de schémas. Soit $K \in D(X_\etale)$.
Supposons que :
\begin{enumerate}
\item $K$ est borné inférieurement, c'est-à-dire $K \in D^+(X_\etale)$;
\item $f$ est localement acyclique relativement à $K$;
\item $f$ est propre;
\item $K$ a des faisceaux de cohomologie de torsion.
\end{enumerate}
Alors, pour tout point géométrique $\overline{s}$ de $S$ et tout point
géométrique $\overline{t}$ de $\Spec(\mathcal{O}^{sh}_{S, \overline{s}})$,
le morphisme de spécialisation
$sp : (Rf_*K)_{\overline{s}} \to (Rf_*K)_{\overline{t}}$
et le morphisme de cospécialisation
$cosp : R\Gamma(X_{\overline{t}}, K) \to R\Gamma(X_{\overline{s}}, K)$
sont des isomorphismes.
\end{lemma}

\begin{proof}
Par le théorème de changement de base propre (sous la forme du
lemme \ref{lemma-proper-base-change-stalk}), nous avons
$(Rf_*K)_{\overline{s}} = R\Gamma(X_{\overline{s}}, K)$
et de même pour $\overline{t}$. La démonstration « correcte » consisterait
à montrer que l'argument du
lemme \ref{lemma-specialization-cospecialization}
établit que $sp$ et $cosp$ sont des isomorphismes inverses dans ce cas.
Nous montrerons plutôt directement que $cosp$ est un isomorphisme.
D'après la discussion ci-dessus, nous voyons que $cosp$ est un isomorphisme
si et seulement si l'image inverse par $i$
$$
R\Gamma(X \times_S \Spec(\mathcal{O}^{sh}_{S, \overline{s}}),
Rh_*(K|_{X_{\overline{t}}}))
\longrightarrow
R\Gamma(X_{\overline{s}}, i^{-1}Rh_*(K|_{X_{\overline{t}}}))
$$
est un isomorphisme dans $D^+(\textit{Ab})$. Cela résulte du théorème de
changement de base propre appliqué au morphisme propre
$f' : X \times_S \Spec(\mathcal{O}^{sh}_{S, \overline{s}}) \to
\Spec(\mathcal{O}^{sh}_{S, \overline{s}})$
et au morphisme $\overline{s} \to \Spec(\mathcal{O}^{sh}_{S, \overline{s}})$,
avec le complexe $K' = Rh_*(K|_{X_{\overline{t}}})$. Le complexe
$K'$ est borné inférieurement et a des faisceaux de cohomologie de torsion
par le lemme \ref{lemma-torsion-direct-image}.
Puisque $\Spec(\mathcal{O}^{sh}_{S, \overline{s}})$ est strictement hensélien,
avec $\overline{s}$ au-dessus du point fermé,
nous voyons que la source de la flèche affichée est égale à
$(Rf'_*K')_{\overline{s}}$, que la cible est égale à
$R\Gamma(X_{\overline{s}}, K')$ et que le morphisme affiché est un isomorphisme
par le lemme déjà utilisé
\ref{lemma-proper-base-change-stalk}.
Nous voyons donc que trois des quatre flèches du diagramme
du lemme \ref{lemma-specialization-cospecialization} sont des isomorphismes,
et nous concluons.
\end{proof}

\begin{lemma}
\label{lemma-sp-isom-proper-loc-cst-torsion}
Soit $f : X \to S$ un morphisme de schémas. Soit $\mathcal{F}$
un faisceau abélien sur $X_\etale$. Supposons que :
\begin{enumerate}
\item $f$ est lisse et propre;
\item $\mathcal{F}$ est localement constant;
\item $\mathcal{F}_{\overline{x}}$ est un groupe de torsion dont
tous les éléments sont d'ordre premier à la caractéristique résiduelle de
$\overline{x}$, pour tout point géométrique $\overline{x}$ de $X$.
\end{enumerate}
Alors, pour tout point géométrique $\overline{s}$ de $S$ et tout point
géométrique $\overline{t}$ de $\Spec(\mathcal{O}^{sh}_{S, \overline{s}})$,
le morphisme de spécialisation
$sp : (Rf_*\mathcal{F})_{\overline{s}} \to (Rf_*\mathcal{F})_{\overline{t}}$
est un isomorphisme.
\end{lemma}

\begin{proof}
Cela découle des lemmes \ref{lemma-sp-isom-proper-torsion-loc-ac}
et \ref{lemma-locally-acyclic-locally-constant}, ainsi que de la
proposition \ref{proposition-smooth-locally-acyclic}.
\end{proof}

















\section{Dimension cohomologique}
\label{section-cd}

\noindent
Nous pouvons déduire certaines bornes sur la dimension cohomologique des
schémas et sur la dimension cohomologique des corps en utilisant
les résultats de la section \ref{section-vanishing-torsion}
et une application, apparemment anodine, du théorème de changement
de base propre (dans la démonstration de la proposition \ref{proposition-cd-affine}).

\begin{definition}
\label{definition-cd}
Soit $X$ un schéma quasi-compact et quasi-séparé.
La {\it dimension cohomologique de $X$} est le plus petit
élément
$$
\text{cd}(X) \in \{0, 1, 2, \ldots\} \cup \{\infty\}
$$
tel que, pour tout faisceau abélien de torsion $\mathcal{F}$
sur $X_\etale$, nous ayons $H^i_\etale(X, \mathcal{F}) = 0$
pour $i > \text{cd}(X)$. Si $X = \Spec(A)$, nous l'appelons parfois
la dimension cohomologique de $A$.
\end{definition}

\noindent
Si le schéma est de caractéristique $p$, nous pouvons souvent
obtenir des bornes plus fines pour l'annulation de la cohomologie des
faisceaux de $p$-torsion. Nous traiterons cela ailleurs
(référence à ajouter ultérieurement).

\begin{lemma}
\label{lemma-cd-limit}
Soit $X = \lim X_i$ la limite projective d'un système filtrant de
schémas quasi-compacts et quasi-séparés dont les morphismes de
transition sont affines. Alors $\text{cd}(X) \leq \max \text{cd}(X_i)$.
\end{lemma}

\begin{proof}
Notons $f_i : X \to X_i$ les projections.
Soit $\mathcal{F}$ un faisceau abélien de torsion sur $X_\etale$.
Nous avons alors $\mathcal{F} = \lim f_i^{-1}f_{i, *}\mathcal{F}$
par le lemme \ref{lemma-linus-hamann}.
Ainsi, $H^q_\etale(X, \mathcal{F}) =
\colim H^q_\etale(X_i, f_{i, *}\mathcal{F})$
par le théorème \ref{theorem-colimit}.
Le lemme en découle.
\end{proof}

\begin{lemma}
\label{lemma-cd-curve-over-field}
Soit $K$ un corps. Soit $X$ un schéma affine de dimension $1$
et de type fini sur $K$. Alors
$\text{cd}(X) \leq 1 + \text{cd}(K)$.
\end{lemma}

\begin{proof}
Soit $\mathcal{F}$ un faisceau abélien de torsion sur $X_\etale$.
Considérons la suite spectrale de Leray du morphisme
$f : X \to \Spec(K)$. Nous obtenons
$$
E_2^{p, q} = H^p(\Spec(K), R^qf_*\mathcal{F})
$$
qui converge vers $H^{p + q}_\etale(X, \mathcal{F})$.
La fibre de $R^qf_*\mathcal{F}$ en un point géométrique
$\Spec(\overline{K}) \to \Spec(K)$ est la cohomologie de l'image
inverse de $\mathcal{F}$ sur $X_{\overline{K}}$.
Elle s'annule donc en degrés $\geq 2$ par le
théorème \ref{theorem-vanishing-affine-curves}.
\end{proof}

\begin{lemma}
\label{lemma-cd-field-extension}
Soit $L/K$ une extension de corps. Nous avons alors
$\text{cd}(L) \leq \text{cd}(K) + \text{trdeg}_K(L)$.
\end{lemma}

\begin{proof}
Si $\text{trdeg}_K(L) = \infty$, c'est clair.
Sinon, nous pouvons trouver une suite d'extensions
$L= L_r/L_{r - 1}/ \ldots /L_1/L_0 = K$ telle que
$\text{trdeg}_{L_i}(L_{i + 1}) = 1$ et $r = \text{trdeg}_K(L)$.
Il suffit donc de démontrer le lemme dans le cas $r = 1$.
Dans ce cas, nous pouvons écrire $L = \colim A_i$
comme limite inductive filtrante de ses sous-$K$-algèbres de type fini.
Par le lemme \ref{lemma-cd-limit}, il suffit de démontrer que
$\text{cd}(A_i) \leq 1 + \text{cd}(K)$. Cela découle
du lemme \ref{lemma-cd-curve-over-field}.
\end{proof}

\begin{lemma}
\label{lemma-strictly-henselian}
Soit $K$ un corps. Soit $X$ un schéma de type fini sur $K$.
Soit $x \in X$. Posons $a = \text{trdeg}_K(\kappa(x))$
et $d = \dim_x(X)$. Il existe alors un morphisme
$$
K(t_1, \ldots, t_a)^{sep} \longrightarrow \mathcal{O}_{X, x}^{sh}
$$
tel que
\begin{enumerate}
\item le corps résiduel de $\mathcal{O}_{X, x}^{sh}$ soit une extension purement inséparable
de $K(t_1, \ldots, t_a)^{sep}$,
\item $\mathcal{O}_{X, x}^{sh}$ soit une limite inductive filtrante de
$K(t_1, \ldots, t_a)^{sep}$-algèbres de type fini et de dimension $\leq d - a$.
\end{enumerate}
\end{lemma}

\begin{proof}
Nous pouvons supposer $X$ affine. Par normalisation de Noether, quitte à
rétrécir encore $X$, nous pouvons choisir un morphisme fini
$\pi : X \to \mathbf{A}^d_K$; voir le
lemme \ref{algebra-lemma-Noether-normalization-at-point} d'Algèbre.
Puisque $\kappa(x)$ est une extension finie du corps résiduel de $\pi(x)$,
ce corps résiduel est lui aussi de degré de transcendance $a$ sur $K$.
Nous pouvons donc trouver un morphisme fini $\pi' : \mathbf{A}^d_K \to \mathbf{A}^d_K$
tel que $\pi'(\pi(x))$ corresponde au point générique du sous-espace
linéaire $\mathbf{A}^a_K \subset \mathbf{A}^d_K$ défini en annulant
les $d - a$ dernières coordonnées. Par conséquent, le composé
$$
X \xrightarrow{\pi' \circ \pi} \mathbf{A}^d_K \xrightarrow{p} \mathbf{A}^a_K
$$
de $\pi' \circ \pi$ et de la projection $p$ sur les $a$ premières coordonnées
envoie $x$ sur le point générique $\eta \in \mathbf{A}^a_K$. Le morphisme induit
$$
K(t_1, \ldots, t_a)^{sep} =
\mathcal{O}_{\mathbf{A}^a_K, \eta}^{sh}
\longrightarrow \mathcal{O}_{X, x}^{sh}
$$
sur les anneaux locaux étales satisfait (1), car le corps résiduel
de $\mathcal{O}_{X, x}^{sh}$ est manifestement une extension algébrique du corps
séparablement clos $K(t_1, \ldots, t_a)^{sep}$. D'autre part, si $X = \Spec(B)$,
alors $\mathcal{O}_{X, x}^{sh} = \colim B_j$ est une limite inductive filtrante
de $B$-algèbres étales $B_j$. Remarquons que $B_j$ est quasi-finie sur
$K[t_1, \ldots, t_d]$, puisque $B$ est finie sur $K[t_1, \ldots, t_d]$.
Nous pouvons de même écrire
$K(t_1, \ldots, t_a)^{sep} = \colim A_i$ comme limite inductive filtrante
de $K[t_1, \ldots, t_a]$-algèbres étales. Pour tout $i$, nous pouvons
trouver un $j$ tel que
$A_i \to K(t_1, \ldots, t_a)^{sep} \to \mathcal{O}_{X, x}^{sh}$
se factorise par un morphisme $\psi_{i, j} : A_i \to B_j$. Alors $B_j$ est quasi-finie
sur $A_i[t_{a + 1}, \ldots, t_d]$. Par suite,
$$
B_{i, j} = B_j \otimes_{\psi_{i, j}, A_i} K(t_1, \ldots, t_a)^{sep}
$$
est de dimension $\leq d - a$, puisqu'elle est quasi-finie sur
$K(t_1, \ldots, t_a)^{sep}[t_{a + 1}, \ldots, t_d]$.
La démonstration de (2) est maintenant achevée, car $\mathcal{O}_{X, x}^{sh}$ est une
limite inductive filtrante\footnote{Soit $R$ un anneau.
Soit $A = \colim_{i \in I} A_i$
une limite inductive filtrante de $R$-algèbres de présentation finie.
Soit $B = \colim_{j \in J} B_j$ une limite inductive filtrante de $R$-algèbres.
Soit $A \to B$ un morphisme de $R$-algèbres.
Supposons que, pour tout $i \in I$, il existe $j \in J$ et un morphisme de $R$-algèbres
$\psi_{i, j} : A_i \to B_j$.
Posons $(i', j', \psi_{i', j'}) \geq (i, j, \psi_{i, j})$ si
$i' \geq i$, $j' \geq j$, et si $\psi_{i, j}$ et $\psi_{i', j'}$
sont compatibles. La collection des triplets forme alors un
ensemble filtrant et $B = \colim B_j \otimes_{\psi_{i, j}, A_i} A$.}
des algèbres $B_{i, j}$. Certains détails sont omis.
\end{proof}

\begin{proposition}
\label{proposition-cd-affine}
Soit $K$ un corps. Soit $X$ un schéma affine de type fini sur $K$.
Nous avons alors $\text{cd}(X) \leq \dim(X) + \text{cd}(K)$.
\end{proposition}

\begin{proof}
Nous allons le démontrer par récurrence sur $\dim(X)$.
Soit $\mathcal{F}$ un faisceau abélien de torsion sur $X_\etale$.

\medskip\noindent
Le cas $\dim(X) = 0$. Dans ce cas, le morphisme structural
$f : X \to \Spec(K)$ est fini. Nous voyons donc que $R^if_*\mathcal{F} = 0$
pour $i > 0$; voir la proposition \ref{proposition-finite-higher-direct-image-zero}.
Ainsi, $H^i_\etale(X, \mathcal{F}) = H^i_\etale(\Spec(K), f_*\mathcal{F})$
par la suite spectrale de Leray de $f$
(lemme \ref{sites-cohomology-lemma-Leray} de Cohomologie sur les sites),
et le résultat est clair.

\medskip\noindent
Le cas $\dim(X) = 1$. C'est le lemme \ref{lemma-cd-curve-over-field}.

\medskip\noindent
Supposons $d = \dim(X) > 1$ et la proposition vraie pour les schémas
affines de type fini et de dimension $< d$ sur des corps.
Par normalisation de Noether, voir par exemple le
lemme \ref{varieties-lemma-noether-normalization-affine} de Variétés,
il existe un morphisme fini $f : X \to \mathbf{A}^d_K$.
Rappelons que $R^if_*\mathcal{F} = 0$ pour $i > 0$ par la
proposition \ref{proposition-finite-higher-direct-image-zero}.
Par la suite spectrale de Leray de $f$
(lemme \ref{sites-cohomology-lemma-Leray} de Cohomologie sur les sites),
nous concluons qu'il suffit de démontrer le résultat pour $f_*\mathcal{F}$
sur $\mathbf{A}^d_K$.

\medskip\noindent
Interlude I. Soit $j : X \to Y$ une immersion ouverte de variétés lisses
de dimension $d$ sur $K$ (non nécessairement affines),
dont le complément est le support d'un diviseur de Cartier effectif $D$.
Les faisceaux $R^qj_*\mathcal{F}$ pour $q > 0$ sont supportés par $D$.
Nous affirmons que $(R^qj_*\mathcal{F})_{\overline{y}} = 0$
pour $a = \text{trdeg}_K(\kappa(y)) > d - q$. En effet, par le
théorème \ref{theorem-higher-direct-images}, nous avons
$$
(R^qj_*\mathcal{F})_{\overline{y}} =
H^q(\Spec(\mathcal{O}_{Y, y}^{sh}) \times_Y X, \mathcal{F})
$$
Choisissons une équation locale $f \in \mathfrak m_y \subset \mathcal{O}_{Y, y}$
de $D$. Nous avons alors
$$
\Spec(\mathcal{O}_{Y, y}^{sh}) \times_Y X =
\Spec(\mathcal{O}_{Y, y}^{sh}[1/f])
$$
À l'aide du lemme \ref{lemma-strictly-henselian}, nous obtenons un plongement
$$
K(t_1, \ldots, t_a)^{sep}(x) =
K(t_1, \ldots, t_a)^{sep}[x]_{(x)}[1/x]
\longrightarrow
\mathcal{O}_{Y, y}^{sh}[1/f]
$$
Puisque le degré de transcendance sur $K$ du corps des fractions de
$\mathcal{O}_{Y, y}^{sh}$ est $d$, nous voyons que $\mathcal{O}_{Y, y}^{sh}[1/f]$
est une limite inductive filtrante d'algèbres de type fini de dimension $(d - a - 1)$ sur
le corps $K(t_1, \ldots, t_a)^{sep}(x)$, qui est lui-même de
dimension cohomologique $1$ par le lemme \ref{lemma-cd-field-extension}. L'hypothèse
de récurrence et le lemme \ref{lemma-cd-limit}
donnent donc l'annulation voulue.

\medskip\noindent
Interlude II. Soit $Z$ une variété lisse sur $K$ de dimension $d - 1$.
Soit $E_a \subset Z$ l'ensemble des points $z \in Z$ tels que
$\text{trdeg}_K(\kappa(z)) \leq a$. Remarquons que $E_a$ est stable
par spécialisation; voir le
lemme \ref{varieties-lemma-dimension-locally-algebraic} de Variétés.
Supposons que $\mathcal{G}$ soit un faisceau abélien de torsion sur $Z$
dont le support est contenu dans $E_a$. Nous affirmons alors que
$H^b_\etale(Z, \mathcal{G}) = 0$ pour $b > a + \text{cd}(K)$.
En effet, nous pouvons écrire $\mathcal{G} = \colim \mathcal{G}_i$
avec $\mathcal{G}_i$ un faisceau abélien de torsion
supporté par un sous-schéma fermé $Z_i$ contenu dans $E_a$; voir le
lemme \ref{lemma-support-in-subset}.
L'hypothèse de récurrence implique alors
l'annulation voulue pour $\mathcal{G}_i$\footnote{Nous employons ici d'abord la
proposition \ref{proposition-closed-immersion-pushforward}
pour écrire $\mathcal{G}_i$ comme l'image directe d'un faisceau sur $Z_i$;
l'hypothèse de récurrence donne l'annulation de ce faisceau sur $Z_i$, et
la suite spectrale de Leray de $Z_i \to Z$ donne l'annulation
de $\mathcal{G}_i$.}. Enfin, nous
concluons par le théorème \ref{theorem-colimit}.

\medskip\noindent
Considérons le diagramme commutatif
$$
\xymatrix{
\mathbf{A}^d_K \ar[rd]_f \ar[rr]_-j & &
\mathbf{P}^1_K \times_K \mathbf{A}^{d - 1}_K \ar[ld]^g \\
& \mathbf{A}^{d - 1}_K
}
$$
Remarquons que $j$ est une immersion ouverte de variétés lisses de dimension $d$
dont le complément est un diviseur de Cartier effectif $D$. Nous
pouvons donc employer les résultats obtenus dans l'interlude I.
Nous allons étudier la suite spectrale de Leray relative
$$
E_2^{p, q} = R^pg_*R^qj_*\mathcal{F} \Rightarrow R^{p + q}f_*\mathcal{F}
$$
Comme $R^qj_*\mathcal{F}$ pour $q > 0$ est supporté par $D$ et que
$g|_D : D \to \mathbf{A}^{d - 1}_K$ est un isomorphisme, nous trouvons
$R^pg_*R^qj_*\mathcal{F} = 0$ pour $p > 0$ et $q > 0$. De plus, nous avons
$R^qj_*\mathcal{F} = 0$ pour $q > d$. D'autre part, $g$ est un
morphisme propre de dimension relative $1$. Ainsi, par le
lemme \ref{lemma-cohomological-dimension-proper},
nous voyons que $R^pg_*j_*\mathcal{F} = 0$ pour $p > 2$. La page $E_2$
de la suite spectrale a donc la forme suivante
$$
\begin{matrix}
g_*R^dj_*\mathcal{F} & 0 & 0 \\
\ldots & \ldots & \ldots \\
g_*R^2j_*\mathcal{F} & 0 & 0 \\
g_*R^1j_*\mathcal{F} & 0 & 0 \\
g_*j_*\mathcal{F} & R^1g_*j_*\mathcal{F} & R^2g_*j_*\mathcal{F}
\end{matrix}
$$
Nous concluons que $R^qf_*\mathcal{F} = g_*R^qj_*\mathcal{F}$ pour $q > 2$.
Par l'interlude I, nous voyons que le support de $R^qf_*\mathcal{F}$ pour $q > 2$
est contenu dans l'ensemble des points de $\mathbf{A}^{d - 1}_K$
dont le corps résiduel est de degré de transcendance $\leq d - q$.
Par l'interlude II,
$$
H^p(\mathbf{A}^{d - 1}_K, R^qf_*\mathcal{F}) = 0
\text{ pour }p > d - q + \text{cd}(K)\text{ et }q > 2
$$
 D'autre part, par le théorème \ref{theorem-higher-direct-images},
nous avons $R^2f_*\mathcal{F}_{\overline{\eta}} =
H^2(\mathbf{A}^1_{\overline{\eta}}, \mathcal{F}) = 0$
(annulation par le cas de dimension $1$), où $\eta$ est le
point générique de $\mathbf{A}^{d - 1}_K$.
Ainsi, de nouveau par l'interlude II, nous voyons que
$$
H^p(\mathbf{A}^{d - 1}_K, R^2f_*\mathcal{F}) = 0
\text{ pour }p > d - 2 + \text{cd}(K)
$$
Enfin, nous avons
$$
H^p(\mathbf{A}^{d - 1}_K, R^qf_*\mathcal{F}) = 0
\text{ pour }p > d - 1 + \text{cd}(K)\text{ et }q = 0, 1
$$
par l'hypothèse de récurrence. En combinant tout ce qui précède
avec la suite spectrale de Leray
$H^p(\mathbf{A}^{d - 1}_K, R^qf_*\mathcal{F}) \Rightarrow
H^{p + q}(\mathbf{A}^d_K, \mathcal{F})$, nous concluons.
\end{proof}

\begin{lemma}
\label{lemma-interlude-II}
Soit $K$ un corps. Soit $X$ un schéma affine de type fini sur $K$.
Soit $E_a \subset X$ l'ensemble des points
$x \in X$ tels que $\text{trdeg}_K(\kappa(x)) \leq a$.
Soit $\mathcal{F}$ un faisceau abélien de torsion sur $X_\etale$
dont le support est contenu dans $E_a$. Alors
$H^b_\etale(X, \mathcal{F}) = 0$ pour $b > a + \text{cd}(K)$.
\end{lemma}

\begin{proof}
Nous pouvons écrire $\mathcal{F} = \colim \mathcal{F}_i$
avec $\mathcal{F}_i$ un faisceau abélien de torsion
supporté par un sous-schéma fermé $Z_i$ contenu dans $E_a$; voir le
lemme \ref{lemma-support-in-subset}.
La proposition \ref{proposition-cd-affine} donne alors
l'annulation voulue pour $\mathcal{F}_i$. Nous omettons les détails;
indications : utiliser d'abord la proposition \ref{proposition-closed-immersion-pushforward}
pour écrire $\mathcal{F}_i$ comme l'image directe d'un faisceau sur $Z_i$,
employer l'annulation de ce faisceau sur $Z_i$, puis la
suite spectrale de Leray de $Z_i \to X$ pour obtenir l'annulation
de $\mathcal{F}_i$. Enfin, nous
concluons par le théorème \ref{theorem-colimit}.
\end{proof}

\begin{lemma}
\label{lemma-interlude-I}
Soit $f : X \to Y$ un morphisme affine de schémas de type
fini sur un corps $K$. Soit $E_a(X)$ l'ensemble des points $x \in X$
tels que $\text{trdeg}_K(\kappa(x)) \leq a$.
Soit $\mathcal{F}$ un faisceau abélien de torsion sur $X_\etale$
dont le support est contenu dans $E_a$. Alors
le support de $R^qf_*\mathcal{F}$ est contenu dans
$E_{a - q}(Y)$.
\end{lemma}

\begin{proof}
La question est locale sur $Y$; nous pouvons donc supposer $Y$ affine.
Alors $X$ est lui aussi affine, et nous pouvons choisir un diagramme
$$
\xymatrix{
X \ar[d]_f \ar[r]_i & \mathbf{A}^{n + m}_K \ar[d]^{\text{pr}} \\
Y \ar[r]^j & \mathbf{A}^n_K
}
$$
où les flèches horizontales sont des immersions fermées et la flèche
verticale de droite est la projection (détails omis).
Alors $j_*R^qf_*\mathcal{F} = R^q\text{pr}_*i_*\mathcal{F}$
par l'annulation des images directes supérieures de $i$ et $j$; voir la
proposition \ref{proposition-finite-higher-direct-image-zero}.
De plus, la description des fibres de $j_*$ dans cette proposition
montre qu'il suffit d'établir l'annulation pour $j_*R^qf_*\mathcal{F}$.
Nous pouvons donc supposer que $f$ est le morphisme de
projection $\text{pr} : \mathbf{A}^{n + m}_K \to \mathbf{A}^n_K$
et que $\mathcal{F}$ est un faisceau abélien de torsion sur $\mathbf{A}^{n + m}_K$
satisfaisant l'hypothèse de l'énoncé du lemme.

\medskip\noindent
Soit $y$ un point de $\mathbf{A}^n_K$.
Par le théorème \ref{theorem-higher-direct-images}, nous avons
$$
(R^q\text{pr}_*\mathcal{F})_{\overline{y}} =
H^q(\mathbf{A}^{n + m}_K \times_{\mathbf{A}^n_K} \Spec(\mathcal{O}_{Y, y}^{sh}),
\mathcal{F}) =
H^q(\mathbf{A}^m_{\mathcal{O}_{Y, y}^{sh}}, \mathcal{F})
$$
Posons $b = \text{trdeg}_K(\kappa(y))$. Le lemme \ref{lemma-strictly-henselian}
nous donne un plongement
$$
L = K(t_1, \ldots, t_b)^{sep} \longrightarrow \mathcal{O}_{Y, y}^{sh}
$$
Écrivons $\mathcal{O}_{Y, y}^{sh} = \colim B_i$ comme la limite
inductive filtrante de sous-$L$-algèbres de type fini $B_i \subset \mathcal{O}_{Y, y}^{sh}$
contenant l'anneau $K[T_1, \ldots, T_n]$ des fonctions régulières
sur $\mathbf{A}^n_K$. Nous obtenons alors
$$
\mathbf{A}^m_{\mathcal{O}_{Y, y}^{sh}} =
\lim \mathbf{A}^m_{B_i}
$$
Si $z \in \mathbf{A}^m_{B_i}$ est un point du support de
$\mathcal{F}$, alors l'image $x$ de $z$ dans $\mathbf{A}^{m + n}_K$
satisfait $\text{trdeg}_K(\kappa(x)) \leq a$ par l'hypothèse du lemme sur
$\mathcal{F}$.
Puisque $\mathcal{O}_{Y, y}^{sh}$ est une limite inductive filtrante
d'algèbres étales sur $K[T_1, \ldots, T_n]$ et que
$B_i \subset \mathcal{O}_{Y, y}^{sh}$
nous voyons que $\kappa(z)/\kappa(x)$ est algébrique
(certains détails sont omis). Alors $\text{trdeg}_K(\kappa(z)) \leq a$
et donc $\text{trdeg}_L(\kappa(z)) \leq a - b$.
Par le lemme \ref{lemma-interlude-II}, nous voyons que
$$
H^q(\mathbf{A}^m_{B_i}, \mathcal{F}) = 0\text{ pour }q > a - b
$$
Ainsi, par le théorème \ref{theorem-colimit}, nous obtenons
$(R^qf_*\mathcal{F})_{\overline{y}} = 0$ pour $q > a - b$,
comme voulu.
\end{proof}








\section{Dimension cohomologique finie}
\label{section-finite-cd}

\noindent
Nous poursuivons la discussion commencée dans la section \ref{section-cd}.

\begin{definition}
\label{definition-cd-f}
Soit $f : X \to Y$ un morphisme de schémas quasi-compact et
quasi-séparé.
La {\it dimension cohomologique de $f$} est le plus petit
élément
$$
\text{cd}(f) \in \{0, 1, 2, \ldots\} \cup \{\infty\}
$$
tel que, pour tout faisceau abélien de torsion $\mathcal{F}$
sur $X_\etale$, nous ayons $R^if_*\mathcal{F} = 0$
pour $i > \text{cd}(f)$.
\end{definition}

\begin{lemma}
\label{lemma-finite-cd}
Soit $K$ un corps.
\begin{enumerate}
\item Si $f : X \to Y$ est un morphisme de schémas de type fini sur $K$,
alors $\text{cd}(f) < \infty$.
\item Si $\text{cd}(K) < \infty$, alors $\text{cd}(X) < \infty$
pour tout schéma $X$ de type fini sur $K$.
\end{enumerate}
\end{lemma}

\begin{proof}
Démonstration de (1). Nous pouvons supposer $Y$ affine. Nous allons employer le
principe de récurrence du
lemme \ref{coherent-lemma-induction-principle} de Cohomologie des schémas
pour le démontrer. Si $X$ est lui aussi affine, le résultat vaut par le
lemme \ref{lemma-interlude-I}. Il suffit donc de montrer que, si
$X = U \cup V$ et si le résultat est vrai pour $U \to Y$, $V \to Y$ et
$U \cap V \to Y$, il est vrai pour $f$. Cela découle de la
suite de Mayer--Vietoris relative; voir le
lemme \ref{lemma-relative-mayer-vietoris}.

\medskip\noindent
Démonstration de (2). Nous allons employer le principe de récurrence du
lemme \ref{coherent-lemma-induction-principle} de Cohomologie des schémas
pour le démontrer. Si $X$ est affine, le résultat vaut par la
proposition \ref{proposition-cd-affine}. Il suffit donc de montrer que, si
$X = U \cup V$ et si le résultat est vrai pour $U$, $V$ et
$U \cap V $, il est vrai pour $X$. Cela découle de la
suite de Mayer--Vietoris; voir le
lemme \ref{lemma-mayer-vietoris}.
\end{proof}

\begin{lemma}
\label{lemma-finite-cd-mod-n-direct-sums}
Cohomologie et sommes directes. Soit $n \geq 1$ un entier.
\begin{enumerate}
\item Soit $f : X \to Y$ un morphisme de schémas quasi-compact et
quasi-séparé tel que $\text{cd}(f) < \infty$. Alors le foncteur
$$
Rf_* :
D(X_\etale, \mathbf{Z}/n\mathbf{Z})
\longrightarrow
D(Y_\etale, \mathbf{Z}/n\mathbf{Z})
$$
commute aux sommes directes.
\item Soit $X$ un schéma quasi-compact et quasi-séparé tel que
$\text{cd}(X) < \infty$. Alors le foncteur
$$
R\Gamma(X, -) :
D(X_\etale, \mathbf{Z}/n\mathbf{Z})
\longrightarrow
D(\mathbf{Z}/n\mathbf{Z})
$$
commute aux sommes directes.
\end{enumerate}
\end{lemma}

\begin{proof}
Démonstration de (1). Puisque $\text{cd}(f) < \infty$, nous voyons que
$$
f_* :
\textit{Mod}(X_\etale, \mathbf{Z}/n\mathbf{Z})
\longrightarrow
\textit{Mod}(Y_\etale, \mathbf{Z}/n\mathbf{Z})
$$
est de dimension cohomologique finie au sens du
lemme \ref{derived-lemma-unbounded-right-derived} de Catégories dérivées.
Soit $I$ un ensemble et, pour $i \in I$, soit $E_i$
un objet de $D(X_\etale, \mathbf{Z}/n\mathbf{Z})$.
Choisissons un complexe K-injectif $\mathcal{I}_i^\bullet$
de $\mathbf{Z}/n\mathbf{Z}$-modules dont chacun des
termes $\mathcal{I}_i^n$ est un faisceau injectif de
$\mathbf{Z}/n\mathbf{Z}$-modules et qui représente $E_i$.
Voir le théorème d'Injectifs
\ref{injectives-theorem-K-injective-embedding-grothendieck}.
Alors $\bigoplus E_i$ est représenté par le complexe
$\bigoplus \mathcal{I}_i^\bullet$ (somme directe terme à terme); voir le
lemme \ref{injectives-lemma-derived-products} d'Injectifs.
Par le lemme \ref{lemma-relative-colimit}, nous avons
$$
R^qf_*(\bigoplus \mathcal{I}_i^n) =
\bigoplus R^qf_*(\mathcal{I}_i^n) = 0
$$
pour $q > 0$ et tout $n$. Nous concluons donc par le
lemme \ref{derived-lemma-unbounded-right-derived} de Catégories dérivées
que nous pouvons calculer $Rf_*(\bigoplus E_i)$ par le complexe
$$
f_*(\bigoplus \mathcal{I}_i^\bullet) =
\bigoplus f_*(\mathcal{I}_i^\bullet)
$$
(égalité donnée de nouveau par le lemme \ref{lemma-relative-colimit}), qui
représente $\bigoplus Rf_*E_i$ par le lemme déjà employé
\ref{injectives-lemma-derived-products} d'Injectifs.

\medskip\noindent
Démonstration de (2). Elle est identique à celle de (1),
et nous l'omettons.
\end{proof}

\begin{lemma}
\label{lemma-proper-mod-n-direct-sums}
Soit $f : X \to Y$ un morphisme propre de schémas. Soit $n \geq 1$
un entier. Alors le foncteur
$$
Rf_* :
D(X_\etale, \mathbf{Z}/n\mathbf{Z})
\longrightarrow
D(Y_\etale, \mathbf{Z}/n\mathbf{Z})
$$
commute aux sommes directes.
\end{lemma}

\begin{proof}
Il suffit de le démontrer lorsque $Y$ est quasi-compact. Par le
lemme \ref{morphisms-lemma-morphism-finite-type-bounded-dimension} de Morphismes,
nous voyons que la dimension des fibres de
$f : X \to Y$ est bornée.
Le lemme \ref{lemma-cohomological-dimension-proper}
implique donc que $\text{cd}(f) < \infty$. Le résultat
découle alors du lemme \ref{lemma-finite-cd-mod-n-direct-sums}.
\end{proof}

\begin{lemma}
\label{lemma-pull-out-constant-mod-n}
Soit $X$ un schéma quasi-compact et quasi-séparé
tel que $\text{cd}(X) < \infty$. Soit $\Lambda$ un anneau de torsion.
Soient $E \in D(X_\etale, \Lambda)$ et $K \in D(\Lambda)$. Alors
$$
R\Gamma(X, E \otimes_\Lambda^\mathbf{L} \underline{K}) =
R\Gamma(X, E) \otimes_\Lambda^\mathbf{L} K
$$
\end{lemma}

\begin{proof}
Il existe un morphisme canonique de droite à gauche par la
section \ref{sites-cohomology-section-projection-formula} de Cohomologie sur les sites.
Soit $T(K)$ la propriété selon laquelle l'énoncé du
lemme vaut pour $K \in D(\Lambda)$.
Nous allons vérifier que les conditions (1), (2) et (3) de la
remarque \ref{more-algebra-remark-P-resolution} de Compléments d'algèbre
valent pour $T$, ce qui permettra de conclure.
La propriété (1) vaut parce que les deux membres de l'égalité
commutent aux sommes directes; voir le
lemme \ref{lemma-finite-cd-mod-n-direct-sums}.
La propriété (2) vaut parce que nous comparons des foncteurs exacts
entre catégories triangulées et que nous pouvons employer le
lemme \ref{derived-lemma-third-isomorphism-triangle} de Catégories dérivées.
La propriété (3) affirme que le lemme vaut lorsque $K = \Lambda[k]$
pour tout décalage $k \in \mathbf{Z}$, ce qui est évident.
\end{proof}

\begin{lemma}
\label{lemma-projection-formula-proper-mod-n}
Soit $f : X \to Y$ un morphisme propre de schémas. Soit $\Lambda$
un anneau de torsion. Soient $E \in D(X_\etale, \Lambda)$ et
$K \in D(Y_\etale, \Lambda)$. Alors
$$
Rf_*E \otimes_\Lambda^\mathbf{L} K =
Rf_*(E \otimes_\Lambda^\mathbf{L} f^{-1}K)
$$
dans $D(Y_\etale, \Lambda)$.
\end{lemma}

\begin{proof}
Il existe un morphisme canonique de gauche à droite par la
section \ref{sites-cohomology-section-projection-formula} de Cohomologie sur les sites.
Nous allons vérifier l'égalité sur les fibres en $\overline{y}$.
Par le changement de base propre (sous la forme du
lemme \ref{lemma-proper-base-change-mod-n}, où $Y' = \overline{y}$),
on se ramène au cas où $Y$ est le spectre d'un
corps algébriquement clos.
C'est le lemme \ref{lemma-pull-out-constant-mod-n},
où nous employons $\text{cd}(X) < \infty$, donné par le
lemme \ref{lemma-cohomological-dimension-proper}.
\end{proof}






\section{Formule de K\"unneth en cohomologie étale}
\label{section-kunneth}

\noindent
Nous démontrons d'abord une formule de K\"unneth lorsque l'un des facteurs
est propre. Nous l'employons ensuite pour démontrer une propriété de changement de base
pour les immersions ouvertes. Celle-ci donne alors un
« changement de base par des morphismes vers des spectres de corps »
(analogue au changement de base lisse).
Enfin, nous l'utilisons pour obtenir une formule de K\"unneth plus générale.

\begin{remark}
\label{remark-define-kunneth-map}
Considérons un diagramme cartésien dans la catégorie des schémas :
$$
\xymatrix{
X \times_S Y \ar[d]_p \ar[r]_q \ar[rd]_c & Y \ar[d]^g \\
X \ar[r]^f & S
}
$$
Soit $\Lambda$ un anneau, et soient $E \in D(X_\etale, \Lambda)$
et $K \in D(Y_\etale, \Lambda)$. Il existe alors un morphisme canonique
$$
Rf_*E \otimes_\Lambda^\mathbf{L} Rg_*K
\longrightarrow
Rc_*(p^{-1}E \otimes_\Lambda^\mathbf{L} q^{-1}K)
$$
Nous pouvons par exemple le définir au moyen des morphismes canoniques
$Rf_*E \to Rc_*p^{-1}E$ et $Rg_*K \to Rc_*q^{-1}K$, ainsi que du
cup-produit relatif défini dans la remarque de Cohomologie sur les sites
\ref{sites-cohomology-remark-cup-product}.
On peut aussi employer l'adjoint du morphisme
$$
c^{-1}(Rf_*E \otimes_\Lambda^\mathbf{L} Rg_*K)
=
p^{-1}f^{-1}Rf_*E \otimes_\Lambda^\mathbf{L} q^{-1} g^{-1}Rg_*K
\to
p^{-1}E \otimes_\Lambda^\mathbf{L} q^{-1}K
$$
qui emploie les morphismes d'adjonction $f^{-1}Rf_*E \to E$ et
$g^{-1}Rg_*K \to K$.
\end{remark}


\begin{lemma}
\label{lemma-kunneth-one-proper}
Soit $k$ un corps séparablement clos. Soit $X$ un schéma propre sur $k$.
Soit $Y$ un schéma quasi-compact et quasi-séparé sur $k$.
\begin{enumerate}
\item Si $E \in D^+(X_\etale)$ a des faisceaux de cohomologie de torsion et
$K \in D^+(Y_\etale)$, alors
$$
R\Gamma(X \times_{\Spec(k)} Y,
\text{pr}_1^{-1}E
\otimes_\mathbf{Z}^\mathbf{L}
\text{pr}_2^{-1}K
)
=
R\Gamma(X, E)
\otimes_\mathbf{Z}^\mathbf{L}
R\Gamma(Y, K)
$$
\item Si $n \geq 1$ est un entier, si $Y$ est de type fini sur $k$,
si $E \in D(X_\etale, \mathbf{Z}/n\mathbf{Z})$ et
$K \in D(Y_\etale, \mathbf{Z}/n\mathbf{Z})$, alors
$$
R\Gamma(X \times_{\Spec(k)} Y,
\text{pr}_1^{-1}E
\otimes_{\mathbf{Z}/n\mathbf{Z}}^\mathbf{L}
\text{pr}_2^{-1}K
)
=
R\Gamma(X, E)
\otimes_{\mathbf{Z}/n\mathbf{Z}}^\mathbf{L}
R\Gamma(Y, K)
$$
\end{enumerate}
\end{lemma}

\begin{proof}
Démonstration de (1). Par le lemme \ref{lemma-projection-formula-proper}, nous avons
$$
R\text{pr}_{2, *}(
\text{pr}_1^{-1}E
\otimes_\mathbf{Z}^\mathbf{L}
\text{pr}_2^{-1}K) =
R\text{pr}_{2, *}(\text{pr}_1^{-1}E)
\otimes_\mathbf{Z}^\mathbf{L}
K
$$
Par changement de base propre (sous la forme du lemme \ref{lemma-proper-base-change}),
cet objet est égal à
$$
\underline{R\Gamma(X, E)}
\otimes_\mathbf{Z}^\mathbf{L}
K
$$
de $D(Y_\etale)$. L'application de $R\Gamma(Y, -)$ à cet objet redonne le
membre de gauche de l'égalité de (1) par la suite spectrale de Leray
de $\text{pr}_2$. Nous concluons donc par le
lemme \ref{lemma-pull-out-constant}.

\medskip\noindent
Démonstration de (2). Elle est exactement la même que celle de (1),
sauf que nous employons les lemmes \ref{lemma-projection-formula-proper-mod-n},
\ref{lemma-proper-base-change-mod-n} et
\ref{lemma-pull-out-constant-mod-n}, ainsi que
$\text{cd}(Y) < \infty$ par le lemme \ref{lemma-finite-cd}.
\end{proof}

\begin{lemma}
\label{lemma-supported-in-closed-points}
Soit $K$ un corps séparablement clos. Soit $X$ un schéma de type
fini sur $K$. Soit $\mathcal{F}$ un faisceau abélien sur $X_\etale$
dont le support est contenu dans l'ensemble des points fermés de $X$.
Alors $H^q(X, \mathcal{F}) = 0$ pour $q > 0$ et $\mathcal{F}$
est engendré par ses sections globales.
\end{lemma}

\begin{proof}
(Si $\mathcal{F}$ est de torsion, l'annulation découle immédiatement
du lemme \ref{lemma-interlude-II}.)
Par le lemme \ref{lemma-support-in-subset}, nous pouvons écrire $\mathcal{F}$
comme limite inductive filtrante de faisceaux constructibles $\mathcal{F}_i$
de $\mathbf{Z}$-modules dont les supports $Z_i \subset X$ sont des ensembles
finis de points fermés. Par la
proposition \ref{proposition-closed-immersion-pushforward},
un tel faisceau est de la forme $(Z_i \to X)_*\mathcal{G}_i$,
où $\mathcal{G}_i$ est un faisceau sur $Z_i$. Comme $K$ est séparablement clos,
le schéma $Z_i$ est une réunion disjointe finie de spectres de corps
séparablement clos. Rappelons que $H^q(Z_i, \mathcal{G}_i) = H^q(X, \mathcal{F}_i)$
par la suite spectrale de Leray de $Z_i \to X$ et l'annulation des images
directes supérieures de ce morphisme
(proposition \ref{proposition-finite-higher-direct-image-zero}).
Par les lemmes \ref{lemma-equivalence-abelian-sheaves-point} et
\ref{lemma-compare-cohomology-point}
nous voyons que $H^q(Z_i, \mathcal{G}_i)$ est nul pour $q > 0$
et que $H^0(Z_i, \mathcal{G}_i)$ engendre $\mathcal{G}_i$.
Nous en déduisons l'annulation de $H^q(X, \mathcal{F}_i)$ pour $q > 0$
et le fait que $\mathcal{F}_i$ est engendré par ses sections globales.
Par le théorème \ref{theorem-colimit}, nous voyons que
$H^q(X, \mathcal{F}) = 0$ pour $q > 0$.
La démonstration est achevée, car une
limite inductive filtrante de faisceaux de groupes abéliens engendrés par leurs sections globales
est engendrée par ses sections globales (détails omis).
\end{proof}

\begin{lemma}
\label{lemma-vanishing-closed-points}
Soit $K$ un corps séparablement clos. Soit $X$ un schéma de type
fini sur $K$. Soit $Q \in D(X_\etale)$. Supposons que $Q_{\overline{x}}$
ne soit non nul que si $x$ est un point fermé de $X$. Alors
$$
Q = 0 \Leftrightarrow H^i(X, Q) = 0 \text{ pour tout }i
$$
\end{lemma}

\begin{proof}
L'implication de gauche à droite est triviale. Nous devons donc
démontrer l'implication réciproque.

\medskip\noindent
Supposons $Q$ borné inférieurement; ce cas suffit pour presque toutes les
applications. Si $Q$ n'est pas nul, considérons le plus petit $i$
tel que le faisceau de cohomologie $H^i(Q)$ soit non nul.
Par le lemme \ref{lemma-supported-in-closed-points}, nous avons
$H^i(X, Q) = H^0(X, H^i(Q)) \not = 0$, et nous concluons.

\medskip\noindent
Cas général. Soit $\mathcal{B} \subset \Ob(X_\etale)$ la sous-catégorie des
objets quasi-compacts. Par le lemme \ref{lemma-supported-in-closed-points},
les hypothèses du lemme de Cohomologie sur les sites
\ref{sites-cohomology-lemma-cohomology-over-U-trivial}
sont satisfaites. Nous concluons que $H^q(U, Q) = H^0(U, H^q(Q))$
pour tout $U \in \mathcal{B}$. En particulier, cela vaut pour $U = X$.
La conclusion découle donc du lemme \ref{lemma-supported-in-closed-points},
puisque $Q$ est nul dans $D(X_\etale)$ si et seulement si $H^q(Q)$ est nul
pour tout $q$.
\end{proof}

\begin{lemma}
\label{lemma-kunneth-localize-on-X}
Soit $K$ un corps. Soit $j : U \to X$ une immersion ouverte de
schémas de type fini sur $K$. Soit $Y$ un schéma de type fini
sur $K$. Considérons le diagramme
$$
\xymatrix{
Y \times_{\Spec(K)} X \ar[d]_q  &
Y \times_{\Spec(K)} U \ar[l]^h \ar[d]^p \\
X & U \ar[l]_j
}
$$
Alors le morphisme de changement de base $q^{-1}Rj_*\mathcal{F} \to Rh_*p^{-1}\mathcal{F}$
est un isomorphisme pour tout faisceau abélien $\mathcal{F}$ sur $U_\etale$
dont les fibres sont de torsion d'ordres inversibles dans $K$.
\end{lemma}

\begin{proof}
Écrivons $\mathcal{F} = \colim \mathcal{F}[n]$, où la limite porte
sur le système multiplicatif des entiers inversibles dans $K$.
Comme la cohomologie commute aux limites inductives filtrantes dans notre situation
(pour un renvoi précis, voir le lemme \ref{lemma-base-change-Rf-star-colim}),
il suffit de démontrer le lemme pour $\mathcal{F}[n]$. Nous pouvons donc
supposer que $\mathcal{F}$ est un faisceau de $\mathbf{Z}/n\mathbf{Z}$-modules
pour un entier $n$ inversible dans $K$ (nous ne l'emploierons qu'à la toute fin de
la démonstration). Dans la démonstration, nous notons brièvement $X \times_K Y$ le produit
fibré au-dessus de $\Spec(K)$.
Nous allons démontrer le lemme par récurrence sur $\dim(X) + \dim(Y)$.
Le lemme est trivial si $\dim(X) \leq 0$, car, dans ce cas,
$U$ est un sous-schéma ouvert et fermé de $X$.
Choisissons un point $z \in X \times_K Y$. Nous allons montrer
que le morphisme sur la fibre en $\overline{z}$ est un isomorphisme.

\medskip\noindent
Supposons que $z \mapsto x \in X$ et que $\text{trdeg}_K(\kappa(x)) > 0$.
Posons $X' = \Spec(\mathcal{O}_{X, x}^{sh})$
et notons $U' \subset X'$ l'image inverse de $U$.
Considérons le changement de base
$$
\xymatrix{
Y \times_K X' \ar[d]_{q'}  &
Y \times_K U' \ar[l]^{h'} \ar[d]^{p'} \\
X' & U' \ar[l]_{j'}
}
$$
de notre diagramme par $X' \to X$. Remarquons que $X' \to X$ est une limite
inductive filtrante de morphismes étales. Par changement de base lisse, sous la forme du
lemme \ref{lemma-smooth-base-change-general}, l'image inverse de
$q^{-1}Rj_*\mathcal{F} \to Rh_*p^{-1}\mathcal{F}$ sur $X'$, puis sur
$Y \times_K X'$, est le morphisme
$(q')^{-1}Rj'_*\mathcal{F}' \to Rh'_*(p')^{-1}\mathcal{F}'$, où
$\mathcal{F}'$ est l'image inverse de $\mathcal{F}$ sur $U'$.
(Il suffirait ici d'employer le changement de base étale, qui est
un résultat essentiellement trivial.) Il suffit donc de montrer
que $(q')^{-1}Rj'_*\mathcal{F}' \to Rh'_*(p')^{-1}\mathcal{F}'$
est un isomorphisme pour prouver que notre morphisme initial
est un isomorphisme sur les fibres en $\overline{z}$.
Par le lemme \ref{lemma-strictly-henselian},
il existe un corps séparablement clos $L/K$ tel que
$X' = \lim X_i$, où $X_i$ est affine de type fini sur $L$
et $\dim(X_i) < \dim(X)$. Pour $i$ assez grand, il
existe un ouvert $U_i \subset X_i$ dont l'image inverse dans $X'$ est $U'$.
Nous pouvons appliquer l'hypothèse de récurrence au diagramme
$$
\vcenter{
\xymatrix{
Y \times_K X_i \ar[d]_{q_i}  &
Y \times_K U_i \ar[l]^{h_i} \ar[d]^{p_i} \\
X_i & U_i \ar[l]_{j_i}
}
}
\quad\text{égal à}\quad
\vcenter{
\xymatrix{
Y_L \times_L X_i \ar[d]_{q_i}  &
Y_L \times_L U_i \ar[l]^{h_i} \ar[d]^{p_i} \\
X_i & U_i \ar[l]_{j_i}
}
}
$$
sur le corps $L$ et à l'image inverse de $\mathcal{F}$ sur ces diagrammes.
Par le lemme \ref{lemma-base-change-Rf-star-colim}, nous concluons que le morphisme
$(q')^{-1}Rj'_*\mathcal{F}' \to Rh'_*(p')^{-1}\mathcal{F}'$ est un isomorphisme.

\medskip\noindent
Supposons que $z \mapsto y \in Y$ et que $\text{trdeg}_K(\kappa(y)) > 0$.
Soit $Y' = \Spec(\mathcal{O}_{Y, y}^{sh})$.
Par le lemme \ref{lemma-strictly-henselian}, il existe un corps séparablement clos
$L/K$ tel que $Y' = \lim Y_i$, où $Y_i$ est affine de type fini sur $L$
et $\dim(Y_i) < \dim(Y)$. En particulier, $Y'$ est un schéma sur $L$.
Notons par un indice $L$ le changement de base des schémas sur $K$
aux schémas sur $L$. Considérons les diagrammes commutatifs
$$
\vcenter{
\xymatrix{
Y' \times_K X \ar[d]_f  &
Y' \times_K U \ar[l]^{h'} \ar[d]^{f'} \\
Y \times_K X \ar[d]_q  &
Y \times_K U \ar[l]^h \ar[d]^p \\
X & U \ar[l]_j
}
}
\quad\text{et}\quad
\vcenter{
\xymatrix{
Y' \times_L X_L \ar[d]_{q'}  &
Y' \times_L U_L \ar[l]^{h'} \ar[d]^{p'} \\
X_L \ar[d] &
U_L \ar[l]^{j_L} \ar[d] \\
X & U \ar[l]_j
}
}
$$
et remarquons que les lignes supérieure et inférieure sont les mêmes à gauche et à droite.
Par changement de base lisse, nous voyons que
$f^{-1}Rh_*p^{-1}\mathcal{F} = Rh'_*(f')^{-1}p^{-1}\mathcal{F}$
(comme au paragraphe précédent).
Par changement de base lisse pour $\Spec(L) \to \Spec(K)$
(lemme \ref{lemma-base-change-field-extension}),
nous voyons que $Rj_{L, *}\mathcal{F}_L$ est l'image inverse de
$Rj_*\mathcal{F}$ sur $X_L$. En combinant ces deux observations,
nous concluons qu'il suffit de montrer que le morphisme de changement de base
du carré supérieur du diagramme de droite
est un isomorphisme pour prouver que notre morphisme initial
est un isomorphisme sur les fibres en $\overline{z}$\footnote{Nous employons ici
le fait qu'une « composée verticale » de morphismes de changement de base est un morphisme de
changement de base, comme l'explique la remarque de Cohomologie sur les sites
\ref{sites-cohomology-remark-compose-base-change}.}.
Puis, en employant $Y' = \lim Y_i$ et en raisonnant exactement
comme au paragraphe précédent, nous voyons que l'hypothèse de
récurrence force notre morphisme sur $Y' \times_K X$ à être
un isomorphisme.

\medskip\noindent
Ainsi, tout contre-exemple pour lequel $\dim(X) + \dim(Y)$ est minimal
ne pourrait avoir de défaut d'isomorphisme de
$q^{-1}Rj_*\mathcal{F} \to Rh_*p^{-1}\mathcal{F}$
que sur les fibres en des points fermés de $X \times_K Y$ (car un point
$z$ de $X \times_K Y$ est fermé si et seulement si
ses images dans $X$ et dans $Y$ sont toutes deux fermées).
Comme il suffit de démontrer l'isomorphisme localement,
nous pouvons supposer $X$ et $Y$ affines. Nous pouvons alors
choisir une immersion ouverte dense $Y \to Y'$ avec $Y'$
projectif. (Choisir une immersion fermée $Y \to \mathbf{A}^n_K$
et prendre pour $Y'$ l'adhérence schématique de $Y$ dans
$\mathbf{P}^n_K$.) Alors $\dim(Y') = \dim(Y)$, et nous
obtenons donc un contre-exemple « minimal » avec $Y$ projectif sur $K$.
Au paragraphe suivant, nous montrons que c'est impossible.

\medskip\noindent
Considérons un diagramme comme dans l'énoncé du lemme, tel que
$q^{-1}Rj_*\mathcal{F} \to Rh_*p^{-1}\mathcal{F}$
soit un isomorphisme en tout point non fermé de $X \times_K Y$
et que $Y$ soit projectif.
La restriction du morphisme à $(X \times_K Y)_{K^{sep}}$
est le morphisme correspondant pour le diagramme du lemme
après changement de base à $K^{sep}$. Nous pouvons donc supposer, et supposons, $K$
séparablement clos. Choisissons un triangle distingué
$$
q^{-1}Rj_*\mathcal{F} \to Rh_*p^{-1}\mathcal{F} \to Q \to
(q^{-1}Rj_*\mathcal{F})[1]
$$
dans $D((X \times_K Y)_\etale)$. Comme $Q$ est supporté par des
points fermés, il suffit de démontrer
$H^i(X \times_K Y, Q) = 0$ pour tout $i$; voir le
lemme \ref{lemma-vanishing-closed-points}.
Il suffit donc de démontrer que
$q^{-1}Rj_*\mathcal{F} \to Rh_*p^{-1}\mathcal{F}$
induit un isomorphisme en cohomologie. Rappelons que $\mathcal{F}$
est annulé par un entier $n$ inversible dans $K$.
Par la formule de K\"unneth du lemme \ref{lemma-kunneth-one-proper},
nous avons
\begin{align*}
R\Gamma(X \times_K Y, q^{-1}Rj_*\mathcal{F})
&=
R\Gamma(X, Rj_*\mathcal{F})
\otimes_{\mathbf{Z}/n\mathbf{Z}}^\mathbf{L}
R\Gamma(Y, \mathbf{Z}/n\mathbf{Z}) \\
& =
R\Gamma(U, \mathcal{F})
\otimes_{\mathbf{Z}/n\mathbf{Z}}^\mathbf{L}
R\Gamma(Y, \mathbf{Z}/n\mathbf{Z})
\end{align*}
et
$$
R\Gamma(X \times_K Y, Rh_*p^{-1}\mathcal{F}) =
R\Gamma(U \times_K Y, p^{-1}\mathcal{F}) =
R\Gamma(U, \mathcal{F})
\otimes_{\mathbf{Z}/n\mathbf{Z}}^\mathbf{L}
R\Gamma(Y, \mathbf{Z}/n\mathbf{Z})
$$
Cela achève la démonstration.
\end{proof}

\begin{lemma}
\label{lemma-punctual-base-change}
Soit $K$ un corps. Pour tout diagramme commutatif
$$
\xymatrix{
X \ar[d] & X' \ar[l] \ar[d]_{f'} & Y \ar[l]^h \ar[d]^e \\
\Spec(K) & S' \ar[l] & T \ar[l]_g
}
$$
de schémas sur $K$ tel que
$X' = X \times_{\Spec(K)} S'$ et $Y = X' \times_{S'} T$, avec
$g$ quasi-compact et quasi-séparé, et pour tout faisceau abélien
$\mathcal{F}$ sur $T_\etale$ dont les fibres sont de torsion d'ordres
inversibles dans $K$, le morphisme de changement de base
$$
(f')^{-1}Rg_*\mathcal{F}
\longrightarrow
Rh_*e^{-1}\mathcal{F}
$$
est un isomorphisme.
\end{lemma}

\begin{proof}
La question est locale sur $X$; nous pouvons donc supposer $X$ affine.
Par le lemme \ref{limits-lemma-relative-approximation} de Limits,
nous pouvons écrire $X = \lim X_i$ comme une limite projective filtrante, à
morphismes de transition affines, de schémas $X_i$ de type fini sur $K$.
Notons $X'_i =  X_i \times_{\Spec(K)} S'$ et $Y_i = X'_i \times_{S'} T$.
Par le lemme \ref{lemma-base-change-Rf-star-colim},
il suffit de démontrer l'énoncé pour les carrés
de sommets $X_i, Y_i, S', T$.
Nous pouvons donc supposer $X$ de type fini sur $K$.
De même, nous pouvons écrire
$\mathcal{F} = \colim \mathcal{F}[n]$, où la limite porte
sur le système multiplicatif des entiers inversibles dans $K$.
Le même lemme employé ci-dessus nous ramène au cas où
$\mathcal{F}$ est un faisceau de $\mathbf{Z}/n\mathbf{Z}$-modules
pour un entier $n$ inversible dans $K$.

\medskip\noindent
Nous pouvons remplacer $K$ par sa clôture algébrique $\overline{K}$.
En effet, la formation de l'image directe commute au changement de base
à $\overline{K}$ d'après le
lemme \ref{lemma-base-change-field-extension} (pour $g$ comme pour
$h$). Et il suffit de démontrer l'égalité après
restriction à $X'_{\overline{K}}$. Nous pouvons ensuite remplacer
$X$ par sa réduction, grâce à l'invariance topologique de la cohomologie
étale; voir la
proposition \ref{proposition-topological-invariance}.
Après ce remplacement, le morphisme $X \to \Spec(K)$
est plat, de présentation finie et à fibres géométriquement réduites;
il en va de même après tout changement de base, notamment pour $X' \to S'$.
Ainsi, $(f')^{-1}Rg_*\mathcal{F} \to Rh_*e^{-1}\mathcal{F}$
est un isomorphisme par le lemme \ref{lemma-fppf-reduced-fibres-base-change-f-star}.

\medskip\noindent
Nous pouvons alors appliquer le
lemme \ref{lemma-base-change-does-not-hold-post}
pour voir qu'il suffit de démontrer ceci : étant donné un diagramme commutatif
$$
\xymatrix{
X \ar[d]_f & X' \ar[d] \ar[l] & Y \ar[l]^h \ar[d] \\
\Spec(K) & S' \ar[l] & \Spec(L) \ar[l]
}
$$
dont les deux carrés sont cartésiens, où $S'$ est affine, intègre et normal,
de corps des fonctions algébriquement clos $K$, alors
$R^qh_*(\mathbf{Z}/d\mathbf{Z})$ est nul pour $q > 0$ et tout $d | n$.
Remarquons que cette annulation équivaut à l'affirmation que
$$
(f')^{-1}R^q(\Spec(L) \to S')_*\mathbf{Z}/d\mathbf{Z}
\longrightarrow
R^qh_*\mathbf{Z}/d\mathbf{Z}
$$
est un isomorphisme, car le membre de gauche est nul, par exemple par le
lemme \ref{lemma-Rf-star-zero-normal-with-alg-closed-function-field}.

\medskip\noindent
Écrivons $S' = \Spec(B)$, de sorte que $L$ soit le corps des fractions de $B$.
Écrivons $B = \bigcup_{i \in I} B_i$ comme la réunion de ses
sous-$K$-algèbres de type fini $B_i$. Soit $J$ l'ensemble des couples $(i, g)$ où
$i \in I$ et $g \in B_i$ est non nul, ordonné par
$(i', g') \geq (i, g)$ si et seulement si $i' \geq i$ et si
l'image de $g$ est inversible dans $(B_{i'})_{g'}$.
Alors $L = \colim_{(i, g) \in J} (B_i)_g$.
Pour $j = (i, g) \in J$, posons $S_j = \Spec(B_i)$
et $U_j = \Spec((B_i)_g)$.
Alors
$$
\vcenter{
\xymatrix{
X' \ar[d] & Y \ar[l]^h \ar[d] \\
S' & \Spec(L) \ar[l]
}
}
\quad\text{est la limite projective des}\quad
\vcenter{
\xymatrix{
X \times_K S_j \ar[d] & X \times_K U_j \ar[l]^{h_j} \ar[d] \\
S_j & U_j \ar[l]
}
}
$$
Nous pouvons donc appliquer le lemme \ref{lemma-base-change-Rf-star-colim}
pour voir qu'il suffit de démontrer
le changement de base dans les diagrammes de droite, ce que nous avons
démontré dans le lemme \ref{lemma-kunneth-localize-on-X}.
\end{proof}

\begin{lemma}
\label{lemma-punctual-base-change-upgrade}
Soit $K$ un corps. Soit $n \geq 1$ inversible dans $K$.
Considérons un diagramme commutatif
$$
\xymatrix{
X \ar[d] & X' \ar[l]^p \ar[d]_{f'} & Y \ar[l]^h \ar[d]^e \\
\Spec(K) & S' \ar[l] & T \ar[l]_g
}
$$
de schémas tel que
$X' = X \times_{\Spec(K)} S'$ et $Y = X' \times_{S'} T$, avec
$g$ quasi-compact et quasi-séparé. Le morphisme canonique
$$
p^{-1}E \otimes_{\mathbf{Z}/n\mathbf{Z}}^\mathbf{L} (f')^{-1}Rg_*F
\longrightarrow
Rh_*(h^{-1}p^{-1}E \otimes_{\mathbf{Z}/n\mathbf{Z}}^\mathbf{L} e^{-1}F)
$$
est un isomorphisme si $E \in D^+(X_\etale, \mathbf{Z}/n\mathbf{Z})$
est d'amplitude de Tor dans $[a, \infty]$ pour un $a \in \mathbf{Z}$ et si
$F \in D^+(T_\etale, \mathbf{Z}/n\mathbf{Z})$.
\end{lemma}

\begin{proof}
Ce lemme généralise le lemme \ref{lemma-punctual-base-change}
aux objets de la catégorie dérivée; son assertion est vraie parce que,
dans le lemme \ref{lemma-punctual-base-change}, le schéma $X$ sur $K$
est arbitraire. Nous recommandons vivement au lecteur d'omettre cette démonstration laborieuse
(ou de n'en lire que le dernier paragraphe).

\medskip\noindent
Nous pouvons représenter $E$ par un complexe K-plat borné inférieurement
$\mathcal{E}^\bullet$ formé de $\mathbf{Z}/n\mathbf{Z}$-modules plats.
Voir le lemme de Cohomologie sur les sites
\ref{sites-cohomology-lemma-bounded-below-tor-amplitude}.
Choisissons un entier $b$ tel que $H^i(F) = 0$ pour $i < b$.
Choisissons un entier $N$ assez grand et considérons la suite exacte courte
$$
0 \to \sigma_{\geq N + 1}\mathcal{E}^\bullet \to
\mathcal{E}^\bullet \to
\sigma_{\leq N}\mathcal{E}^\bullet \to 0
$$
de troncatures bêtes. Elle produit un triangle distingué
$E'' \to E \to E' \to E''[1]$ dans $D(X_\etale, \mathbf{Z}/n\mathbf{Z})$.
Pour $F$ fixé, les deux membres de la flèche
de l'énoncé du lemme sont des foncteurs exacts en $E$. Remarquons que
$$
p^{-1}E'' \otimes_{\mathbf{Z}/n\mathbf{Z}}^\mathbf{L} (f')^{-1}Rg_*F
\quad\text{et}\quad
Rh_*(h^{-1}p^{-1}E'' \otimes_{\mathbf{Z}/n\mathbf{Z}}^\mathbf{L} e^{-1}F)
$$
sont concentrés en degrés $\geq N + b$. Si nous pouvons donc démontrer le lemme
pour l'objet $E'$, nous voyons qu'il vaut en degrés
$\leq N + b$, ce qui permet de conclure. Certains détails sont omis.
Nous pouvons ainsi supposer que $E$ est représenté
par un complexe borné de $\mathbf{Z}/n\mathbf{Z}$-modules plats.
Un autre argument de même nature permet de supposer
que $E$ est donné par un unique $\mathbf{Z}/n\mathbf{Z}$-module plat
$\mathcal{E}$.

\medskip\noindent
Nous appliquons ensuite les mêmes arguments à la variable $F$
pour nous ramener au cas où $F$ est donné par un unique
faisceau de $\mathbf{Z}/n\mathbf{Z}$-modules $\mathcal{F}$.
Supposons $\mathcal{F}$ annulé par un entier $m | n$.
Si $\ell$ est un nombre premier divisant $m$ et $m > \ell$,
nous pouvons considérer la suite exacte courte
$0 \to \mathcal{F}[\ell] \to \mathcal{F} \to
\mathcal{F}/\mathcal{F}[\ell] \to 0$
et nous ramener à un entier $m$ plus petit. Nous sommes finalement ramenés
au cas où $\mathcal{F}$ est annulé par un nombre
premier $\ell$ divisant $n$.
Dans ce cas, remarquons que
$$
p^{-1}\mathcal{E}
\otimes_{\mathbf{Z}/n\mathbf{Z}}^\mathbf{L}
(f')^{-1}Rg_*\mathcal{F}
=
p^{-1}(\mathcal{E}/\ell \mathcal{E})
\otimes_{\mathbf{F}_\ell}^\mathbf{L}
(f')^{-1}Rg_*\mathcal{F}
$$
par platitude de $\mathcal{E}$. Il en va de même pour l'autre terme.
Cela nous ramène au cas où nous travaillons avec des faisceaux
d'espaces vectoriels sur $\mathbf{F}_\ell$, traité ci-dessous.

\medskip\noindent
Supposons que $\ell$ soit un nombre premier inversible dans $K$.
Supposons que $\mathcal{E}$ et $\mathcal{F}$ soient des faisceaux
d'espaces vectoriels sur $\mathbf{F}_\ell$ sur $X_\etale$ et $T_\etale$.
Nous voulons montrer que
$$
p^{-1}\mathcal{E} \otimes_{\mathbf{F}_\ell} (f')^{-1}R^qg_*\mathcal{F}
\longrightarrow
R^qh_*(h^{-1}p^{-1}\mathcal{E} \otimes_{\mathbf{F}_\ell} e^{-1}\mathcal{F})
$$
est un isomorphisme pour tout $q \geq 0$. Cette question est locale sur $X$;
nous pouvons donc supposer $X$ affine. Nous pouvons écrire $\mathcal{E}$
comme limite inductive filtrante de faisceaux constructibles
d'espaces vectoriels sur $\mathbf{F}_\ell$ sur $X_\etale$; voir le
lemme \ref{lemma-torsion-colimit-constructible}.
Comme les produits tensoriels commutent
aux limites inductives filtrantes, et qu'il en va de même des
images directes supérieures (lemme \ref{lemma-relative-colimit}),
nous pouvons supposer que $\mathcal{E}$ est un faisceau constructible
d'espaces vectoriels sur $\mathbf{F}_\ell$ sur $X_\etale$.
Nous pouvons alors choisir un entier $m$ et des morphismes finis de présentation finie
$\pi_i : X_i \to X$, $i = 1, \ldots, m$,
tels qu'il existe un morphisme injectif
$$
\mathcal{E} \to
\bigoplus\nolimits_{i = 1}^m
\pi_{i, *}\mathbf{F}_\ell
$$
Voir le lemme \ref{lemma-constructible-maps-into-constant-general}.
Remarquons que la somme directe est également un faisceau constructible
(lemme \ref{lemma-finite-pushforward-constructible}).
Son conoyau est donc lui aussi constructible
(lemme \ref{lemma-constructible-abelian}).
Par décalage de dimension, c'est-à-dire par récurrence sur $q$,
dans la catégorie des faisceaux constructibles
d'espaces vectoriels sur $\mathbf{F}_\ell$ sur $X_\etale$, il suffit de démontrer le
résultat pour les faisceaux $\pi_{i, *}\mathbf{F}_\ell$
(détails omis; indication : commencer par démontrer l'injectivité pour $q = 0$
pour tout $\mathcal{E}$ constructible).
Pour traiter ce cas, complétons le diagramme du lemme en
$$
\xymatrix{
X_i \ar[d]^{\pi_i} &
X'_i \ar[l]^{p_i} \ar[d]^{\pi'_i} &
Y_i \ar[l]^{h_i} \ar[d]^{\rho_i} \\
X \ar[d] & X' \ar[l]^p \ar[d]_{f'} & Y \ar[l]^h \ar[d]^e \\
\Spec(K) & S' \ar[l] & T \ar[l]_g
}
$$
dont tous les carrés sont cartésiens. Dans les égalités ci-dessous, nous
allons employer $R\pi_{i, *} = \pi_{i, *}$, et les égalités analogues
pour $\pi'_i$ et $\rho_i$; nous allons employer la suite spectrale de Leray,
ainsi que
le lemme \ref{lemma-finite-pushforward-commutes-with-base-change}, et
nous allons employer
le lemme \ref{lemma-projection-formula-proper-mod-n}
(bien que ce dernier soit presque trivial pour les morphismes finis), appliqué à
$\pi_i$, $\pi'_i$ et $\rho_i$.
Nous voyons ainsi que
\begin{align*}
p^{-1}\pi_{i, *}\mathbf{F}_\ell
\otimes_{\mathbf{F}_\ell} (f')^{-1}R^qg_*\mathcal{F}
& =
\pi'_{i, *}\mathbf{F}_\ell
\otimes_{\mathbf{F}_\ell} (f')^{-1}R^qg_*\mathcal{F} \\
& =
\pi'_{i, *}((\pi'_i)^{-1} (f')^{-1}R^qg_*\mathcal{F})
\end{align*}
De même, nous avons
\begin{align*}
R^qh_*(h^{-1}p^{-1} \pi_{i, *}\mathbf{F}_\ell
\otimes_{\mathbf{F}_\ell} e^{-1}\mathcal{F})
& =
R^qh_*(\rho_{i, *}\mathbf{F}_\ell
\otimes_{\mathbf{F}_\ell} e^{-1}\mathcal{F}) \\
& =
R^qh_*\rho_{i, *}(\rho_i^{-1}e^{-1}\mathcal{F}) \\
& =
\pi'_{i, *}R^qh_{i, *} \rho_i^{-1}e^{-1}\mathcal{F}
\end{align*}
Puisque
$R^qh_{i, *} \rho_i^{-1}e^{-1}\mathcal{F} = 
(\pi'_i)^{-1} (f')^{-1}R^qg_*\mathcal{F}$ par le
lemme \ref{lemma-punctual-base-change},
nous concluons.
\end{proof}

\begin{lemma}
\label{lemma-punctual-base-change-upgrade-unbounded}
Soit $K$ un corps. Soit $n \geq 1$ inversible dans $K$.
Considérons un diagramme commutatif
$$
\xymatrix{
X \ar[d] & X' \ar[l]^p \ar[d]_{f'} & Y \ar[l]^h \ar[d]^e \\
\Spec(K) & S' \ar[l] & T \ar[l]_g
}
$$
de schémas de type fini sur $K$ tel que
$X' = X \times_{\Spec(K)} S'$ et $Y = X' \times_{S'} T$.
Le morphisme canonique
$$
p^{-1}E \otimes_{\mathbf{Z}/n\mathbf{Z}}^\mathbf{L} (f')^{-1}Rg_*F
\longrightarrow
Rh_*(h^{-1}p^{-1}E \otimes_{\mathbf{Z}/n\mathbf{Z}}^\mathbf{L} e^{-1}F)
$$
est un isomorphisme pour $E \in D(X_\etale, \mathbf{Z}/n\mathbf{Z})$
et $F \in D(T_\etale, \mathbf{Z}/n\mathbf{Z})$.
\end{lemma}

\begin{proof}
Nous allons nous ramener au
lemme \ref{lemma-punctual-base-change-upgrade}
en employant le fait que nos foncteurs commutent aux sommes directes.
Nous suggérons au lecteur d'omettre cette démonstration.
Rappelons que le produit tensoriel dérivé commute aux
sommes directes. Rappelons aussi que l'image inverse (dérivée) commute aux sommes directes.
Enfin, $Rh_*$ et $Rg_*$ commutent aux sommes directes; voir les
lemmes \ref{lemma-finite-cd} et
\ref{lemma-finite-cd-mod-n-direct-sums}
(c'est ici que nous employons le fait que nos schémas sont de type fini
sur $K$).

\medskip\noindent
Pour achever la démonstration, nous pouvons raisonner comme suit.
Écrivons d'abord $E = \text{hocolim} \tau_{\leq N} E$.
Comme nos foncteurs commutent aux sommes directes, ils commutent
aux colimites homotopiques. Il suffit donc de démontrer
le lemme pour $E$ borné supérieurement. De même, nous pouvons
supposer $F$ borné supérieurement.
Nous pouvons alors représenter $E$ par un complexe borné supérieurement
$\mathcal{E}^\bullet$ de faisceaux de $\mathbf{Z}/n\mathbf{Z}$-modules.
Alors
$$
\mathcal{E}^\bullet = \colim \sigma_{\geq -N}\mathcal{E}^\bullet
$$
(troncatures bêtes).
Nous pouvons donc supposer que $\mathcal{E}^\bullet$ est un complexe
borné de faisceaux de $\mathbf{Z}/n\mathbf{Z}$-modules.
Pour $F$, nous choisissons un complexe borné supérieurement de
faisceaux plats (!) de $\mathbf{Z}/n\mathbf{Z}$-modules.
Nous nous ramenons alors au cas où $F$ est représenté
par un complexe borné de faisceaux plats de $\mathbf{Z}/n\mathbf{Z}$-modules.
Le lemme \ref{lemma-punctual-base-change-upgrade}
s'applique, et nous concluons.
\end{proof}

\begin{lemma}
\label{lemma-kunneth}
Soit $k$ un corps séparablement clos. Soient $X$ et $Y$ des
schémas de type fini sur $k$. Soit $n \geq 1$ un entier
inversible dans $k$. Alors, pour
$E \in D(X_\etale, \mathbf{Z}/n\mathbf{Z})$ et
$K \in D(Y_\etale, \mathbf{Z}/n\mathbf{Z})$
nous avons
$$
R\Gamma(X \times_{\Spec(k)} Y,
\text{pr}_1^{-1}E
\otimes_{\mathbf{Z}/n\mathbf{Z}}^\mathbf{L}
\text{pr}_2^{-1}K
)
=
R\Gamma(X, E)
\otimes_{\mathbf{Z}/n\mathbf{Z}}^\mathbf{L}
R\Gamma(Y, K)
$$
\end{lemma}

\begin{proof}
Par le lemme \ref{lemma-punctual-base-change-upgrade-unbounded}, nous avons
$$
R\text{pr}_{1, *}(
\text{pr}_1^{-1}E
\otimes_{\mathbf{Z}/n\mathbf{Z}}^\mathbf{L}
\text{pr}_2^{-1}K) =
E \otimes_{\mathbf{Z}/n\mathbf{Z}}^\mathbf{L}
\underline{R\Gamma(Y, K)}
$$
Nous concluons par le
lemme \ref{lemma-pull-out-constant-mod-n},
que nous pouvons employer parce que
$\text{cd}(X) < \infty$ par le lemme \ref{lemma-finite-cd}.
\end{proof}














\section{Comparaison des topologies chaotique et de Zariski}
\label{section-compare-chaotic-Zariski}

\noindent
Lorsqu'on construit le faisceau structural d'un schéma affine, on en construit d'abord
les valeurs sur les ouverts affines, puis on l'étend à tous les ouverts.
Une construction analogue est souvent utile pour construire des complexes
de groupes abéliens sur un schéma $X$. Rappelons que $X_{affine, Zar}$
désigne la catégorie des ouverts affines de $X$, munie de la topologie définie
par les recouvrements de Zariski standard; voir la
définition \ref{topologies-definition-big-small-Zariski} de Topologies.
Rappelons au lecteur que le topos de $X_{affine, Zar}$
est le petit topos de Zariski de $X$; voir le lemme de Topologies
\ref{topologies-lemma-alternative-zariski}.
Dans cette section, nous notons $X_{affine}$ la même catégorie
sous-jacente munie de la topologie chaotique, c'est-à-dire telle que les faisceaux
coïncident avec les préfaisceaux. Nous obtenons un morphisme de sites
$$
\epsilon : X_{affine, Zar} \longrightarrow X_{affine}
$$
comme dans la section \ref{sites-cohomology-section-compare} de Cohomologie sur les sites.

\begin{lemma}
\label{lemma-check-zar}
Dans la situation ci-dessus, soit $K$ un objet de $D^+(X_{affine})$.
Alors $K$ appartient à l'image essentielle du foncteur (pleinement fidèle)
$R\epsilon_* ; D(X_{affine, Zar}) \to D(X_{affine})$ si et seulement
si les deux conditions suivantes sont satisfaites :
\begin{enumerate}
\item $R\Gamma(\emptyset, K)$ est nul dans $D(\textit{Ab})$, et
\item si $U = V \cup W$, où $U, V, W \subset X$ sont des ouverts affines et
$V, W \subset U$ des ouverts standard
(définition \ref{algebra-definition-Zariski-topology} d'Algèbre), alors
le morphisme $c^K_{U, V, W, V \cap W}$ du
lemme \ref{sites-cohomology-lemma-c-square} de Cohomologie sur les sites
est un quasi-isomorphisme.
\end{enumerate}
\end{lemma}

\begin{proof}
(Le foncteur $R\epsilon_*$ est pleinement fidèle d'après la discussion de la
section \ref{sites-cohomology-section-compare} de Cohomologie sur les sites.)
À un problème près concernant l'ensemble vide,
cela découle du très général lemme de Cohomologie sur les sites
\ref{sites-cohomology-lemma-descent-squares}, dont les hypothèses valent par le
lemme \ref{schemes-lemma-sheaf-on-affines} de Schémas et le
lemme de Cohomologie sur les sites
\ref{sites-cohomology-lemma-descent-squares-helper}.

\medskip\noindent
Pour contourner ce problème, notons $X_{affine, almost-chaotic}$
le site où l'objet vide $\emptyset$ possède deux recouvrements,
à savoir $\{\emptyset \to \emptyset\}$ et le recouvrement vide
(voir l'exemple \ref{sites-example-site-topological} de Sites pour une
discussion). Nous avons alors des morphismes de sites
$$
X_{affine, Zar} \to X_{affine, almost-chaotic} \to X_{affine}
$$
L'argument ci-dessus vaut pour la première flèche. Nous laissons alors
au lecteur le soin de vérifier qu'un objet $K$ de $D^+(X_{affine})$
appartient à l'image essentielle du foncteur (pleinement fidèle)
$D(X_{affine, almost-chaotic}) \to D(X_{affine})$ si et seulement
si $R\Gamma(\emptyset, K)$ est nul dans $D(\textit{Ab})$.
\end{proof}









\section{Comparaison des grands et petits topos}
\label{section-compare}

\noindent
Soit $S$ un schéma. Dans le
lemme \ref{topologies-lemma-at-the-bottom-etale} de Topologies,
nous avons introduit les morphismes de comparaison
$\pi_S : (\Sch/S)_\etale \to S_\etale$ et
$i_S : \Sh(S_\etale) \to \Sh((\Sch/S)_\etale)$,
avec $\pi_S \circ i_S = \text{id}$ et $\pi_{S, *} = i_S^{-1}$.
Plus généralement, si $f : T \to S$ est un objet de $(\Sch/S)_\etale$,
il existe un morphisme $i_f : \Sh(T_\etale) \to \Sh((\Sch/S)_\etale)$
tel que $f_{small} = \pi_S \circ i_f$; voir les
lemmes \ref{topologies-lemma-put-in-T-etale} et
\ref{topologies-lemma-morphism-big-small-etale} de Topologies. Dans la
remarque \ref{descent-remark-change-topologies-ringed} de Descente,
nous les avons étendus en un morphisme de sites annelés
$$
\pi_S : ((\Sch/S)_\etale, \mathcal{O}) \to (S_\etale, \mathcal{O}_S)
$$
et des morphismes de topos annelés
$$
i_S : (\Sh(S_\etale), \mathcal{O}_S) \to (\Sh((\Sch/S)_\etale), \mathcal{O})
$$
et
$$
i_f : (\Sh(T_\etale), \mathcal{O}_T) \to (\Sh((\Sch/S)_\etale, \mathcal{O}))
$$
Remarquons que la restriction $i_S^{-1} = \pi_{S, *}$ (voir la
définition \ref{topologies-definition-restriction-small-etale} de Topologies)
transforme $\mathcal{O}$ en $\mathcal{O}_S$.
De même, $i_f^{-1}$ transforme $\mathcal{O}$ en $\mathcal{O}_T$.
Voir la remarque \ref{descent-remark-change-topologies-ringed} de Descente.
Ainsi, $i_S^*\mathcal{F} = i_S^{-1}\mathcal{F}$ et
$i_f^*\mathcal{F} = i_f^{-1}\mathcal{F}$ pour tout $\mathcal{O}$-module
$\mathcal{F}$ sur $(\Sch/S)_\etale$. En particulier, $i_S^*$ et $i_f^*$
sont des foncteurs exacts. Le foncteur $i_S^*$ est souvent noté
$\mathcal{F} \mapsto \mathcal{F}|_{S_\etale}$ (ce qui n'entre pas
en conflit avec la notation de la
définition \ref{topologies-definition-restriction-small-etale} de Topologies).

\begin{lemma}
\label{lemma-compare-injectives}
Soit $S$ un schéma. Soit $T$ un objet de $(\Sch/S)_\etale$.
\begin{enumerate}
\item Si $\mathcal{I}$ est injectif dans $\textit{Ab}((\Sch/S)_\etale)$, alors
\begin{enumerate}
\item $i_f^{-1}\mathcal{I}$ est injectif dans $\textit{Ab}(T_\etale)$,
\item $\mathcal{I}|_{S_\etale}$ est injectif dans $\textit{Ab}(S_\etale)$,
\end{enumerate}
\item Si $\mathcal{I}^\bullet$ est un complexe K-injectif
dans $\textit{Ab}((\Sch/S)_\etale)$, alors
\begin{enumerate}
\item $i_f^{-1}\mathcal{I}^\bullet$ est un complexe K-injectif dans
$\textit{Ab}(T_\etale)$,
\item $\mathcal{I}^\bullet|_{S_\etale}$ est un complexe K-injectif dans
$\textit{Ab}(S_\etale)$,
\end{enumerate}
\end{enumerate}
Les énoncés correspondants ne valent pas pour les modules.
\end{lemma}

\begin{proof}
Les parties (1)(b) et (2)(b)
découlent formellement du fait que le foncteur de restriction
$\pi_{S, *} = i_S^{-1}$ est adjoint à droite du foncteur exact
$\pi_S^{-1}$; voir le
lemme \ref{homology-lemma-adjoint-preserve-injectives} d'Homology et le
lemme \ref{derived-lemma-adjoint-preserve-K-injectives} de Catégories dérivées.

\medskip\noindent
Les parties (1)(a) et (2)(a) se voient de deux manières. Première démonstration : on peut utiliser
le fait que $i_f^{-1}$ est adjoint à droite du foncteur exact $i_{f, !}$.
Ce foncteur est construit dans le
lemme \ref{topologies-lemma-put-in-T-etale} de Topologies
pour les faisceaux d'ensembles, et dans le
lemme \ref{sites-modules-lemma-g-shriek-adjoint} de Modules sur les sites pour les faisceaux abéliens.
Le lemme de Modules sur les sites
\ref{sites-modules-lemma-exactness-lower-shriek} montre qu'il est exact.
Seconde démonstration. Nous pouvons employer $i_f = i_T \circ f_{big}$, comme le montre
le lemme \ref{topologies-lemma-morphism-big-small-etale} de Topologies.
Puisque $f_{big}$ est une localisation, son image inverse
conserve les injectifs et les K-injectifs; voir les
lemmes \ref{sites-cohomology-lemma-cohomology-of-open} et
\ref{sites-cohomology-lemma-restrict-K-injective-to-open}.
Nous appliquons alors les parties (1)(b) et (2)(b) déjà démontrées
au foncteur $i_T^{-1}$, et nous concluons.

\medskip\noindent
Soit $S = \Spec(\mathbf{Z})$ et considérons le morphisme
$2 : \mathcal{O}_S \to \mathcal{O}_S$. C'est un morphisme injectif
de $\mathcal{O}_S$-modules sur $S_\etale$. Cependant, l'image inverse
$\pi_S^*(2) : \mathcal{O} \to \mathcal{O}$ n'est pas injective,
comme on le voit en l'évaluant sur $\Spec(\mathbf{F}_2)$. Choisissons maintenant
une injection $\alpha : \mathcal{O} \to \mathcal{I}$ dans un
$\mathcal{O}$-module injectif $\mathcal{I}$ sur $(\Sch/S)_\etale$.
Considérons alors le diagramme
$$
\xymatrix{
\mathcal{O}_S \ar[d]_2 \ar[rr]_{\alpha|_{S_\etale}} & &
\mathcal{I}|_{S_\etale} \\
\mathcal{O}_S \ar@{..>}[rru]
}
$$
La flèche pointillée ne peut exister dans la catégorie des $\mathcal{O}_S$-modules,
car cela signifierait
(par adjonction) que le morphisme injectif $\alpha$ se factorise par
le morphisme non injectif $\pi_S^*(2)$, ce qui est impossible.
Ainsi, $\mathcal{I}|_{S_\etale}$ n'est pas un $\mathcal{O}_S$-module injectif.
\end{proof}

\noindent
Soit $f : T \to S$ un morphisme de schémas. Le diagramme commutatif du
lemme \ref{topologies-lemma-morphism-big-small-etale} (3) de Topologies
donne un diagramme commutatif de sites annelés
$$
\xymatrix{
(T_\etale, \mathcal{O}_T) \ar[d]_{f_{small}} &
((\Sch/T)_\etale, \mathcal{O}) \ar[d]^{f_{big}} \ar[l]^{\pi_T} \\
(S_\etale, \mathcal{O}_S) &
((\Sch/S)_\etale, \mathcal{O}) \ar[l]_{\pi_S}
}
$$
comme on le voit aisément en explicitant les définitions de
$f_{small}^\sharp$, $f_{big}^\sharp$, $\pi_S^\sharp$ et $\pi_T^\sharp$.
Cela signifie en particulier que
\begin{equation}
\label{equation-compare-big-small}
(f_{big, *}\mathcal{F})|_{S_\etale} =
f_{small, *}(\mathcal{F}|_{T_\etale})
\end{equation}
pour tout faisceau $\mathcal{F}$ sur $(\Sch/T)_\etale$; si $\mathcal{F}$
est un faisceau de $\mathcal{O}$-modules, alors (\ref{equation-compare-big-small})
est un isomorphisme de $\mathcal{O}_S$-modules sur $S_\etale$.

\begin{lemma}
\label{lemma-compare-higher-direct-image}
Soit $f : T \to S$ un morphisme de schémas.
\begin{enumerate}
\item Pour $K$ dans $D((\Sch/T)_\etale)$, nous avons
$
(Rf_{big, *}K)|_{S_\etale} = Rf_{small, *}(K|_{T_\etale})
$
dans $D(S_\etale)$.
\item Pour $K$ dans $D((\Sch/T)_\etale, \mathcal{O})$, nous avons
$
(Rf_{big, *}K)|_{S_\etale} = Rf_{small, *}(K|_{T_\etale})
$
dans $D(\textit{Mod}(S_\etale, \mathcal{O}_S))$.
\end{enumerate}
Plus généralement, soit $g : S' \to S$ un objet de $(\Sch/S)_\etale$.
Considérons le produit fibré
$$
\xymatrix{
T' \ar[r]_{g'} \ar[d]_{f'} & T \ar[d]^f \\
S' \ar[r]^g & S
}
$$
Alors
\begin{enumerate}
\item[(3)] Pour $K$ dans $D((\Sch/T)_\etale)$, nous avons
$i_g^{-1}(Rf_{big, *}K) = Rf'_{small, *}(i_{g'}^{-1}K)$
dans $D(S'_\etale)$.
\item[(4)] Pour $K$ dans $D((\Sch/T)_\etale, \mathcal{O})$, nous avons
$i_g^*(Rf_{big, *}K) = Rf'_{small, *}(i_{g'}^*K)$
dans $D(\textit{Mod}(S'_\etale, \mathcal{O}_{S'}))$.
\item[(5)] Pour $K$ dans $D((\Sch/T)_\etale)$, nous avons
$g_{big}^{-1}(Rf_{big, *}K) = Rf'_{big, *}((g'_{big})^{-1}K)$
dans $D((\Sch/S')_\etale)$.
\item[(6)] Pour $K$ dans $D((\Sch/T)_\etale, \mathcal{O})$, nous avons
$g_{big}^*(Rf_{big, *}K) = Rf'_{big, *}((g'_{big})^*K)$
dans $D(\textit{Mod}(S'_\etale, \mathcal{O}_{S'}))$.
\end{enumerate}
\end{lemma}

\begin{proof}
La partie (1) découle du
lemme \ref{lemma-compare-injectives}
et de (\ref{equation-compare-big-small}),
en choisissant un complexe K-injectif de faisceaux abéliens représentant $K$.

\medskip\noindent
La partie (3) découle du lemme \ref{lemma-compare-injectives}
et du lemme de Topologies
\ref{topologies-lemma-morphism-big-small-cartesian-diagram-etale}
en choisissant un complexe K-injectif de faisceaux abéliens représentant $K$.

\medskip\noindent
La partie (5) est le lemme de Cohomologie sur les sites
\ref{sites-cohomology-lemma-localize-cartesian-square}.

\medskip\noindent
La partie (6) est le lemme de Cohomologie sur les sites
\ref{sites-cohomology-lemma-localize-cartesian-square-modules}.

\medskip\noindent
La partie (2) peut se démontrer comme suit. Nous avons vu ci-dessus
que $\pi_S \circ f_{big} = f_{small} \circ \pi_T$ comme morphismes
de sites annelés. Nous obtenons donc
$R\pi_{S, *} \circ Rf_{big, *} = Rf_{small, *} \circ R\pi_{T, *}$
par le lemme de Cohomologie sur les sites
\ref{sites-cohomology-lemma-derived-pushforward-composition}.
Puisque les foncteurs de restriction $\pi_{S, *}$ et $\pi_{T, *}$
sont exacts, nous concluons.

\medskip\noindent
La partie (4) découle de la partie (6) et de la partie (2) appliquée à $f' : T' \to S'$.
\end{proof}

\noindent
Soit $S$ un schéma et soit $\mathcal{H}$ un faisceau abélien sur
$(\Sch/S)_\etale$. Rappelons que $H^n_\etale(U, \mathcal{H})$ désigne
la cohomologie de $\mathcal{H}$ sur un objet $U$ de $(\Sch/S)_\etale$.

\begin{lemma}
\label{lemma-compare-cohomology}
Soit $f : T \to S$ un morphisme de schémas. Alors
\begin{enumerate}
\item Pour $K$ dans $D(S_\etale)$, nous avons
$H^n_\etale(S, \pi_S^{-1}K) = H^n(S_\etale, K)$.
\item Pour $K$ dans $D(S_\etale, \mathcal{O}_S)$, nous avons
$H^n_\etale(S, L\pi_S^*K) = H^n(S_\etale, K)$.
\item Pour $K$ dans $D(S_\etale)$, nous avons
$H^n_\etale(T, \pi_S^{-1}K) = H^n(T_\etale, f_{small}^{-1}K)$.
\item Pour $K$ dans $D(S_\etale, \mathcal{O}_S)$, nous avons
$H^n_\etale(T, L\pi_S^*K) = H^n(T_\etale, Lf_{small}^*K)$.
\item Pour $M$ dans $D((\Sch/S)_\etale)$, nous avons
$H^n_\etale(T, M) = H^n(T_\etale, i_f^{-1}M)$.
\item Pour $M$ dans $D((\Sch/S)_\etale, \mathcal{O})$, nous avons
$H^n_\etale(T, M) = H^n(T_\etale, i_f^*M)$.
\end{enumerate}
\end{lemma}

\begin{proof}
Pour démontrer (5), représentons $M$ par un complexe K-injectif de faisceaux abéliens,
appliquons le lemme \ref{lemma-compare-injectives}
et explicitons les définitions. La partie (3) en découle,
car $i_f^{-1}\pi_S^{-1} = f_{small}^{-1}$. La partie (1) est un cas
particulier de (3).

\medskip\noindent
La partie (6) découle du très général lemme de Cohomologie sur les sites
\ref{sites-cohomology-lemma-pullback-same-cohomology}. La partie
(4) en découle, car $Lf_{small}^* = i_f^* \circ L\pi_S^*$.
La partie (2) est un cas particulier de (4).
\end{proof}

\begin{lemma}
\label{lemma-cohomological-descent-etale}
Soit $S$ un schéma. Pour $K \in D(S_\etale)$, le morphisme
$$
K \longrightarrow R\pi_{S, *}\pi_S^{-1}K
$$
est un isomorphisme.
\end{lemma}

\begin{proof}
C'est vrai parce que $\pi_S^{-1}$ et $\pi_{S, *} = i_S^{-1}$
sont tous deux des foncteurs exacts, et que le composé $\pi_{S, *} \circ \pi_S^{-1}$
est le foncteur identité.
\end{proof}

\begin{lemma}
\label{lemma-compare-higher-direct-image-proper}
Soit $f : T \to S$ un morphisme propre de schémas. Nous avons alors
\begin{enumerate}
\item $\pi_S^{-1} \circ f_{small, *} = f_{big, *} \circ \pi_T^{-1}$
comme foncteurs $\Sh(T_\etale) \to \Sh((\Sch/S)_\etale)$,
\item $\pi_S^{-1}Rf_{small, *}K = Rf_{big, *}\pi_T^{-1}K$
pour $K$ dans $D^+(T_\etale)$ dont les faisceaux de cohomologie sont de torsion,
\item $\pi_S^{-1}Rf_{small, *}K = Rf_{big, *}\pi_T^{-1}K$
pour $K$ dans $D(T_\etale, \mathbf{Z}/n\mathbf{Z})$, et
\item $\pi_S^{-1}Rf_{small, *}K = Rf_{big, *}\pi_T^{-1}K$
pour tout $K$ dans $D(T_\etale)$ si $f$ est fini.
\end{enumerate}
\end{lemma}

\begin{proof}
Démonstration de (1). Soit $\mathcal{F}$ un faisceau sur $T_\etale$.
Soit $g : S' \to S$ un objet de $(\Sch/S)_\etale$. Considérons le
produit fibré
$$
\xymatrix{
T' \ar[r]_{f'} \ar[d]_{g'} & S' \ar[d]^g \\
T \ar[r]^f & S
}
$$
Nous avons alors
$$
(f_{big, *}\pi_T^{-1}\mathcal{F})(S') =
(\pi_T^{-1}\mathcal{F})(T') =
((g'_{small})^{-1}\mathcal{F})(T')  =
(f'_{small, *}(g'_{small})^{-1}\mathcal{F})(S')
$$
où la deuxième égalité vient du lemme \ref{lemma-describe-pullback}.
D'autre part,
$$
(\pi_S^{-1}f_{small, *}\mathcal{F})(S') =
(g_{small}^{-1}f_{small, *}\mathcal{F})(S')
$$
de nouveau par le lemme \ref{lemma-describe-pullback}.
Le changement de base propre pour les faisceaux d'ensembles
(lemme \ref{lemma-proper-base-change-f-star})
montre donc que les deux ensembles sont canoniquement isomorphes.
L'isomorphisme est compatible avec les morphismes de restriction
et définit un isomorphisme
$\pi_S^{-1}f_{small, *}\mathcal{F} = f_{big, *}\pi_T^{-1}\mathcal{F}$.
Il définit ainsi un isomorphisme de foncteurs
$\pi_S^{-1} \circ f_{small, *} = f_{big, *} \circ \pi_T^{-1}$.

\medskip\noindent
Démonstration de (2). Il existe un morphisme canonique de changement de base
$\pi_S^{-1}Rf_{small, *}K \to Rf_{big, *}\pi_T^{-1}K$
pour tout $K$ dans $D(T_\etale)$; voir la
remarque \ref{sites-cohomology-remark-base-change} de Cohomologie sur les sites.
Pour montrer que c'est un isomorphisme, il suffit de montrer que l'image inverse du
morphisme de changement de base par $i_g : \Sh(S'_\etale) \to \Sh((\Sch/S)_\etale)$
est un isomorphisme pour tout objet $g : S' \to S$ de $(\Sch/S)_\etale$.
Soient $T', g', f'$ comme au paragraphe précédent.
L'image inverse du morphisme de changement de base est
\begin{align*}
g_{small}^{-1}Rf_{small, *}K
& =
i_g^{-1}\pi_S^{-1}Rf_{small, *}K \\
& \to
i_g^{-1}Rf_{big, *}\pi_T^{-1}K \\
& =
Rf'_{small, *}(i_{g'}^{-1}\pi_T^{-1}K) \\
& =
Rf'_{small, *}((g'_{small})^{-1}K)
\end{align*}
où nous avons employé $\pi_S \circ i_g = g_{small}$,
$\pi_T \circ i_{g'} = g'_{small}$ et le
lemme \ref{lemma-compare-higher-direct-image}.
Ce morphisme est un isomorphisme par le théorème de changement de base pour un morphisme propre
(lemme \ref{lemma-proper-base-change}), pourvu que $K$ soit borné
inférieurement et que les faisceaux de cohomologie de $K$ soient de torsion.

\medskip\noindent
La démonstration de la partie (3) est la même que celle de la partie (2), sauf que
nous employons le lemme \ref{lemma-proper-base-change-mod-n}
au lieu du lemme \ref{lemma-proper-base-change}.

\medskip\noindent
Démonstration de (4). Si $f$ est fini, les foncteurs
$f_{small, *}$ et $f_{big, *}$ sont exacts. Cela découle
de la proposition \ref{proposition-finite-higher-direct-image-zero}
pour $f_{small}$. Comme tout changement de base $f'$ de $f$ est également fini,
la partie (3) du lemme \ref{lemma-compare-higher-direct-image} montre
que $f_{big, *}$ est lui aussi exact (ses foncteurs dérivés supérieurs étant nuls).
Ce cas découle donc de la partie (1).
\end{proof}




\section{Comparaison des topologies fppf et étale}
\label{section-fppf-etale}

\noindent
Un modèle pour cette section est la section consacrée à la comparaison de la
topologie usuelle et de la topologie qc sur les espaces topologiques
localement compacts, étudiée dans la section de Cohomologie sur les sites
\ref{sites-cohomology-section-cohomology-LC}.
Nous rappelons d'abord certains résultats des
sections de Topologies
\ref{topologies-section-change-topologies} et
\ref{topologies-section-etale}.

\medskip\noindent
Soit $S$ un schéma et soit $(\Sch/S)_{fppf}$ un site fppf.
Sur la même catégorie sous-jacente, nous avons une seconde topologie,
à savoir la topologie étale, et donc un second site
$(\Sch/S)_\etale$. Le foncteur identité
$(\Sch/S)_\etale \to (\Sch/S)_{fppf}$ est continu et définit
un morphisme de sites
$$
\epsilon_S : (\Sch/S)_{fppf} \longrightarrow (\Sch/S)_\etale
$$
Voir la section \ref{sites-cohomology-section-compare} de Cohomologie sur les sites.
Remarquons que $\epsilon_{S, *}$ est le foncteur identité sur les
préfaisceaux sous-jacents et que $\epsilon_S^{-1}$ associe à un faisceau étale son
faisceau fppf associé. Soit $S_\etale$ le petit site étale.
Il existe un morphisme de sites
$$
\pi_S : (\Sch/S)_\etale \longrightarrow S_\etale
$$
induit par le foncteur continu
$S_\etale \to (\Sch/S)_\etale$, $U \mapsto U$.
En effet, $S_\etale$ possède des produits fibrés et un objet final, le
foncteur ci-dessus y commute, et la
proposition \ref{sites-proposition-get-morphism} de Sites s'applique.

\begin{lemma}
\label{lemma-describe-pullback-pi-fppf}
Avec les notations ci-dessus.
Soit $\mathcal{F}$ un faisceau sur $S_\etale$. La règle
$$
(\Sch/S)_{fppf} \longrightarrow \textit{Ensembles},\quad
(f : X \to S) \longmapsto \Gamma(X, f_{small}^{-1}\mathcal{F})
$$
est un faisceau fppf, et a fortiori un faisceau sur $(\Sch/S)_\etale$.
En fait, ce faisceau est égal à
$\pi_S^{-1}\mathcal{F}$ sur $(\Sch/S)_\etale$ et à
$\epsilon_S^{-1}\pi_S^{-1}\mathcal{F}$ sur $(\Sch/S)_{fppf}$.
\end{lemma}

\begin{proof}
L'énoncé relatif à la topologie étale est le contenu
du lemme \ref{lemma-describe-pullback}. Pour achever la démonstration, il
suffit de montrer que $\pi_S^{-1}\mathcal{F}$ est un faisceau pour la topologie
fppf. Cela est démontré dans le lemme \ref{lemma-describe-pullback}
également.
\end{proof}

\noindent
Dans la situation du lemme \ref{lemma-describe-pullback-pi-fppf},
la composée de $\epsilon_S$ et $\pi_S$, ainsi que l'égalité,
déterminent un morphisme de sites
$$
a_S : (\Sch/S)_{fppf} \longrightarrow S_\etale
$$

\begin{lemma}
\label{lemma-push-pull-fppf-etale}
Avec les notations ci-dessus.
Soit $f : X \to Y$ un morphisme de $(\Sch/S)_{fppf}$.
Il existe alors des diagrammes commutatifs de topos
$$
\xymatrix{
\Sh((\Sch/X)_{fppf}) \ar[rr]_{f_{big, fppf}} \ar[d]_{\epsilon_X} & &
\Sh((\Sch/Y)_{fppf}) \ar[d]^{\epsilon_Y} \\
\Sh((\Sch/X)_\etale) \ar[rr]^{f_{big, \etale}} & &
\Sh((\Sch/Y)_\etale)
}
$$
et
$$
\xymatrix{
\Sh((\Sch/X)_{fppf}) \ar[rr]_{f_{big, fppf}} \ar[d]_{a_X} & &
\Sh((\Sch/Y)_{fppf}) \ar[d]^{a_Y} \\
\Sh(X_\etale) \ar[rr]^{f_{small}} & &
\Sh(Y_\etale)
}
$$
avec $a_X = \pi_X \circ \epsilon_X$ et $a_Y = \pi_X \circ \epsilon_X$.
\end{lemma}

\begin{proof}
La commutativité des diagrammes découle de la discussion de la
section \ref{topologies-section-change-topologies} de Topologies.
\end{proof}

\begin{lemma}
\label{lemma-proper-push-pull-fppf-etale}
Dans le lemme \ref{lemma-push-pull-fppf-etale}, si $f$ est propre, alors nous avons
$a_Y^{-1} \circ f_{small, *} = f_{big, fppf, *} \circ a_X^{-1}$.
\end{lemma}

\begin{proof}
On peut le démontrer en répétant la démonstration de la
partie (1) du lemme \ref{lemma-compare-higher-direct-image-proper};
nous allons plutôt l'en déduire.
Comme $\epsilon_{Y, *}$ est le foncteur identité sur les préfaisceaux sous-jacents,
il reflète les isomorphismes. La description
du lemme \ref{lemma-describe-pullback-pi-fppf}
montre que $\epsilon_{Y, *} \circ a_Y^{-1} = \pi_Y^{-1}$,
et de même pour $X$. Pour montrer que le morphisme canonique
$a_Y^{-1}f_{small, *}\mathcal{F} \to f_{big, fppf, *}a_X^{-1}\mathcal{F}$
est un isomorphisme, il suffit de montrer que
\begin{align*}
\pi_Y^{-1}f_{small, *}\mathcal{F}
& =
\epsilon_{Y, *}a_Y^{-1}f_{small, *}\mathcal{F} \\
& \to 
\epsilon_{Y, *}f_{big, fppf, *}a_X^{-1}\mathcal{F} \\
& =
f_{big, \etale, *} \epsilon_{X, *}a_X^{-1}\mathcal{F} \\
& =
f_{big, \etale, *}\pi_X^{-1}\mathcal{F}
\end{align*}
est un isomorphisme. C'est la partie
(1) du lemme \ref{lemma-compare-higher-direct-image-proper}.
\end{proof}

\begin{lemma}
\label{lemma-descent-sheaf-fppf-etale}
Dans le lemme \ref{lemma-push-pull-fppf-etale}, supposons
$f$ plat, localement de présentation finie et surjectif.
Alors le foncteur
$$
\Sh(Y_\etale) \longrightarrow
\left\{
(\mathcal{G}, \mathcal{H}, \alpha)
\middle|
\begin{matrix}
\mathcal{G} \in \Sh(X_\etale),\ \mathcal{H} \in \Sh((\Sch/Y)_{fppf}), \\
\alpha : a_X^{-1}\mathcal{G} \to f_{big, fppf}^{-1}\mathcal{H}
\text{ un isomorphisme}
\end{matrix}
\right\}
$$
qui associe à $\mathcal{F}$
$(f_{small}^{-1}\mathcal{F}, a_Y^{-1}\mathcal{F}, can)$ est une équivalence.
\end{lemma}

\begin{proof}
Le foncteur $a_X^{-1}$ est pleinement fidèle (car $a_{X, *}a_X^{-1} = \text{id}$ par le
lemme \ref{lemma-describe-pullback-pi-fppf}). Le foncteur d'oubli
$(\mathcal{G}, \mathcal{H}, \alpha) \mapsto \mathcal{H}$ identifie donc la
catégorie des triplets à une sous-catégorie pleine de $\Sh((\Sch/Y)_{fppf})$.
De plus, le foncteur $a_Y^{-1}$ est pleinement fidèle; le foncteur
du lemme est donc lui aussi pleinement fidèle.

\medskip\noindent
Supposons donné un recouvrement étale $\{Y_i \to Y\}$.
Soit $f_i : X_i \to Y_i$ le changement de base de $f$.
Notons $f_{ij} = f_i \times f_j : X_i \times_X X_j  \to Y_i \times_Y Y_j$.
Assertion : si le lemme est vrai pour $f_i$ et $f_{ij}$ pour tous $i, j$, alors
il est vrai pour $f$. Pour le voir, remarquons que le recouvrement étale donné
détermine un recouvrement étale de l'objet final dans chacun des
quatre sites $Y_\etale, X_\etale, (\Sch/Y)_{fppf}, (\Sch/X)_{fppf}$.
La catégorie des faisceaux équivaut donc à la catégorie des données de
recollement pour ce recouvrement
(lemme \ref{sites-lemma-mapping-property-glue} de Sites)
dans chacun des quatre cas. Un grand diagramme commutatif de
catégories achève alors la démonstration de l'affirmation. Nous omettons les détails.
L'assertion montre que nous pouvons travailler localement sur $Y$ pour la topologie étale.

\medskip\noindent
Remarquons que $\{X \to Y\}$ est un recouvrement fppf. En travaillant localement sur $Y$,
nous pouvons supposer qu'il existe un morphisme $s : X' \to X$
tel que le composé $f' = f \circ s : X' \to Y$ soit fini localement libre et surjectif ; voir le
lemme de Compléments sur les morphismes
\ref{more-morphisms-lemma-dominate-fppf-etale-locally}.
Assertion : si le lemme est vrai pour $f'$, alors il est vrai pour $f$.
En effet, étant donné un triplet $(\mathcal{G}, \mathcal{H}, \alpha)$
pour $f$, son image inverse par $s$ donne un triplet
$(s_{small}^{-1}\mathcal{G}, \mathcal{H}, s_{big, fppf}^{-1}\alpha)$
pour $f'$. Une solution de ce triplet donne un faisceau
$\mathcal{F}$ sur $Y_\etale$ tel que $a_Y^{-1}\mathcal{F} = \mathcal{H}$.
D'après le premier paragraphe de la démonstration, cela signifie que le triplet
appartient à l'image essentielle. Nous sommes ainsi ramenés
au cas décrit au paragraphe suivant.

\medskip\noindent
Supposons $f$ fini localement libre et surjectif.
Soit $(\mathcal{G}, \mathcal{H}, \alpha)$ un triplet.
Considérons dans ce cas le triplet
$$
(\mathcal{G}_1, \mathcal{H}_1, \alpha_1) =
(f_{small}^{-1}f_{small, *}\mathcal{G},
f_{big, fppf, *}f_{big, fppf}^{-1}\mathcal{H}, \alpha_1)
$$
où $\alpha_1$ est défini à partir des identifications
\begin{align*}
a_X^{-1}f_{small}^{-1}f_{small, *}\mathcal{G}
& =
f_{big, fppf}^{-1}a_Y^{-1}f_{small, *}\mathcal{G} \\
& =
f_{big, fppf}^{-1}f_{big, fppf, *}a_X^{-1}\mathcal{G} \\
& \to
f_{big, fppf}^{-1}f_{big, fppf, *}f_{big, fppf}^{-1}\mathcal{H}
\end{align*}
où la troisième égalité vient du lemme \ref{lemma-proper-push-pull-fppf-etale}
et la flèche est donnée par $\alpha$.
Ce triplet appartient à l'image de notre foncteur, car
$\mathcal{F}_1 = f_{small, *}\mathcal{F}$ constitue une solution
(pour le voir, employer de nouveau le lemme \ref{lemma-proper-push-pull-fppf-etale};
détails omis). Il existe un morphisme canonique de triplets
$$
(\mathcal{G}, \mathcal{H}, \alpha)
\to
(\mathcal{G}_1, \mathcal{H}_1, \alpha_1)
$$
qui emploie l'unité $\text{id} \to f_{big, fppf, *}f_{big, fppf}^{-1}$
sur la seconde composante (il suffit de prescrire les morphismes sur la
seconde composante, d'après le premier paragraphe de la démonstration). Puisque
$\{f : X \to Y\}$ est un recouvrement fppf, le morphisme
$\mathcal{H} \to \mathcal{H}_1$ est injectif (détails omis).
Posons
$$
\mathcal{G}_2 = \mathcal{G}_1 \amalg_\mathcal{G} \mathcal{G}_1\quad
\mathcal{H}_2 = \mathcal{H}_1 \amalg_\mathcal{H} \mathcal{H}_1
$$
et soit $\alpha_2$ l'isomorphisme induit (les foncteurs d'image inverse
sont exacts, donc cette construction a un sens). Alors $\mathcal{H}$ est
l'égalisateur des deux morphismes $\mathcal{H}_1 \to \mathcal{H}_2$.
En répétant les arguments ci-dessus pour le triplet
$(\mathcal{G}_2, \mathcal{H}_2, \alpha_2)$
nous trouvons un morphisme injectif de triplets
$$
(\mathcal{G}_2, \mathcal{H}_2, \alpha_2)
\to
(\mathcal{G}_3, \mathcal{H}_3, \alpha_3)
$$
tel que ce dernier triplet appartienne à l'image de notre foncteur.
Disons qu'il correspond à $\mathcal{F}_3$ dans $\Sh(Y_\etale)$.
Par pleine fidélité, nous obtenons deux morphismes
$\mathcal{F}_1 \to \mathcal{F}_3$, et nous pouvons prendre pour
$\mathcal{F}$ leur égalisateur.
Par exactitude des foncteurs d'image inverse considérés, nous
trouvons $a_Y^{-1}\mathcal{F} = \mathcal{H}$, comme voulu.
\end{proof}

\begin{lemma}
\label{lemma-compare-fppf-etale}
Considérons le morphisme de comparaison
$\epsilon : (\Sch/S)_{fppf} \to (\Sch/S)_\etale$.
Soit $\mathcal{P}$ la classe des morphismes finis de schémas.
Pour $X$ dans $(\Sch/S)_\etale$, notons
$\mathcal{A}'_X \subset \textit{Ab}((\Sch/X)_\etale)$
la sous-catégorie pleine formée des faisceaux de la forme
$\pi_X^{-1}\mathcal{F}$ avec $\mathcal{F}$ dans $\textit{Ab}(X_\etale)$.
Alors les propriétés de Cohomologie sur les sites
(\ref{sites-cohomology-item-base-change-P}),
(\ref{sites-cohomology-item-restriction-A}),
(\ref{sites-cohomology-item-A-sheaf}),
(\ref{sites-cohomology-item-A-and-P}) et
(\ref{sites-cohomology-item-refine-tau-by-P})
énumérées dans la situation de Cohomologie sur les sites
\ref{sites-cohomology-situation-compare} sont satisfaites.
\end{lemma}

\begin{proof}
Montrons d'abord que $\mathcal{A}'_X \subset \textit{Ab}((\Sch/X)_\etale)$
est une sous-catégorie de Serre faible en vérifiant les conditions (1), (2), (3) et (4)
du lemme \ref{homology-lemma-characterize-weak-serre-subcategory} d'Homology.
Les parties (1), (2), (3) sont immédiates, car $\pi_X^{-1}$ est exact et
pleinement fidèle, par exemple par le lemme \ref{lemma-cohomological-descent-etale}. Si
$0 \to \pi_X^{-1}\mathcal{F} \to \mathcal{G} \to \pi_X^{-1}\mathcal{F}' \to 0$
est une suite exacte courte dans $\textit{Ab}((\Sch/X)_\etale)$,
alors $0 \to \mathcal{F} \to \pi_{X, *}\mathcal{G} \to \mathcal{F}' \to 0$
est exacte par le lemme \ref{lemma-cohomological-descent-etale}.
Ainsi, $\mathcal{G} = \pi_X^{-1}\pi_{X, *}\mathcal{G}$ appartient à
$\mathcal{A}'_X$, ce qui vérifie la dernière condition.

\medskip\noindent
La propriété (\ref{sites-cohomology-item-base-change-P}) de Cohomologie sur les sites vaut
par l'existence des produits fibrés de schémas
et le fait que le changement de base d'un morphisme fini de
schémas est un morphisme fini de schémas; voir le
lemme \ref{morphisms-lemma-base-change-finite} de Morphismes.

\medskip\noindent
La propriété (\ref{sites-cohomology-item-restriction-A}) de Cohomologie sur les sites
découle du diagramme commutatif (3) du
lemme \ref{topologies-lemma-morphism-big-small-etale} de Topologies.

\medskip\noindent
La propriété (\ref{sites-cohomology-item-A-sheaf}) de Cohomologie sur les sites est le
lemme \ref{lemma-describe-pullback-pi-fppf}.

\medskip\noindent
La propriété (\ref{sites-cohomology-item-A-and-P}) de Cohomologie sur les sites vaut par la
partie (4) du lemme \ref{lemma-compare-higher-direct-image-proper}.

\medskip\noindent
La propriété (\ref{sites-cohomology-item-refine-tau-by-P}) de Cohomologie sur les sites
est impliquée par le
lemme de Compléments sur les morphismes
\ref{more-morphisms-lemma-dominate-fppf-etale-locally}.
\end{proof}

\begin{lemma}
\label{lemma-V-C-all-n-etale-fppf}
Avec les notations ci-dessus.
\begin{enumerate}
\item Pour $X \in \Ob((\Sch/S)_{fppf})$ et un faisceau abélien $\mathcal{F}$
sur $X_\etale$, nous avons
$\epsilon_{X, *}a_X^{-1}\mathcal{F} = \pi_X^{-1}\mathcal{F}$
et $R^i\epsilon_{X, *}(a_X^{-1}\mathcal{F}) = 0$ pour $i > 0$.
\item Pour un morphisme fini $f : X \to Y$ dans $(\Sch/S)_{fppf}$
et un faisceau abélien $\mathcal{F}$ sur $X$, nous avons
$a_Y^{-1}(R^if_{small, *}\mathcal{F}) =
R^if_{big, fppf, *}(a_X^{-1}\mathcal{F})$
pour tout $i$.
\item Pour un schéma $X$ et $K$ dans $D^+(X_\etale)$, le morphisme
$\pi_X^{-1}K \to R\epsilon_{X, *}(a_X^{-1}K)$ est un isomorphisme.
\item Pour un morphisme fini $f : X \to Y$ de schémas
et $K$ dans $D^+(X_\etale)$, nous avons
$a_Y^{-1}(Rf_{small, *}K) = Rf_{big, fppf, *}(a_X^{-1}K)$.
\item Pour un morphisme propre $f : X \to Y$ de schémas
et $K$ dans $D^+(X_\etale)$ à faisceaux de cohomologie de torsion, nous avons
$a_Y^{-1}(Rf_{small, *}K) = Rf_{big, fppf, *}(a_X^{-1}K)$.
\end{enumerate}
\end{lemma}

\begin{proof}
Par le lemme \ref{lemma-compare-fppf-etale}, tous les lemmes de la
section \ref{sites-cohomology-section-compare-general} de Cohomologie sur les sites
s'appliquent à notre situation. Pour traduire les résultats,
remarquons que la catégorie $\mathcal{A}_X$ du
lemme \ref{sites-cohomology-lemma-A} de Cohomologie sur les sites
est l'image essentielle de
$a_X^{-1} : \textit{Ab}(X_\etale) \to \textit{Ab}((\Sch/X)_{fppf})$.

\medskip\noindent
La partie (1) équivaut à $(V_n)$ pour tout $n$, ce qui vaut par le
lemme \ref{sites-cohomology-lemma-V-C-all-n-general} de Cohomologie sur les sites.

\medskip\noindent
La partie (2) découle de l'application de $\epsilon_Y^{-1}$ à la conclusion du
lemme \ref{sites-cohomology-lemma-V-implies-C-general} de Cohomologie sur les sites.

\medskip\noindent
La partie (3) découle du lemme de Cohomologie sur les sites
\ref{sites-cohomology-lemma-V-C-all-n-general}, partie (1),
car $\pi_X^{-1}K$ appartient à $D^+_{\mathcal{A}'_X}((\Sch/X)_\etale)$
et $a_X^{-1} = \epsilon_X^{-1} \circ a_X^{-1}$.

\medskip\noindent
La partie (4) découle du lemme de Cohomologie sur les sites
\ref{sites-cohomology-lemma-V-C-all-n-general}, partie (2),
pour la même raison.

\medskip\noindent
Partie (5). Nous employons
\begin{align*}
R\epsilon_{Y, *}Rf_{big, fppf, *}a_X^{-1}K
& =
Rf_{big, \etale, *}R\epsilon_{X, *}a_X^{-1}K \\
& =
Rf_{big, \etale, *}\pi_X^{-1}K \\
& =
\pi_Y^{-1}Rf_{small, *}K \\
& =
R\epsilon_{Y, *} a_Y^{-1}Rf_{small, *}K
\end{align*}
La première égalité vient du diagramme commutatif du
lemme \ref{lemma-push-pull-fppf-etale}
et du lemme de Cohomologie sur les sites
\ref{sites-cohomology-lemma-derived-pushforward-composition}.
La deuxième égalité est (3). La troisième est la partie (2) du
lemme \ref{lemma-compare-higher-direct-image-proper}.
La quatrième est de nouveau (3). Le morphisme de changement de base
$a_Y^{-1}(Rf_{small, *}K) \to Rf_{big, fppf, *}(a_X^{-1}K)$
induit donc un isomorphisme
$$
R\epsilon_{Y, *}a_Y^{-1}Rf_{small, *}K \to
R\epsilon_{Y, *}Rf_{big, fppf, *}a_X^{-1}K
$$
La remarque suivante achève la démonstration : un morphisme
$\alpha : a_Y^{-1}L \to M$, avec $L$ dans $D^+(Y_\etale)$
et $M$ dans $D^+((\Sch/Y)_{fppf})$, tel que $R\epsilon_{Y, *}\alpha$
soit un isomorphisme, est un isomorphisme. En effet,
nous montrons par récurrence sur $i$ que $H^i(\alpha)$ est un isomorphisme.
C'est vrai pour tout $i$ assez petit.
Si cela vaut pour $i \leq i_0$, nous voyons que
$R^j\epsilon_{Y, *}H^i(M) = 0$ pour $j > 0$ et $i \leq i_0$
par (1), puisque $H^i(M) = a_Y^{-1}H^i(L)$ dans cet intervalle.
Ainsi, $\epsilon_{Y, *}H^{i_0 + 1}(M) = H^{i_0 + 1}(R\epsilon_{Y, *}M)$
par un argument de suite spectrale.
Ainsi, $\epsilon_{Y, *}H^{i_0 + 1}(M) = \pi_Y^{-1}H^{i_0 + 1}(L) =
\epsilon_{Y, *}a_Y^{-1}H^{i_0 + 1}(L)$.
Cela implique que $H^{i_0 + 1}(\alpha)$ est un isomorphisme
(car $\epsilon_{Y, *}$ reflète les isomorphismes, puisqu'il est le
foncteur identité sur les préfaisceaux sous-jacents), comme voulu.
\end{proof}

\begin{lemma}
\label{lemma-cohomological-descent-etale-fppf}
Soit $X$ un schéma. Pour $K \in D^+(X_\etale)$, le morphisme
$$
K \longrightarrow Ra_{X, *}a_X^{-1}K
$$
est un isomorphisme, où $a_X : \Sh((\Sch/X)_{fppf}) \to \Sh(X_\etale)$
est comme ci-dessus.
\end{lemma}

\begin{proof}
Nous nous ramenons d'abord au cas où
$K$ est donné par un unique faisceau abélien. Représentons en effet $K$
par un complexe borné inférieurement $\mathcal{F}^\bullet$. Le cas d'un
faisceau montre que
$\mathcal{F}^n = a_{X, *} a_X^{-1} \mathcal{F}^n$
et que les faisceaux $R^qa_{X, *}a_X^{-1}\mathcal{F}^n$
sont nuls pour $q > 0$. Le lemme d'acyclicité de Leray
(lemme \ref{derived-lemma-leray-acyclicity} de Catégories dérivées),
appliqué à $a_X^{-1}\mathcal{F}^\bullet$
et au foncteur $a_{X, *}$, permet de conclure. Supposons désormais $K = \mathcal{F}$.

\medskip\noindent
Par le lemme \ref{lemma-describe-pullback-pi-fppf}, nous avons
$a_{X, *}a_X^{-1}\mathcal{F} = \mathcal{F}$. Il suffit donc de montrer que
$R^qa_{X, *}a_X^{-1}\mathcal{F} = 0$ pour $q > 0$.
Pour cela, nous pouvons employer $a_X = \epsilon_X \circ \pi_X$ et
la suite spectrale de Leray
(lemme \ref{sites-cohomology-lemma-relative-Leray} de Cohomologie sur les sites).
Par le lemme \ref{lemma-V-C-all-n-etale-fppf},
nous avons $R^i\epsilon_{X, *}(a_X^{-1}\mathcal{F}) = 0$ pour $i > 0$
et
$\epsilon_{X, *}a_X^{-1}\mathcal{F} = \pi_X^{-1}\mathcal{F}$.
Par le lemme \ref{lemma-cohomological-descent-etale}, nous avons
$R^j\pi_{X, *}(\pi_X^{-1}\mathcal{F}) = 0$ pour $j > 0$.
Cela achève la démonstration.
\end{proof}

\begin{lemma}
\label{lemma-compare-cohomology-etale-fppf}
Pour un schéma $X$ et $a_X : \Sh((\Sch/X)_{fppf}) \to \Sh(X_\etale)$
comme ci-dessus :
\begin{enumerate}
\item $H^q(X_\etale, \mathcal{F}) = H^q_{fppf}(X, a_X^{-1}\mathcal{F})$
pour un faisceau abélien $\mathcal{F}$ sur $X_\etale$,
\item $H^q(X_\etale, K) = H^q_{fppf}(X, a_X^{-1}K)$ pour $K \in D^+(X_\etale)$.
\end{enumerate}
Exemple : si $A$ est un groupe abélien, alors
$H^q_\etale(X, \underline{A}) = H^q_{fppf}(X, \underline{A})$.
\end{lemma}

\begin{proof}
Cela découle du lemme \ref{lemma-cohomological-descent-etale-fppf}
par la remarque \ref{sites-cohomology-remark-before-Leray} de Cohomologie sur les sites.
\end{proof}







\section{Comparaison des topologies fppf et étale : modules}
\label{section-fppf-etale-modules}

\noindent
Nous poursuivons la discussion de la section \ref{section-fppf-etale}, mais nous
examinons ici brièvement ce qui se passe pour les faisceaux de modules.

\medskip\noindent
Soit $S$ un schéma. Les morphismes de sites $\epsilon_S$, $\pi_S$ et
leur composé $a_S$, introduits dans la section \ref{section-fppf-etale},
s'étendent naturellement en morphismes de sites annelés. Le premier
s'écrit
$$
\epsilon_S :
((\Sch/S)_{fppf}, \mathcal{O})
\longrightarrow
((\Sch/S)_\etale, \mathcal{O})
$$
Remarquons que nous pouvons employer le même symbole pour le faisceau structural, car
ces faisceaux ont le même préfaisceau sous-jacent. Le deuxième est
$$
\pi_S :
((\Sch/S)_\etale, \mathcal{O})
\longrightarrow
(S_\etale, \mathcal{O}_S)
$$
Le troisième est le morphisme
$$
a_S :
((\Sch/S)_{fppf}, \mathcal{O})
\longrightarrow
(S_\etale, \mathcal{O}_S)
$$
Nous savons déjà que la catégorie des modules quasi-cohérents sur le
schéma $S$ est la même que celle des modules quasi-cohérents
sur $(S_\etale, \mathcal{O}_S)$; voir la
proposition \ref{descent-proposition-equivalence-quasi-coherent} de Descente.
Comme nous voulons énoncer une comparaison entre les cohomologies
étale et fppf, nous considérerons dans le reste de cette
section les faisceaux quasi-cohérents du point de vue du
petit site étale.
Rappelons ce que nous savons déjà des modules quasi-cohérents
sur ces sites.

\begin{lemma}
\label{lemma-review-quasi-coherent}
Soit $S$ un schéma. Soit $\mathcal{F}$ un
$\mathcal{O}_S$-module quasi-cohérent sur $S_\etale$.
\begin{enumerate}
\item La règle
$$
\mathcal{F}^a : (\Sch/S)_\etale \longrightarrow \textit{Ab},\quad
(f : T \to S) \longmapsto \Gamma(T, f_{small}^*\mathcal{F})
$$
satisfait la condition de faisceau pour les recouvrements fppf, et a fortiori étales,
\item $\mathcal{F}^a = \pi_S^*\mathcal{F}$ sur $(\Sch/S)_\etale$,
\item $\mathcal{F}^a = a_S^*\mathcal{F}$ sur $(\Sch/S)_{fppf}$,
\item la règle $\mathcal{F} \mapsto \mathcal{F}^a$ définit
une équivalence entre les $\mathcal{O}_S$-modules quasi-cohérents
et les modules quasi-cohérents sur
$((\Sch/S)_\etale, \mathcal{O})$,
\item la règle $\mathcal{F} \mapsto \mathcal{F}^a$ définit
une équivalence entre les $\mathcal{O}_S$-modules quasi-cohérents
et les modules quasi-cohérents sur
$((\Sch/S)_{fppf}, \mathcal{O})$,
\item nous avons $\epsilon_{S, *}a_S^*\mathcal{F} = \pi_S^*\mathcal{F}$
et $a_{S, *}a_S^*\mathcal{F} = \mathcal{F}$,
\item nous avons $R^i\epsilon_{S, *}(a_S^*\mathcal{F}) = 0$
et $R^ia_{S, *}(a_S^*\mathcal{F}) = 0$ pour $i > 0$.
\end{enumerate}
\end{lemma}

\begin{proof}
Nous invitons vivement le lecteur à trouver sa propre démonstration de ces résultats
à partir des matériaux des
sections \ref{descent-section-quasi-coherent-sheaves},
\ref{descent-section-quasi-coherent-cohomology} et
\ref{descent-section-quasi-coherent-sheaves-bis}.

\medskip\noindent
Expliquons d'abord pourquoi la notation de ce lemme est compatible avec notre
emploi antérieur de la notation $\mathcal{F}^a$ dans les
sections \ref{section-quasi-coherent} et
\ref{section-cohomology-quasi-coherent},
ainsi que dans la
section \ref{descent-section-quasi-coherent-sheaves} de Descente.
En effet, la
proposition \ref{descent-proposition-equivalence-quasi-coherent} de Descente
montre qu'il existe un module quasi-cohérent
$\mathcal{F}_0$ sur le schéma $S$ (autrement dit sur le petit
site de Zariski) tel que $\mathcal{F}$ soit la restriction de la
règle
$$
\mathcal{F}_0^a : (\Sch/S)_\etale \longrightarrow \textit{Ab},\quad
(f : T \to S) \longmapsto \Gamma(T, f^*\mathcal{F})
$$
à la sous-catégorie $S_\etale \subset (\Sch/S)_\etale$,
où $f^*$ désigne ici l'image inverse usuelle des faisceaux de modules sur les schémas.
Puisque $\mathcal{F}_0^a$ est l'image inverse par le morphisme de
sites annelés
$$
((\Sch/S)_\etale, \mathcal{O}) \longrightarrow (S_{Zar}, \mathcal{O}_{S_{Zar}})
$$
par la remarque \ref{descent-remark-change-topologies-ringed-sites} de Descente,
il découle immédiatement (par composition des images inverses) que
$\mathcal{F}^a = \mathcal{F}_0^a$. Cela démontre la propriété de faisceau,
même pour les recouvrements fpqc, par le
lemme \ref{descent-lemma-sheaf-condition-holds} de Descente (voir aussi la
proposition \ref{proposition-quasi-coherent-sheaf-fpqc}).
Les assertions (2) et (3) découlent alors
de nouveau de la remarque \ref{descent-remark-change-topologies-ringed-sites} de Descente,
et (4) et (5) de la
proposition \ref{descent-proposition-equivalence-quasi-coherent} de Descente
(voir aussi le méta-résultat du
théorème \ref{theorem-quasi-coherent}).

\medskip\noindent
La partie (6) découle immédiatement de la description du faisceau
$\mathcal{F}^a = \pi_S^*\mathcal{F} = a_S^*\mathcal{F}$.

\medskip\noindent
Pour tout faisceau abélien $\mathcal{H}$ sur $(\Sch/S)_{fppf}$, l'
image directe supérieure $R^p\epsilon_{S, *}\mathcal{H}$ est le faisceau
associé au préfaisceau $U \mapsto H^p_{fppf}(U, \mathcal{H})$
sur $(\Sch/S)_\etale$. Voir le
lemme \ref{sites-cohomology-lemma-higher-direct-images} de Cohomologie sur les sites.
Ainsi, pour démontrer
$R^p\epsilon_{S, *}a_S^*\mathcal{F} = R^p\epsilon_{S, *}\mathcal{F}^a = 0$
pour $p > 0$, il suffit de montrer que tout schéma $U$ sur $S$
possède un recouvrement étale $\{U_i \to U\}_{i \in I}$ tel que
$H^p_{fppf}(U_i, \mathcal{F}^a) = 0$ pour $p > 0$.
En prenant un recouvrement ouvert par des affines, l'
annulation requise découle de la comparaison avec la cohomologie usuelle
(proposition \ref{descent-proposition-same-cohomology-quasi-coherent} de Descente
ou
théorème \ref{theorem-zariski-fpqc-quasi-coherent}),
ainsi que de l'annulation de la cohomologie des faisceaux quasi-cohérents
sur les schémas affines donnée par le lemme de Cohomologie des schémas
\ref{coherent-lemma-quasi-coherent-affine-cohomology-zero}.

\medskip\noindent
Pour montrer que $R^pa_{S, *}a_S^{-1}\mathcal{F} = R^pa_{S, *}\mathcal{F}^a = 0$
pour $p > 0$, nous raisonnons exactement de la même manière. Cela achève la démonstration.
\end{proof}

\begin{lemma}
\label{lemma-cohomological-descent-etale-fppf-modules}
Soit $S$ un schéma. Étant donné $\mathcal{F}$, faisceau quasi-cohérent
de $\mathcal{O}_S$-modules sur $S_\etale$, les morphismes
$$
\pi_S^*\mathcal{F} \longrightarrow R\epsilon_{S, *}(a_S^*\mathcal{F})
\quad\text{et}\quad
\mathcal{F} \longrightarrow Ra_{S, *}(a_S^*\mathcal{F})
$$
sont des isomorphismes, où
$a_S : \Sh((\Sch/S)_{fppf}) \to \Sh(S_\etale)$ est comme ci-dessus.
\end{lemma}

\begin{proof}
C'est une conséquence immédiate des
parties (6) et (7) du
lemme \ref{lemma-review-quasi-coherent}.
\end{proof}

\begin{lemma}
\label{lemma-cohomological-descent-complex-modules}
Soit $S = \Spec(A)$ un schéma affine. Soit $M^\bullet$ un complexe
de $A$-modules. Considérons le complexe $\mathcal{F}^\bullet$ de
préfaisceaux de $\mathcal{O}$-modules sur
$(\textit{Aff}/S)_{fppf}$ donné par la règle
$$
(U/S) = (\Spec(B)/\Spec(A)) \longmapsto M^\bullet \otimes_A B
$$
C'est alors un complexe de modules, et le morphisme canonique
$$
M^\bullet \longrightarrow
R\Gamma((\textit{Aff}/S)_{fppf}, \mathcal{F}^\bullet)
$$
est un quasi-isomorphisme.
\end{lemma}

\begin{proof}
Chaque $\mathcal{F}^n$ est un faisceau de modules, car il coïncide avec la
restriction du module $\mathcal{G}^n = (\widetilde{M}^n)^a$ du
lemme \ref{lemma-review-quasi-coherent} à
$(\textit{Aff}/S)_{fppf} \subset (\Sch/S)_{fppf}$.
Puisque cette inclusion définit une équivalence de topos annelés
(lemme \ref{topologies-lemma-affine-big-site-fppf} de Topologies),
nous avons
$$
R\Gamma((\textit{Aff}/S)_{fppf}, \mathcal{F}^\bullet) =
R\Gamma((\Sch/S)_{fppf}, \mathcal{G}^\bullet)
$$
Remarquons que $M^\bullet = R\Gamma(S, \widetilde{M}^\bullet)$,
par exemple par le lemme de Catégories dérivées des schémas
\ref{perfect-lemma-affine-compare-bounded}. Nous cherchons donc
à montrer que le morphisme de comparaison
$$
R\Gamma(S, \widetilde{M}^\bullet)
\longrightarrow
R\Gamma((\Sch/S)_{fppf}, (\widetilde{M}^\bullet)^a)
$$
est un isomorphisme. Si $M^\bullet$ est borné inférieurement, cela
vaut par la proposition de Descente
\ref{descent-proposition-same-cohomology-quasi-coherent}
et la première suite spectrale du lemme de Catégories dérivées
\ref{derived-lemma-two-ss-complex-functor}.
Dans le cas général, écrivons $M^\bullet = \lim M_n^\bullet$
avec $M_n^\bullet = \tau_{\geq -n}M^\bullet$. Le système
$M_n^p$ est alors stationnaire de valeur $M^p$. Nous affirmons que
$$
(\widetilde{M}^\bullet)^a = R\lim (\widetilde{M}_n^\bullet)^a
$$
En effet, il suffit de montrer que le morphisme naturel de gauche à droite
induit un isomorphisme en cohomologie sur tout objet affine
$U = \Spec(B)$ de $(\Sch/S)_{fppf}$. Pour $i \in \mathbf{Z}$ et
$n > |i|$, nous avons
$$
H^i(U, (\widetilde{M}_n^\bullet)^a) =
H^i(\tau_{\geq -n}M^\bullet \otimes_A B) =
H^i(M^\bullet \otimes_A B)
$$
La première égalité vaut par le cas borné inférieurement traité ci-dessus.
Le lemme de Cohomologie sur les sites
\ref{sites-cohomology-lemma-RGamma-commutes-with-Rlim}
montre donc l'affirmation. Enfin, nous obtenons
\begin{align*}
R\Gamma((\Sch/S)_{fppf}, (\widetilde{M}^\bullet)^a)
& =
R\Gamma((\Sch/S)_{fppf}, R\lim (\widetilde{M}_n^\bullet)^a) \\
& =
R\lim R\Gamma((\Sch/S)_{fppf}, (\widetilde{M}_n^\bullet)^a) \\
& =
R\lim M_n^\bullet \\
& =
M^\bullet
\end{align*}
comme voulu. La deuxième égalité vaut parce que $R\lim$ commute à
$R\Gamma$; voir le lemme de Cohomologie sur les sites
\ref{sites-cohomology-lemma-RGamma-commutes-with-Rlim}.
\end{proof}










\section{Comparaison des topologies ph et étale}
\label{section-ph-etale}

\noindent
Un modèle pour cette section est la section consacrée à la comparaison de la
topologie usuelle et de la topologie qc sur les espaces topologiques
localement compacts, étudiée dans la section de Cohomologie sur les sites
\ref{sites-cohomology-section-cohomology-LC}.
Nous rappelons d'abord certains résultats des
sections de Topologies
\ref{topologies-section-change-topologies} et
\ref{topologies-section-etale}.

\medskip\noindent
Soit $S$ un schéma et soit $(\Sch/S)_{ph}$ un site ph.
Sur la même catégorie sous-jacente, nous avons une seconde topologie,
à savoir la topologie étale, et donc un second site
$(\Sch/S)_\etale$. Le foncteur identité
$(\Sch/S)_\etale \to (\Sch/S)_{ph}$ est continu
(par le lemme \ref{more-morphisms-lemma-fppf-ph} de Compléments sur les morphismes
et le lemme de Topologies
\ref{topologies-lemma-zariski-etale-smooth-syntomic-fppf})
et définit un morphisme de sites
$$
\epsilon_S : (\Sch/S)_{ph} \longrightarrow (\Sch/S)_\etale
$$
Voir la section \ref{sites-cohomology-section-compare} de Cohomologie sur les sites.
Remarquons que $\epsilon_{S, *}$ est le foncteur identité sur les
préfaisceaux sous-jacents et que $\epsilon_S^{-1}$ associe à un faisceau étale son
faisceau ph associé.
Soit $S_\etale$ le petit site étale.
Il existe un morphisme de sites
$$
\pi_S : (\Sch/S)_\etale \longrightarrow S_\etale
$$
induit par le foncteur continu
$S_\etale \to (\Sch/S)_\etale$, $U \mapsto U$.
En effet, $S_\etale$ possède des produits fibrés et un objet final, le
foncteur ci-dessus y commute, et la
proposition \ref{sites-proposition-get-morphism} de Sites s'applique.

\begin{lemma}
\label{lemma-describe-pullback-pi-ph}
Avec les notations ci-dessus.
Soit $\mathcal{F}$ un faisceau sur $S_\etale$. La règle
$$
(\Sch/S)_{ph} \longrightarrow \textit{Ensembles},\quad
(f : X \to S) \longmapsto \Gamma(X, f_{small}^{-1}\mathcal{F})
$$
est un faisceau ph, et a fortiori un faisceau sur $(\Sch/S)_\etale$.
En fait, ce faisceau est égal à
$\pi_S^{-1}\mathcal{F}$ sur $(\Sch/S)_\etale$ et à
$\epsilon_S^{-1}\pi_S^{-1}\mathcal{F}$ sur $(\Sch/S)_{ph}$.
\end{lemma}

\begin{proof}
L'énoncé relatif à la topologie étale est le contenu
du lemme \ref{lemma-describe-pullback}. Pour achever la démonstration, il
suffit de montrer que $\pi_S^{-1}\mathcal{F}$ est un faisceau pour la topologie
ph. Par le lemme \ref{topologies-lemma-characterize-sheaf} de Topologies,
il suffit de montrer que, pour tout morphisme propre surjectif
$V \to U$ de schémas sur $S$, nous avons un diagramme égalisateur
$$
\xymatrix{
(\pi_S^{-1}\mathcal{F})(U) \ar[r] &
(\pi_S^{-1}\mathcal{F})(V) \ar@<1ex>[r] \ar@<-1ex>[r] &
(\pi_S^{-1}\mathcal{F})(V \times_U V)
}
$$
Posons $\mathcal{G} = \pi_S^{-1}\mathcal{F}|_{U_\etale}$.
Considérons le diagramme commutatif
$$
\xymatrix{
V \times_U V \ar[r] \ar[rd]_g \ar[d] & V \ar[d]^f \\
V \ar[r]^f & U
}
$$
Nous avons
$$
(\pi_S^{-1}\mathcal{F})(V) = \Gamma(V, f^{-1}\mathcal{G}) =
\Gamma(U, f_*f^{-1}\mathcal{G})
$$
où nous employons $f_*$ et $f^{-1}$ pour désigner les fonctorialités entre
petits sites étales. D'autre part, nous avons
$$
(\pi_S^{-1}\mathcal{F})(V \times_U V) =
\Gamma(V \times_U V, g^{-1}\mathcal{G}) =
\Gamma(U, g_*g^{-1}\mathcal{G})
$$
Les deux morphismes du diagramme égalisateur proviennent des deux morphismes
$$
f_*f^{-1}\mathcal{G} \longrightarrow g_*g^{-1}\mathcal{G}
$$
Il suffit donc de démontrer que $\mathcal{G}$ est
l'égalisateur de ces deux morphismes de faisceaux.
Soit $\overline{u}$ un point géométrique de $U$. Posons
$\Omega = \mathcal{G}_{\overline{u}}$.
En prenant les fibres en $\overline{u}$ par le
lemme \ref{lemma-proper-pushforward-stalk},
nous obtenons les deux morphismes
$$
H^0(V_{\overline{u}}, \underline{\Omega}) \longrightarrow
H^0((V \times_U V)_{\overline{u}}, \underline{\Omega}) =
H^0(V_{\overline{u}} \times_{\overline{u}} V_{\overline{u}},
\underline{\Omega})
$$
où $\underline{\Omega}$ désigne le faisceau constant de valeur
$\Omega$. Ces morphismes sont bien sûr les images inverses par les projections.
Il est alors clair que les sections provenant de l'image inverse
par la projection sur le premier facteur sont constantes sur les fibres de
la première projection, tandis que celles provenant de l'image inverse
par la projection sur le second facteur sont constantes sur les fibres de
la seconde projection. Les sections dans l'intersection des images
de ces morphismes d'image inverse sont constantes sur tout
$V_{\overline{u}} \times_{\overline{u}} V_{\overline{u}}$, c'est-à-dire
qu'elles proviennent d'éléments de $\Omega$, comme voulu.
\end{proof}

\noindent
Dans la situation du lemme \ref{lemma-describe-pullback-pi-ph},
la composée de $\epsilon_S$ et $\pi_S$, ainsi que l'égalité,
déterminent un morphisme de sites
$$
a_S : (\Sch/S)_{ph} \longrightarrow S_\etale
$$

\begin{lemma}
\label{lemma-push-pull-ph-etale}
Avec les notations ci-dessus.
Soit $f : X \to Y$ un morphisme de $(\Sch/S)_{ph}$.
Il existe alors des diagrammes commutatifs de topos
$$
\xymatrix{
\Sh((\Sch/X)_{ph}) \ar[rr]_{f_{big, ph}} \ar[d]_{\epsilon_X} & &
\Sh((\Sch/Y)_{ph}) \ar[d]^{\epsilon_Y} \\
\Sh((\Sch/X)_\etale) \ar[rr]^{f_{big, \etale}} & &
\Sh((\Sch/Y)_\etale)
}
$$
et
$$
\xymatrix{
\Sh((\Sch/X)_{ph}) \ar[rr]_{f_{big, ph}} \ar[d]_{a_X} & &
\Sh((\Sch/Y)_{ph}) \ar[d]^{a_Y} \\
\Sh(X_\etale) \ar[rr]^{f_{small}} & &
\Sh(Y_\etale)
}
$$
avec $a_X = \pi_X \circ \epsilon_X$ et $a_Y = \pi_X \circ \epsilon_X$.
\end{lemma}

\begin{proof}
La commutativité des diagrammes découle de la discussion de la
section \ref{topologies-section-change-topologies} de Topologies.
\end{proof}

\begin{lemma}
\label{lemma-proper-push-pull-ph-etale}
Dans le lemme \ref{lemma-push-pull-ph-etale}, si $f$ est propre, alors nous avons
$a_Y^{-1} \circ f_{small, *} = f_{big, ph, *} \circ a_X^{-1}$.
\end{lemma}

\begin{proof}
On peut le démontrer en répétant la démonstration de la
partie (1) du lemme \ref{lemma-compare-higher-direct-image-proper};
nous allons plutôt l'en déduire.
Comme $\epsilon_{Y, *}$ est le foncteur identité sur les préfaisceaux sous-jacents,
il reflète les isomorphismes. La description
du lemme \ref{lemma-describe-pullback-pi-ph}
montre que $\epsilon_{Y, *} \circ a_Y^{-1} = \pi_Y^{-1}$,
et de même pour $X$. Pour montrer que le morphisme canonique
$a_Y^{-1}f_{small, *}\mathcal{F} \to f_{big, ph, *}a_X^{-1}\mathcal{F}$
est un isomorphisme, il suffit de montrer que
\begin{align*}
\pi_Y^{-1}f_{small, *}\mathcal{F}
& =
\epsilon_{Y, *}a_Y^{-1}f_{small, *}\mathcal{F} \\
& \to 
\epsilon_{Y, *}f_{big, ph, *}a_X^{-1}\mathcal{F} \\
& =
f_{big, \etale, *} \epsilon_{X, *}a_X^{-1}\mathcal{F} \\
& =
f_{big, \etale, *}\pi_X^{-1}\mathcal{F}
\end{align*}
est un isomorphisme. C'est la partie
(1) du lemme \ref{lemma-compare-higher-direct-image-proper}.
\end{proof}

\begin{lemma}
\label{lemma-compare-ph-etale}
Considérons le morphisme de comparaison
$\epsilon : (\Sch/S)_{ph} \to (\Sch/S)_\etale$.
Soit $\mathcal{P}$ la classe des morphismes propres de schémas.
Pour $X \in (\Sch/S)_\etale$, notons
$\mathcal{A}'_X \subset \textit{Ab}((\Sch/X)_\etale)$
la sous-catégorie pleine formée des faisceaux de la forme
$\pi_X^{-1}\mathcal{F}$, où $\mathcal{F}$ est un
faisceau abélien de torsion sur $X_\etale$.
Alors les propriétés de Cohomologie sur les sites
(\ref{sites-cohomology-item-base-change-P}),
(\ref{sites-cohomology-item-restriction-A}),
(\ref{sites-cohomology-item-A-sheaf}),
(\ref{sites-cohomology-item-A-and-P}) et
(\ref{sites-cohomology-item-refine-tau-by-P})
de la situation \ref{sites-cohomology-situation-compare} de Cohomologie sur les sites
sont satisfaites.
\end{lemma}

\begin{proof}
Montrons d'abord que $\mathcal{A}'_X \subset \textit{Ab}((\Sch/X)_\etale)$
est une sous-catégorie de Serre faible en vérifiant les conditions (1), (2), (3) et (4)
du lemme \ref{homology-lemma-characterize-weak-serre-subcategory} d'Homology.
Les parties (1), (2), (3) sont immédiates, car $\pi_X^{-1}$ est exact et
pleinement fidèle, par exemple par le lemme \ref{lemma-cohomological-descent-etale}. Si
$0 \to \pi_X^{-1}\mathcal{F} \to \mathcal{G} \to \pi_X^{-1}\mathcal{F}' \to 0$
est une suite exacte courte dans $\textit{Ab}((\Sch/X)_\etale)$,
alors $0 \to \mathcal{F} \to \pi_{X, *}\mathcal{G} \to \mathcal{F}' \to 0$
est exacte par le lemme \ref{lemma-cohomological-descent-etale}.
Nous voyons en particulier que $\pi_{X, *}\mathcal{G}$ est un faisceau
abélien de torsion sur $X_\etale$.
Ainsi, $\mathcal{G} = \pi_X^{-1}\pi_{X, *}\mathcal{G}$ appartient à
$\mathcal{A}'_X$, ce qui vérifie la dernière condition.

\medskip\noindent
La propriété (\ref{sites-cohomology-item-base-change-P}) de Cohomologie sur les sites vaut
par l'existence des produits fibrés de schémas
et le fait que le changement de base d'un morphisme propre de
schémas est un morphisme propre de schémas; voir le
lemme \ref{morphisms-lemma-base-change-proper} de Morphismes.

\medskip\noindent
La propriété (\ref{sites-cohomology-item-restriction-A}) de Cohomologie sur les sites
découle du diagramme commutatif (3) du
lemme \ref{topologies-lemma-morphism-big-small-etale} de Topologies.

\medskip\noindent
La propriété (\ref{sites-cohomology-item-A-sheaf}) de Cohomologie sur les sites est le
lemme \ref{lemma-describe-pullback-pi-ph}.

\medskip\noindent
La propriété (\ref{sites-cohomology-item-A-and-P}) de Cohomologie sur les sites vaut par la
partie (2) du lemme \ref{lemma-compare-higher-direct-image-proper}
et le fait que $R^if_{small, *}\mathcal{F}$
est de torsion si $\mathcal{F}$ est un faisceau
abélien de torsion sur $X_\etale$; voir le lemme \ref{lemma-torsion-direct-image}.

\medskip\noindent
La propriété (\ref{sites-cohomology-item-refine-tau-by-P}) de Cohomologie sur les sites
découle du lemme de Compléments sur les morphismes
\ref{more-morphisms-lemma-dominate-fppf-etale-locally}
combiné au fait qu'un morphisme fini est propre
et qu'un morphisme propre surjectif définit un recouvrement ph; voir le
lemme \ref{topologies-lemma-surjective-proper-ph} de Topologies.
\end{proof}

\begin{lemma}
\label{lemma-V-C-all-n-etale-ph}
Avec les notations ci-dessus.
\begin{enumerate}
\item Pour $X \in \Ob((\Sch/S)_{ph})$ et un faisceau abélien de torsion $\mathcal{F}$
sur $X_\etale$, nous avons
$\epsilon_{X, *}a_X^{-1}\mathcal{F} = \pi_X^{-1}\mathcal{F}$
et $R^i\epsilon_{X, *}(a_X^{-1}\mathcal{F}) = 0$ pour $i > 0$.
\item Pour un morphisme propre $f : X \to Y$ dans $(\Sch/S)_{ph}$
et un faisceau abélien de torsion $\mathcal{F}$ sur $X$, nous avons
$a_Y^{-1}(R^if_{small, *}\mathcal{F}) =
R^if_{big, ph, *}(a_X^{-1}\mathcal{F})$
pour tout $i$.
\item Pour un schéma $X$ et $K \in D^+(X_\etale)$ à faisceaux de
cohomologie de torsion, le morphisme
$\pi_X^{-1}K \to R\epsilon_{X, *}(a_X^{-1}K)$ est un isomorphisme.
\item Pour un morphisme propre $f : X \to Y$ de schémas
et $K \in D^+(X_\etale)$ à faisceaux de cohomologie de torsion, nous avons
$a_Y^{-1}(Rf_{small, *}K) = Rf_{big, ph, *}(a_X^{-1}K)$.
\end{enumerate}
\end{lemma}

\begin{proof}
Par le lemme \ref{lemma-compare-ph-etale}, tous les lemmes de la
section \ref{sites-cohomology-section-compare-general} de Cohomologie sur les sites
s'appliquent à notre situation. Pour traduire les résultats,
remarquons que la catégorie $\mathcal{A}_X$ du
lemme \ref{sites-cohomology-lemma-A} de Cohomologie sur les sites
est la sous-catégorie pleine de $\textit{Ab}((\Sch/X)_{ph})$
formée des faisceaux de la forme $a_X^{-1}\mathcal{F}$,
où $\mathcal{F}$ est un faisceau abélien de torsion sur $X_\etale$.

\medskip\noindent
La partie (1) équivaut à $(V_n)$ pour tout $n$, ce qui vaut par le
lemme \ref{sites-cohomology-lemma-V-C-all-n-general} de Cohomologie sur les sites.

\medskip\noindent
La partie (2) découle de l'application de $\epsilon_Y^{-1}$ à la conclusion du
lemme \ref{sites-cohomology-lemma-V-implies-C-general} de Cohomologie sur les sites.

\medskip\noindent
La partie (3) découle de la partie (1) du lemme de Cohomologie sur les sites
\ref{sites-cohomology-lemma-V-C-all-n-general},
car $\pi_X^{-1}K$ appartient à $D^+_{\mathcal{A}'_X}((\Sch/X)_\etale)$
et $a_X^{-1} = \epsilon_X^{-1} \circ \pi_X^{-1}$.

\medskip\noindent
La partie (4) découle de la partie (2) du lemme de Cohomologie sur les sites
\ref{sites-cohomology-lemma-V-C-all-n-general},
pour la même raison.
\end{proof}

\begin{lemma}
\label{lemma-cohomological-descent-etale-ph}
Soit $X$ un schéma. Pour $K \in D^+(X_\etale)$ à faisceaux de cohomologie
de torsion, le morphisme
$$
K \longrightarrow Ra_{X, *}a_X^{-1}K
$$
est un isomorphisme, où $a_X : \Sh((\Sch/X)_{ph}) \to \Sh(X_\etale)$ est comme ci-dessus.
\end{lemma}

\begin{proof}
Nous nous ramenons d'abord au cas où
$K$ est donné par un unique faisceau abélien. Représentons en effet $K$
par un complexe borné inférieurement $\mathcal{F}^\bullet$ de faisceaux
abéliens de torsion. C'est possible par le lemme de Cohomologie sur les sites
\ref{sites-cohomology-lemma-torsion}. Le cas d'un
faisceau montre que
$\mathcal{F}^n = a_{X, *} a_X^{-1} \mathcal{F}^n$
et que les faisceaux $R^qa_{X, *}a_X^{-1}\mathcal{F}^n$
sont nuls pour $q > 0$. Le lemme d'acyclicité de Leray
(lemme \ref{derived-lemma-leray-acyclicity} de Catégories dérivées),
appliqué à $a_X^{-1}\mathcal{F}^\bullet$
et au foncteur $a_{X, *}$, permet de conclure. Supposons désormais $K = \mathcal{F}$,
où $\mathcal{F}$ est un faisceau abélien de torsion.

\medskip\noindent
Par le lemme \ref{lemma-describe-pullback-pi-ph}, nous avons
$a_{X, *}a_X^{-1}\mathcal{F} = \mathcal{F}$. Il suffit donc de montrer que
$R^qa_{X, *}a_X^{-1}\mathcal{F} = 0$ pour $q > 0$.
Pour cela, nous pouvons employer $a_X = \pi_X \circ \epsilon_X$ et
la suite spectrale de Leray
(lemme \ref{sites-cohomology-lemma-relative-Leray} de Cohomologie sur les sites).
Par le lemme \ref{lemma-V-C-all-n-etale-ph},
nous avons $R^i\epsilon_{X, *}(a_X^{-1}\mathcal{F}) = 0$ pour $i > 0$
et $\epsilon_{X, *}a_X^{-1}\mathcal{F} = \pi_X^{-1}\mathcal{F}$.
Par le lemme \ref{lemma-cohomological-descent-etale}, nous avons
$R^j\pi_{X, *}(\pi_X^{-1}\mathcal{F}) = 0$ pour $j > 0$.
Cela achève la démonstration.
\end{proof}

\begin{lemma}
\label{lemma-compare-cohomology-etale-ph}
Pour un schéma $X$ et $a_X : \Sh((\Sch/X)_{ph}) \to \Sh(X_\etale)$
comme ci-dessus :
\begin{enumerate}
\item $H^q(X_\etale, \mathcal{F}) = H^q_{ph}(X, a_X^{-1}\mathcal{F})$
pour un faisceau abélien de torsion $\mathcal{F}$ sur $X_\etale$,
\item $H^q(X_\etale, K) = H^q_{ph}(X, a_X^{-1}K)$
pour $K \in D^+(X_\etale)$ à faisceaux de cohomologie de torsion.
\end{enumerate}
Exemple : si $A$ est un groupe abélien de torsion, alors
$H^q_\etale(X, \underline{A}) = H^q_{ph}(X, \underline{A})$.
\end{lemma}

\begin{proof}
Cela découle du lemme \ref{lemma-cohomological-descent-etale-ph}
par la remarque \ref{sites-cohomology-remark-before-Leray} de Cohomologie sur les sites.
\end{proof}












\section{Comparaison des topologies h et étale}
\label{section-h-etale}

\noindent
Un modèle pour cette section est la section consacrée à la comparaison de la
topologie usuelle et de la topologie qc sur les espaces topologiques
localement compacts, étudiée dans la section de Cohomologie sur les sites
\ref{sites-cohomology-section-cohomology-LC}.
De plus, cette section est presque mot pour mot identique à la
section comparant les topologies ph et étale.
Nous rappelons d'abord certains résultats des
sections de Topologies
\ref{topologies-section-change-topologies} et
\ref{topologies-section-etale}, ainsi que de la
section \ref{flat-section-h} de Compléments sur la platitude.

\medskip\noindent
Soit $S$ un schéma et soit $(\Sch/S)_h$ un site h.
Sur la même catégorie sous-jacente, nous avons une seconde topologie,
à savoir la topologie étale, et donc un second site
$(\Sch/S)_\etale$. Le foncteur identité
$(\Sch/S)_\etale \to (\Sch/S)_h$ est continu
(par le lemme \ref{flat-lemma-zariski-h} de Compléments sur la platitude
et le lemme de Topologies
\ref{topologies-lemma-zariski-etale-smooth-syntomic-fppf})
et définit un morphisme de sites
$$
\epsilon_S : (\Sch/S)_h \longrightarrow (\Sch/S)_\etale
$$
Voir la section \ref{sites-cohomology-section-compare} de Cohomologie sur les sites.
Remarquons que $\epsilon_{S, *}$ est le foncteur identité sur les
préfaisceaux sous-jacents et que $\epsilon_S^{-1}$ associe à un faisceau étale son
faisceau h associé. Soit $S_\etale$ le petit site étale.
Il existe un morphisme de sites
$$
\pi_S : (\Sch/S)_\etale \longrightarrow S_\etale
$$
induit par le foncteur continu
$S_\etale \to (\Sch/S)_\etale$, $U \mapsto U$.
En effet, $S_\etale$ possède des produits fibrés et un objet final, le
foncteur ci-dessus y commute, et la
proposition \ref{sites-proposition-get-morphism} de Sites s'applique.

\begin{lemma}
\label{lemma-describe-pullback-pi-h}
Avec les notations ci-dessus.
Soit $\mathcal{F}$ un faisceau sur $S_\etale$. La règle
$$
(\Sch/S)_h \longrightarrow \textit{Ensembles},\quad
(f : X \to S) \longmapsto \Gamma(X, f_{small}^{-1}\mathcal{F})
$$
est un faisceau h, et a fortiori un faisceau sur $(\Sch/S)_\etale$.
En fait, ce faisceau est égal à
$\pi_S^{-1}\mathcal{F}$ sur $(\Sch/S)_\etale$ et à
$\epsilon_S^{-1}\pi_S^{-1}\mathcal{F}$ sur $(\Sch/S)_h$.
\end{lemma}

\begin{proof}
L'énoncé relatif à la topologie étale est le contenu
du lemme \ref{lemma-describe-pullback}. Pour achever la démonstration, il
suffit de montrer que $\pi_S^{-1}\mathcal{F}$ est un faisceau pour la topologie
h. Or, dans le lemme \ref{lemma-describe-pullback-pi-ph},
nous avons montré que
$\pi_S^{-1}\mathcal{F}$ est même un faisceau pour la topologie ph,
qui est plus fine.
\end{proof}

\noindent
Dans la situation du lemme \ref{lemma-describe-pullback-pi-h},
la composée de $\epsilon_S$ et $\pi_S$, ainsi que l'égalité,
déterminent un morphisme de sites
$$
a_S : (\Sch/S)_h \longrightarrow S_\etale
$$

\begin{lemma}
\label{lemma-push-pull-h-etale}
Avec les notations ci-dessus.
Soit $f : X \to Y$ un morphisme de $(\Sch/S)_h$.
Il existe alors des diagrammes commutatifs de topos
$$
\xymatrix{
\Sh((\Sch/X)_h) \ar[rr]_{f_{big, h}} \ar[d]_{\epsilon_X} & &
\Sh((\Sch/Y)_h) \ar[d]^{\epsilon_Y} \\
\Sh((\Sch/X)_\etale) \ar[rr]^{f_{big, \etale}} & &
\Sh((\Sch/Y)_\etale)
}
$$
et
$$
\xymatrix{
\Sh((\Sch/X)_h) \ar[rr]_{f_{big, h}} \ar[d]_{a_X} & &
\Sh((\Sch/Y)_h) \ar[d]^{a_Y} \\
\Sh(X_\etale) \ar[rr]^{f_{small}} & &
\Sh(Y_\etale)
}
$$
avec $a_X = \pi_X \circ \epsilon_X$ et $a_Y = \pi_Y \circ \epsilon_Y$.
\end{lemma}

\begin{proof}
La commutativité des diagrammes découle de la discussion de la
section \ref{topologies-section-change-topologies} de Topologies.
\end{proof}

\begin{lemma}
\label{lemma-proper-push-pull-h-etale}
Dans le lemme \ref{lemma-push-pull-h-etale}, si $f$ est propre, alors nous avons
$a_Y^{-1} \circ f_{small, *} = f_{big, h, *} \circ a_X^{-1}$.
\end{lemma}

\begin{proof}
On peut le démontrer en répétant la démonstration de la
partie (1) du lemme \ref{lemma-compare-higher-direct-image-proper};
nous allons plutôt l'en déduire.
Comme $\epsilon_{Y, *}$ est le foncteur identité sur les préfaisceaux sous-jacents,
il reflète les isomorphismes. La description
du lemme \ref{lemma-describe-pullback-pi-h}
montre que $\epsilon_{Y, *} \circ a_Y^{-1} = \pi_Y^{-1}$,
et de même pour $X$. Pour montrer que le morphisme canonique
$a_Y^{-1}f_{small, *}\mathcal{F} \to f_{big, h, *}a_X^{-1}\mathcal{F}$
est un isomorphisme, il suffit de montrer que
\begin{align*}
\pi_Y^{-1}f_{small, *}\mathcal{F}
& =
\epsilon_{Y, *}a_Y^{-1}f_{small, *}\mathcal{F} \\
& \to 
\epsilon_{Y, *}f_{big, h, *}a_X^{-1}\mathcal{F} \\
& =
f_{big, \etale, *} \epsilon_{X, *}a_X^{-1}\mathcal{F} \\
& =
f_{big, \etale, *}\pi_X^{-1}\mathcal{F}
\end{align*}
est un isomorphisme. C'est la partie
(1) du lemme \ref{lemma-compare-higher-direct-image-proper}.
\end{proof}

\begin{lemma}
\label{lemma-compare-h-etale}
Considérons le morphisme de comparaison $\epsilon : (\Sch/S)_h \to (\Sch/S)_\etale$.
Soit $\mathcal{P}$ la classe des morphismes propres.
Pour $X \in (\Sch/S)_\etale$, notons
$\mathcal{A}'_X \subset \textit{Ab}((\Sch/X)_\etale)$
la sous-catégorie pleine formée des faisceaux de la forme
$\pi_X^{-1}\mathcal{F}$, où $\mathcal{F}$ est un
faisceau abélien de torsion sur $X_\etale$.
Alors les propriétés de Cohomologie sur les sites
(\ref{sites-cohomology-item-base-change-P}),
(\ref{sites-cohomology-item-restriction-A}),
(\ref{sites-cohomology-item-A-sheaf}),
(\ref{sites-cohomology-item-A-and-P}) et
(\ref{sites-cohomology-item-refine-tau-by-P})
de la situation de Cohomologie sur les sites
\ref{sites-cohomology-situation-compare} sont satisfaites.
\end{lemma}

\begin{proof}
Montrons d'abord que $\mathcal{A}'_X \subset \textit{Ab}((\Sch/X)_\etale)$
est une sous-catégorie de Serre faible en vérifiant les conditions (1), (2), (3) et (4)
du lemme \ref{homology-lemma-characterize-weak-serre-subcategory} d'Homology.
Les parties (1), (2), (3) sont immédiates, car $\pi_X^{-1}$ est exact et
pleinement fidèle, par exemple par le lemme \ref{lemma-cohomological-descent-etale}. Si
$0 \to \pi_X^{-1}\mathcal{F} \to \mathcal{G} \to \pi_X^{-1}\mathcal{F}' \to 0$
est une suite exacte courte dans $\textit{Ab}((\Sch/X)_\etale)$,
alors $0 \to \mathcal{F} \to \pi_{X, *}\mathcal{G} \to \mathcal{F}' \to 0$
est exacte par le lemme \ref{lemma-cohomological-descent-etale}.
Nous voyons en particulier que $\pi_{X, *}\mathcal{G}$ est un faisceau
abélien de torsion sur $X_\etale$.
Ainsi, $\mathcal{G} = \pi_X^{-1}\pi_{X, *}\mathcal{G}$ appartient à
$\mathcal{A}'_X$, ce qui vérifie la dernière condition.

\medskip\noindent
La propriété (\ref{sites-cohomology-item-base-change-P}) de Cohomologie sur les sites vaut
par l'existence des produits fibrés de schémas,
le fait que le changement de base d'un morphisme propre de
schémas est un morphisme propre de schémas; voir le
lemme \ref{morphisms-lemma-base-change-proper} de Morphismes; et
le fait que le changement de base d'un morphisme de présentation finie
est un morphisme de présentation finie; voir le
lemme \ref{morphisms-lemma-base-change-finite-presentation} de Morphismes.

\medskip\noindent
La propriété (\ref{sites-cohomology-item-restriction-A}) de Cohomologie sur les sites
découle du diagramme commutatif (3) du
lemme \ref{topologies-lemma-morphism-big-small-etale} de Topologies.

\medskip\noindent
La propriété (\ref{sites-cohomology-item-A-sheaf}) de Cohomologie sur les sites est le
lemme \ref{lemma-describe-pullback-pi-h}.

\medskip\noindent
La propriété (\ref{sites-cohomology-item-A-and-P}) de Cohomologie sur les sites vaut par la
partie (2) du lemme \ref{lemma-compare-higher-direct-image-proper}
et le fait que $R^if_{small,*}\mathcal{F}$
est de torsion si $\mathcal{F}$ est un faisceau
abélien de torsion sur $X_\etale$; voir le lemme \ref{lemma-torsion-direct-image}.

\medskip\noindent
La propriété (\ref{sites-cohomology-item-refine-tau-by-P}) de Cohomologie sur les sites
découle du lemme de Compléments sur les morphismes
\ref{more-morphisms-lemma-dominate-fppf-etale-locally}
combiné au fait qu'un morphisme fini localement libre surjectif
est surjectif, propre et de présentation finie, et définit donc
un recouvrement h par le lemme de Compléments sur la platitude
\ref{flat-lemma-surjective-proper-finite-presentation-h}.
\end{proof}

\begin{lemma}
\label{lemma-V-C-all-n-etale-h}
Avec les notations ci-dessus.
\begin{enumerate}
\item Pour $X \in \Ob((\Sch/S)_{h})$ et un faisceau abélien de torsion $\mathcal{F}$
sur $X_\etale$, nous avons
$\epsilon_{X, *}a_X^{-1}\mathcal{F} = \pi_X^{-1}\mathcal{F}$
et $R^i\epsilon_{X, *}(a_X^{-1}\mathcal{F}) = 0$ pour $i > 0$.
\item Pour un morphisme propre $f : X \to Y$ dans $(\Sch/S)_h$
et un faisceau abélien de torsion $\mathcal{F}$ sur $X$, nous avons
$a_Y^{-1}(R^if_{small, *}\mathcal{F}) =
R^if_{big, h, *}(a_X^{-1}\mathcal{F})$
pour tout $i$.
\item Pour un schéma $X$ et $K \in D^+(X_\etale)$ à faisceaux de
cohomologie de torsion, le morphisme
$\pi_X^{-1}K \to R\epsilon_{X, *}(a_X^{-1}K)$ est un isomorphisme.
\item Pour un morphisme propre $f : X \to Y$ de schémas
et $K \in D^+(X_\etale)$ à faisceaux de cohomologie de torsion, nous avons
$a_Y^{-1}(Rf_{small, *}K) = Rf_{big, h, *}(a_X^{-1}K)$.
\end{enumerate}
\end{lemma}

\begin{proof}
Par le lemme \ref{lemma-compare-h-etale}, tous les lemmes de la
section \ref{sites-cohomology-section-compare-general} de Cohomologie sur les sites
s'appliquent à notre situation. Pour traduire les résultats,
remarquons que la catégorie $\mathcal{A}_X$ du
lemme \ref{sites-cohomology-lemma-A} de Cohomologie sur les sites
est la sous-catégorie pleine de $\textit{Ab}((\Sch/X)_h)$
formée des faisceaux de la forme $a_X^{-1}\mathcal{F}$,
où $\mathcal{F}$ est un faisceau abélien de torsion sur $X_\etale$.

\medskip\noindent
La partie (1) équivaut à $(V_n)$ pour tout $n$, ce qui vaut par le
lemme \ref{sites-cohomology-lemma-V-C-all-n-general} de Cohomologie sur les sites.

\medskip\noindent
La partie (2) découle de l'application de $\epsilon_Y^{-1}$ à la conclusion du
lemme \ref{sites-cohomology-lemma-V-implies-C-general} de Cohomologie sur les sites.

\medskip\noindent
La partie (3) découle de la partie (1) du lemme de Cohomologie sur les sites
\ref{sites-cohomology-lemma-V-C-all-n-general},
car $\pi_X^{-1}K$ appartient à $D^+_{\mathcal{A}'_X}((\Sch/X)_\etale)$
et $a_X^{-1} = \epsilon_X^{-1} \circ \pi_X^{-1}$.

\medskip\noindent
La partie (4) découle de la partie (2) du lemme de Cohomologie sur les sites
\ref{sites-cohomology-lemma-V-C-all-n-general},
pour la même raison.
\end{proof}

\begin{lemma}
\label{lemma-cohomological-descent-etale-h}
Soit $X$ un schéma. Pour $K \in D^+(X_\etale)$ à faisceaux de cohomologie
de torsion, le morphisme
$$
K \longrightarrow Ra_{X, *}a_X^{-1}K
$$
est un isomorphisme, où $a_X : \Sh((\Sch/X)_h) \to \Sh(X_\etale)$ est comme ci-dessus.
\end{lemma}

\begin{proof}
Nous nous ramenons d'abord au cas où
$K$ est donné par un unique faisceau abélien. Représentons en effet $K$
par un complexe borné inférieurement $\mathcal{F}^\bullet$ de faisceaux
abéliens de torsion. C'est possible par le lemme de Cohomologie sur les sites
\ref{sites-cohomology-lemma-torsion}. Le cas d'un
faisceau montre que
$\mathcal{F}^n = a_{X, *} a_X^{-1} \mathcal{F}^n$
et que les faisceaux $R^qa_{X, *}a_X^{-1}\mathcal{F}^n$
sont nuls pour $q > 0$. Le lemme d'acyclicité de Leray
(lemme \ref{derived-lemma-leray-acyclicity} de Catégories dérivées),
appliqué à $a_X^{-1}\mathcal{F}^\bullet$
et au foncteur $a_{X, *}$, permet de conclure. Supposons désormais $K = \mathcal{F}$,
où $\mathcal{F}$ est un faisceau abélien de torsion.

\medskip\noindent
Par le lemme \ref{lemma-describe-pullback-pi-h}, nous avons
$a_{X, *}a_X^{-1}\mathcal{F} = \mathcal{F}$. Il suffit donc de montrer que
$R^qa_{X, *}a_X^{-1}\mathcal{F} = 0$ pour $q > 0$.
Pour cela, nous pouvons employer $a_X = \pi_X \circ \epsilon_X$ et
la suite spectrale de Leray
(lemme \ref{sites-cohomology-lemma-relative-Leray} de Cohomologie sur les sites).
Par le lemme \ref{lemma-V-C-all-n-etale-h},
nous avons $R^i\epsilon_{X, *}(a_X^{-1}\mathcal{F}) = 0$ pour $i > 0$
et $\epsilon_{X, *}a_X^{-1}\mathcal{F} = \pi_X^{-1}\mathcal{F}$.
Par le lemme \ref{lemma-cohomological-descent-etale}, nous avons
$R^j\pi_{X, *}(\pi_X^{-1}\mathcal{F}) = 0$ pour $j > 0$.
Cela achève la démonstration.
\end{proof}

\begin{lemma}
\label{lemma-compare-cohomology-etale-h}
Pour un schéma $X$ et $a_X : \Sh((\Sch/X)_h) \to \Sh(X_\etale)$
comme ci-dessus :
\begin{enumerate}
\item $H^q(X_\etale, \mathcal{F}) = H^q_h(X, a_X^{-1}\mathcal{F})$
pour un faisceau abélien de torsion $\mathcal{F}$ sur $X_\etale$,
\item $H^q(X_\etale, K) = H^q_h(X, a_X^{-1}K)$
pour $K \in D^+(X_\etale)$ à faisceaux de cohomologie de torsion.
\end{enumerate}
Exemple : si $A$ est un groupe abélien de torsion, alors
$H^q_\etale(X, \underline{A}) = H^q_h(X, \underline{A})$.
\end{lemma}

\begin{proof}
Cela découle du lemme \ref{lemma-cohomological-descent-etale-h}
par la remarque \ref{sites-cohomology-remark-before-Leray} de Cohomologie sur les sites.
\end{proof}











\section{Descente des faisceaux étales}
\label{section-glueing-etale}

\noindent
Nous démontrons que les faisceaux étales se recollent pour les topologies
fppf et h, ainsi que des résultats voisins. Nous avons déjà établi les
résultats connexes suivants :
\begin{enumerate}
\item le lemme \ref{lemma-describe-pullback} affirme qu'un faisceau
sur le petit site étale d'un schéma $S$ détermine un faisceau sur le
grand site étale de $S$ satisfaisant la condition de faisceau pour les
recouvrements fpqc (et a fortiori pour les recouvrements de Zariski, étales,
lisses, syntomiques et fppf);
\item le lemme \ref{lemma-describe-pullback-pi-fppf} reformule
le point précédent pour la topologie fppf;
\item le lemme \ref{lemma-describe-pullback-pi-ph} établit le même résultat pour la
topologie ph;
\item le lemme \ref{lemma-describe-pullback-pi-h} établit le même résultat pour la
topologie h;
\item le lemme \ref{lemma-descent-sheaf-fppf-etale} est une version de la
descente fppf pour les faisceaux étales; et
\item la remarque \ref{remark-cohomological-descent-finite} affirme que
les faisceaux étales descendent par les morphismes finis surjectifs
(nous allons préciser et renforcer ce résultat ci-dessous).
\end{enumerate}
Dans le chapitre sur les espaces simpliciaux, nous démontrerons d'autres
résultats à ce sujet; voir par exemple les sections de Espaces simpliciaux
\ref{spaces-simplicial-section-fppf-hypercovering}
et \ref{spaces-simplicial-section-proper-hypercovering-spaces}.

\medskip\noindent
Pour énoncer commodément nos résultats, introduisons quelques notations.
Soit $\mathcal{U} = \{f_i : X_i \to X\}$
une famille de morphismes de schémas de but fixé.
Une {\it donnée de descente} pour les faisceaux étales relativement à
$\mathcal{U}$ est une famille
$((\mathcal{F}_i)_{i \in I}, (\varphi_{ij})_{i, j \in I})$, où
\begin{enumerate}
\item $\mathcal{F}_i$ appartient à $\Sh(X_{i, \etale})$; et
\item $\varphi_{ij} :
\text{pr}_{0, small}^{-1} \mathcal{F}_i
\longrightarrow
\text{pr}_{1, small}^{-1} \mathcal{F}_j$
est un isomorphisme dans $\Sh((X_i \times_X X_j)_\etale)$.
\end{enumerate}
On exige la {\it condition de cocycle} : les diagrammes
$$
\xymatrix{
\text{pr}_{0, small}^{-1}\mathcal{F}_i
\ar[dr]_{\text{pr}_{02, small}^{-1}\varphi_{ik}}
\ar[rr]^{\text{pr}_{01, small}^{-1}\varphi_{ij}} & &
\text{pr}_{1, small}^{-1}\mathcal{F}_j
\ar[dl]^{\text{pr}_{12, small}^{-1}\varphi_{jk}} \\
& \text{pr}_{2, small}^{-1}\mathcal{F}_k
}
$$
commutent dans $\Sh((X_i \times_X X_j \times_X X_k)_\etale)$.
On dispose d'une notion évidente de {\it morphisme de données de descente},
et l'on obtient une catégorie de données de descente.
Une donnée de descente
$((\mathcal{F}_i)_{i \in I}, (\varphi_{ij})_{i, j \in I})$
est dite {\it effective} s'il existe un objet
$\mathcal{F}$ de $\Sh(X_\etale)$ et des isomorphismes
$\varphi_i : f_{i, small}^{-1} \mathcal{F} \to \mathcal{F}_i$
dans $\Sh(X_{i, \etale})$, compatibles aux $\varphi_{ij}$, c'est-à-dire
tels que
$$
\varphi_{ij} =
\text{pr}_{1, small}^{-1} (\varphi_j) \circ
\text{pr}_{0, small}^{-1} (\varphi_i^{-1})
$$
On peut également formuler cela comme suit. Tout objet $\mathcal{F}$
de $\Sh(X_\etale)$ fournit la {\it donnée de descente canonique}
$(f_{i, small}^{-1}\mathcal{F}, c_{ij})$, où $c_{ij}$
est l'isomorphisme canonique
$$
c_{ij} : \text{pr}_{0, small}^{-1} f_{i, small}^{-1}\mathcal{F}
\longrightarrow
\text{pr}_{1, small}^{-1} f_{j, small}^{-1}\mathcal{F}
$$
La donnée de descente
$((\mathcal{F}_i)_{i \in I}, (\varphi_{ij})_{i, j \in I})$
est effective si et seulement si elle est isomorphe à la donnée de descente
canonique associée à un certain $\mathcal{F}$ de $\Sh(X_\etale)$.

\medskip\noindent
Si la famille est réduite à un seul morphisme $\{X \to Y\}$,
nous voyons une donnée de descente comme un couple $(\mathcal{F}, \varphi)$,
où $\mathcal{F}$ est un objet de $\Sh(X_\etale)$ et
$\varphi$ un isomorphisme
$$
\text{pr}_{0, small}^{-1} \mathcal{F}
\longrightarrow
\text{pr}_{1, small}^{-1} \mathcal{F}
$$
dans $\Sh((X \times_Y X)_\etale)$ satisfaisant la condition de cocycle :
$$
\xymatrix{
\text{pr}_{0, small}^{-1}\mathcal{F}
\ar[dr]_{\text{pr}_{02, small}^{-1}\varphi}
\ar[rr]^{\text{pr}_{01, small}^{-1}\varphi} & &
\text{pr}_{1, small}^{-1}\mathcal{F}
\ar[dl]^{\text{pr}_{12, small}^{-1}\varphi} \\
& \text{pr}_{2, small}^{-1}\mathcal{F}
}
$$
commute dans $\Sh((X \times_Y X \times_Y X)_\etale)$.
Les notions de morphisme de données de descente et d'effectivité
sont exactement celles données ci-dessus.

\medskip\noindent
Nous démontrons d'abord la descente effective pour les morphismes entiers surjectifs.

\begin{lemma}
\label{lemma-glue-etale-sheaf-section}
Soit $f : X \to Y$ un morphisme de schémas qui possède une section. Alors le
foncteur
$$
\Sh(Y_\etale)
\longrightarrow
\text{données de descente pour les faisceaux étales relativement à }\{X \to Y\}
$$
qui associe à $\mathcal{G}$ dans $\Sh(Y_\etale)$ sa donnée de descente canonique
est une équivalence de catégories.
\end{lemma}

\begin{proof}
C'est formel et ne dépend que de la fonctorialité des foncteurs
d'image inverse. Nous omettons les détails. Indication : si $s : Y \to X$ est une section,
un quasi-inverse est le foncteur qui associe à $(\mathcal{F}, \varphi)$
l'objet $s_{small}^{-1}\mathcal{F}$.
\end{proof}

\begin{lemma}
\label{lemma-glue-etale-sheaf-integral-surjective}
Soit $f : X \to Y$ un morphisme entier surjectif de schémas.
Le foncteur
$$
\Sh(Y_\etale)
\longrightarrow
\text{données de descente pour les faisceaux étales relativement à }\{X \to Y\}
$$
est une équivalence de catégories.
\end{lemma}

\begin{proof}
Dans cette démonstration, nous omettons l'indice ${}_{small}$ de nos foncteurs
d'image inverse et d'image directe.
Notons $X_1 = X \times_Y X$ et notons $f_1 : X_1 \to Y$ le
morphisme $f \circ \text{pr}_0 = f \circ \text{pr}_1$.
Soit $(\mathcal{F}, \varphi)$ une donnée de descente pour $\{X \to Y\}$.
Posons $\mathcal{F}_1 = \text{pr}_0^{-1}\mathcal{F}$.
Nous pouvons voir $\varphi$ comme définissant un isomorphisme
$\mathcal{F}_1 \to \text{pr}_1^{-1}\mathcal{F}$.
Nous affirmons que la règle qui associe à une donnée de descente
$(\mathcal{F}, \varphi)$
le faisceau
$$
\mathcal{G} =
\text{Égalisateur}\left(
\xymatrix{
f_*\mathcal{F}
\ar@<1ex>[r] \ar@<-1ex>[r] &
f_{1, *}\mathcal{F}_1
}
\right)
$$
est un quasi-inverse du foncteur de l'énoncé du lemme.
Le premier des deux morphismes provient du morphisme
$$
f_*\mathcal{F} \to
f_*\text{pr}_{0, *}\text{pr}_0^{-1}\mathcal{F} =
f_{1, *}\mathcal{F}_1
$$
et le second provient du morphisme
$$
f_*\mathcal{F} \to
f_* \text{pr}_{1, *}\text{pr}_1^{-1}\mathcal{F}
\xleftarrow{\varphi}
f_* \text{pr}_{0, *} \text{pr}_0^{-1}\mathcal{F} =
f_{1, *}\mathcal{F}_1
$$
où le morphisme dirigé vers la gauche est inversible.
Pour justifier cette construction, il faut montrer
que le morphisme canonique $f^{-1}\mathcal{G} \to \mathcal{F}$
est un isomorphisme; nous omettons les détails. Pour ce faire, il
suffit de le vérifier après image inverse par une collection de morphismes
$\Spec(k) \to Y$, où $k$ est un corps algébriquement clos.
En effet, les changements de base correspondants $X_k \to X$ sont conjointement
surjectifs, et l'on peut vérifier qu'un morphisme de faisceaux sur $X_\etale$
est un isomorphisme en examinant ses fibres aux points géométriques; voir le
théorème \ref{theorem-exactness-stalks}.
Par le lemme \ref{lemma-integral-pushforward-commutes-with-base-change},
la construction de $\mathcal{G}$ à partir de la donnée de descente
$(\mathcal{F}, \varphi)$ commute à tout changement de base.
Nous pouvons donc supposer que $Y$ est le spectre d'un corps
algébriquement clos (remarquons que le changement de base préserve les propriétés
du morphisme $f$; voir les lemmes de Morphismes
\ref{morphisms-lemma-base-change-surjective}
et \ref{morphisms-lemma-base-change-finite}).
Dans ce cas, le morphisme $X \to Y$ possède une section, et nous
savons donc, par le lemme \ref{lemma-glue-etale-sheaf-section},
que le foncteur est une équivalence.
Or le lecteur peut vérifier que le foncteur est une équivalence
si et seulement si la construction ci-dessus est un quasi-inverse;
nous omettons les détails. Cela achève la démonstration.
\end{proof}

\begin{lemma}
\label{lemma-glue-etale-sheaf-proper-surjective}
Soit $f : X \to Y$ un morphisme propre surjectif de schémas.
Le foncteur
$$
\Sh(Y_\etale)
\longrightarrow
\text{données de descente pour les faisceaux étales relativement à }\{X \to Y\}
$$
est une équivalence de catégories.
\end{lemma}

\begin{proof}
La démonstration du
lemme \ref{lemma-glue-etale-sheaf-integral-surjective}
s'applique mot pour mot, si ce n'est que l'emploi du
lemme \ref{lemma-integral-pushforward-commutes-with-base-change}
doit être remplacé par celui du
lemme \ref{lemma-proper-base-change-f-star}.
\end{proof}

\begin{lemma}
\label{lemma-glue-etale-sheaf-check-after-base-change}
Soit $f : X \to Y$ un morphisme de schémas. Soit $Z \to Y$ un morphisme
entier surjectif de schémas ou un morphisme propre surjectif de schémas.
Si les foncteurs
$$
\Sh(Z_\etale)
\longrightarrow
\text{données de descente pour les faisceaux étales relativement à }\{X \times_Y Z \to Z\}
$$
et
$$
\Sh((Z \times_Y Z)_\etale)
\longrightarrow
\text{données de descente pour les faisceaux étales relativement à }
\{X \times_Y (Z \times_Y Z) \to Z \times_Y Z\}
$$
sont des équivalences de catégories, alors
$$
\Sh(Y_\etale)
\longrightarrow
\text{données de descente pour les faisceaux étales relativement à }\{X \to Y\}
$$
est une équivalence.
\end{lemma}

\begin{proof}
C'est une conséquence formelle des définitions et des
lemmes \ref{lemma-glue-etale-sheaf-integral-surjective} et
\ref{lemma-glue-etale-sheaf-proper-surjective}. Nous omettons les détails.
\end{proof}

\begin{lemma}
\label{lemma-glue-etale-sheaf-fppf-cover}
Soit $f : X \to Y$ un morphisme de schémas
surjectif, plat et localement de présentation finie.
Le foncteur
$$
\Sh(Y_\etale)
\longrightarrow
\text{données de descente pour les faisceaux étales relativement à }\{X \to Y\}
$$
est une équivalence de catégories.
\end{lemma}

\begin{proof}
Exactement comme dans la démonstration du
lemme \ref{lemma-glue-etale-sheaf-integral-surjective},
nous affirmons qu'un quasi-inverse est le foncteur qui associe
à $(\mathcal{F}, \varphi)$ le faisceau
$$
\mathcal{G} =
\text{Égalisateur}\left(
\xymatrix{
f_*\mathcal{F}
\ar@<1ex>[r] \ar@<-1ex>[r] &
f_{1, *}\mathcal{F}_1
}
\right)
$$
et, pour le démontrer, il suffit de montrer que
$f^{-1}\mathcal{G} \to \mathcal{F}$ est un isomorphisme.
Cette propriété se vérifie localement; nous pouvons donc supposer, et supposons, $Y$
affine. Il existe alors un morphisme fini surjectif
$Z \to Y$ et un recouvrement ouvert
$Z = \bigcup W_i$ tels que $W_i \to Y$ se factorise par $X$.
Voir le
lemme \ref{more-morphisms-lemma-dominate-fppf-finite} de Compléments sur les morphismes.
En appliquant le lemme
\ref{lemma-glue-etale-sheaf-check-after-base-change}
nous voyons qu'il suffit de démontrer
le lemme après avoir remplacé $Y$ par $Z$ et $Z \times_Y Z$, et $f$
par son changement de base. Nous pouvons donc supposer que $f$ possède localement des sections pour Zariski.
Puisque le problème est local sur $Y$, nous nous ramenons alors
au cas où nous disposons d'une section, lequel est le
lemme \ref{lemma-glue-etale-sheaf-section}.
\end{proof}

\begin{lemma}
\label{lemma-glue-etale-sheaf-fppf}
Soit $\{f_i : X_i \to X\}$ un recouvrement fppf de schémas.
Le foncteur
$$
\Sh(X_\etale)
\longrightarrow
\text{données de descente pour les faisceaux étales relativement à }\{f_i : X_i \to X\}
$$
est une équivalence de catégories.
\end{lemma}

\begin{proof}
Appliquons le lemme \ref{lemma-glue-etale-sheaf-fppf-cover}
au morphisme $f : \coprod X_i \to X$.
Un argument formel montre alors que les données de descente pour $f$
sont les mêmes que les données de descente pour le recouvrement; comparer
avec le lemme \ref{descent-lemma-family-is-one} de Descente.
Nous omettons les détails.
\end{proof}

\begin{lemma}
\label{lemma-glue-etale-sheaf-modification}
Soit $f : X' \to X$ un morphisme propre de schémas. Soit $i : Z \to X$
une immersion fermée. Posons $E = Z \times_X X'$. Voici le diagramme :
$$
\xymatrix{
E \ar[d]_g \ar[r]_j & X' \ar[d]^f \\
Z \ar[r]^i & X
}
$$
Si $f$ est un isomorphisme au-dessus de $X \setminus Z$, alors le foncteur
$$
\Sh(X_\etale)
\longrightarrow
\Sh(X'_\etale) \times_{\Sh(E_\etale)} \Sh(Z_\etale)
$$
est une équivalence de catégories.
\end{lemma}

\begin{proof}
Nous travaillons avec le $2$-produit fibré de catégories construit dans
l'exemple \ref{categories-example-2-fibre-product-categories} de Catégories.
Le foncteur associe à $\mathcal{F}$ le triplet
$(f^{-1}\mathcal{F}, i^{-1}\mathcal{F}, c)$, où
$c : j^{-1}f^{-1}\mathcal{F} \to g^{-1}i^{-1}\mathcal{F}$
est l'isomorphisme canonique. Construisons un foncteur quasi-inverse. Soit
$(\mathcal{F}', \mathcal{G}, \alpha)$ un objet
du membre de droite de la flèche.
Nous obtenons un isomorphisme
$$
i^{-1}f_*\mathcal{F}' = g_*j^{-1}\mathcal{F}'
\xrightarrow{g_*\alpha}
g_*g^{-1}\mathcal{G}
$$
La première égalité est le lemme \ref{lemma-proper-base-change-f-star}.
Nous en déduisons des morphismes
$i_*\mathcal{G} \to i_*g_*g^{-1}\mathcal{G}$
et $f_*\mathcal{F}' \to i_*g_*g^{-1}\mathcal{G}$. Posons
$$
\mathcal{F} = f_*\mathcal{F}' \times_{i_*g_*g^{-1}\mathcal{G}} i_*\mathcal{G}
$$
Nous affirmons que $\mathcal{F}$ est un objet du membre de gauche
de la flèche, dont l'image dans le membre de droite est isomorphe au
triplet initial. Calculons la fibre de $\mathcal{F}$
en un point géométrique $\overline{x}$ de $X$.

\medskip\noindent
Si $\overline{x}$ n'appartient pas
à $Z$, alors, d'une part, $\overline{x}$ provient d'un unique
point géométrique $\overline{x}'$ de $X'$ et
$\mathcal{F}'_{\overline{x}'} = (f_*\mathcal{F}')_{\overline{x}}$
et, d'autre part, $(i_*\mathcal{G})_{\overline{x}}$
et $(i_*g_*g^{-1}\mathcal{G})_{\overline{x}}$ sont des singletons.
Nous voyons donc que $\mathcal{F}_{\overline{x}}$ est égal à
$\mathcal{F}'_{\overline{x}'}$.

\medskip\noindent
Si $\overline{x}$ appartient à $Z$, c'est-à-dire si $\overline{x}$ est l'image d'un
point géométrique $\overline{z}$ de $Z$, alors nous obtenons
$(i_*\mathcal{G})_{\overline{x}} = \mathcal{G}_{\overline{z}}$
et
$$
(i_*g_*g^{-1}\mathcal{G})_{\overline{x}} =
(g_*g^{-1}\mathcal{G})_{\overline{z}} =
\Gamma(E_{\overline{z}}, g^{-1}\mathcal{G}|_{E_{\overline{z}}})
$$
(par le changement de base propre pour l'image directe employé ci-dessus),
et de même
$$
(f_*\mathcal{F}')_{\overline{x}} =
\Gamma(X'_{\overline{x}}, \mathcal{F}'|_{X'_{\overline{x}}})
$$
Puisque nous avons l'identification
$E_{\overline{z}} = X'_{\overline{x}}$ et que $\alpha$
définit un isomorphisme entre les faisceaux
$\mathcal{F}'|_{X'_{\overline{x}}}$ et
$g^{-1}\mathcal{G}|_{E_{\overline{z}}}$
nous en déduisons
$$
\mathcal{F}_{\overline{x}} = \mathcal{G}_{\overline{z}}
$$
dans ce cas.

\medskip\noindent
Pour achever la démonstration, observons qu'il existe des morphismes canoniques
$i^{-1}\mathcal{F} \to \mathcal{G}$ et $f^{-1}\mathcal{F} \to \mathcal{F}'$
compatibles avec $\alpha$ et induisent sur les fibres les isomorphismes
obtenus ci-dessus. Nous omettons la construction détaillée de ces morphismes.
\end{proof}

\begin{lemma}
\label{lemma-h-descent-etale-sheaves}
Soit $S$ un schéma. Alors la catégorie fibrée
$$
p : \mathcal{S} \longrightarrow (\Sch/S)_h
$$
dont la catégorie fibre au-dessus de $U$ est la catégorie $\Sh(U_\etale)$
des faisceaux sur le petit site étale de $U$ est un champ.
\end{lemma}

\begin{proof}
Pour démontrer le lemme, nous vérifions les conditions (1), (2) et (3) du
lemme \ref{flat-lemma-refine-check-h-stack} de Compléments sur la platitude.

\medskip\noindent
La condition (1) vaut parce que les faisceaux se recollent (et que les
recouvrements de Zariski sont des recouvrements étales). Voir le
lemme \ref{sites-lemma-glue-sheaves} de Sites.

\medskip\noindent
Pour vérifier la condition (2), supposons que $f : X \to Y$ soit un morphisme
surjectif, plat, propre et de présentation finie sur $S$, avec $Y$ affine.
Nous avons alors la descente pour $\{X \to Y\}$, soit par le
lemme \ref{lemma-glue-etale-sheaf-fppf-cover}, soit par le
lemme \ref{lemma-glue-etale-sheaf-proper-surjective}.

\medskip\noindent
La condition (3) découle immédiatement du résultat plus général qu'est le
lemme \ref{lemma-glue-etale-sheaf-modification}.
\end{proof}



















\section{Carrés d'éclatement et cohomologie étale}
\label{section-blow-up-square}

\noindent
Les carrés d'éclatement sont introduits dans la
section \ref{flat-section-blow-up-ph} de Compléments sur la platitude.
Le théorème de changement de base propre permet d'obtenir
un résultat de type Mayer--Vietoris pour les carrés d'éclatement.

\begin{lemma}
\label{lemma-blow-up-square-cohomological-descent}
Soit $X$ un schéma et soit $Z \subset X$ un sous-schéma fermé
défini par un idéal quasi-cohérent de type fini. Considérons le
carré d'éclatement correspondant
$$
\xymatrix{
E \ar[d]_\pi \ar[r]_j & X' \ar[d]^b \\
Z \ar[r]^i & X
}
$$
Pour $K \in D^+(X_\etale)$ à faisceaux de cohomologie de torsion,
nous avons un triangle distingué
$$
K \to Ri_*(K|_Z) \oplus Rb_*(K|_{X'}) \to Rc_*(K|_E) \to K[1]
$$
dans $D(X_\etale)$, où $c = i \circ \pi = b \circ j$.
\end{lemma}

\begin{proof}
La notation $K|_{X'}$ désigne $b_{small}^{-1}K$.
Choisissons un complexe borné inférieurement $\mathcal{F}^\bullet$
de faisceaux abéliens représentant $K$. Remarquons que
$i_*(\mathcal{F}^\bullet|_Z)$ représente $Ri_*(K|_Z)$,
car $i_*$ est exact
(proposition \ref{proposition-finite-higher-direct-image-zero}).
Choisissons un quasi-isomorphisme
$b_{small}^{-1}\mathcal{F}^\bullet \to \mathcal{I}^\bullet$
où $\mathcal{I}^\bullet$ est un complexe borné inférieurement de faisceaux
abéliens injectifs sur $X'_\etale$. Ce morphisme est adjoint à un morphisme
$\mathcal{F}^\bullet \to b_*(\mathcal{I}^\bullet)$, et
$b_*(\mathcal{I}^\bullet)$ représente $Rb_*(K|_{X'})$.
Nous avons $\pi_*(\mathcal{I}^\bullet|_E) = (b_*\mathcal{I}^\bullet)|_Z$
par le lemme \ref{lemma-proper-base-change-f-star}, et, par le
lemme \ref{lemma-proper-base-change}, ce complexe représente
$R\pi_*(K|_E)$. Par conséquent, le morphisme
$$
Ri_*(K|_Z) \oplus Rb_*(K|_{X'}) \to Rc_*(K|_E)
$$
est représenté par le morphisme surjectif de complexes bornés inférieurement
$$
i_*(\mathcal{F}^\bullet|_Z) \oplus
b_*(\mathcal{I}^\bullet)
\to
i_*\left(b_*(\mathcal{I}^\bullet)|_Z\right)
$$
Pour obtenir le triangle distingué cherché, il suffit de montrer que
le morphisme canonique
$\mathcal{F}^\bullet \to i_*(\mathcal{F}^\bullet|_Z) \oplus
b_*(\mathcal{I}^\bullet)$
induit un quasi-isomorphisme sur le noyau du morphisme de complexes
affiché ci-dessus (en effet, une suite exacte courte
de complexes détermine un triangle distingué dans la catégorie
dérivée; voir la
section \ref{derived-section-canonical-delta-functor} de Catégories dérivées).
On peut le vérifier sur les fibres en un point géométrique $\overline{x}$ de $X$.
Si $\overline{x}$ n'appartient pas à $Z$, alors $X' \to X$ est un isomorphisme
au-dessus d'un voisinage ouvert de $\overline{x}$. Si $\overline{x}'$
désigne le point géométrique correspondant de $X'$ dans ce cas, il
faut donc montrer que
$$
\mathcal{F}^\bullet_{\overline{x}} \to \mathcal{I}^\bullet_{\overline{x}'}
$$
est un quasi-isomorphisme. Cela résulte du choix de $\mathcal{I}^\bullet$.
Si $\overline{x}$ appartient à $Z$, alors
$b_*(\mathcal{I}^\bullet)_{\overline{x}} \to 
i_*\left(b_*(\mathcal{I}^\bullet)|_Z\right)_{\overline{x}}$
est un isomorphisme de complexes de groupes abéliens. Le
noyau est donc égal à
$i_*(\mathcal{F}^\bullet|_Z)_{\overline{x}} =
\mathcal{F}^\bullet_{\overline{x}}$, comme voulu.
\end{proof}

\begin{lemma}
\label{lemma-blow-up-square-etale-cohomology}
Soit $X$ un schéma et soit $K \in D^+(X_\etale)$ à
faisceaux de cohomologie de torsion. Soit $Z \subset X$ un sous-schéma fermé
défini par un idéal quasi-cohérent de type fini. Considérons le
carré d'éclatement correspondant
$$
\xymatrix{
E \ar[d] \ar[r] & X' \ar[d]^b \\
Z \ar[r] & X
}
$$
Il existe alors une suite exacte longue canonique
$$
H^p_\etale(X, K) \to
H^p_\etale(X', K|_{X'}) \oplus
H^p_\etale(Z, K|_Z) \to
H^p_\etale(E, K|_E) \to
H^{p + 1}_\etale(X, K)
$$
\end{lemma}

\begin{proof}[Première démonstration]
Cela découle immédiatement du
lemme \ref{lemma-blow-up-square-cohomological-descent}
et de l'égalité
$$
R\Gamma(X, Rb_*(K|_{X'})) = R\Gamma(X', K|_{X'})
$$
(voir la section \ref{sites-cohomology-section-leray} de Cohomologie sur les sites),
ainsi que des égalités analogues pour les autres termes.
\end{proof}

\begin{proof}[Seconde démonstration]
Par le lemme \ref{lemma-compare-cohomology-etale-ph},
ces groupes de cohomologie sont ceux de
$X, X', E, Z$ à valeurs dans un certain complexe de faisceaux abéliens
sur le site $(\Sch/X)_{ph}$. (Il s'agit de l'objet
$a_X^{-1}K$ de la catégorie dérivée; voir le
lemme \ref{lemma-describe-pullback-pi-ph} ci-dessus
et rappelons que $K|_{X'} = b_{small}^{-1}K$.)
Par le lemme \ref{flat-lemma-blow-up-square-ph} de Compléments sur la platitude,
le faisceau ph associé au diagramme de préfaisceaux
représentables est cocartésien. Le lemme découle donc du résultat très général qu'est le
lemme \ref{sites-cohomology-lemma-square-triangle} de Cohomologie sur les sites,
appliqué au site $(\Sch/X)_{ph}$ et au diagramme commutatif
du lemme.
\end{proof}

\begin{lemma}
\label{lemma-blow-up-square-equivalence}
Soit $X$ un schéma et soit $Z \subset X$ un sous-schéma fermé
défini par un idéal quasi-cohérent de type fini. Considérons le
carré d'éclatement correspondant
$$
\xymatrix{
E \ar[d]_\pi \ar[r]_j & X' \ar[d]^b \\
Z \ar[r]^i & X
}
$$
Supposons donnés
\begin{enumerate}
\item un objet $K'$ de $D^+(X'_\etale)$ à faisceaux de cohomologie de torsion;
\item un objet $L$ de $D^+(Z_\etale)$ à faisceaux de cohomologie de torsion; et
\item un isomorphisme $\gamma : K'|_E \to L|_E$.
\end{enumerate}
Il existe alors un objet $K$ de $D^+(X_\etale)$
et des isomorphismes $f : K|_{X'} \to K'$, $g : K|_Z \to L$ tels
que $\gamma = g|_E \circ f^{-1}|_E$.
De plus, supposons donnés
\begin{enumerate}
\item un objet $M$ de $D^+(X_\etale)$ à faisceaux de cohomologie de torsion;
\item un morphisme $\alpha : K' \to M|_{X'}$ de $D(X'_\etale)$;
\item un morphisme $\beta : L \to M|_Z$ de $D(Z_\etale)$;
\end{enumerate}
tels que
$$
\alpha|_E  = \beta|_E \circ \gamma.
$$
Alors il existe un morphisme $K \to M$ dans $D(X_\etale)$
dont la restriction à $X'$ est $\alpha \circ f$
et dont la restriction à $Z$ est $\beta \circ g$.
\end{lemma}

\begin{proof}
Si $K$ existe, le lemme \ref{lemma-blow-up-square-cohomological-descent}
donne un triangle distingué dans lequel il s'insère. Choisissons donc simplement
un triangle distingué
$$
K \to Ri_*(L) \oplus Rb_*(K') \to Rc_*(L|_E) \to K[1]
$$
où $c = i \circ \pi = b \circ j$. Ici, le morphisme $Ri_*(L) \to Rc_*(L|_E)$
est l'image par $Ri_*$ du morphisme d'adjonction $L \to R\pi_*(L|_E)$.
Le morphisme $Rb_*(K') \to Rc_*(L|_E)$ est la composée du morphisme canonique
$Rb_*(K') \to Rc_*(K'|_E)$ et de $Rc_*(\gamma)$.
Les morphismes $g$ et $f$ de l'énoncé du lemme sont les adjoints
de ces morphismes. Si l'on restreint ce triangle distingué à $Z$,
alors le morphisme $Rb_*(K') \to Rc_*(L|_E)$ devient un isomorphisme
par le théorème de changement de base (lemme \ref{lemma-proper-base-change}), et donc
le morphisme $g : K|_Z \to L$ est un isomorphisme.
L'examen du triangle distingué montre que $f : K|_{X'} \to K'$
est un isomorphisme au-dessus de $X' \setminus E = X \setminus Z$.
De plus, nous avons $\gamma \circ f|_E = g|_E$ par construction.
Comme $\gamma$ et $g$ sont des isomorphismes, nous en déduisons
que $f$ induit également des isomorphismes sur les fibres aux points géométriques de $E$.
Ainsi, $f$ est un isomorphisme.

\medskip\noindent
Pour la dernière assertion, on peut remplacer $K'$ par $K|_{X'}$,
$L$ par $K|_Z$ et $\gamma$ par l'identification canonique.
Observons que $\alpha$ et $\beta$ induisent un carré commutatif
$$
\xymatrix{
K \ar[r] \ar@{..>}[d] &
Ri_*(K|_Z) \oplus Rb_*(K|_{X'}) \ar[r] \ar[d]_{\beta \oplus \alpha} &
Rc_*(K|_E) \ar[r] \ar[d]_{\alpha|_E} &
K[1] \ar@{..>}[d] \\
M \ar[r] &
Ri_*(M|_Z) \oplus Rb_*(M|_{X'}) \ar[r] &
Rc_*(M|_E) \ar[r] &
M[1]
}
$$
Les axiomes d'une catégorie dérivée fournissent donc une flèche
pointillée qui définit un morphisme de triangles distingués.
\end{proof}














\section{Carrés presque d'éclatement et topologie h}
\label{section-blow-up-h}

\noindent
Dans cette section, nous poursuivons la discussion de la
section \ref{flat-section-blow-up-h} de Compléments sur la platitude.
Pour la commodité du lecteur, rappelons qu'un
carré presque d'éclatement est un diagramme commutatif
\begin{equation}
\label{equation-almost-blow-up-square}
\vcenter{
\xymatrix{
E \ar[d] \ar[r] & X' \ar[d]^b \\
Z \ar[r] & X
}
}
\end{equation}
de schémas satisfaisant les conditions suivantes :
\begin{enumerate}
\item $Z \to X$ est une immersion fermée de présentation finie;
\item $E = b^{-1}(Z)$ est un sous-schéma fermé localement principal de $X'$;
\item $b$ est propre et de présentation finie;
\item le sous-schéma fermé $X'' \subset X'$ défini par l'idéal quasi-cohérent
des sections de $\mathcal{O}_{X'}$ à support dans $E$
(lemme \ref{properties-lemma-sections-supported-on-closed-subset} de Propriétés)
est l'éclatement de $X$ le long de $Z$.
\end{enumerate}
Il s'ensuit que le morphisme $b$ induit un isomorphisme
$X' \setminus E \to X \setminus Z$.

\medskip\noindent
Nous allons donner un critère caractérisant le « caractère de faisceau h » des
objets de la catégorie dérivée du grand site fppf
$(\Sch/S)_{fppf}$ d'un schéma $S$. Sur la même catégorie sous-jacente,
nous disposons d'une seconde topologie, à savoir la topologie h
(section \ref{flat-section-h} de Compléments sur la platitude).
Rappelons que les recouvrements fppf sont des recouvrements h
(lemme \ref{flat-lemma-zariski-h} de Compléments sur la platitude). Nous pouvons donc
considérer le morphisme
$$
\epsilon : (\Sch/S)_h \longrightarrow (\Sch/S)_{fppf}
$$
Voir la section \ref{sites-cohomology-section-compare} de Cohomologie sur les sites.
En particulier, nous avons un foncteur pleinement fidèle
$$
R\epsilon_* : D((\Sch/S)_h) \longrightarrow D((\Sch/S)_{fppf})
$$
et nous pouvons demander : quelle est l'image essentielle de ce foncteur?

\begin{lemma}
\label{lemma-blow-up-square-h}
Avec les notations ci-dessus, si $K$ appartient à l'image essentielle
de $R\epsilon_*$, alors les morphismes $c^K_{X, Z, X', E}$ du
lemme \ref{sites-cohomology-lemma-c-square} de Cohomologie sur les sites
sont des quasi-isomorphismes.
\end{lemma}

\begin{proof}
Notons ${}^\#$ le passage au faisceau associé pour la topologie h.
Nous avons vu dans le lemme \ref{flat-lemma-blow-up-square-h} de Compléments sur la platitude
que $h_X^\# = h_Z^\# \amalg_{h_E^\#} h_{X'}^\#$. D'autre part,
le morphisme $h_E^\# \to h_{X'}^\#$ est injectif, car $E \to X'$ est un
monomorphisme. Ce lemme est donc un cas particulier du
lemme \ref{sites-cohomology-lemma-descent-squares-helper} de Cohomologie sur les sites
(qui est lui-même une conséquence formelle du
lemme \ref{sites-cohomology-lemma-square-triangle} de Cohomologie sur les sites).
\end{proof}

\begin{proposition}
\label{proposition-check-h}
Soit $K$ un objet de $D^+((\Sch/S)_{fppf})$.
Alors $K$ appartient à l'image essentielle de
$R\epsilon_* : D((\Sch/S)_h) \to D((\Sch/S)_{fppf})$
si et seulement si $c^K_{X, X', Z, E}$ est un quasi-isomorphisme
pour tout carré presque d'éclatement (\ref{equation-almost-blow-up-square})
dans $(\Sch/S)_h$ tel que $X$ soit affine.
\end{proposition}

\begin{proof}
Cela résulte de l'application du
lemme \ref{sites-cohomology-lemma-descent-squares} de Cohomologie sur les sites,
dont les hypothèses sont satisfaites par le
lemme \ref{lemma-blow-up-square-h} et la
proposition \ref{flat-proposition-check-h} de Compléments sur la platitude.
\end{proof}

\begin{lemma}
\label{lemma-refine-check-h}
Soit $K$ un objet de $D^+((\Sch/S)_{fppf})$. Alors $K$ appartient à
l'image essentielle de $R\epsilon_* : D((\Sch/S)_h) \to D((\Sch/S)_{fppf})$
si et seulement si $c^K_{X, X', Z, E}$ est un quasi-isomorphisme
pour tout carré presque d'éclatement comme dans les
exemples \ref{flat-example-one-generator} et
\ref{flat-example-two-generators}.
\end{lemma}

\begin{proof}
Cela résulte de l'application du
lemme \ref{sites-cohomology-lemma-descent-squares} de Cohomologie sur les sites,
dont les hypothèses sont satisfaites par le
lemme \ref{lemma-blow-up-square-h} et le
lemme \ref{flat-lemma-refine-check-h} de Compléments sur la platitude.
\end{proof}






\section{Cohomologie du faisceau structural pour la topologie h}
\label{section-cohomology-O-h}

\noindent
Soit $p$ un nombre premier. Soit $(\mathcal{C}, \mathcal{O})$ un site annelé
tel que $p\mathcal{O} = 0$. Nous notons alors $\colim_F \mathcal{O}$
la limite inductive, dans la catégorie des faisceaux d'anneaux, du système
$$
\mathcal{O} \xrightarrow{F} \mathcal{O} \xrightarrow{F}
\mathcal{O} \xrightarrow{F} \ldots
$$
où $F : \mathcal{O} \to \mathcal{O}$, $f \mapsto f^p$,
est l'endomorphisme de Frobenius.

\begin{lemma}
\label{lemma-h-sheaf-colim-F}
Soit $p$ un nombre premier. Soit $S$ un schéma sur $\mathbf{F}_p$.
Considérons le faisceau $\mathcal{O}^{perf} = \colim_F \mathcal{O}$
sur $(\Sch/S)_{fppf}$. Alors $\mathcal{O}^{perf}$ appartient à l'image
essentielle de $R\epsilon_* : D((\Sch/S)_h) \to D((\Sch/S)_{fppf})$.
\end{lemma}

\begin{proof}
Nous employons le critère du lemme \ref{lemma-refine-check-h}.
Avant d'en vérifier les conditions, remarquons que, pour un objet
quasi-compact et quasi-séparé $X$ de
$(\Sch/S)_{fppf}$, nous avons
$$
H^i_{fppf}(X, \mathcal{O}^{perf}) = \colim_F H^i_{fppf}(X, \mathcal{O})
$$
Voir le lemme de Cohomologie sur les sites
\ref{sites-cohomology-lemma-colim-works-over-collection}.
Nous emploierons également $H^i_{fppf}(X, \mathcal{O}) = H^i(X, \mathcal{O})$; voir la
proposition \ref{descent-proposition-same-cohomology-quasi-coherent} de Descente.

\medskip\noindent
Soient $A, f, J$ comme dans
l'exemple \ref{flat-example-one-generator} de Compléments sur la platitude,
et considérons le carré presque d'éclatement associé.
Puisque $X$, $X'$, $Z$, $E$ sont affines, la cohomologie
supérieure de $\mathcal{O}$ est nulle. Il suffit donc
de vérifier que
$$
0 \to
\mathcal{O}^{perf}(X) \to
\mathcal{O}^{perf}(X') \oplus \mathcal{O}^{perf}(Z) \to
\mathcal{O}^{perf}(E) \to 0
$$
est une suite exacte courte. Cela a été démontré dans (la démonstration du)
lemme \ref{flat-lemma-h-sheaf-colim-F} de Compléments sur la platitude.

\medskip\noindent
Soient $X, X', Z, E$ comme dans
l'exemple \ref{flat-example-two-generators} de Compléments sur la platitude.
Comme $X$ et $Z$ sont affines, nous avons
$H^i(X, \mathcal{O}_X) = H^i(Z, \mathcal{O}_Z) = 0$ pour $i > 0$.
Par le lemme \ref{flat-lemma-funny-blow-up} de Compléments sur la platitude,
nous avons $H^i(X', \mathcal{O}_{X'}) = 0$ pour $i > 0$.
Comme $E = \mathbf{P}^1_Z$ et que $Z$ est affine, nous avons aussi
$H^i(E, \mathcal{O}_E) = 0$ pour $i > 0$. Comme au paragraphe
précédent, nous nous ramenons à vérifier que
$$
0 \to
\mathcal{O}^{perf}(X) \to
\mathcal{O}^{perf}(X') \oplus \mathcal{O}^{perf}(Z) \to
\mathcal{O}^{perf}(E) \to 0
$$
est une suite exacte courte. Cela a été démontré dans (la démonstration du)
lemme \ref{flat-lemma-h-sheaf-colim-F} de Compléments sur la platitude.
\end{proof}

\begin{proposition}
\label{proposition-h-cohomology-structure-sheaf}
Soit $p$ un nombre premier. Soit $S$ un schéma quasi-compact et quasi-séparé
sur $\mathbf{F}_p$. Alors
$$
H^i((\Sch/S)_h, \mathcal{O}^h) =
\colim_F H^i(S, \mathcal{O})
$$
Dans le membre de gauche, $\mathcal{O}^h$ désigne
le faisceau h associé au faisceau structural.
\end{proposition}

\begin{proof}
C'est une reformulation du lemme \ref{lemma-h-sheaf-colim-F}.
Rappelons que
$\mathcal{O}^h = \mathcal{O}^{perf} = \colim_F \mathcal{O}$; voir le
lemme \ref{flat-lemma-char-p} de Compléments sur la platitude.
Le lemme \ref{lemma-h-sheaf-colim-F} montre que
$\mathcal{O}^{perf}$, vu comme objet de $D((\Sch/S)_{fppf})$,
est de la forme $R\epsilon_*K$ pour un certain $K \in D((\Sch/S)_h)$.
Alors $K = \epsilon^{-1}\mathcal{O}^{perf}$, qui est en fait
égal à $\mathcal{O}^{perf}$ puisque $\mathcal{O}^{perf}$ est un faisceau h. Voir la
section \ref{sites-cohomology-section-compare} de Cohomologie sur les sites.
Ainsi, $R\epsilon_*\mathcal{O}^{perf} = \mathcal{O}^{perf}$
(nous prions le lecteur d'excuser cette notation ambiguë).
Le résultat découle alors de la suite spectrale de Leray
$$
R\Gamma_h(S, \mathcal{O}^{perf}) =
R\Gamma_{fppf}(S, R\epsilon_*\mathcal{O}^{perf}) =
R\Gamma_{fppf}(S, \mathcal{O}^{perf})
$$
et du fait que
$$
H^i_{fppf}(S, \mathcal{O}^{perf}) =
H^i_{fppf}(S, \colim_F \mathcal{O}) =
\colim_F H^i_{fppf}(S, \mathcal{O})
$$
puisque $S$ est quasi-compact et quasi-séparé
(voir la démonstration du lemme \ref{lemma-h-sheaf-colim-F}).
\end{proof}












\begin{multicols}{2}[\section{Autres chapitres}]
\noindent
Préliminaires
\begin{enumerate}
\item \hyperref[introduction-section-phantom]{Introduction}
\item \hyperref[conventions-section-phantom]{Conventions}
\item \hyperref[sets-section-phantom]{Théorie des ensembles}
\item \hyperref[categories-section-phantom]{Catégories}
\item \hyperref[topology-section-phantom]{Topologie}
\item \hyperref[sheaves-section-phantom]{Faisceaux sur les espaces}
\item \hyperref[sites-section-phantom]{Sites et faisceaux}
\item \hyperref[stacks-section-phantom]{Champs}
\item \hyperref[fields-section-phantom]{Corps}
\item \hyperref[algebra-section-phantom]{Algèbre commutative}
\item \hyperref[brauer-section-phantom]{Groupes de Brauer}
\item \hyperref[homology-section-phantom]{Algèbre homologique}
\item \hyperref[derived-section-phantom]{Catégories dérivées}
\item \hyperref[simplicial-section-phantom]{Méthodes simpliciales}
\item \hyperref[more-algebra-section-phantom]{Compléments d'algèbre}
\item \hyperref[smoothing-section-phantom]{Lissage des morphismes d'anneaux}
\item \hyperref[modules-section-phantom]{Faisceaux de modules}
\item \hyperref[sites-modules-section-phantom]{Modules sur les sites}
\item \hyperref[injectives-section-phantom]{Objets injectifs}
\item \hyperref[cohomology-section-phantom]{Cohomologie des faisceaux}
\item \hyperref[sites-cohomology-section-phantom]{Cohomologie sur les sites}
\item \hyperref[dga-section-phantom]{Algèbre différentielle graduée}
\item \hyperref[dpa-section-phantom]{Algèbre à puissances divisées}
\item \hyperref[sdga-section-phantom]{Faisceaux différentiels gradués}
\item \hyperref[hypercovering-section-phantom]{Hyperrecouvrements}
\end{enumerate}
Schémas
\begin{enumerate}
\setcounter{enumi}{25}
\item \hyperref[schemes-section-phantom]{Schémas}
\item \hyperref[constructions-section-phantom]{Constructions de schémas}
\item \hyperref[properties-section-phantom]{Propriétés des schémas}
\item \hyperref[morphisms-section-phantom]{Morphismes de schémas}
\item \hyperref[coherent-section-phantom]{Cohomologie des schémas}
\item \hyperref[divisors-section-phantom]{Diviseurs}
\item \hyperref[limits-section-phantom]{Limites de schémas}
\item \hyperref[varieties-section-phantom]{Variétés}
\item \hyperref[topologies-section-phantom]{Topologies sur les schémas}
\item \hyperref[descent-section-phantom]{Descente}
\item \hyperref[perfect-section-phantom]{Catégories dérivées de schémas}
\item \hyperref[more-morphisms-section-phantom]{Compléments sur les morphismes}
\item \hyperref[flat-section-phantom]{Compléments sur la platitude}
\item \hyperref[groupoids-section-phantom]{Schémas en groupoïdes}
\item \hyperref[more-groupoids-section-phantom]{Compléments sur les schémas en groupoïdes}
\item \hyperref[etale-section-phantom]{Morphismes étales de schémas}
\end{enumerate}
Thèmes de théorie des schémas
\begin{enumerate}
\setcounter{enumi}{41}
\item \hyperref[chow-section-phantom]{Homologie de Chow}
\item \hyperref[intersection-section-phantom]{Théorie de l'intersection}
\item \hyperref[pic-section-phantom]{Schémas de Picard des courbes}
\item \hyperref[weil-section-phantom]{Théories cohomologiques de Weil}
\item \hyperref[adequate-section-phantom]{Modules adéquats}
\item \hyperref[dualizing-section-phantom]{Complexes dualisants}
\item \hyperref[duality-section-phantom]{Dualité pour les schémas}
\item \hyperref[discriminant-section-phantom]{Discriminants et différentes}
\item \hyperref[derham-section-phantom]{Cohomologie de de Rham}
\item \hyperref[local-cohomology-section-phantom]{Cohomologie locale}
\item \hyperref[algebraization-section-phantom]{Géométrie algébrique et formelle}
\item \hyperref[curves-section-phantom]{Courbes algébriques}
\item \hyperref[resolve-section-phantom]{Résolution des surfaces}
\item \hyperref[models-section-phantom]{Réduction semi-stable}
\item \hyperref[functors-section-phantom]{Foncteurs et morphismes}
\item \hyperref[equiv-section-phantom]{Catégories dérivées de variétés}
\item \hyperref[pione-section-phantom]{Groupes fondamentaux de schémas}
\item \hyperref[etale-cohomology-section-phantom]{Cohomologie étale}
\item \hyperref[crystalline-section-phantom]{Cohomologie cristalline}
\item \hyperref[proetale-section-phantom]{Cohomologie pro-étale}
\item \hyperref[relative-cycles-section-phantom]{Cycles relatifs}
\item \hyperref[more-etale-section-phantom]{Compléments de cohomologie étale}
\item \hyperref[trace-section-phantom]{La formule des traces}
\end{enumerate}
Espaces algébriques
\begin{enumerate}
\setcounter{enumi}{64}
\item \hyperref[spaces-section-phantom]{Espaces algébriques}
\item \hyperref[spaces-properties-section-phantom]{Propriétés des espaces algébriques}
\item \hyperref[spaces-morphisms-section-phantom]{Morphismes d'espaces algébriques}
\item \hyperref[decent-spaces-section-phantom]{Espaces algébriques décents}
\item \hyperref[spaces-cohomology-section-phantom]{Cohomologie des espaces algébriques}
\item \hyperref[spaces-limits-section-phantom]{Limites d'espaces algébriques}
\item \hyperref[spaces-divisors-section-phantom]{Diviseurs sur les espaces algébriques}
\item \hyperref[spaces-over-fields-section-phantom]{Espaces algébriques sur des corps}
\item \hyperref[spaces-topologies-section-phantom]{Topologies sur les espaces algébriques}
\item \hyperref[spaces-descent-section-phantom]{Descente et espaces algébriques}
\item \hyperref[spaces-perfect-section-phantom]{Catégories dérivées d'espaces}
\item \hyperref[spaces-more-morphisms-section-phantom]{Compléments sur les morphismes d'espaces}
\item \hyperref[spaces-flat-section-phantom]{Platitude sur les espaces algébriques}
\item \hyperref[spaces-groupoids-section-phantom]{Groupoïdes dans les espaces algébriques}
\item \hyperref[spaces-more-groupoids-section-phantom]{Compléments sur les groupoïdes dans les espaces}
\item \hyperref[bootstrap-section-phantom]{Amorçage}
\item \hyperref[spaces-pushouts-section-phantom]{Sommes amalgamées d'espaces algébriques}
\end{enumerate}
Thèmes de géométrie
\begin{enumerate}
\setcounter{enumi}{81}
\item \hyperref[spaces-chow-section-phantom]{Groupes de Chow des espaces}
\item \hyperref[groupoids-quotients-section-phantom]{Quotients de groupoïdes}
\item \hyperref[spaces-more-cohomology-section-phantom]{Compléments sur la cohomologie des espaces}
\item \hyperref[spaces-simplicial-section-phantom]{Espaces simpliciaux}
\item \hyperref[spaces-duality-section-phantom]{Dualité pour les espaces}
\item \hyperref[formal-spaces-section-phantom]{Espaces algébriques formels}
\item \hyperref[restricted-section-phantom]{Algébrisation des espaces formels}
\item \hyperref[spaces-resolve-section-phantom]{Résolution des surfaces revisitée}
\end{enumerate}
Théorie des déformations
\begin{enumerate}
\setcounter{enumi}{89}
\item \hyperref[formal-defos-section-phantom]{Théorie formelle des déformations}
\item \hyperref[defos-section-phantom]{Théorie des déformations}
\item \hyperref[cotangent-section-phantom]{Le complexe cotangent}
\item \hyperref[examples-defos-section-phantom]{Problèmes de déformation}
\end{enumerate}
Champs algébriques
\begin{enumerate}
\setcounter{enumi}{93}
\item \hyperref[algebraic-section-phantom]{Champs algébriques}
\item \hyperref[examples-stacks-section-phantom]{Exemples de champs}
\item \hyperref[stacks-sheaves-section-phantom]{Faisceaux sur les champs algébriques}
\item \hyperref[criteria-section-phantom]{Critères de représentabilité}
\item \hyperref[artin-section-phantom]{Axiomes d'Artin}
\item \hyperref[quot-section-phantom]{Espaces de Quot et de Hilbert}
\item \hyperref[stacks-properties-section-phantom]{Propriétés des champs algébriques}
\item \hyperref[stacks-morphisms-section-phantom]{Morphismes de champs algébriques}
\item \hyperref[stacks-limits-section-phantom]{Limites de champs algébriques}
\item \hyperref[stacks-cohomology-section-phantom]{Cohomologie des champs algébriques}
\item \hyperref[stacks-perfect-section-phantom]{Catégories dérivées de champs}
\item \hyperref[stacks-introduction-section-phantom]{Introduction aux champs algébriques}
\item \hyperref[stacks-more-morphisms-section-phantom]{Compléments sur les morphismes de champs}
\item \hyperref[stacks-geometry-section-phantom]{La géométrie des champs}
\end{enumerate}
Thèmes de théorie des espaces de modules
\begin{enumerate}
\setcounter{enumi}{107}
\item \hyperref[moduli-section-phantom]{Champs de modules}
\item \hyperref[moduli-curves-section-phantom]{Espaces de modules de courbes}
\end{enumerate}
Divers
\begin{enumerate}
\setcounter{enumi}{109}
\item \hyperref[examples-section-phantom]{Exemples}
\item \hyperref[exercises-section-phantom]{Exercices}
\item \hyperref[guide-section-phantom]{Guide bibliographique}
\item \hyperref[desirables-section-phantom]{Desiderata}
\item \hyperref[coding-section-phantom]{Style de programmation}
\item \hyperref[obsolete-section-phantom]{Obsolète}
\item \hyperref[fdl-section-phantom]{Licence de documentation libre GNU}
\item \hyperref[index-section-phantom]{Index généré automatiquement}
\end{enumerate}
\end{multicols}


\bibliography{my}
\bibliographystyle{amsalpha}

\end{document}
